% generated by GAPDoc2LaTeX from XML source (Frank Luebeck)
\documentclass[a4paper,11pt]{report}

            \usepackage{a4wide}
            \newcommand{\bbZ}{\mathbb{Z}}
        
\usepackage[top=37mm,bottom=37mm,left=27mm,right=27mm]{geometry}
\sloppy
\pagestyle{myheadings}
\usepackage{amssymb}
\usepackage[utf8]{inputenc}
\usepackage{makeidx}
\makeindex
\usepackage{color}
\definecolor{FireBrick}{rgb}{0.5812,0.0074,0.0083}
\definecolor{RoyalBlue}{rgb}{0.0236,0.0894,0.6179}
\definecolor{RoyalGreen}{rgb}{0.0236,0.6179,0.0894}
\definecolor{RoyalRed}{rgb}{0.6179,0.0236,0.0894}
\definecolor{LightBlue}{rgb}{0.8544,0.9511,1.0000}
\definecolor{Black}{rgb}{0.0,0.0,0.0}

\definecolor{linkColor}{rgb}{0.0,0.0,0.554}
\definecolor{citeColor}{rgb}{0.0,0.0,0.554}
\definecolor{fileColor}{rgb}{0.0,0.0,0.554}
\definecolor{urlColor}{rgb}{0.0,0.0,0.554}
\definecolor{promptColor}{rgb}{0.0,0.0,0.589}
\definecolor{brkpromptColor}{rgb}{0.589,0.0,0.0}
\definecolor{gapinputColor}{rgb}{0.589,0.0,0.0}
\definecolor{gapoutputColor}{rgb}{0.0,0.0,0.0}

%%  for a long time these were red and blue by default,
%%  now black, but keep variables to overwrite
\definecolor{FuncColor}{rgb}{0.0,0.0,0.0}
%% strange name because of pdflatex bug:
\definecolor{Chapter }{rgb}{0.0,0.0,0.0}
\definecolor{DarkOlive}{rgb}{0.1047,0.2412,0.0064}


\usepackage{fancyvrb}

\usepackage{mathptmx,helvet}
\usepackage[T1]{fontenc}
\usepackage{textcomp}


\usepackage[
            pdftex=true,
            bookmarks=true,        
            a4paper=true,
            pdftitle={Written with GAPDoc},
            pdfcreator={LaTeX with hyperref package / GAPDoc},
            colorlinks=true,
            backref=page,
            breaklinks=true,
            linkcolor=linkColor,
            citecolor=citeColor,
            filecolor=fileColor,
            urlcolor=urlColor,
            pdfpagemode={UseNone}, 
           ]{hyperref}

\newcommand{\maintitlesize}{\fontsize{50}{55}\selectfont}

% write page numbers to a .pnr log file for online help
\newwrite\pagenrlog
\immediate\openout\pagenrlog =\jobname.pnr
\immediate\write\pagenrlog{PAGENRS := [}
\newcommand{\logpage}[1]{\protect\write\pagenrlog{#1, \thepage,}}
%% were never documented, give conflicts with some additional packages

\newcommand{\GAP}{\textsf{GAP}}

%% nicer description environments, allows long labels
\usepackage{enumitem}
\setdescription{style=nextline}

%% depth of toc
\setcounter{tocdepth}{1}





%% command for ColorPrompt style examples
\newcommand{\gapprompt}[1]{\color{promptColor}{\bfseries #1}}
\newcommand{\gapbrkprompt}[1]{\color{brkpromptColor}{\bfseries #1}}
\newcommand{\gapinput}[1]{\color{gapinputColor}{#1}}


\begin{document}

\logpage{[ 0, 0, 0 ]}
\begin{titlepage}
\mbox{}\vfill

\begin{center}{\maintitlesize \textbf{ Digraphs \mbox{}}}\\
\vfill

\hypersetup{pdftitle= Digraphs }
\markright{\scriptsize \mbox{}\hfill  Digraphs  \hfill\mbox{}}
{\Huge \textbf{ Graphs, digraphs, and multidigraphs in \textsf{GAP} \mbox{}}}\\
\vfill

{\Huge  1.14.0 \mbox{}}\\[1cm]
{ 22 January 2026 \mbox{}}\\[1cm]
\mbox{}\\[2cm]
{\Large \textbf{ Jan De Beule\\
    \mbox{}}}\\
{\Large \textbf{ Julius Jonusas\\
   \mbox{}}}\\
{\Large \textbf{ James Mitchell\\
    \mbox{}}}\\
{\Large \textbf{ Wilf A. Wilson\\
   \mbox{}}}\\
{\Large \textbf{ Michael Young\\
    \mbox{}}}\\
{\Large \textbf{ Marina Anagnostopoulou\texttt{\symbol{45}}Merkouri\\
  \mbox{}}}\\
{\Large \textbf{ Cora Aked\\
  \mbox{}}}\\
{\Large \textbf{ Finn Buck\\
  \mbox{}}}\\
{\Large \textbf{ Stuart Burrell\\
   \mbox{}}}\\
{\Large \textbf{ Graham Campbell\\
 \mbox{}}}\\
{\Large \textbf{ Devansh Chopra\\
  \mbox{}}}\\
{\Large \textbf{ Raiyan Chowdhury\\
 \mbox{}}}\\
{\Large \textbf{ Reinis Cirpons\\
  \mbox{}}}\\
{\Large \textbf{ Ashley Clayton\\
  \mbox{}}}\\
{\Large \textbf{ Tom Conti\texttt{\symbol{45}}Leslie\\
   \mbox{}}}\\
{\Large \textbf{ Joseph Edwards\\
    \mbox{}}}\\
{\Large \textbf{ Luna Elliott\\
   \mbox{}}}\\
{\Large \textbf{ Jan Engelhardt\\
   \mbox{}}}\\
{\Large \textbf{ Fernando Flores Brito\\
  \mbox{}}}\\
{\Large \textbf{ Isuru Fernando\\
  \mbox{}}}\\
{\Large \textbf{ Ewan Gilligan\\
  \mbox{}}}\\
{\Large \textbf{ Gillis Frankie\\
  \mbox{}}}\\
{\Large \textbf{ Sebastian Gutsche\\
  \mbox{}}}\\
{\Large \textbf{ Samantha Harper\\
  \mbox{}}}\\
{\Large \textbf{ Max Horn\\
    \mbox{}}}\\
{\Large \textbf{ Harry Jack\\
  \mbox{}}}\\
{\Large \textbf{ Christopher Jefferson\\
    \mbox{}}}\\
{\Large \textbf{ Malachi Johns\\
  \mbox{}}}\\
{\Large \textbf{ Olexandr Konovalov\\
    \mbox{}}}\\
{\Large \textbf{ Hyeokjun Kwon\\
  \mbox{}}}\\
{\Large \textbf{ Aidan Lau\\
 \mbox{}}}\\
{\Large \textbf{ Andrea Lee\\
  \mbox{}}}\\
{\Large \textbf{ Saffron McIver\\
  \mbox{}}}\\
{\Large \textbf{ Seyyed Ali Mohammadiyeh\\
  \mbox{}}}\\
{\Large \textbf{ Rheya Monerasinghe\\
  \mbox{}}}\\
{\Large \textbf{ Michael Orlitzky\\
   \mbox{}}}\\
{\Large \textbf{ Matthew Pancer\\
  \mbox{}}}\\
{\Large \textbf{ Markus Pfeiffer\\
   \mbox{}}}\\
{\Large \textbf{ Amelia Phillips\\
  \mbox{}}}\\
{\Large \textbf{ Daniel Pointon\\
  \mbox{}}}\\
{\Large \textbf{ Pramoth Ragavan\\
  \mbox{}}}\\
{\Large \textbf{ Lea Racine\\
  \mbox{}}}\\
{\Large \textbf{ Christopher Russell\\
 \mbox{}}}\\
{\Large \textbf{ Artur Schaefer\\
  \mbox{}}}\\
{\Large \textbf{ Isabella Scott\\
  \mbox{}}}\\
{\Large \textbf{ Kamran Sharma\\
  \mbox{}}}\\
{\Large \textbf{ Finn Smith\\
   \mbox{}}}\\
{\Large \textbf{ Ben Spiers\\
  \mbox{}}}\\
{\Large \textbf{ Maria Tsalakou\\
   \mbox{}}}\\
{\Large \textbf{ Agastyaa Vishvanath\\
  \mbox{}}}\\
{\Large \textbf{ Meike Weiss\\
    \mbox{}}}\\
{\Large \textbf{ Murray Whyte\\
   \mbox{}}}\\
{\Large \textbf{ Fabian Zickgraf\\
  \mbox{}}}\\
\hypersetup{pdfauthor= Jan De Beule\\
    ;  Julius Jonusas\\
   ;  James Mitchell\\
    ;  Wilf A. Wilson\\
   ;  Michael Young\\
    ;  Marina Anagnostopoulou\texttt{\symbol{45}}Merkouri\\
  ;  Cora Aked\\
  ;  Finn Buck\\
  ;  Stuart Burrell\\
   ;  Graham Campbell\\
 ;  Devansh Chopra\\
  ;  Raiyan Chowdhury\\
 ;  Reinis Cirpons\\
  ;  Ashley Clayton\\
  ;  Tom Conti\texttt{\symbol{45}}Leslie\\
   ;  Joseph Edwards\\
    ;  Luna Elliott\\
   ;  Jan Engelhardt\\
   ;  Fernando Flores Brito\\
  ;  Isuru Fernando\\
  ;  Ewan Gilligan\\
  ;  Gillis Frankie\\
  ;  Sebastian Gutsche\\
  ;  Samantha Harper\\
  ;  Max Horn\\
    ;  Harry Jack\\
  ;  Christopher Jefferson\\
    ;  Malachi Johns\\
  ;  Olexandr Konovalov\\
    ;  Hyeokjun Kwon\\
  ;  Aidan Lau\\
 ;  Andrea Lee\\
  ;  Saffron McIver\\
  ;  Seyyed Ali Mohammadiyeh\\
  ;  Rheya Monerasinghe\\
  ;  Michael Orlitzky\\
   ;  Matthew Pancer\\
  ;  Markus Pfeiffer\\
   ;  Amelia Phillips\\
  ;  Daniel Pointon\\
  ;  Pramoth Ragavan\\
  ;  Lea Racine\\
  ;  Christopher Russell\\
 ;  Artur Schaefer\\
  ;  Isabella Scott\\
  ;  Kamran Sharma\\
  ;  Finn Smith\\
   ;  Ben Spiers\\
  ;  Maria Tsalakou\\
   ;  Agastyaa Vishvanath\\
  ;  Meike Weiss\\
    ;  Murray Whyte\\
   ;  Fabian Zickgraf\\
  }
\end{center}\vfill

\mbox{}\\
{\mbox{}\\
\small \noindent \textbf{ Jan De Beule\\
    }  Email: \href{mailto://jdebeule@cage.ugent.be} {\texttt{jdebeule@cage.ugent.be}}\\
  Homepage: \href{https://researchportal.vub.be/en/persons/jan-de-beule} {\texttt{https://researchportal.vub.be/en/persons/jan\texttt{\symbol{45}}de\texttt{\symbol{45}}beule}}\\
  Address: \begin{minipage}[t]{8cm}\noindent
 Vrije Universiteit Brussel, Vakgroep Wiskunde, Pleinlaan 2, B
\texttt{\symbol{45}} 1050 Brussels, Belgium\\
 \end{minipage}
}\\
{\mbox{}\\
\small \noindent \textbf{ Julius Jonusas\\
   }  Email: \href{mailto://j.jonusas@gmail.com} {\texttt{j.jonusas@gmail.com}}\\
  Homepage: \href{http://julius.jonusas.work} {\texttt{http://julius.jonusas.work}}}\\
{\mbox{}\\
\small \noindent \textbf{ James Mitchell\\
    }  Email: \href{mailto://jdm3@st-andrews.ac.uk} {\texttt{jdm3@st\texttt{\symbol{45}}andrews.ac.uk}}\\
  Homepage: \href{https://jdbm.me} {\texttt{https://jdbm.me}}\\
  Address: \begin{minipage}[t]{8cm}\noindent
 Mathematical Institute, North Haugh, St Andrews, Fife, KY16 9SS, Scotland\\
 \end{minipage}
}\\
{\mbox{}\\
\small \noindent \textbf{ Wilf A. Wilson\\
   }  Email: \href{mailto://gap@wilf-wilson.net} {\texttt{gap@wilf\texttt{\symbol{45}}wilson.net}}\\
  Homepage: \href{https://wilf.me} {\texttt{https://wilf.me}}}\\
{\mbox{}\\
\small \noindent \textbf{ Michael Young\\
    }  Email: \href{mailto://mct25@st-andrews.ac.uk} {\texttt{mct25@st\texttt{\symbol{45}}andrews.ac.uk}}\\
  Homepage: \href{https://myoung.uk/work/} {\texttt{https://myoung.uk/work/}}\\
  Address: \begin{minipage}[t]{8cm}\noindent
 Jack Cole Building, North Haugh, St Andrews, Fife, KY16 9SX, Scotland\\
 \end{minipage}
}\\
{\mbox{}\\
\small \noindent \textbf{ Marina Anagnostopoulou\texttt{\symbol{45}}Merkouri\\
  }  Email: \href{mailto://mam49@st-andrews.ac.uk} {\texttt{mam49@st\texttt{\symbol{45}}andrews.ac.uk}}}\\
{\mbox{}\\
\small \noindent \textbf{ Cora Aked\\
  }  Email: \href{mailto://ca219@st-andrews.ac.uk} {\texttt{ca219@st\texttt{\symbol{45}}andrews.ac.uk}}}\\
{\mbox{}\\
\small \noindent \textbf{ Finn Buck\\
  }  Email: \href{mailto://finneganlbuck@gmail.com} {\texttt{finneganlbuck@gmail.com}}}\\
{\mbox{}\\
\small \noindent \textbf{ Stuart Burrell\\
   }  Email: \href{mailto://stuartburrell1994@gmail.com} {\texttt{stuartburrell1994@gmail.com}}\\
  Homepage: \href{https://stuartburrell.github.io} {\texttt{https://stuartburrell.github.io}}}\\
{\mbox{}\\
\small \noindent \textbf{ Devansh Chopra\\
  }  Email: \href{mailto://dc268@st-andrews.ac.uk} {\texttt{dc268@st\texttt{\symbol{45}}andrews.ac.uk}}}\\
{\mbox{}\\
\small \noindent \textbf{ Reinis Cirpons\\
  }  Email: \href{mailto://rc234@st-andrews.ac.uk} {\texttt{rc234@st\texttt{\symbol{45}}andrews.ac.uk}}}\\
{\mbox{}\\
\small \noindent \textbf{ Ashley Clayton\\
  }  Email: \href{mailto://ac323@st-andrews.ac.uk} {\texttt{ac323@st\texttt{\symbol{45}}andrews.ac.uk}}}\\
{\mbox{}\\
\small \noindent \textbf{ Tom Conti\texttt{\symbol{45}}Leslie\\
   }  Email: \href{mailto://tom.contileslie@gmail.com} {\texttt{tom.contileslie@gmail.com}}\\
  Homepage: \href{https://tomcontileslie.com} {\texttt{https://tomcontileslie.com}}}\\
{\mbox{}\\
\small \noindent \textbf{ Joseph Edwards\\
    }  Email: \href{mailto://jde1@st-andrews.ac.uk} {\texttt{jde1@st\texttt{\symbol{45}}andrews.ac.uk}}\\
  Homepage: \href{https://github.com/Joseph-Edwards} {\texttt{https://github.com/Joseph\texttt{\symbol{45}}Edwards}}\\
  Address: \begin{minipage}[t]{8cm}\noindent
 Mathematical Institute, North Haugh, St Andrews, Fife, KY16 9SS, Scotland\\
 \end{minipage}
}\\
{\mbox{}\\
\small \noindent \textbf{ Luna Elliott\\
   }  Email: \href{mailto://luna.elliott142857@gmail.com} {\texttt{luna.elliott142857@gmail.com}}\\
  Homepage: \href{https://research.manchester.ac.uk/en/persons/luna-elliott} {\texttt{https://research.manchester.ac.uk/en/persons/luna\texttt{\symbol{45}}elliott}}}\\
{\mbox{}\\
\small \noindent \textbf{ Jan Engelhardt\\
   }  Email: \href{mailto://jengelh@inai.de} {\texttt{jengelh@inai.de}}\\
  Homepage: \href{https://inai.de} {\texttt{https://inai.de}}}\\
{\mbox{}\\
\small \noindent \textbf{ Fernando Flores Brito\\
  }  Email: \href{mailto://ffloresbrito@gmail.com} {\texttt{ffloresbrito@gmail.com}}}\\
{\mbox{}\\
\small \noindent \textbf{ Isuru Fernando\\
  }  Email: \href{mailto://isuruf@gmail.com} {\texttt{isuruf@gmail.com}}}\\
{\mbox{}\\
\small \noindent \textbf{ Ewan Gilligan\\
  }  Email: \href{mailto://eg207@st-andrews.ac.uk} {\texttt{eg207@st\texttt{\symbol{45}}andrews.ac.uk}}}\\
{\mbox{}\\
\small \noindent \textbf{ Gillis Frankie\\
  }  Email: \href{mailto://fotg1@st-andrews.ac.uk} {\texttt{fotg1@st\texttt{\symbol{45}}andrews.ac.uk}}}\\
{\mbox{}\\
\small \noindent \textbf{ Sebastian Gutsche\\
  }  Email: \href{mailto://gutsche@momo.math.rwth-aachen.de} {\texttt{gutsche@momo.math.rwth\texttt{\symbol{45}}aachen.de}}}\\
{\mbox{}\\
\small \noindent \textbf{ Samantha Harper\\
  }  Email: \href{mailto://seh25@st-andrews.ac.uk} {\texttt{seh25@st\texttt{\symbol{45}}andrews.ac.uk}}}\\
{\mbox{}\\
\small \noindent \textbf{ Max Horn\\
    }  Email: \href{mailto://mhorn@rptu.de} {\texttt{mhorn@rptu.de}}\\
  Homepage: \href{https://www.quendi.de/math} {\texttt{https://www.quendi.de/math}}\\
  Address: \begin{minipage}[t]{8cm}\noindent
 Fachbereich Mathematik, RPTU Kaiserslautern\texttt{\symbol{45}}Landau,
Gottlieb\texttt{\symbol{45}}Daimler\texttt{\symbol{45}}Stra{\ss}e 48, 67663
Kaiserslautern, Germany\\
 \end{minipage}
}\\
{\mbox{}\\
\small \noindent \textbf{ Harry Jack\\
  }  Email: \href{mailto://hrj4@st-andrews.ac.uk} {\texttt{hrj4@st\texttt{\symbol{45}}andrews.ac.uk}}}\\
{\mbox{}\\
\small \noindent \textbf{ Christopher Jefferson\\
    }  Email: \href{mailto://caj21@st-andrews.ac.uk} {\texttt{caj21@st\texttt{\symbol{45}}andrews.ac.uk}}\\
  Homepage: \href{https://heather.cafe/} {\texttt{https://heather.cafe/}}\\
  Address: \begin{minipage}[t]{8cm}\noindent
 Jack Cole Building, North Haugh, St Andrews, Fife, KY16 9SX, Scotland\\
 \end{minipage}
}\\
{\mbox{}\\
\small \noindent \textbf{ Malachi Johns\\
  }  Email: \href{mailto://zlj1@st-andrews.ac.uk} {\texttt{zlj1@st\texttt{\symbol{45}}andrews.ac.uk}}}\\
{\mbox{}\\
\small \noindent \textbf{ Olexandr Konovalov\\
    }  Email: \href{mailto://obk1@st-andrews.ac.uk} {\texttt{obk1@st\texttt{\symbol{45}}andrews.ac.uk}}\\
  Homepage: \href{https://olexandr-konovalov.github.io/} {\texttt{https://olexandr\texttt{\symbol{45}}konovalov.github.io/}}\\
  Address: \begin{minipage}[t]{8cm}\noindent
 Jack Cole Building, North Haugh, St Andrews, Fife, KY16 9SX, Scotland\\
 \end{minipage}
}\\
{\mbox{}\\
\small \noindent \textbf{ Hyeokjun Kwon\\
  }  Email: \href{mailto://hk78@st-andrews.ac.uk} {\texttt{hk78@st\texttt{\symbol{45}}andrews.ac.uk}}}\\
{\mbox{}\\
\small \noindent \textbf{ Andrea Lee\\
  }  Email: \href{mailto://ahwl1@st-andrews.ac.uk} {\texttt{ahwl1@st\texttt{\symbol{45}}andrews.ac.uk}}}\\
{\mbox{}\\
\small \noindent \textbf{ Saffron McIver\\
  }  Email: \href{mailto://sm544@st-andrews.ac.uk} {\texttt{sm544@st\texttt{\symbol{45}}andrews.ac.uk}}}\\
{\mbox{}\\
\small \noindent \textbf{ Seyyed Ali Mohammadiyeh\\
  }  Email: \href{mailto://MaxBaseCode@Gmail.Com} {\texttt{MaxBaseCode@Gmail.Com}}}\\
{\mbox{}\\
\small \noindent \textbf{ Rheya Monerasinghe\\
  }  Email: \href{mailto://rm387@st-andrews.ac.uk} {\texttt{rm387@st\texttt{\symbol{45}}andrews.ac.uk}}}\\
{\mbox{}\\
\small \noindent \textbf{ Michael Orlitzky\\
   }  Email: \href{mailto://michael@orlitzky.com} {\texttt{michael@orlitzky.com}}\\
  Homepage: \href{https://michael.orlitzky.com/} {\texttt{https://michael.orlitzky.com/}}}\\
{\mbox{}\\
\small \noindent \textbf{ Matthew Pancer\\
  }  Email: \href{mailto://mp322@st-andrews.ac.uk} {\texttt{mp322@st\texttt{\symbol{45}}andrews.ac.uk}}}\\
{\mbox{}\\
\small \noindent \textbf{ Markus Pfeiffer\\
   }  Email: \href{mailto://markus.pfeiffer@morphism.de} {\texttt{markus.pfeiffer@morphism.de}}\\
  Homepage: \href{https://markusp.morphism.de/} {\texttt{https://markusp.morphism.de/}}}\\
{\mbox{}\\
\small \noindent \textbf{ Amelia Phillips\\
  }  Email: \href{mailto://ap410@st-andrews.ac.uk} {\texttt{ap410@st\texttt{\symbol{45}}andrews.ac.uk}}}\\
{\mbox{}\\
\small \noindent \textbf{ Daniel Pointon\\
  }  Email: \href{mailto://dp211@st-andrews.ac.uk} {\texttt{dp211@st\texttt{\symbol{45}}andrews.ac.uk}}}\\
{\mbox{}\\
\small \noindent \textbf{ Pramoth Ragavan\\
  }  Email: \href{mailto://107881923+pramothragavan@users.noreply.github.com} {\texttt{107881923+pramothragavan@users.noreply.github.com}}}\\
{\mbox{}\\
\small \noindent \textbf{ Lea Racine\\
  }  Email: \href{mailto://lr217@st-andrews.ac.uk} {\texttt{lr217@st\texttt{\symbol{45}}andrews.ac.uk}}}\\
{\mbox{}\\
\small \noindent \textbf{ Artur Schaefer\\
  }  Email: \href{mailto://as305@st-and.ac.uk} {\texttt{as305@st\texttt{\symbol{45}}and.ac.uk}}}\\
{\mbox{}\\
\small \noindent \textbf{ Isabella Scott\\
  }  Email: \href{mailto://iscott@uchicago.edu} {\texttt{iscott@uchicago.edu}}}\\
{\mbox{}\\
\small \noindent \textbf{ Kamran Sharma\\
  }  Email: \href{mailto://kks4@st-andrews.ac.uk} {\texttt{kks4@st\texttt{\symbol{45}}andrews.ac.uk}}}\\
{\mbox{}\\
\small \noindent \textbf{ Finn Smith\\
   }  Email: \href{mailto://fls3@st-andrews.ac.uk} {\texttt{fls3@st\texttt{\symbol{45}}andrews.ac.uk}}\\
  Address: \begin{minipage}[t]{8cm}\noindent
 Mathematical Institute, North Haugh, St Andrews, Fife, KY16 9SS, Scotland\\
 \end{minipage}
}\\
{\mbox{}\\
\small \noindent \textbf{ Ben Spiers\\
  }  Email: \href{mailto://bspiers972@outlook.com} {\texttt{bspiers972@outlook.com}}}\\
{\mbox{}\\
\small \noindent \textbf{ Maria Tsalakou\\
   }  Email: \href{mailto://mt200@st-andrews.ac.uk} {\texttt{mt200@st\texttt{\symbol{45}}andrews.ac.uk}}\\
  Homepage: \href{https://mariatsalakou.github.io/} {\texttt{https://mariatsalakou.github.io/}}}\\
{\mbox{}\\
\small \noindent \textbf{ Agastyaa Vishvanath\\
  }  Email: \href{mailto://av215@st-andrews.ac.uk} {\texttt{av215@st\texttt{\symbol{45}}andrews.ac.uk}}}\\
{\mbox{}\\
\small \noindent \textbf{ Meike Weiss\\
    }  Email: \href{mailto://weiss@art.rwth-aachen.de} {\texttt{weiss@art.rwth\texttt{\symbol{45}}aachen.de}}\\
  Homepage: \href{https://bit.ly/4e6pUeP} {\texttt{https://bit.ly/4e6pUeP}}\\
  Address: \begin{minipage}[t]{8cm}\noindent
 Chair of Algebra and Representation Theory, Pontdriesch
10\texttt{\symbol{45}}16, 52062 Aachen\\
 \end{minipage}
}\\
{\mbox{}\\
\small \noindent \textbf{ Murray Whyte\\
   }  Email: \href{mailto://mw231@st-andrews.ac.uk} {\texttt{mw231@st\texttt{\symbol{45}}andrews.ac.uk}}\\
  Address: \begin{minipage}[t]{8cm}\noindent
 Mathematical Institute, North Haugh, St Andrews, Fife, KY16 9SS, Scotland\\
 \end{minipage}
}\\
{\mbox{}\\
\small \noindent \textbf{ Fabian Zickgraf\\
  }  Email: \href{mailto://f.zickgraf@dashdos.com} {\texttt{f.zickgraf@dashdos.com}}}\\
\end{titlepage}

\newpage\setcounter{page}{2}
{\small 
\section*{Abstract}
\logpage{[ 0, 0, 1 ]}
 The \textsf{Digraphs} package is a \textsf{GAP} package containing methods for graphs, digraphs, and multidigraphs. \mbox{}}\\[1cm]
{\small 
\section*{Copyright}
\logpage{[ 0, 0, 2 ]}
 Jan De Beule, Julius Jonu{\v s}as, James D. Mitchell, Wilf A. Wilson, Michael
Young et al.

 \textsf{Digraphs} is free software; you can redistribute it and/or modify it under the terms of
the \href{ https://www.fsf.org/licenses/gpl.html} {GNU General Public License} as published by the Free Software Foundation; either version 3 of the License,
or (at your option) any later version. \mbox{}}\\[1cm]
{\small 
\section*{Acknowledgements}
\logpage{[ 0, 0, 3 ]}
 We would like to thank Christopher Jefferson for his help in including \href{http://www.tcs.tkk.fi/Software/bliss/} {bliss}  in \textsf{Digraphs}. We also gratefully acknowledge the encouragement and assistance of Leonard
Soicher, and the inspiration of his \href{https://gap-packages.github.io/grape} {GRAPE}  package, at many points throughout the development of \textsf{Digraphs}. This package's methods for computing digraph homomorphisms are based on work
by Max Neunh{\"o}ffer, and independently Artur Sch{\"a}fer. \mbox{}}\\[1cm]
\newpage

\def\contentsname{Contents\logpage{[ 0, 0, 4 ]}}

\tableofcontents
\newpage

 
\chapter{\textcolor{Chapter }{ The \textsf{Digraphs} package }}\label{The Digraphs package}
\logpage{[ 1, 0, 0 ]}
\hyperdef{L}{X7F202ABA780E595D}{}
{
  \index{Digraphs@\textsf{Digraphs} package overview} 
\section{\textcolor{Chapter }{Introduction}}\logpage{[ 1, 1, 0 ]}
\hyperdef{L}{X7DFB63A97E67C0A1}{}
{
 This is the manual for version 1.14.0 of the \textsf{Digraphs} package. This package was developed at the University of St Andrews by: 
\begin{itemize}
\item Jan De Beule,
\item Julius Jonu{\v s}as,
\item James D. Mitchell,
\item Maria Tsalakou,
\item Wilf A. Wilson, and
\item Michael C. Young.
\end{itemize}
 Additional contributions were made by many people, for which the authors are
very grateful. A full list can be found on the title page. The \textsf{Digraphs} package contains a variety of methods for efficiently creating and storing
mutable and immutable digraphs and computing information about them. Full
explanations of all the functions contained in the package are provided below. 

 If the \href{https://gap-packages.github.io/grape} {GRAPE}  package is available, it will be loaded automatically. Digraphs created with
the \textsf{Digraphs} package can be converted to \href{https://gap-packages.github.io/grape} {GRAPE}  graphs with \texttt{Graph} (\ref{Graph}), and conversely \href{https://gap-packages.github.io/grape} {GRAPE}  graphs can be converted to \textsf{Digraphs} objects with \texttt{Digraph} (\ref{Digraph}). \href{https://gap-packages.github.io/grape} {GRAPE}  is not required for \textsf{Digraphs} to run. 

 The \href{http://www.tcs.tkk.fi/Software/bliss/} {bliss}  tool \cite{JK07} is included in this package. It is an open\texttt{\symbol{45}}source tool for
computing automorphism groups and canonical forms of graphs, written by Tommi
Junttila and Petteri Kaski. Several of the methods in the \textsf{Digraphs} package rely on \href{http://www.tcs.tkk.fi/Software/bliss/} {bliss} . If the \textsf{NautyTracesInterface} package for GAP is available then it is also possible to use \href{https://pallini.di.uniroma1.it} {nauty}  \cite{MP14} for computing automorphism groups and canonical forms in \textsf{Digraphs}. See Section \ref{Isomorphisms and canonical labellings} for more details. 

 The \href{https://github.com/graph-algorithms/edge-addition-planarity-suite} {edge-addition-planarity-suite}  is also included in \textsf{Digraphs}; see \cite{BM04}, \cite{B06}, \cite{BM06}, and \cite{B12} . The \href{https://github.com/graph-algorithms/edge-addition-planarity-suite} {edge-addition-planarity-suite}  is an open\texttt{\symbol{45}}source implementation of the edge addition
planar graph embedding algorithm and related algorithms by John M. Boyer. See
Section \ref{Planarity} for more details. 

 From version 1.0.0 of this package, digraphs can be either mutable or
immutable. Mutable digraphs can be changed in\texttt{\symbol{45}}place by many
of the methods in the package, which avoids unnecessary copying. Immutable
digraphs cannot be changed in\texttt{\symbol{45}}place, but their advantage is
that the value of an attribute of an immutable digraph is only ever computed
once. Mutable digraphs can be converted into immutable digraphs
in\texttt{\symbol{45}}place using \texttt{MakeImmutable} (\textbf{Reference: MakeImmutable}). One of the motivations for introducing mutable digraphs in version 1.0.0 was
that in practice the authors often wanted to create a digraph and immediately
modify it (removing certain edges, loops, and so on). Before version 1.0.0,
this involved copying the digraph several times, with each copy being
discarded almost immediately. From version 1.0.0, this unnecessary copying can
be eliminated by first creating a mutable digraph, then changing it
in\texttt{\symbol{45}}place, and finally converting the mutable digraph to an
immutable one in\texttt{\symbol{45}}place (if desirable). 
\subsection{\textcolor{Chapter }{Definitions}}\label{Definitions}
\logpage{[ 1, 1, 1 ]}
\hyperdef{L}{X84541F61810C741D}{}
{
 For the purposes of this package and its documentation, the following
definitions apply: 

 A \emph{digraph} $E=(E^0,E^1,r,s)$, also known as a \emph{directed graph}, consists of a set of vertices $E^0$ and a set of edges $E^1$ together with functions $s, r: E^1 \to E^0$, called the \emph{source} and \emph{range}, respectively. The source and range of an edge is respectively the values of $s, r$ at that edge. An edge is called a \emph{loop} if its source and range are the same. A digraph is called a \emph{multidigraph} if there exist two or more edges with the same source and the same range. 

 A \emph{directed walk} on a digraph is a sequence of alternating vertices and edges $(v_1, e_1, v_2, e_2, ..., e_{n-1}, v_n)$ such that each edge $e_i$ has source $v_i$ and range $v_{i+1}$. A \emph{directed path} is a directed walk where no vertex (and hence no edge) is repeated. A \emph{directed circuit} is a directed walk where $v_1 = v_n$, and a \emph{directed cycle} is a directed circuit where where no vertex is repeated, except for $v_1 = v_n$. 

 The \emph{length} of a directed walk $(v_1, e_1, v_2, e_2, ..., e_{n-1}, v_n)$ is equal to $n-1$, the number of edges it contains. A directed walk (or path) $(v_1, e_1, v_2, e_2, ..., e_{n-1}, v_n)$ is sometimes called a directed walk (or path) \emph{from vertex $v_1$ to vertex $v_n$}. A directed walk of zero length, i.e. a sequence $(v)$ for some vertex $v$, is called \emph{trivial}. A trivial directed walk is considered to be both a circuit and a cycle, as
is the empty directed walk $()$. A \emph{simple circuit} is another name for a non\texttt{\symbol{45}}trivial and
non\texttt{\symbol{45}}empty directed cycle.

 }

 }

 }

 
\chapter{\textcolor{Chapter }{Installing \textsf{Digraphs}}}\label{Installing Digraphs}
\logpage{[ 2, 0, 0 ]}
\hyperdef{L}{X817088B27F90D596}{}
{
  
\section{\textcolor{Chapter }{For those in a hurry}}\label{For those in a hurry}
\logpage{[ 2, 1, 0 ]}
\hyperdef{L}{X7DA3059C79842BF3}{}
{
  In this section we give a brief description of how to start using \textsf{Digraphs}.

 It is assumed that you have a working copy of \textsf{GAP} with version number 4.11.0 or higher. The most
up\texttt{\symbol{45}}to\texttt{\symbol{45}}date version of \textsf{GAP} and instructions on how to install it can be obtained from the main \textsf{GAP} webpage \href{https://www.gap-system.org} {\texttt{https://www.gap\texttt{\symbol{45}}system.org}}.

 The following is a summary of the steps that should lead to a successful
installation of \textsf{Digraphs}: 
\begin{itemize}
\item  ensure that the \href{https://gap-packages.github.io/io} {IO}  package version 4.5.1 or higher is available. \href{https://gap-packages.github.io/io} {IO}  must be compiled before \textsf{Digraphs} can be loaded. 
\item  ensure that the \textsf{orb} package version 4.8.2 or higher is available. \textsf{orb} has better performance when compiled, but although compilation is recommended,
it is not required to be compiled for \textsf{Digraphs} to be loaded. 
\item  ensure that the \textsf{datastructures} package version 0.2.5 or higher is available. 
\item  \textsc{This step is optional:} certain functions in \textsf{Digraphs} require the \href{https://gap-packages.github.io/grape} {GRAPE}  package to be available; see Section \ref{The Grape package} for full details. To use these functions make sure that the \href{https://gap-packages.github.io/grape} {GRAPE}  package version 4.8.1 or higher is available. If \href{https://gap-packages.github.io/grape} {GRAPE}  is not available, then \textsf{Digraphs} can be used as normal with the exception that the functions listed in
Subsection \ref{The Grape package} will not work. 
\item  \textsc{This step is optional:} certain functions in \textsf{Digraphs} require the \textsf{NautyTracesInterface} package to be available.  If you want to make use of these functions, please ensure that the \textsf{NautyTracesInterface} package version 0.2 or higher is available. If \textsf{NautyTracesInterface} is not available, then \textsf{Digraphs} can be used as normal with the exception that functions whose names contain ``Nauty'' will not work. 
\item  download the package archive \texttt{digraphs\texttt{\symbol{45}}1.14.0.tar.gz} from \href{https://digraphs.github.io/Digraphs} {the Digraphs package webpage}. 
\item  unzip and untar the file, this should create a directory called \texttt{digraphs\texttt{\symbol{45}}1.14.0}. 
\item  locate the \texttt{pkg} directory of your \textsf{GAP} directory, which contains the directories \texttt{lib}, \texttt{doc} and so on. Move the directory \texttt{digraphs\texttt{\symbol{45}}1.14.0} into the \texttt{pkg} directory. 
\item  it is necessary to compile the \textsf{Digraphs} package. Inside the \texttt{pkg/digraphs\texttt{\symbol{45}}1.14.0} directory, type 
\begin{Verbatim}[commandchars=!@|,fontsize=\small,frame=single,label=]
  ./configure
  make
\end{Verbatim}
 Further information about this step can be found in Section \ref{Compiling the kernel module}. 
\item  start \textsf{GAP} in the usual way (i.e. type \texttt{gap} at the command line). 
\item  type \texttt{LoadPackage("digraphs");} 
\end{itemize}
 If you want to check that the package is working correctly, you should run
some of the tests described in Section \ref{Testing your installation}. 
\subsection{\textcolor{Chapter }{Configuration options}}\label{Configuration options}
\logpage{[ 2, 1, 1 ]}
\hyperdef{L}{X7C1F2E6D860DBDEC}{}
{
  In addition to the usual autoconf generated configuration flags, the following
flags are provided. \begin{center}
\begin{tabular}{|l|l|}\hline
Option&
Meaning\\
\hline
\hline
\texttt{\texttt{\symbol{45}}\texttt{\symbol{45}}enable\texttt{\symbol{45}}code\texttt{\symbol{45}}coverage}&
 enable code coverage support\\
\texttt{\texttt{\symbol{45}}\texttt{\symbol{45}}enable\texttt{\symbol{45}}compile\texttt{\symbol{45}}warnings}&
 enable compiler warnings\\
\texttt{\texttt{\symbol{45}}\texttt{\symbol{45}}enable\texttt{\symbol{45}}debug}&
 enable debug mode\\
\texttt{\texttt{\symbol{45}}\texttt{\symbol{45}}with\texttt{\symbol{45}}external\texttt{\symbol{45}}bliss}&
 use external \href{http://www.tcs.tkk.fi/Software/bliss/} {bliss} \\
\texttt{\texttt{\symbol{45}}\texttt{\symbol{45}}with\texttt{\symbol{45}}external\texttt{\symbol{45}}planarity}&
 use external \href{https://github.com/graph-algorithms/edge-addition-planarity-suite} {edge-addition-planarity-suite} \\
\texttt{\texttt{\symbol{45}}\texttt{\symbol{45}}with\texttt{\symbol{45}}gaproot}&
 specify root of GAP installation\\
\texttt{\texttt{\symbol{45}}\texttt{\symbol{45}}without\texttt{\symbol{45}}intrinsics}&
 do not use compiler intrinsics even if available\\
\hline
\end{tabular}\\[2mm]
\textbf{Table: }Configuration flags\end{center}

 }

 }

   
\section{\textcolor{Chapter }{Optional package dependencies}}\label{Optional package dependencies}
\logpage{[ 2, 2, 0 ]}
\hyperdef{L}{X780AA9D97EBCA95D}{}
{
   The \textsf{Digraphs} package is written in \textsf{GAP} and C code and requires the \href{https://gap-packages.github.io/io} {IO}  package. The \href{https://gap-packages.github.io/io} {IO}  package is used to read and write transformations, partial permutations, and
bipartitions to a file. 

  
\subsection{\textcolor{Chapter }{The Grape package}}\label{The Grape package}
\logpage{[ 2, 2, 1 ]}
\hyperdef{L}{X8493C7587FCF6D8B}{}
{
  The \href{https://gap-packages.github.io/grape} {GRAPE}  package must be available for the following operations to be available: 
\begin{itemize}
\item \texttt{Graph} (\ref{Graph}) with a digraph argument
\item \texttt{AsGraph} (\ref{AsGraph}) with a digraph argument
\item \texttt{Digraph} (\ref{Digraph}) with a \href{https://gap-packages.github.io/grape} {GRAPE}  graph argument
\end{itemize}
 If \href{https://gap-packages.github.io/grape} {GRAPE}  is not available, then \textsf{Digraphs} can be used as normal with the exception that the functions above will not
work. }

 }

   
\section{\textcolor{Chapter }{Compiling the kernel module}}\label{Compiling the kernel module}
\logpage{[ 2, 3, 0 ]}
\hyperdef{L}{X849F6196875A6DF5}{}
{
  The \textsf{Digraphs} package has a \textsf{GAP} kernel component in C which should be compiled. This component contains
certain low\texttt{\symbol{45}}level functions required by \textsf{Digraphs}. 

 It is not possible to use the \textsf{Digraphs} package without compiling it.

 To compile the kernel component inside the \texttt{pkg/digraphs\texttt{\symbol{45}}1.14.0} directory, type 
\begin{Verbatim}[commandchars=!@|,fontsize=\small,frame=single,label=]
  
  ./configure
  make
\end{Verbatim}
 

 If you installed the package in another 'pkg' directory than the standard
'pkg' directory in your \textsf{GAP} installation, then you have to do two things. Firstly during compilation you
have to use the option
'\texttt{\symbol{45}}\texttt{\symbol{45}}with\texttt{\symbol{45}}gaproot=PATH'
of the 'configure' script where 'PATH' is a path to the main GAP root
directory (if not given the default '../..' is assumed).

 If you installed \textsf{GAP} on several architectures, you must execute the configure/make step for each of
the architectures. You can either do this immediately after configuring and
compiling GAP itself on this architecture, or alternatively set the
environment variable 'CONFIGNAME' to the name of the configuration you used
when compiling GAP before running './configure'. Note however that your
compiler choice and flags (environment variables 'CC' and 'CFLAGS') need to be
chosen to match the setup of the original GAP compilation. For example you
have to specify 32\texttt{\symbol{45}}bit or 64\texttt{\symbol{45}}bit mode
correctly! }

   
\section{\textcolor{Chapter }{Rebuilding the documentation}}\label{Rebuilding the documentation}
\logpage{[ 2, 4, 0 ]}
\hyperdef{L}{X857CBE5484CF703A}{}
{
  The \textsf{Digraphs} package comes complete with pdf, html, and text versions of the documentation.
However, you might find it necessary, at some point, to rebuild the
documentation. To rebuild the documentation, please use the function \texttt{DigraphsMakeDoc} (\ref{DigraphsMakeDoc}). 

\subsection{\textcolor{Chapter }{DigraphsMakeDoc}}
\logpage{[ 2, 4, 1 ]}\nobreak
\hyperdef{L}{X870631C38610AC25}{}
{\noindent\textcolor{FuncColor}{$\triangleright$\enspace\texttt{DigraphsMakeDoc({\mdseries\slshape })\index{DigraphsMakeDoc@\texttt{DigraphsMakeDoc}}
\label{DigraphsMakeDoc}
}\hfill{\scriptsize (function)}}\\
\textbf{\indent Returns:\ }
Nothing



 This function should be called with no argument to compile the \textsf{Digraphs} documentation. }

 }

   
\section{\textcolor{Chapter }{Testing your installation}}\label{Testing your installation}
\logpage{[ 2, 5, 0 ]}
\hyperdef{L}{X7862D3F37C5BBDEF}{}
{
  In this section we describe how to test that \textsf{Digraphs} is working as intended. To test that \textsf{Digraphs} is installed correctly use \texttt{DigraphsTestInstall} (\ref{DigraphsTestInstall}) or for more extensive tests use \texttt{DigraphsTestStandard} (\ref{DigraphsTestStandard}). 

 If something goes wrong, then please review the instructions in Section \ref{For those in a hurry} and ensure that \textsf{Digraphs} has been properly installed. If you continue having problems, please use the \href{https://github.com/digraphs/Digraphs/issues} {issue tracker} to report the issues you are having. 

\subsection{\textcolor{Chapter }{DigraphsTestInstall}}
\logpage{[ 2, 5, 1 ]}\nobreak
\hyperdef{L}{X86AF4DAE80B978DA}{}
{\noindent\textcolor{FuncColor}{$\triangleright$\enspace\texttt{DigraphsTestInstall({\mdseries\slshape })\index{DigraphsTestInstall@\texttt{DigraphsTestInstall}}
\label{DigraphsTestInstall}
}\hfill{\scriptsize (function)}}\\
\textbf{\indent Returns:\ }
\texttt{true} or \texttt{false}.



 This function can be called without arguments to test your installation of \textsf{Digraphs} is working correctly. These tests should take no more than a few seconds to
complete. To test more comprehensively that \textsf{Digraphs} is working correctly, use \texttt{DigraphsTestStandard} (\ref{DigraphsTestStandard}). }

 

\subsection{\textcolor{Chapter }{DigraphsTestStandard}}
\logpage{[ 2, 5, 2 ]}\nobreak
\hyperdef{L}{X7A20088B7C406C4A}{}
{\noindent\textcolor{FuncColor}{$\triangleright$\enspace\texttt{DigraphsTestStandard({\mdseries\slshape })\index{DigraphsTestStandard@\texttt{DigraphsTestStandard}}
\label{DigraphsTestStandard}
}\hfill{\scriptsize (function)}}\\
\textbf{\indent Returns:\ }
\texttt{true} or \texttt{false}.



 This function can be called without arguments to test all of the methods
included in \textsf{Digraphs}. These tests should take less than a minute to complete. 

 To quickly test that \textsf{Digraphs} is installed correctly use \texttt{DigraphsTestInstall} (\ref{DigraphsTestInstall}). For a more thorough test, use \texttt{DigraphsTestExtreme} (\ref{DigraphsTestExtreme}). }

 

\subsection{\textcolor{Chapter }{DigraphsTestExtreme}}
\logpage{[ 2, 5, 3 ]}\nobreak
\hyperdef{L}{X7885EF3E785D4298}{}
{\noindent\textcolor{FuncColor}{$\triangleright$\enspace\texttt{DigraphsTestExtreme({\mdseries\slshape })\index{DigraphsTestExtreme@\texttt{DigraphsTestExtreme}}
\label{DigraphsTestExtreme}
}\hfill{\scriptsize (function)}}\\
\textbf{\indent Returns:\ }
\texttt{true} or \texttt{false}.



 This function should be called with no argument. It executes a series of very
demanding tests, which measure the performance of a variety of functions on
large examples. These tests take a long time to complete, at least several
minutes. 

 For these tests to complete, the digraphs library \texttt{digraphs\texttt{\symbol{45}}lib} must be downloaded and placed in the \texttt{digraphs} directory in a subfolder named \texttt{digraphs\texttt{\symbol{45}}lib}. This library can be found in the \href{http://github.com/digraphs/digraphs-lib} {digraphs-lib}  repository. }

 }

   }

 
\chapter{\textcolor{Chapter }{Creating digraphs}}\label{Creating digraphs}
\logpage{[ 3, 0, 0 ]}
\hyperdef{L}{X7D34861E863A5D93}{}
{
 In this chapter we describe how to create digraphs.

 
\section{\textcolor{Chapter }{Creating digraphs}}\logpage{[ 3, 1, 0 ]}
\hyperdef{L}{X7D34861E863A5D93}{}
{
 

\subsection{\textcolor{Chapter }{IsDigraph}}
\logpage{[ 3, 1, 1 ]}\nobreak
\hyperdef{L}{X7877ADC77F85E630}{}
{\noindent\textcolor{FuncColor}{$\triangleright$\enspace\texttt{IsDigraph\index{IsDigraph@\texttt{IsDigraph}}
\label{IsDigraph}
}\hfill{\scriptsize (Category)}}\\


 Every digraph in \textsf{Digraphs} belongs to the category \texttt{IsDigraph}. Some basic attributes and operations for digraphs are \texttt{DigraphVertices} (\ref{DigraphVertices}), \texttt{DigraphEdges} (\ref{DigraphEdges}), and \texttt{OutNeighbours} (\ref{OutNeighbours}). }

 

\subsection{\textcolor{Chapter }{IsMutableDigraph}}
\logpage{[ 3, 1, 2 ]}\nobreak
\hyperdef{L}{X7D7EDF83820ED6F5}{}
{\noindent\textcolor{FuncColor}{$\triangleright$\enspace\texttt{IsMutableDigraph\index{IsMutableDigraph@\texttt{IsMutableDigraph}}
\label{IsMutableDigraph}
}\hfill{\scriptsize (Category)}}\\


 \texttt{IsMutableDigraph} is a synonym for \texttt{IsDigraph} (\ref{IsDigraph}) and \texttt{IsMutable} (\textbf{Reference: IsMutable}). A mutable digraph may be changed in\texttt{\symbol{45}}place by methods in
the \textsf{Digraphs} package, and is not attribute\texttt{\symbol{45}}storing {\textendash} see \texttt{IsAttributeStoringRep} (\textbf{Reference: IsAttributeStoringRep}). 

 A mutable digraph may be converted into an immutable
attribute\texttt{\symbol{45}}storing digraph by calling \texttt{MakeImmutable} (\textbf{Reference: MakeImmutable}) on the digraph. }

 

\subsection{\textcolor{Chapter }{IsImmutableDigraph}}
\logpage{[ 3, 1, 3 ]}\nobreak
\hyperdef{L}{X7CAFAA89804F80BD}{}
{\noindent\textcolor{FuncColor}{$\triangleright$\enspace\texttt{IsImmutableDigraph\index{IsImmutableDigraph@\texttt{IsImmutableDigraph}}
\label{IsImmutableDigraph}
}\hfill{\scriptsize (Category)}}\\


 \texttt{IsImmutableDigraph} is a subcategory of \texttt{IsDigraph} (\ref{IsDigraph}). Digraphs that lie in \texttt{IsImmutableDigraph} are immutable and attribute\texttt{\symbol{45}}storing. In particular, they
lie in \texttt{IsAttributeStoringRep} (\textbf{Reference: IsAttributeStoringRep}). 

 A mutable digraph may be converted to an immutable digraph that lies in the
category \texttt{IsImmutableDigraph} by calling \texttt{MakeImmutable} (\textbf{Reference: MakeImmutable}) on the digraph.

 The operation \texttt{DigraphMutableCopy} (\ref{DigraphMutableCopy}) can be used to construct a mutable copy of an immutable digraph. It is however
not possible to convert an immutable digraph into a mutable digraph
in\texttt{\symbol{45}}place. }

 

\subsection{\textcolor{Chapter }{IsCayleyDigraph}}
\logpage{[ 3, 1, 4 ]}\nobreak
\hyperdef{L}{X7E749324800B38A5}{}
{\noindent\textcolor{FuncColor}{$\triangleright$\enspace\texttt{IsCayleyDigraph\index{IsCayleyDigraph@\texttt{IsCayleyDigraph}}
\label{IsCayleyDigraph}
}\hfill{\scriptsize (Category)}}\\


 \texttt{IsCayleyDigraph} is a subcategory of \texttt{IsDigraph}. Digraphs that are Cayley digraphs of a group and that are constructed by the
operation \texttt{CayleyDigraph} (\ref{CayleyDigraph}) are constructed in this category, and are always immutable. }

 

\subsection{\textcolor{Chapter }{IsDigraphWithAdjacencyFunction}}
\logpage{[ 3, 1, 5 ]}\nobreak
\hyperdef{L}{X80F1B6D28478D8B9}{}
{\noindent\textcolor{FuncColor}{$\triangleright$\enspace\texttt{IsDigraphWithAdjacencyFunction\index{IsDigraphWithAdjacencyFunction@\texttt{IsDigraphWithAdjacencyFunction}}
\label{IsDigraphWithAdjacencyFunction}
}\hfill{\scriptsize (Category)}}\\


 \texttt{IsDigraphWithAdjacencyFunction} is a subcategory of \texttt{IsDigraph}. Digraphs that are \emph{created} using an adjacency function are constructed in this category. }

 

\subsection{\textcolor{Chapter }{DigraphByOutNeighboursType}}
\logpage{[ 3, 1, 6 ]}\nobreak
\hyperdef{L}{X86E798B779515678}{}
{\noindent\textcolor{FuncColor}{$\triangleright$\enspace\texttt{DigraphByOutNeighboursType\index{DigraphByOutNeighboursType@\texttt{DigraphByOutNeighboursType}}
\label{DigraphByOutNeighboursType}
}\hfill{\scriptsize (global variable)}}\\
\noindent\textcolor{FuncColor}{$\triangleright$\enspace\texttt{DigraphFamily\index{DigraphFamily@\texttt{DigraphFamily}}
\label{DigraphFamily}
}\hfill{\scriptsize (family)}}\\


 The type of all digraphs is \texttt{DigraphByOutNeighboursType}. The family of all digraphs is \texttt{DigraphFamily}. }

 

\subsection{\textcolor{Chapter }{Digraph}}
\logpage{[ 3, 1, 7 ]}\nobreak
\hyperdef{L}{X834843057CE86655}{}
{\noindent\textcolor{FuncColor}{$\triangleright$\enspace\texttt{Digraph({\mdseries\slshape [filt, ]obj[, source, range]})\index{Digraph@\texttt{Digraph}}
\label{Digraph}
}\hfill{\scriptsize (operation)}}\\
\noindent\textcolor{FuncColor}{$\triangleright$\enspace\texttt{Digraph({\mdseries\slshape [filt, ]list, func})\index{Digraph@\texttt{Digraph}!for a list and function}
\label{Digraph:for a list and function}
}\hfill{\scriptsize (operation)}}\\
\noindent\textcolor{FuncColor}{$\triangleright$\enspace\texttt{Digraph({\mdseries\slshape [filt, ]G, list, act, adj})\index{Digraph@\texttt{Digraph}!for a group, list, function, and function}
\label{Digraph:for a group, list, function, and function}
}\hfill{\scriptsize (operation)}}\\
\textbf{\indent Returns:\ }
A digraph.



 If the optional first argument \mbox{\texttt{\mdseries\slshape filt}} is present, then this should specify the category or representation the
digraph being created will belong to. For example, if \mbox{\texttt{\mdseries\slshape filt}} is \texttt{IsMutableDigraph} (\ref{IsMutableDigraph}), then the digraph being created will be mutable, if \mbox{\texttt{\mdseries\slshape filt}} is \texttt{IsImmutableDigraph} (\ref{IsImmutableDigraph}), then the digraph will be immutable. If the optional first argument \mbox{\texttt{\mdseries\slshape filt}} is not present, then \texttt{IsImmutableDigraph} (\ref{IsImmutableDigraph}) is used by default.

 
\begin{description}
\item[{for a list (i.e. an adjacency list)}]  if \mbox{\texttt{\mdseries\slshape obj}} is a list of lists of positive integers in the range from \texttt{1} to \texttt{Length(\mbox{\texttt{\mdseries\slshape obj}})}, then this function returns the digraph with vertices $E ^ 0 = $\texttt{[1 .. Length(\mbox{\texttt{\mdseries\slshape obj}})]}, and edges corresponding to the entries of \mbox{\texttt{\mdseries\slshape obj}}. 

 More precisely, there is an edge from vertex \texttt{i} to \texttt{j} if and only if \texttt{j} is in \texttt{\mbox{\texttt{\mdseries\slshape obj}}[i]}; the source of this edge is \texttt{i} and the range is \texttt{j}. If \texttt{j} occurs in \texttt{\mbox{\texttt{\mdseries\slshape obj}}[i]} with multiplicity \texttt{k}, then there are \texttt{k} edges from \texttt{i} to \texttt{j}. 
\item[{for three lists}]  if \mbox{\texttt{\mdseries\slshape obj}} is a duplicate\texttt{\symbol{45}}free list, and \mbox{\texttt{\mdseries\slshape source}} and \mbox{\texttt{\mdseries\slshape range}} are lists of equal length consisting of positive integers in the list \texttt{[1 .. Length(\mbox{\texttt{\mdseries\slshape obj}})]}, then this function returns a digraph with vertices $E ^ 0 = $\texttt{[1 .. Length(\mbox{\texttt{\mdseries\slshape obj}})]}, and \texttt{Length(\mbox{\texttt{\mdseries\slshape source}})} edges. For each \texttt{i} in \texttt{[1 .. Length(\mbox{\texttt{\mdseries\slshape source}})]} there exists an edge with source vertex \texttt{source[i]} and range vertex \texttt{range[i]}. See \texttt{DigraphSource} (\ref{DigraphSource}) and \texttt{DigraphRange} (\ref{DigraphRange}). 

 The vertices of the digraph will be labelled by the elements of \mbox{\texttt{\mdseries\slshape obj}}. 
\item[{for an integer, and two lists}]  if \mbox{\texttt{\mdseries\slshape obj}} is an integer, and \mbox{\texttt{\mdseries\slshape source}} and \mbox{\texttt{\mdseries\slshape range}} are lists of equal length consisting of positive integers in the list \texttt{[1 .. \mbox{\texttt{\mdseries\slshape obj}}]}, then this function returns a digraph with vertices $E ^ 0 = $\texttt{[1 .. \mbox{\texttt{\mdseries\slshape obj}}]}, and \texttt{Length(\mbox{\texttt{\mdseries\slshape source}})} edges. For each \texttt{i} in \texttt{[1 .. Length(\mbox{\texttt{\mdseries\slshape source}})]} there exists an edge with source vertex \texttt{source[i]} and range vertex \texttt{range[i]}. See \texttt{DigraphSource} (\ref{DigraphSource}) and \texttt{DigraphRange} (\ref{DigraphRange}). 
\item[{for a list and a function}]  if \mbox{\texttt{\mdseries\slshape list}} is a list and \mbox{\texttt{\mdseries\slshape func}} is a function taking 2 arguments that are elements of \mbox{\texttt{\mdseries\slshape list}}, and \mbox{\texttt{\mdseries\slshape func}} returns \texttt{true} or \texttt{false}, then this operation creates a digraph with vertices \texttt{[1 .. Length(\mbox{\texttt{\mdseries\slshape list}})]} and an edge from vertex \texttt{i} to vertex \texttt{j} if and only if \texttt{\mbox{\texttt{\mdseries\slshape func}}(\mbox{\texttt{\mdseries\slshape list}}[i], \mbox{\texttt{\mdseries\slshape list}}[j])} returns \texttt{true}. 
\item[{for a group, a list, and two functions}]  The arguments will be \mbox{\texttt{\mdseries\slshape G, list, act, adj}}. 

 Let \mbox{\texttt{\mdseries\slshape G}} be a group acting on the objects in \mbox{\texttt{\mdseries\slshape list}} via the action \mbox{\texttt{\mdseries\slshape act}}, and let \mbox{\texttt{\mdseries\slshape adj}} be a function taking two objects from \mbox{\texttt{\mdseries\slshape list}} as arguments and returning \texttt{true} or \texttt{false}. The function \mbox{\texttt{\mdseries\slshape adj}} will describe the adjacency between objects from \mbox{\texttt{\mdseries\slshape list}}, which is invariant under the action of \mbox{\texttt{\mdseries\slshape G}}. This variant of the constructor returns a digraph with vertices the objects
of \mbox{\texttt{\mdseries\slshape list}} and directed edges \texttt{[x, y]} when \texttt{f(x, y)} is \texttt{true}. 

 The action of the group \mbox{\texttt{\mdseries\slshape G}} on the objects in \mbox{\texttt{\mdseries\slshape list}} is stored in the attribute \texttt{DigraphGroup} (\ref{DigraphGroup}), and is used to speed up operations like \texttt{DigraphDiameter} (\ref{DigraphDiameter}). 
\item[{for a Grape package graph}]  if \mbox{\texttt{\mdseries\slshape obj}} is a \href{https://gap-packages.github.io/grape} {GRAPE}  package graph (i.e. a record for which the function \texttt{IsGraph} returns \texttt{true}), then this function returns a digraph isomorphic to \mbox{\texttt{\mdseries\slshape obj}}. 
\item[{for a binary relation}]  if \mbox{\texttt{\mdseries\slshape obj}} is a binary relation on the points \texttt{[1 .. n]} for some positive integer $n$, then this function returns the digraph defined by \mbox{\texttt{\mdseries\slshape obj}}. Specifically, this function returns a digraph which has $n$ vertices, and which has an edge with source \texttt{i} and range \texttt{j} if and only if \texttt{[i,j]} is a pair in the binary relation \mbox{\texttt{\mdseries\slshape obj}}. 
\item[{for a string naming a digraph}]  if \mbox{\texttt{\mdseries\slshape obj}} is a non\texttt{\symbol{45}}empty string, then this function returns the
digraph that has name \mbox{\texttt{\mdseries\slshape obj}}. \textsf{Digraphs} comes with a database containing a few hundred common digraph names that can
be loaded in this way. Valid names include \texttt{"folkman"}, \texttt{"diamond"} and \texttt{"brinkmann"}. If the name is commonly followed by the word \texttt{"graph"}, then it is called without writing \texttt{"graph"} at the end. You can explore the available graph names using \texttt{ListNamedDigraphs} (\ref{ListNamedDigraphs}). Digraph names are case and whitespace insensitive. 

 Note that any undirected graphs in the database are stored as symmetric
digraphs, so the resulting digraph will have twice as many edges as its
undirected counterpart. 
\end{description}
 
\begin{Verbatim}[commandchars=!@|,fontsize=\small,frame=single,label=Example]
  !gapprompt@gap>| !gapinput@gr := Digraph([|
  !gapprompt@>| !gapinput@[2, 5, 8, 10], [2, 3, 4, 2, 5, 6, 8, 9, 10], [1],|
  !gapprompt@>| !gapinput@[3, 5, 7, 8, 10], [2, 5, 7], [3, 6, 7, 9, 10], [1, 4],|
  !gapprompt@>| !gapinput@[1, 5, 9], [1, 2, 7, 8], [3, 5]]);|
  <immutable multidigraph with 10 vertices, 38 edges>
  !gapprompt@gap>| !gapinput@gr := Digraph(["a", "b", "c"], ["a"], ["b"]);|
  <immutable digraph with 3 vertices, 1 edge>
  !gapprompt@gap>| !gapinput@gr := Digraph(5, [1, 2, 2, 4, 1, 1], [2, 3, 5, 5, 1, 1]);|
  <immutable multidigraph with 5 vertices, 6 edges>
  !gapprompt@gap>| !gapinput@Petersen := Graph(SymmetricGroup(5), [[1, 2]], OnSets,|
  !gapprompt@>| !gapinput@function(x, y) return Intersection(x, y) = []; end);;|
  !gapprompt@gap>| !gapinput@Digraph(Petersen);|
  <immutable digraph with 10 vertices, 30 edges>
  !gapprompt@gap>| !gapinput@gr := Digraph([1 .. 10], ReturnTrue);|
  <immutable digraph with 10 vertices, 100 edges>
  !gapprompt@gap>| !gapinput@Digraph("Diamond");|
  <immutable digraph with 4 vertices, 10 edges>
\end{Verbatim}
 The next example illustrates the uses of the fourth and fifth variants of this
constructor. The resulting digraph is a strongly regular graph, and it is
actually the point graph of the van Lint\texttt{\symbol{45}}Schrijver partial
geometry, \cite{vLS81}. The algebraic description is taken from the seminal paper of Calderbank and
Kantor \cite{CK86}. 
\begin{Verbatim}[commandchars=!@|,fontsize=\small,frame=single,label=Example]
  !gapprompt@gap>| !gapinput@f := GF(3 ^ 4);|
  GF(3^4)
  !gapprompt@gap>| !gapinput@gamma := First(f, x -> Order(x) = 5);|
  Z(3^4)^64
  !gapprompt@gap>| !gapinput@L := Union([Zero(f)], List(Group(gamma)));|
  [ 0*Z(3), Z(3)^0, Z(3^4)^16, Z(3^4)^32, Z(3^4)^48, Z(3^4)^64 ]
  !gapprompt@gap>| !gapinput@omega := Union(List(L, x -> List(Difference(L, [x]), y -> x - y)));|
  [ Z(3)^0, Z(3), Z(3^4)^5, Z(3^4)^7, Z(3^4)^8, Z(3^4)^13, Z(3^4)^15, 
    Z(3^4)^16, Z(3^4)^21, Z(3^4)^23, Z(3^4)^24, Z(3^4)^29, Z(3^4)^31, 
    Z(3^4)^32, Z(3^4)^37, Z(3^4)^39, Z(3^4)^45, Z(3^4)^47, Z(3^4)^48, 
    Z(3^4)^53, Z(3^4)^55, Z(3^4)^56, Z(3^4)^61, Z(3^4)^63, Z(3^4)^64, 
    Z(3^4)^69, Z(3^4)^71, Z(3^4)^72, Z(3^4)^77, Z(3^4)^79 ]
  !gapprompt@gap>| !gapinput@adj := function(x, y)|
  !gapprompt@>| !gapinput@  return x - y in omega;|
  !gapprompt@>| !gapinput@end;|
  function( x, y ) ... end
  !gapprompt@gap>| !gapinput@digraph := Digraph(AsList(f), adj);|
  <immutable digraph with 81 vertices, 2430 edges>
  !gapprompt@gap>| !gapinput@group := Group(Z(3));;|
  !gapprompt@gap>| !gapinput@act := \*;|
  <Operation "*">
  !gapprompt@gap>| !gapinput@digraph := Digraph(group, List(f), act, adj);|
  <immutable digraph with 81 vertices, 2430 edges>
\end{Verbatim}
 }

 

\subsection{\textcolor{Chapter }{DigraphByAdjacencyMatrix}}
\logpage{[ 3, 1, 8 ]}\nobreak
\hyperdef{L}{X8023FE387A3AB609}{}
{\noindent\textcolor{FuncColor}{$\triangleright$\enspace\texttt{DigraphByAdjacencyMatrix({\mdseries\slshape [filt, ]list})\index{DigraphByAdjacencyMatrix@\texttt{DigraphByAdjacencyMatrix}}
\label{DigraphByAdjacencyMatrix}
}\hfill{\scriptsize (operation)}}\\
\textbf{\indent Returns:\ }
A digraph.



 If the optional first argument \mbox{\texttt{\mdseries\slshape filt}} is present, then this should specify the category or representation the
digraph being created will belong to. For example, if \mbox{\texttt{\mdseries\slshape filt}} is \texttt{IsMutableDigraph} (\ref{IsMutableDigraph}), then the digraph being created will be mutable, if \mbox{\texttt{\mdseries\slshape filt}} is \texttt{IsImmutableDigraph} (\ref{IsImmutableDigraph}), then the digraph will be immutable. If the optional first argument \mbox{\texttt{\mdseries\slshape filt}} is not present, then \texttt{IsImmutableDigraph} (\ref{IsImmutableDigraph}) is used by default.

 If \mbox{\texttt{\mdseries\slshape list}} is the adjacency matrix of a digraph in the sense of \texttt{AdjacencyMatrix} (\ref{AdjacencyMatrix}), then this operation returns the digraph which is defined by \mbox{\texttt{\mdseries\slshape list}}. 

 Alternatively, if \mbox{\texttt{\mdseries\slshape list}} is a square boolean matrix, then this operation returns the digraph with \texttt{Length(}\mbox{\texttt{\mdseries\slshape list}}\texttt{)} vertices which has the edge \texttt{[i,j]} if and only if \mbox{\texttt{\mdseries\slshape list}}\texttt{[i][j]} is \texttt{true}. 

 
\begin{Verbatim}[commandchars=!@|,fontsize=\small,frame=single,label=Example]
  !gapprompt@gap>| !gapinput@DigraphByAdjacencyMatrix([|
  !gapprompt@>| !gapinput@[0, 1, 0, 2, 0],|
  !gapprompt@>| !gapinput@[1, 1, 1, 0, 1],|
  !gapprompt@>| !gapinput@[0, 3, 2, 1, 1],|
  !gapprompt@>| !gapinput@[0, 0, 1, 0, 1],|
  !gapprompt@>| !gapinput@[2, 0, 0, 0, 0]]);|
  <immutable multidigraph with 5 vertices, 18 edges>
  !gapprompt@gap>| !gapinput@D := DigraphByAdjacencyMatrix([|
  !gapprompt@>| !gapinput@[true, false, true],|
  !gapprompt@>| !gapinput@[false, false, true],|
  !gapprompt@>| !gapinput@[false, true, false]]);|
  <immutable digraph with 3 vertices, 4 edges>
  !gapprompt@gap>| !gapinput@OutNeighbours(D);|
  [ [ 1, 3 ], [ 3 ], [ 2 ] ]
  !gapprompt@gap>| !gapinput@D := DigraphByAdjacencyMatrix(IsMutableDigraph, |
  !gapprompt@>| !gapinput@[[true, false, true],|
  !gapprompt@>| !gapinput@ [false, false, true],|
  !gapprompt@>| !gapinput@ [false, true, false]]);|
  <mutable digraph with 3 vertices, 4 edges>
\end{Verbatim}
 }

 

\subsection{\textcolor{Chapter }{DigraphByEdges}}
\logpage{[ 3, 1, 9 ]}\nobreak
\hyperdef{L}{X7F37B6768349E269}{}
{\noindent\textcolor{FuncColor}{$\triangleright$\enspace\texttt{DigraphByEdges({\mdseries\slshape [filt, ]list[, n]})\index{DigraphByEdges@\texttt{DigraphByEdges}}
\label{DigraphByEdges}
}\hfill{\scriptsize (operation)}}\\
\textbf{\indent Returns:\ }
A digraph.



 If the optional first argument \mbox{\texttt{\mdseries\slshape filt}} is present, then this should specify the category or representation the
digraph being created will belong to. For example, if \mbox{\texttt{\mdseries\slshape filt}} is \texttt{IsMutableDigraph} (\ref{IsMutableDigraph}), then the digraph being created will be mutable, if \mbox{\texttt{\mdseries\slshape filt}} is \texttt{IsImmutableDigraph} (\ref{IsImmutableDigraph}), then the digraph will be immutable. If the optional first argument \mbox{\texttt{\mdseries\slshape filt}} is not present, then \texttt{IsImmutableDigraph} (\ref{IsImmutableDigraph}) is used by default.

 If \mbox{\texttt{\mdseries\slshape list}} is list of pairs of positive integers, then this function returns the digraph
with the minimum number of vertices \texttt{m} such that its list equal \mbox{\texttt{\mdseries\slshape list}}.

 If the optional second argument \mbox{\texttt{\mdseries\slshape n}} is a positive integer with \texttt{\mbox{\texttt{\mdseries\slshape n}} {\textgreater}= m} (with \texttt{m} defined as above), then this function returns the digraph with \mbox{\texttt{\mdseries\slshape n}} vertices and list \mbox{\texttt{\mdseries\slshape list}}. 

 See \texttt{DigraphEdges} (\ref{DigraphEdges}). 
\begin{Verbatim}[commandchars=!@|,fontsize=\small,frame=single,label=Example]
  !gapprompt@gap>| !gapinput@DigraphByEdges(|
  !gapprompt@>| !gapinput@[[1, 3], [2, 1], [2, 3], [2, 5], [3, 6],|
  !gapprompt@>| !gapinput@ [4, 6], [5, 2], [5, 4], [5, 6], [6, 6]]);|
  <immutable digraph with 6 vertices, 10 edges>
  !gapprompt@gap>| !gapinput@DigraphByEdges(|
  !gapprompt@>| !gapinput@[[1, 3], [2, 1], [2, 3], [2, 5], [3, 6],|
  !gapprompt@>| !gapinput@ [4, 6], [5, 2], [5, 4], [5, 6], [6, 6]], 12);|
  <immutable digraph with 12 vertices, 10 edges>
  !gapprompt@gap>| !gapinput@DigraphByEdges(IsMutableDigraph, |
  !gapprompt@>| !gapinput@[[1, 3], [2, 1], [2, 3], [2, 5], [3, 6],|
  !gapprompt@>| !gapinput@ [4, 6], [5, 2], [5, 4], [5, 6], [6, 6]], 12);|
  <mutable digraph with 12 vertices, 10 edges>
\end{Verbatim}
 }

 

\subsection{\textcolor{Chapter }{EdgeOrbitsDigraph}}
\logpage{[ 3, 1, 10 ]}\nobreak
\hyperdef{L}{X7B75C1D680757D6F}{}
{\noindent\textcolor{FuncColor}{$\triangleright$\enspace\texttt{EdgeOrbitsDigraph({\mdseries\slshape G, edges[, n]})\index{EdgeOrbitsDigraph@\texttt{EdgeOrbitsDigraph}}
\label{EdgeOrbitsDigraph}
}\hfill{\scriptsize (operation)}}\\
\textbf{\indent Returns:\ }
 An immutable digraph. 



 If \mbox{\texttt{\mdseries\slshape G}} is a permutation group, \mbox{\texttt{\mdseries\slshape edges}} is an edge or list of edges, and \mbox{\texttt{\mdseries\slshape n}} is a non\texttt{\symbol{45}}negative integer such that \mbox{\texttt{\mdseries\slshape G}} fixes \texttt{[1 .. \mbox{\texttt{\mdseries\slshape n}}]} setwise, then this operation returns an immutable digraph with \mbox{\texttt{\mdseries\slshape n}} vertices and the union of the orbits of the edges in \mbox{\texttt{\mdseries\slshape  edges }} under the action of the permutation group \mbox{\texttt{\mdseries\slshape G}}. An edge in this context is simply a pair of positive integers. 

 If the optional third argument \mbox{\texttt{\mdseries\slshape n}} is not present, then the largest moved point of the permutation group \mbox{\texttt{\mdseries\slshape G}} is used by default.

 
\begin{Verbatim}[commandchars=!@|,fontsize=\small,frame=single,label=Example]
  !gapprompt@gap>| !gapinput@digraph := EdgeOrbitsDigraph(Group((1, 3), (1, 2)(3, 4)),|
  !gapprompt@>| !gapinput@                                [[1, 2], [4, 5]], 5);|
  <immutable digraph with 5 vertices, 12 edges>
  !gapprompt@gap>| !gapinput@OutNeighbours(digraph);|
  [ [ 2, 4, 5 ], [ 1, 3, 5 ], [ 2, 4, 5 ], [ 1, 3, 5 ], [  ] ]
  !gapprompt@gap>| !gapinput@RepresentativeOutNeighbours(digraph);|
  [ [ 2, 4, 5 ], [  ] ]
\end{Verbatim}
 }

 

\subsection{\textcolor{Chapter }{DigraphByInNeighbours}}
\logpage{[ 3, 1, 11 ]}\nobreak
\hyperdef{L}{X81BC49B57EAADEFB}{}
{\noindent\textcolor{FuncColor}{$\triangleright$\enspace\texttt{DigraphByInNeighbours({\mdseries\slshape [filt, ]list})\index{DigraphByInNeighbours@\texttt{DigraphByInNeighbours}}
\label{DigraphByInNeighbours}
}\hfill{\scriptsize (operation)}}\\
\noindent\textcolor{FuncColor}{$\triangleright$\enspace\texttt{DigraphByInNeighbors({\mdseries\slshape [filt, ]list})\index{DigraphByInNeighbors@\texttt{DigraphByInNeighbors}}
\label{DigraphByInNeighbors}
}\hfill{\scriptsize (operation)}}\\
\textbf{\indent Returns:\ }
A digraph.



 If the optional first argument \mbox{\texttt{\mdseries\slshape filt}} is present, then this should specify the category or representation the
digraph being created will belong to. For example, if \mbox{\texttt{\mdseries\slshape filt}} is \texttt{IsMutableDigraph} (\ref{IsMutableDigraph}), then the digraph being created will be mutable, if \mbox{\texttt{\mdseries\slshape filt}} is \texttt{IsImmutableDigraph} (\ref{IsImmutableDigraph}), then the digraph will be immutable. If the optional first argument \mbox{\texttt{\mdseries\slshape filt}} is not present, then \texttt{IsImmutableDigraph} (\ref{IsImmutableDigraph}) is used by default.

 If \mbox{\texttt{\mdseries\slshape list}} is a list of lists of positive integers list the range \texttt{[1 .. Length(\mbox{\texttt{\mdseries\slshape list}})]}, then this function returns the digraph with vertices $E^0=$\texttt{[1 .. Length(\mbox{\texttt{\mdseries\slshape list}})]}, and edges corresponding to the entries of \mbox{\texttt{\mdseries\slshape list}}. More precisely, there is an edge with source vertex \texttt{i} and range vertex \texttt{j} if \texttt{i} is in the list \texttt{\mbox{\texttt{\mdseries\slshape list}}[j]}. 

 If \texttt{i} occurs in the list \texttt{\mbox{\texttt{\mdseries\slshape list}}[j]} with multiplicity \texttt{k}, then there are \texttt{k} multiple edges from \texttt{i} to \texttt{j}. 

 See \texttt{InNeighbours} (\ref{InNeighbours}). 
\begin{Verbatim}[commandchars=!@|,fontsize=\small,frame=single,label=Example]
  !gapprompt@gap>| !gapinput@D := DigraphByInNeighbours([|
  !gapprompt@>| !gapinput@[2, 5, 8, 10], [2, 3, 4, 5, 6, 8, 9, 10],|
  !gapprompt@>| !gapinput@[1], [3, 5, 7, 8, 10], [2, 5, 7], [3, 6, 7, 9, 10], [1, 4],|
  !gapprompt@>| !gapinput@[1, 5, 9], [1, 2, 7, 8], [3, 5]]);|
  <immutable digraph with 10 vertices, 37 edges>
  !gapprompt@gap>| !gapinput@D := DigraphByInNeighbours([[2, 3, 2], [1], [1, 2, 3]]);|
  <immutable multidigraph with 3 vertices, 7 edges>
  !gapprompt@gap>| !gapinput@D := DigraphByInNeighbours(IsMutableDigraph, |
  !gapprompt@>| !gapinput@                              [[2, 3, 2], [1], [1, 2, 3]]);|
  <mutable multidigraph with 3 vertices, 7 edges>
\end{Verbatim}
 }

 

\subsection{\textcolor{Chapter }{CayleyDigraph}}
\logpage{[ 3, 1, 12 ]}\nobreak
\hyperdef{L}{X7FCADADC7EC28478}{}
{\noindent\textcolor{FuncColor}{$\triangleright$\enspace\texttt{CayleyDigraph({\mdseries\slshape [filt, ]G[, gens]})\index{CayleyDigraph@\texttt{CayleyDigraph}}
\label{CayleyDigraph}
}\hfill{\scriptsize (operation)}}\\
\textbf{\indent Returns:\ }
A digraph.



 Let \mbox{\texttt{\mdseries\slshape G}} be any group and let \mbox{\texttt{\mdseries\slshape gens}} be a list of elements of \mbox{\texttt{\mdseries\slshape G}}. This operation returns a digraph that corresponds to the Cayley graph of \mbox{\texttt{\mdseries\slshape G}} with respect to \mbox{\texttt{\mdseries\slshape gens}}. 

 The vertices of the digraph correspond to the elements of \mbox{\texttt{\mdseries\slshape G}}, in the order given by \texttt{Set(\mbox{\texttt{\mdseries\slshape G}})}. There exists an edge from vertex \texttt{u} to vertex \texttt{v} if and only if there exists a generator \texttt{g} in \mbox{\texttt{\mdseries\slshape gens}} such that \texttt{Set(\mbox{\texttt{\mdseries\slshape G}})[u] * g = Set(\mbox{\texttt{\mdseries\slshape G}})[v]}. 

 The labels of the vertices \texttt{u}, \texttt{v}, and the edge \texttt{[u, v]} are the corresponding elements \texttt{AsList(\mbox{\texttt{\mdseries\slshape G}})[u]}, \texttt{AsList(\mbox{\texttt{\mdseries\slshape G}})[v]}, and generator \texttt{g}, respectively; see \texttt{DigraphVertexLabel} (\ref{DigraphVertexLabel}) and \texttt{DigraphEdgeLabel} (\ref{DigraphEdgeLabel}). 

 If the optional first argument \mbox{\texttt{\mdseries\slshape filt}} is present, then this should specify the category or representation the
digraph being created will belong to. For example, if \mbox{\texttt{\mdseries\slshape filt}} is \texttt{IsMutableDigraph} (\ref{IsMutableDigraph}), then the digraph being created will be mutable, if \mbox{\texttt{\mdseries\slshape filt}} is \texttt{IsImmutableDigraph} (\ref{IsImmutableDigraph}), then the digraph will be immutable. If the optional first argument \mbox{\texttt{\mdseries\slshape filt}} is not present, then \texttt{IsImmutableDigraph} (\ref{IsImmutableDigraph}) is used by default.

 If the optional third argument \mbox{\texttt{\mdseries\slshape gens}} is not present, then the generators of \mbox{\texttt{\mdseries\slshape G}} are used by default.

 The digraph created by this operation belongs to the category \texttt{IsCayleyDigraph} (\ref{IsCayleyDigraph}), the group \mbox{\texttt{\mdseries\slshape G}} can be recovered from the digraph using \texttt{GroupOfCayleyDigraph} (\ref{GroupOfCayleyDigraph}), and the generators \mbox{\texttt{\mdseries\slshape gens}} can be obtained using \texttt{GeneratorsOfCayleyDigraph} (\ref{GeneratorsOfCayleyDigraph}).

 
\begin{Verbatim}[commandchars=!@|,fontsize=\small,frame=single,label=Example]
  !gapprompt@gap>| !gapinput@G := DihedralGroup(8);|
  <pc group of size 8 with 3 generators>
  !gapprompt@gap>| !gapinput@CayleyDigraph(G);|
  <immutable digraph with 8 vertices, 24 edges>
  !gapprompt@gap>| !gapinput@G := DihedralGroup(IsPermGroup, 8);|
  Group([ (1,2,3,4), (2,4) ])
  !gapprompt@gap>| !gapinput@CayleyDigraph(G);|
  <immutable digraph with 8 vertices, 16 edges>
  !gapprompt@gap>| !gapinput@digraph := CayleyDigraph(G, [()]);|
  <immutable digraph with 8 vertices, 8 edges>
  !gapprompt@gap>| !gapinput@GroupOfCayleyDigraph(digraph) = G;|
  true
  !gapprompt@gap>| !gapinput@GeneratorsOfCayleyDigraph(digraph);|
  [ () ]
  !gapprompt@gap>| !gapinput@digraph := CayleyDigraph(IsMutable, G, [()]);|
  <mutable digraph with 8 vertices, 8 edges>
\end{Verbatim}
 }

 

\subsection{\textcolor{Chapter }{ListNamedDigraphs}}
\logpage{[ 3, 1, 13 ]}\nobreak
\hyperdef{L}{X7BB820C9813F035F}{}
{\noindent\textcolor{FuncColor}{$\triangleright$\enspace\texttt{ListNamedDigraphs({\mdseries\slshape s[, level]})\index{ListNamedDigraphs@\texttt{ListNamedDigraphs}}
\label{ListNamedDigraphs}
}\hfill{\scriptsize (operation)}}\\
\textbf{\indent Returns:\ }
A list of strings representing digraph names.



 This function searches through the list of names that are currently in the \textsf{Digraphs} database of named digraphs. The first argument \mbox{\texttt{\mdseries\slshape s}} should be a partially completed string; this function returns all completions \texttt{str} of the string \mbox{\texttt{\mdseries\slshape s}} such that \texttt{Digraph(str)} will successfully return a digraph. 

 The optional second argument \mbox{\texttt{\mdseries\slshape level}} controls the flexibility of the search. If \texttt{\mbox{\texttt{\mdseries\slshape level}} = 1}, then only strings beginning exactly with \mbox{\texttt{\mdseries\slshape s}} are returned. If \texttt{\mbox{\texttt{\mdseries\slshape level}} = 2}, then all names containing \mbox{\texttt{\mdseries\slshape s}} as a substring are returned. If \texttt{\mbox{\texttt{\mdseries\slshape level}} = 3}, then once again a substring search is carried out, but characters that are
not alphanumeric are ignored in the search. 

 If \mbox{\texttt{\mdseries\slshape level}} is not specified, it is set by default to equal 2. 

 The search is always case and whitespace insensitive, and this is also the
case when applying \texttt{Digraph} (\ref{Digraph}) to a string. }

 }

 
\section{\textcolor{Chapter }{Changing representations}}\logpage{[ 3, 2, 0 ]}
\hyperdef{L}{X83608C407CC8836D}{}
{
 

\subsection{\textcolor{Chapter }{AsBinaryRelation}}
\logpage{[ 3, 2, 1 ]}\nobreak
\hyperdef{L}{X7FBC4BDB82FBEDD2}{}
{\noindent\textcolor{FuncColor}{$\triangleright$\enspace\texttt{AsBinaryRelation({\mdseries\slshape digraph})\index{AsBinaryRelation@\texttt{AsBinaryRelation}}
\label{AsBinaryRelation}
}\hfill{\scriptsize (operation)}}\\
\textbf{\indent Returns:\ }
A binary relation.



 If \mbox{\texttt{\mdseries\slshape digraph}} is a digraph with a positive number of vertices $n$, and no multiple edges, then this operation returns a binary relation on the
points \texttt{[1..n]}. The pair \texttt{[i,j]} is in the binary relation if and only if \texttt{[i,j]} is an edge in \mbox{\texttt{\mdseries\slshape digraph}}. 

 
\begin{Verbatim}[commandchars=!@|,fontsize=\small,frame=single,label=Example]
  !gapprompt@gap>| !gapinput@D := Digraph([[3, 2], [1, 2], [2], [3, 4]]);|
  <immutable digraph with 4 vertices, 7 edges>
  !gapprompt@gap>| !gapinput@AsBinaryRelation(D);|
  Binary Relation on 4 points
\end{Verbatim}
 }

 

\subsection{\textcolor{Chapter }{AsDigraph (for a binary relation)}}
\logpage{[ 3, 2, 2 ]}\nobreak
\hyperdef{L}{X86834E307EACC670}{}
{\noindent\textcolor{FuncColor}{$\triangleright$\enspace\texttt{AsDigraph({\mdseries\slshape [filt, ]f[, n]})\index{AsDigraph@\texttt{AsDigraph}!for a binary relation}
\label{AsDigraph:for a binary relation}
}\hfill{\scriptsize (operation)}}\\
\textbf{\indent Returns:\ }
A digraph, or \texttt{fail}.



 If \mbox{\texttt{\mdseries\slshape f}} is a binary relation represented as one of the following in \textsf{GAP}: 
\begin{description}
\item[{ a transformation }]  satisfying \texttt{IsTransformation} (\textbf{Reference: IsTransformation}); 
\item[{ a permutation }]  satisfying \texttt{IsPerm} (\textbf{Reference: IsPerm}); 
\item[{ a partial perm }]  satisfying \texttt{IsPartialPerm} (\textbf{Reference: IsPartialPerm}); 
\item[{ a binary relation }]  satisfying \texttt{IsBinaryRelation} (\textbf{Reference: IsBinaryRelation}); 
\end{description}
 and \mbox{\texttt{\mdseries\slshape n}} is a non\texttt{\symbol{45}}negative integer, then \texttt{AsDigraph} attempts to construct a digraph with \mbox{\texttt{\mdseries\slshape n}} vertices whose edges are determined by \mbox{\texttt{\mdseries\slshape f}}.

 The digraph returned by \texttt{AsDigraph} has for each vertex \texttt{v} in \texttt{[1 .. \mbox{\texttt{\mdseries\slshape n}}]}, an edge with source \texttt{v} and range \texttt{v \texttt{\symbol{94}} \mbox{\texttt{\mdseries\slshape f}}}. If \texttt{v \texttt{\symbol{94}} \mbox{\texttt{\mdseries\slshape f}}} is greater than \mbox{\texttt{\mdseries\slshape n}} for any \texttt{v}, then \texttt{fail} is returned. 

 If the optional second argument \mbox{\texttt{\mdseries\slshape n}} is not supplied, then the degree of the transformation \mbox{\texttt{\mdseries\slshape f}}, the largest moved point of the permutation \mbox{\texttt{\mdseries\slshape f}}, the maximum of the degree and the codegree of the partial perm \mbox{\texttt{\mdseries\slshape f}}, or as applicable, is used by default. 

 If the optional first argument \mbox{\texttt{\mdseries\slshape filt}} is present, then this should specify the category or representation the
digraph being created will belong to. For example, if \mbox{\texttt{\mdseries\slshape filt}} is \texttt{IsMutableDigraph} (\ref{IsMutableDigraph}), then the digraph being created will be mutable, if \mbox{\texttt{\mdseries\slshape filt}} is \texttt{IsImmutableDigraph} (\ref{IsImmutableDigraph}), then the digraph will be immutable. If the optional first argument \mbox{\texttt{\mdseries\slshape filt}} is not present, then \texttt{IsImmutableDigraph} (\ref{IsImmutableDigraph}) is used by default.

 
\begin{Verbatim}[commandchars=!@|,fontsize=\small,frame=single,label=Example]
  !gapprompt@gap>| !gapinput@f := Transformation([4, 3, 3, 1, 7, 9, 10, 4, 2, 3]);|
  Transformation( [ 4, 3, 3, 1, 7, 9, 10, 4, 2, 3 ] )
  !gapprompt@gap>| !gapinput@AsDigraph(f);|
  <immutable functional digraph with 10 vertices>
  !gapprompt@gap>| !gapinput@AsDigraph(f, 4);|
  <immutable functional digraph with 4 vertices>
  !gapprompt@gap>| !gapinput@AsDigraph(f, 5);|
  fail
  !gapprompt@gap>| !gapinput@AsDigraph((1, 2, 3, 4)) = CycleDigraph(4);|
  true
  !gapprompt@gap>| !gapinput@D := AsDigraph(IsMutableDigraph, (1, 3)(2, 4), 5);|
  <mutable digraph with 5 vertices, 5 edges>
  !gapprompt@gap>| !gapinput@DigraphEdges(D);|
  [ [ 1, 3 ], [ 2, 4 ], [ 3, 1 ], [ 4, 2 ], [ 5, 5 ] ]
  !gapprompt@gap>| !gapinput@b := BinaryRelationOnPoints(|
  !gapprompt@>| !gapinput@[[3], [1, 3, 5], [1], [1, 2, 4], [2, 3, 5]]);|
  Binary Relation on 5 points
  !gapprompt@gap>| !gapinput@D := AsDigraph(b);|
  <immutable digraph with 5 vertices, 11 edges>
\end{Verbatim}
 }

 

\subsection{\textcolor{Chapter }{Graph}}
\logpage{[ 3, 2, 3 ]}\nobreak
\hyperdef{L}{X7B335342839E5146}{}
{\noindent\textcolor{FuncColor}{$\triangleright$\enspace\texttt{Graph({\mdseries\slshape digraph})\index{Graph@\texttt{Graph}}
\label{Graph}
}\hfill{\scriptsize (operation)}}\\
\textbf{\indent Returns:\ }
A \href{https://gap-packages.github.io/grape} {GRAPE}  package graph.



 If \mbox{\texttt{\mdseries\slshape digraph}} is a mutable or immutable digraph without multiple edges, then this operation
returns a \href{https://gap-packages.github.io/grape} {GRAPE}  package graph that is isomorphic to \mbox{\texttt{\mdseries\slshape digraph}}. 

 If \mbox{\texttt{\mdseries\slshape digraph}} is a multidigraph, then since \href{https://gap-packages.github.io/grape} {GRAPE}  does not support multiple edges, the multiple edges will be reduced to a
single edge in the result. In order words, for a multidigraph this operation
will return the same as \texttt{Graph(DigraphRemoveAllMultipleEdges(}\mbox{\texttt{\mdseries\slshape digraph}}\texttt{))}. 
\begin{Verbatim}[commandchars=!@|,fontsize=\small,frame=single,label=Example]
  !gapprompt@gap>| !gapinput@Petersen := Graph(SymmetricGroup(5), [[1, 2]], OnSets,|
  !gapprompt@>| !gapinput@function(x, y) return Intersection(x, y) = []; end);;|
  !gapprompt@gap>| !gapinput@Display(Petersen);|
  rec(
    adjacencies := [ [ 3, 5, 8 ] ],
    group := 
     Group( [ ( 1, 2, 3, 5, 7)( 4, 6, 8, 9,10), ( 2, 4)( 6, 9)( 7,10) 
       ] ),
    isGraph := true,
    names := [ [ 1, 2 ], [ 2, 3 ], [ 3, 4 ], [ 1, 3 ], [ 4, 5 ], 
        [ 2, 4 ], [ 1, 5 ], [ 3, 5 ], [ 1, 4 ], [ 2, 5 ] ],
    order := 10,
    representatives := [ 1 ],
    schreierVector := [ -1, 1, 1, 2, 1, 1, 1, 1, 2, 2 ] )
  !gapprompt@gap>| !gapinput@Digraph(Petersen);|
  <immutable digraph with 10 vertices, 30 edges>
  !gapprompt@gap>| !gapinput@Graph(last) = Petersen;|
  true
\end{Verbatim}
 }

 

\subsection{\textcolor{Chapter }{AsGraph}}
\logpage{[ 3, 2, 4 ]}\nobreak
\hyperdef{L}{X7C4F13E080EC16B0}{}
{\noindent\textcolor{FuncColor}{$\triangleright$\enspace\texttt{AsGraph({\mdseries\slshape digraph})\index{AsGraph@\texttt{AsGraph}}
\label{AsGraph}
}\hfill{\scriptsize (attribute)}}\\
\textbf{\indent Returns:\ }
A \href{https://gap-packages.github.io/grape} {GRAPE}  package graph.



 If \mbox{\texttt{\mdseries\slshape digraph}} is a digraph, then this method returns the same as \texttt{Graph} (\ref{Graph}), except that if \mbox{\texttt{\mdseries\slshape digraph}} is immutable, then the result will be stored as a mutable attribute of \mbox{\texttt{\mdseries\slshape digraph}}. In this latter case, when \texttt{AsGraph(}\mbox{\texttt{\mdseries\slshape digraph}}\texttt{)} is called subsequently, the same \textsf{GAP} object will be returned as before. 
\begin{Verbatim}[commandchars=!@|,fontsize=\small,frame=single,label=Example]
  !gapprompt@gap>| !gapinput@D := Digraph([[1, 2], [3], []]);|
  <immutable digraph with 3 vertices, 3 edges>
  !gapprompt@gap>| !gapinput@G := AsGraph(D);|
  rec( adjacencies := [ [ 1, 2 ], [ 3 ], [  ] ], group := Group(()), 
    isGraph := true, names := [ 1 .. 3 ], order := 3, 
    representatives := [ 1, 2, 3 ], schreierVector := [ -1, -2, -3 ] )
\end{Verbatim}
 }

 

\subsection{\textcolor{Chapter }{AsTransformation}}
\logpage{[ 3, 2, 5 ]}\nobreak
\hyperdef{L}{X7C5360B2799943F3}{}
{\noindent\textcolor{FuncColor}{$\triangleright$\enspace\texttt{AsTransformation({\mdseries\slshape digraph})\index{AsTransformation@\texttt{AsTransformation}}
\label{AsTransformation}
}\hfill{\scriptsize (attribute)}}\\
\textbf{\indent Returns:\ }
A transformation, or \texttt{fail}



 If \mbox{\texttt{\mdseries\slshape digraph}} is a functional digraph, then \texttt{AsTransformation} returns the transformation which is defined by \mbox{\texttt{\mdseries\slshape digraph}}. See \texttt{IsFunctionalDigraph} (\ref{IsFunctionalDigraph}). Otherwise, \texttt{AsTransformation(}\mbox{\texttt{\mdseries\slshape digraph}}\texttt{)} returns \texttt{fail}. 

 If \mbox{\texttt{\mdseries\slshape digraph}} is a functional digraph with $n$ vertices, then \texttt{AsTransformation(}\mbox{\texttt{\mdseries\slshape digraph}}\texttt{)} will return the transformation \texttt{f} of degree at most $n$ where for each $1 \leq i \leq n$, \texttt{i \texttt{\symbol{94}} f} is equal to the unique out\texttt{\symbol{45}}neighbour of vertex \texttt{i} in \mbox{\texttt{\mdseries\slshape digraph}}. 
\begin{Verbatim}[commandchars=!@|,fontsize=\small,frame=single,label=Example]
  !gapprompt@gap>| !gapinput@D := Digraph([[1], [3], [2]]);|
  <immutable digraph with 3 vertices, 3 edges>
  !gapprompt@gap>| !gapinput@AsTransformation(D);|
  Transformation( [ 1, 3, 2 ] )
  !gapprompt@gap>| !gapinput@D := CycleDigraph(IsMutableDigraph, 3);|
  <mutable digraph with 3 vertices, 3 edges>
  !gapprompt@gap>| !gapinput@AsTransformation(D);|
  Transformation( [ 2, 3, 1 ] )
  !gapprompt@gap>| !gapinput@AsPermutation(last);|
  (1,2,3)
  !gapprompt@gap>| !gapinput@D := Digraph([[2, 3], [], []]);|
  <immutable digraph with 3 vertices, 2 edges>
  !gapprompt@gap>| !gapinput@AsTransformation(D);|
  fail
\end{Verbatim}
 }

 }

 
\section{\textcolor{Chapter }{New digraphs from old}}\logpage{[ 3, 3, 0 ]}
\hyperdef{L}{X85126078848B420A}{}
{
 

\subsection{\textcolor{Chapter }{DigraphImmutableCopy}}
\logpage{[ 3, 3, 1 ]}\nobreak
\hyperdef{L}{X83D93A8A8251E6F9}{}
{\noindent\textcolor{FuncColor}{$\triangleright$\enspace\texttt{DigraphImmutableCopy({\mdseries\slshape digraph})\index{DigraphImmutableCopy@\texttt{DigraphImmutableCopy}}
\label{DigraphImmutableCopy}
}\hfill{\scriptsize (operation)}}\\
\noindent\textcolor{FuncColor}{$\triangleright$\enspace\texttt{DigraphMutableCopy({\mdseries\slshape digraph})\index{DigraphMutableCopy@\texttt{DigraphMutableCopy}}
\label{DigraphMutableCopy}
}\hfill{\scriptsize (operation)}}\\
\noindent\textcolor{FuncColor}{$\triangleright$\enspace\texttt{DigraphCopySameMutability({\mdseries\slshape digraph})\index{DigraphCopySameMutability@\texttt{DigraphCopySameMutability}}
\label{DigraphCopySameMutability}
}\hfill{\scriptsize (operation)}}\\
\noindent\textcolor{FuncColor}{$\triangleright$\enspace\texttt{DigraphCopy({\mdseries\slshape digraph})\index{DigraphCopy@\texttt{DigraphCopy}}
\label{DigraphCopy}
}\hfill{\scriptsize (operation)}}\\
\textbf{\indent Returns:\ }
A digraph.



 Each of these operations returns a new copy of \mbox{\texttt{\mdseries\slshape digraph}}, of the appropriate mutability, retaining none of the attributes or
properties of \mbox{\texttt{\mdseries\slshape digraph}}. 

 \texttt{DigraphCopy} is a synonym for \texttt{DigraphCopySameMutability}. 
\begin{Verbatim}[commandchars=!@|,fontsize=\small,frame=single,label=Example]
  !gapprompt@gap>| !gapinput@D := CycleDigraph(10);|
  <immutable cycle digraph with 10 vertices>
  !gapprompt@gap>| !gapinput@DigraphCopy(D) = D;|
  true
  !gapprompt@gap>| !gapinput@IsIdenticalObj(DigraphCopy(D), D);|
  false
  !gapprompt@gap>| !gapinput@DigraphMutableCopy(D);|
  <mutable digraph with 10 vertices, 10 edges>
\end{Verbatim}
 }

 

\subsection{\textcolor{Chapter }{DigraphImmutableCopyIfImmutable}}
\logpage{[ 3, 3, 2 ]}\nobreak
\hyperdef{L}{X8399F7427B227228}{}
{\noindent\textcolor{FuncColor}{$\triangleright$\enspace\texttt{DigraphImmutableCopyIfImmutable({\mdseries\slshape digraph})\index{DigraphImmutableCopyIfImmutable@\texttt{DigraphImmutableCopyIfImmutable}}
\label{DigraphImmutableCopyIfImmutable}
}\hfill{\scriptsize (operation)}}\\
\noindent\textcolor{FuncColor}{$\triangleright$\enspace\texttt{DigraphImmutableCopyIfMutable({\mdseries\slshape digraph})\index{DigraphImmutableCopyIfMutable@\texttt{DigraphImmutableCopyIfMutable}}
\label{DigraphImmutableCopyIfMutable}
}\hfill{\scriptsize (operation)}}\\
\noindent\textcolor{FuncColor}{$\triangleright$\enspace\texttt{DigraphMutableCopyIfMutable({\mdseries\slshape digraph})\index{DigraphMutableCopyIfMutable@\texttt{DigraphMutableCopyIfMutable}}
\label{DigraphMutableCopyIfMutable}
}\hfill{\scriptsize (operation)}}\\
\noindent\textcolor{FuncColor}{$\triangleright$\enspace\texttt{DigraphMutableCopyIfImmutable({\mdseries\slshape digraph})\index{DigraphMutableCopyIfImmutable@\texttt{DigraphMutableCopyIfImmutable}}
\label{DigraphMutableCopyIfImmutable}
}\hfill{\scriptsize (operation)}}\\
\textbf{\indent Returns:\ }
A digraph.



 Each of these operations returns either the original argument \mbox{\texttt{\mdseries\slshape digraph}}, or a new copy of \mbox{\texttt{\mdseries\slshape digraph}} of the appropriate mutability, according to the mutability of \mbox{\texttt{\mdseries\slshape digraph}}. 
\begin{Verbatim}[commandchars=!@|,fontsize=\small,frame=single,label=Example]
  !gapprompt@gap>| !gapinput@C := CycleDigraph(10);|
  <immutable cycle digraph with 10 vertices>
  !gapprompt@gap>| !gapinput@D := DigraphImmutableCopyIfImmutable(C);|
  <immutable digraph with 10 vertices, 10 edges>
  !gapprompt@gap>| !gapinput@IsIdenticalObj(C, D);|
  false
  !gapprompt@gap>| !gapinput@C = D;|
  true
  !gapprompt@gap>| !gapinput@D := DigraphImmutableCopyIfMutable(C);|
  <immutable cycle digraph with 10 vertices>
  !gapprompt@gap>| !gapinput@IsIdenticalObj(C, D);|
  true
  !gapprompt@gap>| !gapinput@C = D;|
  true
  !gapprompt@gap>| !gapinput@D := DigraphMutableCopyIfMutable(C);|
  <immutable cycle digraph with 10 vertices>
  !gapprompt@gap>| !gapinput@IsMutableDigraph(D);|
  false
  !gapprompt@gap>| !gapinput@D := DigraphMutableCopyIfImmutable(C);|
  <mutable digraph with 10 vertices, 10 edges>
  !gapprompt@gap>| !gapinput@IsMutableDigraph(D);|
  true
  !gapprompt@gap>| !gapinput@C := CycleDigraph(IsMutableDigraph, 10);|
  <mutable digraph with 10 vertices, 10 edges>
  !gapprompt@gap>| !gapinput@D := DigraphImmutableCopyIfImmutable(C);|
  <mutable digraph with 10 vertices, 10 edges>
  !gapprompt@gap>| !gapinput@IsIdenticalObj(C, D);|
  true
  !gapprompt@gap>| !gapinput@C = D;|
  true
  !gapprompt@gap>| !gapinput@D := DigraphImmutableCopyIfMutable(C);|
  <immutable digraph with 10 vertices, 10 edges>
  !gapprompt@gap>| !gapinput@IsIdenticalObj(C, D);|
  false
  !gapprompt@gap>| !gapinput@C = D;|
  true
  !gapprompt@gap>| !gapinput@D := DigraphMutableCopyIfMutable(C);|
  <mutable digraph with 10 vertices, 10 edges>
  !gapprompt@gap>| !gapinput@IsMutableDigraph(D);|
  true
  !gapprompt@gap>| !gapinput@D := DigraphMutableCopyIfImmutable(C);|
  <mutable digraph with 10 vertices, 10 edges>
  !gapprompt@gap>| !gapinput@IsIdenticalObj(C, D);|
  true
  !gapprompt@gap>| !gapinput@IsMutableDigraph(D);|
  true
\end{Verbatim}
 }

 

\subsection{\textcolor{Chapter }{InducedSubdigraph}}
\logpage{[ 3, 3, 3 ]}\nobreak
\hyperdef{L}{X83C51DA182CCEA2F}{}
{\noindent\textcolor{FuncColor}{$\triangleright$\enspace\texttt{InducedSubdigraph({\mdseries\slshape digraph, verts})\index{InducedSubdigraph@\texttt{InducedSubdigraph}}
\label{InducedSubdigraph}
}\hfill{\scriptsize (operation)}}\\
\textbf{\indent Returns:\ }
A digraph.



 If \mbox{\texttt{\mdseries\slshape digraph}} is a digraph, and \mbox{\texttt{\mdseries\slshape verts}} is a subset of the vertices of \mbox{\texttt{\mdseries\slshape digraph}}, then this operation returns a digraph constructed from \mbox{\texttt{\mdseries\slshape digraph}} by retaining precisely those vertices in \mbox{\texttt{\mdseries\slshape verts}}, and those edges whose source and range vertices are both contained in \mbox{\texttt{\mdseries\slshape verts}}. 

 The vertices of the induced subdigraph are \texttt{[1..Length(verts)]} but the original vertex labels can be accessed via \texttt{DigraphVertexLabels} (\ref{DigraphVertexLabels}). 

 If \mbox{\texttt{\mdseries\slshape digraph}} belongs to \texttt{IsMutableDigraph} (\ref{IsMutableDigraph}), then \mbox{\texttt{\mdseries\slshape digraph}} is modified in place. If \mbox{\texttt{\mdseries\slshape digraph}} belongs to \texttt{IsImmutableDigraph} (\ref{IsImmutableDigraph}), a new immutable digraph containing the appropriate vertices and edges is
returned.

 
\begin{Verbatim}[commandchars=!@|,fontsize=\small,frame=single,label=Example]
  !gapprompt@gap>| !gapinput@D := Digraph([[1, 1, 2, 3, 4, 4], [1, 3, 4], [3, 1], [1, 1]]);|
  <immutable multidigraph with 4 vertices, 13 edges>
  !gapprompt@gap>| !gapinput@InducedSubdigraph(D, [1, 3, 4]);|
  <immutable multidigraph with 3 vertices, 9 edges>
  !gapprompt@gap>| !gapinput@DigraphVertices(last);|
  [ 1 .. 3 ]
  !gapprompt@gap>| !gapinput@D := DigraphMutableCopy(D);|
  <mutable multidigraph with 4 vertices, 13 edges>
  !gapprompt@gap>| !gapinput@new := InducedSubdigraph(D, [1, 3, 4]);|
  <mutable multidigraph with 3 vertices, 9 edges>
  !gapprompt@gap>| !gapinput@D = new;|
  true
\end{Verbatim}
 }

 

\subsection{\textcolor{Chapter }{ReducedDigraph}}
\logpage{[ 3, 3, 4 ]}\nobreak
\hyperdef{L}{X7CAF093B85A93D2F}{}
{\noindent\textcolor{FuncColor}{$\triangleright$\enspace\texttt{ReducedDigraph({\mdseries\slshape digraph})\index{ReducedDigraph@\texttt{ReducedDigraph}}
\label{ReducedDigraph}
}\hfill{\scriptsize (operation)}}\\
\noindent\textcolor{FuncColor}{$\triangleright$\enspace\texttt{ReducedDigraphAttr({\mdseries\slshape digraph})\index{ReducedDigraphAttr@\texttt{ReducedDigraphAttr}}
\label{ReducedDigraphAttr}
}\hfill{\scriptsize (attribute)}}\\
\textbf{\indent Returns:\ }
A digraph.



 This function returns a digraph isomorphic to the subdigraph of \mbox{\texttt{\mdseries\slshape digraph}} induced by the set of non\texttt{\symbol{45}}isolated vertices, i.e. the set
of those vertices of \mbox{\texttt{\mdseries\slshape digraph}} which are the source or range of some edge in \mbox{\texttt{\mdseries\slshape digraph}}. See \texttt{InducedSubdigraph} (\ref{InducedSubdigraph}). 

 The ordering of the remaining vertices of \mbox{\texttt{\mdseries\slshape digraph}} is preserved, as are the labels of the remaining vertices and edges; see \texttt{DigraphVertexLabels} (\ref{DigraphVertexLabels}) and \texttt{DigraphEdgeLabels} (\ref{DigraphEdgeLabels}). This can allow one to match a vertex in the reduced digraph to the
corresponding vertex in \mbox{\texttt{\mdseries\slshape digraph}}. 

 If \mbox{\texttt{\mdseries\slshape digraph}} is immutable, then a new immutable digraph is returned. Otherwise, the
isolated vertices of the mutable digraph \mbox{\texttt{\mdseries\slshape digraph}} are removed in\texttt{\symbol{45}}place. 
\begin{Verbatim}[commandchars=!@|,fontsize=\small,frame=single,label=Example]
  !gapprompt@gap>| !gapinput@D := Digraph([[1, 2], [], [], [1, 4], []]);|
  <immutable digraph with 5 vertices, 4 edges>
  !gapprompt@gap>| !gapinput@R := ReducedDigraph(D);|
  <immutable digraph with 3 vertices, 4 edges>
  !gapprompt@gap>| !gapinput@OutNeighbours(R);|
  [ [ 1, 2 ], [  ], [ 1, 3 ] ]
  !gapprompt@gap>| !gapinput@DigraphEdges(D);|
  [ [ 1, 1 ], [ 1, 2 ], [ 4, 1 ], [ 4, 4 ] ]
  !gapprompt@gap>| !gapinput@DigraphEdges(R);|
  [ [ 1, 1 ], [ 1, 2 ], [ 3, 1 ], [ 3, 3 ] ]
  !gapprompt@gap>| !gapinput@DigraphVertexLabel(R, 3);|
  4
  !gapprompt@gap>| !gapinput@DigraphVertexLabel(R, 2);|
  2
  !gapprompt@gap>| !gapinput@D := Digraph(IsMutableDigraph, [[], [3], [2]]);|
  <mutable digraph with 3 vertices, 2 edges>
  !gapprompt@gap>| !gapinput@ReducedDigraph(D);|
  <mutable digraph with 2 vertices, 2 edges>
  !gapprompt@gap>| !gapinput@D;|
  <mutable digraph with 2 vertices, 2 edges>
\end{Verbatim}
 }

 

\subsection{\textcolor{Chapter }{MaximalSymmetricSubdigraph}}
\logpage{[ 3, 3, 5 ]}\nobreak
\hyperdef{L}{X829E3EAC7C4B3B1E}{}
{\noindent\textcolor{FuncColor}{$\triangleright$\enspace\texttt{MaximalSymmetricSubdigraph({\mdseries\slshape digraph})\index{MaximalSymmetricSubdigraph@\texttt{MaximalSymmetricSubdigraph}}
\label{MaximalSymmetricSubdigraph}
}\hfill{\scriptsize (operation)}}\\
\noindent\textcolor{FuncColor}{$\triangleright$\enspace\texttt{MaximalSymmetricSubdigraphAttr({\mdseries\slshape digraph})\index{MaximalSymmetricSubdigraphAttr@\texttt{MaximalSymmetricSubdigraphAttr}}
\label{MaximalSymmetricSubdigraphAttr}
}\hfill{\scriptsize (attribute)}}\\
\noindent\textcolor{FuncColor}{$\triangleright$\enspace\texttt{MaximalSymmetricSubdigraphWithoutLoops({\mdseries\slshape digraph})\index{MaximalSymmetricSubdigraphWithoutLoops@\texttt{Maximal}\-\texttt{Symmetric}\-\texttt{Subdigraph}\-\texttt{Without}\-\texttt{Loops}}
\label{MaximalSymmetricSubdigraphWithoutLoops}
}\hfill{\scriptsize (operation)}}\\
\noindent\textcolor{FuncColor}{$\triangleright$\enspace\texttt{MaximalSymmetricSubdigraphWithoutLoopsAttr({\mdseries\slshape digraph})\index{MaximalSymmetricSubdigraphWithoutLoopsAttr@\texttt{Maximal}\-\texttt{Symmetric}\-\texttt{Subdigraph}\-\texttt{Without}\-\texttt{Loops}\-\texttt{Attr}}
\label{MaximalSymmetricSubdigraphWithoutLoopsAttr}
}\hfill{\scriptsize (attribute)}}\\
\textbf{\indent Returns:\ }
A digraph.



 If \mbox{\texttt{\mdseries\slshape digraph}} is a digraph, then \texttt{MaximalSymmetricSubdigraph} returns a symmetric digraph without multiple edges which has the same vertex
set as \mbox{\texttt{\mdseries\slshape digraph}}, and whose edge list is formed from \mbox{\texttt{\mdseries\slshape digraph}} by ignoring the multiplicity of edges, and by ignoring edges \texttt{[u,v]} for which there does not exist an edge \texttt{[v,u]}. 

 The digraph returned by \texttt{MaximalSymmetricSubdigraphWithoutLoops} is the same, except that loops are removed.

 If \mbox{\texttt{\mdseries\slshape digraph}} is immutable, then a new immutable digraph is returned. Otherwise, the mutable
digraph \mbox{\texttt{\mdseries\slshape digraph}} is changed in\texttt{\symbol{45}}place into such a digraph described above.

 See \texttt{IsSymmetricDigraph} (\ref{IsSymmetricDigraph}), \texttt{IsMultiDigraph} (\ref{IsMultiDigraph}), and \texttt{DigraphHasLoops} (\ref{DigraphHasLoops}) for more information. 
\begin{Verbatim}[commandchars=!@|,fontsize=\small,frame=single,label=Example]
  !gapprompt@gap>| !gapinput@D := Digraph([[2, 2], [1, 3], [4], [3, 1]]);|
  <immutable multidigraph with 4 vertices, 7 edges>
  !gapprompt@gap>| !gapinput@not IsSymmetricDigraph(D) and IsMultiDigraph(D);|
  true
  !gapprompt@gap>| !gapinput@OutNeighbours(D);|
  [ [ 2, 2 ], [ 1, 3 ], [ 4 ], [ 3, 1 ] ]
  !gapprompt@gap>| !gapinput@S := MaximalSymmetricSubdigraph(D);|
  <immutable symmetric digraph with 4 vertices, 4 edges>
  !gapprompt@gap>| !gapinput@IsSymmetricDigraph(S) and not IsMultiDigraph(S);|
  true
  !gapprompt@gap>| !gapinput@OutNeighbours(S);|
  [ [ 2 ], [ 1 ], [ 4 ], [ 3 ] ]
  !gapprompt@gap>| !gapinput@D := CycleDigraph(IsMutableDigraph, 3);|
  <mutable digraph with 3 vertices, 3 edges>
  !gapprompt@gap>| !gapinput@MaximalSymmetricSubdigraph(D);|
  <mutable empty digraph with 3 vertices>
  !gapprompt@gap>| !gapinput@D;|
  <mutable empty digraph with 3 vertices>
\end{Verbatim}
 }

 

\subsection{\textcolor{Chapter }{MaximalAntiSymmetricSubdigraph}}
\logpage{[ 3, 3, 6 ]}\nobreak
\hyperdef{L}{X79BA6A66846D5A95}{}
{\noindent\textcolor{FuncColor}{$\triangleright$\enspace\texttt{MaximalAntiSymmetricSubdigraph({\mdseries\slshape digraph})\index{MaximalAntiSymmetricSubdigraph@\texttt{MaximalAntiSymmetricSubdigraph}}
\label{MaximalAntiSymmetricSubdigraph}
}\hfill{\scriptsize (operation)}}\\
\noindent\textcolor{FuncColor}{$\triangleright$\enspace\texttt{MaximalAntiSymmetricSubdigraphAttr({\mdseries\slshape digraph})\index{MaximalAntiSymmetricSubdigraphAttr@\texttt{MaximalAntiSymmetricSubdigraphAttr}}
\label{MaximalAntiSymmetricSubdigraphAttr}
}\hfill{\scriptsize (attribute)}}\\
\textbf{\indent Returns:\ }
A digraph.



 If \mbox{\texttt{\mdseries\slshape digraph}} is a digraph, then \texttt{MaximalAntiSymmetricSubdigraph} returns an anti\texttt{\symbol{45}}symmetric subdigraph of \mbox{\texttt{\mdseries\slshape digraph}} formed by retaining the vertices of \mbox{\texttt{\mdseries\slshape digraph}}, discarding any duplicate edges, and discarding any edge \texttt{[i,j]} of \mbox{\texttt{\mdseries\slshape digraph}} where \texttt{i {\textgreater} j} and the reverse edge \texttt{[j,i]} is an edge of \mbox{\texttt{\mdseries\slshape digraph}}. In other words, for every symmetric pair of edges \texttt{[i,j]} and \texttt{[j,i]} in \mbox{\texttt{\mdseries\slshape digraph}}, where \texttt{i} and \texttt{j} are distinct, it discards the edge $[\max(i,j),\min(i,j)]$. 

 If \mbox{\texttt{\mdseries\slshape digraph}} is immutable, then a new immutable digraph is returned. Otherwise, the mutable
digraph \mbox{\texttt{\mdseries\slshape digraph}} is changed in\texttt{\symbol{45}}place. 

 See \texttt{IsAntisymmetricDigraph} (\ref{IsAntisymmetricDigraph}) for more information. 
\begin{Verbatim}[commandchars=!@|,fontsize=\small,frame=single,label=Example]
  !gapprompt@gap>| !gapinput@D := Digraph([[2, 2], [1, 3], [4], [3, 1]]);|
  <immutable multidigraph with 4 vertices, 7 edges>
  !gapprompt@gap>| !gapinput@not IsAntiSymmetricDigraph(D) and IsMultiDigraph(D);|
  true
  !gapprompt@gap>| !gapinput@OutNeighbours(D);|
  [ [ 2, 2 ], [ 1, 3 ], [ 4 ], [ 3, 1 ] ]
  !gapprompt@gap>| !gapinput@D := MaximalAntiSymmetricSubdigraph(D);|
  <immutable antisymmetric digraph with 4 vertices, 4 edges>
  !gapprompt@gap>| !gapinput@IsAntiSymmetricDigraph(D) and not IsMultiDigraph(D);|
  true
  !gapprompt@gap>| !gapinput@OutNeighbours(D);|
  [ [ 2 ], [ 3 ], [ 4 ], [ 1 ] ]
  !gapprompt@gap>| !gapinput@D := Digraph(IsMutableDigraph, [[2], [1]]);|
  <mutable digraph with 2 vertices, 2 edges>
  !gapprompt@gap>| !gapinput@MaximalAntiSymmetricSubdigraph(D);|
  <mutable digraph with 2 vertices, 1 edge>
  !gapprompt@gap>| !gapinput@D;|
  <mutable digraph with 2 vertices, 1 edge>
\end{Verbatim}
 }

 

\subsection{\textcolor{Chapter }{UndirectedSpanningForest}}
\logpage{[ 3, 3, 7 ]}\nobreak
\hyperdef{L}{X7DD9766C86D3ED20}{}
{\noindent\textcolor{FuncColor}{$\triangleright$\enspace\texttt{UndirectedSpanningForest({\mdseries\slshape digraph})\index{UndirectedSpanningForest@\texttt{UndirectedSpanningForest}}
\label{UndirectedSpanningForest}
}\hfill{\scriptsize (operation)}}\\
\noindent\textcolor{FuncColor}{$\triangleright$\enspace\texttt{UndirectedSpanningForestAttr({\mdseries\slshape digraph})\index{UndirectedSpanningForestAttr@\texttt{UndirectedSpanningForestAttr}}
\label{UndirectedSpanningForestAttr}
}\hfill{\scriptsize (attribute)}}\\
\noindent\textcolor{FuncColor}{$\triangleright$\enspace\texttt{UndirectedSpanningTree({\mdseries\slshape digraph})\index{UndirectedSpanningTree@\texttt{UndirectedSpanningTree}}
\label{UndirectedSpanningTree}
}\hfill{\scriptsize (operation)}}\\
\noindent\textcolor{FuncColor}{$\triangleright$\enspace\texttt{UndirectedSpanningTreeAttr({\mdseries\slshape digraph})\index{UndirectedSpanningTreeAttr@\texttt{UndirectedSpanningTreeAttr}}
\label{UndirectedSpanningTreeAttr}
}\hfill{\scriptsize (attribute)}}\\
\textbf{\indent Returns:\ }
A digraph, or \texttt{fail}.



 If \mbox{\texttt{\mdseries\slshape digraph}} is a digraph with at least one vertex, then \texttt{UndirectedSpanningForest} returns an undirected spanning forest of \mbox{\texttt{\mdseries\slshape digraph}}, otherwise this attribute returns \texttt{fail}. See \texttt{IsUndirectedSpanningForest} (\ref{IsUndirectedSpanningForest}) for the definition of an undirected spanning forest.

 If \mbox{\texttt{\mdseries\slshape digraph}} is a digraph with at least one vertex and whose \texttt{MaximalSymmetricSubdigraph} (\ref{MaximalSymmetricSubdigraph}) is connected (see \texttt{IsConnectedDigraph} (\ref{IsConnectedDigraph})), then \texttt{UndirectedSpanningTree} returns an undirected spanning tree of \mbox{\texttt{\mdseries\slshape digraph}}, otherwise this attribute returns \texttt{fail}. See \texttt{IsUndirectedSpanningTree} (\ref{IsUndirectedSpanningTree}) for the definition of an undirected spanning tree.

 If \mbox{\texttt{\mdseries\slshape digraph}} is immutable, then an immutable digraph is returned. Otherwise, the mutable
digraph \mbox{\texttt{\mdseries\slshape digraph}} is changed in\texttt{\symbol{45}}place into an undirected spanning tree of \mbox{\texttt{\mdseries\slshape digraph}}.

 Note that for an immutable digraph that has known undirected spanning tree,
the attribute \texttt{UndirectedSpanningTree} returns the same digraph as the attribute \texttt{UndirectedSpanningForest}. 
\begin{Verbatim}[commandchars=!@|,fontsize=\small,frame=single,label=Example]
  !gapprompt@gap>| !gapinput@D := Digraph([[1, 2, 1, 3], [1], [4], [3, 4, 3]]);|
  <immutable multidigraph with 4 vertices, 9 edges>
  !gapprompt@gap>| !gapinput@UndirectedSpanningTree(D);|
  fail
  !gapprompt@gap>| !gapinput@forest := UndirectedSpanningForest(D);|
  <immutable undirected forest with 4 vertices>
  !gapprompt@gap>| !gapinput@OutNeighbours(forest);|
  [ [ 2 ], [ 1 ], [ 4 ], [ 3 ] ]
  !gapprompt@gap>| !gapinput@IsUndirectedSpanningForest(D, forest);|
  true
  !gapprompt@gap>| !gapinput@DigraphConnectedComponents(forest).comps;|
  [ [ 1, 2 ], [ 3, 4 ] ]
  !gapprompt@gap>| !gapinput@DigraphConnectedComponents(MaximalSymmetricSubdigraph(D)).comps;|
  [ [ 1, 2 ], [ 3, 4 ] ]
  !gapprompt@gap>| !gapinput@UndirectedSpanningForest(MaximalSymmetricSubdigraph(D))|
  !gapprompt@>| !gapinput@= forest;|
  true
  !gapprompt@gap>| !gapinput@D := CompleteDigraph(4);|
  <immutable complete digraph with 4 vertices>
  !gapprompt@gap>| !gapinput@tree := UndirectedSpanningTree(D);|
  <immutable undirected tree with 4 vertices>
  !gapprompt@gap>| !gapinput@IsUndirectedSpanningTree(D, tree);|
  true
  !gapprompt@gap>| !gapinput@tree = UndirectedSpanningForest(D);|
  true
  !gapprompt@gap>| !gapinput@UndirectedSpanningForest(EmptyDigraph(0));|
  fail
  !gapprompt@gap>| !gapinput@D := PetersenGraph(IsMutableDigraph);|
  <mutable digraph with 10 vertices, 30 edges>
  !gapprompt@gap>| !gapinput@UndirectedSpanningTree(D);|
  <mutable digraph with 10 vertices, 18 edges>
  !gapprompt@gap>| !gapinput@D;|
  <mutable digraph with 10 vertices, 18 edges>
\end{Verbatim}
 }

 

\subsection{\textcolor{Chapter }{DigraphShortestPathSpanningTree}}
\logpage{[ 3, 3, 8 ]}\nobreak
\hyperdef{L}{X79C770918610AD97}{}
{\noindent\textcolor{FuncColor}{$\triangleright$\enspace\texttt{DigraphShortestPathSpanningTree({\mdseries\slshape digraph, v})\index{DigraphShortestPathSpanningTree@\texttt{DigraphShortestPathSpanningTree}}
\label{DigraphShortestPathSpanningTree}
}\hfill{\scriptsize (operation)}}\\
\textbf{\indent Returns:\ }
A digraph, or \texttt{fail}.



 If \mbox{\texttt{\mdseries\slshape v}} is a vertex in \mbox{\texttt{\mdseries\slshape digraph}} and every other vertex of \mbox{\texttt{\mdseries\slshape digraph}} is reachable from \mbox{\texttt{\mdseries\slshape v}}, then this operation returns the shortest path spanning tree of \mbox{\texttt{\mdseries\slshape digraph}} rooted at \mbox{\texttt{\mdseries\slshape v}}. If there exist vertices in \mbox{\texttt{\mdseries\slshape digraph}} (other than \mbox{\texttt{\mdseries\slshape v}}) that are not reachable from \mbox{\texttt{\mdseries\slshape v}}, then \texttt{DigraphShortestPathSpanningTree} returns \texttt{fail}. See \texttt{IsReachable} (\ref{IsReachable}). 

 The \emph{shortest path spanning tree of \mbox{\texttt{\mdseries\slshape digraph}} rooted at \mbox{\texttt{\mdseries\slshape v}}} is a subdigraph of \mbox{\texttt{\mdseries\slshape digraph}} that is a directed tree, with unique source vertex \mbox{\texttt{\mdseries\slshape v}}, and where for each other vertex \mbox{\texttt{\mdseries\slshape u}} in \mbox{\texttt{\mdseries\slshape digraph}}, the unique directed path from \mbox{\texttt{\mdseries\slshape v}} to \mbox{\texttt{\mdseries\slshape u}} in the tree is the lexicographically\texttt{\symbol{45}}least shortest
directed path from \mbox{\texttt{\mdseries\slshape v}} to \mbox{\texttt{\mdseries\slshape u}} in \mbox{\texttt{\mdseries\slshape digraph}}. 

 See \texttt{IsDirectedTree} (\ref{IsDirectedTree}), \texttt{DigraphSources} (\ref{DigraphSources}), and \texttt{DigraphShortestPath} (\ref{DigraphShortestPath}). 

 If \mbox{\texttt{\mdseries\slshape digraph}} belongs to \texttt{IsMutableDigraph} (\ref{IsMutableDigraph}), then \mbox{\texttt{\mdseries\slshape digraph}} is modified in place. If \mbox{\texttt{\mdseries\slshape digraph}} belongs to \texttt{IsImmutableDigraph} (\ref{IsImmutableDigraph}), then the spanning tree is a new immutable digraph. 
\begin{Verbatim}[commandchars=!@|,fontsize=\small,frame=single,label=Example]
  !gapprompt@gap>| !gapinput@D := Digraph([[1, 2], [3], [2, 4], [1], [2, 4]]);|
  <immutable digraph with 5 vertices, 8 edges>
  !gapprompt@gap>| !gapinput@ForAll([2 .. 5], v -> IsReachable(D, 1, v));|
  false
  !gapprompt@gap>| !gapinput@DigraphShortestPathSpanningTree(D, 1);|
  fail
  !gapprompt@gap>| !gapinput@tree := DigraphShortestPathSpanningTree(D, 5);|
  <immutable directed tree with 5 vertices>
  !gapprompt@gap>| !gapinput@OutNeighbours(tree);|
  [ [  ], [ 3 ], [  ], [ 1 ], [ 2, 4 ] ]
  !gapprompt@gap>| !gapinput@ForAll(DigraphVertices(D), v ->|
  !gapprompt@>| !gapinput@DigraphShortestPath(D, 5, v) = DigraphPath(tree, 5, v));|
  true
\end{Verbatim}
 }

 

\subsection{\textcolor{Chapter }{QuotientDigraph}}
\logpage{[ 3, 3, 9 ]}\nobreak
\hyperdef{L}{X7D1D26D27F5B56C2}{}
{\noindent\textcolor{FuncColor}{$\triangleright$\enspace\texttt{QuotientDigraph({\mdseries\slshape digraph, p})\index{QuotientDigraph@\texttt{QuotientDigraph}}
\label{QuotientDigraph}
}\hfill{\scriptsize (operation)}}\\
\textbf{\indent Returns:\ }
A digraph.



 If \mbox{\texttt{\mdseries\slshape digraph}} is a digraph, and \mbox{\texttt{\mdseries\slshape p}} is a partition of the vertices of \mbox{\texttt{\mdseries\slshape digraph}}, then this operation returns a digraph constructed by amalgamating all
vertices of \mbox{\texttt{\mdseries\slshape digraph}} which lie in the same part of \mbox{\texttt{\mdseries\slshape p}}. 

 A partition of the vertices of \mbox{\texttt{\mdseries\slshape digraph}} is a list of non\texttt{\symbol{45}}empty disjoint lists, such that the union
of all the sub\texttt{\symbol{45}}lists is equal to vertex set of \mbox{\texttt{\mdseries\slshape digraph}}. In particular, each vertex must appear in precisely one
sub\texttt{\symbol{45}}list. 

 The vertices of \mbox{\texttt{\mdseries\slshape digraph}} in part \texttt{i} of \mbox{\texttt{\mdseries\slshape p}} will become vertex \texttt{i} in the quotient, and there exists some edge in \mbox{\texttt{\mdseries\slshape digraph}} with source in part \texttt{i} and range in part \texttt{j} if and only if there is an edge from \texttt{i} to \texttt{j} in the quotient. In particular, this means that the quotient of a digraph has
no multiple edges. which was a change introduced in version 1.0.0 of the \textsf{Digraphs} package. 

 If \mbox{\texttt{\mdseries\slshape digraph}} belongs to \texttt{IsMutableDigraph} (\ref{IsMutableDigraph}), then \mbox{\texttt{\mdseries\slshape digraph}} is modified in place. If \mbox{\texttt{\mdseries\slshape digraph}} belongs to \texttt{IsImmutableDigraph} (\ref{IsImmutableDigraph}), a new immutable digraph with the above properties is returned.

 
\begin{Verbatim}[commandchars=!@|,fontsize=\small,frame=single,label=Example]
  !gapprompt@gap>| !gapinput@D := Digraph([[2, 1], [4], [1], [1, 3, 4]]);|
  <immutable digraph with 4 vertices, 7 edges>
  !gapprompt@gap>| !gapinput@DigraphVertices(D);|
  [ 1 .. 4 ]
  !gapprompt@gap>| !gapinput@DigraphEdges(D);|
  [ [ 1, 2 ], [ 1, 1 ], [ 2, 4 ], [ 3, 1 ], [ 4, 1 ], [ 4, 3 ], 
    [ 4, 4 ] ]
  !gapprompt@gap>| !gapinput@p := [[1], [2, 4], [3]];|
  [ [ 1 ], [ 2, 4 ], [ 3 ] ]
  !gapprompt@gap>| !gapinput@quo := QuotientDigraph(D, p);|
  <immutable digraph with 3 vertices, 6 edges>
  !gapprompt@gap>| !gapinput@DigraphVertices(quo);|
  [ 1 .. 3 ]
  !gapprompt@gap>| !gapinput@DigraphEdges(quo);|
  [ [ 1, 1 ], [ 1, 2 ], [ 2, 1 ], [ 2, 2 ], [ 2, 3 ], [ 3, 1 ] ]
  !gapprompt@gap>| !gapinput@QuotientDigraph(EmptyDigraph(0), []);|
  <immutable empty digraph with 0 vertices>
\end{Verbatim}
 }

 

\subsection{\textcolor{Chapter }{DigraphReverse}}
\logpage{[ 3, 3, 10 ]}\nobreak
\hyperdef{L}{X78DECD26811EFD7C}{}
{\noindent\textcolor{FuncColor}{$\triangleright$\enspace\texttt{DigraphReverse({\mdseries\slshape digraph})\index{DigraphReverse@\texttt{DigraphReverse}}
\label{DigraphReverse}
}\hfill{\scriptsize (operation)}}\\
\noindent\textcolor{FuncColor}{$\triangleright$\enspace\texttt{DigraphReverseAttr({\mdseries\slshape digraph})\index{DigraphReverseAttr@\texttt{DigraphReverseAttr}}
\label{DigraphReverseAttr}
}\hfill{\scriptsize (attribute)}}\\
\textbf{\indent Returns:\ }
A digraph.



 The reverse of a digraph is the digraph formed by reversing the orientation of
each of its edges, i.e. for every edge \texttt{[i, j]} of a digraph, the reverse contains the corresponding edge \texttt{[j, i]}.

 \texttt{DigraphReverse} returns the reverse of the digraph \mbox{\texttt{\mdseries\slshape digraph}}. If \mbox{\texttt{\mdseries\slshape digraph}} is immutable, then a new immutable digraph is returned. Otherwise, the mutable
digraph \mbox{\texttt{\mdseries\slshape digraph}} is changed in\texttt{\symbol{45}}place into its reverse. 
\begin{Verbatim}[commandchars=!@|,fontsize=\small,frame=single,label=Example]
  !gapprompt@gap>| !gapinput@D := Digraph([[3], [1, 3, 5], [1], [1, 2, 4], [2, 3, 5]]);|
  <immutable digraph with 5 vertices, 11 edges>
  !gapprompt@gap>| !gapinput@DigraphReverse(D);|
  <immutable digraph with 5 vertices, 11 edges>
  !gapprompt@gap>| !gapinput@OutNeighbours(last);|
  [ [ 2, 3, 4 ], [ 4, 5 ], [ 1, 2, 5 ], [ 4 ], [ 2, 5 ] ]
  !gapprompt@gap>| !gapinput@D := Digraph([[2, 4], [1], [4], [3, 4]]);|
  <immutable digraph with 4 vertices, 6 edges>
  !gapprompt@gap>| !gapinput@DigraphEdges(D);|
  [ [ 1, 2 ], [ 1, 4 ], [ 2, 1 ], [ 3, 4 ], [ 4, 3 ], [ 4, 4 ] ]
  !gapprompt@gap>| !gapinput@DigraphEdges(DigraphReverse(D));|
  [ [ 1, 2 ], [ 2, 1 ], [ 3, 4 ], [ 4, 1 ], [ 4, 3 ], [ 4, 4 ] ]
  !gapprompt@gap>| !gapinput@D := CycleDigraph(IsMutableDigraph, 3);|
  <mutable digraph with 3 vertices, 3 edges>
  !gapprompt@gap>| !gapinput@OutNeighbours(D);|
  [ [ 2 ], [ 3 ], [ 1 ] ]
  !gapprompt@gap>| !gapinput@DigraphReverse(D);|
  <mutable digraph with 3 vertices, 3 edges>
  !gapprompt@gap>| !gapinput@OutNeighbours(D);|
  [ [ 3 ], [ 1 ], [ 2 ] ]
\end{Verbatim}
 }

 

\subsection{\textcolor{Chapter }{DigraphDual}}
\logpage{[ 3, 3, 11 ]}\nobreak
\hyperdef{L}{X7F71D99D852B130F}{}
{\noindent\textcolor{FuncColor}{$\triangleright$\enspace\texttt{DigraphDual({\mdseries\slshape digraph})\index{DigraphDual@\texttt{DigraphDual}}
\label{DigraphDual}
}\hfill{\scriptsize (operation)}}\\
\noindent\textcolor{FuncColor}{$\triangleright$\enspace\texttt{DigraphDualAttr({\mdseries\slshape digraph})\index{DigraphDualAttr@\texttt{DigraphDualAttr}}
\label{DigraphDualAttr}
}\hfill{\scriptsize (attribute)}}\\
\textbf{\indent Returns:\ }
A digraph.



 The \emph{dual} of \mbox{\texttt{\mdseries\slshape digraph}} has the same vertices as \mbox{\texttt{\mdseries\slshape digraph}}, and there is an edge in the dual from \texttt{i} to \texttt{j} whenever there is no edge from \texttt{i} to \texttt{j} in \mbox{\texttt{\mdseries\slshape digraph}}. The \emph{dual} is sometimes called the \emph{complement}.

 \texttt{DigraphDual} returns the dual of the digraph \mbox{\texttt{\mdseries\slshape digraph}}. If \mbox{\texttt{\mdseries\slshape digraph}} is an immutable digraph, then a new immutable digraph is returned. Otherwise,
the mutable digraph \mbox{\texttt{\mdseries\slshape digraph}} is changed in\texttt{\symbol{45}}place into its dual. 
\begin{Verbatim}[commandchars=!@|,fontsize=\small,frame=single,label=Example]
  !gapprompt@gap>| !gapinput@D := Digraph([[2, 3], [], [4, 6], [5], [],|
  !gapprompt@>| !gapinput@[7, 8, 9], [], [], []]);|
  <immutable digraph with 9 vertices, 8 edges>
  !gapprompt@gap>| !gapinput@DigraphDual(D);|
  <immutable digraph with 9 vertices, 73 edges>
  !gapprompt@gap>| !gapinput@D := CycleDigraph(IsMutableDigraph, 3);|
  <mutable digraph with 3 vertices, 3 edges>
  !gapprompt@gap>| !gapinput@DigraphDual(D);|
  <mutable digraph with 3 vertices, 6 edges>
  !gapprompt@gap>| !gapinput@D;|
  <mutable digraph with 3 vertices, 6 edges>
\end{Verbatim}
 }

 

\subsection{\textcolor{Chapter }{DigraphSymmetricClosure}}
\logpage{[ 3, 3, 12 ]}\nobreak
\hyperdef{L}{X874883DD7DD450C4}{}
{\noindent\textcolor{FuncColor}{$\triangleright$\enspace\texttt{DigraphSymmetricClosure({\mdseries\slshape digraph})\index{DigraphSymmetricClosure@\texttt{DigraphSymmetricClosure}}
\label{DigraphSymmetricClosure}
}\hfill{\scriptsize (operation)}}\\
\noindent\textcolor{FuncColor}{$\triangleright$\enspace\texttt{DigraphSymmetricClosureAttr({\mdseries\slshape digraph})\index{DigraphSymmetricClosureAttr@\texttt{DigraphSymmetricClosureAttr}}
\label{DigraphSymmetricClosureAttr}
}\hfill{\scriptsize (attribute)}}\\
\textbf{\indent Returns:\ }
A digraph.



 If \mbox{\texttt{\mdseries\slshape digraph}} is a digraph, then this attribute gives the minimal symmetric digraph which
has the same vertices and contains all the edges of \mbox{\texttt{\mdseries\slshape digraph}}.

 A digraph is \emph{symmetric} if its adjacency matrix \texttt{AdjacencyMatrix} (\ref{AdjacencyMatrix}) is symmetric. For a digraph with multiple edges this means that there are the
same number of edges from a vertex \texttt{u} to a vertex \texttt{v} as there are from \texttt{v} to \texttt{u}; see \texttt{IsSymmetricDigraph} (\ref{IsSymmetricDigraph}).

 If \mbox{\texttt{\mdseries\slshape digraph}} is immutable, then a new immutable digraph is returned. Otherwise, the mutable
digraph \mbox{\texttt{\mdseries\slshape digraph}} is changed in\texttt{\symbol{45}}place into its symmetric closure. 
\begin{Verbatim}[commandchars=!@|,fontsize=\small,frame=single,label=Example]
  !gapprompt@gap>| !gapinput@D := Digraph([[1, 2, 3], [2, 4], [1], [3, 4]]);|
  <immutable digraph with 4 vertices, 8 edges>
  !gapprompt@gap>| !gapinput@D := DigraphSymmetricClosure(D);|
  <immutable symmetric digraph with 4 vertices, 11 edges>
  !gapprompt@gap>| !gapinput@IsSymmetricDigraph(D);|
  true
  !gapprompt@gap>| !gapinput@List(OutNeighbours(D), AsSet);|
  [ [ 1, 2, 3 ], [ 1, 2, 4 ], [ 1, 4 ], [ 2, 3, 4 ] ]
  !gapprompt@gap>| !gapinput@D := Digraph([[2, 2], [1]]);|
  <immutable multidigraph with 2 vertices, 3 edges>
  !gapprompt@gap>| !gapinput@D := DigraphSymmetricClosure(D);|
  <immutable symmetric multidigraph with 2 vertices, 4 edges>
  !gapprompt@gap>| !gapinput@OutNeighbours(D);|
  [ [ 2, 2 ], [ 1, 1 ] ]
  !gapprompt@gap>| !gapinput@D := CycleDigraph(IsMutableDigraph, 3);|
  <mutable digraph with 3 vertices, 3 edges>
  !gapprompt@gap>| !gapinput@DigraphSymmetricClosure(D);|
  <mutable digraph with 3 vertices, 6 edges>
  !gapprompt@gap>| !gapinput@D;|
  <mutable digraph with 3 vertices, 6 edges>
\end{Verbatim}
 }

 

\subsection{\textcolor{Chapter }{DigraphTransitiveClosure}}
\logpage{[ 3, 3, 13 ]}\nobreak
\hyperdef{L}{X7A6C419080AD41DE}{}
{\noindent\textcolor{FuncColor}{$\triangleright$\enspace\texttt{DigraphTransitiveClosure({\mdseries\slshape digraph})\index{DigraphTransitiveClosure@\texttt{DigraphTransitiveClosure}}
\label{DigraphTransitiveClosure}
}\hfill{\scriptsize (operation)}}\\
\noindent\textcolor{FuncColor}{$\triangleright$\enspace\texttt{DigraphTransitiveClosureAttr({\mdseries\slshape digraph})\index{DigraphTransitiveClosureAttr@\texttt{DigraphTransitiveClosureAttr}}
\label{DigraphTransitiveClosureAttr}
}\hfill{\scriptsize (attribute)}}\\
\noindent\textcolor{FuncColor}{$\triangleright$\enspace\texttt{DigraphReflexiveTransitiveClosure({\mdseries\slshape digraph})\index{DigraphReflexiveTransitiveClosure@\texttt{DigraphReflexiveTransitiveClosure}}
\label{DigraphReflexiveTransitiveClosure}
}\hfill{\scriptsize (operation)}}\\
\noindent\textcolor{FuncColor}{$\triangleright$\enspace\texttt{DigraphReflexiveTransitiveClosureAttr({\mdseries\slshape digraph})\index{DigraphReflexiveTransitiveClosureAttr@\texttt{Digraph}\-\texttt{Reflexive}\-\texttt{Transitive}\-\texttt{Closure}\-\texttt{Attr}}
\label{DigraphReflexiveTransitiveClosureAttr}
}\hfill{\scriptsize (attribute)}}\\
\textbf{\indent Returns:\ }
A digraph.



 If \mbox{\texttt{\mdseries\slshape digraph}} is a digraph with no multiple edges, then these attributes return the
(reflexive) transitive closure of \mbox{\texttt{\mdseries\slshape digraph}}.

 A digraph is \emph{reflexive} if it has a loop at every vertex, and it is \emph{transitive} if whenever \texttt{[i,j]} and \texttt{[j,k]} are edges of \mbox{\texttt{\mdseries\slshape digraph}}, \texttt{[i,k]} is also an edge. The \emph{(reflexive) transitive closure} of a digraph \mbox{\texttt{\mdseries\slshape digraph}} is the least (reflexive and) transitive digraph containing \mbox{\texttt{\mdseries\slshape digraph}}.

 If \mbox{\texttt{\mdseries\slshape digraph}} is immutable, then a new immutable digraph is returned. Otherwise, the mutable
digraph \mbox{\texttt{\mdseries\slshape digraph}} is changed in\texttt{\symbol{45}}place into its (reflexive) transitive
closure.

 Let $n$ be the number of vertices of \mbox{\texttt{\mdseries\slshape digraph}}, and let $m$ be the number of edges. For an arbitrary digraph, these attributes will use a
version of the Floyd\texttt{\symbol{45}}Warshall algorithm, with complexity $O(n^3)$. However, for a topologically sortable digraph [see \texttt{DigraphTopologicalSort} (\ref{DigraphTopologicalSort})], these attributes will use methods with complexity $O(m + n + m \cdot n)$ when this is faster.

 
\begin{Verbatim}[commandchars=!@|,fontsize=\small,frame=single,label=Example]
  !gapprompt@gap>| !gapinput@D := DigraphFromDiSparse6String(".H`eOWR`Ul^");|
  <immutable digraph with 9 vertices, 8 edges>
  !gapprompt@gap>| !gapinput@IsReflexiveDigraph(D) or IsTransitiveDigraph(D);|
  false
  !gapprompt@gap>| !gapinput@OutNeighbours(D);|
  [ [ 4, 6 ], [ 1, 3 ], [  ], [ 5 ], [  ], [ 7, 8, 9 ], [  ], [  ], 
    [  ] ]
  !gapprompt@gap>| !gapinput@T := DigraphTransitiveClosure(D);|
  <immutable transitive digraph with 9 vertices, 18 edges>
  !gapprompt@gap>| !gapinput@OutNeighbours(T);|
  [ [ 4, 6, 5, 7, 8, 9 ], [ 1, 3, 4, 5, 6, 7, 8, 9 ], [  ], [ 5 ], 
    [  ], [ 7, 8, 9 ], [  ], [  ], [  ] ]
  !gapprompt@gap>| !gapinput@RT := DigraphReflexiveTransitiveClosure(D);|
  <immutable preorder digraph with 9 vertices, 27 edges>
  !gapprompt@gap>| !gapinput@OutNeighbours(RT);|
  [ [ 4, 6, 5, 7, 8, 9, 1 ], [ 1, 3, 4, 5, 6, 7, 8, 9, 2 ], [ 3 ], 
    [ 5, 4 ], [ 5 ], [ 7, 8, 9, 6 ], [ 7 ], [ 8 ], [ 9 ] ]
  !gapprompt@gap>| !gapinput@D := CycleDigraph(IsMutableDigraph, 3);|
  <mutable digraph with 3 vertices, 3 edges>
  !gapprompt@gap>| !gapinput@DigraphReflexiveTransitiveClosure(D);|
  <mutable digraph with 3 vertices, 9 edges>
  !gapprompt@gap>| !gapinput@D;|
  <mutable digraph with 3 vertices, 9 edges>
\end{Verbatim}
 }

 

\subsection{\textcolor{Chapter }{DigraphTransitiveReduction}}
\logpage{[ 3, 3, 14 ]}\nobreak
\hyperdef{L}{X82AD17517E273600}{}
{\noindent\textcolor{FuncColor}{$\triangleright$\enspace\texttt{DigraphTransitiveReduction({\mdseries\slshape digraph})\index{DigraphTransitiveReduction@\texttt{DigraphTransitiveReduction}}
\label{DigraphTransitiveReduction}
}\hfill{\scriptsize (operation)}}\\
\noindent\textcolor{FuncColor}{$\triangleright$\enspace\texttt{DigraphTransitiveReductionAttr({\mdseries\slshape digraph})\index{DigraphTransitiveReductionAttr@\texttt{DigraphTransitiveReductionAttr}}
\label{DigraphTransitiveReductionAttr}
}\hfill{\scriptsize (attribute)}}\\
\noindent\textcolor{FuncColor}{$\triangleright$\enspace\texttt{DigraphReflexiveTransitiveReduction({\mdseries\slshape digraph})\index{DigraphReflexiveTransitiveReduction@\texttt{DigraphReflexiveTransitiveReduction}}
\label{DigraphReflexiveTransitiveReduction}
}\hfill{\scriptsize (operation)}}\\
\noindent\textcolor{FuncColor}{$\triangleright$\enspace\texttt{DigraphReflexiveTransitiveReductionAttr({\mdseries\slshape digraph})\index{DigraphReflexiveTransitiveReductionAttr@\texttt{Digraph}\-\texttt{Reflexive}\-\texttt{Transitive}\-\texttt{Reduction}\-\texttt{Attr}}
\label{DigraphReflexiveTransitiveReductionAttr}
}\hfill{\scriptsize (attribute)}}\\
\textbf{\indent Returns:\ }
A digraph.



 If \mbox{\texttt{\mdseries\slshape digraph}} is a topologically sortable digraph [see \texttt{DigraphTopologicalSort} (\ref{DigraphTopologicalSort})] with no multiple edges, then these operations return the (reflexive)
transitive reduction of \mbox{\texttt{\mdseries\slshape digraph}}.

 The (reflexive) transitive reduction of such a digraph is the unique least
subgraph such that the (reflexive) transitive closure of the subgraph is equal
to the (reflexive) transitive closure of \mbox{\texttt{\mdseries\slshape digraph}} [see \texttt{DigraphReflexiveTransitiveClosure} (\ref{DigraphReflexiveTransitiveClosure})]. In order words, it is the least subgraph of \mbox{\texttt{\mdseries\slshape digraph}} which retains the same reachability as \mbox{\texttt{\mdseries\slshape digraph}}.

 If \mbox{\texttt{\mdseries\slshape digraph}} is immutable, then a new immutable digraph is returned. Otherwise, the mutable
digraph \mbox{\texttt{\mdseries\slshape digraph}} is changed in\texttt{\symbol{45}}place into its (reflexive) transitive
reduction.

 Let $n$ be the number of vertices of an arbitrary digraph, and let $m$ be the number of edges. Then these operations use methods with complexity $O(m + n + m \cdot n)$. 

 
\begin{Verbatim}[commandchars=!@|,fontsize=\small,frame=single,label=Example]
  !gapprompt@gap>| !gapinput@D := Digraph([[1, 2, 3], [3], [3]]);;|
  !gapprompt@gap>| !gapinput@DigraphHasLoops(D);|
  true
  !gapprompt@gap>| !gapinput@D1 := DigraphReflexiveTransitiveReduction(D);|
  <immutable digraph with 3 vertices, 2 edges>
  !gapprompt@gap>| !gapinput@DigraphHasLoops(D1);|
  false
  !gapprompt@gap>| !gapinput@OutNeighbours(D1);|
  [ [ 2 ], [ 3 ], [  ] ]
  !gapprompt@gap>| !gapinput@D2 := DigraphTransitiveReduction(D);|
  <immutable digraph with 3 vertices, 4 edges>
  !gapprompt@gap>| !gapinput@DigraphHasLoops(D2);|
  true
  !gapprompt@gap>| !gapinput@OutNeighbours(D2);|
  [ [ 2, 1 ], [ 3 ], [ 3 ] ]
  !gapprompt@gap>| !gapinput@DigraphReflexiveTransitiveClosure(D)|
  !gapprompt@>| !gapinput@ = DigraphReflexiveTransitiveClosure(D1);|
  true
  !gapprompt@gap>| !gapinput@DigraphTransitiveClosure(D)|
  !gapprompt@>| !gapinput@ = DigraphTransitiveClosure(D2);|
  true
  !gapprompt@gap>| !gapinput@D := Digraph(IsMutableDigraph, [[1], [1], [1, 2, 3]]);|
  <mutable digraph with 3 vertices, 5 edges>
  !gapprompt@gap>| !gapinput@DigraphReflexiveTransitiveReduction(D);|
  <mutable digraph with 3 vertices, 2 edges>
  !gapprompt@gap>| !gapinput@D;|
  <mutable digraph with 3 vertices, 2 edges>
\end{Verbatim}
 }

 

\subsection{\textcolor{Chapter }{DigraphAddVertex}}
\logpage{[ 3, 3, 15 ]}\nobreak
\hyperdef{L}{X83B2506D79453208}{}
{\noindent\textcolor{FuncColor}{$\triangleright$\enspace\texttt{DigraphAddVertex({\mdseries\slshape digraph[, label]})\index{DigraphAddVertex@\texttt{DigraphAddVertex}}
\label{DigraphAddVertex}
}\hfill{\scriptsize (operation)}}\\
\textbf{\indent Returns:\ }
A digraph.



 The operation returns a digraph constructed from \mbox{\texttt{\mdseries\slshape digraph}} by adding a single new vertex, and no new edges. 

 If the optional second argument \mbox{\texttt{\mdseries\slshape label}} is a \textsf{GAP} object, then the new vertex will be labelled \mbox{\texttt{\mdseries\slshape label}}. 

 If \mbox{\texttt{\mdseries\slshape digraph}} belongs to \texttt{IsMutableDigraph} (\ref{IsMutableDigraph}), then the vertex is added directly to \mbox{\texttt{\mdseries\slshape digraph}}. If \mbox{\texttt{\mdseries\slshape digraph}} belongs to \texttt{IsImmutableDigraph} (\ref{IsImmutableDigraph}), an immutable copy of \mbox{\texttt{\mdseries\slshape digraph}} with the additional vertex is returned.

 
\begin{Verbatim}[commandchars=!@|,fontsize=\small,frame=single,label=Example]
  !gapprompt@gap>| !gapinput@D := CompleteDigraph(3);|
  <immutable complete digraph with 3 vertices>
  !gapprompt@gap>| !gapinput@new := DigraphAddVertex(D);|
  <immutable digraph with 4 vertices, 6 edges>
  !gapprompt@gap>| !gapinput@D = new;|
  false
  !gapprompt@gap>| !gapinput@DigraphVertices(new);|
  [ 1 .. 4 ]
  !gapprompt@gap>| !gapinput@new := DigraphAddVertex(D, Group([(1, 2)]));|
  <immutable digraph with 4 vertices, 6 edges>
  !gapprompt@gap>| !gapinput@DigraphVertexLabels(new);|
  [ 1, 2, 3, Group([ (1,2) ]) ]
  !gapprompt@gap>| !gapinput@D := CompleteBipartiteDigraph(IsMutableDigraph, 2, 3);|
  <mutable digraph with 5 vertices, 12 edges>
  !gapprompt@gap>| !gapinput@new := DigraphAddVertex(D);|
  <mutable digraph with 6 vertices, 12 edges>
  !gapprompt@gap>| !gapinput@D = new;|
  true
\end{Verbatim}
 }

 

\subsection{\textcolor{Chapter }{DigraphAddVertices (for a digraph and an integer)}}
\logpage{[ 3, 3, 16 ]}\nobreak
\hyperdef{L}{X8134BEE7786BD3A7}{}
{\noindent\textcolor{FuncColor}{$\triangleright$\enspace\texttt{DigraphAddVertices({\mdseries\slshape digraph, m})\index{DigraphAddVertices@\texttt{DigraphAddVertices}!for a digraph and an integer}
\label{DigraphAddVertices:for a digraph and an integer}
}\hfill{\scriptsize (operation)}}\\
\noindent\textcolor{FuncColor}{$\triangleright$\enspace\texttt{DigraphAddVertices({\mdseries\slshape digraph, labels})\index{DigraphAddVertices@\texttt{DigraphAddVertices}!for a digraph and a list of labels}
\label{DigraphAddVertices:for a digraph and a list of labels}
}\hfill{\scriptsize (operation)}}\\
\textbf{\indent Returns:\ }
A digraph.



 For a non\texttt{\symbol{45}}negative integer \mbox{\texttt{\mdseries\slshape m}}, this operation returns a digraph constructed from \mbox{\texttt{\mdseries\slshape digraph}} by adding \mbox{\texttt{\mdseries\slshape m}} new vertices. 

 Otherwise, if \mbox{\texttt{\mdseries\slshape labels}} is a list consisting of \texttt{k} \textsf{GAP} objects, then this operation returns a digraph constructed from \mbox{\texttt{\mdseries\slshape digraph}} by adding \texttt{k} new vertices, which are labelled according to this list. 

 If \mbox{\texttt{\mdseries\slshape digraph}} belongs to \texttt{IsMutableDigraph} (\ref{IsMutableDigraph}), then the vertices are added directly to \mbox{\texttt{\mdseries\slshape digraph}}, which is changed in\texttt{\symbol{45}}place. If \mbox{\texttt{\mdseries\slshape digraph}} belongs to \texttt{IsImmutableDigraph} (\ref{IsImmutableDigraph}), then \mbox{\texttt{\mdseries\slshape digraph}} itself is returned if no vertices are added (i.e. \texttt{\mbox{\texttt{\mdseries\slshape m}}=0} or \mbox{\texttt{\mdseries\slshape labels}} is empty), otherwise the result is a new immutable digraph. 

 
\begin{Verbatim}[commandchars=!@|,fontsize=\small,frame=single,label=Example]
  !gapprompt@gap>| !gapinput@D := CompleteDigraph(3);|
  <immutable complete digraph with 3 vertices>
  !gapprompt@gap>| !gapinput@new := DigraphAddVertices(D, 3);|
  <immutable digraph with 6 vertices, 6 edges>
  !gapprompt@gap>| !gapinput@DigraphVertices(new);|
  [ 1 .. 6 ]
  !gapprompt@gap>| !gapinput@new := DigraphAddVertices(D, [Group([(1, 2)]), "d"]);|
  <immutable digraph with 5 vertices, 6 edges>
  !gapprompt@gap>| !gapinput@DigraphVertexLabels(new);|
  [ 1, 2, 3, Group([ (1,2) ]), "d" ]
  !gapprompt@gap>| !gapinput@DigraphAddVertices(D, 0) = D;|
  true
  !gapprompt@gap>| !gapinput@D := CompleteBipartiteDigraph(IsMutableDigraph, 2, 3);|
  <mutable digraph with 5 vertices, 12 edges>
  !gapprompt@gap>| !gapinput@new := DigraphAddVertices(D, 4);|
  <mutable digraph with 9 vertices, 12 edges>
  !gapprompt@gap>| !gapinput@D = new;|
  true
\end{Verbatim}
 }

 

\subsection{\textcolor{Chapter }{DigraphAddEdge (for a digraph and an edge)}}
\logpage{[ 3, 3, 17 ]}\nobreak
\hyperdef{L}{X7E58CC4880627658}{}
{\noindent\textcolor{FuncColor}{$\triangleright$\enspace\texttt{DigraphAddEdge({\mdseries\slshape digraph, edge})\index{DigraphAddEdge@\texttt{DigraphAddEdge}!for a digraph and an edge}
\label{DigraphAddEdge:for a digraph and an edge}
}\hfill{\scriptsize (operation)}}\\
\noindent\textcolor{FuncColor}{$\triangleright$\enspace\texttt{DigraphAddEdge({\mdseries\slshape digraph, src, ran})\index{DigraphAddEdge@\texttt{DigraphAddEdge}!for a digraph, source, and range}
\label{DigraphAddEdge:for a digraph, source, and range}
}\hfill{\scriptsize (operation)}}\\
\textbf{\indent Returns:\ }
A digraph.



 If \mbox{\texttt{\mdseries\slshape edge}} is a pair of vertices of \mbox{\texttt{\mdseries\slshape digraph}}, or \mbox{\texttt{\mdseries\slshape src}} and \mbox{\texttt{\mdseries\slshape ran}} are vertices of \mbox{\texttt{\mdseries\slshape digraph}}, then this operation returns a digraph constructed from \mbox{\texttt{\mdseries\slshape digraph}} by adding a new edge with source \mbox{\texttt{\mdseries\slshape edge}}\texttt{[1]} [\mbox{\texttt{\mdseries\slshape src}}] and range \mbox{\texttt{\mdseries\slshape edge}}\texttt{[2]} [\mbox{\texttt{\mdseries\slshape ran}}]. 

 If \mbox{\texttt{\mdseries\slshape digraph}} belongs to \texttt{IsMutableDigraph} (\ref{IsMutableDigraph}), then the edge is added directly to \mbox{\texttt{\mdseries\slshape digraph}}. If \mbox{\texttt{\mdseries\slshape digraph}} belongs to \texttt{IsImmutableDigraph} (\ref{IsImmutableDigraph}), then an immutable copy of \mbox{\texttt{\mdseries\slshape digraph}} with the additional edge is returned. 

 
\begin{Verbatim}[commandchars=!@|,fontsize=\small,frame=single,label=Example]
  !gapprompt@gap>| !gapinput@D1 := Digraph([[2], [3], []]);|
  <immutable digraph with 3 vertices, 2 edges>
  !gapprompt@gap>| !gapinput@DigraphEdges(D1);|
  [ [ 1, 2 ], [ 2, 3 ] ]
  !gapprompt@gap>| !gapinput@D2 := DigraphAddEdge(D1, [3, 1]);|
  <immutable digraph with 3 vertices, 3 edges>
  !gapprompt@gap>| !gapinput@DigraphEdges(D2);|
  [ [ 1, 2 ], [ 2, 3 ], [ 3, 1 ] ]
  !gapprompt@gap>| !gapinput@D3 := DigraphAddEdge(D2, [2, 3]);|
  <immutable multidigraph with 3 vertices, 4 edges>
  !gapprompt@gap>| !gapinput@DigraphEdges(D3);|
  [ [ 1, 2 ], [ 2, 3 ], [ 2, 3 ], [ 3, 1 ] ]
  !gapprompt@gap>| !gapinput@D := CycleDigraph(IsMutableDigraph, 4);|
  <mutable digraph with 4 vertices, 4 edges>
  !gapprompt@gap>| !gapinput@new := DigraphAddEdge(D, [1, 3]);|
  <mutable digraph with 4 vertices, 5 edges>
  !gapprompt@gap>| !gapinput@DigraphEdges(new);|
  [ [ 1, 2 ], [ 1, 3 ], [ 2, 3 ], [ 3, 4 ], [ 4, 1 ] ]
  !gapprompt@gap>| !gapinput@D = new;|
  true
\end{Verbatim}
 }

 

\subsection{\textcolor{Chapter }{DigraphAddEdgeOrbit}}
\logpage{[ 3, 3, 18 ]}\nobreak
\hyperdef{L}{X7BE5C7028760B053}{}
{\noindent\textcolor{FuncColor}{$\triangleright$\enspace\texttt{DigraphAddEdgeOrbit({\mdseries\slshape digraph, edge})\index{DigraphAddEdgeOrbit@\texttt{DigraphAddEdgeOrbit}}
\label{DigraphAddEdgeOrbit}
}\hfill{\scriptsize (operation)}}\\
\textbf{\indent Returns:\ }
 A new digraph. 



 This operation returns a new digraph with the same vertices and edges as \mbox{\texttt{\mdseries\slshape digraph}} and with additional edges consisting of the orbit of the edge \mbox{\texttt{\mdseries\slshape edge}} under the action of the \texttt{DigraphGroup} (\ref{DigraphGroup}) of \mbox{\texttt{\mdseries\slshape digraph}}. If \mbox{\texttt{\mdseries\slshape edge}} is already an edge in \mbox{\texttt{\mdseries\slshape digraph}}, then \mbox{\texttt{\mdseries\slshape digraph}} is returned unchanged. The argument \mbox{\texttt{\mdseries\slshape digraph}} must be an immutable digraph. 

 An edge is simply a pair of vertices of \mbox{\texttt{\mdseries\slshape digraph}}. 
\begin{Verbatim}[commandchars=!@|,fontsize=\small,frame=single,label=Example]
  !gapprompt@gap>| !gapinput@gr1 := CayleyDigraph(DihedralGroup(8));|
  <immutable digraph with 8 vertices, 24 edges>
  !gapprompt@gap>| !gapinput@gr2 := DigraphAddEdgeOrbit(gr1, [1, 8]);|
  <immutable digraph with 8 vertices, 32 edges>
  !gapprompt@gap>| !gapinput@DigraphEdges(gr1);|
  [ [ 1, 2 ], [ 1, 3 ], [ 1, 4 ], [ 2, 1 ], [ 2, 5 ], [ 2, 6 ], 
    [ 3, 8 ], [ 3, 4 ], [ 3, 7 ], [ 4, 6 ], [ 4, 7 ], [ 4, 1 ], 
    [ 5, 7 ], [ 5, 6 ], [ 5, 8 ], [ 6, 4 ], [ 6, 8 ], [ 6, 2 ], 
    [ 7, 5 ], [ 7, 1 ], [ 7, 3 ], [ 8, 3 ], [ 8, 2 ], [ 8, 5 ] ]
  !gapprompt@gap>| !gapinput@DigraphEdges(gr2);|
  [ [ 1, 2 ], [ 1, 3 ], [ 1, 4 ], [ 1, 8 ], [ 2, 1 ], [ 2, 5 ], 
    [ 2, 6 ], [ 2, 7 ], [ 3, 8 ], [ 3, 4 ], [ 3, 7 ], [ 3, 6 ], 
    [ 4, 6 ], [ 4, 7 ], [ 4, 1 ], [ 4, 5 ], [ 5, 7 ], [ 5, 6 ], 
    [ 5, 8 ], [ 5, 4 ], [ 6, 4 ], [ 6, 8 ], [ 6, 2 ], [ 6, 3 ], 
    [ 7, 5 ], [ 7, 1 ], [ 7, 3 ], [ 7, 2 ], [ 8, 3 ], [ 8, 2 ], 
    [ 8, 5 ], [ 8, 1 ] ]
  !gapprompt@gap>| !gapinput@gr3 := DigraphRemoveEdgeOrbit(gr2, [1, 8]);|
  <immutable digraph with 8 vertices, 24 edges>
  !gapprompt@gap>| !gapinput@gr3 = gr1;|
  true
\end{Verbatim}
 }

 

\subsection{\textcolor{Chapter }{DigraphAddEdges}}
\logpage{[ 3, 3, 19 ]}\nobreak
\hyperdef{L}{X8693A61B7F752C76}{}
{\noindent\textcolor{FuncColor}{$\triangleright$\enspace\texttt{DigraphAddEdges({\mdseries\slshape digraph, edges})\index{DigraphAddEdges@\texttt{DigraphAddEdges}}
\label{DigraphAddEdges}
}\hfill{\scriptsize (operation)}}\\
\textbf{\indent Returns:\ }
A digraph.



 If \mbox{\texttt{\mdseries\slshape edges}} is a (possibly empty) list of pairs of vertices of \mbox{\texttt{\mdseries\slshape digraph}}, then this operation returns a digraph constructed from \mbox{\texttt{\mdseries\slshape digraph}} by adding the edges specified by \mbox{\texttt{\mdseries\slshape edges}}. More precisely, for every \texttt{edge} in \mbox{\texttt{\mdseries\slshape edges}}, a new edge will be added with source \texttt{edge[1]} and range \texttt{edges[2]}. 

 If an edge is included in \mbox{\texttt{\mdseries\slshape edges}} with multiplicity \texttt{k}, then it will be added \texttt{k} times. If \mbox{\texttt{\mdseries\slshape digraph}} belongs to \texttt{IsMutableDigraph} (\ref{IsMutableDigraph}), then the edges are added directly to \mbox{\texttt{\mdseries\slshape digraph}}. If \mbox{\texttt{\mdseries\slshape digraph}} belongs to \texttt{IsImmutableDigraph} (\ref{IsImmutableDigraph}), then the result is returned as an immutable digraph. 

 
\begin{Verbatim}[commandchars=!@|,fontsize=\small,frame=single,label=Example]
  !gapprompt@gap>| !gapinput@func := function(n)|
  !gapprompt@>| !gapinput@ local source, range, i;|
  !gapprompt@>| !gapinput@ source := [];|
  !gapprompt@>| !gapinput@ range  := [];|
  !gapprompt@>| !gapinput@ for i in [1 .. n - 2] do|
  !gapprompt@>| !gapinput@   Add(source, i);|
  !gapprompt@>| !gapinput@   Add(range, i + 1);|
  !gapprompt@>| !gapinput@ od;|
  !gapprompt@>| !gapinput@ return Digraph(n, source, range);|
  !gapprompt@>| !gapinput@end;;|
  !gapprompt@gap>| !gapinput@D := func(1024);|
  <immutable digraph with 1024 vertices, 1022 edges>
  !gapprompt@gap>| !gapinput@new := DigraphAddEdges(D,|
  !gapprompt@>| !gapinput@[[1023, 1024], [1, 1024], [1023, 1024], [1024, 1]]);|
  <immutable multidigraph with 1024 vertices, 1026 edges>
  !gapprompt@gap>| !gapinput@D = new;|
  false
  !gapprompt@gap>| !gapinput@D2 := DigraphMutableCopy(func(1024));|
  <mutable digraph with 1024 vertices, 1022 edges>
  !gapprompt@gap>| !gapinput@new := DigraphAddEdges(D2,|
  !gapprompt@>| !gapinput@[[1023, 1024], [1, 1024], [1023, 1024], [1024, 1]]);|
  <mutable multidigraph with 1024 vertices, 1026 edges>
  !gapprompt@gap>| !gapinput@D2 = new;|
  true
\end{Verbatim}
 }

 

\subsection{\textcolor{Chapter }{DigraphRemoveVertex}}
\logpage{[ 3, 3, 20 ]}\nobreak
\hyperdef{L}{X7B634A2B83C08B16}{}
{\noindent\textcolor{FuncColor}{$\triangleright$\enspace\texttt{DigraphRemoveVertex({\mdseries\slshape digraph, v})\index{DigraphRemoveVertex@\texttt{DigraphRemoveVertex}}
\label{DigraphRemoveVertex}
}\hfill{\scriptsize (operation)}}\\
\textbf{\indent Returns:\ }
A digraph.



 If \mbox{\texttt{\mdseries\slshape v}} is a vertex of \mbox{\texttt{\mdseries\slshape digraph}}, then this operation returns a digraph constructed from \mbox{\texttt{\mdseries\slshape digraph}} by removing vertex \mbox{\texttt{\mdseries\slshape v}}, along with any edge whose source or range vertex is \mbox{\texttt{\mdseries\slshape v}}.

 If \mbox{\texttt{\mdseries\slshape digraph}} has \texttt{n} vertices, then the vertices of the returned digraph are \texttt{[1..n\texttt{\symbol{45}}1]}, but the original labels can be accessed via \texttt{DigraphVertexLabels} (\ref{DigraphVertexLabels}). 

 If \mbox{\texttt{\mdseries\slshape digraph}} belongs to \texttt{IsMutableDigraph} (\ref{IsMutableDigraph}), then the vertex is removed directly from \mbox{\texttt{\mdseries\slshape digraph}}. If \mbox{\texttt{\mdseries\slshape digraph}} belongs to \texttt{IsImmutableDigraph} (\ref{IsImmutableDigraph}), an immutable copy of \mbox{\texttt{\mdseries\slshape digraph}} without the vertex is returned.

 
\begin{Verbatim}[commandchars=!@|,fontsize=\small,frame=single,label=Example]
  !gapprompt@gap>| !gapinput@D := Digraph(["a", "b", "c"],|
  !gapprompt@>| !gapinput@                 ["a", "a", "b", "c", "c"],|
  !gapprompt@>| !gapinput@                 ["b", "c", "a", "a", "c"]);|
  <immutable digraph with 3 vertices, 5 edges>
  !gapprompt@gap>| !gapinput@DigraphVertexLabels(D);|
  [ "a", "b", "c" ]
  !gapprompt@gap>| !gapinput@DigraphEdges(D);|
  [ [ 1, 2 ], [ 1, 3 ], [ 2, 1 ], [ 3, 1 ], [ 3, 3 ] ]
  !gapprompt@gap>| !gapinput@new := DigraphRemoveVertex(D, 2);|
  <immutable digraph with 2 vertices, 3 edges>
  !gapprompt@gap>| !gapinput@DigraphVertexLabels(new);|
  [ "a", "c" ]
  !gapprompt@gap>| !gapinput@D := CycleDigraph(IsMutableDigraph, 5);|
  <mutable digraph with 5 vertices, 5 edges>
  !gapprompt@gap>| !gapinput@new := DigraphRemoveVertex(D, 1);|
  <mutable digraph with 4 vertices, 3 edges>
  !gapprompt@gap>| !gapinput@DigraphVertexLabels(D);|
  [ 2, 3, 4, 5 ]
  !gapprompt@gap>| !gapinput@D = new;|
  true
\end{Verbatim}
 }

 

\subsection{\textcolor{Chapter }{DigraphRemoveVertices}}
\logpage{[ 3, 3, 21 ]}\nobreak
\hyperdef{L}{X7E290E847A5A299A}{}
{\noindent\textcolor{FuncColor}{$\triangleright$\enspace\texttt{DigraphRemoveVertices({\mdseries\slshape digraph, verts})\index{DigraphRemoveVertices@\texttt{DigraphRemoveVertices}}
\label{DigraphRemoveVertices}
}\hfill{\scriptsize (operation)}}\\
\textbf{\indent Returns:\ }
A digraph.



 If \mbox{\texttt{\mdseries\slshape verts}} is a (possibly empty) duplicate\texttt{\symbol{45}}free list of vertices of \mbox{\texttt{\mdseries\slshape digraph}}, then this operation returns a digraph constructed from \mbox{\texttt{\mdseries\slshape digraph}} by removing every vertex in \mbox{\texttt{\mdseries\slshape verts}}, along with any edge whose source or range vertex is in \mbox{\texttt{\mdseries\slshape verts}}.

 If \mbox{\texttt{\mdseries\slshape digraph}} has \texttt{n} vertices, then the vertices of the new digraph are \texttt{[1 .. n\texttt{\symbol{45}}Length(\mbox{\texttt{\mdseries\slshape verts}})]}, but the original labels can be accessed via \texttt{DigraphVertexLabels} (\ref{DigraphVertexLabels}). 

 If \mbox{\texttt{\mdseries\slshape digraph}} belongs to \texttt{IsMutableDigraph} (\ref{IsMutableDigraph}), then the vertices are removed directly from \mbox{\texttt{\mdseries\slshape digraph}}. If \mbox{\texttt{\mdseries\slshape digraph}} belongs to \texttt{IsImmutableDigraph} (\ref{IsImmutableDigraph}), an immutable copy of \mbox{\texttt{\mdseries\slshape digraph}} without the vertices is returned.

 
\begin{Verbatim}[commandchars=!@|,fontsize=\small,frame=single,label=Example]
  !gapprompt@gap>| !gapinput@D := Digraph([[3], [1, 3, 5], [1], [1, 2, 4], [2, 3, 5]]);|
  <immutable digraph with 5 vertices, 11 edges>
  !gapprompt@gap>| !gapinput@SetDigraphVertexLabels(D, ["a", "b", "c", "d", "e"]);|
  !gapprompt@gap>| !gapinput@new := DigraphRemoveVertices(D, [2, 4]);|
  <immutable digraph with 3 vertices, 4 edges>
  !gapprompt@gap>| !gapinput@DigraphVertexLabels(new);|
  [ "a", "c", "e" ]
  !gapprompt@gap>| !gapinput@D := CycleDigraph(IsMutableDigraph, 5);|
  <mutable digraph with 5 vertices, 5 edges>
  !gapprompt@gap>| !gapinput@new := DigraphRemoveVertices(D, [1, 3]);|
  <mutable digraph with 3 vertices, 1 edge>
  !gapprompt@gap>| !gapinput@DigraphVertexLabels(D);|
  [ 2, 4, 5 ]
  !gapprompt@gap>| !gapinput@D = new;|
  true
\end{Verbatim}
 }

 

\subsection{\textcolor{Chapter }{DigraphRemoveEdge (for a digraph and an edge)}}
\logpage{[ 3, 3, 22 ]}\nobreak
\hyperdef{L}{X8433E3BC7E5EA6BF}{}
{\noindent\textcolor{FuncColor}{$\triangleright$\enspace\texttt{DigraphRemoveEdge({\mdseries\slshape digraph, edge})\index{DigraphRemoveEdge@\texttt{DigraphRemoveEdge}!for a digraph and an edge}
\label{DigraphRemoveEdge:for a digraph and an edge}
}\hfill{\scriptsize (operation)}}\\
\noindent\textcolor{FuncColor}{$\triangleright$\enspace\texttt{DigraphRemoveEdge({\mdseries\slshape digraph, src, ran})\index{DigraphRemoveEdge@\texttt{DigraphRemoveEdge}!for a digraph, source, and range}
\label{DigraphRemoveEdge:for a digraph, source, and range}
}\hfill{\scriptsize (operation)}}\\
\textbf{\indent Returns:\ }
A digraph.



 If \mbox{\texttt{\mdseries\slshape digraph}} is a digraph with no multiple edges and \mbox{\texttt{\mdseries\slshape edge}} is a pair of vertices of \mbox{\texttt{\mdseries\slshape digraph}}, or \mbox{\texttt{\mdseries\slshape src}} and \mbox{\texttt{\mdseries\slshape ran}} are vertices of \mbox{\texttt{\mdseries\slshape digraph}}, then this operation returns a digraph constructed from \mbox{\texttt{\mdseries\slshape digraph}} by removing the edge specified by \mbox{\texttt{\mdseries\slshape edge}} or \mbox{\texttt{\mdseries\slshape [src, ran]}}. 

 If \mbox{\texttt{\mdseries\slshape digraph}} belongs to \texttt{IsMutableDigraph} (\ref{IsMutableDigraph}), then the edge is removed directly from \mbox{\texttt{\mdseries\slshape digraph}}. If \mbox{\texttt{\mdseries\slshape digraph}} belongs to \texttt{IsImmutableDigraph} (\ref{IsImmutableDigraph}), an immutable copy of \mbox{\texttt{\mdseries\slshape digraph}} without the edge is returned. 

 Note that if \mbox{\texttt{\mdseries\slshape digraph}} belongs to \texttt{IsImmutableDigraph} (\ref{IsImmutableDigraph}), then a new copy of \mbox{\texttt{\mdseries\slshape digraph}} will be returned even if \mbox{\texttt{\mdseries\slshape edge}} or \mbox{\texttt{\mdseries\slshape [src, ran]}} does not define an edge of \mbox{\texttt{\mdseries\slshape digraph}}.

 
\begin{Verbatim}[commandchars=!@|,fontsize=\small,frame=single,label=Example]
  !gapprompt@gap>| !gapinput@D := CycleDigraph(250000);|
  <immutable cycle digraph with 250000 vertices>
  !gapprompt@gap>| !gapinput@D := DigraphRemoveEdge(D, [250000, 1]);|
  <immutable digraph with 250000 vertices, 249999 edges>
  !gapprompt@gap>| !gapinput@new := DigraphRemoveEdge(D, [25000, 2]);;|
  !gapprompt@gap>| !gapinput@new = D;|
  true
  !gapprompt@gap>| !gapinput@IsIdenticalObj(new, D);|
  false
  !gapprompt@gap>| !gapinput@D := DigraphMutableCopy(D);;|
  !gapprompt@gap>| !gapinput@new := DigraphRemoveEdge(D, 2500, 2);;|
  !gapprompt@gap>| !gapinput@IsIdenticalObj(new, D);|
  true
\end{Verbatim}
 }

 

\subsection{\textcolor{Chapter }{DigraphRemoveEdgeOrbit}}
\logpage{[ 3, 3, 23 ]}\nobreak
\hyperdef{L}{X85981D9187F49018}{}
{\noindent\textcolor{FuncColor}{$\triangleright$\enspace\texttt{DigraphRemoveEdgeOrbit({\mdseries\slshape digraph, edge})\index{DigraphRemoveEdgeOrbit@\texttt{DigraphRemoveEdgeOrbit}}
\label{DigraphRemoveEdgeOrbit}
}\hfill{\scriptsize (operation)}}\\
\textbf{\indent Returns:\ }
 A new digraph. 



 This operation returns a new digraph with the same vertices as \mbox{\texttt{\mdseries\slshape digraph}} and with the orbit of the edge \mbox{\texttt{\mdseries\slshape edge}} (under the action of the \texttt{DigraphGroup} (\ref{DigraphGroup}) of \mbox{\texttt{\mdseries\slshape digraph}}) removed. If \mbox{\texttt{\mdseries\slshape edge}} is not an edge in \mbox{\texttt{\mdseries\slshape digraph}}, then \mbox{\texttt{\mdseries\slshape digraph}} is returned unchanged. The argument \mbox{\texttt{\mdseries\slshape digraph}} must be an immutable digraph. 

 An edge is simply a pair of vertices of \mbox{\texttt{\mdseries\slshape digraph}}. 
\begin{Verbatim}[commandchars=!@|,fontsize=\small,frame=single,label=Example]
  !gapprompt@gap>| !gapinput@gr1 := CayleyDigraph(DihedralGroup(8));|
  <immutable digraph with 8 vertices, 24 edges>
  !gapprompt@gap>| !gapinput@gr2 := DigraphAddEdgeOrbit(gr1, [1, 8]);|
  <immutable digraph with 8 vertices, 32 edges>
  !gapprompt@gap>| !gapinput@DigraphEdges(gr1);|
  [ [ 1, 2 ], [ 1, 3 ], [ 1, 4 ], [ 2, 1 ], [ 2, 5 ], [ 2, 6 ], 
    [ 3, 8 ], [ 3, 4 ], [ 3, 7 ], [ 4, 6 ], [ 4, 7 ], [ 4, 1 ], 
    [ 5, 7 ], [ 5, 6 ], [ 5, 8 ], [ 6, 4 ], [ 6, 8 ], [ 6, 2 ], 
    [ 7, 5 ], [ 7, 1 ], [ 7, 3 ], [ 8, 3 ], [ 8, 2 ], [ 8, 5 ] ]
  !gapprompt@gap>| !gapinput@DigraphEdges(gr2);|
  [ [ 1, 2 ], [ 1, 3 ], [ 1, 4 ], [ 1, 8 ], [ 2, 1 ], [ 2, 5 ], 
    [ 2, 6 ], [ 2, 7 ], [ 3, 8 ], [ 3, 4 ], [ 3, 7 ], [ 3, 6 ], 
    [ 4, 6 ], [ 4, 7 ], [ 4, 1 ], [ 4, 5 ], [ 5, 7 ], [ 5, 6 ], 
    [ 5, 8 ], [ 5, 4 ], [ 6, 4 ], [ 6, 8 ], [ 6, 2 ], [ 6, 3 ], 
    [ 7, 5 ], [ 7, 1 ], [ 7, 3 ], [ 7, 2 ], [ 8, 3 ], [ 8, 2 ], 
    [ 8, 5 ], [ 8, 1 ] ]
  !gapprompt@gap>| !gapinput@gr3 := DigraphRemoveEdgeOrbit(gr2, [1, 8]);|
  <immutable digraph with 8 vertices, 24 edges>
  !gapprompt@gap>| !gapinput@gr3 = gr1;|
  true
\end{Verbatim}
 }

 

\subsection{\textcolor{Chapter }{DigraphRemoveEdges}}
\logpage{[ 3, 3, 24 ]}\nobreak
\hyperdef{L}{X87093FDA7F88E732}{}
{\noindent\textcolor{FuncColor}{$\triangleright$\enspace\texttt{DigraphRemoveEdges({\mdseries\slshape digraph, edges})\index{DigraphRemoveEdges@\texttt{DigraphRemoveEdges}}
\label{DigraphRemoveEdges}
}\hfill{\scriptsize (operation)}}\\
\textbf{\indent Returns:\ }
A digraph.



 If one of the following holds: 
\begin{itemize}
\item  \mbox{\texttt{\mdseries\slshape digraph}} is a digraph with no multiple edges, and \mbox{\texttt{\mdseries\slshape edges}} is a list of pairs of vertices of \mbox{\texttt{\mdseries\slshape digraph}}, or 
\item  \mbox{\texttt{\mdseries\slshape digraph}} is a digraph and \mbox{\texttt{\mdseries\slshape edges}} is an empty list 
\end{itemize}
 then this operation returns a digraph constructed from \mbox{\texttt{\mdseries\slshape digraph}} by removing all of the edges specified by \mbox{\texttt{\mdseries\slshape edges}} (see \texttt{DigraphRemoveEdge} (\ref{DigraphRemoveEdge:for a digraph and an edge})). 

 If \mbox{\texttt{\mdseries\slshape digraph}} belongs to \texttt{IsMutableDigraph} (\ref{IsMutableDigraph}), then the edge is removed directly from \mbox{\texttt{\mdseries\slshape digraph}}. If \mbox{\texttt{\mdseries\slshape digraph}} belongs to \texttt{IsImmutableDigraph} (\ref{IsImmutableDigraph}), the edge is removed from an immutable copy of \mbox{\texttt{\mdseries\slshape digraph}} and this new digraph is returned.

 Note that if \mbox{\texttt{\mdseries\slshape edges}} is empty, then this operation will always return \mbox{\texttt{\mdseries\slshape digraph}} rather than a copy. Also, if any element of \mbox{\texttt{\mdseries\slshape edges}} is invalid (i.e. does not define an edge of \mbox{\texttt{\mdseries\slshape digraph}}) then that element will simply be ignored. 
\begin{Verbatim}[commandchars=!@|,fontsize=\small,frame=single,label=Example]
  !gapprompt@gap>| !gapinput@D := CycleDigraph(250000);|
  <immutable cycle digraph with 250000 vertices>
  !gapprompt@gap>| !gapinput@D := DigraphRemoveEdges(D, [[250000, 1]]);|
  <immutable digraph with 250000 vertices, 249999 edges>
  !gapprompt@gap>| !gapinput@D := DigraphMutableCopy(D);|
  <mutable digraph with 250000 vertices, 249999 edges>
  !gapprompt@gap>| !gapinput@new := DigraphRemoveEdges(D, [[1, 2], [2, 3], [3, 100]]);|
  <mutable digraph with 250000 vertices, 249997 edges>
  !gapprompt@gap>| !gapinput@new = D;|
  true
\end{Verbatim}
 }

 

\subsection{\textcolor{Chapter }{DigraphRemoveLoops}}
\logpage{[ 3, 3, 25 ]}\nobreak
\hyperdef{L}{X79324AF7818C0C02}{}
{\noindent\textcolor{FuncColor}{$\triangleright$\enspace\texttt{DigraphRemoveLoops({\mdseries\slshape digraph})\index{DigraphRemoveLoops@\texttt{DigraphRemoveLoops}}
\label{DigraphRemoveLoops}
}\hfill{\scriptsize (operation)}}\\
\noindent\textcolor{FuncColor}{$\triangleright$\enspace\texttt{DigraphRemoveLoopsAttr({\mdseries\slshape digraph})\index{DigraphRemoveLoopsAttr@\texttt{DigraphRemoveLoopsAttr}}
\label{DigraphRemoveLoopsAttr}
}\hfill{\scriptsize (attribute)}}\\
\textbf{\indent Returns:\ }
A digraph.



 If \mbox{\texttt{\mdseries\slshape digraph}} is a digraph, then this operation returns a digraph constructed from \mbox{\texttt{\mdseries\slshape digraph}} by removing every loop. A loop is an edge with equal source and range.

 If \mbox{\texttt{\mdseries\slshape digraph}} is immutable, then a new immutable digraph is returned. Otherwise, the loops
are removed from the mutable digraph \mbox{\texttt{\mdseries\slshape digraph}} in\texttt{\symbol{45}}place. 
\begin{Verbatim}[commandchars=!@|,fontsize=\small,frame=single,label=Example]
  !gapprompt@gap>| !gapinput@D := Digraph([[1, 2, 4], [1, 4], [3, 4], [1, 4, 5], [1, 5]]);|
  <immutable digraph with 5 vertices, 12 edges>
  !gapprompt@gap>| !gapinput@DigraphRemoveLoops(D);|
  <immutable digraph with 5 vertices, 8 edges>
  !gapprompt@gap>| !gapinput@D := Digraph(IsMutableDigraph, [[1, 2], [1]]);|
  <mutable digraph with 2 vertices, 3 edges>
  !gapprompt@gap>| !gapinput@DigraphRemoveLoops(D);|
  <mutable digraph with 2 vertices, 2 edges>
  !gapprompt@gap>| !gapinput@D;|
  <mutable digraph with 2 vertices, 2 edges>
\end{Verbatim}
 }

 

\subsection{\textcolor{Chapter }{DigraphRemoveAllMultipleEdges}}
\logpage{[ 3, 3, 26 ]}\nobreak
\hyperdef{L}{X7DCCD0247897A3DE}{}
{\noindent\textcolor{FuncColor}{$\triangleright$\enspace\texttt{DigraphRemoveAllMultipleEdges({\mdseries\slshape digraph})\index{DigraphRemoveAllMultipleEdges@\texttt{DigraphRemoveAllMultipleEdges}}
\label{DigraphRemoveAllMultipleEdges}
}\hfill{\scriptsize (operation)}}\\
\noindent\textcolor{FuncColor}{$\triangleright$\enspace\texttt{DigraphRemoveAllMultipleEdgesAttr({\mdseries\slshape digraph})\index{DigraphRemoveAllMultipleEdgesAttr@\texttt{DigraphRemoveAllMultipleEdgesAttr}}
\label{DigraphRemoveAllMultipleEdgesAttr}
}\hfill{\scriptsize (attribute)}}\\
\textbf{\indent Returns:\ }
A digraph.



 If \mbox{\texttt{\mdseries\slshape digraph}} is a digraph, then this operation returns a digraph constructed from \mbox{\texttt{\mdseries\slshape digraph}} by removing all multiple edges. The result is the largest subdigraph of \mbox{\texttt{\mdseries\slshape digraph}} which does not contain multiple edges.

 If \mbox{\texttt{\mdseries\slshape digraph}} is immutable, then a new immutable digraph is returned. Otherwise, the
multiple edges of the mutable digraph \mbox{\texttt{\mdseries\slshape digraph}} are removed in\texttt{\symbol{45}}place. 
\begin{Verbatim}[commandchars=!@|,fontsize=\small,frame=single,label=Example]
  !gapprompt@gap>| !gapinput@D1 := Digraph([[1, 2, 3, 2], [1, 1, 3], [2, 2, 2]]);|
  <immutable multidigraph with 3 vertices, 10 edges>
  !gapprompt@gap>| !gapinput@D2 := DigraphRemoveAllMultipleEdges(D1);|
  <immutable digraph with 3 vertices, 6 edges>
  !gapprompt@gap>| !gapinput@OutNeighbours(D2);|
  [ [ 1, 2, 3 ], [ 1, 3 ], [ 2 ] ]
  !gapprompt@gap>| !gapinput@D := Digraph(IsMutableDigraph, [[2, 2], [1]]);|
  <mutable multidigraph with 2 vertices, 3 edges>
  !gapprompt@gap>| !gapinput@DigraphRemoveAllMultipleEdges(D);|
  <mutable digraph with 2 vertices, 2 edges>
  !gapprompt@gap>| !gapinput@D;|
  <mutable digraph with 2 vertices, 2 edges>
\end{Verbatim}
 }

 

\subsection{\textcolor{Chapter }{DigraphContractEdge (for a digraph and a list of positive integers)}}
\logpage{[ 3, 3, 27 ]}\nobreak
\hyperdef{L}{X792AD1147E2BFCB7}{}
{\noindent\textcolor{FuncColor}{$\triangleright$\enspace\texttt{DigraphContractEdge({\mdseries\slshape digraph, edge})\index{DigraphContractEdge@\texttt{DigraphContractEdge}!for a digraph and a list of positive integers}
\label{DigraphContractEdge:for a digraph and a list of positive integers}
}\hfill{\scriptsize (operation)}}\\
\noindent\textcolor{FuncColor}{$\triangleright$\enspace\texttt{DigraphContractEdge({\mdseries\slshape digraph, src, ran})\index{DigraphContractEdge@\texttt{DigraphContractEdge}!for a digraph and two positive integers}
\label{DigraphContractEdge:for a digraph and two positive integers}
}\hfill{\scriptsize (operation)}}\\
\textbf{\indent Returns:\ }
A digraph.



 If \mbox{\texttt{\mdseries\slshape edge}} is a pair of vertices of \mbox{\texttt{\mdseries\slshape digraph}}, or \mbox{\texttt{\mdseries\slshape src}} and \mbox{\texttt{\mdseries\slshape ran}} are vertices of \mbox{\texttt{\mdseries\slshape digraph}}, where \mbox{\texttt{\mdseries\slshape ran}} {\textless}{\textgreater} \mbox{\texttt{\mdseries\slshape src}}, then then this operation merges the two vertices of the edge given into one.
Edges incident to \mbox{\texttt{\mdseries\slshape src}} and \mbox{\texttt{\mdseries\slshape ran}} will now be incident to \texttt{v}, the new vertex, with their direction preserved. 

 A new digraph constructed from \mbox{\texttt{\mdseries\slshape digraph}} is returned, unless \mbox{\texttt{\mdseries\slshape digraph}} belongs to \texttt{IsMutableDigraph} (\ref{IsMutableDigraph}); in this case changes are made directly to \mbox{\texttt{\mdseries\slshape digraph}}, which is then returned. The \mbox{\texttt{\mdseries\slshape digraph}} must not belong to \texttt{IsMultiDigraph} (\ref{IsMultiDigraph}). 

 The labels of any remaining edges will be preserved. 

 Assigned vertex labels for \mbox{\texttt{\mdseries\slshape src}} and \mbox{\texttt{\mdseries\slshape ran}} are combined into a list, and assigned to the new vertex \texttt{v}. 

 If an edge \mbox{\texttt{\mdseries\slshape [src, src]}} or \mbox{\texttt{\mdseries\slshape [ran, ran]}} exists, a singular edge \texttt{[v, v]} is created. If edge \mbox{\texttt{\mdseries\slshape [ran, src]}} exists, this is also removed. 

 
\begin{Verbatim}[commandchars=!@|,fontsize=\small,frame=single,label=Example]
  !gapprompt@gap>| !gapinput@D := DigraphByEdges([[1, 2], [2, 1]]);|
  <immutable digraph with 2 vertices, 2 edges>
  !gapprompt@gap>| !gapinput@D2 := DigraphContractEdge(D, 1, 2);|
  <immutable empty digraph with 1 vertex>
  !gapprompt@gap>| !gapinput@DigraphEdges(D2);|
  [  ]
  !gapprompt@gap>| !gapinput@D := DigraphByEdges(IsMutableDigraph, [[1, 2], [2, 3], [3, 4]]);|
  <mutable digraph with 4 vertices, 3 edges>
  !gapprompt@gap>| !gapinput@DigraphVertexLabels(D);;  # setting vertex labels|
  !gapprompt@gap>| !gapinput@DigraphContractEdge(D, [2, 3]);|
  <mutable digraph with 3 vertices, 2 edges>
  !gapprompt@gap>| !gapinput@DigraphEdges(D);|
  [ [ 1, 3 ], [ 3, 2 ] ]
  !gapprompt@gap>| !gapinput@DigraphVertexLabels(D);|
  [ 1, 4, [ 2, 3 ] ]
\end{Verbatim}
 }

 

\subsection{\textcolor{Chapter }{DigraphReverseEdges (for a digraph and a list of edges)}}
\logpage{[ 3, 3, 28 ]}\nobreak
\hyperdef{L}{X7821753D85402A8C}{}
{\noindent\textcolor{FuncColor}{$\triangleright$\enspace\texttt{DigraphReverseEdges({\mdseries\slshape digraph, edges})\index{DigraphReverseEdges@\texttt{DigraphReverseEdges}!for a digraph and a list of edges}
\label{DigraphReverseEdges:for a digraph and a list of edges}
}\hfill{\scriptsize (operation)}}\\
\noindent\textcolor{FuncColor}{$\triangleright$\enspace\texttt{DigraphReverseEdge({\mdseries\slshape digraph, edge})\index{DigraphReverseEdge@\texttt{DigraphReverseEdge}!for a digraph and an edge}
\label{DigraphReverseEdge:for a digraph and an edge}
}\hfill{\scriptsize (operation)}}\\
\noindent\textcolor{FuncColor}{$\triangleright$\enspace\texttt{DigraphReverseEdge({\mdseries\slshape digraph, src, ran})\index{DigraphReverseEdge@\texttt{DigraphReverseEdge}!for a digraph, source, and range}
\label{DigraphReverseEdge:for a digraph, source, and range}
}\hfill{\scriptsize (operation)}}\\
\textbf{\indent Returns:\ }
A digraph.



 If \mbox{\texttt{\mdseries\slshape digraph}} is a digraph without multiple edges, and \mbox{\texttt{\mdseries\slshape edges}} is a list of pairs of vertices of \mbox{\texttt{\mdseries\slshape digraph}} (the entries of each pair corresponding to the source and the range of an
edge, respectively), then \texttt{DigraphReverseEdges} returns a digraph constructed from \mbox{\texttt{\mdseries\slshape digraph}} by reversing the orientation of every edge specified by \mbox{\texttt{\mdseries\slshape edges}}. If only one edge is to be reversed, then \texttt{DigraphReverseEdge} can be used instead. In this case, the second argument should just be a single
vertex\texttt{\symbol{45}}pair, or the second and third arguments should be
the source and range of an edge respectively. 

 Note that even though \mbox{\texttt{\mdseries\slshape digraph}} cannot have multiple edges, the output may have multiple edges. 

 If \mbox{\texttt{\mdseries\slshape digraph}} belongs to \texttt{IsMutableDigraph} (\ref{IsMutableDigraph}), then the edges are reversed in \mbox{\texttt{\mdseries\slshape digraph}}. If \mbox{\texttt{\mdseries\slshape digraph}} belongs to \texttt{IsImmutableDigraph} (\ref{IsImmutableDigraph}), an immutable copy of \mbox{\texttt{\mdseries\slshape digraph}} with the specified edges reversed is returned.

 
\begin{Verbatim}[commandchars=!|A,fontsize=\small,frame=single,label=Example]
  !gapprompt|gap>A !gapinput|D := DigraphFromDiSparse6String(".Tg?i@s?t_e?_qEsC");A
  <immutable digraph with 21 vertices, 8 edges>
  !gapprompt|gap>A !gapinput|DigraphEdges(D);A
  [ [ 1, 2 ], [ 1, 7 ], [ 1, 8 ], [ 5, 21 ], [ 7, 19 ], [ 9, 1 ], 
    [ 11, 2 ], [ 21, 1 ] ]
  !gapprompt|gap>A !gapinput|new := DigraphReverseEdge(D, [7, 19]);A
  <immutable digraph with 21 vertices, 8 edges>
  !gapprompt|gap>A !gapinput|DigraphEdges(new);A
  [ [ 1, 2 ], [ 1, 7 ], [ 1, 8 ], [ 5, 21 ], [ 9, 1 ], [ 11, 2 ], 
    [ 19, 7 ], [ 21, 1 ] ]
  !gapprompt|gap>A !gapinput|D2 := DigraphMutableCopy(new);;A
  !gapprompt|gap>A !gapinput|new := DigraphReverseEdges(D2, [[19, 7]]);;A
  !gapprompt|gap>A !gapinput|D2 = new;A
  true
  !gapprompt|gap>A !gapinput|D = new;A
  true
\end{Verbatim}
 }

 

\subsection{\textcolor{Chapter }{DigraphDisjointUnion (for an
    arbitrary number of digraphs)}}
\logpage{[ 3, 3, 29 ]}\nobreak
\hyperdef{L}{X814F1DFC83DB273F}{}
{\noindent\textcolor{FuncColor}{$\triangleright$\enspace\texttt{DigraphDisjointUnion({\mdseries\slshape D1, D2, ...})\index{DigraphDisjointUnion@\texttt{DigraphDisjointUnion}!for an
    arbitrary number of digraphs}
\label{DigraphDisjointUnion:for an
    arbitrary number of digraphs}
}\hfill{\scriptsize (function)}}\\
\noindent\textcolor{FuncColor}{$\triangleright$\enspace\texttt{DigraphDisjointUnion({\mdseries\slshape list})\index{DigraphDisjointUnion@\texttt{DigraphDisjointUnion}!for a list of digraphs}
\label{DigraphDisjointUnion:for a list of digraphs}
}\hfill{\scriptsize (function)}}\\
\textbf{\indent Returns:\ }
A digraph.



 In the first form, if \mbox{\texttt{\mdseries\slshape D1}}, \mbox{\texttt{\mdseries\slshape D2}}, etc. are digraphs, then \texttt{DigraphDisjointUnion} returns their disjoint union. In the second form, if \mbox{\texttt{\mdseries\slshape list}} is a non\texttt{\symbol{45}}empty list of digraphs, then \texttt{DigraphDisjointUnion} returns the disjoint union of the digraphs contained in the list. 

 For a disjoint union of digraphs, the vertex set is the disjoint union of the
vertex sets, and the edge list is the disjoint union of the edge lists. 

 More specifically, for a collection of digraphs \mbox{\texttt{\mdseries\slshape D1}}, \mbox{\texttt{\mdseries\slshape D2}}, \texttt{...}, the disjoint union with have \texttt{DigraphNrVertices(}\mbox{\texttt{\mdseries\slshape D1}}\texttt{)} \texttt{+} \texttt{DigraphNrVertices(}\mbox{\texttt{\mdseries\slshape D2}}\texttt{)} \texttt{+} \texttt{...} vertices. The edges of \mbox{\texttt{\mdseries\slshape D1}} will remain unchanged, whilst the edges of the \texttt{i}th digraph, \mbox{\texttt{\mdseries\slshape D}}\texttt{[i]}, will be changed so that they belong to the vertices of the disjoint union
corresponding to \mbox{\texttt{\mdseries\slshape D}}\texttt{[i]}. In particular, the edges of \mbox{\texttt{\mdseries\slshape D}}\texttt{[i]} will have their source and range increased by \texttt{DigraphNrVertices(}\mbox{\texttt{\mdseries\slshape D1}}\texttt{)} \texttt{+} \texttt{...} \texttt{+} \texttt{DigraphNrVertices(}\mbox{\texttt{\mdseries\slshape D}}\texttt{[i\texttt{\symbol{45}}1])}. 

 Note that previously set \texttt{DigraphVertexLabels} (\ref{DigraphVertexLabels}) will be lost.

 If the first digraph \mbox{\texttt{\mdseries\slshape D1}} [\mbox{\texttt{\mdseries\slshape list[1]}}] belongs to \texttt{IsMutableDigraph} (\ref{IsMutableDigraph}), then \mbox{\texttt{\mdseries\slshape D1}} [\mbox{\texttt{\mdseries\slshape list[1]}}] is modified in place to contain the appropriate vertices and edges. If \mbox{\texttt{\mdseries\slshape digraph}} belongs to \texttt{IsImmutableDigraph} (\ref{IsImmutableDigraph}), a new immutable digraph containing the appropriate vertices and edges is
returned. 
\begin{Verbatim}[commandchars=!@|,fontsize=\small,frame=single,label=Example]
  !gapprompt@gap>| !gapinput@D1 := CycleDigraph(3);|
  <immutable cycle digraph with 3 vertices>
  !gapprompt@gap>| !gapinput@OutNeighbours(D1);|
  [ [ 2 ], [ 3 ], [ 1 ] ]
  !gapprompt@gap>| !gapinput@D2 := CompleteDigraph(3);|
  <immutable complete digraph with 3 vertices>
  !gapprompt@gap>| !gapinput@OutNeighbours(D2);|
  [ [ 2, 3 ], [ 1, 3 ], [ 1, 2 ] ]
  !gapprompt@gap>| !gapinput@union := DigraphDisjointUnion(D1, D2);|
  <immutable digraph with 6 vertices, 9 edges>
  !gapprompt@gap>| !gapinput@OutNeighbours(union);|
  [ [ 2 ], [ 3 ], [ 1 ], [ 5, 6 ], [ 4, 6 ], [ 4, 5 ] ]
\end{Verbatim}
 }

 

\subsection{\textcolor{Chapter }{DigraphEdgeUnion (for a positive number of digraphs)}}
\logpage{[ 3, 3, 30 ]}\nobreak
\hyperdef{L}{X7DA997697D310E44}{}
{\noindent\textcolor{FuncColor}{$\triangleright$\enspace\texttt{DigraphEdgeUnion({\mdseries\slshape D1, D2, ...})\index{DigraphEdgeUnion@\texttt{DigraphEdgeUnion}!for a positive number of digraphs}
\label{DigraphEdgeUnion:for a positive number of digraphs}
}\hfill{\scriptsize (function)}}\\
\noindent\textcolor{FuncColor}{$\triangleright$\enspace\texttt{DigraphEdgeUnion({\mdseries\slshape list})\index{DigraphEdgeUnion@\texttt{DigraphEdgeUnion}!for a list of digraphs}
\label{DigraphEdgeUnion:for a list of digraphs}
}\hfill{\scriptsize (function)}}\\
\textbf{\indent Returns:\ }
A digraph.



 In the first form, if \mbox{\texttt{\mdseries\slshape D1}}, \mbox{\texttt{\mdseries\slshape D2}}, etc. are digraphs, then \texttt{DigraphEdgeUnion} returns their edge union. In the second form, if \mbox{\texttt{\mdseries\slshape list}} is a non\texttt{\symbol{45}}empty list of digraphs, then \texttt{DigraphEdgeUnion} returns the edge union of the digraphs contained in the list. 

 The vertex set of the edge union of a collection of digraphs is the \emph{union} of the vertex sets, whilst the edge list of the edge union is the \emph{concatenation} of the edge lists. The number of vertices of the edge union is equal to the \emph{maximum} number of vertices of one of the digraphs, whilst the number of edges of the
edge union will equal the \emph{sum} of the number of edges of each digraph. 

 Note that previously set \texttt{DigraphVertexLabels} (\ref{DigraphVertexLabels}) will be lost. 

 If the first digraph \mbox{\texttt{\mdseries\slshape D1}} [\mbox{\texttt{\mdseries\slshape list[1]}}] belongs to \texttt{IsMutableDigraph} (\ref{IsMutableDigraph}), then \mbox{\texttt{\mdseries\slshape D1}} [\mbox{\texttt{\mdseries\slshape list[1]}}] is modified in place to contain the appropriate vertices and edges. If \mbox{\texttt{\mdseries\slshape digraph}} belongs to \texttt{IsImmutableDigraph} (\ref{IsImmutableDigraph}), a new immutable digraph containing the appropriate vertices and edges is
returned. 
\begin{Verbatim}[commandchars=!@|,fontsize=\small,frame=single,label=Example]
  !gapprompt@gap>| !gapinput@D := CycleDigraph(10);|
  <immutable cycle digraph with 10 vertices>
  !gapprompt@gap>| !gapinput@DigraphEdgeUnion(D, D);|
  <immutable multidigraph with 10 vertices, 20 edges>
  !gapprompt@gap>| !gapinput@D1 := Digraph([[2], [1]]);|
  <immutable digraph with 2 vertices, 2 edges>
  !gapprompt@gap>| !gapinput@D2 := Digraph([[2, 3], [2], [1]]);|
  <immutable digraph with 3 vertices, 4 edges>
  !gapprompt@gap>| !gapinput@union := DigraphEdgeUnion(D1, D2);|
  <immutable multidigraph with 3 vertices, 6 edges>
  !gapprompt@gap>| !gapinput@OutNeighbours(union);|
  [ [ 2, 2, 3 ], [ 1, 2 ], [ 1 ] ]
  !gapprompt@gap>| !gapinput@union = DigraphByEdges(|
  !gapprompt@>| !gapinput@Concatenation(DigraphEdges(D1), DigraphEdges(D2)));|
  true
\end{Verbatim}
 }

 

\subsection{\textcolor{Chapter }{DigraphJoin (for a positive number of digraphs)}}
\logpage{[ 3, 3, 31 ]}\nobreak
\hyperdef{L}{X7DDFC759860E3390}{}
{\noindent\textcolor{FuncColor}{$\triangleright$\enspace\texttt{DigraphJoin({\mdseries\slshape D1, D2, ...})\index{DigraphJoin@\texttt{DigraphJoin}!for a positive number of digraphs}
\label{DigraphJoin:for a positive number of digraphs}
}\hfill{\scriptsize (function)}}\\
\noindent\textcolor{FuncColor}{$\triangleright$\enspace\texttt{DigraphJoin({\mdseries\slshape list})\index{DigraphJoin@\texttt{DigraphJoin}!for a list of digraphs}
\label{DigraphJoin:for a list of digraphs}
}\hfill{\scriptsize (function)}}\\
\textbf{\indent Returns:\ }
A digraph.



 In the first form, if \mbox{\texttt{\mdseries\slshape D1}}, \mbox{\texttt{\mdseries\slshape D2}}, etc. are digraphs, then \texttt{DigraphJoin} returns their join. In the second form, if \mbox{\texttt{\mdseries\slshape list}} is a non\texttt{\symbol{45}}empty list of digraphs, then \texttt{DigraphJoin} returns the join of the digraphs contained in the list. 

 The join of a collection of digraphs \mbox{\texttt{\mdseries\slshape D1}}, \mbox{\texttt{\mdseries\slshape D2}}, \texttt{...} is formed by first taking the \texttt{DigraphDisjointUnion} (\ref{DigraphDisjointUnion:for a list of digraphs}) of the collection. In the disjoint union, if $i \neq j$ then there are no edges between vertices corresponding to digraphs \mbox{\texttt{\mdseries\slshape D}}\texttt{[i]} and \mbox{\texttt{\mdseries\slshape D}}\texttt{[j]} in the collection; the join is created by including all such edges.

 For example, the join of two empty digraphs is a complete bipartite digraph.

 Note that previously set \texttt{DigraphVertexLabels} (\ref{DigraphVertexLabels}) will be lost. 

 If the first digraph \mbox{\texttt{\mdseries\slshape D1}} [\mbox{\texttt{\mdseries\slshape list[1]}}] belongs to \texttt{IsMutableDigraph} (\ref{IsMutableDigraph}), then \mbox{\texttt{\mdseries\slshape D1}} [\mbox{\texttt{\mdseries\slshape list[1]}}] is modified in place to contain the appropriate vertices and edges. If \mbox{\texttt{\mdseries\slshape digraph}} belongs to \texttt{IsImmutableDigraph} (\ref{IsImmutableDigraph}), a new immutable digraph containing the appropriate vertices and edges is
returned. 
\begin{Verbatim}[commandchars=!@|,fontsize=\small,frame=single,label=Example]
  !gapprompt@gap>| !gapinput@D := CompleteDigraph(3);|
  <immutable complete digraph with 3 vertices>
  !gapprompt@gap>| !gapinput@IsCompleteDigraph(DigraphJoin(D, D));|
  true
  !gapprompt@gap>| !gapinput@D2 := CycleDigraph(3);|
  <immutable cycle digraph with 3 vertices>
  !gapprompt@gap>| !gapinput@DigraphJoin(D, D2);|
  <immutable digraph with 6 vertices, 27 edges>
\end{Verbatim}
 }

 

\subsection{\textcolor{Chapter }{DigraphCartesianProduct (for a positive number of digraphs)}}
\logpage{[ 3, 3, 32 ]}\nobreak
\hyperdef{L}{X7D625CC87DBFFDED}{}
{\noindent\textcolor{FuncColor}{$\triangleright$\enspace\texttt{DigraphCartesianProduct({\mdseries\slshape gr1, gr2, ...})\index{DigraphCartesianProduct@\texttt{DigraphCartesianProduct}!for a positive number of digraphs}
\label{DigraphCartesianProduct:for a positive number of digraphs}
}\hfill{\scriptsize (function)}}\\
\noindent\textcolor{FuncColor}{$\triangleright$\enspace\texttt{DigraphCartesianProduct({\mdseries\slshape list})\index{DigraphCartesianProduct@\texttt{DigraphCartesianProduct}!for a list of digraphs}
\label{DigraphCartesianProduct:for a list of digraphs}
}\hfill{\scriptsize (function)}}\\
\textbf{\indent Returns:\ }
A digraph.



 In the first form, if \mbox{\texttt{\mdseries\slshape gr1}}, \mbox{\texttt{\mdseries\slshape gr2}}, etc. are digraphs, then \texttt{DigraphCartesianProduct} returns a digraph isomorphic to their Cartesian product. 

 In the second form, if \mbox{\texttt{\mdseries\slshape list}} is a non\texttt{\symbol{45}}empty list of digraphs, then \texttt{DigraphCartesianProduct} returns a digraph isomorphic to the Cartesian product of the digraphs
contained in the list. 

 Mathematically, the Cartesian product of two digraphs \texttt{G}, \texttt{H} is a digraph with vertex set \texttt{Cartesian(DigraphVertices(G), DigraphVertices(H))} such that there is an edge from \texttt{[u, u']} to \texttt{[v, v']} iff \texttt{ u = v } and there is an edge from \texttt{u'} to \texttt{v'} in \texttt{H} or \texttt{ u' = v'} and there is an edge from \texttt{u} to \texttt{v} in \texttt{G}. 

 Due to technical reasons, the digraph \texttt{D} returned by \texttt{DigraphCartesianProduct} has vertex set \texttt{[1 .. DigraphNrVertices(G)*DigraphNrVertices(H)]} instead, and the exact method of encoding pairs of vertices into integers is
implementation specific. The original vertex pair can be somewhat regained by
using \texttt{DigraphCartesianProductProjections} (\ref{DigraphCartesianProductProjections}). In addition, \texttt{DigraphVertexLabels} (\ref{DigraphVertexLabels}) are preserved: if vertex pair \texttt{[u,u']} was encoded as \texttt{i} then the vertex label of \texttt{i} will be the pair of vertex labels of \texttt{u} and \texttt{u'} i.e. \texttt{DigraphVertexLabel(D,i) = [DigraphVertexLabel(G,u), DigraphVertexLabel(H,u')]}. 

 As the Cartesian product is associative, the Cartesian product of a collection
of digraphs \mbox{\texttt{\mdseries\slshape gr1}}, \mbox{\texttt{\mdseries\slshape gr2}}, \texttt{...} is computed in the obvious fashion. 

 
\begin{Verbatim}[commandchars=!@|,fontsize=\small,frame=single,label=Example]
  !gapprompt@gap>| !gapinput@gr := ChainDigraph(4);|
  <immutable chain digraph with 4 vertices>
  !gapprompt@gap>| !gapinput@gr2 := CycleDigraph(3);|
  <immutable cycle digraph with 3 vertices>
  !gapprompt@gap>| !gapinput@gr3 := DigraphCartesianProduct(gr, gr2);|
  <immutable digraph with 12 vertices, 21 edges>
  !gapprompt@gap>| !gapinput@IsIsomorphicDigraph(gr3, |
  !gapprompt@>| !gapinput@Digraph([[2, 5], [3, 6], [4, 7], [8], |
  !gapprompt@>| !gapinput@         [6, 9], [7, 10], [8, 11], [12],|
  !gapprompt@>| !gapinput@         [10, 1], [11, 2], [12, 3], [4]]));|
  true
\end{Verbatim}
 }

 

\subsection{\textcolor{Chapter }{DigraphDirectProduct (for a positive number of digraphs)}}
\logpage{[ 3, 3, 33 ]}\nobreak
\hyperdef{L}{X84D24DC9833B54A5}{}
{\noindent\textcolor{FuncColor}{$\triangleright$\enspace\texttt{DigraphDirectProduct({\mdseries\slshape gr1, gr2, ...})\index{DigraphDirectProduct@\texttt{DigraphDirectProduct}!for a positive number of digraphs}
\label{DigraphDirectProduct:for a positive number of digraphs}
}\hfill{\scriptsize (function)}}\\
\noindent\textcolor{FuncColor}{$\triangleright$\enspace\texttt{DigraphDirectProduct({\mdseries\slshape list})\index{DigraphDirectProduct@\texttt{DigraphDirectProduct}!for a list of digraphs}
\label{DigraphDirectProduct:for a list of digraphs}
}\hfill{\scriptsize (function)}}\\
\textbf{\indent Returns:\ }
A digraph.



 In the first form, if \mbox{\texttt{\mdseries\slshape gr1}}, \mbox{\texttt{\mdseries\slshape gr2}}, etc. are digraphs, then \texttt{DigraphDirectProduct} returns a digraph isomorphic to their direct product. 

 In the second form, if \mbox{\texttt{\mdseries\slshape list}} is a non\texttt{\symbol{45}}empty list of digraphs, then \texttt{DigraphDirectProduct} returns a digraph isomorphic to the direct product of the digraphs contained
in the list. 

 Mathematically, the direct product of two digraphs \texttt{G}, \texttt{H} is a digraph with vertex set \texttt{Cartesian(DigraphVertices(G), DigraphVertices(H))} such that there is an edge from \texttt{[u, u']} to \texttt{[v, v']} iff there is an edge from \texttt{u} to \texttt{v} in \texttt{G} and an edge from \texttt{u'} to \texttt{v'} in \texttt{H}. 

 Due to technical reasons, the digraph \texttt{D} returned by \texttt{DigraphDirectProduct} has vertex set \texttt{[1 .. DigraphNrVertices(G)*DigraphNrVertices(H)]} instead, and the exact method of encoding pairs of vertices into integers is
implementation specific. The original vertex pair can be somewhat regained by
using \texttt{DigraphDirectProductProjections} (\ref{DigraphDirectProductProjections}). In addition \texttt{DigraphVertexLabels} (\ref{DigraphVertexLabels}) are preserved: if vertex pair \texttt{[u,u']} was encoded as \texttt{i} then the vertex label of \texttt{i} will be the pair of vertex labels of \texttt{u} and \texttt{u'} i.e. \texttt{DigraphVertexLabel(D,i) = [DigraphVertexLabel(G,u), DigraphVertexLabel(H,u')]}. 

 As the direct product is associative, the direct product of a collection of
digraphs \mbox{\texttt{\mdseries\slshape gr1}}, \mbox{\texttt{\mdseries\slshape gr2}}, \texttt{...} is computed in the obvious fashion. 

 
\begin{Verbatim}[commandchars=!@|,fontsize=\small,frame=single,label=Example]
  !gapprompt@gap>| !gapinput@gr := ChainDigraph(4);|
  <immutable chain digraph with 4 vertices>
  !gapprompt@gap>| !gapinput@gr2 := CycleDigraph(3);|
  <immutable cycle digraph with 3 vertices>
  !gapprompt@gap>| !gapinput@gr3 := DigraphDirectProduct(gr, gr2);|
  <immutable digraph with 12 vertices, 9 edges>
  !gapprompt@gap>| !gapinput@IsIsomorphicDigraph(gr3, |
  !gapprompt@>| !gapinput@Digraph([[6], [7], [8], [], |
  !gapprompt@>| !gapinput@         [10], [11], [12], [],|
  !gapprompt@>| !gapinput@         [2], [3], [4], []]));|
  true
\end{Verbatim}
 }

 

\subsection{\textcolor{Chapter }{ConormalProduct}}
\logpage{[ 3, 3, 34 ]}\nobreak
\hyperdef{L}{X8160BCC378AF000F}{}
{\noindent\textcolor{FuncColor}{$\triangleright$\enspace\texttt{ConormalProduct({\mdseries\slshape D1, D2})\index{ConormalProduct@\texttt{ConormalProduct}}
\label{ConormalProduct}
}\hfill{\scriptsize (operation)}}\\
\textbf{\indent Returns:\ }
A digraph.



 If \mbox{\texttt{\mdseries\slshape D1}} and \mbox{\texttt{\mdseries\slshape D2}} are digraphs without multiple edges, then \texttt{ConormalProduct} calculates the \emph{co\texttt{\symbol{45}}normal product digraph} (\texttt{CNPD}) of \mbox{\texttt{\mdseries\slshape D1}} and \mbox{\texttt{\mdseries\slshape D2}}. \texttt{CNPD} has vertex set \texttt{V1 x V2} where \texttt{V1} is the vertex set of \mbox{\texttt{\mdseries\slshape D1}} and \texttt{V2} is the vertex set of \mbox{\texttt{\mdseries\slshape D2}} (a vertex \texttt{[a, b]} has label \texttt{(a \texttt{\symbol{45}} 1) * |V2| + b} in the output). There is an edge from \texttt{[a, b]} to \texttt{[c, d]} when at least one of the following two conditions are satisfied: 
\begin{itemize}
\item  There is an edge from \texttt{a} to \texttt{c} in \mbox{\texttt{\mdseries\slshape D1}}. 
\item  There is an edge from \texttt{b} to \texttt{d} in \mbox{\texttt{\mdseries\slshape D2}}. 
\end{itemize}
 
\begin{Verbatim}[commandchars=!@|,fontsize=\small,frame=single,label=Example]
  !gapprompt@gap>| !gapinput@ConormalProduct(DigraphSymmetricClosure(CycleDigraph(7)),|
  !gapprompt@>| !gapinput@DigraphSymmetricClosure(CycleDigraph(4)));|
  <immutable digraph with 28 vertices, 504 edges>
  !gapprompt@gap>| !gapinput@ConormalProduct(NullDigraph(0), CompleteDigraph(10));|
  <immutable empty digraph with 0 vertices>
\end{Verbatim}
 }

 

\subsection{\textcolor{Chapter }{HomomorphicProduct}}
\logpage{[ 3, 3, 35 ]}\nobreak
\hyperdef{L}{X79FD2AF279F20A72}{}
{\noindent\textcolor{FuncColor}{$\triangleright$\enspace\texttt{HomomorphicProduct({\mdseries\slshape D1, D2})\index{HomomorphicProduct@\texttt{HomomorphicProduct}}
\label{HomomorphicProduct}
}\hfill{\scriptsize (operation)}}\\
\textbf{\indent Returns:\ }
A digraph.



 If \mbox{\texttt{\mdseries\slshape D1}} and \mbox{\texttt{\mdseries\slshape D2}} are digraphs without multiple edges, then \texttt{HomomorphicProduct} calculates the \emph{homomorphic product digraph} (\texttt{HPD}) of \mbox{\texttt{\mdseries\slshape D1}} and \mbox{\texttt{\mdseries\slshape D2}}. \texttt{HPD} has vertex set \texttt{V1 x V2} where \texttt{V1} is the vertex set of \mbox{\texttt{\mdseries\slshape D1}} and \texttt{V2} is the vertex set of \mbox{\texttt{\mdseries\slshape D2}} (a vertex \texttt{[a, b]} has label \texttt{(a \texttt{\symbol{45}} 1) * |V2| + b} in the output). There is an edge from \texttt{[a, b]} to \texttt{[c, d]} when at least one of the following two conditions are satisfied: 
\begin{itemize}
\item  The vertices \texttt{a} and \texttt{c} of \mbox{\texttt{\mdseries\slshape D1}} are equal. 
\item  There is an edge from \texttt{a} to \texttt{c} in \mbox{\texttt{\mdseries\slshape D1}} and no edge from \texttt{b} to \texttt{d} in \mbox{\texttt{\mdseries\slshape D2}}. 
\end{itemize}
 
\begin{Verbatim}[commandchars=!@|,fontsize=\small,frame=single,label=Example]
  !gapprompt@gap>| !gapinput@HomomorphicProduct(PetersenGraph(),|
  !gapprompt@>| !gapinput@DigraphSymmetricClosure(ChainDigraph(4)));|
  <immutable digraph with 40 vertices, 460 edges>
  !gapprompt@gap>| !gapinput@D1 := Digraph([[2], [1, 3, 4], [2, 5], [2, 5], [3, 4]]);|
  <immutable digraph with 5 vertices, 10 edges>
  !gapprompt@gap>| !gapinput@D2 := Digraph([[2], [1, 3], [2, 4], [3]]);|
  <immutable digraph with 4 vertices, 6 edges>
  !gapprompt@gap>| !gapinput@HomomorphicProduct(D1, D2);|
  <immutable digraph with 20 vertices, 180 edges>
\end{Verbatim}
 }

 

\subsection{\textcolor{Chapter }{LexicographicProduct}}
\logpage{[ 3, 3, 36 ]}\nobreak
\hyperdef{L}{X8151176882BA9901}{}
{\noindent\textcolor{FuncColor}{$\triangleright$\enspace\texttt{LexicographicProduct({\mdseries\slshape D1, D2})\index{LexicographicProduct@\texttt{LexicographicProduct}}
\label{LexicographicProduct}
}\hfill{\scriptsize (operation)}}\\
\textbf{\indent Returns:\ }
A digraph.



 If \mbox{\texttt{\mdseries\slshape D1}} and \mbox{\texttt{\mdseries\slshape D2}} are digraphs without multiple edges, then \texttt{LexicographicProduct} calculates the \emph{lexicographic product digraph} (\texttt{LPD}) of \mbox{\texttt{\mdseries\slshape D1}} and \mbox{\texttt{\mdseries\slshape D2}}. \texttt{LPD} has vertex set \texttt{V1 x V2} where \texttt{V1} is the vertex set of \mbox{\texttt{\mdseries\slshape D1}} and \texttt{V2} is the vertex set of \mbox{\texttt{\mdseries\slshape D2}} (a vertex \texttt{[a, b]} has label \texttt{(a \texttt{\symbol{45}} 1) * |V2| + b} in the output). There is an edge from \texttt{[a, b]} to \texttt{[c, d]} when at least one of the following two conditions are satisfied: 
\begin{itemize}
\item  There is an edge from \texttt{a} to \texttt{c} in \mbox{\texttt{\mdseries\slshape D1}}. 
\item  The vertices \texttt{a} and \texttt{c} of \mbox{\texttt{\mdseries\slshape D1}} are equal and there is an edge from \texttt{b} to \texttt{d} in \mbox{\texttt{\mdseries\slshape D2}}. 
\end{itemize}
 
\begin{Verbatim}[commandchars=!@|,fontsize=\small,frame=single,label=Example]
  !gapprompt@gap>| !gapinput@LexicographicProduct(DigraphSymmetricClosure(CycleDigraph(3)),|
  !gapprompt@>| !gapinput@DigraphSymmetricClosure(ChainDigraph(2)));|
  <immutable digraph with 6 vertices, 30 edges>
  !gapprompt@gap>| !gapinput@OutNeighbours(last);|
  [ [ 2, 3, 4, 5, 6 ], [ 1, 3, 4, 5, 6 ], [ 1, 2, 4, 5, 6 ], 
    [ 1, 2, 3, 5, 6 ], [ 1, 2, 3, 4, 6 ], [ 1, 2, 3, 4, 5 ] ]
  !gapprompt@gap>| !gapinput@LexicographicProduct(NullDigraph(0), CompleteDigraph(10));|
  <immutable empty digraph with 0 vertices>
\end{Verbatim}
 }

 

\subsection{\textcolor{Chapter }{ModularProduct}}
\logpage{[ 3, 3, 37 ]}\nobreak
\hyperdef{L}{X807F95057F9DF576}{}
{\noindent\textcolor{FuncColor}{$\triangleright$\enspace\texttt{ModularProduct({\mdseries\slshape D1, D2})\index{ModularProduct@\texttt{ModularProduct}}
\label{ModularProduct}
}\hfill{\scriptsize (operation)}}\\
\textbf{\indent Returns:\ }
A digraph.



 If \mbox{\texttt{\mdseries\slshape D1}} and \mbox{\texttt{\mdseries\slshape D2}} are digraphs without multiple edges, then \texttt{ModularProduct} calculates the \emph{modular product digraph} (\texttt{MPD}) of \mbox{\texttt{\mdseries\slshape D1}} and \mbox{\texttt{\mdseries\slshape D2}}. \texttt{MPD} has vertex set \texttt{V1 x V2} where \texttt{V1} is the vertex set of \mbox{\texttt{\mdseries\slshape D1}} and \texttt{V2} is the vertex set of \mbox{\texttt{\mdseries\slshape D2}} (a vertex \texttt{[a, b]} has label \texttt{(a \texttt{\symbol{45}} 1) * |V2| + b} in the output). There is an edge from \texttt{[a, b]} to \texttt{[c, d]} precisely when the following two conditions are satisfied: 
\begin{itemize}
\item  The vertices \texttt{a} and \texttt{c} of \mbox{\texttt{\mdseries\slshape D1}} are equal if and only if the vertices \texttt{b} and \texttt{d} of \mbox{\texttt{\mdseries\slshape D2}} are equal. 
\item  There is an edge from \texttt{a} to \texttt{c} in \mbox{\texttt{\mdseries\slshape D1}} if and only if there is an edge from \texttt{b} to \texttt{d} in \mbox{\texttt{\mdseries\slshape D2}}. 
\end{itemize}
 Notably, the complete (with loops) subdigraphs of \texttt{MPD} are precisely the partial isomorphisms from \mbox{\texttt{\mdseries\slshape D1}} to \mbox{\texttt{\mdseries\slshape D2}}. 
\begin{Verbatim}[commandchars=!@|,fontsize=\small,frame=single,label=Example]
  !gapprompt@gap>| !gapinput@ModularProduct(Digraph([[1], [1, 2]]), Digraph([[], [2]]));|
  <immutable digraph with 4 vertices, 4 edges>
  !gapprompt@gap>| !gapinput@OutNeighbours(last);|
  [ [ 4 ], [ 2, 3 ], [  ], [ 4 ] ]
  !gapprompt@gap>| !gapinput@ModularProduct(PetersenGraph(),|
  !gapprompt@>| !gapinput@DigraphSymmetricClosure(CycleDigraph(5)));|
  <immutable digraph with 50 vertices, 950 edges>
  !gapprompt@gap>| !gapinput@ModularProduct(NullDigraph(0), CompleteDigraph(10));|
  <immutable empty digraph with 0 vertices>
\end{Verbatim}
 }

 

\subsection{\textcolor{Chapter }{StrongProduct}}
\logpage{[ 3, 3, 38 ]}\nobreak
\hyperdef{L}{X82E200F07FEFAF27}{}
{\noindent\textcolor{FuncColor}{$\triangleright$\enspace\texttt{StrongProduct({\mdseries\slshape D1, D2})\index{StrongProduct@\texttt{StrongProduct}}
\label{StrongProduct}
}\hfill{\scriptsize (operation)}}\\
\textbf{\indent Returns:\ }
A digraph.



 If \mbox{\texttt{\mdseries\slshape D1}} and \mbox{\texttt{\mdseries\slshape D2}} are digraphs without multiple edges, then \texttt{StrongProduct} calculates the \emph{strong product digraph} (\texttt{SPD}) of \mbox{\texttt{\mdseries\slshape D1}} and \mbox{\texttt{\mdseries\slshape D2}}. \texttt{SPD} has vertex set \texttt{V1 x V2} where \texttt{V1} is the vertex set of \mbox{\texttt{\mdseries\slshape D1}} and \texttt{V2} is the vertex set of \mbox{\texttt{\mdseries\slshape D2}} (a vertex \texttt{[a, b]} has label \texttt{(a \texttt{\symbol{45}} 1) * |V2| + b} in the output). There is an edge from \texttt{[a, b]} to \texttt{[c, d]} when at least one of the following three conditions are satisfied: 
\begin{itemize}
\item  The vertices \texttt{a} and \texttt{c} of \mbox{\texttt{\mdseries\slshape D1}} are equal and there is an edge from \texttt{b} to \texttt{d} in \mbox{\texttt{\mdseries\slshape D2}}. 
\item  The vertices \texttt{b} and \texttt{d} of \mbox{\texttt{\mdseries\slshape D2}} are equal and there is an edge from \texttt{a} to \texttt{c} in \mbox{\texttt{\mdseries\slshape D1}}. 
\item  There is an edge from \texttt{a} to \texttt{c} in \mbox{\texttt{\mdseries\slshape D1}} and there is an edge from \texttt{b} to \texttt{d} in \mbox{\texttt{\mdseries\slshape D2}}. 
\end{itemize}
 The SPD of two paths of lengths \texttt{m} and \texttt{n} is also the king's graph for an \texttt{m} by \texttt{n} board. 
\begin{Verbatim}[commandchars=!@|,fontsize=\small,frame=single,label=Example]
  !gapprompt@gap>| !gapinput@D1 := Digraph([[2], [1, 3, 4], [2, 5], [2, 5], [3, 4]]);|
  <immutable digraph with 5 vertices, 10 edges>
  !gapprompt@gap>| !gapinput@D2 := Digraph([[2], [1, 3, 4], [2], [2]]);|
  <immutable digraph with 4 vertices, 6 edges>
  !gapprompt@gap>| !gapinput@StrongProduct(D1, D2);|
  <immutable digraph with 20 vertices, 130 edges>
  !gapprompt@gap>| !gapinput@StrongProduct(DigraphSymmetricClosure(ChainDigraph(3)),|
  !gapprompt@>| !gapinput@DigraphSymmetricClosure(ChainDigraph(4)));|
  <immutable digraph with 12 vertices, 58 edges>
\end{Verbatim}
 }

 

\subsection{\textcolor{Chapter }{DigraphCartesianProductProjections}}
\logpage{[ 3, 3, 39 ]}\nobreak
\hyperdef{L}{X8795087A78FE7D54}{}
{\noindent\textcolor{FuncColor}{$\triangleright$\enspace\texttt{DigraphCartesianProductProjections({\mdseries\slshape digraph})\index{DigraphCartesianProductProjections@\texttt{DigraphCartesianProductProjections}}
\label{DigraphCartesianProductProjections}
}\hfill{\scriptsize (attribute)}}\\
\textbf{\indent Returns:\ }
A list of transformations.



 If \mbox{\texttt{\mdseries\slshape digraph}} is a Cartesian product digraph, \texttt{digraph = DigraphCartesianProduct(gr{\textunderscore}1, gr{\textunderscore}2,
... )}, then \texttt{DigraphCartesianProductProjections} returns a list \texttt{proj} such that \texttt{proj[i]} is the projection onto the \texttt{i}\texttt{\symbol{45}}th coordinate of the product. 

 A projection is an idempotent endomorphism of \mbox{\texttt{\mdseries\slshape digraph}}. If \texttt{gr1, gr2, ...} are all loopless digraphs, then the induced subdigraph of \mbox{\texttt{\mdseries\slshape digraph}} on the image of \texttt{proj[i]} is isomorphic to \texttt{gr{\textunderscore}i}. 

 Currently this attribute is only set upon creating an immutable digraph via \texttt{DigraphCartesianProduct} and there is no way of calculating it for other digraphs. 

 For more information see \texttt{DigraphCartesianProduct} (\ref{DigraphCartesianProduct:for a list of digraphs}) 
\begin{Verbatim}[commandchars=!@|,fontsize=\small,frame=single,label=Example]
  !gapprompt@gap>| !gapinput@D := DigraphCartesianProduct(ChainDigraph(3), CycleDigraph(4),|
  !gapprompt@>| !gapinput@Digraph([[2], [2]]));;|
  !gapprompt@gap>| !gapinput@HasDigraphCartesianProductProjections(D);|
  true
  !gapprompt@gap>| !gapinput@proj := DigraphCartesianProductProjections(D);; Length(proj);|
  3
  !gapprompt@gap>| !gapinput@IsIdempotent(proj[2]);|
  true
  !gapprompt@gap>| !gapinput@RankOfTransformation(proj[3]);|
  2
  !gapprompt@gap>| !gapinput@S := ImageSetOfTransformation(proj[2]);;|
  !gapprompt@gap>| !gapinput@IsIsomorphicDigraph(CycleDigraph(4), InducedSubdigraph(D, S));|
  true
\end{Verbatim}
 }

 

\subsection{\textcolor{Chapter }{DigraphDirectProductProjections}}
\logpage{[ 3, 3, 40 ]}\nobreak
\hyperdef{L}{X84FB64F185B804C2}{}
{\noindent\textcolor{FuncColor}{$\triangleright$\enspace\texttt{DigraphDirectProductProjections({\mdseries\slshape digraph})\index{DigraphDirectProductProjections@\texttt{DigraphDirectProductProjections}}
\label{DigraphDirectProductProjections}
}\hfill{\scriptsize (attribute)}}\\
\textbf{\indent Returns:\ }
A list of transformations.



 If \mbox{\texttt{\mdseries\slshape digraph}} is a direct product digraph, \texttt{digraph = DigraphDirectProduct(gr{\textunderscore}1, gr{\textunderscore}2, ...
)}, then \texttt{DigraphDirectProductProjections} returns a list \texttt{proj} such that \texttt{proj[i]} is the projection onto the \texttt{i}\texttt{\symbol{45}}th coordinate of the product. 

 A projection is an idempotent endomorphism of \mbox{\texttt{\mdseries\slshape digraph}}. If \texttt{gr1, gr2, ...} are all loopless digraphs, then the image of \mbox{\texttt{\mdseries\slshape digraph}} under \texttt{proj[i]} is isomorphic to \texttt{gr{\textunderscore}i}. 

 Currently this attribute is only set upon creating an immutable digraph via \texttt{DigraphDirectProduct} and there is no way of calculating it for other digraphs. 

 For more information, see \texttt{DigraphDirectProduct} (\ref{DigraphDirectProduct:for a list of digraphs}) 
\begin{Verbatim}[commandchars=!@|,fontsize=\small,frame=single,label=Example]
  !gapprompt@gap>| !gapinput@D := DigraphDirectProduct(ChainDigraph(3), CycleDigraph(4),|
  !gapprompt@>| !gapinput@Digraph([[2], [2]]));;|
  !gapprompt@gap>| !gapinput@HasDigraphDirectProductProjections(D);|
  true
  !gapprompt@gap>| !gapinput@proj := DigraphDirectProductProjections(D);; Length(proj);|
  3
  !gapprompt@gap>| !gapinput@IsIdempotent(proj[2]);|
  true
  !gapprompt@gap>| !gapinput@RankOfTransformation(proj[3]);|
  2
  !gapprompt@gap>| !gapinput@P := DigraphRemoveAllMultipleEdges(|
  !gapprompt@>| !gapinput@ReducedDigraph(OnDigraphs(D, proj[2])));; |
  !gapprompt@gap>| !gapinput@IsIsomorphicDigraph(CycleDigraph(4), P);|
  true
\end{Verbatim}
 }

 

\subsection{\textcolor{Chapter }{LineDigraph}}
\logpage{[ 3, 3, 41 ]}\nobreak
\hyperdef{L}{X8595BF937B749F22}{}
{\noindent\textcolor{FuncColor}{$\triangleright$\enspace\texttt{LineDigraph({\mdseries\slshape digraph})\index{LineDigraph@\texttt{LineDigraph}}
\label{LineDigraph}
}\hfill{\scriptsize (operation)}}\\
\noindent\textcolor{FuncColor}{$\triangleright$\enspace\texttt{EdgeDigraph({\mdseries\slshape digraph})\index{EdgeDigraph@\texttt{EdgeDigraph}}
\label{EdgeDigraph}
}\hfill{\scriptsize (operation)}}\\
\textbf{\indent Returns:\ }
A digraph.



 Given a digraph \mbox{\texttt{\mdseries\slshape digraph}}, the operation returns the digraph obtained by associating a vertex with each
edge of \mbox{\texttt{\mdseries\slshape digraph}}, and creating an edge from a vertex \texttt{v} to a vertex \texttt{u} if and only if the terminal vertex of the edge associated with \texttt{v} is the start vertex of the edge associated with \texttt{u}.

 Note that the returned digraph is always a new immutable digraph, and the
argument \mbox{\texttt{\mdseries\slshape digraph}} is never modified. 
\begin{Verbatim}[commandchars=!@|,fontsize=\small,frame=single,label=Example]
  !gapprompt@gap>| !gapinput@LineDigraph(CompleteDigraph(3));|
  <immutable digraph with 6 vertices, 12 edges>
  !gapprompt@gap>| !gapinput@LineDigraph(ChainDigraph(3));|
  <immutable digraph with 2 vertices, 1 edge>
\end{Verbatim}
 }

 

\subsection{\textcolor{Chapter }{LineUndirectedDigraph}}
\logpage{[ 3, 3, 42 ]}\nobreak
\hyperdef{L}{X8364C6F17A1680CB}{}
{\noindent\textcolor{FuncColor}{$\triangleright$\enspace\texttt{LineUndirectedDigraph({\mdseries\slshape digraph})\index{LineUndirectedDigraph@\texttt{LineUndirectedDigraph}}
\label{LineUndirectedDigraph}
}\hfill{\scriptsize (operation)}}\\
\noindent\textcolor{FuncColor}{$\triangleright$\enspace\texttt{EdgeUndirectedDigraph({\mdseries\slshape digraph})\index{EdgeUndirectedDigraph@\texttt{EdgeUndirectedDigraph}}
\label{EdgeUndirectedDigraph}
}\hfill{\scriptsize (operation)}}\\
\textbf{\indent Returns:\ }
A digraph.



 Given a symmetric digraph \mbox{\texttt{\mdseries\slshape digraph}}, the operation returns the symmetric digraph obtained by associating a vertex
with each edge of \mbox{\texttt{\mdseries\slshape digraph}}, ignoring directions and multiplicities, and adding an edge between two
vertices if and only if the corresponding edges have a vertex in common.

 Note that the returned digraph is always a new immutable digraph, and the
argument \mbox{\texttt{\mdseries\slshape digraph}} is never modified. 
\begin{Verbatim}[commandchars=!@|,fontsize=\small,frame=single,label=Example]
  !gapprompt@gap>| !gapinput@LineUndirectedDigraph(CompleteDigraph(3));|
  <immutable digraph with 3 vertices, 6 edges>
  !gapprompt@gap>| !gapinput@LineUndirectedDigraph(DigraphSymmetricClosure(ChainDigraph(3)));|
  <immutable digraph with 2 vertices, 2 edges>
\end{Verbatim}
 }

 

\subsection{\textcolor{Chapter }{DoubleDigraph}}
\logpage{[ 3, 3, 43 ]}\nobreak
\hyperdef{L}{X7FB8B48C87C0ED16}{}
{\noindent\textcolor{FuncColor}{$\triangleright$\enspace\texttt{DoubleDigraph({\mdseries\slshape digraph})\index{DoubleDigraph@\texttt{DoubleDigraph}}
\label{DoubleDigraph}
}\hfill{\scriptsize (operation)}}\\
\textbf{\indent Returns:\ }
A digraph.



 Let \mbox{\texttt{\mdseries\slshape digraph}} be a digraph with vertex set \texttt{V}. This function returns the double digraph of \mbox{\texttt{\mdseries\slshape digraph}}. The vertex set of the double digraph is the original vertex set together
with a duplicate. The edges are \texttt{[u{\textunderscore}1, v{\textunderscore}2]} and \texttt{[u{\textunderscore}2, v{\textunderscore}1]} if and only if \texttt{[u, v]} is an edge in \mbox{\texttt{\mdseries\slshape digraph}}, together with the original edges and their duplicates.

 If \mbox{\texttt{\mdseries\slshape digraph}} is mutable, then \mbox{\texttt{\mdseries\slshape digraph}} is modified in\texttt{\symbol{45}}place. If \mbox{\texttt{\mdseries\slshape digraph}} is immutable, then a new immutable digraph constructed as described above is
returned. 
\begin{Verbatim}[commandchars=!@|,fontsize=\small,frame=single,label=Example]
  !gapprompt@gap>| !gapinput@gamma := Digraph([[2], [3], [1]]);|
  <immutable digraph with 3 vertices, 3 edges>
  !gapprompt@gap>| !gapinput@DoubleDigraph(gamma);|
  <immutable digraph with 6 vertices, 12 edges>
\end{Verbatim}
 }

 

\subsection{\textcolor{Chapter }{BipartiteDoubleDigraph}}
\logpage{[ 3, 3, 44 ]}\nobreak
\hyperdef{L}{X7C6E6CB284982C7A}{}
{\noindent\textcolor{FuncColor}{$\triangleright$\enspace\texttt{BipartiteDoubleDigraph({\mdseries\slshape digraph})\index{BipartiteDoubleDigraph@\texttt{BipartiteDoubleDigraph}}
\label{BipartiteDoubleDigraph}
}\hfill{\scriptsize (operation)}}\\
\textbf{\indent Returns:\ }
A digraph.



 Let \mbox{\texttt{\mdseries\slshape digraph}} be a digraph with vertex set \texttt{V}. This function returns the bipartite double digraph of \mbox{\texttt{\mdseries\slshape digraph}}. The vertex set of the double digraph is the original vertex set together
with a duplicate. The edges are \texttt{[u{\textunderscore}1, v{\textunderscore}2]} and \texttt{[u{\textunderscore}2, v{\textunderscore}1]} if and only if \texttt{[u, v]} is an edge in \mbox{\texttt{\mdseries\slshape digraph}}. The resulting graph is bipartite, since the original edges are not included
in the resulting digraph.

 If \mbox{\texttt{\mdseries\slshape digraph}} is mutable, then \mbox{\texttt{\mdseries\slshape digraph}} is modified in\texttt{\symbol{45}}place. If \mbox{\texttt{\mdseries\slshape digraph}} is immutable, then a new immutable digraph constructed as described above is
returned. 
\begin{Verbatim}[commandchars=!@|,fontsize=\small,frame=single,label=Example]
  !gapprompt@gap>| !gapinput@gamma := Digraph([[2], [3], [1]]);|
  <immutable digraph with 3 vertices, 3 edges>
  !gapprompt@gap>| !gapinput@BipartiteDoubleDigraph(gamma);|
  <immutable digraph with 6 vertices, 6 edges>
\end{Verbatim}
 }

 

\subsection{\textcolor{Chapter }{DigraphAddAllLoops}}
\logpage{[ 3, 3, 45 ]}\nobreak
\hyperdef{L}{X8167A50A83256ED1}{}
{\noindent\textcolor{FuncColor}{$\triangleright$\enspace\texttt{DigraphAddAllLoops({\mdseries\slshape digraph})\index{DigraphAddAllLoops@\texttt{DigraphAddAllLoops}}
\label{DigraphAddAllLoops}
}\hfill{\scriptsize (operation)}}\\
\noindent\textcolor{FuncColor}{$\triangleright$\enspace\texttt{DigraphAddAllLoopsAttr({\mdseries\slshape digraph})\index{DigraphAddAllLoopsAttr@\texttt{DigraphAddAllLoopsAttr}}
\label{DigraphAddAllLoopsAttr}
}\hfill{\scriptsize (attribute)}}\\
\textbf{\indent Returns:\ }
A digraph.



 For a digraph \mbox{\texttt{\mdseries\slshape digraph}} this operation returns a new digraph constructed from \mbox{\texttt{\mdseries\slshape digraph}}, such that a loop is added for every vertex which did not have a loop in \mbox{\texttt{\mdseries\slshape digraph}}.

 If \mbox{\texttt{\mdseries\slshape digraph}} is immutable, then a new immutable digraph is returned. Otherwise, the loops
are added to the loopless vertices of the mutable digraph \mbox{\texttt{\mdseries\slshape digraph}} in\texttt{\symbol{45}}place. 
\begin{Verbatim}[commandchars=!@|,fontsize=\small,frame=single,label=Example]
  !gapprompt@gap>| !gapinput@D := EmptyDigraph(13);|
  <immutable empty digraph with 13 vertices>
  !gapprompt@gap>| !gapinput@D := DigraphAddAllLoops(D);|
  <immutable reflexive digraph with 13 vertices, 13 edges>
  !gapprompt@gap>| !gapinput@OutNeighbours(D);|
  [ [ 1 ], [ 2 ], [ 3 ], [ 4 ], [ 5 ], [ 6 ], [ 7 ], [ 8 ], [ 9 ], 
    [ 10 ], [ 11 ], [ 12 ], [ 13 ] ]
  !gapprompt@gap>| !gapinput@D := Digraph([[1, 2, 3], [1, 3], [1]]);|
  <immutable digraph with 3 vertices, 6 edges>
  !gapprompt@gap>| !gapinput@D := DigraphAddAllLoops(D);|
  <immutable reflexive digraph with 3 vertices, 8 edges>
  !gapprompt@gap>| !gapinput@OutNeighbours(D);|
  [ [ 1, 2, 3 ], [ 1, 3, 2 ], [ 1, 3 ] ]
  !gapprompt@gap>| !gapinput@D := CycleDigraph(IsMutableDigraph, 3);|
  <mutable digraph with 3 vertices, 3 edges>
  !gapprompt@gap>| !gapinput@DigraphAddAllLoops(D);|
  <mutable digraph with 3 vertices, 6 edges>
  !gapprompt@gap>| !gapinput@D;|
  <mutable digraph with 3 vertices, 6 edges>
\end{Verbatim}
 }

 

\subsection{\textcolor{Chapter }{DistanceDigraph (for digraph and int)}}
\logpage{[ 3, 3, 46 ]}\nobreak
\hyperdef{L}{X865436437DF95FEF}{}
{\noindent\textcolor{FuncColor}{$\triangleright$\enspace\texttt{DistanceDigraph({\mdseries\slshape digraph, i})\index{DistanceDigraph@\texttt{DistanceDigraph}!for digraph and int}
\label{DistanceDigraph:for digraph and int}
}\hfill{\scriptsize (operation)}}\\
\noindent\textcolor{FuncColor}{$\triangleright$\enspace\texttt{DistanceDigraph({\mdseries\slshape digraph, list})\index{DistanceDigraph@\texttt{DistanceDigraph}!for digraph and list}
\label{DistanceDigraph:for digraph and list}
}\hfill{\scriptsize (operation)}}\\
\textbf{\indent Returns:\ }
A digraph.



 The first argument is a digraph, the second argument is a
non\texttt{\symbol{45}}negative integer or a list of positive integers. This
operation returns a digraph on the same set of vertices as \mbox{\texttt{\mdseries\slshape digraph}}, with two vertices being adjacent if and only if the distance between them in \mbox{\texttt{\mdseries\slshape digraph}} equals \mbox{\texttt{\mdseries\slshape i}} or is a number in \mbox{\texttt{\mdseries\slshape list}}. See \texttt{DigraphShortestDistance} (\ref{DigraphShortestDistance:for a digraph and two vertices}). 

 If \mbox{\texttt{\mdseries\slshape digraph}} is mutable, then \mbox{\texttt{\mdseries\slshape digraph}} is modified in\texttt{\symbol{45}}place. If \mbox{\texttt{\mdseries\slshape digraph}} is immutable, then a new immutable digraph constructed as described above is
returned. 
\begin{Verbatim}[commandchars=!|N,fontsize=\small,frame=single,label=Example]
  !gapprompt|gap>N !gapinput|digraph := DigraphFromSparse6String(N
  !gapprompt|>N !gapinput|":]n?AL`BC_DEbEF`GIaGHdIJeGKcKL_@McDHfILaBJfHMjKM");N
  <immutable symmetric digraph with 30 vertices, 90 edges>
  !gapprompt|gap>N !gapinput|DistanceDigraph(digraph, 1);N
  <immutable digraph with 30 vertices, 90 edges>
  !gapprompt|gap>N !gapinput|DistanceDigraph(digraph, [1, 2]);N
  <immutable digraph with 30 vertices, 270 edges>
\end{Verbatim}
 }

 

\subsection{\textcolor{Chapter }{DigraphClosure}}
\logpage{[ 3, 3, 47 ]}\nobreak
\hyperdef{L}{X86F9CCEA839ABC48}{}
{\noindent\textcolor{FuncColor}{$\triangleright$\enspace\texttt{DigraphClosure({\mdseries\slshape digraph, k})\index{DigraphClosure@\texttt{DigraphClosure}}
\label{DigraphClosure}
}\hfill{\scriptsize (operation)}}\\
\textbf{\indent Returns:\ }
A digraph.



 Given a symmetric loopless digraph with no multiple edges \mbox{\texttt{\mdseries\slshape digraph}}, the \emph{\mbox{\texttt{\mdseries\slshape k}}\texttt{\symbol{45}}closure of \mbox{\texttt{\mdseries\slshape digraph}}} is defined to be the unique smallest symmetric loopless digraph \texttt{C} with no multiple edges on the vertices of \mbox{\texttt{\mdseries\slshape digraph}} that contains all the edges of \mbox{\texttt{\mdseries\slshape digraph}} and satisfies the property that the sum of the degrees of every two
non\texttt{\symbol{45}}adjacenct vertices in \texttt{C} is less than \mbox{\texttt{\mdseries\slshape k}}. See \texttt{IsSymmetricDigraph} (\ref{IsSymmetricDigraph}), \texttt{DigraphHasLoops} (\ref{DigraphHasLoops}), \texttt{IsMultiDigraph} (\ref{IsMultiDigraph}), and \texttt{OutDegreeOfVertex} (\ref{OutDegreeOfVertex}). 

 The operation \texttt{DigraphClosure} returns the \mbox{\texttt{\mdseries\slshape k}}\texttt{\symbol{45}}closure of \mbox{\texttt{\mdseries\slshape digraph}}. 
\begin{Verbatim}[commandchars=!@|,fontsize=\small,frame=single,label=Example]
  !gapprompt@gap>| !gapinput@D := CompleteDigraph(6);|
  <immutable complete digraph with 6 vertices>
  !gapprompt@gap>| !gapinput@D := DigraphRemoveEdges(D, [[1, 2], [2, 1]]);|
  <immutable digraph with 6 vertices, 28 edges>
  !gapprompt@gap>| !gapinput@closure := DigraphClosure(D, 6);|
  <immutable digraph with 6 vertices, 30 edges>
  !gapprompt@gap>| !gapinput@IsCompleteDigraph(closure);|
  true
\end{Verbatim}
 }

 

\subsection{\textcolor{Chapter }{DigraphMycielskian}}
\logpage{[ 3, 3, 48 ]}\nobreak
\hyperdef{L}{X7B5AC5FE859F4D80}{}
{\noindent\textcolor{FuncColor}{$\triangleright$\enspace\texttt{DigraphMycielskian({\mdseries\slshape digraph})\index{DigraphMycielskian@\texttt{DigraphMycielskian}}
\label{DigraphMycielskian}
}\hfill{\scriptsize (operation)}}\\
\noindent\textcolor{FuncColor}{$\triangleright$\enspace\texttt{DigraphMycielskianAttr({\mdseries\slshape digraph})\index{DigraphMycielskianAttr@\texttt{DigraphMycielskianAttr}}
\label{DigraphMycielskianAttr}
}\hfill{\scriptsize (attribute)}}\\
\textbf{\indent Returns:\ }
A digraph.



 If \mbox{\texttt{\mdseries\slshape digraph}} is a symmetric digraph, then \texttt{DigraphMycielskian} returns the Mycielskian of \mbox{\texttt{\mdseries\slshape digraph}}.

 The Mycielskian of a symmetric digraph is a larger symmetric digraph
constructed from it, which has a larger chromatic number. For further
information, see \href{https://en.wikipedia.org/wiki/Mycielskian} {\texttt{https://en.wikipedia.org/wiki/Mycielskian}}.

 If \mbox{\texttt{\mdseries\slshape digraph}} is immutable, then a new immutable digraph is returned. Otherwise, the mutable
digraph \mbox{\texttt{\mdseries\slshape digraph}} is changed in\texttt{\symbol{45}}place into its Mycielskian.

 
\begin{Verbatim}[commandchars=!@|,fontsize=\small,frame=single,label=Example]
  !gapprompt@gap>| !gapinput@D := CycleDigraph(2);|
  <immutable cycle digraph with 2 vertices>
  !gapprompt@gap>| !gapinput@ChromaticNumber(D);|
  2
  !gapprompt@gap>| !gapinput@D := DigraphMycielskian(D);|
  <immutable digraph with 5 vertices, 10 edges>
  !gapprompt@gap>| !gapinput@ChromaticNumber(D);|
  3
  !gapprompt@gap>| !gapinput@D := DigraphMycielskian(D);|
  <immutable digraph with 11 vertices, 40 edges>
  !gapprompt@gap>| !gapinput@ChromaticNumber(D);|
  4
  !gapprompt@gap>| !gapinput@D := CompleteBipartiteDigraph(IsMutable, 2, 3);|
  <mutable digraph with 5 vertices, 12 edges>
  !gapprompt@gap>| !gapinput@DigraphMycielskian(D);|
  <mutable digraph with 11 vertices, 46 edges>
  !gapprompt@gap>| !gapinput@D;|
  <mutable digraph with 11 vertices, 46 edges>
\end{Verbatim}
 }

 }

 
\section{\textcolor{Chapter }{Random digraphs}}\logpage{[ 3, 4, 0 ]}
\hyperdef{L}{X85B5D9B97F5187B7}{}
{
 

\subsection{\textcolor{Chapter }{RandomDigraph}}
\logpage{[ 3, 4, 1 ]}\nobreak
\hyperdef{L}{X86CF9F66788B2A24}{}
{\noindent\textcolor{FuncColor}{$\triangleright$\enspace\texttt{RandomDigraph({\mdseries\slshape [filt, ]n[, p]})\index{RandomDigraph@\texttt{RandomDigraph}}
\label{RandomDigraph}
}\hfill{\scriptsize (operation)}}\\
\textbf{\indent Returns:\ }
A digraph.



 If the optional first argument \mbox{\texttt{\mdseries\slshape filt}} is present, then this should specify the category or representation the
digraph being created will belong to. For example, if \mbox{\texttt{\mdseries\slshape filt}} is \texttt{IsMutableDigraph} (\ref{IsMutableDigraph}), then the digraph being created will be mutable, if \mbox{\texttt{\mdseries\slshape filt}} is \texttt{IsImmutableDigraph} (\ref{IsImmutableDigraph}), then the digraph will be immutable. If the optional first argument \mbox{\texttt{\mdseries\slshape filt}} is not present, then \texttt{IsImmutableDigraph} (\ref{IsImmutableDigraph}) is used by default.

 The other implemented filters are as follows: \texttt{IsConnectedDigraph} (\ref{IsConnectedDigraph}), \texttt{IsSymmetricDigraph} (\ref{IsSymmetricDigraph}), \texttt{IsAcyclicDigraph} (\ref{IsAcyclicDigraph}), \texttt{IsEulerianDigraph} (\ref{IsEulerianDigraph}), \texttt{IsHamiltonianDigraph} (\ref{IsHamiltonianDigraph}).

 For \texttt{IsConnectedDigraph} (\ref{IsConnectedDigraph}), a random tree is first created independent of the value of \mbox{\texttt{\mdseries\slshape p}}, guaranteeing connectivity (with $\mbox{\texttt{\mdseries\slshape n}}-1$ edges), and then edges are added between the remaining pairs of vertices with
probability approximately \mbox{\texttt{\mdseries\slshape p}}.

 For \texttt{IsHamiltonianDigraph} (\ref{IsHamiltonianDigraph}), a random Hamiltonian cycle is first created independent of the value of \mbox{\texttt{\mdseries\slshape p}} (with \mbox{\texttt{\mdseries\slshape n}} edges), and then edges are added between the remaining pairs of vertices with
probability approximately \mbox{\texttt{\mdseries\slshape p}}.

 For \texttt{IsEulerianDigraph} (\ref{IsEulerianDigraph}), a random Eulerian cycle is created where \mbox{\texttt{\mdseries\slshape p}} influences how long the cycle will be. The cycle grows by randomly considering
edges that extend the cycle, and adding an edge with probability approximately \mbox{\texttt{\mdseries\slshape p}}. The cycle stops when we get back to the start vertex and have no more edges
left to consider from it that extend the cycle further (any possible edge from
the start vertex has either been added to the cycle, or rejected, leaving no
more edges to consider). Thus $\mbox{\texttt{\mdseries\slshape p}} = 1$ does not necessarily guarantee a complete digraph. Instead, it guarantees that
all edges considered up to the point where the cycle stops, are added.

 For \texttt{IsAcyclicDigraph} (\ref{IsAcyclicDigraph}) and \texttt{IsSymmetricDigraph} (\ref{IsSymmetricDigraph}), edges are added between any pairs of vertices with probability approximately \mbox{\texttt{\mdseries\slshape p}}.

 If \mbox{\texttt{\mdseries\slshape n}} is a positive integer, then this function returns a random digraph with \mbox{\texttt{\mdseries\slshape n}} vertices and without multiple edges. The result may or may not have loops. If
using \texttt{IsAcyclicDigraph} (\ref{IsAcyclicDigraph}), the resulting graph will not have any loops by definition.

 If the optional second argument \mbox{\texttt{\mdseries\slshape p}} is a float with value $0 \leq $ \mbox{\texttt{\mdseries\slshape  p }} $ \leq 1$, then an edge will exist between each pair of vertices with probability
approximately \mbox{\texttt{\mdseries\slshape p}}. If \mbox{\texttt{\mdseries\slshape p}} is not specified, then a random probability will be assumed (chosen with
uniform probability).

 
\begin{Verbatim}[commandchars=!@|,fontsize=\small,frame=single,label=Example]
  !gapprompt@gap>| !gapinput@RandomDigraph(1000);|
  <immutable digraph with 1000 vertices, 364444 edges>
  !gapprompt@gap>| !gapinput@RandomDigraph(10000, 0.023);|
  <immutable digraph with 10000 vertices, 2300438 edges>
  !gapprompt@gap>| !gapinput@RandomDigraph(IsMutableDigraph, 1000, 1 / 2);|
  <mutable digraph with 1000 vertices, 499739 edges>
  !gapprompt@gap>| !gapinput@RandomDigraph(IsConnectedDigraph, 1000, 0.75);|
  <immutable digraph with 1000 vertices, 750265 edges>
  !gapprompt@gap>| !gapinput@RandomDigraph(IsSymmetricDigraph, 1000);|
  <immutable digraph with 1000 vertices, 329690 edges>
  !gapprompt@gap>| !gapinput@RandomDigraph(IsAcyclicDigraph, 1000, 0.25);|
  <immutable digraph with 1000 vertices, 125070 edges>
  !gapprompt@gap>| !gapinput@RandomDigraph(IsHamiltonianDigraph, 1000, 0.5);|
  <immutable digraph with 1000 vertices, 500327 edges>
  !gapprompt@gap>| !gapinput@RandomDigraph(IsEulerianDigraph, 1000, 0.5);|
  <immutable digraph with 1000 vertices, 433869 edges>
\end{Verbatim}
 }

 

\subsection{\textcolor{Chapter }{RandomMultiDigraph}}
\logpage{[ 3, 4, 2 ]}\nobreak
\hyperdef{L}{X78FE275E7E77D56F}{}
{\noindent\textcolor{FuncColor}{$\triangleright$\enspace\texttt{RandomMultiDigraph({\mdseries\slshape n[, m]})\index{RandomMultiDigraph@\texttt{RandomMultiDigraph}}
\label{RandomMultiDigraph}
}\hfill{\scriptsize (operation)}}\\
\textbf{\indent Returns:\ }
A digraph.



 If \mbox{\texttt{\mdseries\slshape n}} is a positive integer, then this function returns a random digraph with \mbox{\texttt{\mdseries\slshape n}} vertices. If the optional second argument \mbox{\texttt{\mdseries\slshape m}} is a positive integer, then the digraph will have \mbox{\texttt{\mdseries\slshape m}} edges. If \mbox{\texttt{\mdseries\slshape m}} is not specified, then the number of edges will be chosen randomly (with
uniform probability) from the range \texttt{[1 .. }${n \choose 2}$\texttt{]}. 

 The method used by this function chooses each edge from the set of all
possible edges with uniform probability. No effort is made to avoid creating
multiple edges, so it is possible (but not guaranteed) that the result will
have multiple edges. The result may or may not have loops.

 
\begin{Verbatim}[commandchars=!@|,fontsize=\small,frame=single,label=Example]
  !gapprompt@gap>| !gapinput@RandomMultiDigraph(1000);|
  <immutable multidigraph with 1000 vertices, 216659 edges>
  !gapprompt@gap>| !gapinput@RandomMultiDigraph(1000, 950);|
  <immutable multidigraph with 1000 vertices, 950 edges>
\end{Verbatim}
 }

 

\subsection{\textcolor{Chapter }{RandomTournament}}
\logpage{[ 3, 4, 3 ]}\nobreak
\hyperdef{L}{X7D36B5E57F055051}{}
{\noindent\textcolor{FuncColor}{$\triangleright$\enspace\texttt{RandomTournament({\mdseries\slshape [filt, ]n})\index{RandomTournament@\texttt{RandomTournament}}
\label{RandomTournament}
}\hfill{\scriptsize (operation)}}\\
\textbf{\indent Returns:\ }
A digraph.



 If the optional first argument \mbox{\texttt{\mdseries\slshape filt}} is present, then this should specify the category or representation the
digraph being created will belong to. For example, if \mbox{\texttt{\mdseries\slshape filt}} is \texttt{IsMutableDigraph} (\ref{IsMutableDigraph}), then the digraph being created will be mutable, if \mbox{\texttt{\mdseries\slshape filt}} is \texttt{IsImmutableDigraph} (\ref{IsImmutableDigraph}), then the digraph will be immutable. If the optional first argument \mbox{\texttt{\mdseries\slshape filt}} is not present, then \texttt{IsImmutableDigraph} (\ref{IsImmutableDigraph}) is used by default.

 If \mbox{\texttt{\mdseries\slshape n}} is a non\texttt{\symbol{45}}negative integer, this function returns a random
tournament with \mbox{\texttt{\mdseries\slshape n}} vertices. See \texttt{IsTournament} (\ref{IsTournament}). 

 
\begin{Verbatim}[commandchars=!@|,fontsize=\small,frame=single,label=Example]
  !gapprompt@gap>| !gapinput@RandomTournament(10);|
  <immutable tournament with 10 vertices>
  !gapprompt@gap>| !gapinput@RandomTournament(IsMutableDigraph, 10);|
  <mutable digraph with 1000 vertices, 500601 edges>
\end{Verbatim}
 }

 

\subsection{\textcolor{Chapter }{RandomLattice}}
\logpage{[ 3, 4, 4 ]}\nobreak
\hyperdef{L}{X7A023E4787682475}{}
{\noindent\textcolor{FuncColor}{$\triangleright$\enspace\texttt{RandomLattice({\mdseries\slshape n})\index{RandomLattice@\texttt{RandomLattice}}
\label{RandomLattice}
}\hfill{\scriptsize (operation)}}\\
\textbf{\indent Returns:\ }
A digraph.



 If the optional first argument \mbox{\texttt{\mdseries\slshape filt}} is present, then this should specify the category or representation the
digraph being created will belong to. For example, if \mbox{\texttt{\mdseries\slshape filt}} is \texttt{IsMutableDigraph} (\ref{IsMutableDigraph}), then the digraph being created will be mutable, if \mbox{\texttt{\mdseries\slshape filt}} is \texttt{IsImmutableDigraph} (\ref{IsImmutableDigraph}), then the digraph will be immutable. If the optional first argument \mbox{\texttt{\mdseries\slshape filt}} is not present, then \texttt{IsImmutableDigraph} (\ref{IsImmutableDigraph}) is used by default.

 If \mbox{\texttt{\mdseries\slshape n}} is a positive integer, this function return a random lattice with \texttt{m} vertices, where it is guaranteed that \texttt{m} is between \mbox{\texttt{\mdseries\slshape n}} and \texttt{2 * \mbox{\texttt{\mdseries\slshape n}}}. See \texttt{IsLatticeDigraph} (\ref{IsLatticeDigraph}). 

 
\begin{Verbatim}[commandchars=!@|,fontsize=\small,frame=single,label=Example]
  !gapprompt@gap>| !gapinput@RandomLattice(10);|
  <immutable lattice digraph with 10 vertices, 39 edges>
  !gapprompt@gap>| !gapinput@RandomLattice(IsMutableDigraph, 10);|
  <mutable digraph with 12 vertices, 52 edges>
\end{Verbatim}
 }

 }

 
\section{\textcolor{Chapter }{Standard examples}}\logpage{[ 3, 5, 0 ]}
\hyperdef{L}{X7C76D1DC7DAF03D3}{}
{
 

\subsection{\textcolor{Chapter }{AndrasfaiGraph}}
\logpage{[ 3, 5, 1 ]}\nobreak
\hyperdef{L}{X7D5E3E337D03EDFF}{}
{\noindent\textcolor{FuncColor}{$\triangleright$\enspace\texttt{AndrasfaiGraph({\mdseries\slshape [filt, ]n})\index{AndrasfaiGraph@\texttt{AndrasfaiGraph}}
\label{AndrasfaiGraph}
}\hfill{\scriptsize (operation)}}\\
\textbf{\indent Returns:\ }
A digraph.



 If \mbox{\texttt{\mdseries\slshape n}} is an integer greater than 0, then this operation returns the \mbox{\texttt{\mdseries\slshape n}}th \emph{Andrasfai graph}. The Andrasfai graph is a circulant graph with \texttt{3\mbox{\texttt{\mdseries\slshape n}} \texttt{\symbol{45}} 1} vertices. The indices of the Andrasfai graph are given by the numbers between \texttt{1} and \texttt{3\mbox{\texttt{\mdseries\slshape n}} \texttt{\symbol{45}} 1} that are congruent to \texttt{1} mod \texttt{3} (that is, for each index $j$, vertex $i$ is adjacent to the $i + j$th and $i - j$ vertices). The graph has \texttt{6(3\mbox{\texttt{\mdseries\slshape n}} \texttt{\symbol{45}} 1)} edges. The graph is triangle free.

 As a circulant graph, the Andrasfai graph is biconnected, cyclic, Hamiltonian,
regular, and vertex transitive.

 See \href{https://mathworld.wolfram.com/AndrasfaiGraph.html} {\texttt{https://mathworld.wolfram.com/AndrasfaiGraph.html}} for further details.

 If the optional first argument \mbox{\texttt{\mdseries\slshape filt}} is present, then this should specify the category or representation the
digraph being created will belong to. For example, if \mbox{\texttt{\mdseries\slshape filt}} is \texttt{IsMutableDigraph} (\ref{IsMutableDigraph}), then the digraph being created will be mutable, if \mbox{\texttt{\mdseries\slshape filt}} is \texttt{IsImmutableDigraph} (\ref{IsImmutableDigraph}), then the digraph will be immutable. If the optional first argument \mbox{\texttt{\mdseries\slshape filt}} is not present, then \texttt{IsImmutableDigraph} (\ref{IsImmutableDigraph}) is used by default.

 
\begin{Verbatim}[commandchars=!@|,fontsize=\small,frame=single,label=Example]
  !gapprompt@gap>| !gapinput@D := AndrasfaiGraph(4);|
  <immutable Hamiltonian biconnected vertex-transitive symmetric digraph\
   with 11 vertices, 44 edges>
  !gapprompt@gap>| !gapinput@IsBiconnectedDigraph(D);|
  true
  !gapprompt@gap>| !gapinput@IsIsomorphicDigraph(D, CirculantGraph(11, [1, 4, 7, 10]));|
  true
\end{Verbatim}
 }

 

\subsection{\textcolor{Chapter }{BananaTree}}
\logpage{[ 3, 5, 2 ]}\nobreak
\hyperdef{L}{X8636C2898395B7DF}{}
{\noindent\textcolor{FuncColor}{$\triangleright$\enspace\texttt{BananaTree({\mdseries\slshape [filt, ]n, k})\index{BananaTree@\texttt{BananaTree}}
\label{BananaTree}
}\hfill{\scriptsize (operation)}}\\
\textbf{\indent Returns:\ }
A digraph



 If \mbox{\texttt{\mdseries\slshape n}} and \mbox{\texttt{\mdseries\slshape k}} are positive integers with \mbox{\texttt{\mdseries\slshape k}} greater than $1$, then this operation returns the \emph{banana tree} with parameters \mbox{\texttt{\mdseries\slshape n}} and \mbox{\texttt{\mdseries\slshape k}}, as defined below. 

 From \href{https://mathworld.wolfram.com/BananaTree.html} {\texttt{https://mathworld.wolfram.com/BananaTree.html}}: 

 ``An \texttt{(\mbox{\texttt{\mdseries\slshape n}},\mbox{\texttt{\mdseries\slshape k}})}\texttt{\symbol{45}}banana tree, as defined by Chen et al. (1997), is a graph
obtained by connecting one leaf of each of \mbox{\texttt{\mdseries\slshape n}} copies of an \mbox{\texttt{\mdseries\slshape k}}\texttt{\symbol{45}}star graph with a single root vertex that is distinct from
all the stars.'' 

 Specifically, in the resulting digraph, vertex \texttt{1} is the 'root', and for each \texttt{m} in \texttt{[1 .. \mbox{\texttt{\mdseries\slshape k}}]}, the \texttt{m}th star is on the vertices \texttt{[((m \texttt{\symbol{45}} 1) * n) + 2 .. (m * n) + 1]}, with the first of these being the 'centre' of the star, and the second being
the 'leaf' adjacent \texttt{1}. 

 See also \texttt{StarGraph} (\ref{StarGraph}) and \texttt{IsUndirectedTree} (\ref{IsUndirectedTree}). 

 If the optional first argument \mbox{\texttt{\mdseries\slshape filt}} is present, then this should specify the category or representation the
digraph being created will belong to. For example, if \mbox{\texttt{\mdseries\slshape filt}} is \texttt{IsMutableDigraph} (\ref{IsMutableDigraph}), then the digraph being created will be mutable, if \mbox{\texttt{\mdseries\slshape filt}} is \texttt{IsImmutableDigraph} (\ref{IsImmutableDigraph}), then the digraph will be immutable. If the optional first argument \mbox{\texttt{\mdseries\slshape filt}} is not present, then \texttt{IsImmutableDigraph} (\ref{IsImmutableDigraph}) is used by default.

 
\begin{Verbatim}[commandchars=!@|,fontsize=\small,frame=single,label=Example]
  !gapprompt@gap>| !gapinput@D := BananaTree(2, 4);|
  <immutable undirected tree with 9 vertices>
  !gapprompt@gap>| !gapinput@D := BananaTree(3, 3);|
  <immutable undirected tree with 10 vertices>
  !gapprompt@gap>| !gapinput@D := BananaTree(5, 2);|
  <immutable undirected tree with 11 vertices>
  !gapprompt@gap>| !gapinput@D := BananaTree(3, 4);|
  <immutable undirected tree with 13 vertices>
\end{Verbatim}
 }

 

\subsection{\textcolor{Chapter }{BinaryTree}}
\logpage{[ 3, 5, 3 ]}\nobreak
\hyperdef{L}{X7FC309427BB170D8}{}
{\noindent\textcolor{FuncColor}{$\triangleright$\enspace\texttt{BinaryTree({\mdseries\slshape [filt, ]m})\index{BinaryTree@\texttt{BinaryTree}}
\label{BinaryTree}
}\hfill{\scriptsize (operation)}}\\
\textbf{\indent Returns:\ }
A digraph.



 This function returns a binary tree of depth \mbox{\texttt{\mdseries\slshape m}} with \texttt{2 \texttt{\symbol{94}} \mbox{\texttt{\mdseries\slshape m}} \texttt{\symbol{45}} 1} vertices. All edges are directed towards the root of the tree, which is vertex \texttt{1}. 

 Note that \texttt{BinaryTree(\mbox{\texttt{\mdseries\slshape m}})} is the induced subdigraph of \texttt{BinaryTree(\mbox{\texttt{\mdseries\slshape m}}+1)} on the vertices \texttt{[1..2\texttt{\symbol{94}}(\mbox{\texttt{\mdseries\slshape m}}\texttt{\symbol{45}}1)]}. 

 If the optional first argument \mbox{\texttt{\mdseries\slshape filt}} is present, then this should specify the category or representation the
digraph being created will belong to. For example, if \mbox{\texttt{\mdseries\slshape filt}} is \texttt{IsMutableDigraph} (\ref{IsMutableDigraph}), then the digraph being created will be mutable, if \mbox{\texttt{\mdseries\slshape filt}} is \texttt{IsImmutableDigraph} (\ref{IsImmutableDigraph}), then the digraph will be immutable. If the optional first argument \mbox{\texttt{\mdseries\slshape filt}} is not present, then \texttt{IsImmutableDigraph} (\ref{IsImmutableDigraph}) is used by default.

 
\begin{Verbatim}[commandchars=!@|,fontsize=\small,frame=single,label=Example]
  !gapprompt@gap>| !gapinput@BinaryTree(1);|
  <immutable empty digraph with 1 vertex>
  !gapprompt@gap>| !gapinput@BinaryTree(8);|
  <immutable digraph with 255 vertices, 254 edges>
  !gapprompt@gap>| !gapinput@BinaryTree(IsMutableDigraph, 8);|
  <mutable digraph with 255 vertices, 254 edges>
\end{Verbatim}
 }

 

\subsection{\textcolor{Chapter }{BinomialTreeGraph}}
\logpage{[ 3, 5, 4 ]}\nobreak
\hyperdef{L}{X845C374280D6EAA4}{}
{\noindent\textcolor{FuncColor}{$\triangleright$\enspace\texttt{BinomialTreeGraph({\mdseries\slshape [filt, ]n})\index{BinomialTreeGraph@\texttt{BinomialTreeGraph}}
\label{BinomialTreeGraph}
}\hfill{\scriptsize (operation)}}\\
\textbf{\indent Returns:\ }
A digraph.



 If \mbox{\texttt{\mdseries\slshape n}} is a positive integer then this operation returns the \mbox{\texttt{\mdseries\slshape n}}th \emph{binomial tree graph}. The binomial tree graph has \mbox{\texttt{\mdseries\slshape n}} vertices and \texttt{\mbox{\texttt{\mdseries\slshape n}}\texttt{\symbol{45}}1} undirected edges. The vertices of the binomial tree graph are the numbers from
1 to \mbox{\texttt{\mdseries\slshape n}} in binary representation, with a vertex \texttt{v} having as a direct parent the vertex with binary representation the same as \texttt{v} but with the lowest 1\texttt{\symbol{45}}bit cleared. For example, the vertex $011$ has parent $010$, and the vertex $010$ has parent $000$.

 The binomial tree graph is an undirected tree, and is symmetric as a digraph.

 See \href{https://metacpan.org/pod/Graph::Maker::BinomialTree} {\texttt{https://metacpan.org/pod/Graph::Maker::BinomialTree}} for further details.

 If the optional first argument \mbox{\texttt{\mdseries\slshape filt}} is present, then this should specify the category or representation the
digraph being created will belong to. For example, if \mbox{\texttt{\mdseries\slshape filt}} is \texttt{IsMutableDigraph} (\ref{IsMutableDigraph}), then the digraph being created will be mutable, if \mbox{\texttt{\mdseries\slshape filt}} is \texttt{IsImmutableDigraph} (\ref{IsImmutableDigraph}), then the digraph will be immutable. If the optional first argument \mbox{\texttt{\mdseries\slshape filt}} is not present, then \texttt{IsImmutableDigraph} (\ref{IsImmutableDigraph}) is used by default.

 
\begin{Verbatim}[commandchars=!@|,fontsize=\small,frame=single,label=Example]
  !gapprompt@gap>| !gapinput@D := BinomialTreeGraph(9);|
  <immutable undirected tree with 9 vertices>
\end{Verbatim}
 }

 

\subsection{\textcolor{Chapter }{BishopsGraph}}
\logpage{[ 3, 5, 5 ]}\nobreak
\hyperdef{L}{X7E3240047C92733F}{}
{\noindent\textcolor{FuncColor}{$\triangleright$\enspace\texttt{BishopsGraph({\mdseries\slshape [filt, ][color, ]m, n})\index{BishopsGraph@\texttt{BishopsGraph}}
\label{BishopsGraph}
}\hfill{\scriptsize (operation)}}\\
\noindent\textcolor{FuncColor}{$\triangleright$\enspace\texttt{BishopGraph({\mdseries\slshape [filt, ][color, ]m, n})\index{BishopGraph@\texttt{BishopGraph}}
\label{BishopGraph}
}\hfill{\scriptsize (operation)}}\\
\textbf{\indent Returns:\ }
A digraph.



 If \mbox{\texttt{\mdseries\slshape m}} and \mbox{\texttt{\mdseries\slshape n}} are positive integers, then this operation returns the \emph{bishop's graph} of an \mbox{\texttt{\mdseries\slshape m}} by \mbox{\texttt{\mdseries\slshape n}} chessboard, as a symmetric digraph. 

 A bishop's graph represents all possible moves of the bishop chess piece
across a chessboard. An \mbox{\texttt{\mdseries\slshape m}} by \mbox{\texttt{\mdseries\slshape n}} chessboard is a grid of \mbox{\texttt{\mdseries\slshape m}} columns ("files") and \mbox{\texttt{\mdseries\slshape n}} rows ("ranks") that intersect in squares. Orthogonally adjacent squares are
alternately colored light and dark, with the square in the first rank and file
being dark.

 The \texttt{\mbox{\texttt{\mdseries\slshape m}} * \mbox{\texttt{\mdseries\slshape n}}} vertices of the bishop's graph can be placed onto the \texttt{\mbox{\texttt{\mdseries\slshape m}} * \mbox{\texttt{\mdseries\slshape n}}} squares of an \mbox{\texttt{\mdseries\slshape m}} by \mbox{\texttt{\mdseries\slshape n}} chessboard, such that two vertices are adjacent in the digraph if and only if
a bishop can move between the corresponding squares in a single turn. A legal
bishop's move is defined as any move which moves the bishop piece to a
diagonally adjacent square or to a square which can be reached through a
series of diagonally adjacent squares, with all of these small moves being in
the same direction.

 The chosen correspondence between vertices and chess squares is given by \texttt{DigraphVertexLabels} (\ref{DigraphVertexLabels}). In more detail, the vertices of the digraph are labelled by elements of the
Cartesian product \texttt{[1..\mbox{\texttt{\mdseries\slshape m}}] x [1..\mbox{\texttt{\mdseries\slshape n}}]}, where the first entry indexes the column (file) of the square in the
chessboard, and the second entry indexes the row (rank) of the square. (Note
that the files are traditionally indexed by the lowercase letters of the
alphabet). The vertices are sorted in ascending order, first by row (second
component) and then column (first component). If the optional second argument \mbox{\texttt{\mdseries\slshape color}} is present, then this should be one of the strings \texttt{"dark"}, \texttt{"light"}, or \texttt{"both"}. The default is \texttt{"both"}. A bishop on a light square can only move to light squares, and a bishop on a
dark square can only move to dark squares. This optional argument controls
which bishops are represented in the resulting digraph. If \mbox{\texttt{\mdseries\slshape color}} is \texttt{"both"}, then the resulting digraph will show all the vertices of an \mbox{\texttt{\mdseries\slshape m}} by \mbox{\texttt{\mdseries\slshape n}} chessboard, and will be disconnected (unless \texttt{\mbox{\texttt{\mdseries\slshape m}} = \mbox{\texttt{\mdseries\slshape n}} = 1}). Otherwise, \texttt{BishopsGraph} returns the induced subdigraph of this on the vertices that correspond to
either the dark squares or the light squares, according to the value of \mbox{\texttt{\mdseries\slshape color}}. 

 If the optional first argument \mbox{\texttt{\mdseries\slshape filt}} is present, then this should specify the category or representation the
digraph being created will belong to. For example, if \mbox{\texttt{\mdseries\slshape filt}} is \texttt{IsMutableDigraph} (\ref{IsMutableDigraph}), then the digraph being created will be mutable, if \mbox{\texttt{\mdseries\slshape filt}} is \texttt{IsImmutableDigraph} (\ref{IsImmutableDigraph}), then the digraph will be immutable. If the optional first argument \mbox{\texttt{\mdseries\slshape filt}} is not present, then \texttt{IsImmutableDigraph} (\ref{IsImmutableDigraph}) is used by default.

 
\begin{Verbatim}[commandchars=!@|,fontsize=\small,frame=single,label=Example]
  !gapprompt@gap>| !gapinput@BishopsGraph(8, 8);|
  <immutable symmetric digraph with 64 vertices, 560 edges>
  !gapprompt@gap>| !gapinput@D := BishopsGraph("dark", 3, 5);|
  <immutable connected symmetric digraph with 8 vertices, 24 edges>
  !gapprompt@gap>| !gapinput@IsConnectedDigraph(D);       |
  true
  !gapprompt@gap>| !gapinput@BishopsGraph("light", 4, 4);|
  <immutable connected symmetric digraph with 8 vertices, 28 edges>
  !gapprompt@gap>| !gapinput@D := BishopsGraph("both", 1, 5);|
  <immutable empty digraph with 5 vertices>
\end{Verbatim}
 }

 

\subsection{\textcolor{Chapter }{BondyGraph}}
\logpage{[ 3, 5, 6 ]}\nobreak
\hyperdef{L}{X81A0AC37816D287B}{}
{\noindent\textcolor{FuncColor}{$\triangleright$\enspace\texttt{BondyGraph({\mdseries\slshape [filt, ]n})\index{BondyGraph@\texttt{BondyGraph}}
\label{BondyGraph}
}\hfill{\scriptsize (operation)}}\\
\textbf{\indent Returns:\ }
A digraph.



 If \mbox{\texttt{\mdseries\slshape n}} is a non\texttt{\symbol{45}}negative integer then this operation returns the \mbox{\texttt{\mdseries\slshape n}}th \emph{Bondy graph}. The Bondy graphs are a family of hypohamiltonian graphs: a graph which is
not Hamiltonian itself but the removal of any vertex produces a Hamiltonian
graph. The Bondy graphs are the \texttt{(3(2\mbox{\texttt{\mdseries\slshape n}} + 1) + 2, 2)}\texttt{\symbol{45}}th Generalised Petersen graphs, and have \texttt{12\mbox{\texttt{\mdseries\slshape n}} + 10} vertices and \texttt{15 + 18\mbox{\texttt{\mdseries\slshape n}}} undirected edges.

 See \href{https://mathworld.wolfram.com/HypohamiltonianGraph.html} {\texttt{https://mathworld.wolfram.com/HypohamiltonianGraph.html}} for further details.

 If the optional first argument \mbox{\texttt{\mdseries\slshape filt}} is present, then this should specify the category or representation the
digraph being created will belong to. For example, if \mbox{\texttt{\mdseries\slshape filt}} is \texttt{IsMutableDigraph} (\ref{IsMutableDigraph}), then the digraph being created will be mutable, if \mbox{\texttt{\mdseries\slshape filt}} is \texttt{IsImmutableDigraph} (\ref{IsImmutableDigraph}), then the digraph will be immutable. If the optional first argument \mbox{\texttt{\mdseries\slshape filt}} is not present, then \texttt{IsImmutableDigraph} (\ref{IsImmutableDigraph}) is used by default.

 
\begin{Verbatim}[commandchars=!@|,fontsize=\small,frame=single,label=Example]
  !gapprompt@gap>| !gapinput@D := BondyGraph(1);|
  <immutable symmetric digraph with 22 vertices, 66 edges>
  !gapprompt@gap>| !gapinput@IsHamiltonianDigraph(D);|
  false
  !gapprompt@gap>| !gapinput@G := List([1 .. 22], x -> DigraphRemoveVertex(D, x));;|
  !gapprompt@gap>| !gapinput@ForAll(G, IsHamiltonianDigraph);|
  true
\end{Verbatim}
 }

 

\subsection{\textcolor{Chapter }{BookGraph}}
\logpage{[ 3, 5, 7 ]}\nobreak
\hyperdef{L}{X861F493382FA7C0B}{}
{\noindent\textcolor{FuncColor}{$\triangleright$\enspace\texttt{BookGraph({\mdseries\slshape [filt, ]m})\index{BookGraph@\texttt{BookGraph}}
\label{BookGraph}
}\hfill{\scriptsize (operation)}}\\
\textbf{\indent Returns:\ }
A digraph.



 The \emph{book graph} is the Cartesian product of a complete digraph with two vertices (as the "book
spine") and the $\mbox{\texttt{\mdseries\slshape m}}+1$ star graph (as the "pages"). For more details on the book graph please refer
to \href{https://mathworld.wolfram.com/BookGraph.html} {\texttt{https://mathworld.wolfram.com/BookGraph.html}}.

 See \texttt{DigraphCartesianProduct} (\ref{DigraphCartesianProduct:for a positive number of digraphs}), \texttt{CompleteDigraph} (\ref{CompleteDigraph}), and \texttt{StarGraph} (\ref{StarGraph}). 

 If the optional first argument \mbox{\texttt{\mdseries\slshape filt}} is present, then this should specify the category or representation the
digraph being created will belong to. For example, if \mbox{\texttt{\mdseries\slshape filt}} is \texttt{IsMutableDigraph} (\ref{IsMutableDigraph}), then the digraph being created will be mutable, if \mbox{\texttt{\mdseries\slshape filt}} is \texttt{IsImmutableDigraph} (\ref{IsImmutableDigraph}), then the digraph will be immutable. If the optional first argument \mbox{\texttt{\mdseries\slshape filt}} is not present, then \texttt{IsImmutableDigraph} (\ref{IsImmutableDigraph}) is used by default.

 
\begin{Verbatim}[commandchars=!@|,fontsize=\small,frame=single,label=Example]
  !gapprompt@gap>| !gapinput@BookGraph(1);|
  <immutable bipartite symmetric digraph with bicomponents of size 2>
  !gapprompt@gap>| !gapinput@BookGraph(2);|
  <immutable bipartite symmetric digraph with bicomponents of size 3>
  !gapprompt@gap>| !gapinput@BookGraph(IsMutable, 12);|
  <mutable digraph with 26 vertices, 74 edges>
  !gapprompt@gap>| !gapinput@BookGraph(7);|
  <immutable bipartite symmetric digraph with bicomponents of size 8>
  !gapprompt@gap>| !gapinput@IsSymmetricDigraph(BookGraph(24));|
  true
  !gapprompt@gap>| !gapinput@IsBipartiteDigraph(BookGraph(24));|
  true
\end{Verbatim}
 }

 

\subsection{\textcolor{Chapter }{BurntPancakeGraph}}
\logpage{[ 3, 5, 8 ]}\nobreak
\hyperdef{L}{X8542287E81BDB55E}{}
{\noindent\textcolor{FuncColor}{$\triangleright$\enspace\texttt{BurntPancakeGraph({\mdseries\slshape [filt, ]n})\index{BurntPancakeGraph@\texttt{BurntPancakeGraph}}
\label{BurntPancakeGraph}
}\hfill{\scriptsize (operation)}}\\
\textbf{\indent Returns:\ }
A digraph.



 If \mbox{\texttt{\mdseries\slshape n}} is a positive integer, then this operation returns the \emph{burnt pancake graph} with $n!$ vertices and $n2^{n}n!$ directed edges. The $n$th burnt pancake graph is the Cayley graph of the hyperoctahedral group acting
on \texttt{[\mbox{\texttt{\mdseries\slshape \texttt{\symbol{45}}n}} .. \texttt{\symbol{45}}1, 1 .. \mbox{\texttt{\mdseries\slshape n}}]} with respect to the generating set consisting of the ``prefix reversals'', which are defined in exactly the same way as in \texttt{PancakeGraph} (\ref{PancakeGraph}). The hyperoctahedral group consists of permutations $p$ acting on \texttt{[\mbox{\texttt{\mdseries\slshape \texttt{\symbol{45}}n}} .. \texttt{\symbol{45}}1, 1 .. \mbox{\texttt{\mdseries\slshape n}}]}, where the image of every point $i$ in \texttt{[\mbox{\texttt{\mdseries\slshape \texttt{\symbol{45}}n}} .. \texttt{\symbol{45}}1, 1 .. \mbox{\texttt{\mdseries\slshape n}}]} is equal to the negative of the image of $-i$ under $p$. GAP only works with permutations of positive integers and so \texttt{BurntPancakeGraph} returns the Cayley graph of the hyperoctahedral group acting on \texttt{[1 .. \mbox{\texttt{\mdseries\slshape 2n}}]} instead of \texttt{[\mbox{\texttt{\mdseries\slshape \texttt{\symbol{45}}n}} .. \texttt{\symbol{45}}1, 1 .. \mbox{\texttt{\mdseries\slshape n}}]}. If the optional first argument \mbox{\texttt{\mdseries\slshape filt}} is not present, then \texttt{IsImmutableDigraph} (\ref{IsImmutableDigraph}) is used by default.

 See \href{https://en.wikipedia.org/wiki/Pancake_graph} {\texttt{https://en.wikipedia.org/wiki/Pancake{\textunderscore}graph}} for further details. 
\begin{Verbatim}[commandchars=!@|,fontsize=\small,frame=single,label=Example]
  !gapprompt@gap>| !gapinput@BurntPancakeGraph(3);|
  <immutable symmetric digraph with 48 vertices, 144 edges>
  !gapprompt@gap>| !gapinput@BurntPancakeGraph(4);|
  <immutable symmetric digraph with 384 vertices, 1536 edges>
  !gapprompt@gap>| !gapinput@BurntPancakeGraph(5);|
  <immutable symmetric digraph with 3840 vertices, 19200 edges>
  !gapprompt@gap>| !gapinput@BurntPancakeGraph(IsMutableDigraph, 1);|
  <mutable digraph with 1 vertex, 1 edge>
\end{Verbatim}
 }

 

\subsection{\textcolor{Chapter }{PancakeGraph}}
\logpage{[ 3, 5, 9 ]}\nobreak
\hyperdef{L}{X84C1D3D67E3979A5}{}
{\noindent\textcolor{FuncColor}{$\triangleright$\enspace\texttt{PancakeGraph({\mdseries\slshape [filt, ]n})\index{PancakeGraph@\texttt{PancakeGraph}}
\label{PancakeGraph}
}\hfill{\scriptsize (operation)}}\\
\textbf{\indent Returns:\ }
A digraph.



 If \mbox{\texttt{\mdseries\slshape n}} is a positive integer, then this operation returns the \emph{pancake graph} with $n!$ vertices and $n!(n - 1)$ directed edges. The $n$th pancake graph is the Cayley graph of the symmetric group acting on \texttt{[1 .. \mbox{\texttt{\mdseries\slshape n}}]} with respect to the generating set consisting of the ``prefix reversals''. This generating set consists of the permutations \texttt{p2}, \texttt{p3}, ..., \texttt{p\mbox{\texttt{\mdseries\slshape n}}} where \texttt{ListPerm(pi, \mbox{\texttt{\mdseries\slshape n}})} is the concatenation of \texttt{[i, i \texttt{\symbol{45}} 1 .. 1]} and \texttt{[i + 1 .. \mbox{\texttt{\mdseries\slshape n}}]}. 

 If the optional first argument \mbox{\texttt{\mdseries\slshape filt}} is not present, then \texttt{IsImmutableDigraph} (\ref{IsImmutableDigraph}) is used by default.

 See \href{https://en.wikipedia.org/wiki/Pancake_graph} {\texttt{https://en.wikipedia.org/wiki/Pancake{\textunderscore}graph}} for further details. 
\begin{Verbatim}[commandchars=!@|,fontsize=\small,frame=single,label=Example]
  !gapprompt@gap>| !gapinput@D := PancakeGraph(5);|
  <immutable Hamiltonian symmetric digraph with 120 vertices, 480 edges>
  !gapprompt@gap>| !gapinput@DigraphUndirectedGirth(D);|
  6
  !gapprompt@gap>| !gapinput@ChromaticNumber(D);|
  3
  !gapprompt@gap>| !gapinput@IsHamiltonianDigraph(D);|
  true
  !gapprompt@gap>| !gapinput@IsCayleyDigraph(D);|
  true
  !gapprompt@gap>| !gapinput@IsVertexTransitive(D);|
  true
\end{Verbatim}
 }

 

\subsection{\textcolor{Chapter }{StackedBookGraph}}
\logpage{[ 3, 5, 10 ]}\nobreak
\hyperdef{L}{X7B5E9D857D47F5C2}{}
{\noindent\textcolor{FuncColor}{$\triangleright$\enspace\texttt{StackedBookGraph({\mdseries\slshape [filt, ]m, n})\index{StackedBookGraph@\texttt{StackedBookGraph}}
\label{StackedBookGraph}
}\hfill{\scriptsize (operation)}}\\
\textbf{\indent Returns:\ }
A digraph.



 The \emph{stacked book graph} is the Cartesian product of the symmetric closure of the chain digraph with \mbox{\texttt{\mdseries\slshape n}} vertices (as the "book spine") and the $\mbox{\texttt{\mdseries\slshape m}}+1$ star graph (as the "pages"). For more details on the stacked book graph please
refer to \href{https://mathworld.wolfram.com/StackedBookGraph.html} {\texttt{https://mathworld.wolfram.com/StackedBookGraph.html}}.

 See \texttt{DigraphCartesianProduct} (\ref{DigraphCartesianProduct:for a positive number of digraphs}), \texttt{DigraphSymmetricClosure} (\ref{DigraphSymmetricClosure}), \texttt{ChainDigraph} (\ref{ChainDigraph}), and \texttt{StarGraph} (\ref{StarGraph}). 

 If the optional first argument \mbox{\texttt{\mdseries\slshape filt}} is present, then this should specify the category or representation the
digraph being created will belong to. For example, if \mbox{\texttt{\mdseries\slshape filt}} is \texttt{IsMutableDigraph} (\ref{IsMutableDigraph}), then the digraph being created will be mutable, if \mbox{\texttt{\mdseries\slshape filt}} is \texttt{IsImmutableDigraph} (\ref{IsImmutableDigraph}), then the digraph will be immutable. If the optional first argument \mbox{\texttt{\mdseries\slshape filt}} is not present, then \texttt{IsImmutableDigraph} (\ref{IsImmutableDigraph}) is used by default.

 
\begin{Verbatim}[commandchars=!@|,fontsize=\small,frame=single,label=Example]
  !gapprompt@gap>| !gapinput@StackedBookGraph(1, 1);|
  <immutable bipartite symmetric digraph with bicomponents of size 1>
  !gapprompt@gap>| !gapinput@StackedBookGraph(1, 2);|
  <immutable bipartite symmetric digraph with bicomponents of size 2>
  !gapprompt@gap>| !gapinput@StackedBookGraph(3, 4);|
  <immutable bipartite symmetric digraph with bicomponents of size 8>
  !gapprompt@gap>| !gapinput@StackedBookGraph(IsMutable, 12, 5);|
  <mutable digraph with 65 vertices, 224 edges>
  !gapprompt@gap>| !gapinput@StackedBookGraph(5, 5);|
  <immutable bipartite symmetric digraph with bicomponent sizes 13 and 1\
  7>
  !gapprompt@gap>| !gapinput@IsSymmetricDigraph(StackedBookGraph(24, 8));|
  true
  !gapprompt@gap>| !gapinput@IsBipartiteDigraph(StackedBookGraph(24, 8));|
  true
\end{Verbatim}
 }

 

\subsection{\textcolor{Chapter }{ChainDigraph}}
\logpage{[ 3, 5, 11 ]}\nobreak
\hyperdef{L}{X870594FC866AC88E}{}
{\noindent\textcolor{FuncColor}{$\triangleright$\enspace\texttt{ChainDigraph({\mdseries\slshape [filt, ]n})\index{ChainDigraph@\texttt{ChainDigraph}}
\label{ChainDigraph}
}\hfill{\scriptsize (operation)}}\\
\textbf{\indent Returns:\ }
A digraph.



 If \mbox{\texttt{\mdseries\slshape n}} is a positive integer, this function returns a chain with \mbox{\texttt{\mdseries\slshape n}} vertices and \texttt{\mbox{\texttt{\mdseries\slshape n}} \texttt{\symbol{45}} 1} edges. Specifically, for each vertex \texttt{i} (with \texttt{i} {\textless} \texttt{n}), there is a directed edge with source \texttt{i} and range \texttt{i + 1}. 

 The \texttt{DigraphReflexiveTransitiveClosure} (\ref{DigraphReflexiveTransitiveClosure}) of a chain represents a total order. 

 If the optional first argument \mbox{\texttt{\mdseries\slshape filt}} is present, then this should specify the category or representation the
digraph being created will belong to. For example, if \mbox{\texttt{\mdseries\slshape filt}} is \texttt{IsMutableDigraph} (\ref{IsMutableDigraph}), then the digraph being created will be mutable, if \mbox{\texttt{\mdseries\slshape filt}} is \texttt{IsImmutableDigraph} (\ref{IsImmutableDigraph}), then the digraph will be immutable. If the optional first argument \mbox{\texttt{\mdseries\slshape filt}} is not present, then \texttt{IsImmutableDigraph} (\ref{IsImmutableDigraph}) is used by default.

 
\begin{Verbatim}[commandchars=!@|,fontsize=\small,frame=single,label=Example]
  !gapprompt@gap>| !gapinput@ChainDigraph(42);|
  <immutable chain digraph with 42 vertices>
  !gapprompt@gap>| !gapinput@ChainDigraph(IsMutableDigraph, 10);|
  <mutable digraph with 10 vertices, 9 edges>
\end{Verbatim}
 }

 

\subsection{\textcolor{Chapter }{CirculantGraph}}
\logpage{[ 3, 5, 12 ]}\nobreak
\hyperdef{L}{X7DB5AB657A797CF2}{}
{\noindent\textcolor{FuncColor}{$\triangleright$\enspace\texttt{CirculantGraph({\mdseries\slshape [filt, ]n, par})\index{CirculantGraph@\texttt{CirculantGraph}}
\label{CirculantGraph}
}\hfill{\scriptsize (operation)}}\\
\textbf{\indent Returns:\ }
A digraph.



 If \mbox{\texttt{\mdseries\slshape n}} is an integer greater than 1, and \mbox{\texttt{\mdseries\slshape par}} is a list of integers that are contained in \texttt{[1..\mbox{\texttt{\mdseries\slshape n}}]} then this operation returns a \emph{circulant graph}. The circulant graph is a graph on \mbox{\texttt{\mdseries\slshape n}} vertices, where for each element $j$ of \mbox{\texttt{\mdseries\slshape par}}, the $i$th vertex is adjacent to the $(i + j)$th and $(i - j)$th vertices.

 If \mbox{\texttt{\mdseries\slshape par}} is $[1]$, then the graph is the \mbox{\texttt{\mdseries\slshape n}}th cyclic graph. If \mbox{\texttt{\mdseries\slshape par}} is \texttt{[1,2,...,Int(\mbox{\texttt{\mdseries\slshape n}}/2)]}, then the graph is the complete graph on \mbox{\texttt{\mdseries\slshape n}} vertices. If \mbox{\texttt{\mdseries\slshape n}} is at least 4 and \mbox{\texttt{\mdseries\slshape par}} is \texttt{[1,\mbox{\texttt{\mdseries\slshape n}}]} then the graph is the \mbox{\texttt{\mdseries\slshape n}}th Mobius ladder graph.

 A circulant graph is vertex transitive, but is not necessarily connected
(consider \texttt{CirculantGraph(4, [2])}, for example). However, a \emph{connected} circulant graph is also Hamiltonian and biconnected.

 See \href{https://mathworld.wolfram.com/CirculantGraph.html} {\texttt{https://mathworld.wolfram.com/CirculantGraph.html}} for further details.

 If the optional first argument \mbox{\texttt{\mdseries\slshape filt}} is present, then this should specify the category or representation the
digraph being created will belong to. For example, if \mbox{\texttt{\mdseries\slshape filt}} is \texttt{IsMutableDigraph} (\ref{IsMutableDigraph}), then the digraph being created will be mutable, if \mbox{\texttt{\mdseries\slshape filt}} is \texttt{IsImmutableDigraph} (\ref{IsImmutableDigraph}), then the digraph will be immutable. If the optional first argument \mbox{\texttt{\mdseries\slshape filt}} is not present, then \texttt{IsImmutableDigraph} (\ref{IsImmutableDigraph}) is used by default.

 
\begin{Verbatim}[commandchars=!@|,fontsize=\small,frame=single,label=Example]
  !gapprompt@gap>| !gapinput@D := CirculantGraph(6, [2]);|
  <immutable vertex-transitive symmetric digraph with 6 vertices, 12 edg\
  es>
  !gapprompt@gap>| !gapinput@DigraphNrConnectedComponents(D);|
  2
  !gapprompt@gap>| !gapinput@D := CirculantGraph(6, [2, 3]);|
  <immutable Hamiltonian biconnected vertex-transitive symmetric digraph\
   with 6 vertices, 18 edges>
  !gapprompt@gap>| !gapinput@AutomorphismGroup(D) = DihedralGroup(IsPermGroup, 12);|
  true
  !gapprompt@gap>| !gapinput@HamiltonianPath(D);|
  [ 1, 3, 5, 2, 6, 4 ]
  !gapprompt@gap>| !gapinput@IsCompleteDigraph(CirculantGraph(6, [1, 2, 3]));|
  true
\end{Verbatim}
 }

 

\subsection{\textcolor{Chapter }{CompleteDigraph}}
\logpage{[ 3, 5, 13 ]}\nobreak
\hyperdef{L}{X812417E278198D9C}{}
{\noindent\textcolor{FuncColor}{$\triangleright$\enspace\texttt{CompleteDigraph({\mdseries\slshape [filt, ]n})\index{CompleteDigraph@\texttt{CompleteDigraph}}
\label{CompleteDigraph}
}\hfill{\scriptsize (operation)}}\\
\textbf{\indent Returns:\ }
A digraph.



 If \mbox{\texttt{\mdseries\slshape n}} is a non\texttt{\symbol{45}}negative integer, this function returns the
complete digraph with \mbox{\texttt{\mdseries\slshape n}} vertices. See \texttt{IsCompleteDigraph} (\ref{IsCompleteDigraph}). 

 If the optional first argument \mbox{\texttt{\mdseries\slshape filt}} is present, then this should specify the category or representation the
digraph being created will belong to. For example, if \mbox{\texttt{\mdseries\slshape filt}} is \texttt{IsMutableDigraph} (\ref{IsMutableDigraph}), then the digraph being created will be mutable, if \mbox{\texttt{\mdseries\slshape filt}} is \texttt{IsImmutableDigraph} (\ref{IsImmutableDigraph}), then the digraph will be immutable. If the optional first argument \mbox{\texttt{\mdseries\slshape filt}} is not present, then \texttt{IsImmutableDigraph} (\ref{IsImmutableDigraph}) is used by default.

 
\begin{Verbatim}[commandchars=!@|,fontsize=\small,frame=single,label=Example]
  !gapprompt@gap>| !gapinput@CompleteDigraph(20);|
  <immutable complete digraph with 20 vertices>
  !gapprompt@gap>| !gapinput@CompleteDigraph(IsMutableDigraph, 10);|
  <mutable digraph with 10 vertices, 90 edges>
\end{Verbatim}
 }

 

\subsection{\textcolor{Chapter }{CompleteBipartiteDigraph}}
\logpage{[ 3, 5, 14 ]}\nobreak
\hyperdef{L}{X8795B0AD856014FA}{}
{\noindent\textcolor{FuncColor}{$\triangleright$\enspace\texttt{CompleteBipartiteDigraph({\mdseries\slshape [filt, ]m, n})\index{CompleteBipartiteDigraph@\texttt{CompleteBipartiteDigraph}}
\label{CompleteBipartiteDigraph}
}\hfill{\scriptsize (operation)}}\\
\textbf{\indent Returns:\ }
A digraph.



 A complete bipartite digraph is a digraph whose vertices can be partitioned
into two non\texttt{\symbol{45}}empty vertex sets, such there exists a unique
edge with source \texttt{i} and range \texttt{j} if and only if \texttt{i} and \texttt{j} lie in different vertex sets. 

 If \mbox{\texttt{\mdseries\slshape m}} and \mbox{\texttt{\mdseries\slshape n}} are positive integers, this function returns the complete bipartite digraph
with vertex sets of sizes \mbox{\texttt{\mdseries\slshape m}} (containing the vertices \texttt{[1 .. m]}) and \mbox{\texttt{\mdseries\slshape n}} (containing the vertices \texttt{[m + 1 .. m + n]}). 

 If the optional first argument \mbox{\texttt{\mdseries\slshape filt}} is present, then this should specify the category or representation the
digraph being created will belong to. For example, if \mbox{\texttt{\mdseries\slshape filt}} is \texttt{IsMutableDigraph} (\ref{IsMutableDigraph}), then the digraph being created will be mutable, if \mbox{\texttt{\mdseries\slshape filt}} is \texttt{IsImmutableDigraph} (\ref{IsImmutableDigraph}), then the digraph will be immutable. If the optional first argument \mbox{\texttt{\mdseries\slshape filt}} is not present, then \texttt{IsImmutableDigraph} (\ref{IsImmutableDigraph}) is used by default.

 
\begin{Verbatim}[commandchars=!@|,fontsize=\small,frame=single,label=Example]
  !gapprompt@gap>| !gapinput@CompleteBipartiteDigraph(2, 3);|
  <immutable complete bipartite digraph with bicomponent sizes 2 and 3>
  !gapprompt@gap>| !gapinput@CompleteBipartiteDigraph(IsMutableDigraph, 3, 2);|
  <mutable digraph with 5 vertices, 12 edges>
\end{Verbatim}
 }

 

\subsection{\textcolor{Chapter }{CompleteMultipartiteDigraph}}
\logpage{[ 3, 5, 15 ]}\nobreak
\hyperdef{L}{X873F29CC863241F8}{}
{\noindent\textcolor{FuncColor}{$\triangleright$\enspace\texttt{CompleteMultipartiteDigraph({\mdseries\slshape [filt, ]orders})\index{CompleteMultipartiteDigraph@\texttt{CompleteMultipartiteDigraph}}
\label{CompleteMultipartiteDigraph}
}\hfill{\scriptsize (operation)}}\\
\textbf{\indent Returns:\ }
A digraph.



 For a list \mbox{\texttt{\mdseries\slshape orders}} of \texttt{n} positive integers, this function returns the digraph containing \texttt{n} independent sets of vertices of orders \texttt{[\mbox{\texttt{\mdseries\slshape l}}[1] .. \mbox{\texttt{\mdseries\slshape l}}[n]]}. Moreover, each vertex is adjacent to every other not contained in the same
independent set. 

 If the optional first argument \mbox{\texttt{\mdseries\slshape filt}} is present, then this should specify the category or representation the
digraph being created will belong to. For example, if \mbox{\texttt{\mdseries\slshape filt}} is \texttt{IsMutableDigraph} (\ref{IsMutableDigraph}), then the digraph being created will be mutable, if \mbox{\texttt{\mdseries\slshape filt}} is \texttt{IsImmutableDigraph} (\ref{IsImmutableDigraph}), then the digraph will be immutable. If the optional first argument \mbox{\texttt{\mdseries\slshape filt}} is not present, then \texttt{IsImmutableDigraph} (\ref{IsImmutableDigraph}) is used by default.

 
\begin{Verbatim}[commandchars=!@|,fontsize=\small,frame=single,label=Example]
  !gapprompt@gap>| !gapinput@CompleteMultipartiteDigraph([5, 4, 2]);|
  <immutable complete multipartite digraph with 11 vertices, 76 edges>
  !gapprompt@gap>| !gapinput@CompleteMultipartiteDigraph(IsMutableDigraph, [5, 4, 2]);|
  <mutable digraph with 11 vertices, 76 edges>
\end{Verbatim}
 }

 

\subsection{\textcolor{Chapter }{CycleDigraph}}
\logpage{[ 3, 5, 16 ]}\nobreak
\hyperdef{L}{X80C29DDE876FFBEB}{}
{\noindent\textcolor{FuncColor}{$\triangleright$\enspace\texttt{CycleDigraph({\mdseries\slshape [filt, ]n})\index{CycleDigraph@\texttt{CycleDigraph}}
\label{CycleDigraph}
}\hfill{\scriptsize (operation)}}\\
\noindent\textcolor{FuncColor}{$\triangleright$\enspace\texttt{DigraphCycle({\mdseries\slshape [filt, ]n})\index{DigraphCycle@\texttt{DigraphCycle}}
\label{DigraphCycle}
}\hfill{\scriptsize (operation)}}\\
\textbf{\indent Returns:\ }
A digraph.



 If \mbox{\texttt{\mdseries\slshape n}} is a positive integer, then these functions return a \emph{cycle digraph} with \mbox{\texttt{\mdseries\slshape n}} vertices and \mbox{\texttt{\mdseries\slshape n}} edges. Specifically, for each vertex \texttt{i} (with \texttt{i} {\textless} \texttt{n}), there is a directed edge with source \texttt{i} and range \texttt{i + 1}. In addition, there is an edge with source \texttt{n} and range \texttt{1}. 

 If the optional first argument \mbox{\texttt{\mdseries\slshape filt}} is present, then this should specify the category or representation the
digraph being created will belong to. For example, if \mbox{\texttt{\mdseries\slshape filt}} is \texttt{IsMutableDigraph} (\ref{IsMutableDigraph}), then the digraph being created will be mutable, if \mbox{\texttt{\mdseries\slshape filt}} is \texttt{IsImmutableDigraph} (\ref{IsImmutableDigraph}), then the digraph will be immutable. If the optional first argument \mbox{\texttt{\mdseries\slshape filt}} is not present, then \texttt{IsImmutableDigraph} (\ref{IsImmutableDigraph}) is used by default.

 
\begin{Verbatim}[commandchars=!@|,fontsize=\small,frame=single,label=Example]
  !gapprompt@gap>| !gapinput@CycleDigraph(1);|
  <immutable digraph with 1 vertex, 1 edge>
  !gapprompt@gap>| !gapinput@CycleDigraph(123);|
  <immutable cycle digraph with 123 vertices>
  !gapprompt@gap>| !gapinput@CycleDigraph(IsMutableDigraph, 10);|
  <mutable digraph with 10 vertices, 10 edges>
  !gapprompt@gap>| !gapinput@DigraphCycle(4) = CycleDigraph(4);|
  true
\end{Verbatim}
 }

 

\subsection{\textcolor{Chapter }{CycleGraph}}
\logpage{[ 3, 5, 17 ]}\nobreak
\hyperdef{L}{X7F6EC0AE81531C3C}{}
{\noindent\textcolor{FuncColor}{$\triangleright$\enspace\texttt{CycleGraph({\mdseries\slshape [filt, ]n})\index{CycleGraph@\texttt{CycleGraph}}
\label{CycleGraph}
}\hfill{\scriptsize (operation)}}\\
\textbf{\indent Returns:\ }
A digraph.



 If \mbox{\texttt{\mdseries\slshape n}} is an integer greater than 2 then this operation returns the \mbox{\texttt{\mdseries\slshape n}}th \emph{cycle graph}, consisting of the cycle on \mbox{\texttt{\mdseries\slshape n}} vertices. The cycle graph, unlike the cycle digraph, is symmetric. The cycle
graph has \mbox{\texttt{\mdseries\slshape n}} vertices and \mbox{\texttt{\mdseries\slshape n}} undirected edges. The cycle graph is simple so the
non\texttt{\symbol{45}}simple graphs with a single vertex and single loop and
with two vertices and two edges between them are excluded.

 See \href{https://mathworld.wolfram.com/CycleGraph.html} {\texttt{https://mathworld.wolfram.com/CycleGraph.html}} for further details.

 If the optional first argument \mbox{\texttt{\mdseries\slshape filt}} is present, then this should specify the category or representation the
digraph being created will belong to. For example, if \mbox{\texttt{\mdseries\slshape filt}} is \texttt{IsMutableDigraph} (\ref{IsMutableDigraph}), then the digraph being created will be mutable, if \mbox{\texttt{\mdseries\slshape filt}} is \texttt{IsImmutableDigraph} (\ref{IsImmutableDigraph}), then the digraph will be immutable. If the optional first argument \mbox{\texttt{\mdseries\slshape filt}} is not present, then \texttt{IsImmutableDigraph} (\ref{IsImmutableDigraph}) is used by default.

 
\begin{Verbatim}[commandchars=!@|,fontsize=\small,frame=single,label=Example]
  !gapprompt@gap>| !gapinput@D := CycleGraph(7);|
  <immutable symmetric digraph with 7 vertices, 14 edges>
\end{Verbatim}
 }

 

\subsection{\textcolor{Chapter }{EmptyDigraph}}
\logpage{[ 3, 5, 18 ]}\nobreak
\hyperdef{L}{X80DAE31A79FEFD40}{}
{\noindent\textcolor{FuncColor}{$\triangleright$\enspace\texttt{EmptyDigraph({\mdseries\slshape [filt, ]n})\index{EmptyDigraph@\texttt{EmptyDigraph}}
\label{EmptyDigraph}
}\hfill{\scriptsize (operation)}}\\
\noindent\textcolor{FuncColor}{$\triangleright$\enspace\texttt{NullDigraph({\mdseries\slshape [filt, ]n})\index{NullDigraph@\texttt{NullDigraph}}
\label{NullDigraph}
}\hfill{\scriptsize (operation)}}\\
\textbf{\indent Returns:\ }
A digraph.



 If \mbox{\texttt{\mdseries\slshape n}} is a non\texttt{\symbol{45}}negative integer, this function returns the \emph{empty} or \emph{null} digraph with \mbox{\texttt{\mdseries\slshape n}} vertices. An empty digraph is one with no edges. 

 \texttt{NullDigraph} is a synonym for \texttt{EmptyDigraph}. 

 If the optional first argument \mbox{\texttt{\mdseries\slshape filt}} is present, then this should specify the category or representation the
digraph being created will belong to. For example, if \mbox{\texttt{\mdseries\slshape filt}} is \texttt{IsMutableDigraph} (\ref{IsMutableDigraph}), then the digraph being created will be mutable, if \mbox{\texttt{\mdseries\slshape filt}} is \texttt{IsImmutableDigraph} (\ref{IsImmutableDigraph}), then the digraph will be immutable. If the optional first argument \mbox{\texttt{\mdseries\slshape filt}} is not present, then \texttt{IsImmutableDigraph} (\ref{IsImmutableDigraph}) is used by default.

 
\begin{Verbatim}[commandchars=!@|,fontsize=\small,frame=single,label=Example]
  !gapprompt@gap>| !gapinput@EmptyDigraph(20);|
  <immutable empty digraph with 20 vertices>
  !gapprompt@gap>| !gapinput@NullDigraph(10);|
  <immutable empty digraph with 10 vertices>
  !gapprompt@gap>| !gapinput@EmptyDigraph(IsMutableDigraph, 10);|
  <mutable empty digraph with 10 vertices>
\end{Verbatim}
 }

 

\subsection{\textcolor{Chapter }{GearGraph}}
\logpage{[ 3, 5, 19 ]}\nobreak
\hyperdef{L}{X7CE45E2B782ADE9A}{}
{\noindent\textcolor{FuncColor}{$\triangleright$\enspace\texttt{GearGraph({\mdseries\slshape [filt, ]n})\index{GearGraph@\texttt{GearGraph}}
\label{GearGraph}
}\hfill{\scriptsize (operation)}}\\
\textbf{\indent Returns:\ }
A digraph.



 If \mbox{\texttt{\mdseries\slshape n}} is a positive integer at least 3, then this operation returns the \emph{gear graph} with \texttt{2\mbox{\texttt{\mdseries\slshape n}} + 1} vertices and \texttt{3\mbox{\texttt{\mdseries\slshape n}}} undirected edges. The \mbox{\texttt{\mdseries\slshape n}}th gear graph is the \texttt{2\mbox{\texttt{\mdseries\slshape }}}th cycle graph with one additional central vertex, to which every other vertex
of the cycle is connected. The gear graph is a symmetric digraph. A gear graph
is a Matchstick graph, that is it is simple with a planar graph embedding, and
is a unit\texttt{\symbol{45}}distance graph (that is, it can be embedded into
the Euclidean plane with vertices being distinct points and edges having
length 1).

 See \href{https://mathworld.wolfram.com/GearGraph.html} {\texttt{https://mathworld.wolfram.com/GearGraph.html}} for further details.

 If the optional first argument \mbox{\texttt{\mdseries\slshape filt}} is present, then this should specify the category or representation the
digraph being created will belong to. For example, if \mbox{\texttt{\mdseries\slshape filt}} is \texttt{IsMutableDigraph} (\ref{IsMutableDigraph}), then the digraph being created will be mutable, if \mbox{\texttt{\mdseries\slshape filt}} is \texttt{IsImmutableDigraph} (\ref{IsImmutableDigraph}), then the digraph will be immutable. If the optional first argument \mbox{\texttt{\mdseries\slshape filt}} is not present, then \texttt{IsImmutableDigraph} (\ref{IsImmutableDigraph}) is used by default.

 
\begin{Verbatim}[commandchars=!@|,fontsize=\small,frame=single,label=Example]
  !gapprompt@gap>| !gapinput@D := GearGraph(4);|
  <immutable symmetric digraph with 9 vertices, 24 edges>
  !gapprompt@gap>| !gapinput@ChromaticNumber(D);|
  2
  !gapprompt@gap>| !gapinput@IsVertexTransitive(D);|
  false
\end{Verbatim}
 }

 

\subsection{\textcolor{Chapter }{HaarGraph}}
\logpage{[ 3, 5, 20 ]}\nobreak
\hyperdef{L}{X795B62398767E313}{}
{\noindent\textcolor{FuncColor}{$\triangleright$\enspace\texttt{HaarGraph({\mdseries\slshape [filt, ]n})\index{HaarGraph@\texttt{HaarGraph}}
\label{HaarGraph}
}\hfill{\scriptsize (operation)}}\\
\textbf{\indent Returns:\ }
A digraph.



 If \mbox{\texttt{\mdseries\slshape n}} is a positive integer then this operation returns the \emph{Haar graph} $H(\mbox{\texttt{\mdseries\slshape n}})$. 

 The number of vertices in the Haar graph $H(\mbox{\texttt{\mdseries\slshape n}})$ is equal to twice $m$, where $m$ is the number of digits required to represent \mbox{\texttt{\mdseries\slshape n}} in binary. These vertices are arranged into bicomponents \texttt{[1..m]} and \texttt{[m+1..2*m]}. Vertices $i$ and $j$ in different bicomponents are adjacent by a symmetric pair of edges if and
only if the binary representation of \mbox{\texttt{\mdseries\slshape n}} has a 1 in position ($j-i$ modulo $m$) + 1 from the left. 

 If the optional first argument \mbox{\texttt{\mdseries\slshape filt}} is present, then this should specify the category or representation the
digraph being created will belong to. For example, if \mbox{\texttt{\mdseries\slshape filt}} is \texttt{IsMutableDigraph} (\ref{IsMutableDigraph}), then the digraph being created will be mutable, if \mbox{\texttt{\mdseries\slshape filt}} is \texttt{IsImmutableDigraph} (\ref{IsImmutableDigraph}), then the digraph will be immutable. If the optional first argument \mbox{\texttt{\mdseries\slshape filt}} is not present, then \texttt{IsImmutableDigraph} (\ref{IsImmutableDigraph}) is used by default.

 
\begin{Verbatim}[commandchars=!@|,fontsize=\small,frame=single,label=Example]
  !gapprompt@gap>| !gapinput@HaarGraph(3);|
  <immutable bipartite vertex-transitive symmetric digraph with bicompon\
  ents of size 2>
  !gapprompt@gap>| !gapinput@D := HaarGraph(16);|
  <immutable bipartite vertex-transitive symmetric digraph with bicompon\
  ents of size 5>
\end{Verbatim}
 }

 

\subsection{\textcolor{Chapter }{HalvedCubeGraph}}
\logpage{[ 3, 5, 21 ]}\nobreak
\hyperdef{L}{X801024F57DDC8A39}{}
{\noindent\textcolor{FuncColor}{$\triangleright$\enspace\texttt{HalvedCubeGraph({\mdseries\slshape [filt, ]n})\index{HalvedCubeGraph@\texttt{HalvedCubeGraph}}
\label{HalvedCubeGraph}
}\hfill{\scriptsize (operation)}}\\
\textbf{\indent Returns:\ }
A digraph.



 If \mbox{\texttt{\mdseries\slshape n}} is a positive integer at least 1, then this operation returns the \mbox{\texttt{\mdseries\slshape n}}th \emph{halved cube graph}, the graph of the \mbox{\texttt{\mdseries\slshape n}}\texttt{\symbol{45}} demihypercube. The vertices of the graph are those of the \mbox{\texttt{\mdseries\slshape n}}th\texttt{\symbol{45}} hypercube, with two vertices adjacent if and only if
they are at distance 1 or 2 from each other. Equivalent constructions are as
the second graph power of the \texttt{\mbox{\texttt{\mdseries\slshape n}}\texttt{\symbol{45}}1}th hypercube graph, or as with vertices labelled as the binary numbers where
two vertices are adjacent if they differ in a single bit, or with vertices
labelled with the subset of binary numbers with even Hamming weight, with
edges connecting vertices with Hamming distance exactly 2. The Halved Cube
graph is distance regular and contains a Hamiltonian cycle.

 See \href{https://mathworld.wolfram.com/HalvedCubeGraph.html} {\texttt{https://mathworld.wolfram.com/HalvedCubeGraph.html}} for further details.

 If the optional first argument \mbox{\texttt{\mdseries\slshape filt}} is present, then this should specify the category or representation the
digraph being created will belong to. For example, if \mbox{\texttt{\mdseries\slshape filt}} is \texttt{IsMutableDigraph} (\ref{IsMutableDigraph}), then the digraph being created will be mutable, if \mbox{\texttt{\mdseries\slshape filt}} is \texttt{IsImmutableDigraph} (\ref{IsImmutableDigraph}), then the digraph will be immutable. If the optional first argument \mbox{\texttt{\mdseries\slshape filt}} is not present, then \texttt{IsImmutableDigraph} (\ref{IsImmutableDigraph}) is used by default.

 
\begin{Verbatim}[commandchars=!@|,fontsize=\small,frame=single,label=Example]
  !gapprompt@gap>| !gapinput@D := HalvedCubeGraph(3);|
  <immutable Hamiltonian symmetric digraph with 4 vertices, 12 edges>
  !gapprompt@gap>| !gapinput@IsDistanceRegularDigraph(D);|
  true
  !gapprompt@gap>| !gapinput@IsHamiltonianDigraph(D);|
  true
\end{Verbatim}
 }

 

\subsection{\textcolor{Chapter }{HanoiGraph}}
\logpage{[ 3, 5, 22 ]}\nobreak
\hyperdef{L}{X7C43BDE47DF6553A}{}
{\noindent\textcolor{FuncColor}{$\triangleright$\enspace\texttt{HanoiGraph({\mdseries\slshape [filt, ]n})\index{HanoiGraph@\texttt{HanoiGraph}}
\label{HanoiGraph}
}\hfill{\scriptsize (operation)}}\\
\textbf{\indent Returns:\ }
A digraph.



 If \mbox{\texttt{\mdseries\slshape n}} is positive integer then this operation returns the \mbox{\texttt{\mdseries\slshape n}}th \emph{Hanoi graph}. The Hanoi graph's vertices represent the possible states of the 'Tower of
Hanoi' puzzle on three 'towers', while its edges represent possible moves. The
Hanoi graph has \texttt{3\texttt{\symbol{94}}\mbox{\texttt{\mdseries\slshape n}}} vertices, and \texttt{3 * (3\texttt{\symbol{94}}\mbox{\texttt{\mdseries\slshape n}} \texttt{\symbol{45}} 1) / 2} undirected edges.

 The Hanoi graph is Hamiltonian. The graph superficially resembles the
Sierpinski triangle. The graph is also a 'penny graph' \texttt{\symbol{45}} a
graph whose vertices can be considered to be
non\texttt{\symbol{45}}overlapping unit circles on a flat surface, with two
vertices adjacent only if the unit circles touch at a single point. Thus the
Hanoi graph is planar.

 See \href{https://mathworld.wolfram.com/HanoiGraph.html} {\texttt{https://mathworld.wolfram.com/HanoiGraph.html}} for further details.

 If the optional first argument \mbox{\texttt{\mdseries\slshape filt}} is present, then this should specify the category or representation the
digraph being created will belong to. For example, if \mbox{\texttt{\mdseries\slshape filt}} is \texttt{IsMutableDigraph} (\ref{IsMutableDigraph}), then the digraph being created will be mutable, if \mbox{\texttt{\mdseries\slshape filt}} is \texttt{IsImmutableDigraph} (\ref{IsImmutableDigraph}), then the digraph will be immutable. If the optional first argument \mbox{\texttt{\mdseries\slshape filt}} is not present, then \texttt{IsImmutableDigraph} (\ref{IsImmutableDigraph}) is used by default.

 
\begin{Verbatim}[commandchars=!@|,fontsize=\small,frame=single,label=Example]
  !gapprompt@gap>| !gapinput@D := HanoiGraph(5);|
  <immutable planar Hamiltonian symmetric digraph with 243 vertices, 726\
   edges>
  !gapprompt@gap>| !gapinput@IsPlanarDigraph(D);|
  true
\end{Verbatim}
 }

 

\subsection{\textcolor{Chapter }{HelmGraph}}
\logpage{[ 3, 5, 23 ]}\nobreak
\hyperdef{L}{X782ABFCE812B020A}{}
{\noindent\textcolor{FuncColor}{$\triangleright$\enspace\texttt{HelmGraph({\mdseries\slshape [filt, ]n})\index{HelmGraph@\texttt{HelmGraph}}
\label{HelmGraph}
}\hfill{\scriptsize (operation)}}\\
\textbf{\indent Returns:\ }
A digraph.



 If \mbox{\texttt{\mdseries\slshape n}} is a positive integer at least 3, then this operation returns the \mbox{\texttt{\mdseries\slshape n}}th \emph{helm graph}. The helm graph is the \mbox{\texttt{\mdseries\slshape n}}\texttt{\symbol{45}}1th wheel graph with, for each external vertex of the
'wheel', adjoining a new vertex incident only to the first vertex. That is,
the graph looks similar to a ship's helm.

 See \href{https://mathworld.wolfram.com/WheelGraph.html} {\texttt{https://mathworld.wolfram.com/WheelGraph.html}} for further details.

 If the optional first argument \mbox{\texttt{\mdseries\slshape filt}} is present, then this should specify the category or representation the
digraph being created will belong to. For example, if \mbox{\texttt{\mdseries\slshape filt}} is \texttt{IsMutableDigraph} (\ref{IsMutableDigraph}), then the digraph being created will be mutable, if \mbox{\texttt{\mdseries\slshape filt}} is \texttt{IsImmutableDigraph} (\ref{IsImmutableDigraph}), then the digraph will be immutable. If the optional first argument \mbox{\texttt{\mdseries\slshape filt}} is not present, then \texttt{IsImmutableDigraph} (\ref{IsImmutableDigraph}) is used by default.

 
\begin{Verbatim}[commandchars=!@|,fontsize=\small,frame=single,label=Example]
  !gapprompt@gap>| !gapinput@D := HelmGraph(4);|
  <immutable symmetric digraph with 9 vertices, 24 edges>
\end{Verbatim}
 }

 

\subsection{\textcolor{Chapter }{HypercubeGraph}}
\logpage{[ 3, 5, 24 ]}\nobreak
\hyperdef{L}{X7EE552F88609B1A2}{}
{\noindent\textcolor{FuncColor}{$\triangleright$\enspace\texttt{HypercubeGraph({\mdseries\slshape [filt, ]n})\index{HypercubeGraph@\texttt{HypercubeGraph}}
\label{HypercubeGraph}
}\hfill{\scriptsize (operation)}}\\
\textbf{\indent Returns:\ }
A digraph.



 If \mbox{\texttt{\mdseries\slshape n}} is a non\texttt{\symbol{45}}negative integer, then this operation returns the \mbox{\texttt{\mdseries\slshape n}}th \emph{hypercube graph}. The graph has \texttt{2\texttt{\symbol{94}}\mbox{\texttt{\mdseries\slshape n}}} vertices and \texttt{2\texttt{\symbol{94}}(\mbox{\texttt{\mdseries\slshape n}}\texttt{\symbol{45}}1)\mbox{\texttt{\mdseries\slshape n}}} edges. It is formed from the vertices and edges of the \mbox{\texttt{\mdseries\slshape n}}\texttt{\symbol{45}}dimensional hypercube. Alternatively, the graph can be
constructed by labelling each vertex with the binary numbers, with two
vertices adjacent if they have Hamming distance exactly one. The hypercube
graphs are Hamiltonian, distance\texttt{\symbol{45}}transitive and therefore
distance\texttt{\symbol{45}}regular, and bipartite.

 See \href{https://mathworld.wolfram.com/HypercubeGraph.html} {\texttt{https://mathworld.wolfram.com/HypercubeGraph.html}} for further details.

 If the optional first argument \mbox{\texttt{\mdseries\slshape filt}} is present, then this should specify the category or representation the
digraph being created will belong to. For example, if \mbox{\texttt{\mdseries\slshape filt}} is \texttt{IsMutableDigraph} (\ref{IsMutableDigraph}), then the digraph being created will be mutable, if \mbox{\texttt{\mdseries\slshape filt}} is \texttt{IsImmutableDigraph} (\ref{IsImmutableDigraph}), then the digraph will be immutable. If the optional first argument \mbox{\texttt{\mdseries\slshape filt}} is not present, then \texttt{IsImmutableDigraph} (\ref{IsImmutableDigraph}) is used by default.

 
\begin{Verbatim}[commandchars=!@|,fontsize=\small,frame=single,label=Example]
  !gapprompt@gap>| !gapinput@D := HypercubeGraph(5);|
  <immutable Hamiltonian bipartite symmetric digraph with bicomponents o\
  f size 16>
\end{Verbatim}
 }

 

\subsection{\textcolor{Chapter }{JohnsonDigraph}}
\logpage{[ 3, 5, 25 ]}\nobreak
\hyperdef{L}{X80ED9CE785819607}{}
{\noindent\textcolor{FuncColor}{$\triangleright$\enspace\texttt{JohnsonDigraph({\mdseries\slshape [filt, ]n, k})\index{JohnsonDigraph@\texttt{JohnsonDigraph}}
\label{JohnsonDigraph}
}\hfill{\scriptsize (operation)}}\\
\textbf{\indent Returns:\ }
A digraph.



 If \mbox{\texttt{\mdseries\slshape n}} and \mbox{\texttt{\mdseries\slshape k}} are non\texttt{\symbol{45}}negative integers, then this operation returns a
symmetric digraph which corresponds to the undirected \emph{Johnson graph} $J(n, k)$. 

 The \emph{Johnson graph} $J(n, k)$ has vertices given by all the \mbox{\texttt{\mdseries\slshape k}}\texttt{\symbol{45}}subsets of the range \texttt{[1 .. \mbox{\texttt{\mdseries\slshape n}}]}, and two vertices are connected by an edge if and only if their intersection
has size $\mbox{\texttt{\mdseries\slshape k}} - 1$. 

 If the optional first argument \mbox{\texttt{\mdseries\slshape filt}} is present, then this should specify the category or representation the
digraph being created will belong to. For example, if \mbox{\texttt{\mdseries\slshape filt}} is \texttt{IsMutableDigraph} (\ref{IsMutableDigraph}), then the digraph being created will be mutable, if \mbox{\texttt{\mdseries\slshape filt}} is \texttt{IsImmutableDigraph} (\ref{IsImmutableDigraph}), then the digraph will be immutable. If the optional first argument \mbox{\texttt{\mdseries\slshape filt}} is not present, then \texttt{IsImmutableDigraph} (\ref{IsImmutableDigraph}) is used by default.

 
\begin{Verbatim}[commandchars=!@|,fontsize=\small,frame=single,label=Example]
  !gapprompt@gap>| !gapinput@gr := JohnsonDigraph(3, 1);|
  <immutable symmetric digraph with 3 vertices, 6 edges>
  !gapprompt@gap>| !gapinput@OutNeighbours(gr);|
  [ [ 2, 3 ], [ 1, 3 ], [ 1, 2 ] ]
  !gapprompt@gap>| !gapinput@gr := JohnsonDigraph(4, 2);|
  <immutable symmetric digraph with 6 vertices, 24 edges>
  !gapprompt@gap>| !gapinput@OutNeighbours(gr);|
  [ [ 2, 3, 4, 5 ], [ 1, 3, 4, 6 ], [ 1, 2, 5, 6 ], [ 1, 2, 5, 6 ], 
    [ 1, 3, 4, 6 ], [ 2, 3, 4, 5 ] ]
  !gapprompt@gap>| !gapinput@JohnsonDigraph(1, 0);|
  <immutable empty digraph with 1 vertex>
  !gapprompt@gap>| !gapinput@JohnsonDigraph(IsMutableDigraph, 1, 0);|
  <mutable empty digraph with 1 vertex>
\end{Verbatim}
 }

 

\subsection{\textcolor{Chapter }{KellerGraph}}
\logpage{[ 3, 5, 26 ]}\nobreak
\hyperdef{L}{X79483C677AF65688}{}
{\noindent\textcolor{FuncColor}{$\triangleright$\enspace\texttt{KellerGraph({\mdseries\slshape [filt, ]n})\index{KellerGraph@\texttt{KellerGraph}}
\label{KellerGraph}
}\hfill{\scriptsize (operation)}}\\
\textbf{\indent Returns:\ }
A digraph.



 If \mbox{\texttt{\mdseries\slshape n}} is a nonnegative integer then this operation returns the \mbox{\texttt{\mdseries\slshape n}}th or \mbox{\texttt{\mdseries\slshape n}}\texttt{\symbol{45}}dimensional \emph{Keller graph}. The graph has vertices given by the \mbox{\texttt{\mdseries\slshape n}}\texttt{\symbol{45}}tuples on the set $[0, 1, 2, 3]$. Two vertices are adjacent if their respective tuples are such that they
differ in at least two coordinates and in at least one coordinate the
difference between the two is \texttt{2} mod \texttt{4}. The Keller graph has \texttt{4\texttt{\symbol{94}}\mbox{\texttt{\mdseries\slshape n}}} vertices.

 The Keller graphs were constructed with the intention of finding
counterexamples to Keller's conjecture (\href{https://mathworld.wolfram.com/KellersConjecture.html} {\texttt{https://mathworld.wolfram.com/KellersConjecture.html}}), and has been used since for testing maximum clique algorithms.

 If \mbox{\texttt{\mdseries\slshape n}} is 1 then the graph is empty, for \mbox{\texttt{\mdseries\slshape n}} greater than 1 the chromatic number of the Keller graph is \texttt{2\texttt{\symbol{94}}\mbox{\texttt{\mdseries\slshape n}}} and the graph is Hamiltonian.

 See \href{https://mathworld.wolfram.com/KellerGraph.html} {\texttt{https://mathworld.wolfram.com/KellerGraph.html}} for further details.

 If the optional first argument \mbox{\texttt{\mdseries\slshape filt}} is present, then this should specify the category or representation the
digraph being created will belong to. For example, if \mbox{\texttt{\mdseries\slshape filt}} is \texttt{IsMutableDigraph} (\ref{IsMutableDigraph}), then the digraph being created will be mutable, if \mbox{\texttt{\mdseries\slshape filt}} is \texttt{IsImmutableDigraph} (\ref{IsImmutableDigraph}), then the digraph will be immutable. If the optional first argument \mbox{\texttt{\mdseries\slshape filt}} is not present, then \texttt{IsImmutableDigraph} (\ref{IsImmutableDigraph}) is used by default.

 
\begin{Verbatim}[commandchars=!@|,fontsize=\small,frame=single,label=Example]
  !gapprompt@gap>| !gapinput@D := KellerGraph(3);|
  <immutable Hamiltonian symmetric digraph with 64 vertices, 2176 edges>
  !gapprompt@gap>| !gapinput@ChromaticNumber(D);|
  8
\end{Verbatim}
 }

 

\subsection{\textcolor{Chapter }{KingsGraph}}
\logpage{[ 3, 5, 27 ]}\nobreak
\hyperdef{L}{X80576A8C861512FD}{}
{\noindent\textcolor{FuncColor}{$\triangleright$\enspace\texttt{KingsGraph({\mdseries\slshape [filt, ]m, n})\index{KingsGraph@\texttt{KingsGraph}}
\label{KingsGraph}
}\hfill{\scriptsize (operation)}}\\
\textbf{\indent Returns:\ }
A digraph.



 If \mbox{\texttt{\mdseries\slshape m}} and \mbox{\texttt{\mdseries\slshape n}} are positive integers, then this operation returns the \emph{king's graph} of an \mbox{\texttt{\mdseries\slshape m}} by \mbox{\texttt{\mdseries\slshape n}} chessboard, as a symmetric digraph. 

 The king's graph represents all possible moves of the king chess piece across
a chessboard. An \mbox{\texttt{\mdseries\slshape m}} by \mbox{\texttt{\mdseries\slshape n}} chessboard is a grid of \mbox{\texttt{\mdseries\slshape m}} columns ("files") and \mbox{\texttt{\mdseries\slshape n}} rows ("ranks") that intersect in squares. Orthogonally adjacent squares are
alternately colored light and dark, with the square in the first rank and file
being dark.

 The king can move only to any orthogonally or diagonally adjacent square. Thus
the \texttt{\mbox{\texttt{\mdseries\slshape m}} * \mbox{\texttt{\mdseries\slshape n}}} vertices of the king's graph can be placed onto the \texttt{\mbox{\texttt{\mdseries\slshape m}} * \mbox{\texttt{\mdseries\slshape n}}} squares of an \mbox{\texttt{\mdseries\slshape m}} by \mbox{\texttt{\mdseries\slshape n}} chessboard, such that two vertices are adjacent in the digraph if and only if
the corresponding squares are orthogonally or diagonally adjacent on the
chessboard. 

 The chosen correspondence between vertices and chess squares is given by \texttt{DigraphVertexLabels} (\ref{DigraphVertexLabels}). In more detail, the vertices of the digraph are labelled by elements of the
Cartesian product \texttt{[1..\mbox{\texttt{\mdseries\slshape m}}] x [1..\mbox{\texttt{\mdseries\slshape n}}]}, where the first entry indexes the column (file) of the square in the
chessboard, and the second entry indexes the row (rank) of the square. (Note
that the files are traditionally indexed by the lowercase letters of the
alphabet). The vertices are sorted in ascending order, first by row (second
component) and then column (first component). See \href{https://en.wikipedia.org/wiki/King's_graph} {Wikipedia} for further information. See also \texttt{SquareGridGraph} (\ref{SquareGridGraph}), \texttt{TriangularGridGraph} (\ref{TriangularGridGraph}), and \texttt{StrongProduct} (\ref{StrongProduct}). 

 If the optional first argument \mbox{\texttt{\mdseries\slshape filt}} is present, then this should specify the category or representation the
digraph being created will belong to. For example, if \mbox{\texttt{\mdseries\slshape filt}} is \texttt{IsMutableDigraph} (\ref{IsMutableDigraph}), then the digraph being created will be mutable, if \mbox{\texttt{\mdseries\slshape filt}} is \texttt{IsImmutableDigraph} (\ref{IsImmutableDigraph}), then the digraph will be immutable. If the optional first argument \mbox{\texttt{\mdseries\slshape filt}} is not present, then \texttt{IsImmutableDigraph} (\ref{IsImmutableDigraph}) is used by default.

 
\begin{Verbatim}[commandchars=!@|,fontsize=\small,frame=single,label=Example]
  !gapprompt@gap>| !gapinput@KingsGraph(8, 8);|
  <immutable connected symmetric digraph with 64 vertices, 420 edges>
  !gapprompt@gap>| !gapinput@D := KingsGraph(IsMutable, 2, 7);|
  <mutable digraph with 14 vertices, 62 edges>
  !gapprompt@gap>| !gapinput@IsPlanarDigraph(D);|
  true
  !gapprompt@gap>| !gapinput@D := KingsGraph(3, 3);|
  <immutable planar connected symmetric digraph with 9 vertices, 40 edge\
  s>
  !gapprompt@gap>| !gapinput@OutNeighbors(D);|
  [ [ 2, 4, 5 ], [ 1, 3, 5, 4, 6 ], [ 2, 6, 5 ], [ 5, 1, 7, 2, 8 ], 
    [ 4, 6, 2, 8, 3, 7, 1, 9 ], [ 5, 3, 9, 8, 2 ], [ 8, 4, 5 ], 
    [ 7, 9, 5, 6, 4 ], [ 8, 6, 5 ] ]
  !gapprompt@gap>| !gapinput@IsSubdigraph(QueensGraph(3, 4), KingsGraph(3, 4));|
  true
\end{Verbatim}
 }

 

\subsection{\textcolor{Chapter }{KneserGraph}}
\logpage{[ 3, 5, 28 ]}\nobreak
\hyperdef{L}{X8655BA8584B3ACD0}{}
{\noindent\textcolor{FuncColor}{$\triangleright$\enspace\texttt{KneserGraph({\mdseries\slshape [filt, ]n, k})\index{KneserGraph@\texttt{KneserGraph}}
\label{KneserGraph}
}\hfill{\scriptsize (operation)}}\\
\textbf{\indent Returns:\ }
A digraph.



 If \mbox{\texttt{\mdseries\slshape n}} and \mbox{\texttt{\mdseries\slshape k}} are integers greater than 0, with \mbox{\texttt{\mdseries\slshape k}} less than \mbox{\texttt{\mdseries\slshape n}}, then this operation returns the \texttt{(\mbox{\texttt{\mdseries\slshape n}},\mbox{\texttt{\mdseries\slshape k}})}\texttt{\symbol{45}}th \emph{Kneser graph}. The Kneser graph's vertices correspond to the \mbox{\texttt{\mdseries\slshape k}}\texttt{\symbol{45}} element subsets of a set of \mbox{\texttt{\mdseries\slshape n}} elements, with two vertices being adjacent if and only if the subsets are
disjoint. The graph has \texttt{Binomial(\mbox{\texttt{\mdseries\slshape n}}, \mbox{\texttt{\mdseries\slshape k}})} vertices and \texttt{Binomial(\mbox{\texttt{\mdseries\slshape n}}, \mbox{\texttt{\mdseries\slshape k}}) * Binomial(\mbox{\texttt{\mdseries\slshape n}} \texttt{\symbol{45}} \mbox{\texttt{\mdseries\slshape k}}, \mbox{\texttt{\mdseries\slshape k}})} edges. Kneser graphs are regular, edge transitive, and vertex transitive. If \mbox{\texttt{\mdseries\slshape k}} is 1, then the graph is the complete graph on \mbox{\texttt{\mdseries\slshape n}} vertices. If \texttt{(\mbox{\texttt{\mdseries\slshape n}}, \mbox{\texttt{\mdseries\slshape k}})} is \texttt{(2m\texttt{\symbol{45}}1, \mbox{\texttt{\mdseries\slshape n}}\texttt{\symbol{45}}1)}, then the graph is the \texttt{m}th Odd graph. The Petersen graph is the $(5, 2)$th Kneser graph.

 If \texttt{\mbox{\texttt{\mdseries\slshape n}} {\textgreater}= 2\mbox{\texttt{\mdseries\slshape k}}} then the graph's chromatic number is \texttt{\mbox{\texttt{\mdseries\slshape n}} \texttt{\symbol{45}} 2\mbox{\texttt{\mdseries\slshape k}} + 2}, and otherwise is 1. The Kneser graph contains a Hamiltonian cycle if \texttt{\mbox{\texttt{\mdseries\slshape n}} {\textgreater}= ((3 + 5 \texttt{\symbol{94}} 0.5) / 2)\mbox{\texttt{\mdseries\slshape k}} + 1}. The graph has clique number equal to the floor of \texttt{\mbox{\texttt{\mdseries\slshape n}} / \mbox{\texttt{\mdseries\slshape k}}}.

 See \href{https://mathworld.wolfram.com/KneserGraph.html} {\texttt{https://mathworld.wolfram.com/KneserGraph.html}} for further details.

 If the optional first argument \mbox{\texttt{\mdseries\slshape filt}} is present, then this should specify the category or representation the
digraph being created will belong to. For example, if \mbox{\texttt{\mdseries\slshape filt}} is \texttt{IsMutableDigraph} (\ref{IsMutableDigraph}), then the digraph being created will be mutable, if \mbox{\texttt{\mdseries\slshape filt}} is \texttt{IsImmutableDigraph} (\ref{IsImmutableDigraph}), then the digraph will be immutable. If the optional first argument \mbox{\texttt{\mdseries\slshape filt}} is not present, then \texttt{IsImmutableDigraph} (\ref{IsImmutableDigraph}) is used by default.

 
\begin{Verbatim}[commandchars=!@|,fontsize=\small,frame=single,label=Example]
  !gapprompt@gap>| !gapinput@D := KneserGraph(7, 3);|
  <immutable edge- and vertex-transitive symmetric digraph with 35 verti\
  ces, 140 edges>
  !gapprompt@gap>| !gapinput@IsRegularDigraph(D);|
  true
  !gapprompt@gap>| !gapinput@ChromaticNumber(D);|
  3
  !gapprompt@gap>| !gapinput@CliqueNumber(D);|
  2
\end{Verbatim}
 }

 

\subsection{\textcolor{Chapter }{KnightsGraph}}
\logpage{[ 3, 5, 29 ]}\nobreak
\hyperdef{L}{X84609DCA79FD9B56}{}
{\noindent\textcolor{FuncColor}{$\triangleright$\enspace\texttt{KnightsGraph({\mdseries\slshape [filt, ]m, n})\index{KnightsGraph@\texttt{KnightsGraph}}
\label{KnightsGraph}
}\hfill{\scriptsize (operation)}}\\
\textbf{\indent Returns:\ }
A digraph.



 If \mbox{\texttt{\mdseries\slshape m}} and \mbox{\texttt{\mdseries\slshape n}} are positive integers, then this operation returns the \emph{knight's graph} of an \mbox{\texttt{\mdseries\slshape m}} by \mbox{\texttt{\mdseries\slshape n}} chessboard, as a symmetric digraph. 

 A knight's graph represents all possible moves of the knight chess piece
across a chessboard. An \mbox{\texttt{\mdseries\slshape m}} by \mbox{\texttt{\mdseries\slshape n}} chessboard is a grid of \mbox{\texttt{\mdseries\slshape m}} columns ("files") and \mbox{\texttt{\mdseries\slshape n}} rows ("ranks") that intersect in squares. Orthogonally adjacent squares are
alternately colored light and dark, with the square in the first rank and file
being dark.

 The \texttt{\mbox{\texttt{\mdseries\slshape m}} * \mbox{\texttt{\mdseries\slshape n}}} vertices of the knight's graph can be placed onto the \texttt{\mbox{\texttt{\mdseries\slshape m}} * \mbox{\texttt{\mdseries\slshape n}}} squares of an \mbox{\texttt{\mdseries\slshape m}} by \mbox{\texttt{\mdseries\slshape n}} chessboard, such that two vertices are adjacent in the digraph if and only if
a knight can move between the corresponding squares in a single turn. 

 The chosen correspondence between vertices and chess squares is given by \texttt{DigraphVertexLabels} (\ref{DigraphVertexLabels}). In more detail, the vertices of the digraph are labelled by elements of the
Cartesian product \texttt{[1..\mbox{\texttt{\mdseries\slshape m}}] x [1..\mbox{\texttt{\mdseries\slshape n}}]}, where the first entry indexes the column (file) of the square in the
chessboard, and the second entry indexes the row (rank) of the square. (Note
that the files are traditionally indexed by the lowercase letters of the
alphabet). The vertices are sorted in ascending order, first by row (second
component) and then column (first component). See \href{https://en.wikipedia.org/wiki/Knight's_graph} {Wikipedia} for further information. 

 If the optional first argument \mbox{\texttt{\mdseries\slshape filt}} is present, then this should specify the category or representation the
digraph being created will belong to. For example, if \mbox{\texttt{\mdseries\slshape filt}} is \texttt{IsMutableDigraph} (\ref{IsMutableDigraph}), then the digraph being created will be mutable, if \mbox{\texttt{\mdseries\slshape filt}} is \texttt{IsImmutableDigraph} (\ref{IsImmutableDigraph}), then the digraph will be immutable. If the optional first argument \mbox{\texttt{\mdseries\slshape filt}} is not present, then \texttt{IsImmutableDigraph} (\ref{IsImmutableDigraph}) is used by default.

 
\begin{Verbatim}[commandchars=!@|,fontsize=\small,frame=single,label=Example]
  !gapprompt@gap>| !gapinput@D := KnightsGraph(8, 8);|
  <immutable connected symmetric digraph with 64 vertices, 336 edges>
  !gapprompt@gap>| !gapinput@IsConnectedDigraph(D);|
  true
  !gapprompt@gap>| !gapinput@D := KnightsGraph(3, 3);|
  <immutable symmetric digraph with 9 vertices, 16 edges>
  !gapprompt@gap>| !gapinput@IsConnectedDigraph(D);|
  false
  !gapprompt@gap>| !gapinput@KnightsGraph(IsMutable, 3, 9);|
  <mutable digraph with 27 vertices, 88 edges>
\end{Verbatim}
 }

 

\subsection{\textcolor{Chapter }{LindgrenSousselierGraph}}
\logpage{[ 3, 5, 30 ]}\nobreak
\hyperdef{L}{X7F61140C822880DA}{}
{\noindent\textcolor{FuncColor}{$\triangleright$\enspace\texttt{LindgrenSousselierGraph({\mdseries\slshape [filt, ]n})\index{LindgrenSousselierGraph@\texttt{LindgrenSousselierGraph}}
\label{LindgrenSousselierGraph}
}\hfill{\scriptsize (operation)}}\\
\textbf{\indent Returns:\ }
A digraph.



 If \mbox{\texttt{\mdseries\slshape n}} is an integer greater than 0, then this operation returns the \mbox{\texttt{\mdseries\slshape n}}th \emph{Lindgren\texttt{\symbol{45}}Sousselier graph}, an infinite family of hyophamiltonian graphs \texttt{\symbol{45}} a graph
that is non\texttt{\symbol{45}}Hamiltonian but removing any vector gives a
Hamiltonian graph. The graph has \texttt{6\mbox{\texttt{\mdseries\slshape n}}+4} vertices and \texttt{15 + 10(\mbox{\texttt{\mdseries\slshape n}}\texttt{\symbol{45}}1)} undirected edges. The first Lindgren\texttt{\symbol{45}}Sousselier graph is
the Petersen graph, and is in fact the smallest hyophamiltonian graph.

 See \href{https://mathworld.wolfram.com/HypohamiltonianGraph.html} {\texttt{https://mathworld.wolfram.com/HypohamiltonianGraph.html}} for further details.

 If the optional first argument \mbox{\texttt{\mdseries\slshape filt}} is present, then this should specify the category or representation the
digraph being created will belong to. For example, if \mbox{\texttt{\mdseries\slshape filt}} is \texttt{IsMutableDigraph} (\ref{IsMutableDigraph}), then the digraph being created will be mutable, if \mbox{\texttt{\mdseries\slshape filt}} is \texttt{IsImmutableDigraph} (\ref{IsImmutableDigraph}), then the digraph will be immutable. If the optional first argument \mbox{\texttt{\mdseries\slshape filt}} is not present, then \texttt{IsImmutableDigraph} (\ref{IsImmutableDigraph}) is used by default.

 
\begin{Verbatim}[commandchars=!@|,fontsize=\small,frame=single,label=Example]
  !gapprompt@gap>| !gapinput@D := LindgrenSousselierGraph(3);|
  <immutable symmetric digraph with 22 vertices, 70 edges>
  !gapprompt@gap>| !gapinput@IsHamiltonianDigraph(D);|
  false
  !gapprompt@gap>| !gapinput@IsHamiltonianDigraph(DigraphRemoveVertex(D, 1));|
  true
  !gapprompt@gap>| !gapinput@IsIsomorphicDigraph(LindgrenSousselierGraph(1), PetersenGraph());|
  true
\end{Verbatim}
 }

 

\subsection{\textcolor{Chapter }{LollipopGraph}}
\logpage{[ 3, 5, 31 ]}\nobreak
\hyperdef{L}{X832E82CF87BF5D43}{}
{\noindent\textcolor{FuncColor}{$\triangleright$\enspace\texttt{LollipopGraph({\mdseries\slshape [filt, ]m, n})\index{LollipopGraph@\texttt{LollipopGraph}}
\label{LollipopGraph}
}\hfill{\scriptsize (operation)}}\\
\textbf{\indent Returns:\ }
A digraph.



 If \mbox{\texttt{\mdseries\slshape m}} and \mbox{\texttt{\mdseries\slshape n}} are positive integers, then this operation returns the \emph{(\mbox{\texttt{\mdseries\slshape m}},\mbox{\texttt{\mdseries\slshape n}})\texttt{\symbol{45}}lollipop graph}. As defined at \href{https://en.wikipedia.org/wiki/Lollipop_graph} {\texttt{https://en.wikipedia.org/wiki/Lollipop{\textunderscore}graph}}, this consists of a complete digraph on the vertices \texttt{[1..\mbox{\texttt{\mdseries\slshape m}}]} (the 'head' of the lollipop), and the symmetric closure of a chain digraph on
the remaining \mbox{\texttt{\mdseries\slshape n}} vertices (the 'stick'), connected by a bridge (the edge \texttt{[\mbox{\texttt{\mdseries\slshape m}}, \mbox{\texttt{\mdseries\slshape m}}+1]} and its reverse). 

 If the optional first argument \mbox{\texttt{\mdseries\slshape filt}} is present, then this should specify the category or representation the
digraph being created will belong to. For example, if \mbox{\texttt{\mdseries\slshape filt}} is \texttt{IsMutableDigraph} (\ref{IsMutableDigraph}), then the digraph being created will be mutable, if \mbox{\texttt{\mdseries\slshape filt}} is \texttt{IsImmutableDigraph} (\ref{IsImmutableDigraph}), then the digraph will be immutable. If the optional first argument \mbox{\texttt{\mdseries\slshape filt}} is not present, then \texttt{IsImmutableDigraph} (\ref{IsImmutableDigraph}) is used by default.

 
\begin{Verbatim}[commandchars=!@|,fontsize=\small,frame=single,label=Example]
  !gapprompt@gap>| !gapinput@D := LollipopGraph(5, 3);|
  <immutable connected symmetric digraph with 8 vertices, 26 edges>
  !gapprompt@gap>| !gapinput@CliqueNumber(D);|
  5
  !gapprompt@gap>| !gapinput@DigraphUndirectedGirth(D);|
  3
  !gapprompt@gap>| !gapinput@LollipopGraph(IsMutableDigraph, 3, 8);|
  <mutable digraph with 11 vertices, 22 edges>
\end{Verbatim}
 }

 

\subsection{\textcolor{Chapter }{MobiusLadderGraph}}
\logpage{[ 3, 5, 32 ]}\nobreak
\hyperdef{L}{X853613228110588E}{}
{\noindent\textcolor{FuncColor}{$\triangleright$\enspace\texttt{MobiusLadderGraph({\mdseries\slshape [filt, ]n})\index{MobiusLadderGraph@\texttt{MobiusLadderGraph}}
\label{MobiusLadderGraph}
}\hfill{\scriptsize (operation)}}\\
\textbf{\indent Returns:\ }
A digraph.



 If \mbox{\texttt{\mdseries\slshape n}} is a positive integer at least 4, then this operation returns the \emph{Mobius ladder graph} that is obtained by introducing a 'twist' in the \mbox{\texttt{\mdseries\slshape n}}th prism graph, similar to the construction of a Mobius strip. The Mobius
ladder graph is isomorphic to the circulant graph \texttt{Ci(2\mbox{\texttt{\mdseries\slshape n}}, [1, \mbox{\texttt{\mdseries\slshape n}}])}. The Mobius ladders are cubic, symmetric, Hamiltonian,
vertex\texttt{\symbol{45}}transitive, and graceful. They are also
non\texttt{\symbol{45}}planar and apex, meaning removing a single vertex
produces a planar graph.

 See \href{https://mathworld.wolfram.com/MoebiusLadder.html} {\texttt{https://mathworld.wolfram.com/MoebiusLadder.html}} for further details.

 If the optional first argument \mbox{\texttt{\mdseries\slshape filt}} is present, then this should specify the category or representation the
digraph being created will belong to. For example, if \mbox{\texttt{\mdseries\slshape filt}} is \texttt{IsMutableDigraph} (\ref{IsMutableDigraph}), then the digraph being created will be mutable, if \mbox{\texttt{\mdseries\slshape filt}} is \texttt{IsImmutableDigraph} (\ref{IsImmutableDigraph}), then the digraph will be immutable. If the optional first argument \mbox{\texttt{\mdseries\slshape filt}} is not present, then \texttt{IsImmutableDigraph} (\ref{IsImmutableDigraph}) is used by default.

 
\begin{Verbatim}[commandchars=!@|,fontsize=\small,frame=single,label=Example]
  !gapprompt@gap>| !gapinput@D := MobiusLadderGraph(7);|
  <immutable symmetric digraph with 14 vertices, 42 edges>
  !gapprompt@gap>| !gapinput@IsHamiltonianDigraph(D);|
  true
  !gapprompt@gap>| !gapinput@IsVertexTransitive(D);|
  true
  !gapprompt@gap>| !gapinput@IsPlanarDigraph(D);|
  false
  !gapprompt@gap>| !gapinput@D2 := DigraphRemoveVertex(D, 1);|
  <immutable digraph with 13 vertices, 36 edges>
  !gapprompt@gap>| !gapinput@IsPlanarDigraph(D2);|
  true
\end{Verbatim}
 }

 

\subsection{\textcolor{Chapter }{MycielskiGraph}}
\logpage{[ 3, 5, 33 ]}\nobreak
\hyperdef{L}{X825943547FD7A687}{}
{\noindent\textcolor{FuncColor}{$\triangleright$\enspace\texttt{MycielskiGraph({\mdseries\slshape [filt, ]n})\index{MycielskiGraph@\texttt{MycielskiGraph}}
\label{MycielskiGraph}
}\hfill{\scriptsize (operation)}}\\
\textbf{\indent Returns:\ }
A digraph.



 If \mbox{\texttt{\mdseries\slshape n}} is an integer greater than 1, then this operation returns the \mbox{\texttt{\mdseries\slshape n}}th \emph{Mycielski graph}. The Mycielskian of a triangle\texttt{\symbol{45}}free graph is a
construction that adds vertices and edges to produce a new graph that is still
triangle\texttt{\symbol{45}}free but has a larger chromatic number. The
Mycielski graphs are a series of graphs with this construction repeated,
starting with the complete graph on two vertices. The graph has \texttt{3 * 2\texttt{\symbol{94}}(\mbox{\texttt{\mdseries\slshape n}}\texttt{\symbol{45}}2) \texttt{\symbol{45}} 1} vertices, with the number of edges being \texttt{(18 \texttt{\symbol{45}} 27 * 2 \texttt{\symbol{94}} \mbox{\texttt{\mdseries\slshape n}} + 14 * 3 \texttt{\symbol{94}} \mbox{\texttt{\mdseries\slshape n}}) / 36}.

 The Mycielski graph has chromatic number equal to \mbox{\texttt{\mdseries\slshape n}}, clique number equal to $2$, and is Hamiltonian. The graph is in fact the graph with chromatic number \mbox{\texttt{\mdseries\slshape n}} with the least possible vertices.

 See \href{https://mathworld.wolfram.com/MycielskiGraph.html} {\texttt{https://mathworld.wolfram.com/MycielskiGraph.html}} for further details.

 If the optional first argument \mbox{\texttt{\mdseries\slshape filt}} is present, then this should specify the category or representation the
digraph being created will belong to. For example, if \mbox{\texttt{\mdseries\slshape filt}} is \texttt{IsMutableDigraph} (\ref{IsMutableDigraph}), then the digraph being created will be mutable, if \mbox{\texttt{\mdseries\slshape filt}} is \texttt{IsImmutableDigraph} (\ref{IsImmutableDigraph}), then the digraph will be immutable. If the optional first argument \mbox{\texttt{\mdseries\slshape filt}} is not present, then \texttt{IsImmutableDigraph} (\ref{IsImmutableDigraph}) is used by default.

 
\begin{Verbatim}[commandchars=!@|,fontsize=\small,frame=single,label=Example]
  !gapprompt@gap>| !gapinput@D := MycielskiGraph(4);|
  <immutable Hamiltonian symmetric digraph with 11 vertices, 40 edges>
  !gapprompt@gap>| !gapinput@ChromaticNumber(D);|
  4
  !gapprompt@gap>| !gapinput@CliqueNumber(D);|
  2
\end{Verbatim}
 }

 

\subsection{\textcolor{Chapter }{OddGraph}}
\logpage{[ 3, 5, 34 ]}\nobreak
\hyperdef{L}{X7904BB2982014ADA}{}
{\noindent\textcolor{FuncColor}{$\triangleright$\enspace\texttt{OddGraph({\mdseries\slshape [filt, ]n})\index{OddGraph@\texttt{OddGraph}}
\label{OddGraph}
}\hfill{\scriptsize (operation)}}\\
\textbf{\indent Returns:\ }
A digraph.



 If \mbox{\texttt{\mdseries\slshape n}} is an integer greater than 0, then this operation returns the \mbox{\texttt{\mdseries\slshape n}}th \emph{odd graph}. The odd graph has vertices labelled with the \texttt{\mbox{\texttt{\mdseries\slshape n}}\texttt{\symbol{45}}1}\texttt{\symbol{45}}subsets of a \texttt{2\mbox{\texttt{\mdseries\slshape n}}\texttt{\symbol{45}}1}\texttt{\symbol{45}}set, with two vertices adjacent if and only if their
subsets are disjoint. The \mbox{\texttt{\mdseries\slshape n}}th odd graph is the \texttt{(2\mbox{\texttt{\mdseries\slshape n}}\texttt{\symbol{45}}1, \mbox{\texttt{\mdseries\slshape n}}\texttt{\symbol{45}}1)}\texttt{\symbol{45}}th Kneser graph. The graph has \texttt{Binomial(2\mbox{\texttt{\mdseries\slshape n}}\texttt{\symbol{45}}1, \mbox{\texttt{\mdseries\slshape n}}\texttt{\symbol{45}}1)} vertices and \texttt{\mbox{\texttt{\mdseries\slshape n}} * Binomial(2\mbox{\texttt{\mdseries\slshape n}}\texttt{\symbol{45}}1, \mbox{\texttt{\mdseries\slshape n}}\texttt{\symbol{45}}1) / 2} edges.

 The odd graph is regular and distance transitive (and therefore distance
regular). They have chromatic number equal to 3, and all Odd graphs with \mbox{\texttt{\mdseries\slshape n}} greater than 3 are Hamiltonian. They are also vertex and edge transitive.

 See \href{https://mathworld.wolfram.com/OddGraph.html} {\texttt{https://mathworld.wolfram.com/OddGraph.html}} for further details.

 If the optional first argument \mbox{\texttt{\mdseries\slshape filt}} is present, then this should specify the category or representation the
digraph being created will belong to. For example, if \mbox{\texttt{\mdseries\slshape filt}} is \texttt{IsMutableDigraph} (\ref{IsMutableDigraph}), then the digraph being created will be mutable, if \mbox{\texttt{\mdseries\slshape filt}} is \texttt{IsImmutableDigraph} (\ref{IsImmutableDigraph}), then the digraph will be immutable. If the optional first argument \mbox{\texttt{\mdseries\slshape filt}} is not present, then \texttt{IsImmutableDigraph} (\ref{IsImmutableDigraph}) is used by default.

 
\begin{Verbatim}[commandchars=!@|,fontsize=\small,frame=single,label=Example]
  !gapprompt@gap>| !gapinput@D := OddGraph(4);|
  <immutable edge- and vertex-transitive symmetric digraph with 35 verti\
  ces, 140 edges>
  !gapprompt@gap>| !gapinput@IsIsomorphicDigraph(D, KneserGraph(7, 3));|
  true
  !gapprompt@gap>| !gapinput@ChromaticNumber(D);|
  3
\end{Verbatim}
 }

 

\subsection{\textcolor{Chapter }{PathGraph}}
\logpage{[ 3, 5, 35 ]}\nobreak
\hyperdef{L}{X815055168405B7F0}{}
{\noindent\textcolor{FuncColor}{$\triangleright$\enspace\texttt{PathGraph({\mdseries\slshape [filt, ]n})\index{PathGraph@\texttt{PathGraph}}
\label{PathGraph}
}\hfill{\scriptsize (operation)}}\\
\textbf{\indent Returns:\ }
A digraph.



 If \mbox{\texttt{\mdseries\slshape n}} is a non\texttt{\symbol{45}}negative integer then this operation returns the \mbox{\texttt{\mdseries\slshape n}}th \emph{path graph}, consisting of the path on \mbox{\texttt{\mdseries\slshape n}} vertices. This is the symmetric closure of the \texttt{ChainDigraph} (\ref{ChainDigraph}). The path graph has \mbox{\texttt{\mdseries\slshape n}} vertices and \texttt{\mbox{\texttt{\mdseries\slshape n}}\texttt{\symbol{45}}1} edges. The path graph is an undirected tree.

 See \href{https://mathworld.wolfram.com/PathGraph.html} {\texttt{https://mathworld.wolfram.com/PathGraph.html}} for further details.

 If the optional first argument \mbox{\texttt{\mdseries\slshape filt}} is present, then this should specify the category or representation the
digraph being created will belong to. For example, if \mbox{\texttt{\mdseries\slshape filt}} is \texttt{IsMutableDigraph} (\ref{IsMutableDigraph}), then the digraph being created will be mutable, if \mbox{\texttt{\mdseries\slshape filt}} is \texttt{IsImmutableDigraph} (\ref{IsImmutableDigraph}), then the digraph will be immutable. If the optional first argument \mbox{\texttt{\mdseries\slshape filt}} is not present, then \texttt{IsImmutableDigraph} (\ref{IsImmutableDigraph}) is used by default.

 
\begin{Verbatim}[commandchars=!@|,fontsize=\small,frame=single,label=Example]
  !gapprompt@gap>| !gapinput@D := PathGraph(12);|
  <immutable undirected tree with 12 vertices>
\end{Verbatim}
 }

 

\subsection{\textcolor{Chapter }{PermutationStarGraph}}
\logpage{[ 3, 5, 36 ]}\nobreak
\hyperdef{L}{X7A38DFC47AAE4A96}{}
{\noindent\textcolor{FuncColor}{$\triangleright$\enspace\texttt{PermutationStarGraph({\mdseries\slshape [filt, ]n, k})\index{PermutationStarGraph@\texttt{PermutationStarGraph}}
\label{PermutationStarGraph}
}\hfill{\scriptsize (operation)}}\\
\textbf{\indent Returns:\ }
A digraph.



 If \mbox{\texttt{\mdseries\slshape n}} is an integer at greater than 0 and \mbox{\texttt{\mdseries\slshape k}} is an integer greater than 1, and \mbox{\texttt{\mdseries\slshape k}} less than or equal to \mbox{\texttt{\mdseries\slshape n}}, then this operation returns the \texttt{(\mbox{\texttt{\mdseries\slshape n}}, \mbox{\texttt{\mdseries\slshape k}})}\texttt{\symbol{45}}th \emph{permutation star graph}. The vertices of the graph are given by the \mbox{\texttt{\mdseries\slshape k}}\texttt{\symbol{45}}length ordered subsets of an \mbox{\texttt{\mdseries\slshape n}}\texttt{\symbol{45}}set, with two vertices being adjacent if one is labelled \texttt{p1 p2 p3 ... pi ... p\mbox{\texttt{\mdseries\slshape k}}}, and the other is either labelled \texttt{pi p2 p3 ... p1 ... p\mbox{\texttt{\mdseries\slshape k}}}, or labelled \texttt{x p2 p3 ... pi ... p\mbox{\texttt{\mdseries\slshape k}}} where \texttt{x} is in \texttt{[1..\mbox{\texttt{\mdseries\slshape n}}]} and is not equal to \texttt{p1}. The graph has \texttt{\mbox{\texttt{\mdseries\slshape n}}! / (\mbox{\texttt{\mdseries\slshape n}} \texttt{\symbol{45}} \mbox{\texttt{\mdseries\slshape k}})!} vertices.

 The permutation star graph is regular and vertex transitive. It has diameter \texttt{2\mbox{\texttt{\mdseries\slshape k}} \texttt{\symbol{45}} 1} if \mbox{\texttt{\mdseries\slshape k}} less than or equal to \texttt{Int(\mbox{\texttt{\mdseries\slshape n}} / 2)}, and diameter \texttt{Int((\mbox{\texttt{\mdseries\slshape n}} \texttt{\symbol{45}} 1) / 2) + \mbox{\texttt{\mdseries\slshape k}}} if \texttt{\mbox{\texttt{\mdseries\slshape k}} {\textgreater}= Int(\mbox{\texttt{\mdseries\slshape n}} / 2) + 1}.

 See \href{https://mathworld.wolfram.com/PermutationStarGraph.html} {\texttt{https://mathworld.wolfram.com/PermutationStarGraph.html}} for further details.

 If the optional first argument \mbox{\texttt{\mdseries\slshape filt}} is present, then this should specify the category or representation the
digraph being created will belong to. For example, if \mbox{\texttt{\mdseries\slshape filt}} is \texttt{IsMutableDigraph} (\ref{IsMutableDigraph}), then the digraph being created will be mutable, if \mbox{\texttt{\mdseries\slshape filt}} is \texttt{IsImmutableDigraph} (\ref{IsImmutableDigraph}), then the digraph will be immutable. If the optional first argument \mbox{\texttt{\mdseries\slshape filt}} is not present, then \texttt{IsImmutableDigraph} (\ref{IsImmutableDigraph}) is used by default.

 
\begin{Verbatim}[commandchars=!@|,fontsize=\small,frame=single,label=Example]
  !gapprompt@gap>| !gapinput@D := PermutationStarGraph(6, 2);|
  <immutable vertex-transitive symmetric digraph with 30 vertices, 150 e\
  dges>
  !gapprompt@gap>| !gapinput@DigraphDiameter(D);|
  3
\end{Verbatim}
 }

 

\subsection{\textcolor{Chapter }{PetersenGraph}}
\logpage{[ 3, 5, 37 ]}\nobreak
\hyperdef{L}{X823F43217A6C375D}{}
{\noindent\textcolor{FuncColor}{$\triangleright$\enspace\texttt{PetersenGraph({\mdseries\slshape [filt]})\index{PetersenGraph@\texttt{PetersenGraph}}
\label{PetersenGraph}
}\hfill{\scriptsize (operation)}}\\
\textbf{\indent Returns:\ }
A digraph.



 From \href{https://en.wikipedia.org/wiki/Petersen_graph} {\texttt{https://en.wikipedia.org/wiki/Petersen{\textunderscore}graph}}:

 ``The Petersen graph is an undirected graph with 10 vertices and 15 edges. It is
a small graph that serves as a useful example and counterexample for many
problems in graph theory. The Petersen graph is named after Julius Petersen,
who in 1898 constructed it to be the smallest bridgeless cubic graph with no
three\texttt{\symbol{45}}edge\texttt{\symbol{45}}coloring.''

 See also \texttt{GeneralisedPetersenGraph} (\ref{GeneralisedPetersenGraph}). 

 If the optional first argument \mbox{\texttt{\mdseries\slshape filt}} is present, then this should specify the category or representation the
digraph being created will belong to. For example, if \mbox{\texttt{\mdseries\slshape filt}} is \texttt{IsMutableDigraph} (\ref{IsMutableDigraph}), then the digraph being created will be mutable, if \mbox{\texttt{\mdseries\slshape filt}} is \texttt{IsImmutableDigraph} (\ref{IsImmutableDigraph}), then the digraph will be immutable. If the optional first argument \mbox{\texttt{\mdseries\slshape filt}} is not present, then \texttt{IsImmutableDigraph} (\ref{IsImmutableDigraph}) is used by default.

 
\begin{Verbatim}[commandchars=!@|,fontsize=\small,frame=single,label=Example]
  !gapprompt@gap>| !gapinput@ChromaticNumber(PetersenGraph());|
  3
  !gapprompt@gap>| !gapinput@PetersenGraph(IsMutableDigraph);|
  <mutable digraph with 10 vertices, 30 edges>
\end{Verbatim}
 }

 

\subsection{\textcolor{Chapter }{GeneralisedPetersenGraph}}
\logpage{[ 3, 5, 38 ]}\nobreak
\hyperdef{L}{X7B5441F386BD105E}{}
{\noindent\textcolor{FuncColor}{$\triangleright$\enspace\texttt{GeneralisedPetersenGraph({\mdseries\slshape [filt, ]n, k})\index{GeneralisedPetersenGraph@\texttt{GeneralisedPetersenGraph}}
\label{GeneralisedPetersenGraph}
}\hfill{\scriptsize (operation)}}\\
\textbf{\indent Returns:\ }
A digraph.



 If \mbox{\texttt{\mdseries\slshape n}} is a positive integer and \mbox{\texttt{\mdseries\slshape k}} is a non\texttt{\symbol{45}}negative integer less than \texttt{\mbox{\texttt{\mdseries\slshape n}} / 2}, then this operation returns the \emph{generalised Petersen graph} $GPG(\mbox{\texttt{\mdseries\slshape n}}, k)$. 

 From \href{https://en.wikipedia.org/wiki/Generalized_Petersen_graph} {\texttt{https://en.wikipedia.org/wiki/Generalized{\textunderscore}Petersen{\textunderscore}graph}}: 

 ``The generalized Petersen graphs are a family of cubic graphs formed by
connecting the vertices of a regular polygon to the corresponding vertices of
a star polygon. They include the Petersen graph and generalize one of the ways
of constructing the Petersen graph. The generalized Petersen graph family was
introduced in 1950 by H. S. M. Coxeter and was given its name in 1969 by Mark
Watkins.'' 

 See also \texttt{PetersenGraph} (\ref{PetersenGraph}). 

 If the optional first argument \mbox{\texttt{\mdseries\slshape filt}} is present, then this should specify the category or representation the
digraph being created will belong to. For example, if \mbox{\texttt{\mdseries\slshape filt}} is \texttt{IsMutableDigraph} (\ref{IsMutableDigraph}), then the digraph being created will be mutable, if \mbox{\texttt{\mdseries\slshape filt}} is \texttt{IsImmutableDigraph} (\ref{IsImmutableDigraph}), then the digraph will be immutable. If the optional first argument \mbox{\texttt{\mdseries\slshape filt}} is not present, then \texttt{IsImmutableDigraph} (\ref{IsImmutableDigraph}) is used by default.

 
\begin{Verbatim}[commandchars=!@|,fontsize=\small,frame=single,label=Example]
  !gapprompt@gap>| !gapinput@GeneralisedPetersenGraph(7, 2);|
  <immutable symmetric digraph with 14 vertices, 42 edges>
  !gapprompt@gap>| !gapinput@GeneralisedPetersenGraph(40, 1);|
  <immutable symmetric digraph with 80 vertices, 240 edges>
  !gapprompt@gap>| !gapinput@D := GeneralisedPetersenGraph(5, 2);|
  <immutable symmetric digraph with 10 vertices, 30 edges>
  !gapprompt@gap>| !gapinput@IsIsomorphicDigraph(D, PetersenGraph());|
  true
  !gapprompt@gap>| !gapinput@GeneralisedPetersenGraph(IsMutableDigraph, 9, 4);|
  <mutable digraph with 18 vertices, 54 edges>
\end{Verbatim}
 }

 

\subsection{\textcolor{Chapter }{PrismGraph}}
\logpage{[ 3, 5, 39 ]}\nobreak
\hyperdef{L}{X85425DC5847E6D20}{}
{\noindent\textcolor{FuncColor}{$\triangleright$\enspace\texttt{PrismGraph({\mdseries\slshape [filt, ]n})\index{PrismGraph@\texttt{PrismGraph}}
\label{PrismGraph}
}\hfill{\scriptsize (operation)}}\\
\textbf{\indent Returns:\ }
A digraph.



 If \mbox{\texttt{\mdseries\slshape n}} is a positive integer at least 3, then this operation returns the \emph{prism graph} that is the skeleton of the \mbox{\texttt{\mdseries\slshape n}}\texttt{\symbol{45}}prism. It has \texttt{2\mbox{\texttt{\mdseries\slshape n}}} vertices and \texttt{3\mbox{\texttt{\mdseries\slshape n}}} undirected edges. The prism graph is a symmetric digraph. The \mbox{\texttt{\mdseries\slshape n}}th prism graph is isomorphic to the graph Cartesian product of the second path
graph and the \mbox{\texttt{\mdseries\slshape n}}th cycle graph, isomorphic to the generalised Petersen graph \texttt{GP(\mbox{\texttt{\mdseries\slshape n}},1)}. If \mbox{\texttt{\mdseries\slshape n}} is odd then the prism graph is isomorphic to the Circulant graph \texttt{Ci(2\mbox{\texttt{\mdseries\slshape n}}, [2,\mbox{\texttt{\mdseries\slshape n}}])}. The prism graph is Hamiltonian.

 See \href{https://mathworld.wolfram.com/PrismGraph.html} {\texttt{https://mathworld.wolfram.com/PrismGraph.html}} for further details.

 If the optional first argument \mbox{\texttt{\mdseries\slshape filt}} is present, then this should specify the category or representation the
digraph being created will belong to. For example, if \mbox{\texttt{\mdseries\slshape filt}} is \texttt{IsMutableDigraph} (\ref{IsMutableDigraph}), then the digraph being created will be mutable, if \mbox{\texttt{\mdseries\slshape filt}} is \texttt{IsImmutableDigraph} (\ref{IsImmutableDigraph}), then the digraph will be immutable. If the optional first argument \mbox{\texttt{\mdseries\slshape filt}} is not present, then \texttt{IsImmutableDigraph} (\ref{IsImmutableDigraph}) is used by default.

 
\begin{Verbatim}[commandchars=!@|,fontsize=\small,frame=single,label=Example]
  !gapprompt@gap>| !gapinput@D := PrismGraph(4);|
  <immutable symmetric digraph with 8 vertices, 24 edges>
  !gapprompt@gap>| !gapinput@IsHamiltonianDigraph(D);|
  true
  !gapprompt@gap>| !gapinput@D := PrismGraph(5);|
  <immutable symmetric digraph with 10 vertices, 30 edges>
  !gapprompt@gap>| !gapinput@IsIsomorphicDigraph(D, CirculantGraph(10, [2, 5]));|
  true
\end{Verbatim}
 }

 

\subsection{\textcolor{Chapter }{StackedPrismGraph}}
\logpage{[ 3, 5, 40 ]}\nobreak
\hyperdef{L}{X817982877B48D5BD}{}
{\noindent\textcolor{FuncColor}{$\triangleright$\enspace\texttt{StackedPrismGraph({\mdseries\slshape [filt, ]n, k})\index{StackedPrismGraph@\texttt{StackedPrismGraph}}
\label{StackedPrismGraph}
}\hfill{\scriptsize (operation)}}\\
\textbf{\indent Returns:\ }
A digraph.



 If \mbox{\texttt{\mdseries\slshape n}} is an integer at least 3 and \mbox{\texttt{\mdseries\slshape k}} is a positive integer then this operation returns the \texttt{(\mbox{\texttt{\mdseries\slshape n}},\mbox{\texttt{\mdseries\slshape k}})}\texttt{\symbol{45}}th \emph{stacked prism graph}. The graph is \mbox{\texttt{\mdseries\slshape k}} concentric \mbox{\texttt{\mdseries\slshape n}}\texttt{\symbol{45}}Cycle graphs connected by spokes. The stacked prism is the
graph Cartesian product of the \mbox{\texttt{\mdseries\slshape n}}th cycle graph and the \mbox{\texttt{\mdseries\slshape k}}th path graph. The graph has \texttt{\mbox{\texttt{\mdseries\slshape n}}\mbox{\texttt{\mdseries\slshape k}}} vertices and \texttt{\mbox{\texttt{\mdseries\slshape n}}(2\mbox{\texttt{\mdseries\slshape k}} \texttt{\symbol{45}} 1)} undirected edges. If \mbox{\texttt{\mdseries\slshape k}} is 1 then the graph is the \mbox{\texttt{\mdseries\slshape n}}th cycle graph, if \mbox{\texttt{\mdseries\slshape k}} is 2 then the graph is the prism graph. 

 See \href{https://mathworld.wolfram.com/StackedPrismGraph.html} {\texttt{https://mathworld.wolfram.com/StackedPrismGraph.html}} for further details.

 If the optional first argument \mbox{\texttt{\mdseries\slshape filt}} is present, then this should specify the category or representation the
digraph being created will belong to. For example, if \mbox{\texttt{\mdseries\slshape filt}} is \texttt{IsMutableDigraph} (\ref{IsMutableDigraph}), then the digraph being created will be mutable, if \mbox{\texttt{\mdseries\slshape filt}} is \texttt{IsImmutableDigraph} (\ref{IsImmutableDigraph}), then the digraph will be immutable. If the optional first argument \mbox{\texttt{\mdseries\slshape filt}} is not present, then \texttt{IsImmutableDigraph} (\ref{IsImmutableDigraph}) is used by default.

 
\begin{Verbatim}[commandchars=!@|,fontsize=\small,frame=single,label=Example]
  !gapprompt@gap>| !gapinput@D := StackedPrismGraph(5, 2);|
  <immutable symmetric digraph with 10 vertices, 30 edges>
  !gapprompt@gap>| !gapinput@IsIsomorphicDigraph(D, PrismGraph(5));|
  true
\end{Verbatim}
 }

 

\subsection{\textcolor{Chapter }{QueensGraph}}
\logpage{[ 3, 5, 41 ]}\nobreak
\hyperdef{L}{X785C3F1F7D690151}{}
{\noindent\textcolor{FuncColor}{$\triangleright$\enspace\texttt{QueensGraph({\mdseries\slshape [filt, ]m, n})\index{QueensGraph@\texttt{QueensGraph}}
\label{QueensGraph}
}\hfill{\scriptsize (operation)}}\\
\noindent\textcolor{FuncColor}{$\triangleright$\enspace\texttt{QueenGraph({\mdseries\slshape [filt, ]m, n})\index{QueenGraph@\texttt{QueenGraph}}
\label{QueenGraph}
}\hfill{\scriptsize (operation)}}\\
\textbf{\indent Returns:\ }
A digraph.



 If \mbox{\texttt{\mdseries\slshape m}} and \mbox{\texttt{\mdseries\slshape n}} are positive integers, then this operation returns the \emph{queen's graph} of an \mbox{\texttt{\mdseries\slshape m}} by \mbox{\texttt{\mdseries\slshape n}} chessboard, as a symmetric digraph. 

 The queen's graph represents all possible moves of the queen chess piece
across a chessboard. An \mbox{\texttt{\mdseries\slshape m}} by \mbox{\texttt{\mdseries\slshape n}} chessboard is a grid of \mbox{\texttt{\mdseries\slshape m}} columns ("files") and \mbox{\texttt{\mdseries\slshape n}} rows ("ranks") that intersect in squares. Orthogonally adjacent squares are
alternately colored light and dark, with the square in the first rank and file
being dark.

 The \texttt{\mbox{\texttt{\mdseries\slshape m}} * \mbox{\texttt{\mdseries\slshape n}}} vertices of the queen's graph can be placed onto the \texttt{\mbox{\texttt{\mdseries\slshape m}} * \mbox{\texttt{\mdseries\slshape n}}} squares of an \mbox{\texttt{\mdseries\slshape m}} by \mbox{\texttt{\mdseries\slshape n}} chessboard, such that two vertices are adjacent in the digraph if and only if
the queen can move between the corresponding squares in a single turn. A legal
queen's move is defined as one which moves the queen to an (orthogonally or
diagonally) adjacent square, or to a square which can be reached through a
series of such moves, with all of the small moves being in the same direction. 

 Note that the \texttt{QueensGraph} is the \texttt{DigraphEdgeUnion} (\ref{DigraphEdgeUnion:for a list of digraphs}) of the \texttt{RooksGraph} (\ref{RooksGraph}) and the \texttt{BishopsGraph} (\ref{BishopsGraph}) of the same dimensions. 

 The chosen correspondence between vertices and chess squares is given by \texttt{DigraphVertexLabels} (\ref{DigraphVertexLabels}). In more detail, the vertices of the digraph are labelled by elements of the
Cartesian product \texttt{[1..\mbox{\texttt{\mdseries\slshape m}}] x [1..\mbox{\texttt{\mdseries\slshape n}}]}, where the first entry indexes the column (file) of the square in the
chessboard, and the second entry indexes the row (rank) of the square. (Note
that the files are traditionally indexed by the lowercase letters of the
alphabet). The vertices are sorted in ascending order, first by row (second
component) and then column (first component). If the optional first argument \mbox{\texttt{\mdseries\slshape filt}} is present, then this should specify the category or representation the
digraph being created will belong to. For example, if \mbox{\texttt{\mdseries\slshape filt}} is \texttt{IsMutableDigraph} (\ref{IsMutableDigraph}), then the digraph being created will be mutable, if \mbox{\texttt{\mdseries\slshape filt}} is \texttt{IsImmutableDigraph} (\ref{IsImmutableDigraph}), then the digraph will be immutable. If the optional first argument \mbox{\texttt{\mdseries\slshape filt}} is not present, then \texttt{IsImmutableDigraph} (\ref{IsImmutableDigraph}) is used by default.

 
\begin{Verbatim}[commandchars=!@|,fontsize=\small,frame=single,label=Example]
  !gapprompt@gap>| !gapinput@QueensGraph(2, 5);|
  <immutable connected symmetric digraph with 10 vertices, 66 edges>
  !gapprompt@gap>| !gapinput@D := QueensGraph(4, 3);|
  <immutable connected symmetric digraph with 12 vertices, 92 edges>
  !gapprompt@gap>| !gapinput@IsRegularDigraph(D);|
  false
  !gapprompt@gap>| !gapinput@QueensGraph(6, 9) =|
  !gapprompt@>| !gapinput@   DigraphEdgeUnion(RooksGraph(6, 9), BishopsGraph(6, 9));|
  true
\end{Verbatim}
 }

 

\subsection{\textcolor{Chapter }{RooksGraph}}
\logpage{[ 3, 5, 42 ]}\nobreak
\hyperdef{L}{X7A6DD11881874F51}{}
{\noindent\textcolor{FuncColor}{$\triangleright$\enspace\texttt{RooksGraph({\mdseries\slshape [filt, ]m, n})\index{RooksGraph@\texttt{RooksGraph}}
\label{RooksGraph}
}\hfill{\scriptsize (operation)}}\\
\noindent\textcolor{FuncColor}{$\triangleright$\enspace\texttt{RookGraph({\mdseries\slshape [filt, ]m, n})\index{RookGraph@\texttt{RookGraph}}
\label{RookGraph}
}\hfill{\scriptsize (operation)}}\\
\textbf{\indent Returns:\ }
A digraph.



 If \mbox{\texttt{\mdseries\slshape m}} and \mbox{\texttt{\mdseries\slshape n}} are positive integers, then this operation returns the \emph{rook's graph} of an \mbox{\texttt{\mdseries\slshape m}} by \mbox{\texttt{\mdseries\slshape n}} chessboard, as a symmetric digraph. 

 A rook's graph represents all possible moves of the rook chess piece across a
chessboard. An \mbox{\texttt{\mdseries\slshape m}} by \mbox{\texttt{\mdseries\slshape n}} chessboard is a grid of \mbox{\texttt{\mdseries\slshape m}} columns ("files") and \mbox{\texttt{\mdseries\slshape n}} rows ("ranks") that intersect in squares. Orthogonally adjacent squares are
alternately colored light and dark, with the square in the first rank and file
being dark.

 The \texttt{\mbox{\texttt{\mdseries\slshape m}} * \mbox{\texttt{\mdseries\slshape n}}} vertices of the rook's graph can be placed onto the \texttt{\mbox{\texttt{\mdseries\slshape m}} * \mbox{\texttt{\mdseries\slshape n}}} squares of an \mbox{\texttt{\mdseries\slshape m}} by \mbox{\texttt{\mdseries\slshape n}} chessboard, such that two vertices are adjacent in the digraph if and only if
a rook can move between the corresponding squares in a single turn. A legal
rook's move is defined as one which moves the rook to an orthogonally adjacent
square, or to a square which can be reached through a series of such moves,
with all of the small moves being in the same direction. 

 The chosen correspondence between vertices and chess squares is given by \texttt{DigraphVertexLabels} (\ref{DigraphVertexLabels}). In more detail, the vertices of the digraph are labelled by elements of the
Cartesian product \texttt{[1..\mbox{\texttt{\mdseries\slshape m}}] x [1..\mbox{\texttt{\mdseries\slshape n}}]}, where the first entry indexes the column (file) of the square in the
chessboard, and the second entry indexes the row (rank) of the square. (Note
that the files are traditionally indexed by the lowercase letters of the
alphabet). The vertices are sorted in ascending order, first by row (second
component) and then column (first component). See \href{https://en.wikipedia.org/wiki/Rook's_graph} {Wikipedia} for further information. See also \texttt{DigraphCartesianProduct} (\ref{DigraphCartesianProduct:for a list of digraphs}) and \texttt{LineDigraph} (\ref{LineDigraph}). 

 If the optional first argument \mbox{\texttt{\mdseries\slshape filt}} is present, then this should specify the category or representation the
digraph being created will belong to. For example, if \mbox{\texttt{\mdseries\slshape filt}} is \texttt{IsMutableDigraph} (\ref{IsMutableDigraph}), then the digraph being created will be mutable, if \mbox{\texttt{\mdseries\slshape filt}} is \texttt{IsImmutableDigraph} (\ref{IsImmutableDigraph}), then the digraph will be immutable. If the optional first argument \mbox{\texttt{\mdseries\slshape filt}} is not present, then \texttt{IsImmutableDigraph} (\ref{IsImmutableDigraph}) is used by default.

 
\begin{Verbatim}[commandchars=!@|,fontsize=\small,frame=single,label=Example]
  !gapprompt@gap>| !gapinput@D := RooksGraph(7, 4);|
  <immutable connected regular symmetric digraph with 28 vertices, 252 e\
  dges>
  !gapprompt@gap>| !gapinput@RooksGraph(1, 8);|
  <immutable connected regular symmetric digraph with 8 vertices, 56 edg\
  es>
\end{Verbatim}
 }

 

\subsection{\textcolor{Chapter }{SquareGridGraph}}
\logpage{[ 3, 5, 43 ]}\nobreak
\hyperdef{L}{X8404987F849D7CF2}{}
{\noindent\textcolor{FuncColor}{$\triangleright$\enspace\texttt{SquareGridGraph({\mdseries\slshape [filt, ]n, k})\index{SquareGridGraph@\texttt{SquareGridGraph}}
\label{SquareGridGraph}
}\hfill{\scriptsize (operation)}}\\
\noindent\textcolor{FuncColor}{$\triangleright$\enspace\texttt{GridGraph({\mdseries\slshape [filt, ]n, k})\index{GridGraph@\texttt{GridGraph}}
\label{GridGraph}
}\hfill{\scriptsize (operation)}}\\
\textbf{\indent Returns:\ }
A digraph.



 If \mbox{\texttt{\mdseries\slshape n}} and \mbox{\texttt{\mdseries\slshape k}} are positive integers, then this operation returns a square grid graph of
dimension \mbox{\texttt{\mdseries\slshape n}} by \mbox{\texttt{\mdseries\slshape k}}. 

 A \emph{square grid graph} of dimension \mbox{\texttt{\mdseries\slshape n}} by \mbox{\texttt{\mdseries\slshape k}} is the \texttt{DigraphCartesianProduct} (\ref{DigraphCartesianProduct:for a positive number of digraphs}) of the symmetric closures of the chain digraphs with \mbox{\texttt{\mdseries\slshape n}} and \mbox{\texttt{\mdseries\slshape k}} vertices; see \texttt{DigraphSymmetricClosure} (\ref{DigraphSymmetricClosure}) and \texttt{ChainDigraph} (\ref{ChainDigraph}). 

 In particular, the \texttt{\mbox{\texttt{\mdseries\slshape n}} * \mbox{\texttt{\mdseries\slshape k}}} vertices can be arranged into an \mbox{\texttt{\mdseries\slshape n}} by \mbox{\texttt{\mdseries\slshape k}} grid such that two vertices are adjacent in the digraph if and only if they
are orthogonally adjacent in the grid. The correspondence between vertices and
grid positions is given by \texttt{DigraphVertexLabels} (\ref{DigraphVertexLabels}). 

 See \href{https://en.wikipedia.org/wiki/Lattice_graph} {\texttt{https://en.wikipedia.org/wiki/Lattice{\textunderscore}graph}} for more information. 

 If the optional first argument \mbox{\texttt{\mdseries\slshape filt}} is present, then this should specify the category or representation the
digraph being created will belong to. For example, if \mbox{\texttt{\mdseries\slshape filt}} is \texttt{IsMutableDigraph} (\ref{IsMutableDigraph}), then the digraph being created will be mutable, if \mbox{\texttt{\mdseries\slshape filt}} is \texttt{IsImmutableDigraph} (\ref{IsImmutableDigraph}), then the digraph will be immutable. If the optional first argument \mbox{\texttt{\mdseries\slshape filt}} is not present, then \texttt{IsImmutableDigraph} (\ref{IsImmutableDigraph}) is used by default.

 
\begin{Verbatim}[commandchars=!@|,fontsize=\small,frame=single,label=Example]
  !gapprompt@gap>| !gapinput@SquareGridGraph(5, 5);|
  <immutable planar connected bipartite symmetric digraph with bicompone\
  nt sizes 13 and 12>
  !gapprompt@gap>| !gapinput@GridGraph(IsMutable, 3, 4);|
  <mutable digraph with 12 vertices, 34 edges>
\end{Verbatim}
 }

 

\subsection{\textcolor{Chapter }{TriangularGridGraph}}
\logpage{[ 3, 5, 44 ]}\nobreak
\hyperdef{L}{X8234361278E8816F}{}
{\noindent\textcolor{FuncColor}{$\triangleright$\enspace\texttt{TriangularGridGraph({\mdseries\slshape [filt, ]n, k})\index{TriangularGridGraph@\texttt{TriangularGridGraph}}
\label{TriangularGridGraph}
}\hfill{\scriptsize (operation)}}\\
\textbf{\indent Returns:\ }
A digraph.



 If \mbox{\texttt{\mdseries\slshape n}} and \mbox{\texttt{\mdseries\slshape k}} are positive integers, then this operation returns a triangular grid graph of
dimension \mbox{\texttt{\mdseries\slshape n}} by \mbox{\texttt{\mdseries\slshape k}}. 

 A \emph{triangular grid graph} of dimension \mbox{\texttt{\mdseries\slshape n}} by \mbox{\texttt{\mdseries\slshape k}} is a symmetric digraph constructed from the \texttt{SquareGridGraph} (\ref{SquareGridGraph}) of the same dimensions, where additionally two vertices are adjacent in the
digraph if they are diagonally adjacent in the grid, on a particular one of
the diagonals. The correspondence between vertices and grid positions is given
by \texttt{DigraphVertexLabels} (\ref{DigraphVertexLabels}). More specifically, the particular diagonal is the one such that, the
vertices corresponding to the grid positions \texttt{[2,1]} and \texttt{[1,2]} are adjacent (if they exist), but those corresponding to \texttt{[1,1]} and \texttt{[2,2]} are not. 

 See \href{https://en.wikipedia.org/wiki/Lattice_graph#Other_kinds} {\texttt{https://en.wikipedia.org/wiki/Lattice{\textunderscore}graph\#Other{\textunderscore}kinds}} for more information. 

 If the optional first argument \mbox{\texttt{\mdseries\slshape filt}} is present, then this should specify the category or representation the
digraph being created will belong to. For example, if \mbox{\texttt{\mdseries\slshape filt}} is \texttt{IsMutableDigraph} (\ref{IsMutableDigraph}), then the digraph being created will be mutable, if \mbox{\texttt{\mdseries\slshape filt}} is \texttt{IsImmutableDigraph} (\ref{IsImmutableDigraph}), then the digraph will be immutable. If the optional first argument \mbox{\texttt{\mdseries\slshape filt}} is not present, then \texttt{IsImmutableDigraph} (\ref{IsImmutableDigraph}) is used by default.

 
\begin{Verbatim}[commandchars=!@|,fontsize=\small,frame=single,label=Example]
  !gapprompt@gap>| !gapinput@TriangularGridGraph(3, 3);|
  <immutable planar connected symmetric digraph with 9 vertices, 32 edge\
  s>
  !gapprompt@gap>| !gapinput@TriangularGridGraph(IsMutable, 3, 3);|
  <mutable digraph with 9 vertices, 32 edges>
\end{Verbatim}
 }

 

\subsection{\textcolor{Chapter }{StarGraph}}
\logpage{[ 3, 5, 45 ]}\nobreak
\hyperdef{L}{X78F78C077CBAE1EC}{}
{\noindent\textcolor{FuncColor}{$\triangleright$\enspace\texttt{StarGraph({\mdseries\slshape [filt, ]k})\index{StarGraph@\texttt{StarGraph}}
\label{StarGraph}
}\hfill{\scriptsize (operation)}}\\
\textbf{\indent Returns:\ }
A digraph.



 If \mbox{\texttt{\mdseries\slshape k}} is a positive integer, then this operation returns the \emph{star graph} with \mbox{\texttt{\mdseries\slshape k}} vertices, which is the undirected tree in which vertex \texttt{1} is adjacent to all other vertices. If \mbox{\texttt{\mdseries\slshape k}} is at least \texttt{2}, then this is the complete bipartite digraph with bicomponents \texttt{[1]} and \texttt{[2 .. \mbox{\texttt{\mdseries\slshape k}}]}. 

 See \texttt{IsUndirectedTree} (\ref{IsUndirectedTree}), \texttt{IsCompleteBipartiteDigraph} (\ref{IsCompleteBipartiteDigraph}), and \texttt{DigraphBicomponents} (\ref{DigraphBicomponents}). 

 If the optional first argument \mbox{\texttt{\mdseries\slshape filt}} is present, then this should specify the category or representation the
digraph being created will belong to. For example, if \mbox{\texttt{\mdseries\slshape filt}} is \texttt{IsMutableDigraph} (\ref{IsMutableDigraph}), then the digraph being created will be mutable, if \mbox{\texttt{\mdseries\slshape filt}} is \texttt{IsImmutableDigraph} (\ref{IsImmutableDigraph}), then the digraph will be immutable. If the optional first argument \mbox{\texttt{\mdseries\slshape filt}} is not present, then \texttt{IsImmutableDigraph} (\ref{IsImmutableDigraph}) is used by default.

 
\begin{Verbatim}[commandchars=!@|,fontsize=\small,frame=single,label=Example]
  !gapprompt@gap>| !gapinput@StarGraph(IsMutable, 10);|
  <mutable digraph with 10 vertices, 18 edges>
  !gapprompt@gap>| !gapinput@StarGraph(5);|
  <immutable complete bipartite digraph with bicomponent sizes 1 and 4>
  !gapprompt@gap>| !gapinput@IsSymmetricDigraph(StarGraph(3));|
  true
  !gapprompt@gap>| !gapinput@IsUndirectedTree(StarGraph(3));|
  true
\end{Verbatim}
 }

 

\subsection{\textcolor{Chapter }{TadpoleGraph}}
\logpage{[ 3, 5, 46 ]}\nobreak
\hyperdef{L}{X81D59C4D809F4AD3}{}
{\noindent\textcolor{FuncColor}{$\triangleright$\enspace\texttt{TadpoleGraph({\mdseries\slshape [filt, ]m, n})\index{TadpoleGraph@\texttt{TadpoleGraph}}
\label{TadpoleGraph}
}\hfill{\scriptsize (operation)}}\\
\textbf{\indent Returns:\ }
A digraph.



 The \emph{tadpole graph} is the symmetric closure of the disjoint union of the cycle digraph on \texttt{[1..\mbox{\texttt{\mdseries\slshape m}}]} (the 'head' of the tadpole) and the chain digraph on \texttt{[\mbox{\texttt{\mdseries\slshape m}}+1..\mbox{\texttt{\mdseries\slshape m}}+\mbox{\texttt{\mdseries\slshape n}}]} (the 'tail' of the tadpole), along with the additional edges \texttt{[1, \mbox{\texttt{\mdseries\slshape m}}+1]} and \texttt{[1, \mbox{\texttt{\mdseries\slshape m}}+1]} which connect the 'head' and the 'tail'. For more details on the tadpole graph
please refer to \href{https://en.wikipedia.org/wiki/Tadpole_graph} {\texttt{https://en.wikipedia.org/wiki/Tadpole{\textunderscore}graph}}.

 See \texttt{DigraphSymmetricClosure} (\ref{DigraphSymmetricClosure}), \texttt{DigraphDisjointUnion} (\ref{DigraphDisjointUnion:for a list of digraphs}), \texttt{CycleDigraph} (\ref{CycleDigraph}), and \texttt{ChainDigraph} (\ref{ChainDigraph}). 

 If the optional first argument \mbox{\texttt{\mdseries\slshape filt}} is present, then this should specify the category or representation the
digraph being created will belong to. For example, if \mbox{\texttt{\mdseries\slshape filt}} is \texttt{IsMutableDigraph} (\ref{IsMutableDigraph}), then the digraph being created will be mutable, if \mbox{\texttt{\mdseries\slshape filt}} is \texttt{IsImmutableDigraph} (\ref{IsImmutableDigraph}), then the digraph will be immutable. If the optional first argument \mbox{\texttt{\mdseries\slshape filt}} is not present, then \texttt{IsImmutableDigraph} (\ref{IsImmutableDigraph}) is used by default.

 
\begin{Verbatim}[commandchars=!@|,fontsize=\small,frame=single,label=Example]
  !gapprompt@gap>| !gapinput@TadpoleGraph(10, 15);|
  <immutable symmetric digraph with 25 vertices, 50 edges>
  !gapprompt@gap>| !gapinput@TadpoleGraph(IsMutableDigraph, 5, 6);|
  <mutable digraph with 11 vertices, 22 edges>
  !gapprompt@gap>| !gapinput@IsSymmetricDigraph(TadpoleGraph(3, 5));|
  true
\end{Verbatim}
 }

 

\subsection{\textcolor{Chapter }{WalshHadamardGraph}}
\logpage{[ 3, 5, 47 ]}\nobreak
\hyperdef{L}{X7BFA33067F83F8B0}{}
{\noindent\textcolor{FuncColor}{$\triangleright$\enspace\texttt{WalshHadamardGraph({\mdseries\slshape [filt, ]n})\index{WalshHadamardGraph@\texttt{WalshHadamardGraph}}
\label{WalshHadamardGraph}
}\hfill{\scriptsize (operation)}}\\
\textbf{\indent Returns:\ }
A digraph.



 If \mbox{\texttt{\mdseries\slshape n}} is a positive integer at least 1, then this operation returns the \emph{Hadamard graph} constructed from the \mbox{\texttt{\mdseries\slshape n}}th Hadamard matrix (of dimension \texttt{2\texttt{\symbol{94}}\mbox{\texttt{\mdseries\slshape n}}}) as constructed by Joseph Walsh. A Hadamard matrix is a square matrix with
entries either $1$ or $-1$, such that all the rows are mutually orthogonal. The \mbox{\texttt{\mdseries\slshape n}}th Walsh Hadamard graph is a graph on \texttt{4\mbox{\texttt{\mdseries\slshape n}}} matrices split into four categories $r_i+, r_i-, c_i+, c_i-$. If $h_ij$ are the elements of the \mbox{\texttt{\mdseries\slshape n}}th Walsh matrix, then if $h_ij = 1$ then $(r_i+, c_j+)$ and $(r_i-, c_j-)$ are edges, if $h_ij = -1$ then $(r_i+, c_j-)$ and $(r_i-, c_j+)$ are edges. Walsh Hadamard graphs are distance transitive and distance regular.

 See \href{https://mathworld.wolfram.com/HadamardGraph.html} {\texttt{https://mathworld.wolfram.com/HadamardGraph.html}} for further details.

 If the optional first argument \mbox{\texttt{\mdseries\slshape filt}} is present, then this should specify the category or representation the
digraph being created will belong to. For example, if \mbox{\texttt{\mdseries\slshape filt}} is \texttt{IsMutableDigraph} (\ref{IsMutableDigraph}), then the digraph being created will be mutable, if \mbox{\texttt{\mdseries\slshape filt}} is \texttt{IsImmutableDigraph} (\ref{IsImmutableDigraph}), then the digraph will be immutable. If the optional first argument \mbox{\texttt{\mdseries\slshape filt}} is not present, then \texttt{IsImmutableDigraph} (\ref{IsImmutableDigraph}) is used by default.

 
\begin{Verbatim}[commandchars=!@|,fontsize=\small,frame=single,label=Example]
  !gapprompt@gap>| !gapinput@D := WalshHadamardGraph(5);|
  <immutable symmetric digraph with 64 vertices, 1024 edges>
  !gapprompt@gap>| !gapinput@IsDistanceRegularDigraph(D);|
  true
\end{Verbatim}
 }

 

\subsection{\textcolor{Chapter }{WebGraph}}
\logpage{[ 3, 5, 48 ]}\nobreak
\hyperdef{L}{X84F3B70A82EEE780}{}
{\noindent\textcolor{FuncColor}{$\triangleright$\enspace\texttt{WebGraph({\mdseries\slshape [filt, ]n})\index{WebGraph@\texttt{WebGraph}}
\label{WebGraph}
}\hfill{\scriptsize (operation)}}\\
\textbf{\indent Returns:\ }
A digraph.



 If \mbox{\texttt{\mdseries\slshape n}} is an integer at least 3 then this operation returns the \mbox{\texttt{\mdseries\slshape n}}th \emph{web graph}. The web graph is the \texttt{(\mbox{\texttt{\mdseries\slshape n}},3)}\texttt{\symbol{45}}th stacked prism graph with the edges of the outer cycle
removed. The graph has \texttt{3\mbox{\texttt{\mdseries\slshape n}}} vertices and \texttt{4\mbox{\texttt{\mdseries\slshape n}}} undirected edges.

 See \href{https://mathworld.wolfram.com/WebGraph.html} {\texttt{https://mathworld.wolfram.com/WebGraph.html}} for further details.

 If the optional first argument \mbox{\texttt{\mdseries\slshape filt}} is present, then this should specify the category or representation the
digraph being created will belong to. For example, if \mbox{\texttt{\mdseries\slshape filt}} is \texttt{IsMutableDigraph} (\ref{IsMutableDigraph}), then the digraph being created will be mutable, if \mbox{\texttt{\mdseries\slshape filt}} is \texttt{IsImmutableDigraph} (\ref{IsImmutableDigraph}), then the digraph will be immutable. If the optional first argument \mbox{\texttt{\mdseries\slshape filt}} is not present, then \texttt{IsImmutableDigraph} (\ref{IsImmutableDigraph}) is used by default.

 
\begin{Verbatim}[commandchars=!@|,fontsize=\small,frame=single,label=Example]
  !gapprompt@gap>| !gapinput@D := WebGraph(5);|
  <immutable symmetric digraph with 15 vertices, 40 edges>
\end{Verbatim}
 }

 

\subsection{\textcolor{Chapter }{WheelGraph}}
\logpage{[ 3, 5, 49 ]}\nobreak
\hyperdef{L}{X817EA60D828A765E}{}
{\noindent\textcolor{FuncColor}{$\triangleright$\enspace\texttt{WheelGraph({\mdseries\slshape [filt, ]n})\index{WheelGraph@\texttt{WheelGraph}}
\label{WheelGraph}
}\hfill{\scriptsize (operation)}}\\
\textbf{\indent Returns:\ }
A digraph.



 If \mbox{\texttt{\mdseries\slshape n}} is a positive integer at least 4, then this operation returns the \mbox{\texttt{\mdseries\slshape n}}th \emph{wheel graph}. The Wheel graph is formed from an \texttt{\mbox{\texttt{\mdseries\slshape n}}\texttt{\symbol{45}}1} cycle graph with a single additional vertex adjacent to all vertices of the
cycle. The graph has \mbox{\texttt{\mdseries\slshape n}} vertices and \texttt{2(\mbox{\texttt{\mdseries\slshape n}}\texttt{\symbol{45}}1)} edges. Wheel graphs are the skeletons of \texttt{\mbox{\texttt{\mdseries\slshape n}}\texttt{\symbol{45}}1} pyramids, and are self\texttt{\symbol{45}}dual. If \mbox{\texttt{\mdseries\slshape n}} is odd, then the chromatic number of the wheel graph is 3 and the Wheel graph
is perfect, and it is 4 otherwise. The wheel graph is also Hamiltonian and
planar.

 See \href{https://mathworld.wolfram.com/WheelGraph.html} {\texttt{https://mathworld.wolfram.com/WheelGraph.html}} for further details.

 If the optional first argument \mbox{\texttt{\mdseries\slshape filt}} is present, then this should specify the category or representation the
digraph being created will belong to. For example, if \mbox{\texttt{\mdseries\slshape filt}} is \texttt{IsMutableDigraph} (\ref{IsMutableDigraph}), then the digraph being created will be mutable, if \mbox{\texttt{\mdseries\slshape filt}} is \texttt{IsImmutableDigraph} (\ref{IsImmutableDigraph}), then the digraph will be immutable. If the optional first argument \mbox{\texttt{\mdseries\slshape filt}} is not present, then \texttt{IsImmutableDigraph} (\ref{IsImmutableDigraph}) is used by default.

 
\begin{Verbatim}[commandchars=!@|,fontsize=\small,frame=single,label=Example]
  !gapprompt@gap>| !gapinput@D := WheelGraph(8);;|
  !gapprompt@gap>| !gapinput@ChromaticNumber(D);|
  4
  !gapprompt@gap>| !gapinput@IsHamiltonianDigraph(D);|
  true
\end{Verbatim}
 }

 

\subsection{\textcolor{Chapter }{WindmillGraph}}
\logpage{[ 3, 5, 50 ]}\nobreak
\hyperdef{L}{X7BE44CA27AA5F8DB}{}
{\noindent\textcolor{FuncColor}{$\triangleright$\enspace\texttt{WindmillGraph({\mdseries\slshape [filt, ]n, m})\index{WindmillGraph@\texttt{WindmillGraph}}
\label{WindmillGraph}
}\hfill{\scriptsize (operation)}}\\
\textbf{\indent Returns:\ }
A digraph.



 If \mbox{\texttt{\mdseries\slshape n}} and \mbox{\texttt{\mdseries\slshape m}} are integers greater than 1 then this operation returns the \texttt{(\mbox{\texttt{\mdseries\slshape n}}, \mbox{\texttt{\mdseries\slshape m}})}\texttt{\symbol{45}}th \emph{windmill graph}. The windmill graph is formed from \mbox{\texttt{\mdseries\slshape m}} copies of the complete graph on \mbox{\texttt{\mdseries\slshape n}} vertices with one shared vertex. The graph has \texttt{\mbox{\texttt{\mdseries\slshape m}}(\mbox{\texttt{\mdseries\slshape n}} \texttt{\symbol{45}} 1) + 1} vertices and \texttt{\mbox{\texttt{\mdseries\slshape m}} * \mbox{\texttt{\mdseries\slshape n}} * (\mbox{\texttt{\mdseries\slshape n}} \texttt{\symbol{45}} 1) / 2} undirected edges. The windmill graph has chromatic number \mbox{\texttt{\mdseries\slshape n}} and diameter 2.

 See \href{https://mathworld.wolfram.com/WindmillGraph.html} {\texttt{https://mathworld.wolfram.com/WindmillGraph.html}} for further details.

 If the optional first argument \mbox{\texttt{\mdseries\slshape filt}} is present, then this should specify the category or representation the
digraph being created will belong to. For example, if \mbox{\texttt{\mdseries\slshape filt}} is \texttt{IsMutableDigraph} (\ref{IsMutableDigraph}), then the digraph being created will be mutable, if \mbox{\texttt{\mdseries\slshape filt}} is \texttt{IsImmutableDigraph} (\ref{IsImmutableDigraph}), then the digraph will be immutable. If the optional first argument \mbox{\texttt{\mdseries\slshape filt}} is not present, then \texttt{IsImmutableDigraph} (\ref{IsImmutableDigraph}) is used by default.

 
\begin{Verbatim}[commandchars=!@|,fontsize=\small,frame=single,label=Example]
  !gapprompt@gap>| !gapinput@D := WindmillGraph(4, 3);|
  <immutable symmetric digraph with 10 vertices, 36 edges>
  !gapprompt@gap>| !gapinput@ChromaticNumber(D);|
  4
\end{Verbatim}
 }

 }

 }

 
\chapter{\textcolor{Chapter }{Operators}}\label{Operators}
\logpage{[ 4, 0, 0 ]}
\hyperdef{L}{X7AD6F77E7D95C996}{}
{
 
\section{\textcolor{Chapter }{Operators for digraphs}}\logpage{[ 4, 1, 0 ]}
\hyperdef{L}{X84E0B5B88358C96B}{}
{
 
\begin{description}
\item[{\texttt{\mbox{\texttt{\mdseries\slshape digraph1}} = \mbox{\texttt{\mdseries\slshape digraph2}}}}]  \index{=@\texttt{=} (for digraphs)} returns \texttt{true} if \mbox{\texttt{\mdseries\slshape digraph1}} and \mbox{\texttt{\mdseries\slshape digraph2}} have the same vertices, and \texttt{DigraphEdges(\mbox{\texttt{\mdseries\slshape digraph1}}) = DigraphEdges(\mbox{\texttt{\mdseries\slshape digraph2}})}, up to some re\texttt{\symbol{45}}ordering of the edge lists. 

 Note that this operator does not compare the vertex labels of \mbox{\texttt{\mdseries\slshape digraph1}} and \mbox{\texttt{\mdseries\slshape digraph2}}. 
\item[{\texttt{\mbox{\texttt{\mdseries\slshape digraph1}} {\textless} \mbox{\texttt{\mdseries\slshape digraph2}}}}]  \index{<@\texttt{{\textless}} (for digraphs)} This operator returns \texttt{true} if one of the following holds: 
\begin{itemize}
\item  The number $n_1$ of vertices in \mbox{\texttt{\mdseries\slshape digraph1}} is less than the number $n_2$ of vertices in \mbox{\texttt{\mdseries\slshape digraph2}}; 
\item  $n_1 = n_2$, and the number $m_1$ of edges in \mbox{\texttt{\mdseries\slshape digraph1}} is less than the number $m_2$ of edges in \mbox{\texttt{\mdseries\slshape digraph2}}; 
\item  $n_1 = n_2$, $m_1 = m_2$, and \texttt{DigraphEdges(\mbox{\texttt{\mdseries\slshape digraph1}})} is less than \texttt{DigraphEdges(\mbox{\texttt{\mdseries\slshape digraph2}})} after having both of these sets have been sorted with respect to the
lexicographical order. 
\end{itemize}
 
\end{description}
 

\subsection{\textcolor{Chapter }{IsSubdigraph}}
\logpage{[ 4, 1, 1 ]}\nobreak
\hyperdef{L}{X829B911D7EFD2D85}{}
{\noindent\textcolor{FuncColor}{$\triangleright$\enspace\texttt{IsSubdigraph({\mdseries\slshape super, sub})\index{IsSubdigraph@\texttt{IsSubdigraph}}
\label{IsSubdigraph}
}\hfill{\scriptsize (operation)}}\\
\textbf{\indent Returns:\ }
\texttt{true} or \texttt{false}.



 If \mbox{\texttt{\mdseries\slshape super}} and \mbox{\texttt{\mdseries\slshape sub}} are digraphs, then this operation returns \texttt{true} if \mbox{\texttt{\mdseries\slshape sub}} is a subdigraph of \mbox{\texttt{\mdseries\slshape super}}, and \texttt{false} if it is not. 

 A digraph \mbox{\texttt{\mdseries\slshape sub}} is a \emph{subdigraph} of a digraph \mbox{\texttt{\mdseries\slshape super}} if \mbox{\texttt{\mdseries\slshape sub}} and \mbox{\texttt{\mdseries\slshape super}} share the same number of vertices, and the collection of edges of \mbox{\texttt{\mdseries\slshape super}} (including repeats) contains the collection of edges of \mbox{\texttt{\mdseries\slshape sub}} (including repeats). 

 In other words, \mbox{\texttt{\mdseries\slshape sub}} is a subdigraph of \mbox{\texttt{\mdseries\slshape super}} if and only if \texttt{DigraphNrVertices(\mbox{\texttt{\mdseries\slshape sub}}) = DigraphNrVertices(\mbox{\texttt{\mdseries\slshape super}})}, and for each pair of vertices \texttt{i} and \texttt{j}, there are at least as many edges of the form \texttt{[i, j]} in \mbox{\texttt{\mdseries\slshape super}} as there are in \mbox{\texttt{\mdseries\slshape sub}}. 

 
\begin{Verbatim}[commandchars=!@|,fontsize=\small,frame=single,label=Example]
  !gapprompt@gap>| !gapinput@g := Digraph([[2, 3], [1], [2, 3]]);|
  <immutable digraph with 3 vertices, 5 edges>
  !gapprompt@gap>| !gapinput@h := Digraph([[2, 3], [], [2]]);|
  <immutable digraph with 3 vertices, 3 edges>
  !gapprompt@gap>| !gapinput@IsSubdigraph(g, h);|
  true
  !gapprompt@gap>| !gapinput@IsSubdigraph(h, g);|
  false
  !gapprompt@gap>| !gapinput@IsSubdigraph(CompleteDigraph(4), CycleDigraph(4));|
  true
  !gapprompt@gap>| !gapinput@IsSubdigraph(CycleDigraph(4), ChainDigraph(4));|
  true
  !gapprompt@gap>| !gapinput@g := Digraph([[2, 2], [1]]);|
  <immutable multidigraph with 2 vertices, 3 edges>
  !gapprompt@gap>| !gapinput@h := Digraph([[2], [1]]);|
  <immutable digraph with 2 vertices, 2 edges>
  !gapprompt@gap>| !gapinput@IsSubdigraph(g, h);|
  true
  !gapprompt@gap>| !gapinput@IsSubdigraph(h, g);|
  false
\end{Verbatim}
 }

 

\subsection{\textcolor{Chapter }{IsUndirectedSpanningTree}}
\logpage{[ 4, 1, 2 ]}\nobreak
\hyperdef{L}{X833C3299787E2309}{}
{\noindent\textcolor{FuncColor}{$\triangleright$\enspace\texttt{IsUndirectedSpanningTree({\mdseries\slshape super, sub})\index{IsUndirectedSpanningTree@\texttt{IsUndirectedSpanningTree}}
\label{IsUndirectedSpanningTree}
}\hfill{\scriptsize (operation)}}\\
\noindent\textcolor{FuncColor}{$\triangleright$\enspace\texttt{IsUndirectedSpanningForest({\mdseries\slshape super, sub})\index{IsUndirectedSpanningForest@\texttt{IsUndirectedSpanningForest}}
\label{IsUndirectedSpanningForest}
}\hfill{\scriptsize (operation)}}\\
\textbf{\indent Returns:\ }
\texttt{true} or \texttt{false}.



 The operation \texttt{IsUndirectedSpanningTree} returns \texttt{true} if the digraph \mbox{\texttt{\mdseries\slshape sub}} is an undirected spanning tree of the digraph \mbox{\texttt{\mdseries\slshape super}}, and the operation \texttt{IsUndirectedSpanningForest} returns \texttt{true} if the digraph \mbox{\texttt{\mdseries\slshape sub}} is an undirected spanning forest of the digraph \mbox{\texttt{\mdseries\slshape super}}. 

 An \emph{undirected spanning tree} of a digraph \mbox{\texttt{\mdseries\slshape super}} is a subdigraph of \mbox{\texttt{\mdseries\slshape super}} that is an undirected tree (see \texttt{IsSubdigraph} (\ref{IsSubdigraph}) and \texttt{IsUndirectedTree} (\ref{IsUndirectedTree})). Note that a digraph whose \texttt{MaximalSymmetricSubdigraph} (\ref{MaximalSymmetricSubdigraph}) is not connected has no undirected spanning trees (see \texttt{IsConnectedDigraph} (\ref{IsConnectedDigraph})). 

 An \emph{undirected spanning forest} of a digraph \mbox{\texttt{\mdseries\slshape super}} is a subdigraph of \mbox{\texttt{\mdseries\slshape super}} that is an undirected forest (see \texttt{IsSubdigraph} (\ref{IsSubdigraph}) and \texttt{IsUndirectedForest} (\ref{IsUndirectedForest})), and is not contained in any larger such subdigraph of \mbox{\texttt{\mdseries\slshape super}}. Equivalently, an undirected spanning forest is a subdigraph of \mbox{\texttt{\mdseries\slshape super}} whose connected components coincide with those of the \texttt{MaximalSymmetricSubdigraph} (\ref{MaximalSymmetricSubdigraph}) of \mbox{\texttt{\mdseries\slshape super}} (see \texttt{DigraphConnectedComponents} (\ref{DigraphConnectedComponents})). 

 Note that an undirected spanning tree is an undirected spanning forest that is
connected. 
\begin{Verbatim}[commandchars=!@|,fontsize=\small,frame=single,label=Example]
  !gapprompt@gap>| !gapinput@D := CompleteDigraph(4);|
  <immutable complete digraph with 4 vertices>
  !gapprompt@gap>| !gapinput@tree := Digraph([[3], [4], [1, 4], [2, 3]]);|
  <immutable digraph with 4 vertices, 6 edges>
  !gapprompt@gap>| !gapinput@IsSubdigraph(D, tree) and IsUndirectedTree(tree);|
  true
  !gapprompt@gap>| !gapinput@IsUndirectedSpanningTree(D, tree);|
  true
  !gapprompt@gap>| !gapinput@forest := EmptyDigraph(4);|
  <immutable empty digraph with 4 vertices>
  !gapprompt@gap>| !gapinput@IsSubdigraph(D, forest) and IsUndirectedForest(forest);|
  true
  !gapprompt@gap>| !gapinput@IsUndirectedSpanningForest(D, forest);|
  false
  !gapprompt@gap>| !gapinput@IsSubdigraph(tree, forest);|
  true
  !gapprompt@gap>| !gapinput@D := DigraphDisjointUnion(CycleDigraph(2), CycleDigraph(2));|
  <immutable digraph with 4 vertices, 4 edges>
  !gapprompt@gap>| !gapinput@IsUndirectedTree(D);|
  false
  !gapprompt@gap>| !gapinput@IsUndirectedForest(D) and IsUndirectedSpanningForest(D, D);|
  true
\end{Verbatim}
 }

 }

 }

 
\chapter{\textcolor{Chapter }{Attributes and operations}}\label{Attributes and operations}
\logpage{[ 5, 0, 0 ]}
\hyperdef{L}{X8739F6CD78C90B14}{}
{
 
\section{\textcolor{Chapter }{Vertices and edges}}\logpage{[ 5, 1, 0 ]}
\hyperdef{L}{X7E814B6478F7D015}{}
{
 

\subsection{\textcolor{Chapter }{DigraphVertices}}
\logpage{[ 5, 1, 1 ]}\nobreak
\hyperdef{L}{X7C45F7D878D896AC}{}
{\noindent\textcolor{FuncColor}{$\triangleright$\enspace\texttt{DigraphVertices({\mdseries\slshape digraph})\index{DigraphVertices@\texttt{DigraphVertices}}
\label{DigraphVertices}
}\hfill{\scriptsize (attribute)}}\\
\textbf{\indent Returns:\ }
A list of positive integers.



 Returns the vertices of the digraph \mbox{\texttt{\mdseries\slshape digraph}}. 

 Note that the vertices of a digraph are always the range of positive integers
from \texttt{1} to the number of vertices of the graph, \texttt{DigraphNrVertices} (\ref{DigraphNrVertices}). Arbitrary \emph{labels} can be assigned to the vertices of a digraph; see \texttt{DigraphVertexLabels} (\ref{DigraphVertexLabels}) for more information about this. 
\begin{Verbatim}[commandchars=!@|,fontsize=\small,frame=single,label=Example]
  !gapprompt@gap>| !gapinput@gr := Digraph(["a", "b", "c"],|
  !gapprompt@>| !gapinput@                 ["a", "b", "b"],|
  !gapprompt@>| !gapinput@                 ["b", "c", "a"]);|
  <immutable digraph with 3 vertices, 3 edges>
  !gapprompt@gap>| !gapinput@DigraphVertices(gr);|
  [ 1 .. 3 ]
  !gapprompt@gap>| !gapinput@gr := Digraph([1, 2, 3, 4, 5, 7],|
  !gapprompt@>| !gapinput@                 [1, 2, 2, 4, 4],|
  !gapprompt@>| !gapinput@                 [2, 7, 5, 3, 7]);|
  <immutable digraph with 6 vertices, 5 edges>
  !gapprompt@gap>| !gapinput@DigraphVertices(gr);|
  [ 1 .. 6 ]
  !gapprompt@gap>| !gapinput@DigraphVertices(RandomDigraph(100));|
  [ 1 .. 100 ]
  !gapprompt@gap>| !gapinput@D := CycleDigraph(IsMutableDigraph, 3);|
  <mutable digraph with 3 vertices, 3 edges>
  !gapprompt@gap>| !gapinput@DigraphVertices(D);|
  [ 1 .. 3 ]
\end{Verbatim}
 }

 

\subsection{\textcolor{Chapter }{DigraphNrVertices}}
\logpage{[ 5, 1, 2 ]}\nobreak
\hyperdef{L}{X7C6F19B57CB2E882}{}
{\noindent\textcolor{FuncColor}{$\triangleright$\enspace\texttt{DigraphNrVertices({\mdseries\slshape digraph})\index{DigraphNrVertices@\texttt{DigraphNrVertices}}
\label{DigraphNrVertices}
}\hfill{\scriptsize (attribute)}}\\
\textbf{\indent Returns:\ }
An non\texttt{\symbol{45}}negative integer.



 Returns the number of vertices of the digraph \mbox{\texttt{\mdseries\slshape digraph}}. 

 
\begin{Verbatim}[commandchars=!@|,fontsize=\small,frame=single,label=Example]
  !gapprompt@gap>| !gapinput@gr := Digraph(["a", "b", "c"],|
  !gapprompt@>| !gapinput@                 ["a", "b", "b"],|
  !gapprompt@>| !gapinput@                 ["b", "c", "a"]);|
  <immutable digraph with 3 vertices, 3 edges>
  !gapprompt@gap>| !gapinput@DigraphNrVertices(gr);|
  3
  !gapprompt@gap>| !gapinput@gr := Digraph([1, 2, 3, 4, 5, 7],|
  !gapprompt@>| !gapinput@                 [1, 2, 2, 4, 4],|
  !gapprompt@>| !gapinput@                 [2, 7, 5, 3, 7]);|
  <immutable digraph with 6 vertices, 5 edges>
  !gapprompt@gap>| !gapinput@DigraphNrVertices(gr);|
  6
  !gapprompt@gap>| !gapinput@DigraphNrVertices(RandomDigraph(100));|
  100
  !gapprompt@gap>| !gapinput@D := CycleDigraph(IsMutableDigraph, 3);|
  <mutable digraph with 3 vertices, 3 edges>
  !gapprompt@gap>| !gapinput@DigraphNrVertices(D);|
  3
\end{Verbatim}
 }

 

\subsection{\textcolor{Chapter }{DigraphEdges}}
\logpage{[ 5, 1, 3 ]}\nobreak
\hyperdef{L}{X7D1C6A4D7ECEC317}{}
{\noindent\textcolor{FuncColor}{$\triangleright$\enspace\texttt{DigraphEdges({\mdseries\slshape digraph})\index{DigraphEdges@\texttt{DigraphEdges}}
\label{DigraphEdges}
}\hfill{\scriptsize (attribute)}}\\
\textbf{\indent Returns:\ }
A list of lists of two positive integers.



 Returns a list of edges of the digraph \mbox{\texttt{\mdseries\slshape digraph}}, where each edge is a pair of elements of \texttt{DigraphVertices} (\ref{DigraphVertices}) of the form \texttt{[source,range]}. 

  The entries of \texttt{DigraphEdges(}\mbox{\texttt{\mdseries\slshape digraph}}\texttt{)} are in one\texttt{\symbol{45}}to\texttt{\symbol{45}}one correspondence with
the edges of \mbox{\texttt{\mdseries\slshape digraph}}. Hence \texttt{DigraphEdges(}\mbox{\texttt{\mdseries\slshape digraph}}\texttt{)} is duplicate\texttt{\symbol{45}}free if and only if \mbox{\texttt{\mdseries\slshape digraph}} contains no multiple edges. 

 The entries of \texttt{DigraphEdges} are guaranteed to be sorted by their first component (i.e. by the source of
each edge), but they are not necessarily then sorted by the second component. 
\begin{Verbatim}[commandchars=!|B,fontsize=\small,frame=single,label=Example]
  !gapprompt|gap>B !gapinput|gr := DigraphFromDiSparse6String(".DaXbOe?EAM@G~");B
  <immutable multidigraph with 5 vertices, 16 edges>
  !gapprompt|gap>B !gapinput|edges := ShallowCopy(DigraphEdges(gr));; Sort(edges);B
  !gapprompt|gap>B !gapinput|edges;B
  [ [ 1, 1 ], [ 1, 3 ], [ 1, 3 ], [ 1, 4 ], [ 1, 5 ], [ 2, 1 ], 
    [ 2, 2 ], [ 2, 3 ], [ 2, 5 ], [ 3, 2 ], [ 3, 4 ], [ 3, 5 ], 
    [ 4, 2 ], [ 4, 4 ], [ 4, 5 ], [ 5, 1 ] ]
  !gapprompt|gap>B !gapinput|D := CycleDigraph(IsMutableDigraph, 3);B
  <mutable digraph with 3 vertices, 3 edges>
  !gapprompt|gap>B !gapinput|DigraphEdges(D);B
  [ [ 1, 2 ], [ 2, 3 ], [ 3, 1 ] ]
\end{Verbatim}
 }

 

\subsection{\textcolor{Chapter }{DigraphNrEdges}}
\logpage{[ 5, 1, 4 ]}\nobreak
\hyperdef{L}{X85E1CFDD7E164AD0}{}
{\noindent\textcolor{FuncColor}{$\triangleright$\enspace\texttt{DigraphNrEdges({\mdseries\slshape digraph})\index{DigraphNrEdges@\texttt{DigraphNrEdges}}
\label{DigraphNrEdges}
}\hfill{\scriptsize (attribute)}}\\
\textbf{\indent Returns:\ }
An integer.



 Returns the number of edges of the digraph \mbox{\texttt{\mdseries\slshape digraph}}. 
\begin{Verbatim}[commandchars=!@|,fontsize=\small,frame=single,label=Example]
  !gapprompt@gap>| !gapinput@gr := Digraph([|
  !gapprompt@>| !gapinput@[1, 3, 4, 5], [1, 2, 3, 5], [2, 4, 5], [2, 4, 5], [1]]);;|
  !gapprompt@gap>| !gapinput@DigraphNrEdges(gr);|
  15
  !gapprompt@gap>| !gapinput@gr := Digraph(["a", "b", "c"],|
  !gapprompt@>| !gapinput@                 ["a", "b", "b"],|
  !gapprompt@>| !gapinput@                 ["b", "a", "a"]);|
  <immutable multidigraph with 3 vertices, 3 edges>
  !gapprompt@gap>| !gapinput@DigraphNrEdges(gr);|
  3
  !gapprompt@gap>| !gapinput@D := CycleDigraph(IsMutableDigraph, 3);|
  <mutable digraph with 3 vertices, 3 edges>
  !gapprompt@gap>| !gapinput@DigraphNrEdges(D);|
  3
\end{Verbatim}
 }

 

\subsection{\textcolor{Chapter }{DigraphNrAdjacencies}}
\logpage{[ 5, 1, 5 ]}\nobreak
\hyperdef{L}{X7BD5D255809C859E}{}
{\noindent\textcolor{FuncColor}{$\triangleright$\enspace\texttt{DigraphNrAdjacencies({\mdseries\slshape digraph})\index{DigraphNrAdjacencies@\texttt{DigraphNrAdjacencies}}
\label{DigraphNrAdjacencies}
}\hfill{\scriptsize (attribute)}}\\
\textbf{\indent Returns:\ }
An integer.



 Returns the number of sets $\{u, v\}$ of vertices of the digraph \mbox{\texttt{\mdseries\slshape digraph}}, such that either $(u, v)$ or $(v, u)$ is an edge. The following equality holds for any digraph \texttt{D} with no multiple edges: \texttt{DigraphNrAdjacencies(D) * 2 \texttt{\symbol{45}} DigraphNrLoops(D) =
DigraphNrEdges(DigraphSymmetricClosure(D))}. 
\begin{Verbatim}[commandchars=!@|,fontsize=\small,frame=single,label=Example]
  !gapprompt@gap>| !gapinput@gr := Digraph([|
  !gapprompt@>| !gapinput@[1, 3, 4, 5], [1, 2, 3, 5], [2, 4, 5], [2, 4, 5], [1]]);;|
  !gapprompt@gap>| !gapinput@DigraphNrAdjacencies(gr);|
  13
  !gapprompt@gap>| !gapinput@DigraphNrAdjacencies(gr) * 2 - DigraphNrLoops(gr) =|
  !gapprompt@>| !gapinput@DigraphNrEdges(DigraphSymmetricClosure(gr));|
  true
\end{Verbatim}
 }

 

\subsection{\textcolor{Chapter }{DigraphNrAdjacenciesWithoutLoops}}
\logpage{[ 5, 1, 6 ]}\nobreak
\hyperdef{L}{X7E53E728795BB862}{}
{\noindent\textcolor{FuncColor}{$\triangleright$\enspace\texttt{DigraphNrAdjacenciesWithoutLoops({\mdseries\slshape digraph})\index{DigraphNrAdjacenciesWithoutLoops@\texttt{DigraphNrAdjacenciesWithoutLoops}}
\label{DigraphNrAdjacenciesWithoutLoops}
}\hfill{\scriptsize (attribute)}}\\
\textbf{\indent Returns:\ }
An integer.



 Returns the number of sets $\{u, v\}$ of vertices of the digraph \mbox{\texttt{\mdseries\slshape digraph}}, such that $u \neq v$ and either $(u, v)$ or $(v, u)$ is an edge. The following equality holds for any digraph \texttt{D} with no multiple edges: \texttt{DigraphNrAdjacenciesWithoutLoops(D) * 2 + DigraphNrLoops(D) =
DigraphNrEdges(DigraphSymmetricClosure(D))}. 
\begin{Verbatim}[commandchars=!@|,fontsize=\small,frame=single,label=Example]
  !gapprompt@gap>| !gapinput@gr := Digraph([|
  !gapprompt@>| !gapinput@[1, 3, 4, 5], [1, 2, 3, 5], [2, 4, 5], [2, 4, 5], [1]]);;|
  !gapprompt@gap>| !gapinput@DigraphNrAdjacenciesWithoutLoops(gr);|
  10
  !gapprompt@gap>| !gapinput@DigraphNrAdjacenciesWithoutLoops(gr) * 2 + DigraphNrLoops(gr) =|
  !gapprompt@>| !gapinput@DigraphNrEdges(DigraphSymmetricClosure(gr));|
  true
\end{Verbatim}
 }

 

\subsection{\textcolor{Chapter }{DigraphNrLoops}}
\logpage{[ 5, 1, 7 ]}\nobreak
\hyperdef{L}{X7BDABAF07917462B}{}
{\noindent\textcolor{FuncColor}{$\triangleright$\enspace\texttt{DigraphNrLoops({\mdseries\slshape digraph})\index{DigraphNrLoops@\texttt{DigraphNrLoops}}
\label{DigraphNrLoops}
}\hfill{\scriptsize (attribute)}}\\
\textbf{\indent Returns:\ }
An integer.



 This function returns the number of loops of the digraph \mbox{\texttt{\mdseries\slshape digraph}}. See \texttt{DigraphHasLoops} (\ref{DigraphHasLoops}). 

 
\begin{Verbatim}[commandchars=!@|,fontsize=\small,frame=single,label=Example]
  !gapprompt@gap>| !gapinput@D := Digraph([[2, 3], [1, 4], [3, 3, 5], [], [2, 5]]);|
  <immutable multidigraph with 5 vertices, 9 edges>
  !gapprompt@gap>| !gapinput@DigraphNrLoops(D);|
  3
  !gapprompt@gap>| !gapinput@D := EmptyDigraph(5);|
  <immutable empty digraph with 5 vertices>
  !gapprompt@gap>| !gapinput@DigraphNrLoops(D);|
  0
  !gapprompt@gap>| !gapinput@D := CompleteDigraph(5);|
  <immutable complete digraph with 5 vertices>
  !gapprompt@gap>| !gapinput@DigraphNrLoops(D);|
  0
\end{Verbatim}
 }

 

\subsection{\textcolor{Chapter }{DigraphSinks}}
\logpage{[ 5, 1, 8 ]}\nobreak
\hyperdef{L}{X85D5E08280914EE4}{}
{\noindent\textcolor{FuncColor}{$\triangleright$\enspace\texttt{DigraphSinks({\mdseries\slshape digraph})\index{DigraphSinks@\texttt{DigraphSinks}}
\label{DigraphSinks}
}\hfill{\scriptsize (attribute)}}\\
\textbf{\indent Returns:\ }
A list of vertices.



 This function returns a list of the sinks of the digraph \mbox{\texttt{\mdseries\slshape digraph}}. A sink of a digraph is a vertex with out\texttt{\symbol{45}}degree zero. See \texttt{OutDegreeOfVertex} (\ref{OutDegreeOfVertex}). 
\begin{Verbatim}[commandchars=!@|,fontsize=\small,frame=single,label=Example]
  !gapprompt@gap>| !gapinput@gr := Digraph([[3, 5, 2, 2], [3], [], [5, 2, 5, 3], []]);|
  <immutable multidigraph with 5 vertices, 9 edges>
  !gapprompt@gap>| !gapinput@DigraphSinks(gr);|
  [ 3, 5 ]
  !gapprompt@gap>| !gapinput@D := CycleDigraph(IsMutableDigraph, 3);|
  <mutable digraph with 3 vertices, 3 edges>
  !gapprompt@gap>| !gapinput@DigraphSinks(D);|
  [  ]
\end{Verbatim}
 }

 

\subsection{\textcolor{Chapter }{DigraphSources}}
\logpage{[ 5, 1, 9 ]}\nobreak
\hyperdef{L}{X7F5C6268839BE98C}{}
{\noindent\textcolor{FuncColor}{$\triangleright$\enspace\texttt{DigraphSources({\mdseries\slshape digraph})\index{DigraphSources@\texttt{DigraphSources}}
\label{DigraphSources}
}\hfill{\scriptsize (attribute)}}\\
\textbf{\indent Returns:\ }
A list of vertices.



 This function returns an immutable list of the sources of the digraph \mbox{\texttt{\mdseries\slshape digraph}}. A source of a digraph is a vertex with in\texttt{\symbol{45}}degree zero.
See \texttt{InDegreeOfVertex} (\ref{InDegreeOfVertex}). 
\begin{Verbatim}[commandchars=!@|,fontsize=\small,frame=single,label=Example]
  !gapprompt@gap>| !gapinput@gr := Digraph([[3, 5, 2, 2], [3], [], [5, 2, 5, 3], []]);|
  <immutable multidigraph with 5 vertices, 9 edges>
  !gapprompt@gap>| !gapinput@DigraphSources(gr);|
  [ 1, 4 ]
  !gapprompt@gap>| !gapinput@D := CycleDigraph(IsMutableDigraph, 3);|
  <mutable digraph with 3 vertices, 3 edges>
  !gapprompt@gap>| !gapinput@DigraphSources(D);|
  [  ]
\end{Verbatim}
 }

 

\subsection{\textcolor{Chapter }{DigraphTopologicalSort}}
\logpage{[ 5, 1, 10 ]}\nobreak
\hyperdef{L}{X785C30378064CF47}{}
{\noindent\textcolor{FuncColor}{$\triangleright$\enspace\texttt{DigraphTopologicalSort({\mdseries\slshape digraph})\index{DigraphTopologicalSort@\texttt{DigraphTopologicalSort}}
\label{DigraphTopologicalSort}
}\hfill{\scriptsize (attribute)}}\\
\textbf{\indent Returns:\ }
A list of positive integers, or \texttt{fail}.



 If \mbox{\texttt{\mdseries\slshape digraph}} is a digraph whose only directed cycles are loops, then \texttt{DigraphTopologicalSort} returns the vertices of \mbox{\texttt{\mdseries\slshape digraph}} ordered so that every edge's source appears no earlier in the list than its
range. If the digraph \mbox{\texttt{\mdseries\slshape digraph}} contains directed cycles of length greater than $1$, then this operation returns \texttt{fail}. 

 See Section \ref{Definitions} for the definition of a directed cycle, and the definition of a loop. 

 The method used for this attribute has complexity $O(m+n)$ where $m$ is the number of edges (counting multiple edges as one) and $n$ is the number of vertices in the digraph. 

 
\begin{Verbatim}[commandchars=!@|,fontsize=\small,frame=single,label=Example]
  !gapprompt@gap>| !gapinput@D := Digraph([|
  !gapprompt@>| !gapinput@[2, 3], [], [4, 6], [5], [], [7, 8, 9], [], [], []]);|
  <immutable digraph with 9 vertices, 8 edges>
  !gapprompt@gap>| !gapinput@DigraphTopologicalSort(D);|
  [ 2, 5, 4, 7, 8, 9, 6, 3, 1 ]
  !gapprompt@gap>| !gapinput@D := Digraph(IsMutableDigraph, [[2, 3], [3], [4], []]);|
  <mutable digraph with 4 vertices, 4 edges>
  !gapprompt@gap>| !gapinput@DigraphTopologicalSort(D);|
  [ 4, 3, 2, 1 ]
\end{Verbatim}
 }

 

\subsection{\textcolor{Chapter }{DigraphVertexLabel}}
\logpage{[ 5, 1, 11 ]}\nobreak
\hyperdef{L}{X7CA91E4B7904F793}{}
{\noindent\textcolor{FuncColor}{$\triangleright$\enspace\texttt{DigraphVertexLabel({\mdseries\slshape digraph, i})\index{DigraphVertexLabel@\texttt{DigraphVertexLabel}}
\label{DigraphVertexLabel}
}\hfill{\scriptsize (operation)}}\\
\noindent\textcolor{FuncColor}{$\triangleright$\enspace\texttt{SetDigraphVertexLabel({\mdseries\slshape digraph, i, obj})\index{SetDigraphVertexLabel@\texttt{SetDigraphVertexLabel}}
\label{SetDigraphVertexLabel}
}\hfill{\scriptsize (operation)}}\\


 If \mbox{\texttt{\mdseries\slshape digraph}} is a digraph, then the first operation returns the label of the vertex \mbox{\texttt{\mdseries\slshape i}}. The second operation can be used to set the label of the vertex \mbox{\texttt{\mdseries\slshape i}} in \mbox{\texttt{\mdseries\slshape digraph}} to the arbitrary \textsf{GAP} object \mbox{\texttt{\mdseries\slshape obj}}. 

 The label of a vertex can be changed an arbitrary number of times. If no label
has been set for the vertex \mbox{\texttt{\mdseries\slshape i}}, then the default value is \mbox{\texttt{\mdseries\slshape i}}. 

 If \mbox{\texttt{\mdseries\slshape digraph}} is a digraph created from a record with a component \texttt{DigraphVertices}, then the labels of the vertices are set to the value of this component.

 Induced subdigraphs, and some other operations which create new digraphs from
old ones, inherit their labels from their parents. 
\begin{Verbatim}[commandchars=!@|,fontsize=\small,frame=single,label=Example]
  !gapprompt@gap>| !gapinput@D := DigraphFromDigraph6String("&DHUEe_");|
  <immutable digraph with 5 vertices, 11 edges>
  !gapprompt@gap>| !gapinput@DigraphVertexLabel(D, 3);|
  3
  !gapprompt@gap>| !gapinput@D := Digraph(["a", "b", "c"], [], []);|
  <immutable empty digraph with 3 vertices>
  !gapprompt@gap>| !gapinput@DigraphVertexLabel(D, 2);|
  "b"
  !gapprompt@gap>| !gapinput@SetDigraphVertexLabel(D, 2, "d");|
  !gapprompt@gap>| !gapinput@DigraphVertexLabel(D, 2);|
  "d"
  !gapprompt@gap>| !gapinput@D := InducedSubdigraph(D, [1, 2]);|
  <immutable empty digraph with 2 vertices>
  !gapprompt@gap>| !gapinput@DigraphVertexLabel(D, 2);|
  "d"
  !gapprompt@gap>| !gapinput@D := Digraph(IsMutableDigraph, ["e", "f", "g"], [], []);|
  <mutable empty digraph with 3 vertices>
  !gapprompt@gap>| !gapinput@DigraphVertexLabel(D, 1);|
  "e"
  !gapprompt@gap>| !gapinput@SetDigraphVertexLabel(D, 1, "h");|
  !gapprompt@gap>| !gapinput@DigraphVertexLabel(D, 1);|
  "h"
  !gapprompt@gap>| !gapinput@InducedSubdigraph(D, [1, 2]);|
  <mutable empty digraph with 2 vertices>
  !gapprompt@gap>| !gapinput@DigraphVertexLabel(D, 1);|
  "h"
\end{Verbatim}
 }

 

\subsection{\textcolor{Chapter }{DigraphVertexLabels}}
\logpage{[ 5, 1, 12 ]}\nobreak
\hyperdef{L}{X7E51F2FE87140B32}{}
{\noindent\textcolor{FuncColor}{$\triangleright$\enspace\texttt{DigraphVertexLabels({\mdseries\slshape digraph})\index{DigraphVertexLabels@\texttt{DigraphVertexLabels}}
\label{DigraphVertexLabels}
}\hfill{\scriptsize (operation)}}\\
\noindent\textcolor{FuncColor}{$\triangleright$\enspace\texttt{SetDigraphVertexLabels({\mdseries\slshape digraph, list})\index{SetDigraphVertexLabels@\texttt{SetDigraphVertexLabels}}
\label{SetDigraphVertexLabels}
}\hfill{\scriptsize (operation)}}\\


 If \mbox{\texttt{\mdseries\slshape digraph}} is a digraph, then \texttt{DigraphVertexLabels} returns a copy of the labels of the vertices in \mbox{\texttt{\mdseries\slshape digraph}}. \texttt{SetDigraphVertexLabels} can be used to set the labels of the vertices in \mbox{\texttt{\mdseries\slshape digraph}} to the list of arbitrary \textsf{GAP} objects \mbox{\texttt{\mdseries\slshape list}}, which must be of the same length as the number of vertices of \mbox{\texttt{\mdseries\slshape digraph}}. 

 If the list \mbox{\texttt{\mdseries\slshape list}} is immutable, then the vertex labels are set to a mutable copy of \mbox{\texttt{\mdseries\slshape list}}. Otherwise, the labels are set to exactly \mbox{\texttt{\mdseries\slshape list}}. 

 The label of a vertex can be changed an arbitrary number of times. If no label
has been set for the vertex \mbox{\texttt{\mdseries\slshape i}}, then the default value is \mbox{\texttt{\mdseries\slshape i}}. 

 If \mbox{\texttt{\mdseries\slshape digraph}} is a digraph created from a record with a component \texttt{DigraphVertices}, then the labels of the vertices are set to the value of this component. As
in the above, if the component is immutable then the digraph's vertex labels
are set to a mutable copy of \texttt{DigraphVertices}. Otherwise, they are set to exactly \texttt{DigraphVertices}. 

 Induced subdigraphs, and other operations which create new digraphs from old
ones, inherit their labels from their parents. 
\begin{Verbatim}[commandchars=!@|,fontsize=\small,frame=single,label=Example]
  !gapprompt@gap>| !gapinput@D := DigraphFromDigraph6String("&DHUEe_");|
  <immutable digraph with 5 vertices, 11 edges>
  !gapprompt@gap>| !gapinput@DigraphVertexLabels(D);|
  [ 1 .. 5 ]
  !gapprompt@gap>| !gapinput@D := Digraph(["a", "b", "c"], [], []);|
  <immutable empty digraph with 3 vertices>
  !gapprompt@gap>| !gapinput@DigraphVertexLabels(D);|
  [ "a", "b", "c" ]
  !gapprompt@gap>| !gapinput@SetDigraphVertexLabel(D, 2, "d");|
  !gapprompt@gap>| !gapinput@DigraphVertexLabels(D);|
  [ "a", "d", "c" ]
  !gapprompt@gap>| !gapinput@D := InducedSubdigraph(D, [1, 3]);|
  <immutable empty digraph with 2 vertices>
  !gapprompt@gap>| !gapinput@DigraphVertexLabels(D);|
  [ "a", "c" ]
  !gapprompt@gap>| !gapinput@D := Digraph(IsMutableDigraph, ["e", "f", "g"], [], []);|
  <mutable empty digraph with 3 vertices>
  !gapprompt@gap>| !gapinput@SetDigraphVertexLabels(D, ["h", "i", "j"]);|
  !gapprompt@gap>| !gapinput@DigraphVertexLabels(D);|
  [ "h", "i", "j" ]
  !gapprompt@gap>| !gapinput@InducedSubdigraph(D, [1, 3]);|
  <mutable empty digraph with 2 vertices>
  !gapprompt@gap>| !gapinput@DigraphVertexLabels(D);|
  [ "h", "j" ]
\end{Verbatim}
 }

 

\subsection{\textcolor{Chapter }{DigraphEdgeLabel}}
\logpage{[ 5, 1, 13 ]}\nobreak
\hyperdef{L}{X79FAEACC7F438C2F}{}
{\noindent\textcolor{FuncColor}{$\triangleright$\enspace\texttt{DigraphEdgeLabel({\mdseries\slshape digraph, i, j})\index{DigraphEdgeLabel@\texttt{DigraphEdgeLabel}}
\label{DigraphEdgeLabel}
}\hfill{\scriptsize (operation)}}\\
\noindent\textcolor{FuncColor}{$\triangleright$\enspace\texttt{SetDigraphEdgeLabel({\mdseries\slshape digraph, i, j, obj})\index{SetDigraphEdgeLabel@\texttt{SetDigraphEdgeLabel}}
\label{SetDigraphEdgeLabel}
}\hfill{\scriptsize (operation)}}\\


 If \mbox{\texttt{\mdseries\slshape digraph}} is a digraph without multiple edges, then the first operation returns the
label of the edge from vertex \mbox{\texttt{\mdseries\slshape i}} to vertex \mbox{\texttt{\mdseries\slshape j}}. The second operation can be used to set the label of the edge between vertex \mbox{\texttt{\mdseries\slshape i}} and vertex \mbox{\texttt{\mdseries\slshape j}} to the arbitrary \textsf{GAP} object \mbox{\texttt{\mdseries\slshape obj}}. 

 The label of an edge can be changed an arbitrary number of times. If no label
has been set for the edge, then the default value is \mbox{\texttt{\mdseries\slshape 1}}. 

 Induced subdigraphs, and some other operations which create new digraphs from
old ones, inherit their edge labels from their parents. See also \texttt{DigraphEdgeLabels} (\ref{DigraphEdgeLabels}). 
\begin{Verbatim}[commandchars=!@|,fontsize=\small,frame=single,label=Example]
  !gapprompt@gap>| !gapinput@D := DigraphFromDigraph6String("&DHUEe_");|
  <immutable digraph with 5 vertices, 11 edges>
  !gapprompt@gap>| !gapinput@DigraphEdgeLabel(D, 3, 1);|
  1
  !gapprompt@gap>| !gapinput@SetDigraphEdgeLabel(D, 2, 5, [42]);|
  !gapprompt@gap>| !gapinput@DigraphEdgeLabel(D, 2, 5);|
  [ 42 ]
  !gapprompt@gap>| !gapinput@D := InducedSubdigraph(D, [2, 5]);|
  <immutable digraph with 2 vertices, 3 edges>
  !gapprompt@gap>| !gapinput@DigraphEdgeLabel(D, 1, 2);|
  [ 42 ]
  !gapprompt@gap>| !gapinput@D := ChainDigraph(IsMutableDigraph, 5);|
  <mutable digraph with 5 vertices, 4 edges>
  !gapprompt@gap>| !gapinput@DigraphEdgeLabel(D, 2, 3);|
  1
  !gapprompt@gap>| !gapinput@SetDigraphEdgeLabel(D, 4, 5, [1729]);|
  !gapprompt@gap>| !gapinput@DigraphEdgeLabel(D, 4, 5);|
  [ 1729 ]
  !gapprompt@gap>| !gapinput@InducedSubdigraph(D, [4, 5]);|
  <mutable digraph with 2 vertices, 1 edge>
  !gapprompt@gap>| !gapinput@DigraphEdgeLabel(D, 1, 2);|
  [ 1729 ]
\end{Verbatim}
 }

 

\subsection{\textcolor{Chapter }{DigraphEdgeLabels}}
\logpage{[ 5, 1, 14 ]}\nobreak
\hyperdef{L}{X7C24851087D4A8FB}{}
{\noindent\textcolor{FuncColor}{$\triangleright$\enspace\texttt{DigraphEdgeLabels({\mdseries\slshape digraph})\index{DigraphEdgeLabels@\texttt{DigraphEdgeLabels}}
\label{DigraphEdgeLabels}
}\hfill{\scriptsize (operation)}}\\
\noindent\textcolor{FuncColor}{$\triangleright$\enspace\texttt{SetDigraphEdgeLabels({\mdseries\slshape digraph, labels})\index{SetDigraphEdgeLabels@\texttt{SetDigraphEdgeLabels}!for a digraph and a list of lists}
\label{SetDigraphEdgeLabels:for a digraph and a list of lists}
}\hfill{\scriptsize (operation)}}\\
\noindent\textcolor{FuncColor}{$\triangleright$\enspace\texttt{SetDigraphEdgeLabels({\mdseries\slshape digraph, func})\index{SetDigraphEdgeLabels@\texttt{SetDigraphEdgeLabels}!for a digraph and a function}
\label{SetDigraphEdgeLabels:for a digraph and a function}
}\hfill{\scriptsize (operation)}}\\


 If \mbox{\texttt{\mdseries\slshape digraph}} is a digraph without multiple edges, then \texttt{DigraphEdgeLabels} returns a copy of the labels of the edges in \mbox{\texttt{\mdseries\slshape digraph}} as a list of lists \texttt{labels} such that \texttt{labels[i][j]} is the label on the edge from vertex \texttt{i} to vertex \texttt{OutNeighbours(digraph)[i][j]}. \texttt{SetDigraphEdgeLabels} can be used to set the labels of the edges in \mbox{\texttt{\mdseries\slshape digraph}} without multiple edges to the list \mbox{\texttt{\mdseries\slshape labels}} of lists of arbitrary \textsf{GAP} objects such that \texttt{list[i][j]} is the label on the edge from vertex \texttt{i} to the vertex \texttt{OutNeighbours(digraph{\textgreater}[i][j]}. Alternatively \texttt{SetDigraphEdgeLabels} can be called with binary function \mbox{\texttt{\mdseries\slshape func}} that as its second argument that when passed two vertices \texttt{i} and \texttt{j} returns the label for the edge between vertex \texttt{i} and vertex \texttt{j}. 

 The label of an edge can be changed an arbitrary number of times. If no label
has been set for an edge, then the default value is \texttt{1}. 

 Induced subdigraphs, and some other operations which create new digraphs from
old ones, inherit their labels from their parents. 
\begin{Verbatim}[commandchars=!@|,fontsize=\small,frame=single,label=Example]
  !gapprompt@gap>| !gapinput@D := DigraphFromDigraph6String("&DHUEe_");|
  <immutable digraph with 5 vertices, 11 edges>
  !gapprompt@gap>| !gapinput@DigraphEdgeLabels(D);|
  [ [ 1 ], [ 1, 1, 1 ], [ 1 ], [ 1, 1, 1 ], [ 1, 1, 1 ] ]
  !gapprompt@gap>| !gapinput@SetDigraphEdgeLabel(D, 2, 1, "d");|
  !gapprompt@gap>| !gapinput@DigraphEdgeLabels(D);|
  [ [ 1 ], [ "d", 1, 1 ], [ 1 ], [ 1, 1, 1 ], [ 1, 1, 1 ] ]
  !gapprompt@gap>| !gapinput@D := InducedSubdigraph(D, [1, 2, 3]);|
  <immutable digraph with 3 vertices, 4 edges>
  !gapprompt@gap>| !gapinput@DigraphEdgeLabels(D);|
  [ [ 1 ], [ "d", 1 ], [ 1 ] ]
  !gapprompt@gap>| !gapinput@OutNeighbours(D);|
  [ [ 3 ], [ 1, 3 ], [ 1 ] ]
  !gapprompt@gap>| !gapinput@D := CompleteBipartiteDigraph(IsMutableDigraph, 2, 3);|
  <mutable digraph with 5 vertices, 12 edges>
  !gapprompt@gap>| !gapinput@DigraphEdgeLabels(D);|
  [ [ 1, 1, 1 ], [ 1, 1, 1 ], [ 1, 1 ], [ 1, 1 ], [ 1, 1 ] ]
  !gapprompt@gap>| !gapinput@SetDigraphEdgeLabel(D, 2, 4, "a");|
  !gapprompt@gap>| !gapinput@DigraphEdgeLabels(D);|
  [ [ 1, 1, 1 ], [ 1, "a", 1 ], [ 1, 1 ], [ 1, 1 ], [ 1, 1 ] ]
  !gapprompt@gap>| !gapinput@InducedSubdigraph(D, [1, 2, 3, 4]);|
  <mutable digraph with 4 vertices, 8 edges>
  !gapprompt@gap>| !gapinput@DigraphEdgeLabels(D);|
  [ [ 1, 1 ], [ 1, "a" ], [ 1, 1 ], [ 1, 1 ] ]
  !gapprompt@gap>| !gapinput@OutNeighbors(D);|
  [ [ 3, 4 ], [ 3, 4 ], [ 1, 2 ], [ 1, 2 ] ]
\end{Verbatim}
 }

 

\subsection{\textcolor{Chapter }{DigraphInEdges}}
\logpage{[ 5, 1, 15 ]}\nobreak
\hyperdef{L}{X7EFAF01B7A155157}{}
{\noindent\textcolor{FuncColor}{$\triangleright$\enspace\texttt{DigraphInEdges({\mdseries\slshape digraph, vertex})\index{DigraphInEdges@\texttt{DigraphInEdges}}
\label{DigraphInEdges}
}\hfill{\scriptsize (operation)}}\\
\textbf{\indent Returns:\ }
A list of edges.



 \texttt{DigraphInEdges} returns the list of all edges of \mbox{\texttt{\mdseries\slshape digraph}} which have \mbox{\texttt{\mdseries\slshape vertex}} as their range. 
\begin{Verbatim}[commandchars=!@|,fontsize=\small,frame=single,label=Example]
  !gapprompt@gap>| !gapinput@D := Digraph([[2, 2], [3, 3], [4, 4], [1, 1]]);|
  <immutable multidigraph with 4 vertices, 8 edges>
  !gapprompt@gap>| !gapinput@DigraphInEdges(D, 2);|
  [ [ 1, 2 ], [ 1, 2 ] ]
\end{Verbatim}
 }

 

\subsection{\textcolor{Chapter }{DigraphOutEdges}}
\logpage{[ 5, 1, 16 ]}\nobreak
\hyperdef{L}{X7BECFE6687ECD028}{}
{\noindent\textcolor{FuncColor}{$\triangleright$\enspace\texttt{DigraphOutEdges({\mdseries\slshape digraph, vertex})\index{DigraphOutEdges@\texttt{DigraphOutEdges}}
\label{DigraphOutEdges}
}\hfill{\scriptsize (operation)}}\\
\textbf{\indent Returns:\ }
A list of edges.



 \texttt{DigraphOutEdges} returns the list of all edges of \mbox{\texttt{\mdseries\slshape digraph}} which have \mbox{\texttt{\mdseries\slshape vertex}} as their source. 
\begin{Verbatim}[commandchars=!@|,fontsize=\small,frame=single,label=Example]
  !gapprompt@gap>| !gapinput@D := Digraph([[2, 2], [3, 3], [4, 4], [1, 1]]);|
  <immutable multidigraph with 4 vertices, 8 edges>
  !gapprompt@gap>| !gapinput@DigraphOutEdges(D, 2);|
  [ [ 2, 3 ], [ 2, 3 ] ]
\end{Verbatim}
 }

 

\subsection{\textcolor{Chapter }{IsDigraphEdge (for digraph and list)}}
\logpage{[ 5, 1, 17 ]}\nobreak
\hyperdef{L}{X7BB8ED88835F07B4}{}
{\noindent\textcolor{FuncColor}{$\triangleright$\enspace\texttt{IsDigraphEdge({\mdseries\slshape digraph, list})\index{IsDigraphEdge@\texttt{IsDigraphEdge}!for digraph and list}
\label{IsDigraphEdge:for digraph and list}
}\hfill{\scriptsize (operation)}}\\
\noindent\textcolor{FuncColor}{$\triangleright$\enspace\texttt{IsDigraphEdge({\mdseries\slshape digraph, u, v})\index{IsDigraphEdge@\texttt{IsDigraphEdge}!for digraph and two pos ints}
\label{IsDigraphEdge:for digraph and two pos ints}
}\hfill{\scriptsize (operation)}}\\
\textbf{\indent Returns:\ }
\texttt{true} or \texttt{false}.



 In the first form, this function returns \texttt{true} if and only if the list \mbox{\texttt{\mdseries\slshape list}} specifies an edge in the digraph \mbox{\texttt{\mdseries\slshape digraph}}. Specifically, this operation returns \texttt{true} if \mbox{\texttt{\mdseries\slshape list}} is a pair of positive integers where \mbox{\texttt{\mdseries\slshape list}}\texttt{[1]} is the source and \mbox{\texttt{\mdseries\slshape list}}\texttt{[2]} is the range of an edge in \mbox{\texttt{\mdseries\slshape digraph}}, and \texttt{false} otherwise. 

 The second form simply returns \texttt{true} if \texttt{[\mbox{\texttt{\mdseries\slshape u}}, \mbox{\texttt{\mdseries\slshape v}}]} is an edge in \mbox{\texttt{\mdseries\slshape digraph}}, and \texttt{false} otherwise. 
\begin{Verbatim}[commandchars=!@|,fontsize=\small,frame=single,label=Example]
  !gapprompt@gap>| !gapinput@D := Digraph([[2, 2], [6], [], [3], [], [1]]);|
  <immutable multidigraph with 6 vertices, 5 edges>
  !gapprompt@gap>| !gapinput@IsDigraphEdge(D, [1, 1]);|
  false
  !gapprompt@gap>| !gapinput@IsDigraphEdge(D, [1, 2]);|
  true
  !gapprompt@gap>| !gapinput@IsDigraphEdge(D, [1, 8]);|
  false
\end{Verbatim}
 }

 

\subsection{\textcolor{Chapter }{IsMatching}}
\logpage{[ 5, 1, 18 ]}\nobreak
\hyperdef{L}{X848FED0B7B4ACD1F}{}
{\noindent\textcolor{FuncColor}{$\triangleright$\enspace\texttt{IsMatching({\mdseries\slshape digraph, list})\index{IsMatching@\texttt{IsMatching}}
\label{IsMatching}
}\hfill{\scriptsize (operation)}}\\
\noindent\textcolor{FuncColor}{$\triangleright$\enspace\texttt{IsMaximalMatching({\mdseries\slshape digraph, list})\index{IsMaximalMatching@\texttt{IsMaximalMatching}}
\label{IsMaximalMatching}
}\hfill{\scriptsize (operation)}}\\
\noindent\textcolor{FuncColor}{$\triangleright$\enspace\texttt{IsMaximumMatching({\mdseries\slshape digraph, list})\index{IsMaximumMatching@\texttt{IsMaximumMatching}}
\label{IsMaximumMatching}
}\hfill{\scriptsize (operation)}}\\
\noindent\textcolor{FuncColor}{$\triangleright$\enspace\texttt{IsPerfectMatching({\mdseries\slshape digraph, list})\index{IsPerfectMatching@\texttt{IsPerfectMatching}}
\label{IsPerfectMatching}
}\hfill{\scriptsize (operation)}}\\
\textbf{\indent Returns:\ }
\texttt{true} or \texttt{false}.



 If \mbox{\texttt{\mdseries\slshape digraph}} is a digraph and \mbox{\texttt{\mdseries\slshape list}} is a list of pairs of vertices of \mbox{\texttt{\mdseries\slshape digraph}}, then \texttt{IsMatching} returns \texttt{true} if \mbox{\texttt{\mdseries\slshape list}} is a matching of \mbox{\texttt{\mdseries\slshape digraph}}. The operation \texttt{IsMaximalMatching} returns \texttt{true} if \mbox{\texttt{\mdseries\slshape list}} is a maximal matching, \texttt{IsMaximumMatching} returns \texttt{true} if \mbox{\texttt{\mdseries\slshape list}} is a maximum matching and \texttt{IsPerfectMatching} returns \texttt{true} if \mbox{\texttt{\mdseries\slshape list}} is a perfect, matching of \mbox{\texttt{\mdseries\slshape digraph}}, respectively. Otherwise, each of these operations return \texttt{false}. 

 A \emph{matching} \texttt{M} of a digraph \mbox{\texttt{\mdseries\slshape digraph}} is a subset of the edges of \mbox{\texttt{\mdseries\slshape digraph}}, i.e. \texttt{DigraphEdges(\mbox{\texttt{\mdseries\slshape digraph}})}, such that no pair of distinct edges in \texttt{M} are incident to the same vertex of \mbox{\texttt{\mdseries\slshape digraph}}. Note that this definition allows a matching to contain loops. See \texttt{DigraphHasLoops} (\ref{DigraphHasLoops}). The matching \texttt{M} is \emph{maximal} if it is contained in no larger matching of the digraph, is \emph{maximum} if it has the greatest cardinality among all matchings and is \emph{perfect} if every vertex of the digraph is incident to an edge in the matching. Every
maximum or perfect matching is maximal. Note, however, that not every perfect
matching of digraphs with loops is maximum. 
\begin{Verbatim}[commandchars=!@|,fontsize=\small,frame=single,label=Example]
  !gapprompt@gap>| !gapinput@D := Digraph([[1, 2], [1, 2], [2, 3, 4], [3, 5], [1]]);|
  <immutable digraph with 5 vertices, 10 edges>
  !gapprompt@gap>| !gapinput@IsMatching(D, [[2, 1], [3, 2]]);|
  false
  !gapprompt@gap>| !gapinput@edges := [[3, 2]];;|
  !gapprompt@gap>| !gapinput@IsMatching(D, edges);|
  true
  !gapprompt@gap>| !gapinput@IsMaximalMatching(D, edges);|
  false
  !gapprompt@gap>| !gapinput@edges := [[2, 1], [3, 4]];;|
  !gapprompt@gap>| !gapinput@IsMaximalMatching(D, edges);|
  true
  !gapprompt@gap>| !gapinput@IsPerfectMatching(D, edges);|
  false
  !gapprompt@gap>| !gapinput@edges := [[1, 2], [3, 3], [4, 5]];;|
  !gapprompt@gap>| !gapinput@IsPerfectMatching(D, edges);|
  true
  !gapprompt@gap>| !gapinput@IsMaximumMatching(D, edges);|
  false
  !gapprompt@gap>| !gapinput@edges := [[1, 1], [2, 2], [3, 3], [4, 5]];;|
  !gapprompt@gap>| !gapinput@IsMaximumMatching(D, edges);|
  true
\end{Verbatim}
 }

 

\subsection{\textcolor{Chapter }{DigraphMaximalMatching}}
\logpage{[ 5, 1, 19 ]}\nobreak
\hyperdef{L}{X7B7A67277B1C9A02}{}
{\noindent\textcolor{FuncColor}{$\triangleright$\enspace\texttt{DigraphMaximalMatching({\mdseries\slshape digraph})\index{DigraphMaximalMatching@\texttt{DigraphMaximalMatching}}
\label{DigraphMaximalMatching}
}\hfill{\scriptsize (attribute)}}\\
\textbf{\indent Returns:\ }
A list of pairs of vertices.



 This function returns a maximal matching of the digraph \mbox{\texttt{\mdseries\slshape digraph}}. 

 For the definition of a maximal matching, see \texttt{IsMaximalMatching} (\ref{IsMaximalMatching}). 
\begin{Verbatim}[commandchars=!BE,fontsize=\small,frame=single,label=Example]
  !gappromptBgap>E !gapinputBD := DigraphFromDiSparse6String(".IeAoXCJU@|SHAe?d");E
  <immutable digraph with 10 vertices, 13 edges>
  !gappromptBgap>E !gapinputBM := DigraphMaximalMatching(D);; IsMaximalMatching(D, M);E
  true
  !gappromptBgap>E !gapinputBD := RandomDigraph(100);;E
  !gappromptBgap>E !gapinputBIsMaximalMatching(D, DigraphMaximalMatching(D));E
  true
  !gappromptBgap>E !gapinputBD := GeneralisedPetersenGraph(IsMutableDigraph, 9, 2);E
  <mutable digraph with 18 vertices, 54 edges>
  !gappromptBgap>E !gapinputBIsMaximalMatching(D, DigraphMaximalMatching(D));E
  true
\end{Verbatim}
 }

 

\subsection{\textcolor{Chapter }{DigraphMaximumMatching}}
\logpage{[ 5, 1, 20 ]}\nobreak
\hyperdef{L}{X78E9847A858788D1}{}
{\noindent\textcolor{FuncColor}{$\triangleright$\enspace\texttt{DigraphMaximumMatching({\mdseries\slshape digraph})\index{DigraphMaximumMatching@\texttt{DigraphMaximumMatching}}
\label{DigraphMaximumMatching}
}\hfill{\scriptsize (attribute)}}\\
\textbf{\indent Returns:\ }
A list of pairs of vertices.



 This function returns a maximum matching of the digraph \mbox{\texttt{\mdseries\slshape digraph}}. 

 For the definition of a maximum matching, see \texttt{IsMaximumMatching} (\ref{IsMaximumMatching}). If \mbox{\texttt{\mdseries\slshape digraph}} is bipartite (see \texttt{IsBipartiteDigraph} (\ref{IsBipartiteDigraph})), then the algorithm used has complexity \texttt{O(m*sqrt(n))}. Otherwise for general graphs the complexity is \texttt{O(m*n*log(n))}. Here \texttt{n} is the number of vertices and \texttt{m} is the number of edges. 
\begin{Verbatim}[commandchars=!|B,fontsize=\small,frame=single,label=Example]
  !gapprompt|gap>B !gapinput|D := DigraphFromDigraph6String("&I@EA_A?AdDp[_c??OO");B
  <immutable digraph with 10 vertices, 23 edges>
  !gapprompt|gap>B !gapinput|M := DigraphMaximumMatching(D);; IsMaximalMatching(D, M);B
  true
  !gapprompt|gap>B !gapinput|Length(M);B
  5
  !gapprompt|gap>B !gapinput|D := Digraph([[5, 6, 7, 8], [6, 7, 8], [7, 8], [8], B
  !gapprompt|>B !gapinput|                 [], [], [], []]);;B
  !gapprompt|gap>B !gapinput|M := DigraphMaximumMatching(D);B
  [ [ 1, 5 ], [ 2, 6 ], [ 3, 7 ], [ 4, 8 ] ]
  !gapprompt|gap>B !gapinput|D := GeneralisedPetersenGraph(IsMutableDigraph, 9, 2);B
  <mutable digraph with 18 vertices, 54 edges>
  !gapprompt|gap>B !gapinput|M := DigraphMaximumMatching(D);;B
  !gapprompt|gap>B !gapinput|IsMaximalMatching(D, M);B
  true
  !gapprompt|gap>B !gapinput|Length(M);B
  9
\end{Verbatim}
 }

 }

 
\section{\textcolor{Chapter }{Neighbours and degree}}\logpage{[ 5, 2, 0 ]}
\hyperdef{L}{X7D7CE8328187D0DF}{}
{
 

\subsection{\textcolor{Chapter }{AdjacencyMatrix}}
\logpage{[ 5, 2, 1 ]}\nobreak
\hyperdef{L}{X7DC2CD70830BEE60}{}
{\noindent\textcolor{FuncColor}{$\triangleright$\enspace\texttt{AdjacencyMatrix({\mdseries\slshape digraph})\index{AdjacencyMatrix@\texttt{AdjacencyMatrix}}
\label{AdjacencyMatrix}
}\hfill{\scriptsize (attribute)}}\\
\noindent\textcolor{FuncColor}{$\triangleright$\enspace\texttt{AdjacencyMatrixMutableCopy({\mdseries\slshape digraph})\index{AdjacencyMatrixMutableCopy@\texttt{AdjacencyMatrixMutableCopy}}
\label{AdjacencyMatrixMutableCopy}
}\hfill{\scriptsize (operation)}}\\
\textbf{\indent Returns:\ }
A square matrix of non\texttt{\symbol{45}}negative integers.



 This function returns the adjacency matrix \texttt{mat} of the digraph \mbox{\texttt{\mdseries\slshape digraph}}.  The value of the matrix entry \texttt{mat[i][j]} is the number of edges in \mbox{\texttt{\mdseries\slshape digraph}} with source \texttt{i} and range \texttt{j}. If \mbox{\texttt{\mdseries\slshape digraph}} has no vertices, then the empty list is returned. 

 The function \texttt{AdjacencyMatrix} returns an immutable list of lists, whereas the function \texttt{AdjacencyMatrixMutableCopy} returns a copy of \texttt{AdjacencyMatrix} that is a mutable list of mutable lists. 

 
\begin{Verbatim}[commandchars=!@|,fontsize=\small,frame=single,label=Example]
  !gapprompt@gap>| !gapinput@gr := Digraph([|
  !gapprompt@>| !gapinput@[2, 2, 2], [1, 3, 6, 8, 9, 10], [4, 6, 8],|
  !gapprompt@>| !gapinput@[1, 2, 3, 9], [3, 3], [3, 5, 6, 10], [1, 2, 7],|
  !gapprompt@>| !gapinput@[1, 2, 3, 10, 5, 6, 10], [1, 3, 4, 5, 8, 10],|
  !gapprompt@>| !gapinput@[2, 3, 4, 6, 7, 10]]);|
  <immutable multidigraph with 10 vertices, 44 edges>
  !gapprompt@gap>| !gapinput@mat := AdjacencyMatrix(gr);;|
  !gapprompt@gap>| !gapinput@Display(mat);|
  [ [  0,  3,  0,  0,  0,  0,  0,  0,  0,  0 ],
    [  1,  0,  1,  0,  0,  1,  0,  1,  1,  1 ],
    [  0,  0,  0,  1,  0,  1,  0,  1,  0,  0 ],
    [  1,  1,  1,  0,  0,  0,  0,  0,  1,  0 ],
    [  0,  0,  2,  0,  0,  0,  0,  0,  0,  0 ],
    [  0,  0,  1,  0,  1,  1,  0,  0,  0,  1 ],
    [  1,  1,  0,  0,  0,  0,  1,  0,  0,  0 ],
    [  1,  1,  1,  0,  1,  1,  0,  0,  0,  2 ],
    [  1,  0,  1,  1,  1,  0,  0,  1,  0,  1 ],
    [  0,  1,  1,  1,  0,  1,  1,  0,  0,  1 ] ]
  !gapprompt@gap>| !gapinput@D := CycleDigraph(IsMutableDigraph, 3);|
  <mutable digraph with 3 vertices, 3 edges>
  !gapprompt@gap>| !gapinput@Display(AdjacencyMatrix(D));|
  [ [  0,  1,  0 ],
    [  0,  0,  1 ],
    [  1,  0,  0 ] ]
\end{Verbatim}
 }

 

\subsection{\textcolor{Chapter }{CharacteristicPolynomial}}
\logpage{[ 5, 2, 2 ]}\nobreak
\hyperdef{L}{X87FA0A727CDB060B}{}
{\noindent\textcolor{FuncColor}{$\triangleright$\enspace\texttt{CharacteristicPolynomial({\mdseries\slshape digraph})\index{CharacteristicPolynomial@\texttt{CharacteristicPolynomial}}
\label{CharacteristicPolynomial}
}\hfill{\scriptsize (attribute)}}\\
\textbf{\indent Returns:\ }
A polynomial with integer coefficients.



 This function returns the characteristic polynomial of the digraph \mbox{\texttt{\mdseries\slshape digraph}}. That is it returns the characteristic polynomial of the adjacency matrix of
the digraph \mbox{\texttt{\mdseries\slshape digraph}} 
\begin{Verbatim}[commandchars=!@|,fontsize=\small,frame=single,label=Example]
  !gapprompt@gap>| !gapinput@D := Digraph([|
  !gapprompt@>| !gapinput@[2, 2, 2], [1, 3, 6, 8, 9, 10], [4, 6, 8],|
  !gapprompt@>| !gapinput@[1, 2, 3, 9], [3, 3], [3, 5, 6, 10], [1, 2, 7],|
  !gapprompt@>| !gapinput@[1, 2, 3, 10, 5, 6, 10], [1, 3, 4, 5, 8, 10],|
  !gapprompt@>| !gapinput@[2, 3, 4, 6, 7, 10]]);|
  <immutable multidigraph with 10 vertices, 44 edges>
  !gapprompt@gap>| !gapinput@CharacteristicPolynomial(D);|
  x_1^10-3*x_1^9-7*x_1^8-x_1^7+14*x_1^6+x_1^5-26*x_1^4+51*x_1^3-10*x_1^2\
  +18*x_1-30
  !gapprompt@gap>| !gapinput@D := CompleteDigraph(5);|
  <immutable complete digraph with 5 vertices>
  !gapprompt@gap>| !gapinput@CharacteristicPolynomial(D);|
  x_1^5-10*x_1^3-20*x_1^2-15*x_1-4
  !gapprompt@gap>| !gapinput@D := CycleDigraph(IsMutableDigraph, 3);|
  <mutable digraph with 3 vertices, 3 edges>
  !gapprompt@gap>| !gapinput@CharacteristicPolynomial(D);|
  x_1^3-1
\end{Verbatim}
 }

 

\subsection{\textcolor{Chapter }{BooleanAdjacencyMatrix}}
\logpage{[ 5, 2, 3 ]}\nobreak
\hyperdef{L}{X8507DC4F794780C1}{}
{\noindent\textcolor{FuncColor}{$\triangleright$\enspace\texttt{BooleanAdjacencyMatrix({\mdseries\slshape digraph})\index{BooleanAdjacencyMatrix@\texttt{BooleanAdjacencyMatrix}}
\label{BooleanAdjacencyMatrix}
}\hfill{\scriptsize (attribute)}}\\
\noindent\textcolor{FuncColor}{$\triangleright$\enspace\texttt{BooleanAdjacencyMatrixMutableCopy({\mdseries\slshape digraph})\index{BooleanAdjacencyMatrixMutableCopy@\texttt{BooleanAdjacencyMatrixMutableCopy}}
\label{BooleanAdjacencyMatrixMutableCopy}
}\hfill{\scriptsize (operation)}}\\
\textbf{\indent Returns:\ }
A square matrix of booleans.



 If \mbox{\texttt{\mdseries\slshape digraph}} is a digraph with a positive number of vertices \texttt{n}, then \texttt{BooleanAdjacencyMatrix(}\mbox{\texttt{\mdseries\slshape digraph}}\texttt{)} returns the boolean adjacency matrix \texttt{mat} of \mbox{\texttt{\mdseries\slshape digraph}}.  The value of the matrix entry \texttt{mat[j][i]} is \texttt{true} if and only if there exists an edge in \mbox{\texttt{\mdseries\slshape digraph}} with source \texttt{j} and range \texttt{i}. If \mbox{\texttt{\mdseries\slshape digraph}} has no vertices, then the empty list is returned. 

  Note that the boolean adjacency matrix loses information about multiple edges. 

 The attribute \texttt{BooleanAdjacencyMatrix} returns an immutable list of immutable lists, whereas the function \texttt{BooleanAdjacencyMatrixMutableCopy} returns a copy of the \texttt{BooleanAdjacencyMatrix} that is a mutable list of mutable lists. 

 
\begin{Verbatim}[commandchars=!@|,fontsize=\small,frame=single,label=Example]
  !gapprompt@gap>| !gapinput@gr := Digraph([[3, 4], [2, 3], [1, 2, 4], [4]]);|
  <immutable digraph with 4 vertices, 8 edges>
  !gapprompt@gap>| !gapinput@PrintArray(BooleanAdjacencyMatrix(gr));|
  [ [  false,  false,   true,   true ],
    [  false,   true,   true,  false ],
    [   true,   true,  false,   true ],
    [  false,  false,  false,   true ] ]
  !gapprompt@gap>| !gapinput@gr := CycleDigraph(4);;|
  !gapprompt@gap>| !gapinput@PrintArray(BooleanAdjacencyMatrix(gr));|
  [ [  false,   true,  false,  false ],
    [  false,  false,   true,  false ],
    [  false,  false,  false,   true ],
    [   true,  false,  false,  false ] ]
  !gapprompt@gap>| !gapinput@BooleanAdjacencyMatrix(EmptyDigraph(0));|
  [  ]
  !gapprompt@gap>| !gapinput@D := CycleDigraph(IsMutableDigraph, 3);|
  <mutable digraph with 3 vertices, 3 edges>
  !gapprompt@gap>| !gapinput@PrintArray(BooleanAdjacencyMatrix(D));|
  [ [  false,   true,  false ],
    [  false,  false,   true ],
    [   true,  false,  false ] ]
\end{Verbatim}
 }

 

\subsection{\textcolor{Chapter }{DigraphAdjacencyFunction}}
\logpage{[ 5, 2, 4 ]}\nobreak
\hyperdef{L}{X7AFCE34A7A04D5C1}{}
{\noindent\textcolor{FuncColor}{$\triangleright$\enspace\texttt{DigraphAdjacencyFunction({\mdseries\slshape digraph})\index{DigraphAdjacencyFunction@\texttt{DigraphAdjacencyFunction}}
\label{DigraphAdjacencyFunction}
}\hfill{\scriptsize (attribute)}}\\
\textbf{\indent Returns:\ }
A function.



 If \mbox{\texttt{\mdseries\slshape digraph}} is a digraph, then \texttt{DigraphAdjacencyFunction} returns a function which takes two integer parameters \texttt{x, y} and returns \texttt{true} if there exists an edge from vertex \texttt{x} to vertex \texttt{y} in \mbox{\texttt{\mdseries\slshape digraph}} and \texttt{false} if not. 
\begin{Verbatim}[commandchars=!@|,fontsize=\small,frame=single,label=Example]
  !gapprompt@gap>| !gapinput@digraph := Digraph([[1, 2], [3], []]);|
  <immutable digraph with 3 vertices, 3 edges>
  !gapprompt@gap>| !gapinput@foo := DigraphAdjacencyFunction(digraph);|
  function( u, v ) ... end
  !gapprompt@gap>| !gapinput@foo(1, 1);|
  true
  !gapprompt@gap>| !gapinput@foo(1, 2);|
  true
  !gapprompt@gap>| !gapinput@foo(1, 3);|
  false
  !gapprompt@gap>| !gapinput@foo(3, 1);|
  false
  !gapprompt@gap>| !gapinput@gr := Digraph(["a", "b", "c"],|
  !gapprompt@>| !gapinput@                 ["a", "b", "b"],|
  !gapprompt@>| !gapinput@                 ["b", "a", "a"]);|
  <immutable multidigraph with 3 vertices, 3 edges>
  !gapprompt@gap>| !gapinput@foo := DigraphAdjacencyFunction(gr);|
  function( u, v ) ... end
  !gapprompt@gap>| !gapinput@foo(1, 2);|
  true
  !gapprompt@gap>| !gapinput@foo(3, 2);|
  false
  !gapprompt@gap>| !gapinput@foo(3, 1);|
  false
  !gapprompt@gap>| !gapinput@D := CycleDigraph(IsMutableDigraph, 3);|
  <mutable digraph with 3 vertices, 3 edges>
  !gapprompt@gap>| !gapinput@foo := DigraphAdjacencyFunction(D);|
  function( u, v ) ... end
  !gapprompt@gap>| !gapinput@foo(1, 2);|
  true
  !gapprompt@gap>| !gapinput@foo(2, 1);|
  false
\end{Verbatim}
 }

 

\subsection{\textcolor{Chapter }{DigraphRange}}
\logpage{[ 5, 2, 5 ]}\nobreak
\hyperdef{L}{X7FDEBF3279759961}{}
{\noindent\textcolor{FuncColor}{$\triangleright$\enspace\texttt{DigraphRange({\mdseries\slshape digraph})\index{DigraphRange@\texttt{DigraphRange}}
\label{DigraphRange}
}\hfill{\scriptsize (attribute)}}\\
\noindent\textcolor{FuncColor}{$\triangleright$\enspace\texttt{DigraphSource({\mdseries\slshape digraph})\index{DigraphSource@\texttt{DigraphSource}}
\label{DigraphSource}
}\hfill{\scriptsize (attribute)}}\\
\textbf{\indent Returns:\ }
A list of positive integers.



 \texttt{DigraphRange} and \texttt{DigraphSource} return the range and source of the digraph \mbox{\texttt{\mdseries\slshape digraph}}. More precisely, position \texttt{i} in \texttt{DigraphSource(\mbox{\texttt{\mdseries\slshape digraph}})} and \texttt{DigraphRange(\mbox{\texttt{\mdseries\slshape digraph}})} give, respectively, the source and range of the \texttt{i}th edge of \mbox{\texttt{\mdseries\slshape digraph}}. 
\begin{Verbatim}[commandchars=!@|,fontsize=\small,frame=single,label=Example]
  !gapprompt@gap>| !gapinput@gr := Digraph([|
  !gapprompt@>| !gapinput@[2, 1, 3, 5], [1, 3, 4], [2, 3], [2], [1, 2, 3, 4]]);|
  <immutable digraph with 5 vertices, 14 edges>
  !gapprompt@gap>| !gapinput@DigraphSource(gr);|
  [ 1, 1, 1, 1, 2, 2, 2, 3, 3, 4, 5, 5, 5, 5 ]
  !gapprompt@gap>| !gapinput@DigraphRange(gr);|
  [ 2, 1, 3, 5, 1, 3, 4, 2, 3, 2, 1, 2, 3, 4 ]
  !gapprompt@gap>| !gapinput@DigraphEdges(gr);|
  [ [ 1, 2 ], [ 1, 1 ], [ 1, 3 ], [ 1, 5 ], [ 2, 1 ], [ 2, 3 ], 
    [ 2, 4 ], [ 3, 2 ], [ 3, 3 ], [ 4, 2 ], [ 5, 1 ], [ 5, 2 ], 
    [ 5, 3 ], [ 5, 4 ] ]
  !gapprompt@gap>| !gapinput@D := CycleDigraph(IsMutableDigraph, 3);|
  <mutable digraph with 3 vertices, 3 edges>
  !gapprompt@gap>| !gapinput@DigraphRange(D);|
  [ 2, 3, 1 ]
  !gapprompt@gap>| !gapinput@DigraphSource(D);|
  [ 1, 2, 3 ]
\end{Verbatim}
 }

 

\subsection{\textcolor{Chapter }{OutNeighbours}}
\logpage{[ 5, 2, 6 ]}\nobreak
\hyperdef{L}{X7E9880767AE68E00}{}
{\noindent\textcolor{FuncColor}{$\triangleright$\enspace\texttt{OutNeighbours({\mdseries\slshape digraph})\index{OutNeighbours@\texttt{OutNeighbours}}
\label{OutNeighbours}
}\hfill{\scriptsize (attribute)}}\\
\noindent\textcolor{FuncColor}{$\triangleright$\enspace\texttt{OutNeighbors({\mdseries\slshape digraph})\index{OutNeighbors@\texttt{OutNeighbors}}
\label{OutNeighbors}
}\hfill{\scriptsize (attribute)}}\\
\noindent\textcolor{FuncColor}{$\triangleright$\enspace\texttt{OutNeighboursMutableCopy({\mdseries\slshape digraph})\index{OutNeighboursMutableCopy@\texttt{OutNeighboursMutableCopy}}
\label{OutNeighboursMutableCopy}
}\hfill{\scriptsize (operation)}}\\
\noindent\textcolor{FuncColor}{$\triangleright$\enspace\texttt{OutNeighborsMutableCopy({\mdseries\slshape digraph})\index{OutNeighborsMutableCopy@\texttt{OutNeighborsMutableCopy}}
\label{OutNeighborsMutableCopy}
}\hfill{\scriptsize (operation)}}\\
\textbf{\indent Returns:\ }
The adjacencies of a digraph.



 \texttt{OutNeighbours} returns the list \texttt{out} of out\texttt{\symbol{45}}neighbours of each vertex of the digraph \mbox{\texttt{\mdseries\slshape digraph}}.  More specifically, a vertex \texttt{j} appears in \texttt{out[i]} each time there exists the edge \texttt{[i, j]} in \mbox{\texttt{\mdseries\slshape digraph}}. 

 The function \texttt{OutNeighbours} returns an immutable list of lists, whereas the function \texttt{OutNeighboursMutableCopy} returns a copy of \texttt{OutNeighbours} which is a mutable list of mutable lists. 

 Note that the entries of \texttt{out} are not guaranteed to be sorted in any particular order. 
\begin{Verbatim}[commandchars=!@|,fontsize=\small,frame=single,label=Example]
  !gapprompt@gap>| !gapinput@gr := Digraph(["a", "b", "c"],|
  !gapprompt@>| !gapinput@                 ["a", "b", "b"],|
  !gapprompt@>| !gapinput@                 ["b", "a", "c"]);|
  <immutable digraph with 3 vertices, 3 edges>
  !gapprompt@gap>| !gapinput@OutNeighbours(gr);|
  [ [ 2 ], [ 1, 3 ], [  ] ]
  !gapprompt@gap>| !gapinput@gr := Digraph([[1, 2, 3], [2, 1], [3]]);|
  <immutable digraph with 3 vertices, 6 edges>
  !gapprompt@gap>| !gapinput@OutNeighbours(gr);|
  [ [ 1, 2, 3 ], [ 2, 1 ], [ 3 ] ]
  !gapprompt@gap>| !gapinput@gr := DigraphByAdjacencyMatrix([|
  !gapprompt@>| !gapinput@ [1, 2, 1],|
  !gapprompt@>| !gapinput@ [1, 1, 0],|
  !gapprompt@>| !gapinput@ [0, 0, 1]]);|
  <immutable multidigraph with 3 vertices, 7 edges>
  !gapprompt@gap>| !gapinput@OutNeighbours(gr);|
  [ [ 1, 2, 2, 3 ], [ 1, 2 ], [ 3 ] ]
  !gapprompt@gap>| !gapinput@OutNeighboursMutableCopy(gr);|
  [ [ 1, 2, 2, 3 ], [ 1, 2 ], [ 3 ] ]
  !gapprompt@gap>| !gapinput@D := CycleDigraph(IsMutableDigraph, 3);|
  <mutable digraph with 3 vertices, 3 edges>
  !gapprompt@gap>| !gapinput@OutNeighbours(D);|
  [ [ 2 ], [ 3 ], [ 1 ] ]
\end{Verbatim}
 }

 

\subsection{\textcolor{Chapter }{InNeighbours}}
\logpage{[ 5, 2, 7 ]}\nobreak
\hyperdef{L}{X85C7AA5A81DA6E11}{}
{\noindent\textcolor{FuncColor}{$\triangleright$\enspace\texttt{InNeighbours({\mdseries\slshape digraph})\index{InNeighbours@\texttt{InNeighbours}}
\label{InNeighbours}
}\hfill{\scriptsize (attribute)}}\\
\noindent\textcolor{FuncColor}{$\triangleright$\enspace\texttt{InNeighbors({\mdseries\slshape digraph})\index{InNeighbors@\texttt{InNeighbors}}
\label{InNeighbors}
}\hfill{\scriptsize (attribute)}}\\
\noindent\textcolor{FuncColor}{$\triangleright$\enspace\texttt{InNeighboursMutableCopy({\mdseries\slshape digraph})\index{InNeighboursMutableCopy@\texttt{InNeighboursMutableCopy}}
\label{InNeighboursMutableCopy}
}\hfill{\scriptsize (operation)}}\\
\noindent\textcolor{FuncColor}{$\triangleright$\enspace\texttt{InNeighborsMutableCopy({\mdseries\slshape digraph})\index{InNeighborsMutableCopy@\texttt{InNeighborsMutableCopy}}
\label{InNeighborsMutableCopy}
}\hfill{\scriptsize (operation)}}\\
\textbf{\indent Returns:\ }
A list of lists of vertices.



 \texttt{InNeighbours} returns the list \texttt{inn} of in\texttt{\symbol{45}}neighbours of each vertex of the digraph \mbox{\texttt{\mdseries\slshape digraph}}.  More specifically, a vertex \texttt{j} appears in \texttt{inn[i]} each time there exists an edge \texttt{[j,i]} in \mbox{\texttt{\mdseries\slshape digraph}}. 

 The function \texttt{InNeighbours} returns an immutable list of lists, whereas the function \texttt{InNeighboursMutableCopy} returns a copy of \texttt{InNeighbours} which is a mutable list of mutable lists. 

 Note that the entries of \texttt{inn} are not necessarily sorted into ascending order, particularly if \mbox{\texttt{\mdseries\slshape digraph}} was constructed via \texttt{DigraphByInNeighbours} (\ref{DigraphByInNeighbours}). 
\begin{Verbatim}[commandchars=!@|,fontsize=\small,frame=single,label=Example]
  !gapprompt@gap>| !gapinput@gr := Digraph(["a", "b", "c"],|
  !gapprompt@>| !gapinput@                 ["a", "b", "b"],|
  !gapprompt@>| !gapinput@                 ["b", "a", "c"]);|
  <immutable digraph with 3 vertices, 3 edges>
  !gapprompt@gap>| !gapinput@InNeighbours(gr);|
  [ [ 2 ], [ 1 ], [ 2 ] ]
  !gapprompt@gap>| !gapinput@gr := Digraph([[1, 2, 3], [2, 1], [3]]);|
  <immutable digraph with 3 vertices, 6 edges>
  !gapprompt@gap>| !gapinput@InNeighbours(gr);|
  [ [ 1, 2 ], [ 1, 2 ], [ 1, 3 ] ]
  !gapprompt@gap>| !gapinput@gr := DigraphByAdjacencyMatrix([|
  !gapprompt@>| !gapinput@ [1, 2, 1],|
  !gapprompt@>| !gapinput@ [1, 1, 0],|
  !gapprompt@>| !gapinput@ [0, 0, 1]]);|
  <immutable multidigraph with 3 vertices, 7 edges>
  !gapprompt@gap>| !gapinput@InNeighbours(gr);|
  [ [ 1, 2 ], [ 1, 1, 2 ], [ 1, 3 ] ]
  !gapprompt@gap>| !gapinput@InNeighboursMutableCopy(gr);|
  [ [ 1, 2 ], [ 1, 1, 2 ], [ 1, 3 ] ]
  !gapprompt@gap>| !gapinput@D := CycleDigraph(IsMutableDigraph, 3);|
  <mutable digraph with 3 vertices, 3 edges>
  !gapprompt@gap>| !gapinput@InNeighbours(D);|
  [ [ 3 ], [ 1 ], [ 2 ] ]
\end{Verbatim}
 }

 

\subsection{\textcolor{Chapter }{OutDegrees}}
\logpage{[ 5, 2, 8 ]}\nobreak
\hyperdef{L}{X7F5ACE807D1BC2E2}{}
{\noindent\textcolor{FuncColor}{$\triangleright$\enspace\texttt{OutDegrees({\mdseries\slshape digraph})\index{OutDegrees@\texttt{OutDegrees}}
\label{OutDegrees}
}\hfill{\scriptsize (attribute)}}\\
\noindent\textcolor{FuncColor}{$\triangleright$\enspace\texttt{OutDegreeSequence({\mdseries\slshape digraph})\index{OutDegreeSequence@\texttt{OutDegreeSequence}}
\label{OutDegreeSequence}
}\hfill{\scriptsize (attribute)}}\\
\noindent\textcolor{FuncColor}{$\triangleright$\enspace\texttt{OutDegreeSet({\mdseries\slshape digraph})\index{OutDegreeSet@\texttt{OutDegreeSet}}
\label{OutDegreeSet}
}\hfill{\scriptsize (attribute)}}\\
\textbf{\indent Returns:\ }
A list of non\texttt{\symbol{45}}negative integers.



 Given a digraph \mbox{\texttt{\mdseries\slshape digraph}} with $n$ vertices, the function \texttt{OutDegrees} returns an immutable list \texttt{out} of length $n$, such that for a vertex \texttt{i} in \mbox{\texttt{\mdseries\slshape digraph}}, the value of \texttt{out[i]} is the out\texttt{\symbol{45}}degree of vertex \texttt{i}. See \texttt{OutDegreeOfVertex} (\ref{OutDegreeOfVertex}). 

 The function \texttt{OutDegreeSequence} returns the same list, after it has been sorted into
non\texttt{\symbol{45}}increasing order. 

 The function \texttt{OutDegreeSet} returns the same list, sorted into increasing order with duplicate entries
removed. 

 
\begin{Verbatim}[commandchars=!@|,fontsize=\small,frame=single,label=Example]
  !gapprompt@gap>| !gapinput@D := Digraph([[1, 3, 2, 2], [], [2, 1], []]);|
  <immutable multidigraph with 4 vertices, 6 edges>
  !gapprompt@gap>| !gapinput@OutDegrees(D);|
  [ 4, 0, 2, 0 ]
  !gapprompt@gap>| !gapinput@OutDegreeSequence(D);|
  [ 4, 2, 0, 0 ]
  !gapprompt@gap>| !gapinput@OutDegreeSet(D);|
  [ 0, 2, 4 ]
  !gapprompt@gap>| !gapinput@D := CycleDigraph(IsMutableDigraph, 3);|
  <mutable digraph with 3 vertices, 3 edges>
  !gapprompt@gap>| !gapinput@OutDegrees(D);|
  [ 1, 1, 1 ]
\end{Verbatim}
 }

 

\subsection{\textcolor{Chapter }{InDegrees}}
\logpage{[ 5, 2, 9 ]}\nobreak
\hyperdef{L}{X7ADDFBFD7A365775}{}
{\noindent\textcolor{FuncColor}{$\triangleright$\enspace\texttt{InDegrees({\mdseries\slshape digraph})\index{InDegrees@\texttt{InDegrees}}
\label{InDegrees}
}\hfill{\scriptsize (attribute)}}\\
\noindent\textcolor{FuncColor}{$\triangleright$\enspace\texttt{InDegreeSequence({\mdseries\slshape digraph})\index{InDegreeSequence@\texttt{InDegreeSequence}}
\label{InDegreeSequence}
}\hfill{\scriptsize (attribute)}}\\
\noindent\textcolor{FuncColor}{$\triangleright$\enspace\texttt{InDegreeSet({\mdseries\slshape digraph})\index{InDegreeSet@\texttt{InDegreeSet}}
\label{InDegreeSet}
}\hfill{\scriptsize (attribute)}}\\
\textbf{\indent Returns:\ }
A list of non\texttt{\symbol{45}}negative integers.



 Given a digraph \mbox{\texttt{\mdseries\slshape digraph}} with $n$ vertices, the function \texttt{InDegrees} returns an immutable list \texttt{inn} of length $n$, such that for a vertex \texttt{i} in \mbox{\texttt{\mdseries\slshape digraph}}, the value of \texttt{inn[i]} is the in\texttt{\symbol{45}}degree of vertex \texttt{i}. See \texttt{InDegreeOfVertex} (\ref{InDegreeOfVertex}). 

 The function \texttt{InDegreeSequence} returns the same list, after it has been sorted into
non\texttt{\symbol{45}}increasing order. 

 The function \texttt{InDegreeSet} returns the same list, sorted into increasing order with duplicate entries
removed. 

 
\begin{Verbatim}[commandchars=!@|,fontsize=\small,frame=single,label=Example]
  !gapprompt@gap>| !gapinput@D := Digraph([[1, 3, 2, 2, 4], [], [2, 1, 4], []]);|
  <immutable multidigraph with 4 vertices, 8 edges>
  !gapprompt@gap>| !gapinput@InDegrees(D);|
  [ 2, 3, 1, 2 ]
  !gapprompt@gap>| !gapinput@InDegreeSequence(D);|
  [ 3, 2, 2, 1 ]
  !gapprompt@gap>| !gapinput@InDegreeSet(D);|
  [ 1, 2, 3 ]
  !gapprompt@gap>| !gapinput@D := CycleDigraph(IsMutableDigraph, 3);|
  <mutable digraph with 3 vertices, 3 edges>
  !gapprompt@gap>| !gapinput@InDegrees(D);|
  [ 1, 1, 1 ]
\end{Verbatim}
 }

 

\subsection{\textcolor{Chapter }{OutDegreeOfVertex}}
\logpage{[ 5, 2, 10 ]}\nobreak
\hyperdef{L}{X7A09EB648070276D}{}
{\noindent\textcolor{FuncColor}{$\triangleright$\enspace\texttt{OutDegreeOfVertex({\mdseries\slshape digraph, vertex})\index{OutDegreeOfVertex@\texttt{OutDegreeOfVertex}}
\label{OutDegreeOfVertex}
}\hfill{\scriptsize (operation)}}\\
\textbf{\indent Returns:\ }
The non\texttt{\symbol{45}}negative integer.



 This operation returns the out\texttt{\symbol{45}}degree of the vertex \mbox{\texttt{\mdseries\slshape vertex}} in the digraph \mbox{\texttt{\mdseries\slshape digraph}}. The out\texttt{\symbol{45}}degree of \mbox{\texttt{\mdseries\slshape vertex}} is the number of edges in \mbox{\texttt{\mdseries\slshape digraph}} whose source is \mbox{\texttt{\mdseries\slshape vertex}}. 

 
\begin{Verbatim}[commandchars=!@|,fontsize=\small,frame=single,label=Example]
  !gapprompt@gap>| !gapinput@D := Digraph([|
  !gapprompt@>| !gapinput@[2, 2, 1], [1, 4], [2, 2, 4, 2], [1, 1, 2, 2, 1, 2, 2]]);|
  <immutable multidigraph with 4 vertices, 16 edges>
  !gapprompt@gap>| !gapinput@OutDegreeOfVertex(D, 1);|
  3
  !gapprompt@gap>| !gapinput@OutDegreeOfVertex(D, 2);|
  2
  !gapprompt@gap>| !gapinput@OutDegreeOfVertex(D, 3);|
  4
  !gapprompt@gap>| !gapinput@OutDegreeOfVertex(D, 4);|
  7
\end{Verbatim}
 }

 

\subsection{\textcolor{Chapter }{OutNeighboursOfVertex}}
\logpage{[ 5, 2, 11 ]}\nobreak
\hyperdef{L}{X83315B0186850806}{}
{\noindent\textcolor{FuncColor}{$\triangleright$\enspace\texttt{OutNeighboursOfVertex({\mdseries\slshape digraph, vertex})\index{OutNeighboursOfVertex@\texttt{OutNeighboursOfVertex}}
\label{OutNeighboursOfVertex}
}\hfill{\scriptsize (operation)}}\\
\noindent\textcolor{FuncColor}{$\triangleright$\enspace\texttt{OutNeighborsOfVertex({\mdseries\slshape digraph, vertex})\index{OutNeighborsOfVertex@\texttt{OutNeighborsOfVertex}}
\label{OutNeighborsOfVertex}
}\hfill{\scriptsize (operation)}}\\
\textbf{\indent Returns:\ }
A list of vertices.



 This operation returns the list \texttt{out} of vertices of the digraph \mbox{\texttt{\mdseries\slshape digraph}}. A vertex \texttt{i} appears in the list \texttt{out} each time there exists an edge with source \mbox{\texttt{\mdseries\slshape vertex}} and range \texttt{i} in \mbox{\texttt{\mdseries\slshape digraph}}; in particular, this means that \texttt{out} may contain duplicates.

 
\begin{Verbatim}[commandchars=!@|,fontsize=\small,frame=single,label=Example]
  !gapprompt@gap>| !gapinput@D := Digraph([|
  !gapprompt@>| !gapinput@[2, 2, 3], [1, 3, 4], [2, 2, 3], [1, 1, 2, 2, 1, 2, 2]]);|
  <immutable multidigraph with 4 vertices, 16 edges>
  !gapprompt@gap>| !gapinput@OutNeighboursOfVertex(D, 1);|
  [ 2, 2, 3 ]
  !gapprompt@gap>| !gapinput@OutNeighboursOfVertex(D, 3);|
  [ 2, 2, 3 ]
\end{Verbatim}
 }

 

\subsection{\textcolor{Chapter }{InDegreeOfVertex}}
\logpage{[ 5, 2, 12 ]}\nobreak
\hyperdef{L}{X7C9CD0527CB9E6EF}{}
{\noindent\textcolor{FuncColor}{$\triangleright$\enspace\texttt{InDegreeOfVertex({\mdseries\slshape digraph, vertex})\index{InDegreeOfVertex@\texttt{InDegreeOfVertex}}
\label{InDegreeOfVertex}
}\hfill{\scriptsize (operation)}}\\
\textbf{\indent Returns:\ }
A non\texttt{\symbol{45}}negative integer.



 This operation returns the in\texttt{\symbol{45}}degree of the vertex \mbox{\texttt{\mdseries\slshape vertex}} in the digraph \mbox{\texttt{\mdseries\slshape digraph}}. The in\texttt{\symbol{45}}degree of \mbox{\texttt{\mdseries\slshape vertex}} is the number of edges in \mbox{\texttt{\mdseries\slshape digraph}} whose range is \mbox{\texttt{\mdseries\slshape vertex}}. 
\begin{Verbatim}[commandchars=!@|,fontsize=\small,frame=single,label=Example]
  !gapprompt@gap>| !gapinput@D := Digraph([|
  !gapprompt@>| !gapinput@[2, 2, 1], [1, 4], [2, 2, 4, 2], [1, 1, 2, 2, 1, 2, 2]]);|
  <immutable multidigraph with 4 vertices, 16 edges>
  !gapprompt@gap>| !gapinput@InDegreeOfVertex(D, 1);|
  5
  !gapprompt@gap>| !gapinput@InDegreeOfVertex(D, 2);|
  9
  !gapprompt@gap>| !gapinput@InDegreeOfVertex(D, 3);|
  0
  !gapprompt@gap>| !gapinput@InDegreeOfVertex(D, 4);|
  2
\end{Verbatim}
 }

 

\subsection{\textcolor{Chapter }{InNeighboursOfVertex}}
\logpage{[ 5, 2, 13 ]}\nobreak
\hyperdef{L}{X7C0DA18B8291F302}{}
{\noindent\textcolor{FuncColor}{$\triangleright$\enspace\texttt{InNeighboursOfVertex({\mdseries\slshape digraph, vertex})\index{InNeighboursOfVertex@\texttt{InNeighboursOfVertex}}
\label{InNeighboursOfVertex}
}\hfill{\scriptsize (operation)}}\\
\noindent\textcolor{FuncColor}{$\triangleright$\enspace\texttt{InNeighborsOfVertex({\mdseries\slshape digraph, vertex})\index{InNeighborsOfVertex@\texttt{InNeighborsOfVertex}}
\label{InNeighborsOfVertex}
}\hfill{\scriptsize (operation)}}\\
\textbf{\indent Returns:\ }
A list of positive vertices.



 This operation returns the list \texttt{inn} of vertices of the digraph \mbox{\texttt{\mdseries\slshape digraph}}. A vertex \texttt{i} appears in the list \texttt{inn} each time there exists an edge with source \texttt{i} and range \mbox{\texttt{\mdseries\slshape vertex}} in \mbox{\texttt{\mdseries\slshape digraph}}; in particular, this means that \texttt{inn} may contain duplicates. 

 
\begin{Verbatim}[commandchars=!@|,fontsize=\small,frame=single,label=Example]
  !gapprompt@gap>| !gapinput@D := Digraph([|
  !gapprompt@>| !gapinput@[2, 2, 3], [1, 3, 4], [2, 2, 3], [1, 1, 2, 2, 1, 2, 2]]);|
  <immutable multidigraph with 4 vertices, 16 edges>
  !gapprompt@gap>| !gapinput@InNeighboursOfVertex(D, 1);|
  [ 2, 4, 4, 4 ]
  !gapprompt@gap>| !gapinput@InNeighboursOfVertex(D, 2);|
  [ 1, 1, 3, 3, 4, 4, 4, 4 ]
\end{Verbatim}
 }

 

\subsection{\textcolor{Chapter }{DigraphLoops}}
\logpage{[ 5, 2, 14 ]}\nobreak
\hyperdef{L}{X83271F607BD809CF}{}
{\noindent\textcolor{FuncColor}{$\triangleright$\enspace\texttt{DigraphLoops({\mdseries\slshape digraph})\index{DigraphLoops@\texttt{DigraphLoops}}
\label{DigraphLoops}
}\hfill{\scriptsize (attribute)}}\\
\textbf{\indent Returns:\ }
A list of vertices.



 If \mbox{\texttt{\mdseries\slshape digraph}} is a digraph, then \texttt{DigraphLoops} returns the list consisting of the \texttt{DigraphVertices} (\ref{DigraphVertices}) of \mbox{\texttt{\mdseries\slshape digraph}} at which there is a loop. See \texttt{DigraphHasLoops} (\ref{DigraphHasLoops}). 

 
\begin{Verbatim}[commandchars=!@|,fontsize=\small,frame=single,label=Example]
  !gapprompt@gap>| !gapinput@D := Digraph([[2], [3], []]);|
  <immutable digraph with 3 vertices, 2 edges>
  !gapprompt@gap>| !gapinput@DigraphHasLoops(D);|
  false
  !gapprompt@gap>| !gapinput@DigraphLoops(D);|
  [  ]
  !gapprompt@gap>| !gapinput@D := Digraph([[3, 5], [1], [2, 4, 3], [4], [2, 1]]);|
  <immutable digraph with 5 vertices, 9 edges>
  !gapprompt@gap>| !gapinput@DigraphLoops(D);|
  [ 3, 4 ]
  !gapprompt@gap>| !gapinput@D := Digraph(IsMutableDigraph, [[1], [1]]);|
  <mutable digraph with 2 vertices, 2 edges>
  !gapprompt@gap>| !gapinput@DigraphLoops(D);|
  [ 1 ]
\end{Verbatim}
 }

 

\subsection{\textcolor{Chapter }{DegreeMatrix}}
\logpage{[ 5, 2, 15 ]}\nobreak
\hyperdef{L}{X7BEAE1C78267F54D}{}
{\noindent\textcolor{FuncColor}{$\triangleright$\enspace\texttt{DegreeMatrix({\mdseries\slshape digraph})\index{DegreeMatrix@\texttt{DegreeMatrix}}
\label{DegreeMatrix}
}\hfill{\scriptsize (attribute)}}\\
\textbf{\indent Returns:\ }
A square matrix of non\texttt{\symbol{45}}negative integers.



 Returns the out\texttt{\symbol{45}}degree matrix \texttt{mat} of the digraph \mbox{\texttt{\mdseries\slshape digraph}}. The value of the \texttt{i}th diagonal matrix entry is the out\texttt{\symbol{45}}degree of the vertex \texttt{i} in \mbox{\texttt{\mdseries\slshape digraph}}. All off\texttt{\symbol{45}}diagonal entries are \texttt{0}. If \mbox{\texttt{\mdseries\slshape digraph}} has no vertices, then the empty list is returned. 

 See \texttt{OutDegrees} (\ref{OutDegrees}) for more information. 
\begin{Verbatim}[commandchars=!@|,fontsize=\small,frame=single,label=Example]
  !gapprompt@gap>| !gapinput@D := Digraph([[1, 2, 3], [4], [1, 3, 4], []]);|
  <immutable digraph with 4 vertices, 7 edges>
  !gapprompt@gap>| !gapinput@PrintArray(DegreeMatrix(D));|
  [ [  3,  0,  0,  0 ],
    [  0,  1,  0,  0 ],
    [  0,  0,  3,  0 ],
    [  0,  0,  0,  0 ] ]
  !gapprompt@gap>| !gapinput@D := CycleDigraph(5);;|
  !gapprompt@gap>| !gapinput@PrintArray(DegreeMatrix(D));|
  [ [  1,  0,  0,  0,  0 ],
    [  0,  1,  0,  0,  0 ],
    [  0,  0,  1,  0,  0 ],
    [  0,  0,  0,  1,  0 ],
    [  0,  0,  0,  0,  1 ] ]
  !gapprompt@gap>| !gapinput@DegreeMatrix(EmptyDigraph(0));|
  [  ]
  !gapprompt@gap>| !gapinput@D := CycleDigraph(IsMutableDigraph, 3);|
  <mutable digraph with 3 vertices, 3 edges>
  !gapprompt@gap>| !gapinput@Display(DegreeMatrix(D));|
  [ [  1,  0,  0 ],
    [  0,  1,  0 ],
    [  0,  0,  1 ] ]
\end{Verbatim}
 }

 

\subsection{\textcolor{Chapter }{LaplacianMatrix}}
\logpage{[ 5, 2, 16 ]}\nobreak
\hyperdef{L}{X865390B08331936B}{}
{\noindent\textcolor{FuncColor}{$\triangleright$\enspace\texttt{LaplacianMatrix({\mdseries\slshape digraph})\index{LaplacianMatrix@\texttt{LaplacianMatrix}}
\label{LaplacianMatrix}
}\hfill{\scriptsize (attribute)}}\\
\textbf{\indent Returns:\ }
A square matrix of integers.



 Returns the out\texttt{\symbol{45}}degree Laplacian matrix \texttt{mat} of the digraph \mbox{\texttt{\mdseries\slshape digraph}}. The out\texttt{\symbol{45}}degree Laplacian matrix is defined as \texttt{DegreeMatrix(digraph) \texttt{\symbol{45}} AdjacencyMatrix(digraph)}. If \mbox{\texttt{\mdseries\slshape digraph}} has no vertices, then the empty list is returned. 

 See \texttt{DegreeMatrix} (\ref{DegreeMatrix}) and \texttt{AdjacencyMatrix} (\ref{AdjacencyMatrix}) for more information. 
\begin{Verbatim}[commandchars=!@|,fontsize=\small,frame=single,label=Example]
  !gapprompt@gap>| !gapinput@gr := Digraph([[1, 2, 3], [4], [1, 3, 4], []]);|
  <immutable digraph with 4 vertices, 7 edges>
  !gapprompt@gap>| !gapinput@PrintArray(LaplacianMatrix(gr));|
  [ [   2,  -1,  -1,   0 ],
    [   0,   1,   0,  -1 ],
    [  -1,   0,   2,  -1 ],
    [   0,   0,   0,   0 ] ]
  !gapprompt@gap>| !gapinput@LaplacianMatrix(EmptyDigraph(0));|
  [  ]
  !gapprompt@gap>| !gapinput@D := CycleDigraph(IsMutableDigraph, 3);|
  <mutable digraph with 3 vertices, 3 edges>
  !gapprompt@gap>| !gapinput@Display(LaplacianMatrix(D));|
  [ [   1,  -1,   0 ],
    [   0,   1,  -1 ],
    [  -1,   0,   1 ] ]
\end{Verbatim}
 }

 }

 
\section{\textcolor{Chapter }{Orders}}\logpage{[ 5, 3, 0 ]}
\hyperdef{L}{X86424F167BD4F629}{}
{
 

\subsection{\textcolor{Chapter }{PartialOrderDigraphMeetOfVertices}}
\logpage{[ 5, 3, 1 ]}\nobreak
\hyperdef{L}{X7DDB33B686B3A2C6}{}
{\noindent\textcolor{FuncColor}{$\triangleright$\enspace\texttt{PartialOrderDigraphMeetOfVertices({\mdseries\slshape digraph, u, v})\index{PartialOrderDigraphMeetOfVertices@\texttt{PartialOrderDigraphMeetOfVertices}}
\label{PartialOrderDigraphMeetOfVertices}
}\hfill{\scriptsize (operation)}}\\
\noindent\textcolor{FuncColor}{$\triangleright$\enspace\texttt{PartialOrderDigraphJoinOfVertices({\mdseries\slshape digraph, u, v})\index{PartialOrderDigraphJoinOfVertices@\texttt{PartialOrderDigraphJoinOfVertices}}
\label{PartialOrderDigraphJoinOfVertices}
}\hfill{\scriptsize (operation)}}\\
\textbf{\indent Returns:\ }
A positive integer or \texttt{fail}



 If the first argument is a partial order digraph \texttt{IsPartialOrderDigraph} (\ref{IsPartialOrderDigraph}) then these operations return the meet, or the join, of the two input vertices.
If the meet (or join) is does not exist then \texttt{fail} is returned. The meet (or join) is guaranteed to exist when the first argument
satisfies \texttt{IsMeetSemilatticeDigraph} (\ref{IsMeetSemilatticeDigraph}) (or \texttt{IsJoinSemilatticeDigraph} (\ref{IsJoinSemilatticeDigraph})) \texttt{\symbol{45}} see the documentation for these properties for the
definition of the meet (or the join). 
\begin{Verbatim}[commandchars=!@|,fontsize=\small,frame=single,label=Example]
  !gapprompt@gap>| !gapinput@D := Digraph([[1], [1, 2], [1, 3], [1, 2, 3, 4]]);|
  <immutable digraph with 4 vertices, 9 edges>
  !gapprompt@gap>| !gapinput@PartialOrderDigraphMeetOfVertices(D, 2, 3);|
  4
  !gapprompt@gap>| !gapinput@PartialOrderDigraphJoinOfVertices(D, 2, 3);|
  1
  !gapprompt@gap>| !gapinput@PartialOrderDigraphMeetOfVertices(D, 1, 2);|
  2
  !gapprompt@gap>| !gapinput@PartialOrderDigraphJoinOfVertices(D, 3, 4);|
  3
  !gapprompt@gap>| !gapinput@D := Digraph([[1], [2], [1, 2, 3], [1, 2, 4]]);|
  <immutable digraph with 4 vertices, 8 edges>
  !gapprompt@gap>| !gapinput@PartialOrderDigraphMeetOfVertices(D, 3, 4);|
  fail
  !gapprompt@gap>| !gapinput@PartialOrderDigraphJoinOfVertices(D, 3, 4);|
  fail
  !gapprompt@gap>| !gapinput@PartialOrderDigraphMeetOfVertices(D, 1, 2);|
  fail
  !gapprompt@gap>| !gapinput@PartialOrderDigraphJoinOfVertices(D, 1, 2);|
  fail
\end{Verbatim}
 }

 

\subsection{\textcolor{Chapter }{NonUpperSemimodularPair}}
\logpage{[ 5, 3, 2 ]}\nobreak
\hyperdef{L}{X824B9896798530F6}{}
{\noindent\textcolor{FuncColor}{$\triangleright$\enspace\texttt{NonUpperSemimodularPair({\mdseries\slshape D})\index{NonUpperSemimodularPair@\texttt{NonUpperSemimodularPair}}
\label{NonUpperSemimodularPair}
}\hfill{\scriptsize (attribute)}}\\
\noindent\textcolor{FuncColor}{$\triangleright$\enspace\texttt{NonLowerSemimodularPair({\mdseries\slshape D})\index{NonLowerSemimodularPair@\texttt{NonLowerSemimodularPair}}
\label{NonLowerSemimodularPair}
}\hfill{\scriptsize (attribute)}}\\
\textbf{\indent Returns:\ }
A pair of vertices or \texttt{fail}.



 \texttt{NonUpperSemimodularPair} returns a pair of vertices in the digraph \mbox{\texttt{\mdseries\slshape D}} that witnesses the fact that \mbox{\texttt{\mdseries\slshape D}} does not represent an upper semimodular lattice, if such a pair exists. 

 If the digraph \mbox{\texttt{\mdseries\slshape D}} does not satisfy \texttt{IsLatticeDigraph} (\ref{IsLatticeDigraph}), then an error is given. Otherwise if the digraph \mbox{\texttt{\mdseries\slshape D}} does satisfy \texttt{IsLatticeDigraph} (\ref{IsLatticeDigraph}), then either a non\texttt{\symbol{45}}upper semimodular pair of vertices is
returned, or \texttt{fail} is returned if no such pair exists (meaning that \mbox{\texttt{\mdseries\slshape D}} is an upper semimodular lattice. 

 \texttt{NonLowerSemimodularPair} behaves in the analogous way to \texttt{NonUpperSemimodularPair} with respect to lower semimodularity. 

 See \texttt{IsUpperSemimodularDigraph} (\ref{IsUpperSemimodularDigraph}) and \texttt{IsLowerSemimodularDigraph} (\ref{IsLowerSemimodularDigraph}) for the definition of upper semimodularity of a lattice. 
\begin{Verbatim}[commandchars=!|E,fontsize=\small,frame=single,label=Example]
  !gapprompt|gap>E !gapinput|D := DigraphFromDigraph6String(E
  !gapprompt|>E !gapinput|"&M~~sc`lYUZO__KIBboC_@h?U_?_GL?A_?c");E
  <immutable digraph with 14 vertices, 66 edges>
  !gapprompt|gap>E !gapinput|NonLowerSemimodularPair(D);E
  [ 10, 9 ]
  !gapprompt|gap>E !gapinput|NonUpperSemimodularPair(D);E
  fail
\end{Verbatim}
 }

 }

 
\section{\textcolor{Chapter }{Reachability and connectivity}}\logpage{[ 5, 4, 0 ]}
\hyperdef{L}{X8537F4088400DC48}{}
{
 

\subsection{\textcolor{Chapter }{DigraphDiameter}}
\logpage{[ 5, 4, 1 ]}\nobreak
\hyperdef{L}{X7F16B9EB8398459C}{}
{\noindent\textcolor{FuncColor}{$\triangleright$\enspace\texttt{DigraphDiameter({\mdseries\slshape digraph})\index{DigraphDiameter@\texttt{DigraphDiameter}}
\label{DigraphDiameter}
}\hfill{\scriptsize (attribute)}}\\
\textbf{\indent Returns:\ }
An integer or \texttt{fail}.



 This function returns the diameter of the digraph \mbox{\texttt{\mdseries\slshape digraph}}. 

 If a digraph \mbox{\texttt{\mdseries\slshape digraph}} is strongly connected and has at least 1 vertex, then the \emph{diameter} is the maximum shortest distance between any pair of distinct vertices.
Otherwise then the diameter of \mbox{\texttt{\mdseries\slshape digraph}} is undefined, and this function returns the value \texttt{fail}. 

 See \texttt{DigraphShortestDistances} (\ref{DigraphShortestDistances}). 

 
\begin{Verbatim}[commandchars=!@|,fontsize=\small,frame=single,label=Example]
  !gapprompt@gap>| !gapinput@D := Digraph([[2], [3], [4, 5], [5], [1, 2, 3, 4, 5]]);|
  <immutable digraph with 5 vertices, 10 edges>
  !gapprompt@gap>| !gapinput@DigraphDiameter(D);|
  3
  !gapprompt@gap>| !gapinput@D := ChainDigraph(2);|
  <immutable chain digraph with 2 vertices>
  !gapprompt@gap>| !gapinput@DigraphDiameter(D);|
  fail
  !gapprompt@gap>| !gapinput@IsStronglyConnectedDigraph(D);|
  false
  !gapprompt@gap>| !gapinput@D := GeneralisedPetersenGraph(IsMutableDigraph, 6, 2);|
  <mutable digraph with 12 vertices, 36 edges>
  !gapprompt@gap>| !gapinput@DigraphDiameter(D);|
  4
\end{Verbatim}
 }

 

\subsection{\textcolor{Chapter }{DigraphShortestDistance (for a digraph and two vertices)}}
\logpage{[ 5, 4, 2 ]}\nobreak
\hyperdef{L}{X8104A9D37BCD8A05}{}
{\noindent\textcolor{FuncColor}{$\triangleright$\enspace\texttt{DigraphShortestDistance({\mdseries\slshape digraph, u, v})\index{DigraphShortestDistance@\texttt{DigraphShortestDistance}!for a digraph and two vertices}
\label{DigraphShortestDistance:for a digraph and two vertices}
}\hfill{\scriptsize (operation)}}\\
\noindent\textcolor{FuncColor}{$\triangleright$\enspace\texttt{DigraphShortestDistance({\mdseries\slshape digraph, list})\index{DigraphShortestDistance@\texttt{DigraphShortestDistance}!for a digraph and a list}
\label{DigraphShortestDistance:for a digraph and a list}
}\hfill{\scriptsize (operation)}}\\
\noindent\textcolor{FuncColor}{$\triangleright$\enspace\texttt{DigraphShortestDistance({\mdseries\slshape digraph, list1, list2})\index{DigraphShortestDistance@\texttt{DigraphShortestDistance}!for a digraph, a list, and a list}
\label{DigraphShortestDistance:for a digraph, a list, and a list}
}\hfill{\scriptsize (operation)}}\\
\textbf{\indent Returns:\ }
An integer or \texttt{fail}



 If there is a directed path in the digraph \mbox{\texttt{\mdseries\slshape digraph}} between vertex \mbox{\texttt{\mdseries\slshape u}} and vertex \mbox{\texttt{\mdseries\slshape v}}, then this operation returns the length of the shortest such directed path.
If no such directed path exists, then this operation returns \texttt{fail}. See Section \ref{Definitions} for the definition of a directed path. 

 If the second form is used, then \mbox{\texttt{\mdseries\slshape list}} should be a list of length two, containing two positive integers which
correspond to the vertices \mbox{\texttt{\mdseries\slshape u}} and \mbox{\texttt{\mdseries\slshape v}}. 

 Note that as usual, a vertex is considered to be at distance 0 from itself .

 If the third form is used, then \mbox{\texttt{\mdseries\slshape list1}} and \mbox{\texttt{\mdseries\slshape list2}} are both lists of vertices. The shortest directed path between \mbox{\texttt{\mdseries\slshape list1}} and \mbox{\texttt{\mdseries\slshape list2}} is then the length of the shortest directed path which starts with a vertex in \mbox{\texttt{\mdseries\slshape list1}} and terminates at a vertex in \mbox{\texttt{\mdseries\slshape list2}}, if such directed path exists. If \mbox{\texttt{\mdseries\slshape list1}} and \mbox{\texttt{\mdseries\slshape list2}} have non\texttt{\symbol{45}}empty intersection, the operation returns \texttt{0}.

 
\begin{Verbatim}[commandchars=!@|,fontsize=\small,frame=single,label=Example]
  !gapprompt@gap>| !gapinput@D := Digraph([[2], [3], [1, 4], [1, 3], [5]]);|
  <immutable digraph with 5 vertices, 7 edges>
  !gapprompt@gap>| !gapinput@DigraphShortestDistance(D, 1, 3);|
  2
  !gapprompt@gap>| !gapinput@DigraphShortestDistance(D, [3, 3]);|
  0
  !gapprompt@gap>| !gapinput@DigraphShortestDistance(D, 5, 2);|
  fail
  !gapprompt@gap>| !gapinput@DigraphShortestDistance(D, [1, 2], [4, 5]);|
  2
  !gapprompt@gap>| !gapinput@DigraphShortestDistance(D, [1, 3], [3, 5]);|
  0
\end{Verbatim}
 }

 

\subsection{\textcolor{Chapter }{DigraphShortestDistances}}
\logpage{[ 5, 4, 3 ]}\nobreak
\hyperdef{L}{X81F99BC67E9D050F}{}
{\noindent\textcolor{FuncColor}{$\triangleright$\enspace\texttt{DigraphShortestDistances({\mdseries\slshape digraph})\index{DigraphShortestDistances@\texttt{DigraphShortestDistances}}
\label{DigraphShortestDistances}
}\hfill{\scriptsize (attribute)}}\\
\textbf{\indent Returns:\ }
A square matrix.



 If \mbox{\texttt{\mdseries\slshape digraph}} is a digraph with $n$ vertices, then this function returns an $n \times n$ matrix \texttt{mat}, where each entry is either a non\texttt{\symbol{45}}negative integer, or \texttt{fail}. If $n = 0$, then an empty list is returned. 

 If there is a directed path from vertex \texttt{i} to vertex \texttt{j}, then the value of \texttt{mat[i][j]} is the length of the shortest such directed path. If no such directed path
exists, then the value of \texttt{mat[i][j]} is \texttt{fail}. We use the convention that the distance from every vertex to itself is \texttt{0}, i.e. \texttt{mat[i][i] = 0} for all vertices \texttt{i}. 

 The method used in this function is a version of the
Floyd\texttt{\symbol{45}}Warshall algorithm, and has complexity $O(n^3)$. 
\begin{Verbatim}[commandchars=!@|,fontsize=\small,frame=single,label=Example]
  !gapprompt@gap>| !gapinput@D := Digraph([[1, 2], [3], [1, 2], [4]]);|
  <immutable digraph with 4 vertices, 6 edges>
  !gapprompt@gap>| !gapinput@mat := DigraphShortestDistances(D);;|
  !gapprompt@gap>| !gapinput@PrintArray(mat);|
  [ [     0,     1,     2,  fail ],
    [     2,     0,     1,  fail ],
    [     1,     1,     0,  fail ],
    [  fail,  fail,  fail,     0 ] ]
  !gapprompt@gap>| !gapinput@D := CycleDigraph(IsMutableDigraph, 3);|
  <mutable digraph with 3 vertices, 3 edges>
  !gapprompt@gap>| !gapinput@DigraphShortestDistances(D);|
  [ [ 0, 1, 2 ], [ 2, 0, 1 ], [ 1, 2, 0 ] ]
\end{Verbatim}
 }

 

\subsection{\textcolor{Chapter }{DigraphLongestDistanceFromVertex}}
\logpage{[ 5, 4, 4 ]}\nobreak
\hyperdef{L}{X8223718079D98A82}{}
{\noindent\textcolor{FuncColor}{$\triangleright$\enspace\texttt{DigraphLongestDistanceFromVertex({\mdseries\slshape digraph, v})\index{DigraphLongestDistanceFromVertex@\texttt{DigraphLongestDistanceFromVertex}}
\label{DigraphLongestDistanceFromVertex}
}\hfill{\scriptsize (operation)}}\\
\textbf{\indent Returns:\ }
An integer, or \texttt{infinity}.



 If \mbox{\texttt{\mdseries\slshape digraph}} is a digraph and \mbox{\texttt{\mdseries\slshape v}} is a vertex in \mbox{\texttt{\mdseries\slshape digraph}}, then this operation returns the length of the longest directed walk in \mbox{\texttt{\mdseries\slshape digraph}} which begins at vertex \mbox{\texttt{\mdseries\slshape v}}. See Section \ref{Definitions} for the definitions of directed walk, directed cycle, and loop. 

 
\begin{itemize}
\item  If there exists a directed walk starting at vertex \mbox{\texttt{\mdseries\slshape v}} which traverses a loop or a directed cycle, then we consider there to be a
walk of infinite length from \mbox{\texttt{\mdseries\slshape v}} (realised by repeatedly traversing the loop/directed cycle), and so the result
is \texttt{infinity}. To disallow walks using loops, try using \texttt{DigraphRemoveLoops} (\ref{DigraphRemoveLoops}):

 \texttt{DigraphLongestDistanceFromVertex(DigraphRemoveLoops(\mbox{\texttt{\mdseries\slshape digraph}},\mbox{\texttt{\mdseries\slshape v}}))}. 
\item  Otherwise, if all directed walks starting at vertex \mbox{\texttt{\mdseries\slshape v}} have finite length, then the length of the longest such walk is returned. 
\end{itemize}
 Note that the result is \texttt{0} if and only if \mbox{\texttt{\mdseries\slshape v}} is a sink of \mbox{\texttt{\mdseries\slshape digraph}}. See \texttt{DigraphSinks} (\ref{DigraphSinks}). 
\begin{Verbatim}[commandchars=!@|,fontsize=\small,frame=single,label=Example]
  !gapprompt@gap>| !gapinput@D := Digraph([[2], [3, 4], [], [5], [], [6]]);|
  <immutable digraph with 6 vertices, 5 edges>
  !gapprompt@gap>| !gapinput@DigraphLongestDistanceFromVertex(D, 1);|
  3
  !gapprompt@gap>| !gapinput@DigraphLongestDistanceFromVertex(D, 3);|
  0
  !gapprompt@gap>| !gapinput@3 in DigraphSinks(D);|
  true
  !gapprompt@gap>| !gapinput@DigraphLongestDistanceFromVertex(D, 6);|
  infinity
\end{Verbatim}
 }

 

\subsection{\textcolor{Chapter }{DigraphDistanceSet (for a digraph, a pos int, and an int)}}
\logpage{[ 5, 4, 5 ]}\nobreak
\hyperdef{L}{X7CB7DDCD84621D38}{}
{\noindent\textcolor{FuncColor}{$\triangleright$\enspace\texttt{DigraphDistanceSet({\mdseries\slshape digraph, vertex, distance})\index{DigraphDistanceSet@\texttt{DigraphDistanceSet}!for a digraph, a pos int, and an int}
\label{DigraphDistanceSet:for a digraph, a pos int, and an int}
}\hfill{\scriptsize (operation)}}\\
\noindent\textcolor{FuncColor}{$\triangleright$\enspace\texttt{DigraphDistanceSet({\mdseries\slshape digraph, vertex, distances})\index{DigraphDistanceSet@\texttt{DigraphDistanceSet}!for a digraph, a pos int, and a list}
\label{DigraphDistanceSet:for a digraph, a pos int, and a list}
}\hfill{\scriptsize (operation)}}\\
\textbf{\indent Returns:\ }
A list of vertices



 This operation returns the list of all vertices in digraph \mbox{\texttt{\mdseries\slshape digraph}} such that the shortest distance to a vertex \mbox{\texttt{\mdseries\slshape vertex}} is \mbox{\texttt{\mdseries\slshape distance}} or is in the list \mbox{\texttt{\mdseries\slshape distances}}. 

 \mbox{\texttt{\mdseries\slshape digraph}} should be a digraph, \mbox{\texttt{\mdseries\slshape vertex}} should be a positive integer, \mbox{\texttt{\mdseries\slshape distance}} should be a non\texttt{\symbol{45}}negative integer, and \mbox{\texttt{\mdseries\slshape distances}} should be a list of non\texttt{\symbol{45}}negative integers. 
\begin{Verbatim}[commandchars=!@|,fontsize=\small,frame=single,label=Example]
  !gapprompt@gap>| !gapinput@D := Digraph([[2], [3], [1, 4], [1, 3]]);|
  <immutable digraph with 4 vertices, 6 edges>
  !gapprompt@gap>| !gapinput@DigraphDistanceSet(D, 2, [1, 2]);|
  [ 3, 1, 4 ]
  !gapprompt@gap>| !gapinput@DigraphDistanceSet(D, 3, 1);|
  [ 1, 4 ]
  !gapprompt@gap>| !gapinput@DigraphDistanceSet(D, 2, 0);|
  [ 2 ]
\end{Verbatim}
 }

 

\subsection{\textcolor{Chapter }{DigraphGirth}}
\logpage{[ 5, 4, 6 ]}\nobreak
\hyperdef{L}{X79A3DA4078CF3C90}{}
{\noindent\textcolor{FuncColor}{$\triangleright$\enspace\texttt{DigraphGirth({\mdseries\slshape digraph})\index{DigraphGirth@\texttt{DigraphGirth}}
\label{DigraphGirth}
}\hfill{\scriptsize (attribute)}}\\
\textbf{\indent Returns:\ }
An integer, or \texttt{infinity}.



 This attribute returns the \emph{girth} of the digraph \mbox{\texttt{\mdseries\slshape digraph}}. The \emph{girth} of a digraph is the length of its shortest simple circuit. See Section \ref{Definitions} for the definitions of simple circuit, directed cycle, and loop. 

 If \mbox{\texttt{\mdseries\slshape digraph}} has no directed cycles, then this function will return \texttt{infinity}. If \mbox{\texttt{\mdseries\slshape digraph}} contains a loop, then this function will return \texttt{1}. 

 In the worst case, the method used in this function is a version of the
Floyd\texttt{\symbol{45}}Warshall algorithm, and has complexity \texttt{O(\mbox{\texttt{\mdseries\slshape n}} \texttt{\symbol{94}} 3)}, where \mbox{\texttt{\mdseries\slshape n}} is the number of vertices in \mbox{\texttt{\mdseries\slshape digraph}}. If the digraph has known automorphisms [see \texttt{DigraphGroup} (\ref{DigraphGroup})], then the performance is likely to be better. 

 For symmetric digraphs, see also \texttt{DigraphUndirectedGirth} (\ref{DigraphUndirectedGirth}). 

 
\begin{Verbatim}[commandchars=!@|,fontsize=\small,frame=single,label=Example]
  !gapprompt@gap>| !gapinput@D := Digraph([[1], [1]]);|
  <immutable digraph with 2 vertices, 2 edges>
  !gapprompt@gap>| !gapinput@DigraphGirth(D);|
  1
  !gapprompt@gap>| !gapinput@D := Digraph([[2, 3], [3], [4], []]);|
  <immutable digraph with 4 vertices, 4 edges>
  !gapprompt@gap>| !gapinput@DigraphGirth(D);|
  infinity
  !gapprompt@gap>| !gapinput@D := Digraph([[2, 3], [3], [4], [1]]);|
  <immutable digraph with 4 vertices, 5 edges>
  !gapprompt@gap>| !gapinput@DigraphGirth(D);|
  3
  !gapprompt@gap>| !gapinput@D := GeneralisedPetersenGraph(IsMutableDigraph, 6, 2);|
  <mutable digraph with 12 vertices, 36 edges>
  !gapprompt@gap>| !gapinput@DigraphGirth(D);|
  2
\end{Verbatim}
 }

 

\subsection{\textcolor{Chapter }{DigraphOddGirth}}
\logpage{[ 5, 4, 7 ]}\nobreak
\hyperdef{L}{X8374B7357EC189C1}{}
{\noindent\textcolor{FuncColor}{$\triangleright$\enspace\texttt{DigraphOddGirth({\mdseries\slshape digraph})\index{DigraphOddGirth@\texttt{DigraphOddGirth}}
\label{DigraphOddGirth}
}\hfill{\scriptsize (attribute)}}\\
\textbf{\indent Returns:\ }
An integer, or \texttt{infinity}.



 This attribute returns the \emph{odd girth} of the digraph \mbox{\texttt{\mdseries\slshape digraph}}. The \emph{odd girth} of a digraph is the length of its shortest simple circuit of odd length. See
Section \ref{Definitions} for the definitions of simple circuit, directed cycle, and loop. 

 If \mbox{\texttt{\mdseries\slshape digraph}} has no directed cycles of odd length, then this function will return \texttt{infinity}, even if \mbox{\texttt{\mdseries\slshape digraph}} has a directed cycle of even length. If \mbox{\texttt{\mdseries\slshape digraph}} contains a loop, then this function will return \texttt{1}. 

 See also \texttt{DigraphGirth} (\ref{DigraphGirth}). 

 
\begin{Verbatim}[commandchars=!@|,fontsize=\small,frame=single,label=Example]
  !gapprompt@gap>| !gapinput@D := Digraph([[2], [3, 1], [1]]);|
  <immutable digraph with 3 vertices, 4 edges>
  !gapprompt@gap>| !gapinput@DigraphOddGirth(D);|
  3
  !gapprompt@gap>| !gapinput@D := CycleDigraph(4);|
  <immutable cycle digraph with 4 vertices>
  !gapprompt@gap>| !gapinput@DigraphOddGirth(D);|
  infinity
  !gapprompt@gap>| !gapinput@D := Digraph([[2], [3], [], [3], [4]]);|
  <immutable digraph with 5 vertices, 4 edges>
  !gapprompt@gap>| !gapinput@DigraphOddGirth(D);|
  infinity
  !gapprompt@gap>| !gapinput@D := CycleDigraph(IsMutableDigraph, 4);|
  <mutable digraph with 4 vertices, 4 edges>
  !gapprompt@gap>| !gapinput@DigraphDisjointUnion(D, CycleDigraph(5));|
  <mutable digraph with 9 vertices, 9 edges>
  !gapprompt@gap>| !gapinput@DigraphOddGirth(D);|
  5
\end{Verbatim}
 }

 

\subsection{\textcolor{Chapter }{DigraphUndirectedGirth}}
\logpage{[ 5, 4, 8 ]}\nobreak
\hyperdef{L}{X84688B337BDDBB09}{}
{\noindent\textcolor{FuncColor}{$\triangleright$\enspace\texttt{DigraphUndirectedGirth({\mdseries\slshape digraph})\index{DigraphUndirectedGirth@\texttt{DigraphUndirectedGirth}}
\label{DigraphUndirectedGirth}
}\hfill{\scriptsize (attribute)}}\\
\textbf{\indent Returns:\ }
An integer or \texttt{infinity}.



 If \mbox{\texttt{\mdseries\slshape digraph}} is a symmetric digraph, then this function returns the girth of \mbox{\texttt{\mdseries\slshape digraph}} when treated as an undirected graph (i.e. each pair of edges $[i, j]$ and $[j, i]$ is treated as a single edge between $i$ and $j$). 

 The \emph{girth} of an undirected graph is the length of its shortest simple cycle, i.e. the
shortest non\texttt{\symbol{45}}trivial path starting and ending at the same
vertex and passing through no vertex or edge more than once. 

 If \mbox{\texttt{\mdseries\slshape digraph}} has no cycles, then this function will return \texttt{infinity}. 

 
\begin{Verbatim}[commandchars=!@|,fontsize=\small,frame=single,label=Example]
  !gapprompt@gap>| !gapinput@D := Digraph([[2, 4], [1, 3], [2, 4], [1, 3]]);|
  <immutable digraph with 4 vertices, 8 edges>
  !gapprompt@gap>| !gapinput@DigraphUndirectedGirth(D);|
  4
  !gapprompt@gap>| !gapinput@D := Digraph([[2], [1, 3], [2]]);|
  <immutable digraph with 3 vertices, 4 edges>
  !gapprompt@gap>| !gapinput@DigraphUndirectedGirth(D);|
  infinity
  !gapprompt@gap>| !gapinput@D := Digraph([[1], [], [4], [3]]);|
  <immutable digraph with 4 vertices, 3 edges>
  !gapprompt@gap>| !gapinput@DigraphUndirectedGirth(D);|
  1
  !gapprompt@gap>| !gapinput@D := GeneralisedPetersenGraph(IsMutableDigraph, 9, 2);|
  <mutable digraph with 18 vertices, 54 edges>
  !gapprompt@gap>| !gapinput@DigraphUndirectedGirth(D);|
  5
\end{Verbatim}
 }

 

\subsection{\textcolor{Chapter }{DigraphConnectedComponents}}
\logpage{[ 5, 4, 9 ]}\nobreak
\hyperdef{L}{X842FAD6A7B835977}{}
{\noindent\textcolor{FuncColor}{$\triangleright$\enspace\texttt{DigraphConnectedComponents({\mdseries\slshape digraph})\index{DigraphConnectedComponents@\texttt{DigraphConnectedComponents}}
\label{DigraphConnectedComponents}
}\hfill{\scriptsize (attribute)}}\\
\noindent\textcolor{FuncColor}{$\triangleright$\enspace\texttt{DigraphNrConnectedComponents({\mdseries\slshape digraph})\index{DigraphNrConnectedComponents@\texttt{DigraphNrConnectedComponents}}
\label{DigraphNrConnectedComponents}
}\hfill{\scriptsize (attribute)}}\\
\textbf{\indent Returns:\ }
A record.



 This function returns the record \texttt{wcc} corresponding to the weakly connected components of the digraph \mbox{\texttt{\mdseries\slshape digraph}}. Two vertices of \mbox{\texttt{\mdseries\slshape digraph}} are in the same weakly connected component whenever they are equal, or there
exists a directed path (ignoring the orientation of edges) between them. More
formally, two vertices are in the same weakly connected component of \mbox{\texttt{\mdseries\slshape digraph}} if and only if they are in the same strongly connected component (see \texttt{DigraphStronglyConnectedComponents} (\ref{DigraphStronglyConnectedComponents})) of the \texttt{DigraphSymmetricClosure} (\ref{DigraphSymmetricClosure}) of \mbox{\texttt{\mdseries\slshape digraph}}. 

 The set of weakly connected components is a partition of the vertex set of \mbox{\texttt{\mdseries\slshape digraph}}. 

 The record \texttt{wcc} has 2 components: \texttt{comps} and \texttt{id}. The component \texttt{comps} is a list of the weakly connected components of \mbox{\texttt{\mdseries\slshape digraph}} (each of which is a list of vertices). The component \texttt{id} is a list such that the vertex \texttt{i} is an element of the weakly connected component \texttt{comps[id[i]]}. 

 The method used in this function has complexity $O(m+n)$, where $m$ is the number of edges and $n$ is the number of vertices in the digraph. 

 \texttt{DigraphNrConnectedComponents(\mbox{\texttt{\mdseries\slshape digraph}})} is simply a shortcut for \texttt{Length(DigraphConnectedComponents(\mbox{\texttt{\mdseries\slshape digraph}}).comps)}, and is no more efficient. 
\begin{Verbatim}[commandchars=!@|,fontsize=\small,frame=single,label=Example]
  !gapprompt@gap>| !gapinput@gr := Digraph([[2], [3, 1], []]);|
  <immutable digraph with 3 vertices, 3 edges>
  !gapprompt@gap>| !gapinput@DigraphConnectedComponents(gr);|
  rec( comps := [ [ 1, 2, 3 ] ], id := [ 1, 1, 1 ] )
  !gapprompt@gap>| !gapinput@gr := Digraph([[1], [1, 2], []]);|
  <immutable digraph with 3 vertices, 3 edges>
  !gapprompt@gap>| !gapinput@DigraphConnectedComponents(gr);|
  rec( comps := [ [ 1, 2 ], [ 3 ] ], id := [ 1, 1, 2 ] )
  !gapprompt@gap>| !gapinput@DigraphNrConnectedComponents(gr);|
  2
  !gapprompt@gap>| !gapinput@gr := EmptyDigraph(0);|
  <immutable empty digraph with 0 vertices>
  !gapprompt@gap>| !gapinput@DigraphConnectedComponents(gr);|
  rec( comps := [  ], id := [  ] )
  !gapprompt@gap>| !gapinput@D := CycleDigraph(IsMutableDigraph, 2);|
  <mutable digraph with 2 vertices, 2 edges>
  !gapprompt@gap>| !gapinput@G := CycleDigraph(3);|
  <immutable cycle digraph with 3 vertices>
  !gapprompt@gap>| !gapinput@DigraphDisjointUnion(D, G);|
  <mutable digraph with 5 vertices, 5 edges>
  !gapprompt@gap>| !gapinput@DigraphConnectedComponents(D);|
  rec( comps := [ [ 1, 2 ], [ 3, 4, 5 ] ], id := [ 1, 1, 2, 2, 2 ] )
\end{Verbatim}
 }

 

\subsection{\textcolor{Chapter }{DigraphConnectedComponent}}
\logpage{[ 5, 4, 10 ]}\nobreak
\hyperdef{L}{X8484EC557810CD31}{}
{\noindent\textcolor{FuncColor}{$\triangleright$\enspace\texttt{DigraphConnectedComponent({\mdseries\slshape digraph, vertex})\index{DigraphConnectedComponent@\texttt{DigraphConnectedComponent}}
\label{DigraphConnectedComponent}
}\hfill{\scriptsize (operation)}}\\
\textbf{\indent Returns:\ }
A list of vertices.



 If \mbox{\texttt{\mdseries\slshape vertex}} is a vertex in the digraph \mbox{\texttt{\mdseries\slshape digraph}}, then this operation returns the connected component of \mbox{\texttt{\mdseries\slshape vertex}} in \mbox{\texttt{\mdseries\slshape digraph}}. See \texttt{DigraphConnectedComponents} (\ref{DigraphConnectedComponents}) for more information. 
\begin{Verbatim}[commandchars=!@|,fontsize=\small,frame=single,label=Example]
  !gapprompt@gap>| !gapinput@D := Digraph([[3], [2], [1, 2], [4]]);|
  <immutable digraph with 4 vertices, 5 edges>
  !gapprompt@gap>| !gapinput@DigraphConnectedComponent(D, 3);|
  [ 1, 2, 3 ]
  !gapprompt@gap>| !gapinput@DigraphConnectedComponent(D, 2);|
  [ 1, 2, 3 ]
  !gapprompt@gap>| !gapinput@DigraphConnectedComponent(D, 4);|
  [ 4 ]
\end{Verbatim}
 }

 

\subsection{\textcolor{Chapter }{DigraphStronglyConnectedComponents}}
\logpage{[ 5, 4, 11 ]}\nobreak
\hyperdef{L}{X833ECD6B7A84944C}{}
{\noindent\textcolor{FuncColor}{$\triangleright$\enspace\texttt{DigraphStronglyConnectedComponents({\mdseries\slshape digraph})\index{DigraphStronglyConnectedComponents@\texttt{DigraphStronglyConnectedComponents}}
\label{DigraphStronglyConnectedComponents}
}\hfill{\scriptsize (attribute)}}\\
\noindent\textcolor{FuncColor}{$\triangleright$\enspace\texttt{DigraphNrStronglyConnectedComponents({\mdseries\slshape digraph})\index{DigraphNrStronglyConnectedComponents@\texttt{Digraph}\-\texttt{Nr}\-\texttt{Strongly}\-\texttt{Connected}\-\texttt{Components}}
\label{DigraphNrStronglyConnectedComponents}
}\hfill{\scriptsize (attribute)}}\\
\textbf{\indent Returns:\ }
A record.



 This function returns the record \texttt{scc} corresponding to the strongly connected components of the digraph \mbox{\texttt{\mdseries\slshape digraph}}. Two vertices of \mbox{\texttt{\mdseries\slshape digraph}} are in the same strongly connected component whenever they are equal, or there
is a directed path from each vertex to the other. The set of strongly
connected components is a partition of the vertex set of \mbox{\texttt{\mdseries\slshape digraph}}. 

 The record \texttt{scc} has 2 components: \texttt{comps} and \texttt{id}. The component \texttt{comps} is a list of the strongly connected components of \mbox{\texttt{\mdseries\slshape digraph}} (each of which is a list of vertices). The component \texttt{id} is a list such that the vertex \texttt{i} is an element of the strongly connected component \texttt{comps[id[i]]}. 

 The method used in this function is a non\texttt{\symbol{45}}recursive version
of Gabow's Algorithm \cite{Gab00} and has complexity $O(m+n)$ where $m$ is the number of edges (counting multiple edges as one) and $n$ is the number of vertices in the digraph. 

 \texttt{DigraphNrStronglyConnectedComponents(\mbox{\texttt{\mdseries\slshape digraph}})} is simply a shortcut for \texttt{Length(DigraphStronglyConnectedComponents(\mbox{\texttt{\mdseries\slshape digraph}}).comps)}, and is no more efficient. 
\begin{Verbatim}[commandchars=!@|,fontsize=\small,frame=single,label=Example]
  !gapprompt@gap>| !gapinput@gr := Digraph([[2], [3, 1], []]);|
  <immutable digraph with 3 vertices, 3 edges>
  !gapprompt@gap>| !gapinput@DigraphStronglyConnectedComponents(gr);|
  rec( comps := [ [ 3 ], [ 1, 2 ] ], id := [ 2, 2, 1 ] )
  !gapprompt@gap>| !gapinput@DigraphNrStronglyConnectedComponents(gr);|
  2
  !gapprompt@gap>| !gapinput@D := DigraphDisjointUnion(CycleDigraph(4), CycleDigraph(5));|
  <immutable digraph with 9 vertices, 9 edges>
  !gapprompt@gap>| !gapinput@DigraphStronglyConnectedComponents(D);|
  rec( comps := [ [ 1, 2, 3, 4 ], [ 5, 6, 7, 8, 9 ] ], 
    id := [ 1, 1, 1, 1, 2, 2, 2, 2, 2 ] )
  !gapprompt@gap>| !gapinput@DigraphNrStronglyConnectedComponents(D);|
  2
  !gapprompt@gap>| !gapinput@D := CycleDigraph(IsMutableDigraph, 2);|
  <mutable digraph with 2 vertices, 2 edges>
  !gapprompt@gap>| !gapinput@G := CycleDigraph(3);|
  <immutable cycle digraph with 3 vertices>
  !gapprompt@gap>| !gapinput@DigraphDisjointUnion(D, G);|
  <mutable digraph with 5 vertices, 5 edges>
  !gapprompt@gap>| !gapinput@DigraphStronglyConnectedComponents(D);|
  rec( comps := [ [ 1, 2 ], [ 3, 4, 5 ] ], id := [ 1, 1, 2, 2, 2 ] )
\end{Verbatim}
 }

 

\subsection{\textcolor{Chapter }{DigraphStronglyConnectedComponent}}
\logpage{[ 5, 4, 12 ]}\nobreak
\hyperdef{L}{X7EFCB5017D662254}{}
{\noindent\textcolor{FuncColor}{$\triangleright$\enspace\texttt{DigraphStronglyConnectedComponent({\mdseries\slshape digraph, vertex})\index{DigraphStronglyConnectedComponent@\texttt{DigraphStronglyConnectedComponent}}
\label{DigraphStronglyConnectedComponent}
}\hfill{\scriptsize (operation)}}\\
\textbf{\indent Returns:\ }
A list of vertices.



 If \mbox{\texttt{\mdseries\slshape vertex}} is a vertex in the digraph \mbox{\texttt{\mdseries\slshape digraph}}, then this operation returns the strongly connected component of \mbox{\texttt{\mdseries\slshape vertex}} in \mbox{\texttt{\mdseries\slshape digraph}}. See \texttt{DigraphStronglyConnectedComponents} (\ref{DigraphStronglyConnectedComponents}) for more information. 
\begin{Verbatim}[commandchars=!@|,fontsize=\small,frame=single,label=Example]
  !gapprompt@gap>| !gapinput@D := Digraph([[3], [2], [1, 2], [3]]);|
  <immutable digraph with 4 vertices, 5 edges>
  !gapprompt@gap>| !gapinput@DigraphStronglyConnectedComponent(D, 3);|
  [ 1, 3 ]
  !gapprompt@gap>| !gapinput@DigraphStronglyConnectedComponent(D, 2);|
  [ 2 ]
  !gapprompt@gap>| !gapinput@DigraphStronglyConnectedComponent(D, 4);|
  [ 4 ]
\end{Verbatim}
 }

 

\subsection{\textcolor{Chapter }{DigraphBicomponents}}
\logpage{[ 5, 4, 13 ]}\nobreak
\hyperdef{L}{X7F1B5A2782F598B1}{}
{\noindent\textcolor{FuncColor}{$\triangleright$\enspace\texttt{DigraphBicomponents({\mdseries\slshape digraph})\index{DigraphBicomponents@\texttt{DigraphBicomponents}}
\label{DigraphBicomponents}
}\hfill{\scriptsize (attribute)}}\\
\textbf{\indent Returns:\ }
A pair of lists of vertices, or \texttt{fail}.



 If \mbox{\texttt{\mdseries\slshape digraph}} is a bipartite digraph, i.e. if it satisfies \texttt{IsBipartiteDigraph} (\ref{IsBipartiteDigraph}), then \texttt{DigraphBicomponents} returns a pair of bicomponents of \mbox{\texttt{\mdseries\slshape digraph}}. Otherwise, \texttt{DigraphBicomponents} returns \texttt{fail}.

 For a bipartite digraph, the vertices can be partitioned into two
non\texttt{\symbol{45}}empty sets such that the source and range of any edge
are in distinct sets. The parts of this partition are called \emph{bicomponents} of \mbox{\texttt{\mdseries\slshape digraph}}. Equivalently, a pair of bicomponents of \mbox{\texttt{\mdseries\slshape digraph}} consists of the color\texttt{\symbol{45}}classes of a
2\texttt{\symbol{45}}coloring of \mbox{\texttt{\mdseries\slshape digraph}}. 

 For a bipartite digraph with at least 3 vertices, there is a unique pair of
bicomponents of bipartite if and only if the digraph is connected. See \texttt{IsConnectedDigraph} (\ref{IsConnectedDigraph}) for more information. 
\begin{Verbatim}[commandchars=!@|,fontsize=\small,frame=single,label=Example]
  !gapprompt@gap>| !gapinput@D := CycleDigraph(3);|
  <immutable cycle digraph with 3 vertices>
  !gapprompt@gap>| !gapinput@DigraphBicomponents(D);|
  fail
  !gapprompt@gap>| !gapinput@D := ChainDigraph(5);|
  <immutable chain digraph with 5 vertices>
  !gapprompt@gap>| !gapinput@DigraphBicomponents(D);|
  [ [ 1, 3, 5 ], [ 2, 4 ] ]
  !gapprompt@gap>| !gapinput@D := Digraph([[5], [1, 4], [5], [5], []]);|
  <immutable digraph with 5 vertices, 5 edges>
  !gapprompt@gap>| !gapinput@DigraphBicomponents(D);|
  [ [ 1, 3, 4 ], [ 2, 5 ] ]
  !gapprompt@gap>| !gapinput@D := CompleteBipartiteDigraph(IsMutableDigraph, 2, 3);|
  <mutable digraph with 5 vertices, 12 edges>
  !gapprompt@gap>| !gapinput@DigraphBicomponents(D);|
  [ [ 1, 2 ], [ 3, 4, 5 ] ]
\end{Verbatim}
 }

 

\subsection{\textcolor{Chapter }{ArticulationPoints}}
\logpage{[ 5, 4, 14 ]}\nobreak
\hyperdef{L}{X7DDE06E47E605DD7}{}
{\noindent\textcolor{FuncColor}{$\triangleright$\enspace\texttt{ArticulationPoints({\mdseries\slshape digraph})\index{ArticulationPoints@\texttt{ArticulationPoints}}
\label{ArticulationPoints}
}\hfill{\scriptsize (attribute)}}\\
\textbf{\indent Returns:\ }
A list of vertices.



 A connected digraph is \emph{biconnected} if it is still connected (in the sense of \texttt{IsConnectedDigraph} (\ref{IsConnectedDigraph})) when any vertex is removed. If the digraph \mbox{\texttt{\mdseries\slshape digraph}} is not biconnected but is connected, then any vertex \texttt{v} of \mbox{\texttt{\mdseries\slshape digraph}} whose removal makes the resulting digraph disconnected is called an \emph{articulation point}.

 \texttt{ArticulationPoints} returns a list of the articulation points of \mbox{\texttt{\mdseries\slshape digraph}}, if any, and, in particular, returns the empty list if \mbox{\texttt{\mdseries\slshape digraph}} is not connected. 

 Multiple edges are ignored by this method. 

 The method used in this operation has complexity $O(m+n)$ where $m$ is the number of edges and $n$ is the number of vertices in the digraph.

 If \mbox{\texttt{\mdseries\slshape D}} has a bridge (see \texttt{Bridges} (\ref{Bridges})), then a node incident to the bridge is an articulation point if and only if
it has degree at least $2$. It follows that if \mbox{\texttt{\mdseries\slshape D}} has a bridge and at least $3$ nodes, then at least one of the nodes incident to the bridge is an
articulation point. The converse does not hold, there are digraphs with
articulation points, but no bridges. 

 See also \texttt{IsBiconnectedDigraph} (\ref{IsBiconnectedDigraph}) and \texttt{IsBridgelessDigraph} (\ref{IsBridgelessDigraph}). 
\begin{Verbatim}[commandchars=!@|,fontsize=\small,frame=single,label=Example]
  !gapprompt@gap>| !gapinput@ArticulationPoints(CycleDigraph(5));|
  [  ]
  !gapprompt@gap>| !gapinput@D := Digraph([[2, 7], [3, 5], [4], [2], [6], [1], []]);;|
  !gapprompt@gap>| !gapinput@ArticulationPoints(D);|
  [ 1, 2 ]
  !gapprompt@gap>| !gapinput@ArticulationPoints(ChainDigraph(5));|
  [ 2, 3, 4 ]
  !gapprompt@gap>| !gapinput@ArticulationPoints(NullDigraph(5));|
  [  ]
  !gapprompt@gap>| !gapinput@D := ChainDigraph(IsMutableDigraph, 4);|
  <mutable digraph with 4 vertices, 3 edges>
  !gapprompt@gap>| !gapinput@ArticulationPoints(D);|
  [ 2, 3 ]
\end{Verbatim}
 }

 

\subsection{\textcolor{Chapter }{MinimalCyclicEdgeCut}}
\logpage{[ 5, 4, 15 ]}\nobreak
\hyperdef{L}{X803459AB86AB9BE2}{}
{\noindent\textcolor{FuncColor}{$\triangleright$\enspace\texttt{MinimalCyclicEdgeCut({\mdseries\slshape digraph})\index{MinimalCyclicEdgeCut@\texttt{MinimalCyclicEdgeCut}}
\label{MinimalCyclicEdgeCut}
}\hfill{\scriptsize (attribute)}}\\
\textbf{\indent Returns:\ }
A list of edges or \texttt{fail}.



 A cyclic edge cut of \mbox{\texttt{\mdseries\slshape digraph}} is a set of edges such that deleting these edges results in at least two
connected components having a cycle. This method computes the cyclic edge cut
with minimal cardinality for cubic graphs with at least 8 vertices. The cyclic
edge cut is returned as a list of undirected edges.

 If the given digraph is not cubic, not connected, has less than 8 vertices or
does not have a cyclic edge cut, the method returns \texttt{fail}. Multiple edges of \mbox{\texttt{\mdseries\slshape digraph}} are ignored by this method and note that \mbox{\texttt{\mdseries\slshape digraph}} is identified as an undirected graph. 

 The method used in this method has complexity $O(n^2*log^2(n)) $ where $n$ is the number of vertices in the digraph.

 
\begin{Verbatim}[commandchars=!@|,fontsize=\small,frame=single,label=Example]
  !gapprompt@gap>| !gapinput@MinimalCyclicEdgeCut(HypercubeGraph(3));|
  [ [ 1, 5 ], [ 2, 6 ], [ 4, 8 ], [ 3, 7 ] ]
  !gapprompt@gap>| !gapinput@MinimalCyclicEdgeCut(CompleteDigraph(4));|
  fail
\end{Verbatim}
 }

 

\subsection{\textcolor{Chapter }{Bridges}}
\logpage{[ 5, 4, 16 ]}\nobreak
\hyperdef{L}{X84D5A125848BD800}{}
{\noindent\textcolor{FuncColor}{$\triangleright$\enspace\texttt{Bridges({\mdseries\slshape D})\index{Bridges@\texttt{Bridges}}
\label{Bridges}
}\hfill{\scriptsize (attribute)}}\\
\textbf{\indent Returns:\ }
A (possibly empty) list of edges.



 A connected digraph is \emph{2\texttt{\symbol{45}}edge\texttt{\symbol{45}}connected} if it is still connected (in the sense of \texttt{IsConnectedDigraph} (\ref{IsConnectedDigraph})) when any edge is removed. If the digraph \mbox{\texttt{\mdseries\slshape D}} is not 2\texttt{\symbol{45}}edge\texttt{\symbol{45}}connected but is
connected, then any edge \texttt{[u, v]} of \mbox{\texttt{\mdseries\slshape D}} whose removal makes the resulting digraph disconnected is called a \emph{bridge}.

 \texttt{Bridges} returns a list of the bridges of \mbox{\texttt{\mdseries\slshape D}}, if any, and, in particular, returns the empty list if \mbox{\texttt{\mdseries\slshape D}} is not connected. 

 Multiple edges are ignored by this method. 

 The method used in this operation has complexity $O(m+n)$ where $m$ is the number of edges and $n$ is the number of vertices in the digraph. 

 If \mbox{\texttt{\mdseries\slshape D}} has a bridge, then a node incident to the bridge is an articulation point (see \texttt{ArticulationPoints} (\ref{ArticulationPoints})) if and only if it has degree at least $2$. It follows that if \mbox{\texttt{\mdseries\slshape D}} has a bridge and at least $3$ nodes, then at least one of the nodes incident to the bridge is an
articulation point. The converse does not hold, there are digraphs with
articulation points, but no bridges. 

 See also \texttt{IsBiconnectedDigraph} (\ref{IsBiconnectedDigraph}) and \texttt{IsBridgelessDigraph} (\ref{IsBridgelessDigraph}). 
\begin{Verbatim}[commandchars=!@|,fontsize=\small,frame=single,label=Example]
  !gapprompt@gap>| !gapinput@D := Digraph([[2, 5], [1, 3, 4, 5], [2, 4], [2, 3], [1, 2]]);|
  <immutable digraph with 5 vertices, 12 edges>
  !gapprompt@gap>| !gapinput@Bridges(D);|
  [  ]
  !gapprompt@gap>| !gapinput@D := Digraph([[2], [3], [4], [2]]);|
  <immutable digraph with 4 vertices, 4 edges>
  !gapprompt@gap>| !gapinput@Bridges(D);|
  [ [ 1, 2 ] ]
  !gapprompt@gap>| !gapinput@Bridges(ChainDigraph(2));|
  [ [ 1, 2 ] ]
  !gapprompt@gap>| !gapinput@ArticulationPoints(ChainDigraph(2));|
  [  ]
\end{Verbatim}
 }

 

\subsection{\textcolor{Chapter }{StrongOrientation}}
\logpage{[ 5, 4, 17 ]}\nobreak
\hyperdef{L}{X865590147BD1C507}{}
{\noindent\textcolor{FuncColor}{$\triangleright$\enspace\texttt{StrongOrientation({\mdseries\slshape D})\index{StrongOrientation@\texttt{StrongOrientation}}
\label{StrongOrientation}
}\hfill{\scriptsize (operation)}}\\
\noindent\textcolor{FuncColor}{$\triangleright$\enspace\texttt{StrongOrientationAttr({\mdseries\slshape D})\index{StrongOrientationAttr@\texttt{StrongOrientationAttr}}
\label{StrongOrientationAttr}
}\hfill{\scriptsize (attribute)}}\\
\textbf{\indent Returns:\ }
A digraph or \texttt{fail}.



 A \emph{strong orientation} of a connected symmetric digraph \mbox{\texttt{\mdseries\slshape D}} (if it exists) is a strongly connected subdigraph \texttt{C} of \mbox{\texttt{\mdseries\slshape D}} such that for every edge \texttt{[u, v]} of \mbox{\texttt{\mdseries\slshape D}} either \texttt{[u, v]} or \texttt{[v, u]} is an edge of \texttt{C} but not both. Robbin's Theorem states that a digraph admits a strong
orientation if and only if it is bridgeless (see \texttt{IsBridgelessDigraph} (\ref{IsBridgelessDigraph})). 

 This operation returns a strong orientation of the digraph \mbox{\texttt{\mdseries\slshape D}} if \mbox{\texttt{\mdseries\slshape D}} is symmetric and \mbox{\texttt{\mdseries\slshape D}} admits a strong orientation. If \mbox{\texttt{\mdseries\slshape D}} is symmetric but does not admit a strong orientation, then \texttt{fail} is returned. If \mbox{\texttt{\mdseries\slshape D}} is not symmetric, then an error is given. 

 If \mbox{\texttt{\mdseries\slshape D}} is immutable, \texttt{StrongOrientation(\mbox{\texttt{\mdseries\slshape D}})} returns an immutable digraph, and if \mbox{\texttt{\mdseries\slshape D}} is mutable, then \texttt{StrongOrientation(\mbox{\texttt{\mdseries\slshape D}})} returns a mutable digraph. 

 The method used in this operation has complexity $O(m+n)$ where $m$ is the number of edges and $n$ is the number of vertices in the digraph. 
\begin{Verbatim}[commandchars=!@|,fontsize=\small,frame=single,label=Example]
  !gapprompt@gap>| !gapinput@StrongOrientation(DigraphSymmetricClosure(CycleDigraph(5))) |
  !gapprompt@>| !gapinput@= CycleDigraph(5);|
  true
  !gapprompt@gap>| !gapinput@D := DigraphSymmetricClosure(Digraph(|
  !gapprompt@>| !gapinput@[[2, 7], [3, 5], [4], [2], [6], [1], []]));;|
  !gapprompt@gap>| !gapinput@IsBridgelessDigraph(D);|
  false
  !gapprompt@gap>| !gapinput@StrongOrientation(D);|
  fail
  !gapprompt@gap>| !gapinput@StrongOrientation(NullDigraph(0));|
  <immutable empty digraph with 0 vertices>
  !gapprompt@gap>| !gapinput@StrongOrientation(DigraphDisjointUnion(CompleteDigraph(3), |
  !gapprompt@>| !gapinput@                                          CompleteDigraph(3)));|
  fail
\end{Verbatim}
 }

 

\subsection{\textcolor{Chapter }{DigraphPeriod}}
\logpage{[ 5, 4, 18 ]}\nobreak
\hyperdef{L}{X853D0B0981A33433}{}
{\noindent\textcolor{FuncColor}{$\triangleright$\enspace\texttt{DigraphPeriod({\mdseries\slshape digraph})\index{DigraphPeriod@\texttt{DigraphPeriod}}
\label{DigraphPeriod}
}\hfill{\scriptsize (attribute)}}\\
\textbf{\indent Returns:\ }
An integer.



 This function returns the period of the digraph \mbox{\texttt{\mdseries\slshape digraph}}.

 If a digraph \mbox{\texttt{\mdseries\slshape digraph}} has at least one directed cycle, then the period is the greatest positive
integer which divides the lengths of all directed cycles of \mbox{\texttt{\mdseries\slshape digraph}}. If \mbox{\texttt{\mdseries\slshape digraph}} has no directed cycles, then this function returns $0$. See Section \ref{Definitions} for the definition of a directed cycle. 

 A digraph with a period of $1$ is said to be \emph{aperiodic}. See \texttt{IsAperiodicDigraph} (\ref{IsAperiodicDigraph}). 

 
\begin{Verbatim}[commandchars=!@|,fontsize=\small,frame=single,label=Example]
  !gapprompt@gap>| !gapinput@D := Digraph([[6], [1], [2], [3], [4, 4], [5]]);|
  <immutable multidigraph with 6 vertices, 7 edges>
  !gapprompt@gap>| !gapinput@DigraphPeriod(D);|
  6
  !gapprompt@gap>| !gapinput@D := Digraph([[2], [3, 5], [4], [5], [1, 2]]);|
  <immutable digraph with 5 vertices, 7 edges>
  !gapprompt@gap>| !gapinput@DigraphPeriod(D);|
  1
  !gapprompt@gap>| !gapinput@D := ChainDigraph(2);|
  <immutable chain digraph with 2 vertices>
  !gapprompt@gap>| !gapinput@DigraphPeriod(D);|
  0
  !gapprompt@gap>| !gapinput@IsAcyclicDigraph(D);|
  true
  !gapprompt@gap>| !gapinput@D := GeneralisedPetersenGraph(IsMutableDigraph, 5, 2);|
  <mutable digraph with 10 vertices, 30 edges>
  !gapprompt@gap>| !gapinput@DigraphPeriod(D);|
  1
\end{Verbatim}
 }

 

\subsection{\textcolor{Chapter }{DigraphFloydWarshall}}
\logpage{[ 5, 4, 19 ]}\nobreak
\hyperdef{L}{X864A31A8809F61C2}{}
{\noindent\textcolor{FuncColor}{$\triangleright$\enspace\texttt{DigraphFloydWarshall({\mdseries\slshape digraph, func, nopath, edge})\index{DigraphFloydWarshall@\texttt{DigraphFloydWarshall}}
\label{DigraphFloydWarshall}
}\hfill{\scriptsize (operation)}}\\
\textbf{\indent Returns:\ }
A matrix.



 If \mbox{\texttt{\mdseries\slshape digraph}} is a digraph with $n$ vertices, then this operation returns an $n \times n$ matrix \texttt{mat} containing the output of a generalised version of the
Floyd\texttt{\symbol{45}}Warshall algorithm, applied to \mbox{\texttt{\mdseries\slshape digraph}}. 

 The operation \texttt{DigraphFloydWarshall} is customised by the arguments \mbox{\texttt{\mdseries\slshape func}}, \mbox{\texttt{\mdseries\slshape nopath}}, and \mbox{\texttt{\mdseries\slshape edge}}. The arguments \mbox{\texttt{\mdseries\slshape nopath}} and \mbox{\texttt{\mdseries\slshape edge}} can be arbitrary \textsf{GAP} objects. The argument \mbox{\texttt{\mdseries\slshape func}} must be a function which accepts 4 arguments: the matrix \texttt{mat}, followed by 3 positive integers. The function \mbox{\texttt{\mdseries\slshape func}} is where the work to calculate the desired outcome must be performed. 

 This method initialises the matrix \texttt{mat} by setting entry \texttt{mat[i][j]} to equal \mbox{\texttt{\mdseries\slshape edge}} if there is an edge with source \texttt{i} and range \texttt{j}, and by setting entry \texttt{mat[i][j]} to equal \mbox{\texttt{\mdseries\slshape nopath}} otherwise. The final part of \texttt{DigraphFloydWarshall} then calls the function \mbox{\texttt{\mdseries\slshape func}} inside three nested for loops, over the vertices of \mbox{\texttt{\mdseries\slshape digraph}}: 

 
\begin{Verbatim}[commandchars=!@|,fontsize=\small,frame=single,label=]
  for i in DigraphsVertices(digraph) do
    for j in DigraphsVertices(digraph) do
      for k in DigraphsVertices(digraph) do
        func(mat, i, j, k);
      od;
    od;
  od;
\end{Verbatim}
 The matrix \texttt{mat} is then returned as the result. An example of using \texttt{DigraphFloydWarshall} to calculate the shortest (non\texttt{\symbol{45}}zero) distances between the
vertices of a digraph is shown below: 

 
\begin{Verbatim}[commandchars=!@|,fontsize=\small,frame=single,label=Example]
  !gapprompt@gap>| !gapinput@D := DigraphFromDigraph6String("&EAHQeDB");|
  <immutable digraph with 6 vertices, 12 edges>
  !gapprompt@gap>| !gapinput@func := function(mat, i, j, k)|
  !gapprompt@>| !gapinput@  if mat[i][k] <> -1 and mat[k][j] <> -1 then|
  !gapprompt@>| !gapinput@    if (mat[i][j] = -1) or (mat[i][j] > mat[i][k] + mat[k][j]) then|
  !gapprompt@>| !gapinput@      mat[i][j] := mat[i][k] + mat[k][j];|
  !gapprompt@>| !gapinput@    fi;|
  !gapprompt@>| !gapinput@  fi;|
  !gapprompt@>| !gapinput@end;|
  function( mat, i, j, k ) ... end
  !gapprompt@gap>| !gapinput@shortest_distances := DigraphFloydWarshall(D, func, -1, 1);;|
  !gapprompt@gap>| !gapinput@Display(shortest_distances);|
  [ [   3,  -1,  -1,   2,   1,   2 ],
    [   4,   2,   1,   3,   2,   1 ],
    [   3,   1,   2,   2,   1,   2 ],
    [   1,  -1,  -1,   1,   1,   2 ],
    [   2,  -1,  -1,   1,   2,   1 ],
    [   3,  -1,  -1,   2,   1,   1 ] ]
\end{Verbatim}
 }

 

\subsection{\textcolor{Chapter }{IsReachable}}
\logpage{[ 5, 4, 20 ]}\nobreak
\hyperdef{L}{X7FBAB09E7C0BE5CF}{}
{\noindent\textcolor{FuncColor}{$\triangleright$\enspace\texttt{IsReachable({\mdseries\slshape digraph, u, v})\index{IsReachable@\texttt{IsReachable}}
\label{IsReachable}
}\hfill{\scriptsize (operation)}}\\
\textbf{\indent Returns:\ }
\texttt{true} or \texttt{false}.



 This operation returns \texttt{true} if there exists a non\texttt{\symbol{45}}trivial directed walk from vertex \mbox{\texttt{\mdseries\slshape u}} to vertex \mbox{\texttt{\mdseries\slshape v}} in the digraph \mbox{\texttt{\mdseries\slshape digraph}}, and \texttt{false} if there does not exist such a directed walk. See Section \ref{Definitions} for the definition of a non\texttt{\symbol{45}}trivial directed walk. 

 The method for \texttt{IsReachable} has worst case complexity of $O(m + n)$ where $m$ is the number of edges and $n$ the number of vertices in \mbox{\texttt{\mdseries\slshape digraph}}. 
\begin{Verbatim}[commandchars=!@|,fontsize=\small,frame=single,label=Example]
  !gapprompt@gap>| !gapinput@D := Digraph([[2], [3], [2, 3]]);|
  <immutable digraph with 3 vertices, 4 edges>
  !gapprompt@gap>| !gapinput@IsReachable(D, 1, 3);|
  true
  !gapprompt@gap>| !gapinput@IsReachable(D, 2, 1);|
  false
  !gapprompt@gap>| !gapinput@IsReachable(D, 3, 3);|
  true
  !gapprompt@gap>| !gapinput@IsReachable(D, 1, 1);|
  false
\end{Verbatim}
 }

 

\subsection{\textcolor{Chapter }{IsDigraphPath}}
\logpage{[ 5, 4, 21 ]}\nobreak
\hyperdef{L}{X7ECD22877AEA89CC}{}
{\noindent\textcolor{FuncColor}{$\triangleright$\enspace\texttt{IsDigraphPath({\mdseries\slshape D, v, a})\index{IsDigraphPath@\texttt{IsDigraphPath}}
\label{IsDigraphPath}
}\hfill{\scriptsize (operation)}}\\
\noindent\textcolor{FuncColor}{$\triangleright$\enspace\texttt{IsDigraphPath({\mdseries\slshape D, list})\index{IsDigraphPath@\texttt{IsDigraphPath}!for a digraph and list}
\label{IsDigraphPath:for a digraph and list}
}\hfill{\scriptsize (operation)}}\\
\textbf{\indent Returns:\ }
\texttt{true} or \texttt{false}.



 This function returns \texttt{true} if the arguments \mbox{\texttt{\mdseries\slshape v}} and \mbox{\texttt{\mdseries\slshape a }} describe a path in the digraph \mbox{\texttt{\mdseries\slshape D}}. A directed path (or directed cycle) of non\texttt{\symbol{45}}zero length \texttt{n\texttt{\symbol{45}}1}, $(v_1, e_1, v_2, e_2, ..., e_{n-1}, v_n)$, is represented by a pair of lists \texttt{[\mbox{\texttt{\mdseries\slshape v}}, \mbox{\texttt{\mdseries\slshape a}}]} as follows: 
\begin{itemize}
\item  \mbox{\texttt{\mdseries\slshape v}} is the list $[v_1, v_2, ..., v_n]$. 
\item  \mbox{\texttt{\mdseries\slshape a}} is the list of positive integers $[a_1, a_2, ..., a_{n-1}]$ where for each each $i < n$, $a_i$ is the position of $v_{i+1}$ in \texttt{OutNeighboursOfVertex(}\mbox{\texttt{\mdseries\slshape D}}\texttt{,}$v_i$\texttt{)} corresponding to the edge $e_i$. This can be useful if the position of a vertex in a list of
out\texttt{\symbol{45}}neighours is significant, for example in orbit
digraphs. 
\end{itemize}
 If the arguments to \texttt{IsDigraphPath} are a digraph \mbox{\texttt{\mdseries\slshape D}} and list \mbox{\texttt{\mdseries\slshape list}}, then this is equivalent to calling \texttt{IsDigraphPath(\mbox{\texttt{\mdseries\slshape D}}, \mbox{\texttt{\mdseries\slshape list}}[1], \mbox{\texttt{\mdseries\slshape list}}[2])}. 
\begin{Verbatim}[commandchars=!@|,fontsize=\small,frame=single,label=Example]
  !gapprompt@gap>| !gapinput@D := Digraph(IsMutableDigraph, Combinations([1 .. 5]), IsSubset);|
  <mutable digraph with 32 vertices, 243 edges>
  !gapprompt@gap>| !gapinput@DigraphReflexiveTransitiveReduction(D);|
  <mutable digraph with 32 vertices, 80 edges>
  !gapprompt@gap>| !gapinput@MakeImmutable(D);|
  <immutable digraph with 32 vertices, 80 edges>
  !gapprompt@gap>| !gapinput@IsDigraphPath(D, [32, 31, 33], [1, 1]);|
  false
  !gapprompt@gap>| !gapinput@IsDigraphPath(D, [1], []);|
  true
  !gapprompt@gap>| !gapinput@IsDigraphPath(D, [6, 9, 16, 17], [3, 3, 2]);|
  true
  !gapprompt@gap>| !gapinput@IsDigraphPath(D, DigraphPath(D, 6, 1));|
  true
\end{Verbatim}
 }

 

\subsection{\textcolor{Chapter }{VerticesReachableFrom (for a digraph and vertex)}}
\logpage{[ 5, 4, 22 ]}\nobreak
\hyperdef{L}{X7F255A2A84CB868C}{}
{\noindent\textcolor{FuncColor}{$\triangleright$\enspace\texttt{VerticesReachableFrom({\mdseries\slshape digraph, root})\index{VerticesReachableFrom@\texttt{VerticesReachableFrom}!for a digraph and vertex}
\label{VerticesReachableFrom:for a digraph and vertex}
}\hfill{\scriptsize (operation)}}\\
\noindent\textcolor{FuncColor}{$\triangleright$\enspace\texttt{VerticesReachableFrom({\mdseries\slshape digraph, list})\index{VerticesReachableFrom@\texttt{VerticesReachableFrom}!for a digraph and a list of vertices}
\label{VerticesReachableFrom:for a digraph and a list of vertices}
}\hfill{\scriptsize (operation)}}\\
\textbf{\indent Returns:\ }
A list.



 This operation returns a list of the vertices \texttt{v}, for which there exists a non\texttt{\symbol{45}}trivial directed walk from
the vertex \mbox{\texttt{\mdseries\slshape root}}, or any of the list of vertices \mbox{\texttt{\mdseries\slshape list}}, to vertex \texttt{v} in the digraph \mbox{\texttt{\mdseries\slshape digraph}}. See Section \ref{Definitions} for the definition of a non\texttt{\symbol{45}}trivial directed walk. 

 The method for \texttt{VerticesReachableFrom} has worst case complexity of $O(m + n)$ where $m$ is the number of edges and $n$ the number of vertices in \mbox{\texttt{\mdseries\slshape digraph}}. 

 This function returns an error if \mbox{\texttt{\mdseries\slshape root}}, or any vertex in \mbox{\texttt{\mdseries\slshape list}}, is not a vertices of \mbox{\texttt{\mdseries\slshape digraph}}. 
\begin{Verbatim}[commandchars=!@|,fontsize=\small,frame=single,label=Example]
  !gapprompt@gap>| !gapinput@D := CompleteDigraph(5);|
  <immutable complete digraph with 5 vertices>
  !gapprompt@gap>| !gapinput@VerticesReachableFrom(D, 1);|
  [ 1, 2, 3, 4, 5 ]
  !gapprompt@gap>| !gapinput@VerticesReachableFrom(D, 3);|
  [ 1, 2, 3, 4, 5 ]
  !gapprompt@gap>| !gapinput@D := EmptyDigraph(5);|
  <immutable empty digraph with 5 vertices>
  !gapprompt@gap>| !gapinput@VerticesReachableFrom(D, 1);|
  [  ]
  !gapprompt@gap>| !gapinput@VerticesReachableFrom(D, 3);|
  [  ]
  !gapprompt@gap>| !gapinput@VerticesReachableFrom(D, [1, 2, 3, 4]);|
  [  ]
  !gapprompt@gap>| !gapinput@VerticesReachableFrom(D, [3, 4, 5]);|
  [  ]
  !gapprompt@gap>| !gapinput@D := CycleDigraph(4);|
  <immutable cycle digraph with 4 vertices>
  !gapprompt@gap>| !gapinput@VerticesReachableFrom(D, 1);|
  [ 1, 2, 3, 4 ]
  !gapprompt@gap>| !gapinput@VerticesReachableFrom(D, 3);|
  [ 1, 2, 3, 4 ]
  !gapprompt@gap>| !gapinput@D := ChainDigraph(5);|
  <immutable chain digraph with 5 vertices>
  !gapprompt@gap>| !gapinput@VerticesReachableFrom(D, 1);|
  [ 2, 3, 4, 5 ]
  !gapprompt@gap>| !gapinput@VerticesReachableFrom(D, 3);|
  [ 4, 5 ]
  !gapprompt@gap>| !gapinput@VerticesReachableFrom(D, 5);|
  [  ]
  !gapprompt@gap>| !gapinput@VerticesReachableFrom(D, [3, 4]);|
  [ 4, 5 ]
\end{Verbatim}
 }

 

\subsection{\textcolor{Chapter }{DigraphPath}}
\logpage{[ 5, 4, 23 ]}\nobreak
\hyperdef{L}{X8039170B82A32257}{}
{\noindent\textcolor{FuncColor}{$\triangleright$\enspace\texttt{DigraphPath({\mdseries\slshape digraph, u, v})\index{DigraphPath@\texttt{DigraphPath}}
\label{DigraphPath}
}\hfill{\scriptsize (operation)}}\\
\textbf{\indent Returns:\ }
A pair of lists, or \texttt{fail}.



 If there exists a non\texttt{\symbol{45}}trivial directed path (or a
non\texttt{\symbol{45}}trivial cycle, in the case that \mbox{\texttt{\mdseries\slshape u}} \texttt{=} \mbox{\texttt{\mdseries\slshape v}}) from vertex \mbox{\texttt{\mdseries\slshape u}} to vertex \mbox{\texttt{\mdseries\slshape v}} in the digraph \mbox{\texttt{\mdseries\slshape digraph}}, then this operation returns such a directed path (or directed cycle).
Otherwise, this operation returns \texttt{fail}. See Section \hyperref[Definitions]{`Definitions'} for the definition of a directed path and a directed cycle. 

 A directed path (or directed cycle) of non\texttt{\symbol{45}}zero length \texttt{n\texttt{\symbol{45}}1}, $(v_1, e_1, v_2, e_2, ..., e_{n-1}, v_n)$, is represented by a pair of lists \texttt{[v,a]} as follows: 
\begin{itemize}
\item  \texttt{v} is the list $[v_1, v_2, ..., v_n]$. 
\item  \texttt{a} is the list of positive integers $[a_1, a_2, ..., a_{n-1}]$ where for each each $i < n$, $a_i$ is the position of $v_{i+1}$ in \texttt{OutNeighboursOfVertex(}\mbox{\texttt{\mdseries\slshape digraph}}\texttt{,}$v_i$\texttt{)} corresponding to the edge $e_i$. This can be useful if the position of a vertex in a list of
out\texttt{\symbol{45}}neighours is significant, for example in orbit
digraphs. 
\end{itemize}
 The method for \texttt{DigraphPath} has worst case complexity of $O(m + n)$ where $m$ is the number of edges and $n$ the number of vertices in \mbox{\texttt{\mdseries\slshape digraph}}.

 See also \texttt{IsDigraphPath} (\ref{IsDigraphPath}). 
\begin{Verbatim}[commandchars=!@|,fontsize=\small,frame=single,label=Example]
  !gapprompt@gap>| !gapinput@D := Digraph([[2], [3], [2, 3]]);|
  <immutable digraph with 3 vertices, 4 edges>
  !gapprompt@gap>| !gapinput@DigraphPath(D, 1, 3);|
  [ [ 1, 2, 3 ], [ 1, 1 ] ]
  !gapprompt@gap>| !gapinput@DigraphPath(D, 2, 1);|
  fail
  !gapprompt@gap>| !gapinput@DigraphPath(D, 3, 3);|
  [ [ 3, 3 ], [ 2 ] ]
  !gapprompt@gap>| !gapinput@DigraphPath(D, 1, 1);|
  fail
\end{Verbatim}
 }

 

\subsection{\textcolor{Chapter }{DigraphShortestPath}}
\logpage{[ 5, 4, 24 ]}\nobreak
\hyperdef{L}{X80E9D645843973A6}{}
{\noindent\textcolor{FuncColor}{$\triangleright$\enspace\texttt{DigraphShortestPath({\mdseries\slshape digraph, u, v})\index{DigraphShortestPath@\texttt{DigraphShortestPath}}
\label{DigraphShortestPath}
}\hfill{\scriptsize (operation)}}\\
\textbf{\indent Returns:\ }
A pair of lists, or \texttt{fail}.



 Returns a shortest directed path in the digraph \mbox{\texttt{\mdseries\slshape digraph}} from vertex \mbox{\texttt{\mdseries\slshape u}} to vertex \mbox{\texttt{\mdseries\slshape v}}, if such a path exists. If \texttt{\mbox{\texttt{\mdseries\slshape u}} = \mbox{\texttt{\mdseries\slshape v}}}, then the shortest non\texttt{\symbol{45}}trivial cycle is returned, again,
if it exists. Otherwise, this operation returns \texttt{fail}. See Section \hyperref[Definitions]{`Definitions'} for the definition of a directed path and a directed cycle. 

 See \texttt{DigraphPath} (\ref{DigraphPath}) for details on the output. The method for \texttt{DigraphShortestPath} has worst case complexity of $O(m + n)$ where $m$ is the number of edges and $n$ the number of vertices in \mbox{\texttt{\mdseries\slshape digraph}}. 
\begin{Verbatim}[commandchars=!@|,fontsize=\small,frame=single,label=Example]
  !gapprompt@gap>| !gapinput@D := Digraph([[1, 2], [3], [2, 4], [1], [2, 4]]);|
  <immutable digraph with 5 vertices, 8 edges>
  !gapprompt@gap>| !gapinput@DigraphShortestPath(D, 5, 1);|
  [ [ 5, 4, 1 ], [ 2, 1 ] ]
  !gapprompt@gap>| !gapinput@DigraphShortestPath(D, 3, 3);|
  [ [ 3, 2, 3 ], [ 1, 1 ] ]
  !gapprompt@gap>| !gapinput@DigraphShortestPath(D, 5, 5);|
  fail
  !gapprompt@gap>| !gapinput@DigraphShortestPath(D, 1, 1);|
  [ [ 1, 1 ], [ 1 ] ]
\end{Verbatim}
 }

 

\subsection{\textcolor{Chapter }{DigraphRandomWalk}}
\logpage{[ 5, 4, 25 ]}\nobreak
\hyperdef{L}{X7C45396C808308C4}{}
{\noindent\textcolor{FuncColor}{$\triangleright$\enspace\texttt{DigraphRandomWalk({\mdseries\slshape digraph, v, t})\index{DigraphRandomWalk@\texttt{DigraphRandomWalk}}
\label{DigraphRandomWalk}
}\hfill{\scriptsize (operation)}}\\
\textbf{\indent Returns:\ }
A pair of lists.



 Returns a directed path corresponding to a \emph{random walk} in the digraph \mbox{\texttt{\mdseries\slshape digraph}}, starting at vertex \mbox{\texttt{\mdseries\slshape v}} and having length no more than \mbox{\texttt{\mdseries\slshape t}}. 

 A random walk is defined as follows. The path begins at \mbox{\texttt{\mdseries\slshape v}}, and at each step it follows a random edge leaving the current vertex. It
continues through the digraph in this way until it has traversed \mbox{\texttt{\mdseries\slshape t}} edges, or until it reaches a vertex with no out\texttt{\symbol{45}}edges (a \emph{sink}) and therefore cannot continue. 

 The output has the same form as that of \texttt{DigraphPath} (\ref{DigraphPath}). 

 
\begin{Verbatim}[commandchars=!@|,fontsize=\small,frame=single,label=Example]
  !gapprompt@gap>| !gapinput@D := Digraph([[1, 2], [3], [2, 4], [1], [2, 4]]);                |
  <immutable digraph with 5 vertices, 8 edges>
  !gapprompt@gap>| !gapinput@DigraphRandomWalk(D, 1, 4);|
  [ [ 1, 2, 3, 2, 3 ], [ 2, 1, 1, 1 ] ]
\end{Verbatim}
 }

 

\subsection{\textcolor{Chapter }{DigraphAbsorptionProbabilities}}
\logpage{[ 5, 4, 26 ]}\nobreak
\hyperdef{L}{X787459247B2005E6}{}
{\noindent\textcolor{FuncColor}{$\triangleright$\enspace\texttt{DigraphAbsorptionProbabilities({\mdseries\slshape digraph})\index{DigraphAbsorptionProbabilities@\texttt{DigraphAbsorptionProbabilities}}
\label{DigraphAbsorptionProbabilities}
}\hfill{\scriptsize (attribute)}}\\
\textbf{\indent Returns:\ }
A matrix of rational numbers.



 A random walk of infinite length through a finite digraph may pass through
several different strongly connected components (SCCs). However, with
probability 1 it must eventually reach an SCC which it can never leave,
because the SCC has no out\texttt{\symbol{45}}edges leading to other SCCs. We
may say that such an SCC has \emph{absorbed} the random walk, and we can calculate the probability of each SCC absorbing a
walk starting at a given vertex. 

 \texttt{DigraphAbsorptionProbabilities} returns an $m \times n$ matrix \texttt{mat}, where $m$ is the number of vertices in \texttt{digraph} and $n$ is the number of strongly connected components. Each entry \texttt{mat[i][j]} is a rational number representing the probability that an unbounded random
walk starting at vertex \texttt{i} will be absorbed by strongly connected component \texttt{j} {\textendash} that is, the probability that it will reach the component and
never leave. 

 Strongly connected components are indexed in the order given by \texttt{DigraphStronglyConnectedComponents} (\ref{DigraphStronglyConnectedComponents}). See also \texttt{DigraphRandomWalk} (\ref{DigraphRandomWalk}) and \texttt{DigraphAbsorptionExpectedSteps} (\ref{DigraphAbsorptionExpectedSteps}). 

 
\begin{Verbatim}[commandchars=!@|,fontsize=\small,frame=single,label=Example]
  !gapprompt@gap>| !gapinput@gr := Digraph([[2, 3, 4], [3], [2], []]);|
  <immutable digraph with 4 vertices, 5 edges>
  !gapprompt@gap>| !gapinput@DigraphStronglyConnectedComponents(gr).comps;|
  [ [ 2, 3 ], [ 4 ], [ 1 ] ]
  !gapprompt@gap>| !gapinput@DigraphAbsorptionProbabilities(gr);|
  [ [ 2/3, 1/3, 0 ], [ 1, 0, 0 ], [ 1, 0, 0 ], [ 0, 1, 0 ] ]
\end{Verbatim}
 }

 

\subsection{\textcolor{Chapter }{DigraphAbsorptionExpectedSteps}}
\logpage{[ 5, 4, 27 ]}\nobreak
\hyperdef{L}{X7B0A62CF7F6F23E9}{}
{\noindent\textcolor{FuncColor}{$\triangleright$\enspace\texttt{DigraphAbsorptionExpectedSteps({\mdseries\slshape digraph})\index{DigraphAbsorptionExpectedSteps@\texttt{DigraphAbsorptionExpectedSteps}}
\label{DigraphAbsorptionExpectedSteps}
}\hfill{\scriptsize (attribute)}}\\
\textbf{\indent Returns:\ }
A list of rational numbers.



 A random walk of unbounded length through a finite digraph may pass through
several different strongly connected components (SCCs). However, with
probability 1 it must eventually reach an SCC which it can never leave,
because the SCC has no out\texttt{\symbol{45}}edges leading to other SCCs.
When this happens, we say the walk has been \emph{absorbed}. 

 \texttt{DigraphAbsorptionExpectedSteps} returns a list \texttt{L} with length equal to the number of vertices of \mbox{\texttt{\mdseries\slshape digraph}}, where \texttt{L[v]} is the expected number of steps a random walk starting at vertex \texttt{v} should take before it is absorbed {\textendash} that is, the expected number
of steps to reach an SCC that it can never leave. 

 See also \texttt{DigraphRandomWalk} (\ref{DigraphRandomWalk}) and \texttt{DigraphAbsorptionProbabilities} (\ref{DigraphAbsorptionProbabilities}). 

 
\begin{Verbatim}[commandchars=!@|,fontsize=\small,frame=single,label=Example]
  !gapprompt@gap>| !gapinput@gr := Digraph([[2], [3, 4], [], [2]]);|
  <immutable digraph with 4 vertices, 4 edges>
  !gapprompt@gap>| !gapinput@DigraphAbsorptionExpectedSteps(gr);|
  [ 4, 3, 0, 4 ]
\end{Verbatim}
 }

 

\subsection{\textcolor{Chapter }{Dominators}}
\logpage{[ 5, 4, 28 ]}\nobreak
\hyperdef{L}{X7FE79CB278CE6991}{}
{\noindent\textcolor{FuncColor}{$\triangleright$\enspace\texttt{Dominators({\mdseries\slshape digraph, root})\index{Dominators@\texttt{Dominators}}
\label{Dominators}
}\hfill{\scriptsize (operation)}}\\
\textbf{\indent Returns:\ }
A list of lists.



 \texttt{Dominators} takes a \mbox{\texttt{\mdseries\slshape digraph}} and a root \mbox{\texttt{\mdseries\slshape root}} and returns the dominators of each vertex with respect to the root. The output
is returned as a list of length \texttt{DigraphNrVertices(\mbox{\texttt{\mdseries\slshape Digraph}})}, whose \mbox{\texttt{\mdseries\slshape i}}th entry is a list with the dominators of vertex \mbox{\texttt{\mdseries\slshape i}} of the \mbox{\texttt{\mdseries\slshape digraph}}. If there is no path from the root to a specific vertex, the output will
contain a hole in the corresponding position. The \emph{dominators} of a vertex $u$ are the vertices that are contained in every path from the $root$ to $u$, not including $u$ itself. The method for this operation is an implementation of an algorithm by
Thomas Lengauer and Robert Endre Tarjan \cite{LT79}. The complexity of this algorithm is $O(mlog n)$ where $m$ is the number of edges and $n$ is the number of nodes in the subdigraph induced by the nodes in \mbox{\texttt{\mdseries\slshape digraph}} reachable from \mbox{\texttt{\mdseries\slshape root}}. 
\begin{Verbatim}[commandchars=!@|,fontsize=\small,frame=single,label=Example]
  !gapprompt@gap>| !gapinput@D := Digraph([[2], [3, 6], [2, 4], [1], [], [3]]);|
  <immutable digraph with 6 vertices, 7 edges>
  !gapprompt@gap>| !gapinput@Dominators(D, 1);|
  [ , [ 1 ], [ 2, 1 ], [ 3, 2, 1 ],, [ 2, 1 ] ]
  !gapprompt@gap>| !gapinput@Dominators(D, 2);|
  [ [ 4, 3, 2 ],, [ 2 ], [ 3, 2 ],, [ 2 ] ]
  !gapprompt@gap>| !gapinput@Dominators(D, 3);|
  [ [ 4, 3 ], [ 3 ],, [ 3 ],, [ 2, 3 ] ]
  !gapprompt@gap>| !gapinput@Dominators(D, 4);|
  [ [ 4 ], [ 1, 4 ], [ 2, 1, 4 ],,, [ 2, 1, 4 ] ]
  !gapprompt@gap>| !gapinput@Dominators(D, 5);|
  [  ]
  !gapprompt@gap>| !gapinput@Dominators(D, 6);|
  [ [ 4, 3, 6 ], [ 3, 6 ], [ 6 ], [ 3, 6 ] ]
\end{Verbatim}
 }

 

\subsection{\textcolor{Chapter }{DominatorTree}}
\logpage{[ 5, 4, 29 ]}\nobreak
\hyperdef{L}{X7D66A0FB7F6100FB}{}
{\noindent\textcolor{FuncColor}{$\triangleright$\enspace\texttt{DominatorTree({\mdseries\slshape digraph, root})\index{DominatorTree@\texttt{DominatorTree}}
\label{DominatorTree}
}\hfill{\scriptsize (operation)}}\\
\textbf{\indent Returns:\ }
A record.



 \texttt{DominatorTree} takes a \mbox{\texttt{\mdseries\slshape digraph}} and a \mbox{\texttt{\mdseries\slshape root}} vertex and returns a record with the following components: 
\begin{description}
\item[{idom}]  the immediate dominators of the vertices with respect to the root. 
\item[{preorder}]  the preorder values of the vertices defined by the depth first search executed
on the digraph. 
\end{description}
 The \emph{immediate dominator} of a vertex $u$ is the unique dominator of $u$ that is dominated by all other dominators of $u$. The algorithm is an implementation of the fast algorithm written by Thomas
Lengauer and Robert Endre Tarjan \cite{LT79}. The complexity of this algorithm is $O(mlog n)$ where $m$ is the number of edges and $n$ is the number of nodes in the subdigraph induced by the nodes in \mbox{\texttt{\mdseries\slshape digraph}} reachable from \mbox{\texttt{\mdseries\slshape root}}. 
\begin{Verbatim}[commandchars=!@|,fontsize=\small,frame=single,label=Example]
  !gapprompt@gap>| !gapinput@D := Digraph([[2, 3], [4, 6], [4, 5], [3, 5], [1, 6], [2, 3]]);|
  <immutable digraph with 6 vertices, 12 edges>
  !gapprompt@gap>| !gapinput@DominatorTree(D, 1);|
  rec( idom := [ fail, 1, 1, 1, 1, 1 ], 
    preorder := [ 1, 2, 4, 3, 5, 6 ] )
  !gapprompt@gap>| !gapinput@DominatorTree(D, 5);|
  rec( idom := [ 5, 5, 5, 5, fail, 5 ], 
    preorder := [ 5, 1, 2, 4, 3, 6 ] )
  !gapprompt@gap>| !gapinput@D := CompleteDigraph(5);|
  <immutable complete digraph with 5 vertices>
  !gapprompt@gap>| !gapinput@DominatorTree(D, 1);|
  rec( idom := [ fail, 1, 1, 1, 1 ], preorder := [ 1, 2, 3, 4, 5 ] )
  !gapprompt@gap>| !gapinput@DominatorTree(D, 2);|
  rec( idom := [ 2, fail, 2, 2, 2 ], preorder := [ 2, 1, 3, 4, 5 ] )
\end{Verbatim}
 }

 

\subsection{\textcolor{Chapter }{IteratorOfPaths}}
\logpage{[ 5, 4, 30 ]}\nobreak
\hyperdef{L}{X7C0416FE7A69CA2C}{}
{\noindent\textcolor{FuncColor}{$\triangleright$\enspace\texttt{IteratorOfPaths({\mdseries\slshape digraph, u, v})\index{IteratorOfPaths@\texttt{IteratorOfPaths}}
\label{IteratorOfPaths}
}\hfill{\scriptsize (operation)}}\\
\textbf{\indent Returns:\ }
An iterator.



 If \mbox{\texttt{\mdseries\slshape digraph}} is a digraph or a list of adjacencies which defines a digraph
\texttt{\symbol{45}} see \texttt{OutNeighbours} (\ref{OutNeighbours}) \texttt{\symbol{45}} then this operation returns an iterator of the
non\texttt{\symbol{45}}trivial directed paths (or directed cycles, in the case
that \mbox{\texttt{\mdseries\slshape u}} \texttt{=} \mbox{\texttt{\mdseries\slshape v}}) in \mbox{\texttt{\mdseries\slshape digraph}} from the vertex \mbox{\texttt{\mdseries\slshape u}} to the vertex \mbox{\texttt{\mdseries\slshape v}}. 

 See \texttt{DigraphPath} (\ref{DigraphPath}) for more information about the representation of a directed path or directed
cycle which is used, and see  (\textbf{Reference: Iterators}) for more information about iterators. See Section \hyperref[Definitions]{`Definitions'} for the definition of a directed path and a directed cycle. 

 
\begin{Verbatim}[commandchars=!@|,fontsize=\small,frame=single,label=Example]
  !gapprompt@gap>| !gapinput@D := Digraph([[1, 4, 4, 2], [3, 5], [2, 3], [1, 2], [4]]);|
  <immutable multidigraph with 5 vertices, 11 edges>
  !gapprompt@gap>| !gapinput@iter := IteratorOfPaths(D, 1, 4);|
  <iterator>
  !gapprompt@gap>| !gapinput@NextIterator(iter);|
  [ [ 1, 4 ], [ 2 ] ]
  !gapprompt@gap>| !gapinput@NextIterator(iter);|
  [ [ 1, 4 ], [ 3 ] ]
  !gapprompt@gap>| !gapinput@NextIterator(iter);|
  [ [ 1, 2, 5, 4 ], [ 4, 2, 1 ] ]
  !gapprompt@gap>| !gapinput@IsDoneIterator(iter);|
  true
  !gapprompt@gap>| !gapinput@iter := IteratorOfPaths(D, 4, 3);|
  <iterator>
  !gapprompt@gap>| !gapinput@NextIterator(iter);|
  [ [ 4, 1, 2, 3 ], [ 1, 4, 1 ] ]
\end{Verbatim}
 }

 

\subsection{\textcolor{Chapter }{DigraphAllSimpleCircuits}}
\logpage{[ 5, 4, 31 ]}\nobreak
\hyperdef{L}{X7ECD16838704FAAA}{}
{\noindent\textcolor{FuncColor}{$\triangleright$\enspace\texttt{DigraphAllSimpleCircuits({\mdseries\slshape digraph})\index{DigraphAllSimpleCircuits@\texttt{DigraphAllSimpleCircuits}}
\label{DigraphAllSimpleCircuits}
}\hfill{\scriptsize (attribute)}}\\
\textbf{\indent Returns:\ }
A list of lists of vertices.



 If \mbox{\texttt{\mdseries\slshape digraph}} is a digraph, then \texttt{DigraphAllSimpleCircuits} returns a list of the \emph{simple circuits} in \mbox{\texttt{\mdseries\slshape digraph}}. 

 See Section \ref{Definitions} for the definition of a simple circuit, and related notions. Note that a loop
is a simple circuit. 

 For a digraph without multiple edges, a simple circuit is uniquely determined
by its subsequence of vertices. However this is not the case for a
multidigraph. The attribute \texttt{DigraphAllSimpleCircuits} ignores multiple edges, and identifies a simple circuit using only its
subsequence of vertices. For example, although the simple circuits $(v, e, v)$ and $(v, e', v)$ (for distinct edges $e$ and $e'$) are mathematically distinct, \texttt{DigraphAllSimpleCircuits} considers them to be the same. 

 With this approach, a directed circuit of length \texttt{n} can be determined by a list of its first \texttt{n} vertices. Thus a simple circuit $(v_1, e_1, v_2, e_2, ..., e_{n-1}, v_n, e_{n+1}, v_1)$ can be represented as the list $[v_1, \ldots, v_n]$, or any cyclic permutation thereof. For each simple circuit of \mbox{\texttt{\mdseries\slshape digraph}}, \texttt{DigraphAllSimpleCircuits(\mbox{\texttt{\mdseries\slshape digraph}})} includes precisely one such list to represent the circuit. 

 
\begin{Verbatim}[commandchars=!@|,fontsize=\small,frame=single,label=Example]
  !gapprompt@gap>| !gapinput@D := Digraph([[], [3], [2, 4], [5, 4], [4]]);|
  <immutable digraph with 5 vertices, 6 edges>
  !gapprompt@gap>| !gapinput@DigraphAllSimpleCircuits(D);|
  [ [ 4 ], [ 4, 5 ], [ 2, 3 ] ]
  !gapprompt@gap>| !gapinput@D := ChainDigraph(10);;|
  !gapprompt@gap>| !gapinput@DigraphAllSimpleCircuits(D);|
  [  ]
  !gapprompt@gap>| !gapinput@D := Digraph([[3], [1], [1]]);|
  <immutable digraph with 3 vertices, 3 edges>
  !gapprompt@gap>| !gapinput@DigraphAllSimpleCircuits(D);|
  [ [ 1, 3 ] ]
  !gapprompt@gap>| !gapinput@D := Digraph([[1, 1]]);|
  <immutable multidigraph with 1 vertex, 2 edges>
  !gapprompt@gap>| !gapinput@DigraphAllSimpleCircuits(D);|
  [ [ 1 ] ]
  !gapprompt@gap>| !gapinput@D := CycleDigraph(IsMutableDigraph, 3);|
  <mutable digraph with 3 vertices, 3 edges>
  !gapprompt@gap>| !gapinput@DigraphAllSimpleCircuits(D);|
  [ [ 1, 2, 3 ] ]
\end{Verbatim}
 }

 

\subsection{\textcolor{Chapter }{DigraphLongestSimpleCircuit}}
\logpage{[ 5, 4, 32 ]}\nobreak
\hyperdef{L}{X7C735C4E86BDD5F6}{}
{\noindent\textcolor{FuncColor}{$\triangleright$\enspace\texttt{DigraphLongestSimpleCircuit({\mdseries\slshape digraph})\index{DigraphLongestSimpleCircuit@\texttt{DigraphLongestSimpleCircuit}}
\label{DigraphLongestSimpleCircuit}
}\hfill{\scriptsize (attribute)}}\\
\textbf{\indent Returns:\ }
A list of vertices, or \texttt{fail}.



 If \mbox{\texttt{\mdseries\slshape digraph}} is a digraph, then \texttt{DigraphLongestSimpleCircuit} returns the longest \emph{simple circuit} in \mbox{\texttt{\mdseries\slshape digraph}}. See Section \ref{Definitions} for the definition of simple circuit, and the definition of length for a
simple circuit.

 This attribute computes \texttt{DigraphAllSimpleCircuits(}\mbox{\texttt{\mdseries\slshape digraph}}\texttt{)} to find all the simple circuits of \mbox{\texttt{\mdseries\slshape digraph}}, and returns one of maximal length. A simple circuit is represented as a list
of vertices, in the same way as described in \texttt{DigraphAllSimpleCircuits} (\ref{DigraphAllSimpleCircuits}).

 If \mbox{\texttt{\mdseries\slshape digraph}} has no simple circuits, then this attribute returns \texttt{fail}. If \mbox{\texttt{\mdseries\slshape digraph}} has multiple simple circuits of maximal length, then this attribute returns
one of them.

 
\begin{Verbatim}[commandchars=!@|,fontsize=\small,frame=single,label=Example]
  !gapprompt@gap>| !gapinput@D := Digraph([[], [3], [2, 4], [5, 4], [4]]);;|
  !gapprompt@gap>| !gapinput@DigraphLongestSimpleCircuit(D);|
  [ 4, 5 ]
  !gapprompt@gap>| !gapinput@D := ChainDigraph(10);;|
  !gapprompt@gap>| !gapinput@DigraphLongestSimpleCircuit(D);|
  fail
  !gapprompt@gap>| !gapinput@D := Digraph([[3], [1], [1, 4], [1, 1]]);;|
  !gapprompt@gap>| !gapinput@DigraphLongestSimpleCircuit(D);|
  [ 1, 3, 4 ]
  !gapprompt@gap>| !gapinput@D := GeneralisedPetersenGraph(IsMutableDigraph, 4, 1);|
  <mutable digraph with 8 vertices, 24 edges>
  !gapprompt@gap>| !gapinput@DigraphLongestSimpleCircuit(D);|
  [ 1, 2, 3, 4, 8, 7, 6, 5 ]
\end{Verbatim}
 }

 

\subsection{\textcolor{Chapter }{DigraphAllUndirectedSimpleCircuits}}
\logpage{[ 5, 4, 33 ]}\nobreak
\hyperdef{L}{X7F9EDAA8854FB538}{}
{\noindent\textcolor{FuncColor}{$\triangleright$\enspace\texttt{DigraphAllUndirectedSimpleCircuits({\mdseries\slshape digraph})\index{DigraphAllUndirectedSimpleCircuits@\texttt{DigraphAllUndirectedSimpleCircuits}}
\label{DigraphAllUndirectedSimpleCircuits}
}\hfill{\scriptsize (attribute)}}\\
\textbf{\indent Returns:\ }
A list of lists of vertices.



 If \mbox{\texttt{\mdseries\slshape digraph}} is a digraph, then \texttt{DigraphAllUndirectedSimpleCircuits} returns a list of the \emph{ undirected simple circuits} in \mbox{\texttt{\mdseries\slshape digraph}}. 

 See Section \ref{Definitions} for the definition of a simple circuit, and related notions. Note that a loop
is a simple circuit. A simple circuit is undirected if the orientation of the
edges in the circuit does not matter.

 The attribute \texttt{DigraphAllUndirectedSimpleCircuitss} ignores multiple edges, and identifies a undirected simple circuit using only
its subsequence of vertices. See Section \ref{DigraphAllSimpleCircuits} for more details on simple circuits.

 
\begin{Verbatim}[commandchars=!@|,fontsize=\small,frame=single,label=Example]
  !gapprompt@gap>| !gapinput@D := ChainDigraph(10);;|
  !gapprompt@gap>| !gapinput@DigraphAllUndirectedSimpleCircuits(D);|
  [  ]
  !gapprompt@gap>| !gapinput@D := Digraph([[1, 1]]);|
  <immutable multidigraph with 1 vertex, 2 edges>
  !gapprompt@gap>| !gapinput@DigraphAllUndirectedSimpleCircuits(D);|
  [ [ 1 ] ]
  !gapprompt@gap>| !gapinput@D := CycleDigraph(IsMutableDigraph, 3);|
  <mutable digraph with 3 vertices, 3 edges>
  !gapprompt@gap>| !gapinput@DigraphAllUndirectedSimpleCircuits(D);|
  [ [ 1, 2, 3 ] ]
  !gapprompt@gap>| !gapinput@D := DigraphSymmetricClosure(CycleDigraph(IsMutableDigraph, 3));|
  <mutable digraph with 3 vertices, 6 edges>
  !gapprompt@gap>| !gapinput@DigraphAllUndirectedSimpleCircuits(D);|
  [ [ 1, 2, 3 ] ]
\end{Verbatim}
 }

 

\subsection{\textcolor{Chapter }{DigraphAllChordlessCycles}}
\logpage{[ 5, 4, 34 ]}\nobreak
\hyperdef{L}{X7F0AC5617CE5E9DD}{}
{\noindent\textcolor{FuncColor}{$\triangleright$\enspace\texttt{DigraphAllChordlessCycles({\mdseries\slshape digraph})\index{DigraphAllChordlessCycles@\texttt{DigraphAllChordlessCycles}}
\label{DigraphAllChordlessCycles}
}\hfill{\scriptsize (attribute)}}\\
\noindent\textcolor{FuncColor}{$\triangleright$\enspace\texttt{DigraphAllChordlessCyclesOfMaximalLength({\mdseries\slshape digraph, maxLength})\index{DigraphAllChordlessCyclesOfMaximalLength@\texttt{Digraph}\-\texttt{All}\-\texttt{Chordless}\-\texttt{Cycles}\-\texttt{Of}\-\texttt{Maximal}\-\texttt{Length}}
\label{DigraphAllChordlessCyclesOfMaximalLength}
}\hfill{\scriptsize (attribute)}}\\
\textbf{\indent Returns:\ }
A list of lists of vertices.



 If \mbox{\texttt{\mdseries\slshape digraph}} is a digraph, then \texttt{DigraphAllChordlessCycles} returns a list of the \emph{chordless cycles} in \mbox{\texttt{\mdseries\slshape digraph}}. The method \texttt{DigraphAllChordlessCyclesOfMaximalLength} returns a list of all \emph{chordless cycles} in \mbox{\texttt{\mdseries\slshape digraph}} of length at most \mbox{\texttt{\mdseries\slshape maxLength}}

 A chordless cycle $C$ is a undirected simple circuit (see Section \ref{Definitions} ) where each pair of vertices in $C$ are not connected by an edge not in $C$. Here, cycles of length two are ignored.

 For a digraph without multiple edges, a simple circuit is uniquely determined
by its subsequence of vertices. However this is not the case for a
multidigraph. The attribute \texttt{DigraphAllChordlessCycles} ignores multiple edges, and identifies a simple circuit using only its
subsequence of vertices. See Section \ref{DigraphAllSimpleCircuits} for more details.

 This method uses the algorithms described in \cite{US14}.

 
\begin{Verbatim}[commandchars=!@|,fontsize=\small,frame=single,label=Example]
  !gapprompt@gap>| !gapinput@D := ChainDigraph(10);;|
  !gapprompt@gap>| !gapinput@DigraphAllChordlessCycles(D);|
  [  ]
  !gapprompt@gap>| !gapinput@D := Digraph([[1, 1]]);|
  <immutable multidigraph with 1 vertex, 2 edges>
  !gapprompt@gap>| !gapinput@DigraphAllChordlessCycles(D);|
  [  ]
  !gapprompt@gap>| !gapinput@D := CycleDigraph(3);;|
  !gapprompt@gap>| !gapinput@DigraphAllChordlessCycles(D);|
  [ [ 2, 1, 3 ] ]
  !gapprompt@gap>| !gapinput@D := CompleteDigraph(4);|
  <immutable complete digraph with 4 vertices>
  !gapprompt@gap>| !gapinput@DigraphAllChordlessCycles(D);|
  [ [ 2, 1, 3 ], [ 2, 1, 4 ], [ 3, 1, 4 ], [ 3, 2, 4 ] ]
\end{Verbatim}
 }

 

\subsection{\textcolor{Chapter }{FacialWalks}}
\logpage{[ 5, 4, 35 ]}\nobreak
\hyperdef{L}{X805F15398735AD7D}{}
{\noindent\textcolor{FuncColor}{$\triangleright$\enspace\texttt{FacialWalks({\mdseries\slshape digraph, list})\index{FacialWalks@\texttt{FacialWalks}}
\label{FacialWalks}
}\hfill{\scriptsize (operation)}}\\
\textbf{\indent Returns:\ }
A list of lists of vertices.



 If \mbox{\texttt{\mdseries\slshape digraph}} is an Eulerian digraph and \mbox{\texttt{\mdseries\slshape list}} is a rotation system of \mbox{\texttt{\mdseries\slshape digraph}}, then \texttt{FacialWalks} returns a list of the \emph{facial walks} in \mbox{\texttt{\mdseries\slshape digraph}}. 

 A rotation system defines for each vertex the ordering of the
out\texttt{\symbol{45}}neighbours. For example, the method \ref{PlanarEmbedding} computes for a planar digraph \mbox{\texttt{\mdseries\slshape D}} the rotation system of a planar embedding of \mbox{\texttt{\mdseries\slshape D}}. The facial walks of \mbox{\texttt{\mdseries\slshape digraph}} are closed walks and they are defined by the rotation system \mbox{\texttt{\mdseries\slshape list}}. They describe the boundaries of the faces of the embedding of \mbox{\texttt{\mdseries\slshape digraph}} given by the rotation system \mbox{\texttt{\mdseries\slshape list}}. The operation \texttt{FacialWalks} ignores multiple edges and loops.

 Here are some examples for planar embeddings: 
\begin{Verbatim}[commandchars=!@|,fontsize=\small,frame=single,label=Example]
  !gapprompt@gap>| !gapinput@D1 := CycleDigraph(4);;|
  !gapprompt@gap>| !gapinput@planar := PlanarEmbedding(D1);|
  [ [ 2 ], [ 3 ], [ 4 ], [ 1 ] ]
  !gapprompt@gap>| !gapinput@FacialWalks(D1, planar);|
  [ [ 1, 2, 3, 4 ] ]
  !gapprompt@gap>| !gapinput@nonPlanar := [[2, 4], [1, 3], [2, 4], [1, 3]];;|
  !gapprompt@gap>| !gapinput@FacialWalks(D1, nonPlanar);|
  [ [ 1, 2, 3, 4 ] ]
  !gapprompt@gap>| !gapinput@D2 := CompleteMultipartiteDigraph([2, 2, 2]);;|
  !gapprompt@gap>| !gapinput@rotationSystem := PlanarEmbedding(D2);|
  [ [ 3, 5, 4, 6 ], [ 6, 4, 5, 3 ], [ 6, 2, 5, 1 ], [ 1, 5, 2, 6 ],
    [ 1, 3, 2, 4 ], [ 1, 4, 2, 3 ] ]
  !gapprompt@gap>| !gapinput@FacialWalks(D2, rotationSystem);|
  [ [ 1, 3, 6 ], [ 1, 4, 5 ], [ 1, 5, 3 ], [ 1, 6, 4 ], [ 2, 3, 5 ],
    [ 2, 4, 6 ], [ 2, 5, 4 ], [ 2, 6, 3 ] ]
\end{Verbatim}
 Here is an example of a non\texttt{\symbol{45}}planar digraph with a
corresponding rotation system: 
\begin{Verbatim}[commandchars=!@|,fontsize=\small,frame=single,label=Example]
  !gapprompt@gap>| !gapinput@D3 := CompleteMultipartiteDigraph([3, 3]);;|
  !gapprompt@gap>| !gapinput@rot := [[6, 5, 4], [6, 5, 4], [6, 5, 4], [1, 2, 3],|
  !gapprompt@>| !gapinput@           [1, 2, 3], [1, 2, 3]];|
  [ [ 6, 5, 4 ], [ 6, 5, 4 ], [ 6, 5, 4 ], [ 1, 2, 3 ], [ 1, 2, 3 ],
    [ 1, 2, 3 ] ]
  !gapprompt@gap>| !gapinput@FacialWalks(D3, rot);|
  [ [ 1, 4, 2, 6, 3, 5 ], [ 1, 5, 2, 4, 3, 6 ], [ 1, 6, 2, 5, 3, 4 ] ]
\end{Verbatim}
 }

 

\subsection{\textcolor{Chapter }{DigraphLayers}}
\logpage{[ 5, 4, 36 ]}\nobreak
\hyperdef{L}{X870E04307C5F213F}{}
{\noindent\textcolor{FuncColor}{$\triangleright$\enspace\texttt{DigraphLayers({\mdseries\slshape digraph, vertex})\index{DigraphLayers@\texttt{DigraphLayers}}
\label{DigraphLayers}
}\hfill{\scriptsize (operation)}}\\
\textbf{\indent Returns:\ }
A list.



 This operation returns a list \texttt{list} such that \texttt{list[i]} is the list of vertices whose minimum distance from the vertex \mbox{\texttt{\mdseries\slshape vertex}} in \mbox{\texttt{\mdseries\slshape digraph}} is \texttt{i \texttt{\symbol{45}} 1}. Vertex \mbox{\texttt{\mdseries\slshape vertex}} is assumed to be at distance \texttt{0} from itself. 
\begin{Verbatim}[commandchars=!@|,fontsize=\small,frame=single,label=Example]
  !gapprompt@gap>| !gapinput@D := CompleteDigraph(4);;|
  !gapprompt@gap>| !gapinput@DigraphLayers(D, 1);|
  [ [ 1 ], [ 2, 3, 4 ] ]
\end{Verbatim}
 }

 

\subsection{\textcolor{Chapter }{DigraphDegeneracy}}
\logpage{[ 5, 4, 37 ]}\nobreak
\hyperdef{L}{X7B2E42327DA118E0}{}
{\noindent\textcolor{FuncColor}{$\triangleright$\enspace\texttt{DigraphDegeneracy({\mdseries\slshape digraph})\index{DigraphDegeneracy@\texttt{DigraphDegeneracy}}
\label{DigraphDegeneracy}
}\hfill{\scriptsize (attribute)}}\\
\textbf{\indent Returns:\ }
A non\texttt{\symbol{45}}negative integer, or \texttt{fail}.



 If \mbox{\texttt{\mdseries\slshape digraph}} is a symmetric digraph without multiple edges \texttt{\symbol{45}} see \texttt{IsSymmetricDigraph} (\ref{IsSymmetricDigraph}) and \texttt{IsMultiDigraph} (\ref{IsMultiDigraph}) \texttt{\symbol{45}} then this attribute returns the degeneracy of \mbox{\texttt{\mdseries\slshape digraph}}. 

 The degeneracy of a digraph is the least integer \texttt{k} such that every induced of \mbox{\texttt{\mdseries\slshape digraph}} contains a vertex whose number of neighbours (excluding itself) is at most \texttt{k}. Note that this means that loops are ignored.

 If \mbox{\texttt{\mdseries\slshape digraph}} is not symmetric or has multiple edges then this attribute returns \texttt{fail}. 
\begin{Verbatim}[commandchars=!@|,fontsize=\small,frame=single,label=Example]
  !gapprompt@gap>| !gapinput@D := DigraphSymmetricClosure(ChainDigraph(5));;|
  !gapprompt@gap>| !gapinput@DigraphDegeneracy(D);|
  1
  !gapprompt@gap>| !gapinput@D := CompleteDigraph(5);;|
  !gapprompt@gap>| !gapinput@DigraphDegeneracy(D);|
  4
  !gapprompt@gap>| !gapinput@D := Digraph([[1], [2, 4, 5], [3, 4], [2, 3, 4], [2], []]);|
  <immutable digraph with 6 vertices, 10 edges>
  !gapprompt@gap>| !gapinput@DigraphDegeneracy(D);|
  1
  !gapprompt@gap>| !gapinput@D := GeneralisedPetersenGraph(IsMutableDigraph, 10, 3);|
  <mutable digraph with 20 vertices, 60 edges>
  !gapprompt@gap>| !gapinput@DigraphDegeneracy(D);|
  3
\end{Verbatim}
 }

 

\subsection{\textcolor{Chapter }{DigraphDegeneracyOrdering}}
\logpage{[ 5, 4, 38 ]}\nobreak
\hyperdef{L}{X827C2BD17A4547E3}{}
{\noindent\textcolor{FuncColor}{$\triangleright$\enspace\texttt{DigraphDegeneracyOrdering({\mdseries\slshape digraph})\index{DigraphDegeneracyOrdering@\texttt{DigraphDegeneracyOrdering}}
\label{DigraphDegeneracyOrdering}
}\hfill{\scriptsize (attribute)}}\\
\textbf{\indent Returns:\ }
A list of integers, or \texttt{fail}.



 If \mbox{\texttt{\mdseries\slshape digraph}} is a digraph for which \texttt{DigraphDegeneracy(}\mbox{\texttt{\mdseries\slshape digraph}}\texttt{)} is a non\texttt{\symbol{45}}negative integer \texttt{k} \texttt{\symbol{45}} see \texttt{DigraphDegeneracy} (\ref{DigraphDegeneracy}) \texttt{\symbol{45}} then this attribute returns a degeneracy ordering of the
vertices of the vertices of \mbox{\texttt{\mdseries\slshape digraph}}.

 A degeneracy ordering of \mbox{\texttt{\mdseries\slshape digraph}} is a list \texttt{ordering} of the vertices of \mbox{\texttt{\mdseries\slshape digraph}} ordered such that for any position \texttt{i} of the list, the vertex \texttt{ordering[i]} has at most \texttt{k} neighbours in later position of the list.

 If \texttt{DigraphDegeneracy(}\mbox{\texttt{\mdseries\slshape digraph}}\texttt{)} returns \texttt{fail}, then this attribute returns \texttt{fail}. 
\begin{Verbatim}[commandchars=!@|,fontsize=\small,frame=single,label=Example]
  !gapprompt@gap>| !gapinput@D := DigraphSymmetricClosure(ChainDigraph(5));;|
  !gapprompt@gap>| !gapinput@DigraphDegeneracyOrdering(D);|
  [ 5, 4, 3, 2, 1 ]
  !gapprompt@gap>| !gapinput@D := CompleteDigraph(5);;|
  !gapprompt@gap>| !gapinput@DigraphDegeneracyOrdering(D);|
  [ 5, 4, 3, 2, 1 ]
  !gapprompt@gap>| !gapinput@D := Digraph([[1], [2, 4, 5], [3, 4], [2, 3, 4], [2], []]);|
  <immutable digraph with 6 vertices, 10 edges>
  !gapprompt@gap>| !gapinput@DigraphDegeneracyOrdering(D);|
  [ 1, 6, 5, 2, 4, 3 ]
  !gapprompt@gap>| !gapinput@D := GeneralisedPetersenGraph(IsMutableDigraph, 3, 1);|
  <mutable digraph with 6 vertices, 18 edges>
  !gapprompt@gap>| !gapinput@DigraphDegeneracyOrdering(D);|
  [ 6, 5, 4, 1, 3, 2 ]
\end{Verbatim}
 }

 

\subsection{\textcolor{Chapter }{HamiltonianPath}}
\logpage{[ 5, 4, 39 ]}\nobreak
\hyperdef{L}{X863FDFC4839A3B82}{}
{\noindent\textcolor{FuncColor}{$\triangleright$\enspace\texttt{HamiltonianPath({\mdseries\slshape digraph})\index{HamiltonianPath@\texttt{HamiltonianPath}}
\label{HamiltonianPath}
}\hfill{\scriptsize (attribute)}}\\
\textbf{\indent Returns:\ }
A list or \texttt{fail}.



 Returns a Hamiltonian path if one exists, \texttt{fail} if not.

 A \emph{Hamiltonian path} of a digraph with \texttt{n} vertices is directed cycle of length \texttt{n}. If \mbox{\texttt{\mdseries\slshape digraph}} is a digraph that contains a Hamiltonian path, then this function returns one,
described in the form used by \texttt{DigraphAllSimpleCircuits} (\ref{DigraphAllSimpleCircuits}). Note if \mbox{\texttt{\mdseries\slshape digraph}} has 0 or 1 vertices, then \texttt{HamiltonianPath} returns \texttt{[]} or \texttt{[1]}, respectively.

 The method used in this attribute has the same worst case complexity as \texttt{DigraphMonomorphism} (\ref{DigraphMonomorphism}). 
\begin{Verbatim}[commandchars=!@|,fontsize=\small,frame=single,label=Example]
  !gapprompt@gap>| !gapinput@D := Digraph([[]]);|
  <immutable empty digraph with 1 vertex>
  !gapprompt@gap>| !gapinput@HamiltonianPath(D);|
  [ 1 ]
  !gapprompt@gap>| !gapinput@D := Digraph([[2], [1]]);|
  <immutable digraph with 2 vertices, 2 edges>
  !gapprompt@gap>| !gapinput@HamiltonianPath(D);|
  [ 1, 2 ]
  !gapprompt@gap>| !gapinput@D := Digraph([[3], [], [2]]);|
  <immutable digraph with 3 vertices, 2 edges>
  !gapprompt@gap>| !gapinput@HamiltonianPath(D);|
  fail
  !gapprompt@gap>| !gapinput@D := Digraph([[2], [3], [1]]);|
  <immutable digraph with 3 vertices, 3 edges>
  !gapprompt@gap>| !gapinput@HamiltonianPath(D);|
  [ 1, 2, 3 ]
  !gapprompt@gap>| !gapinput@D := GeneralisedPetersenGraph(IsMutableDigraph, 5, 2);|
  <mutable digraph with 10 vertices, 30 edges>
  !gapprompt@gap>| !gapinput@HamiltonianPath(D);|
  fail
\end{Verbatim}
 }

 

\subsection{\textcolor{Chapter }{NrSpanningTrees}}
\logpage{[ 5, 4, 40 ]}\nobreak
\hyperdef{L}{X82F30D5681466BC6}{}
{\noindent\textcolor{FuncColor}{$\triangleright$\enspace\texttt{NrSpanningTrees({\mdseries\slshape digraph})\index{NrSpanningTrees@\texttt{NrSpanningTrees}}
\label{NrSpanningTrees}
}\hfill{\scriptsize (attribute)}}\\
\textbf{\indent Returns:\ }
An integer.



 Returns the number of spanning trees of the symmetric digraph \mbox{\texttt{\mdseries\slshape digraph}}. \texttt{NrSpanningTrees} will return an error if \mbox{\texttt{\mdseries\slshape digraph}} is not a symmetric digraph. 

 See \texttt{IsSymmetricDigraph} (\ref{IsSymmetricDigraph}) and \texttt{IsUndirectedSpanningTree} (\ref{IsUndirectedSpanningTree}) for more information. 
\begin{Verbatim}[commandchars=!@|,fontsize=\small,frame=single,label=Example]
  !gapprompt@gap>| !gapinput@D := CompleteDigraph(5);|
  <immutable complete digraph with 5 vertices>
  !gapprompt@gap>| !gapinput@NrSpanningTrees(D);|
  125
  !gapprompt@gap>| !gapinput@D := DigraphSymmetricClosure(CycleDigraph(24));;|
  !gapprompt@gap>| !gapinput@NrSpanningTrees(D);|
  24
  !gapprompt@gap>| !gapinput@NrSpanningTrees(EmptyDigraph(0));|
  0
  !gapprompt@gap>| !gapinput@D := GeneralisedPetersenGraph(IsMutableDigraph, 9, 2);|
  <mutable digraph with 18 vertices, 54 edges>
  !gapprompt@gap>| !gapinput@NrSpanningTrees(D);|
  1134225
\end{Verbatim}
 }

 

\subsection{\textcolor{Chapter }{DigraphDijkstra (for a source and target)}}
\logpage{[ 5, 4, 41 ]}\nobreak
\hyperdef{L}{X79352A8286D1D8F6}{}
{\noindent\textcolor{FuncColor}{$\triangleright$\enspace\texttt{DigraphDijkstra({\mdseries\slshape digraph, source, target})\index{DigraphDijkstra@\texttt{DigraphDijkstra}!for a source and target}
\label{DigraphDijkstra:for a source and target}
}\hfill{\scriptsize (operation)}}\\
\noindent\textcolor{FuncColor}{$\triangleright$\enspace\texttt{DigraphDijkstra({\mdseries\slshape digraph, source})\index{DigraphDijkstra@\texttt{DigraphDijkstra}!for a source}
\label{DigraphDijkstra:for a source}
}\hfill{\scriptsize (operation)}}\\
\textbf{\indent Returns:\ }
Two lists.



 If \mbox{\texttt{\mdseries\slshape digraph}} is a digraph and \mbox{\texttt{\mdseries\slshape source}} and \mbox{\texttt{\mdseries\slshape target}} are vertices of \mbox{\texttt{\mdseries\slshape digraph}}, then \texttt{DigraphDijkstra} calculates the length of the shortest path from \mbox{\texttt{\mdseries\slshape source}} to \mbox{\texttt{\mdseries\slshape target}} and returns two lists. Each element of the first list is the distance of the
corresponding element from \mbox{\texttt{\mdseries\slshape source}}. If a vertex was not visited in the process of calculating the shortest
distance to \mbox{\texttt{\mdseries\slshape target}} or if there is no path connecting that vertex with \mbox{\texttt{\mdseries\slshape source}}, then the corresponding distance is \texttt{infinity}. Each element of the second list gives the previous vertex in the shortest
path from \mbox{\texttt{\mdseries\slshape source}} to the corresponding vertex. For \mbox{\texttt{\mdseries\slshape source}} and for any vertices that remained unvisited this will be \texttt{\texttt{\symbol{45}}1}.

 If the optional second argument \mbox{\texttt{\mdseries\slshape target}} is not present, then \texttt{DigraphDijkstra} returns the shortest path from \mbox{\texttt{\mdseries\slshape source}} to every vertex that is reachable from \mbox{\texttt{\mdseries\slshape source}}. 
\begin{Verbatim}[commandchars=!@|,fontsize=\small,frame=single,label=Example]
  !gapprompt@gap>| !gapinput@mat := [[0, 1, 1], [0, 0, 1], [0, 0, 0]];|
  [ [ 0, 1, 1 ], [ 0, 0, 1 ], [ 0, 0, 0 ] ]
  !gapprompt@gap>| !gapinput@D := DigraphByAdjacencyMatrix(mat);|
  <immutable digraph with 3 vertices, 3 edges>
  !gapprompt@gap>| !gapinput@DigraphDijkstra(D, 2, 3);|
  [ [ infinity, 0, 1 ], [ -1, -1, 2 ] ]
  !gapprompt@gap>| !gapinput@DigraphDijkstra(D, 1, 3);|
  [ [ 0, 1, 1 ], [ -1, 1, 1 ] ]
  !gapprompt@gap>| !gapinput@DigraphDijkstra(D, 1, 2);|
  [ [ 0, 1, 1 ], [ -1, 1, 1 ] ]
\end{Verbatim}
 }

 

\subsection{\textcolor{Chapter }{DigraphVertexConnectivity}}
\logpage{[ 5, 4, 42 ]}\nobreak
\hyperdef{L}{X87ADACAB7A81805A}{}
{\noindent\textcolor{FuncColor}{$\triangleright$\enspace\texttt{DigraphVertexConnectivity({\mdseries\slshape digraph})\index{DigraphVertexConnectivity@\texttt{DigraphVertexConnectivity}}
\label{DigraphVertexConnectivity}
}\hfill{\scriptsize (attribute)}}\\
\textbf{\indent Returns:\ }
An non\texttt{\symbol{45}}negative integer.



 This function returns the vertex connectivity of the digraph \mbox{\texttt{\mdseries\slshape digraph}}. This is also sometimes called the weak vertex connectivity. 

 The vertex connectivity of a connected digraph (in the sense of \texttt{IsConnectedDigraph} (\ref{IsConnectedDigraph})) is the largest number $k$ such that: 
\begin{itemize}
\item  the digraph has at least $k + 1$ vertices, and
\item  the digraph remains connected after the removal of any set of at most $k - 1$ vertices.
\end{itemize}
 If the digraph is not connected, then its vertex connectivity is 0. For a
non\texttt{\symbol{45}}empty digraph whose symmetric closure is not complete,
the vertex connectivity is equal to the size of the smallest set of vertices
whose removal would disconnect the digraph. 

 Vertex connectivity of small digraphs may be counterintuitive with respect to
the notions of connectivity, as defined by \texttt{IsConnectedDigraph} (\ref{IsConnectedDigraph}), and biconnecivity, as defined by \texttt{IsBiconnectedDigraph} (\ref{IsBiconnectedDigraph}). Namely 
\begin{itemize}
\item  the empty digraph is connected, but its vertex connectivity is 0; 
\item  the singleton digraph is connected, but its vertex connectivity is 0;
\item  the 2\texttt{\symbol{45}}cycle digraph is biconnected, but its vertex
connectivity is only 1.
\end{itemize}
 However, for a digraph with at least 3 vertices, having vertex connectivity
greater than or equal to 1 is equivalent to being connected, and having vertex
connectivity greater than or equal to 2 is equivalent to being biconnected. 

 The algorithm makes $n - d - 1 + d \cdot (d - 1) / 2$ calls to the \texttt{DigraphMaximumFlow} (\ref{DigraphMaximumFlow}) maximum flow algorithm, where $n$ is the number of vertices of \mbox{\texttt{\mdseries\slshape digraph}}, and $d$ is the minimum degree of the symmetric closure of \mbox{\texttt{\mdseries\slshape digraph}}, see \texttt{DigraphSymmetricClosure} (\ref{DigraphSymmetricClosure}). 
\begin{Verbatim}[commandchars=!@|,fontsize=\small,frame=single,label=Example]
  !gapprompt@gap>| !gapinput@D := CompleteBipartiteDigraph(4, 5);|
  <immutable complete bipartite digraph with bicomponent sizes 4 and 5>
  !gapprompt@gap>| !gapinput@DigraphVertexConnectivity(D);|
  4
  !gapprompt@gap>| !gapinput@DigraphVertexConnectivity(PancakeGraph(4));|
  3
  !gapprompt@gap>| !gapinput@D := Digraph([[2, 4, 5], [1, 4], [4, 7], [1, 2, 3, 5, 6, 7],|
  !gapprompt@>| !gapinput@                 [1, 4], [4, 7], [3, 4, 6]]);|
  <immutable digraph with 7 vertices, 20 edges>
  !gapprompt@gap>| !gapinput@DigraphVertexConnectivity(D);|
  1
  !gapprompt@gap>| !gapinput@J := JohnsonDigraph(9, 2);|
  <immutable symmetric digraph with 36 vertices, 504 edges>
  !gapprompt@gap>| !gapinput@DigraphVertexConnectivity(J);|
  14
  !gapprompt@gap>| !gapinput@DigraphVertexConnectivity(Digraph([]));|
  0
  !gapprompt@gap>| !gapinput@DigraphVertexConnectivity(Digraph([[1]]));|
  0
  !gapprompt@gap>| !gapinput@DigraphVertexConnectivity(CycleDigraph(2));|
  1
  !gapprompt@gap>| !gapinput@DigraphVertexConnectivity(CompleteDigraph(5));|
  4
\end{Verbatim}
 }

 

\subsection{\textcolor{Chapter }{DigraphCycleBasis}}
\logpage{[ 5, 4, 43 ]}\nobreak
\hyperdef{L}{X794CD6037D4CF58C}{}
{\noindent\textcolor{FuncColor}{$\triangleright$\enspace\texttt{DigraphCycleBasis({\mdseries\slshape digraph})\index{DigraphCycleBasis@\texttt{DigraphCycleBasis}}
\label{DigraphCycleBasis}
}\hfill{\scriptsize (operation)}}\\
\textbf{\indent Returns:\ }
A list and a list of GF(2) vectors



 If \mbox{\texttt{\mdseries\slshape digraph}} is a symmetric loopless digraph with no multiple edges, then \texttt{DigraphCycleBasis} calculates the \emph{fundamental cycle basis} of \mbox{\texttt{\mdseries\slshape digraph}} and returns a list of edges and a list of a basis vectors. If \mbox{\texttt{\mdseries\slshape digraph}} contains a loop, this function will return an error. If \mbox{\texttt{\mdseries\slshape digraph}} is a multi\texttt{\symbol{45}}digraph, then multiple edges incident to the
same vertices are treated as a single edge by this function. See \texttt{IsSymmetricDigraph} (\ref{IsSymmetricDigraph}), \texttt{DigraphHasLoops} (\ref{DigraphHasLoops}), and \texttt{IsMultiDigraph} (\ref{IsMultiDigraph}). The first list returned contains the out\texttt{\symbol{45}}neighours that
provide the ordering on the edges with the symmetric edges appearing only
once; these are the same as \emph{OutNeighbours(MaximalAntiSymmetricSubdigraph(G))}. See \texttt{MaximalAntiSymmetricSubdigraph} (\ref{MaximalAntiSymmetricSubdigraph}) \texttt{OutNeighbours} (\ref{OutNeighbours}). The second list returned contains the basis vectors of the cycle space of
the digraph. These vectors belongs to $(GF(2))^m$ where $m$ is the length of the list of edges.

 A graph is \emph{eulerian} if every vertex has an even degree. A \emph{cycle space} of a graph is a subspace of $(GF(2))^m$ representing the set of all \emph{eulerian} subgraphs without isolated vertices. An eulerian subgraph without isolated
vertices can be uniquely identified using the edges that is contains.
Therefore, given some ordering on the edges of the graph, a subgraph can be
represented by a vector in $(GF(2))^m$ where the $i$\texttt{\symbol{45}}th entry is $1$ if the $i$\texttt{\symbol{45}}th edge is in the subgraph and $0$ otherwise. In this function, the ordering of the edges is returned as the
first list which corresponds to \texttt{OutNeighbours(MaximalAntiSymmetricSubdigraph(G))}. See \texttt{MaximalAntiSymmetricSubdigraph} (\ref{MaximalAntiSymmetricSubdigraph}) and \texttt{OutNeighbours} (\ref{OutNeighbours}). The first basis vector for the complete digraph with 4 vertices shown below
represents the edges \texttt{[1, 2]}, \texttt{[1, 3]} and \texttt{[2, 3]} i.e. cycle subgraph between the vertices \texttt{1}, \texttt{2} and \texttt{3} The cycle space is closed under the symmetric difference of the edges of the
graph. This nicely corresponds to the addition of vectors in $(GF(2))^m$ which makes the vector space formulation of the cycle space very natural.

 A \emph{cycle basis} is a basis of the cycle space. A \emph{fundamental cycle basis} is a special kind of basis of the cycle space where there is a specific
spanning tree (forest, when the graph is disconnected) of the graph and each
basis corresponds to the unique cycle created by adding an edge outside the
spanning tree of the graph to the spanning tree. In this function, the
spanning forest is rooted in the vertex with the smallest label for each
connected component of the graph. 

 The fundamental cycle basis is unique up to reordering of the basis vectors.
The number of basis vectors in the fundamental cycle basis is $m - n + c$, where $m$ is the number of edges, $n$ is the number of vertices, and $c$ is the number of connected components. See \texttt{DigraphConnectedComponents} (\ref{DigraphConnectedComponents}).

 This function performs a depth first traversal of the input digraph with
complexity $O(m + n)$ and the complexity of the computation of the basis is $O(m^2)$ where $m$ is the number of edges in the input digraph. 
\begin{Verbatim}[commandchars=!@|,fontsize=\small,frame=single,label=Example]
  !gapprompt@gap>| !gapinput@D := CycleGraph(4);|
  <immutable symmetric digraph with 4 vertices, 8 edges>
  !gapprompt@gap>| !gapinput@res := DigraphCycleBasis(D);|
  [ [ [ 2, 4 ], [ 3 ], [ 4 ], [  ] ], [ <a GF2 vector of length 4> ] ]
  !gapprompt@gap>| !gapinput@List(res[2][1]);|
  [ Z(2)^0, Z(2)^0, Z(2)^0, Z(2)^0 ]
  !gapprompt@gap>| !gapinput@D := CompleteDigraph(4);                                     |
  <immutable complete digraph with 4 vertices>
  !gapprompt@gap>| !gapinput@res := DigraphCycleBasis(D);|
  [ [ [ 2, 3, 4 ], [ 3, 4 ], [ 4 ], [  ] ], 
    [ <a GF2 vector of length 6>, <a GF2 vector of length 6>, 
        <a GF2 vector of length 6> ] ]
  !gapprompt@gap>| !gapinput@List(res[2][1]);            |
  [ Z(2)^0, 0*Z(2), Z(2)^0, 0*Z(2), Z(2)^0, 0*Z(2) ]
  !gapprompt@gap>| !gapinput@List(res[2][2]);|
  [ 0*Z(2), Z(2)^0, Z(2)^0, 0*Z(2), 0*Z(2), Z(2)^0 ]
  !gapprompt@gap>| !gapinput@List(res[2][3]);|
  [ Z(2)^0, Z(2)^0, 0*Z(2), Z(2)^0, 0*Z(2), 0*Z(2) ]
\end{Verbatim}
 }

 

\subsection{\textcolor{Chapter }{DigraphIsKing}}
\logpage{[ 5, 4, 44 ]}\nobreak
\hyperdef{L}{X835AAF9085BC9D84}{}
{\noindent\textcolor{FuncColor}{$\triangleright$\enspace\texttt{DigraphIsKing({\mdseries\slshape D, v, k})\index{DigraphIsKing@\texttt{DigraphIsKing}}
\label{DigraphIsKing}
}\hfill{\scriptsize (operation)}}\\
\textbf{\indent Returns:\ }
\texttt{true} or \texttt{false}.



 If \mbox{\texttt{\mdseries\slshape D}} is a tournament and \mbox{\texttt{\mdseries\slshape v}} is a vertex in the tournament, then this operation returns \texttt{true} if every other vertex of \mbox{\texttt{\mdseries\slshape D}} is reachable from \mbox{\texttt{\mdseries\slshape v}} by a path of length at most \mbox{\texttt{\mdseries\slshape k}}. Otherwise, an error is given. If \texttt{true} is returned, then the vertex, \mbox{\texttt{\mdseries\slshape v}}, is a \mbox{\texttt{\mdseries\slshape k}}\texttt{\symbol{45}}king. 

 
\begin{Verbatim}[commandchars=!@|,fontsize=\small,frame=single,label=Example]
  !gapprompt@gap>| !gapinput@gr := Digraph([[2, 3, 4], [3, 5], [5], [2, 3], [1, 4]]);|
  <immutable digraph with 5 vertices, 10 edges>
  !gapprompt@gap>| !gapinput@DigraphIsKing(gr, 2, 2);|
  true
  !gapprompt@gap>| !gapinput@DigraphIsKing(gr, 3, 2);|
  false
  !gapprompt@gap>| !gapinput@OutNeighboursOfVertex(gr, 3);|
  [ 5 ]
  !gapprompt@gap>| !gapinput@OutNeighboursOfVertex(gr, 5);|
  [ 1, 4 ]
  !gapprompt@gap>| !gapinput@Union(last, last2);|
  [ 1, 4, 5 ]
  !gapprompt@gap>| !gapinput@DigraphIsKing(gr, 3, 4);|
  true
\end{Verbatim}
 }

 

\subsection{\textcolor{Chapter }{DigraphKings}}
\logpage{[ 5, 4, 45 ]}\nobreak
\hyperdef{L}{X7A17B9B67AC50561}{}
{\noindent\textcolor{FuncColor}{$\triangleright$\enspace\texttt{DigraphKings({\mdseries\slshape D, n})\index{DigraphKings@\texttt{DigraphKings}}
\label{DigraphKings}
}\hfill{\scriptsize (operation)}}\\
\textbf{\indent Returns:\ }
A list.



 If \mbox{\texttt{\mdseries\slshape D}} is a tournament, then this operation returns a list of the \mbox{\texttt{\mdseries\slshape n}}\texttt{\symbol{45}}kings in the tournament (see \texttt{DigraphIsKing} (\ref{DigraphIsKing})).

 If \mbox{\texttt{\mdseries\slshape D}} is not tournament, then an error is given (see \texttt{IsTournament} (\ref{IsTournament})). If the tournament contains a source, then the source is the only
2\texttt{\symbol{45}}king (see \texttt{DigraphSources} (\ref{DigraphSources})). The number of 2\texttt{\symbol{45}}kings in a tournament without a source
is at least three. A tournament cannot have exactly two
2\texttt{\symbol{45}}kings. 
\begin{Verbatim}[commandchars=!@|,fontsize=\small,frame=single,label=Example]
  !gapprompt@gap>| !gapinput@gr := Digraph([[2, 3, 4], [3, 5], [5], [2, 3], [1, 4]]);|
  <immutable digraph with 5 vertices, 10 edges>
  !gapprompt@gap>| !gapinput@DigraphKings(gr, 2);|
  [ 1, 2, 5 ]
  !gapprompt@gap>| !gapinput@gr := Digraph([[2, 3, 4, 5], [3, 5], [5], [2, 3], [4]]);|
  <immutable digraph with 5 vertices, 10 edges>
  !gapprompt@gap>| !gapinput@DigraphSources(gr);|
  [ 1 ]
  !gapprompt@gap>| !gapinput@DigraphKings(gr, 2);|
  [ 1 ]
\end{Verbatim}
 }

 }

 
\section{\textcolor{Chapter }{Cayley graphs of groups}}\logpage{[ 5, 5, 0 ]}
\hyperdef{L}{X82F900777D677F55}{}
{
 

\subsection{\textcolor{Chapter }{GroupOfCayleyDigraph}}
\logpage{[ 5, 5, 1 ]}\nobreak
\hyperdef{L}{X7A000B1D7CCF7093}{}
{\noindent\textcolor{FuncColor}{$\triangleright$\enspace\texttt{GroupOfCayleyDigraph({\mdseries\slshape digraph})\index{GroupOfCayleyDigraph@\texttt{GroupOfCayleyDigraph}}
\label{GroupOfCayleyDigraph}
}\hfill{\scriptsize (attribute)}}\\
\noindent\textcolor{FuncColor}{$\triangleright$\enspace\texttt{SemigroupOfCayleyDigraph({\mdseries\slshape digraph})\index{SemigroupOfCayleyDigraph@\texttt{SemigroupOfCayleyDigraph}}
\label{SemigroupOfCayleyDigraph}
}\hfill{\scriptsize (attribute)}}\\
\textbf{\indent Returns:\ }
A group or semigroup.



 If \mbox{\texttt{\mdseries\slshape digraph}} is an immutable Cayley graph of a group \texttt{G} and \mbox{\texttt{\mdseries\slshape digraph}} belongs to the category \texttt{IsCayleyDigraph} (\ref{IsCayleyDigraph}), then \texttt{GroupOfCayleyDigraph} returns \texttt{G}. 

 If \mbox{\texttt{\mdseries\slshape digraph}} is a Cayley graph of a semigroup \texttt{S} and \mbox{\texttt{\mdseries\slshape digraph}} belongs to the category \texttt{IsCayleyDigraph} (\ref{IsCayleyDigraph}), then \texttt{SemigroupOfCayleyDigraph} returns \texttt{S}. 

 See also \texttt{GeneratorsOfCayleyDigraph} (\ref{GeneratorsOfCayleyDigraph}). 
\begin{Verbatim}[commandchars=!@|,fontsize=\small,frame=single,label=Example]
  !gapprompt@gap>| !gapinput@G := DihedralGroup(IsPermGroup, 8);|
  Group([ (1,2,3,4), (2,4) ])
  !gapprompt@gap>| !gapinput@digraph := CayleyDigraph(G);|
  <immutable digraph with 8 vertices, 16 edges>
  !gapprompt@gap>| !gapinput@GroupOfCayleyDigraph(digraph) = G;|
  true
\end{Verbatim}
 }

 

\subsection{\textcolor{Chapter }{GeneratorsOfCayleyDigraph}}
\logpage{[ 5, 5, 2 ]}\nobreak
\hyperdef{L}{X8528455987D7D2BF}{}
{\noindent\textcolor{FuncColor}{$\triangleright$\enspace\texttt{GeneratorsOfCayleyDigraph({\mdseries\slshape digraph})\index{GeneratorsOfCayleyDigraph@\texttt{GeneratorsOfCayleyDigraph}}
\label{GeneratorsOfCayleyDigraph}
}\hfill{\scriptsize (attribute)}}\\
\textbf{\indent Returns:\ }
A list of generators.



 If \mbox{\texttt{\mdseries\slshape digraph}} is an immutable Cayley graph of a group or semigroup with respect to a set of
generators \texttt{gens} and \mbox{\texttt{\mdseries\slshape digraph}} belongs to the category \texttt{IsCayleyDigraph} (\ref{IsCayleyDigraph}), then \texttt{GeneratorsOfCayleyDigraph} return the list of generators \texttt{gens} over which \mbox{\texttt{\mdseries\slshape digraph}} is defined. 

 See also \texttt{GroupOfCayleyDigraph} (\ref{GroupOfCayleyDigraph}) or \texttt{SemigroupOfCayleyDigraph} (\ref{SemigroupOfCayleyDigraph}). 
\begin{Verbatim}[commandchars=!@|,fontsize=\small,frame=single,label=Example]
  !gapprompt@gap>| !gapinput@G := DihedralGroup(IsPermGroup, 8);|
  Group([ (1,2,3,4), (2,4) ])
  !gapprompt@gap>| !gapinput@digraph := CayleyDigraph(G);|
  <immutable digraph with 8 vertices, 16 edges>
  !gapprompt@gap>| !gapinput@GeneratorsOfCayleyDigraph(digraph) = GeneratorsOfGroup(G);|
  true
  !gapprompt@gap>| !gapinput@digraph := CayleyDigraph(G, [()]);|
  <immutable digraph with 8 vertices, 8 edges>
  !gapprompt@gap>| !gapinput@GeneratorsOfCayleyDigraph(digraph) = [()];|
  true
\end{Verbatim}
 }

 }

 
\section{\textcolor{Chapter }{Associated semigroups}}\logpage{[ 5, 6, 0 ]}
\hyperdef{L}{X790FD6647ECCAE3C}{}
{
 

\subsection{\textcolor{Chapter }{AsSemigroup (for a filter and a digraph)}}
\logpage{[ 5, 6, 1 ]}\nobreak
\hyperdef{L}{X87D5C60D7B0C1309}{}
{\noindent\textcolor{FuncColor}{$\triangleright$\enspace\texttt{AsSemigroup({\mdseries\slshape filt, digraph})\index{AsSemigroup@\texttt{AsSemigroup}!for a filter and a digraph}
\label{AsSemigroup:for a filter and a digraph}
}\hfill{\scriptsize (operation)}}\\
\noindent\textcolor{FuncColor}{$\triangleright$\enspace\texttt{AsMonoid({\mdseries\slshape filt, digraph})\index{AsMonoid@\texttt{AsMonoid}}
\label{AsMonoid}
}\hfill{\scriptsize (operation)}}\\
\textbf{\indent Returns:\ }
A semilattice of partial perms.



 The operation \texttt{AsSemigroup} requires that \mbox{\texttt{\mdseries\slshape filt}} be equal to \texttt{IsPartialPermSemigroup} (\textbf{Reference: IsPartialPermSemigroup}). If \mbox{\texttt{\mdseries\slshape digraph}} is a \texttt{IsJoinSemilatticeDigraph} (\ref{IsJoinSemilatticeDigraph}) or \texttt{IsLatticeDigraph} (\ref{IsLatticeDigraph}) then \texttt{AsSemigroup} returns a semigroup of partial perms which is isomorphic to the semigroup
whose elements are the vertices of \mbox{\texttt{\mdseries\slshape digraph}} with the binary operation \texttt{PartialOrderDigraphJoinOfVertices} (\ref{PartialOrderDigraphJoinOfVertices}). If \mbox{\texttt{\mdseries\slshape digraph}} satisfies \texttt{IsMeetSemilatticeDigraph} (\ref{IsMeetSemilatticeDigraph}) but not \texttt{IsJoinSemilatticeDigraph} (\ref{IsJoinSemilatticeDigraph}) then \texttt{AsSemigroup} returns a semigroup of partial perms which is isomorphic to the semigroup
whose elements are the vertices of \mbox{\texttt{\mdseries\slshape digraph}} with the binary operation \texttt{PartialOrderDigraphMeetOfVertices} (\ref{PartialOrderDigraphMeetOfVertices}). 

 The operation \texttt{AsMonoid} behaves similarly to \texttt{AsSemigroup} except that \mbox{\texttt{\mdseries\slshape filt}} may also be equal to \texttt{IsPartialPermMonoid} (\textbf{Reference: IsPartialPermMonoid}), \mbox{\texttt{\mdseries\slshape digraph}} must satisfy \texttt{IsLatticeDigraph} (\ref{IsLatticeDigraph}), and the output satisfies \texttt{IsMonoid} (\textbf{Reference: IsMonoid}). 

 The output of both of these operations is guaranteed to be of minimal degree
(see \texttt{DegreeOfPartialPermSemigroup} (\textbf{Reference: DegreeOfPartialPermSemigroup})). Furthermore the \texttt{GeneratorsOfSemigroup} (\textbf{Reference: GeneratorsOfSemigroup}) of the output is guaranteed to be the unique generating set of minimal size. 
\begin{Verbatim}[commandchars=!@|,fontsize=\small,frame=single,label=Example]
  !gapprompt@gap>| !gapinput@di := Digraph([[1], [1, 2], [1, 3], [1, 4], [1, 2, 3, 5]]);|
  <immutable digraph with 5 vertices, 11 edges>
  !gapprompt@gap>| !gapinput@S := AsSemigroup(IsPartialPermSemigroup, di);|
  <partial perm semigroup of rank 3 with 4 generators>
  !gapprompt@gap>| !gapinput@ForAll(Elements(S), IsIdempotent);|
  true
  !gapprompt@gap>| !gapinput@IsInverseSemigroup(S);|
  true
  !gapprompt@gap>| !gapinput@Size(S);|
  5
  !gapprompt@gap>| !gapinput@di := Digraph([[1], [1, 2], [1, 2, 3]]);|
  <immutable digraph with 3 vertices, 6 edges>
  !gapprompt@gap>| !gapinput@M := AsMonoid(IsPartialPermMonoid, di);|
  <partial perm monoid of rank 2 with 3 generators>
  !gapprompt@gap>| !gapinput@Size(M);|
  3
\end{Verbatim}
 }

 

\subsection{\textcolor{Chapter }{AsSemigroup (for a filter, semilattice digraph, and two lists)}}
\logpage{[ 5, 6, 2 ]}\nobreak
\hyperdef{L}{X7C6D5EC27C51066B}{}
{\noindent\textcolor{FuncColor}{$\triangleright$\enspace\texttt{AsSemigroup({\mdseries\slshape filt, Y, gps, homs})\index{AsSemigroup@\texttt{AsSemigroup}!for a filter, semilattice digraph, and two lists}
\label{AsSemigroup:for a filter, semilattice digraph, and two lists}
}\hfill{\scriptsize (operation)}}\\
\textbf{\indent Returns:\ }
 A Clifford semigroup of partial perms. 



 The operation \texttt{AsSemigroup} requires that \mbox{\texttt{\mdseries\slshape filt}} be equal to \texttt{IsPartialPermSemigroup} (\textbf{Reference: IsPartialPermSemigroup}). If \mbox{\texttt{\mdseries\slshape Y}} is a \texttt{IsJoinSemilatticeDigraph} (\ref{IsJoinSemilatticeDigraph}) or \texttt{IsMeetSemilatticeDigraph} (\ref{IsMeetSemilatticeDigraph}), \mbox{\texttt{\mdseries\slshape gps}} is a list of groups corresponding to each vertex, and \mbox{\texttt{\mdseries\slshape homs}} is a list containing for each edge \texttt{(i, j)} in the transitive reduction of \mbox{\texttt{\mdseries\slshape digraph}} a triple \texttt{[i, j, hom]} where \texttt{hom} is a group homomorphism from \texttt{gps[i]} to \texttt{gps[j]}, and the diagram of homomorphisms commutes, then \texttt{AsSemigroup} returns a semigroup of partial perms which is isomorphic to the strong
semilattice of groups $S[Y; gps; homs]$. 
\begin{Verbatim}[commandchars=!@|,fontsize=\small,frame=single,label=Example]
  !gapprompt@gap>| !gapinput@G1 := AlternatingGroup(4);;|
  !gapprompt@gap>| !gapinput@G2 := SymmetricGroup(2);;|
  !gapprompt@gap>| !gapinput@G3 := SymmetricGroup(3);;|
  !gapprompt@gap>| !gapinput@gr := Digraph([[1, 3], [2, 3], [3]]);;|
  !gapprompt@gap>| !gapinput@sgn := function(x)|
  !gapprompt@>| !gapinput@if SignPerm(x) = 1 then|
  !gapprompt@>| !gapinput@return ();|
  !gapprompt@>| !gapinput@fi;|
  !gapprompt@>| !gapinput@return (1, 2);|
  !gapprompt@>| !gapinput@end;;|
  !gapprompt@gap>| !gapinput@hom13 := GroupHomomorphismByFunction(G1, G3, sgn);;|
  !gapprompt@gap>| !gapinput@hom23 := GroupHomomorphismByFunction(G2, G3, sgn);;|
  !gapprompt@gap>| !gapinput@T := AsSemigroup(IsPartialPermSemigroup,|
  !gapprompt@>| !gapinput@gr,|
  !gapprompt@>| !gapinput@[G1, G2, G3], [[1, 3, hom13], [2, 3, hom23]]);;|
  !gapprompt@gap>| !gapinput@Size(T);|
  20
  !gapprompt@gap>| !gapinput@D := GreensDClasses(T);;|
  !gapprompt@gap>| !gapinput@List(D, Size);|
  [ 6, 12, 2 ]
\end{Verbatim}
 }

 }

 
\section{\textcolor{Chapter }{Planarity}}\logpage{[ 5, 7, 0 ]}
\hyperdef{L}{X7E2305528492DDC0}{}
{
 

\subsection{\textcolor{Chapter }{KuratowskiPlanarSubdigraph}}
\logpage{[ 5, 7, 1 ]}\nobreak
\hyperdef{L}{X7DC478637E8C190D}{}
{\noindent\textcolor{FuncColor}{$\triangleright$\enspace\texttt{KuratowskiPlanarSubdigraph({\mdseries\slshape digraph})\index{KuratowskiPlanarSubdigraph@\texttt{KuratowskiPlanarSubdigraph}}
\label{KuratowskiPlanarSubdigraph}
}\hfill{\scriptsize (attribute)}}\\
\textbf{\indent Returns:\ }
A list or \texttt{fail}.



 \texttt{KuratowskiPlanarSubdigraph} returns the immutable list of lists of out\texttt{\symbol{45}}neighbours of an
induced subdigraph (excluding multiple edges and loops) of the digraph \mbox{\texttt{\mdseries\slshape digraph}} that witnesses the fact that \mbox{\texttt{\mdseries\slshape digraph}} is not planar, or \texttt{fail} if \mbox{\texttt{\mdseries\slshape digraph}} is planar. In other words, \texttt{KuratowskiPlanarSubdigraph} returns the out\texttt{\symbol{45}}neighbours of a subdigraph of \mbox{\texttt{\mdseries\slshape digraph}} that is homeomorphic to the complete graph with \texttt{5} vertices, or to the complete bipartite graph with vertex sets of sizes \texttt{3} and \texttt{3}. 

 The directions and multiplicities of any edges in \mbox{\texttt{\mdseries\slshape digraph}} are ignored when considering whether or not \mbox{\texttt{\mdseries\slshape digraph}} is planar. 

 See also \texttt{IsPlanarDigraph} (\ref{IsPlanarDigraph}) and \texttt{SubdigraphHomeomorphicToK33} (\ref{SubdigraphHomeomorphicToK33}). 

 This method uses the reference implementation in \href{https://github.com/graph-algorithms/edge-addition-planarity-suite} {edge-addition-planarity-suite}  by John Boyer of the algorithms described in \cite{BM06}. 
\begin{Verbatim}[commandchars=!@|,fontsize=\small,frame=single,label=Example]
  !gapprompt@gap>| !gapinput@D := Digraph([[3, 5, 10], [8, 9, 10], [1, 4], [3, 6], |
  !gapprompt@>| !gapinput@[1, 7, 11], [4, 7], [6, 8], [2, 7], [2, 11], [1, 2], [5, 9]]);|
  <immutable digraph with 11 vertices, 25 edges>
  !gapprompt@gap>| !gapinput@KuratowskiPlanarSubdigraph(D);|
  fail
  !gapprompt@gap>| !gapinput@D := Digraph([[2, 4, 7, 9, 10], [1, 3, 4, 6, 9, 10], [6, 10], |
  !gapprompt@>| !gapinput@[2, 5, 8, 9], [1, 2, 3, 4, 6, 7, 9, 10], [3, 4, 5, 7, 9, 10], |
  !gapprompt@>| !gapinput@[3, 4, 5, 6, 9, 10], [3, 4, 5, 7, 9], [2, 3, 5, 6, 7, 8], [3, 5]]);|
  <immutable digraph with 10 vertices, 50 edges>
  !gapprompt@gap>| !gapinput@IsPlanarDigraph(D);|
  false
  !gapprompt@gap>| !gapinput@KuratowskiPlanarSubdigraph(D);|
  [ [ 2, 9, 7 ], [ 1, 3 ], [ 6 ], [ 5, 9 ], [ 6, 4 ], [ 3, 5 ], [ 4 ], 
    [ 7, 9, 3 ], [ 8 ], [  ] ]
  !gapprompt@gap>| !gapinput@D := Digraph(IsMutableDigraph, [[3, 5, 10], [8, 9, 10], [1, 4],|
  !gapprompt@>| !gapinput@[3, 6], [1, 7, 11], [4, 7], [6, 8], [2, 7], [2, 11], [1, 2], [5, 9]]);|
  <mutable digraph with 11 vertices, 25 edges>
  !gapprompt@gap>| !gapinput@KuratowskiPlanarSubdigraph(D);|
  fail
  !gapprompt@gap>| !gapinput@D := Digraph(IsMutableDigraph, [[2, 4, 7, 9, 10],|
  !gapprompt@>| !gapinput@[1, 3, 4, 6, 9, 10], [6, 10], [2, 5, 8, 9],|
  !gapprompt@>| !gapinput@[1, 2, 3, 4, 6, 7, 9, 10], [3, 4, 5, 7, 9, 10],|
  !gapprompt@>| !gapinput@[3, 4, 5, 6, 9, 10], [3, 4, 5, 7, 9], [2, 3, 5, 6, 7, 8], [3, 5]]);|
  <mutable digraph with 10 vertices, 50 edges>
  !gapprompt@gap>| !gapinput@IsPlanarDigraph(D);|
  false
  !gapprompt@gap>| !gapinput@KuratowskiPlanarSubdigraph(D);|
  [ [ 2, 9, 7 ], [ 1, 3 ], [ 6 ], [ 5, 9 ], [ 6, 4 ], [ 3, 5 ], [ 4 ], 
    [ 7, 9, 3 ], [ 8 ], [  ] ]
\end{Verbatim}
 }

 

\subsection{\textcolor{Chapter }{KuratowskiOuterPlanarSubdigraph}}
\logpage{[ 5, 7, 2 ]}\nobreak
\hyperdef{L}{X78E8F09A8286501B}{}
{\noindent\textcolor{FuncColor}{$\triangleright$\enspace\texttt{KuratowskiOuterPlanarSubdigraph({\mdseries\slshape digraph})\index{KuratowskiOuterPlanarSubdigraph@\texttt{KuratowskiOuterPlanarSubdigraph}}
\label{KuratowskiOuterPlanarSubdigraph}
}\hfill{\scriptsize (attribute)}}\\
\textbf{\indent Returns:\ }
A list or \texttt{fail}.



 \texttt{KuratowskiOuterPlanarSubdigraph} returns the immutable list of immutable lists of
out\texttt{\symbol{45}}neighbours of an induced subdigraph (excluding multiple
edges and loops) of the digraph \mbox{\texttt{\mdseries\slshape digraph}} that witnesses the fact that \mbox{\texttt{\mdseries\slshape digraph}} is not outer planar, or \texttt{fail} if \mbox{\texttt{\mdseries\slshape digraph}} is outer planar. In other words, \texttt{KuratowskiOuterPlanarSubdigraph} returns the out\texttt{\symbol{45}}neighbours of a subdigraph of \mbox{\texttt{\mdseries\slshape digraph}} that is homeomorphic to the complete graph with \texttt{4} vertices, or to the complete bipartite graph with vertex sets of sizes \texttt{2} and \texttt{3}. 

 The directions and multiplicities of any edges in \mbox{\texttt{\mdseries\slshape digraph}} are ignored when considering whether or not \mbox{\texttt{\mdseries\slshape digraph}} is outer planar. 

 See also \texttt{IsOuterPlanarDigraph} (\ref{IsOuterPlanarDigraph}), \texttt{SubdigraphHomeomorphicToK4} (\ref{SubdigraphHomeomorphicToK4}), and \texttt{SubdigraphHomeomorphicToK23} (\ref{SubdigraphHomeomorphicToK23}). 

 This method uses the reference implementation in \href{https://github.com/graph-algorithms/edge-addition-planarity-suite} {edge-addition-planarity-suite}  by John Boyer of the algorithms described in \cite{BM06}. 
\begin{Verbatim}[commandchars=!@|,fontsize=\small,frame=single,label=Example]
  !gapprompt@gap>| !gapinput@D := Digraph([[3, 5, 10], [8, 9, 10], [1, 4], [3, 6], |
  !gapprompt@>| !gapinput@[1, 7, 11], [4, 7], [6, 8], [2, 7], [2, 11], [1, 2], [5, 9]]);|
  <immutable digraph with 11 vertices, 25 edges>
  !gapprompt@gap>| !gapinput@KuratowskiOuterPlanarSubdigraph(D);|
  [ [ 3, 5, 10 ], [ 9, 8, 10 ], [ 4, 1 ], [ 6, 3 ], [ 1, 11 ], 
    [ 7, 4 ], [ 8, 6 ], [ 2, 7 ], [ 11, 2 ], [ 2, 1 ], [ 5, 9 ] ]
  !gapprompt@gap>| !gapinput@D := Digraph([[2, 4, 7, 9, 10], [1, 3, 4, 6, 9, 10], [6, 10], |
  !gapprompt@>| !gapinput@[2, 5, 8, 9], [1, 2, 3, 4, 6, 7, 9, 10], [3, 4, 5, 7, 9, 10], |
  !gapprompt@>| !gapinput@[3, 4, 5, 6, 9, 10], [3, 4, 5, 7, 9], [2, 3, 5, 6, 7, 8], [3, 5]]);|
  <immutable digraph with 10 vertices, 50 edges>
  !gapprompt@gap>| !gapinput@IsOuterPlanarDigraph(D);|
  false
  !gapprompt@gap>| !gapinput@KuratowskiOuterPlanarSubdigraph(D);|
  [ [  ], [  ], [  ], [ 8, 9 ], [  ], [  ], [ 9, 4 ], [ 7, 9, 4 ], 
    [ 8, 7 ], [  ] ]
  !gapprompt@gap>| !gapinput@D := Digraph(IsMutableDigraph, [[3, 5, 10], [8, 9, 10], [1, 4],|
  !gapprompt@>| !gapinput@[3, 6], [1, 7, 11], [4, 7], [6, 8], [2, 7], [2, 11], [1, 2], [5, 9]]);|
  <mutable digraph with 11 vertices, 25 edges>
  !gapprompt@gap>| !gapinput@KuratowskiOuterPlanarSubdigraph(D);|
  [ [ 3, 5, 10 ], [ 9, 8, 10 ], [ 4, 1 ], [ 6, 3 ], [ 1, 11 ], 
    [ 7, 4 ], [ 8, 6 ], [ 2, 7 ], [ 11, 2 ], [ 2, 1 ], [ 5, 9 ] ]
  !gapprompt@gap>| !gapinput@D := Digraph(IsMutableDigraph, [[2, 4, 7, 9, 10],|
  !gapprompt@>| !gapinput@[1, 3, 4, 6, 9, 10], [6, 10], [2, 5, 8, 9],|
  !gapprompt@>| !gapinput@[1, 2, 3, 4, 6, 7, 9, 10], [3, 4, 5, 7, 9, 10],|
  !gapprompt@>| !gapinput@[3, 4, 5, 6, 9, 10], [3, 4, 5, 7, 9], [2, 3, 5, 6, 7, 8], [3, 5]]);|
  <mutable digraph with 10 vertices, 50 edges>
  !gapprompt@gap>| !gapinput@IsOuterPlanarDigraph(D);|
  false
  !gapprompt@gap>| !gapinput@KuratowskiOuterPlanarSubdigraph(D);|
  [ [  ], [  ], [  ], [ 8, 9 ], [  ], [  ], [ 9, 4 ], [ 7, 9, 4 ], 
    [ 8, 7 ], [  ] ]
\end{Verbatim}
 }

 

\subsection{\textcolor{Chapter }{PlanarEmbedding}}
\logpage{[ 5, 7, 3 ]}\nobreak
\hyperdef{L}{X84E3947E7D39BA64}{}
{\noindent\textcolor{FuncColor}{$\triangleright$\enspace\texttt{PlanarEmbedding({\mdseries\slshape digraph})\index{PlanarEmbedding@\texttt{PlanarEmbedding}}
\label{PlanarEmbedding}
}\hfill{\scriptsize (attribute)}}\\
\textbf{\indent Returns:\ }
A list or \texttt{fail}.



 If \mbox{\texttt{\mdseries\slshape digraph}} is a planar digraph, then \texttt{PlanarEmbedding} returns the immutable list of lists of out\texttt{\symbol{45}}neighbours of \mbox{\texttt{\mdseries\slshape digraph}} (excluding multiple edges and loops) such that each vertex's neighbours are
given in clockwise order. If \mbox{\texttt{\mdseries\slshape digraph}} is not planar, then \texttt{fail} is returned. 

 The directions and multiplicities of any edges in \mbox{\texttt{\mdseries\slshape digraph}} are ignored by \texttt{PlanarEmbedding}. 

 See also \texttt{IsPlanarDigraph} (\ref{IsPlanarDigraph}). 

 This method uses the reference implementation in \href{https://github.com/graph-algorithms/edge-addition-planarity-suite} {edge-addition-planarity-suite}  by John Boyer of the algorithms described in \cite{BM06}. 
\begin{Verbatim}[commandchars=!@|,fontsize=\small,frame=single,label=Example]
  !gapprompt@gap>| !gapinput@D := Digraph([[3, 5, 10], [8, 9, 10], [1, 4], [3, 6], |
  !gapprompt@>| !gapinput@[1, 7, 11], [4, 7], [6, 8], [2, 7], [2, 11], [1, 2], [5, 9]]);|
  <immutable digraph with 11 vertices, 25 edges>
  !gapprompt@gap>| !gapinput@PlanarEmbedding(D);|
  [ [ 3, 10, 5 ], [ 10, 8, 9 ], [ 4, 1 ], [ 6, 3 ], [ 1, 11, 7 ], 
    [ 7, 4 ], [ 8, 6 ], [ 7, 2 ], [ 2, 11 ], [ 1, 2 ], [ 9, 5 ] ]
  !gapprompt@gap>| !gapinput@D := Digraph([[2, 4, 7, 9, 10], [1, 3, 4, 6, 9, 10], [6, 10], |
  !gapprompt@>| !gapinput@[2, 5, 8, 9], [1, 2, 3, 4, 6, 7, 9, 10], [3, 4, 5, 7, 9, 10], |
  !gapprompt@>| !gapinput@[3, 4, 5, 6, 9, 10], [3, 4, 5, 7, 9], [2, 3, 5, 6, 7, 8], [3, 5]]);|
  <immutable digraph with 10 vertices, 50 edges>
  !gapprompt@gap>| !gapinput@PlanarEmbedding(D);|
  fail
  !gapprompt@gap>| !gapinput@D := Digraph(IsMutableDigraph, [[3, 5, 10], [8, 9, 10], [1, 4],|
  !gapprompt@>| !gapinput@[3, 6], [1, 7, 11], [4, 7], [6, 8], [2, 7], [2, 11], [1, 2], [5, 9]]);|
  <mutable digraph with 11 vertices, 25 edges>
  !gapprompt@gap>| !gapinput@PlanarEmbedding(D);|
  [ [ 3, 10, 5 ], [ 10, 8, 9 ], [ 4, 1 ], [ 6, 3 ], [ 1, 11, 7 ], 
    [ 7, 4 ], [ 8, 6 ], [ 7, 2 ], [ 2, 11 ], [ 1, 2 ], [ 9, 5 ] ]
  !gapprompt@gap>| !gapinput@D := Digraph(IsMutableDigraph, [[2, 4, 7, 9, 10],|
  !gapprompt@>| !gapinput@[1, 3, 4, 6, 9, 10], [6, 10], [2, 5, 8, 9],|
  !gapprompt@>| !gapinput@[1, 2, 3, 4, 6, 7, 9, 10], [3, 4, 5, 7, 9, 10],|
  !gapprompt@>| !gapinput@[3, 4, 5, 6, 9, 10], [3, 4, 5, 7, 9], [2, 3, 5, 6, 7, 8], [3, 5]]);|
  <mutable digraph with 10 vertices, 50 edges>
  !gapprompt@gap>| !gapinput@PlanarEmbedding(D);|
  fail
\end{Verbatim}
 }

 

\subsection{\textcolor{Chapter }{OuterPlanarEmbedding}}
\logpage{[ 5, 7, 4 ]}\nobreak
\hyperdef{L}{X85DFB8C18088711F}{}
{\noindent\textcolor{FuncColor}{$\triangleright$\enspace\texttt{OuterPlanarEmbedding({\mdseries\slshape digraph})\index{OuterPlanarEmbedding@\texttt{OuterPlanarEmbedding}}
\label{OuterPlanarEmbedding}
}\hfill{\scriptsize (attribute)}}\\
\textbf{\indent Returns:\ }
A list or \texttt{fail}.



 If \mbox{\texttt{\mdseries\slshape digraph}} is an outer planar digraph, then \texttt{OuterPlanarEmbedding} returns the immutable list of lists of out\texttt{\symbol{45}}neighbours of \mbox{\texttt{\mdseries\slshape digraph}} (excluding multiple edges and loops) such that each vertex's neighbours are
given in clockwise order. If \mbox{\texttt{\mdseries\slshape digraph}} is not outer planar, then \texttt{fail} is returned. 

 The directions and multiplicities of any edges in \mbox{\texttt{\mdseries\slshape digraph}} are ignored by \texttt{OuterPlanarEmbedding}. 

 See also \texttt{IsOuterPlanarDigraph} (\ref{IsOuterPlanarDigraph}). 

 This method uses the reference implementation in \href{https://github.com/graph-algorithms/edge-addition-planarity-suite} {edge-addition-planarity-suite}  by John Boyer of the algorithms described in \cite{BM06}. 
\begin{Verbatim}[commandchars=!@|,fontsize=\small,frame=single,label=Example]
  !gapprompt@gap>| !gapinput@D := Digraph([[3, 5, 10], [8, 9, 10], [1, 4], [3, 6], |
  !gapprompt@>| !gapinput@[1, 7, 11], [4, 7], [6, 8], [2, 7], [2, 11], [1, 2], [5, 9]]);|
  <immutable digraph with 11 vertices, 25 edges>
  !gapprompt@gap>| !gapinput@OuterPlanarEmbedding(D);|
  fail
  !gapprompt@gap>| !gapinput@D := Digraph([[2, 4, 7, 9, 10], [1, 3, 4, 6, 9, 10], [6, 10], |
  !gapprompt@>| !gapinput@[2, 5, 8, 9], [1, 2, 3, 4, 6, 7, 9, 10], [3, 4, 5, 7, 9, 10], |
  !gapprompt@>| !gapinput@[3, 4, 5, 6, 9, 10], [3, 4, 5, 7, 9], [2, 3, 5, 6, 7, 8], [3, 5]]);|
  <immutable digraph with 10 vertices, 50 edges>
  !gapprompt@gap>| !gapinput@OuterPlanarEmbedding(D);|
  fail
  !gapprompt@gap>| !gapinput@OuterPlanarEmbedding(CompleteBipartiteDigraph(2, 2));|
  [ [ 3, 4 ], [ 4, 3 ], [ 2, 1 ], [ 1, 2 ] ]
  !gapprompt@gap>| !gapinput@D := Digraph(IsMutableDigraph, [[3, 5, 10], [8, 9, 10], [1, 4],|
  !gapprompt@>| !gapinput@[3, 6], [1, 7, 11], [4, 7], [6, 8], [2, 7], [2, 11], [1, 2], [5, 9]]);|
  <mutable digraph with 11 vertices, 25 edges>
  !gapprompt@gap>| !gapinput@OuterPlanarEmbedding(D);|
  fail
  !gapprompt@gap>| !gapinput@D := Digraph(IsMutableDigraph, [[2, 4, 7, 9, 10],|
  !gapprompt@>| !gapinput@[1, 3, 4, 6, 9, 10], [6, 10], [2, 5, 8, 9],|
  !gapprompt@>| !gapinput@[1, 2, 3, 4, 6, 7, 9, 10], [3, 4, 5, 7, 9, 10],|
  !gapprompt@>| !gapinput@[3, 4, 5, 6, 9, 10], [3, 4, 5, 7, 9], [2, 3, 5, 6, 7, 8], [3, 5]]);|
  <mutable digraph with 10 vertices, 50 edges>
  !gapprompt@gap>| !gapinput@OuterPlanarEmbedding(D);|
  fail
  !gapprompt@gap>| !gapinput@OuterPlanarEmbedding(CompleteBipartiteDigraph(2, 2));|
  [ [ 3, 4 ], [ 4, 3 ], [ 2, 1 ], [ 1, 2 ] ]
\end{Verbatim}
 }

 

\subsection{\textcolor{Chapter }{SubdigraphHomeomorphicToK23}}
\logpage{[ 5, 7, 5 ]}\nobreak
\hyperdef{L}{X806D2D6B85E0B269}{}
{\noindent\textcolor{FuncColor}{$\triangleright$\enspace\texttt{SubdigraphHomeomorphicToK23({\mdseries\slshape digraph})\index{SubdigraphHomeomorphicToK23@\texttt{SubdigraphHomeomorphicToK23}}
\label{SubdigraphHomeomorphicToK23}
}\hfill{\scriptsize (attribute)}}\\
\noindent\textcolor{FuncColor}{$\triangleright$\enspace\texttt{SubdigraphHomeomorphicToK33({\mdseries\slshape digraph})\index{SubdigraphHomeomorphicToK33@\texttt{SubdigraphHomeomorphicToK33}}
\label{SubdigraphHomeomorphicToK33}
}\hfill{\scriptsize (attribute)}}\\
\noindent\textcolor{FuncColor}{$\triangleright$\enspace\texttt{SubdigraphHomeomorphicToK4({\mdseries\slshape digraph})\index{SubdigraphHomeomorphicToK4@\texttt{SubdigraphHomeomorphicToK4}}
\label{SubdigraphHomeomorphicToK4}
}\hfill{\scriptsize (attribute)}}\\
\textbf{\indent Returns:\ }
A list or \texttt{fail}.



 These attributes return the immutable list of lists of
out\texttt{\symbol{45}}neighbours of a subdigraph of the digraph \mbox{\texttt{\mdseries\slshape digraph}} which is homeomorphic to one of the following: the complete bipartite graph
with vertex sets of sizes \texttt{2} and \texttt{3}; the complete bipartite graph with vertex sets of sizes \texttt{3} and \texttt{3}; or the complete graph with \texttt{4} vertices. If \mbox{\texttt{\mdseries\slshape digraph}} has no such subdigraphs, then \texttt{fail} is returned. 

 See also \texttt{IsPlanarDigraph} (\ref{IsPlanarDigraph}) and \texttt{IsOuterPlanarDigraph} (\ref{IsOuterPlanarDigraph}) for more details.

 This method uses the reference implementation in \href{https://github.com/graph-algorithms/edge-addition-planarity-suite} {edge-addition-planarity-suite}  by John Boyer of the algorithms described in \cite{BM06}. 
\begin{Verbatim}[commandchars=!@|,fontsize=\small,frame=single,label=Example]
  !gapprompt@gap>| !gapinput@D := Digraph([[3, 5, 10], [8, 9, 10], [1, 4], [3, 6], [1, 7, 11], |
  !gapprompt@>| !gapinput@[4, 7], [6, 8], [2, 7], [2, 11], [1, 2], [5, 9]]);|
  <immutable digraph with 11 vertices, 25 edges>
  !gapprompt@gap>| !gapinput@SubdigraphHomeomorphicToK4(D);|
  [ [ 3, 5, 10 ], [ 9, 8, 10 ], [ 4, 1 ], [ 6, 3 ], [ 1, 7, 11 ], 
    [ 7, 4 ], [ 8, 6 ], [ 2, 7 ], [ 11, 2 ], [ 2, 1 ], [ 5, 9 ] ]
  !gapprompt@gap>| !gapinput@SubdigraphHomeomorphicToK23(D);|
  [ [ 3, 5, 10 ], [ 9, 8, 10 ], [ 4, 1 ], [ 6, 3 ], [ 1, 11 ], 
    [ 7, 4 ], [ 8, 6 ], [ 2, 7 ], [ 11, 2 ], [ 2, 1 ], [ 5, 9 ] ]
  !gapprompt@gap>| !gapinput@D := Digraph([[3, 5, 10], [8, 9, 10], [1, 4], [3, 6], [1, 11], |
  !gapprompt@>| !gapinput@ [4, 7], [6, 8], [2, 7], [2, 11], [1, 2], [5, 9]]);|
  <immutable digraph with 11 vertices, 24 edges>
  !gapprompt@gap>| !gapinput@SubdigraphHomeomorphicToK4(D);|
  fail
  !gapprompt@gap>| !gapinput@SubdigraphHomeomorphicToK23(D);|
  [ [ 3, 10, 5 ], [ 10, 8, 9 ], [ 4, 1 ], [ 6, 3 ], [ 11, 1 ], 
    [ 7, 4 ], [ 8, 6 ], [ 2, 7 ], [ 2, 11 ], [ 1, 2 ], [ 9, 5 ] ]
  !gapprompt@gap>| !gapinput@SubdigraphHomeomorphicToK33(D);|
  fail
  !gapprompt@gap>| !gapinput@SubdigraphHomeomorphicToK23(NullDigraph(0));|
  fail
  !gapprompt@gap>| !gapinput@SubdigraphHomeomorphicToK33(CompleteDigraph(5));|
  fail
  !gapprompt@gap>| !gapinput@SubdigraphHomeomorphicToK33(CompleteBipartiteDigraph(3, 3));|
  [ [ 4, 6, 5 ], [ 4, 5, 6 ], [ 6, 5, 4 ], [ 1, 2, 3 ], [ 3, 2, 1 ], 
    [ 2, 3, 1 ] ]
  !gapprompt@gap>| !gapinput@SubdigraphHomeomorphicToK4(CompleteDigraph(3));|
  fail
  !gapprompt@gap>| !gapinput@D := Digraph(IsMutableDigraph, [[3, 5, 10], [8, 9, 10], [1, 4],|
  !gapprompt@>| !gapinput@[3, 6], [1, 7, 11], [4, 7], [6, 8], [2, 7], [2, 11], [1, 2], [5, 9]]);|
  <mutable digraph with 11 vertices, 25 edges>
  !gapprompt@gap>| !gapinput@SubdigraphHomeomorphicToK4(D);|
  [ [ 3, 5, 10 ], [ 9, 8, 10 ], [ 4, 1 ], [ 6, 3 ], [ 1, 7, 11 ], 
    [ 7, 4 ], [ 8, 6 ], [ 2, 7 ], [ 11, 2 ], [ 2, 1 ], [ 5, 9 ] ]
  !gapprompt@gap>| !gapinput@SubdigraphHomeomorphicToK23(D);|
  [ [ 3, 5, 10 ], [ 9, 8, 10 ], [ 4, 1 ], [ 6, 3 ], [ 1, 11 ], 
    [ 7, 4 ], [ 8, 6 ], [ 2, 7 ], [ 11, 2 ], [ 2, 1 ], [ 5, 9 ] ]
  !gapprompt@gap>| !gapinput@D := Digraph(IsMutableDigraph, [[3, 5, 10], [8, 9, 10], [1, 4],|
  !gapprompt@>| !gapinput@[3, 6], [1, 11], [4, 7], [6, 8], [2, 7], [2, 11], [1, 2], [5, 9]]);|
  <mutable digraph with 11 vertices, 24 edges>
  !gapprompt@gap>| !gapinput@SubdigraphHomeomorphicToK4(D);|
  fail
  !gapprompt@gap>| !gapinput@SubdigraphHomeomorphicToK23(D);|
  [ [ 3, 10, 5 ], [ 10, 8, 9 ], [ 4, 1 ], [ 6, 3 ], [ 11, 1 ], 
    [ 7, 4 ], [ 8, 6 ], [ 2, 7 ], [ 2, 11 ], [ 1, 2 ], [ 9, 5 ] ]
  !gapprompt@gap>| !gapinput@SubdigraphHomeomorphicToK33(D);|
  fail
  !gapprompt@gap>| !gapinput@SubdigraphHomeomorphicToK23(NullDigraph(0));|
  fail
  !gapprompt@gap>| !gapinput@SubdigraphHomeomorphicToK33(CompleteDigraph(5));|
  fail
  !gapprompt@gap>| !gapinput@SubdigraphHomeomorphicToK33(CompleteBipartiteDigraph(3, 3));|
  [ [ 4, 6, 5 ], [ 4, 5, 6 ], [ 6, 5, 4 ], [ 1, 2, 3 ], [ 3, 2, 1 ], 
    [ 2, 3, 1 ] ]
  !gapprompt@gap>| !gapinput@SubdigraphHomeomorphicToK4(CompleteDigraph(3));|
  fail
\end{Verbatim}
 }

 

\subsection{\textcolor{Chapter }{DualPlanarGraph}}
\logpage{[ 5, 7, 6 ]}\nobreak
\hyperdef{L}{X876803517C34E1CD}{}
{\noindent\textcolor{FuncColor}{$\triangleright$\enspace\texttt{DualPlanarGraph({\mdseries\slshape digraph})\index{DualPlanarGraph@\texttt{DualPlanarGraph}}
\label{DualPlanarGraph}
}\hfill{\scriptsize (attribute)}}\\
\textbf{\indent Returns:\ }
A digraph or \texttt{fail}.



 If \mbox{\texttt{\mdseries\slshape digraph}} is a planar digraph, then \texttt{DualPlanarGraph} returns the the dual graph of \mbox{\texttt{\mdseries\slshape digraph}}. If \mbox{\texttt{\mdseries\slshape digraph}} is not planar, then \texttt{fail} is returned.

 The dual graph of a planar digraph \mbox{\texttt{\mdseries\slshape digraph}} has a vertex for each face of \mbox{\texttt{\mdseries\slshape digraph}} and an edge for each pair of faces that are separated by an edge from each
other. Vertex \mbox{\texttt{\mdseries\slshape i}} of the dual graph corresponds to the facial walk at the \mbox{\texttt{\mdseries\slshape i}}\texttt{\symbol{45}}th position calling \texttt{FacialWalks} (\ref{FacialWalks}) of \mbox{\texttt{\mdseries\slshape digraph}} with the rotation system returned by \texttt{PlanarEmbedding} (\ref{PlanarEmbedding}). 

 Note that \texttt{PlanarEmbedding} (\ref{PlanarEmbedding}), and therefore \texttt{DualPlanarGraph}, uses the reference implementation in \href{https://github.com/graph-algorithms/edge-addition-planarity-suite} {edge-addition-planarity-suite}  by John Boyer of the algorithms described in \cite{BM06}. 
\begin{Verbatim}[commandchars=!@|,fontsize=\small,frame=single,label=Example]
  !gapprompt@gap>| !gapinput@D := CompleteDigraph(4);;|
  !gapprompt@gap>| !gapinput@dualD := DualPlanarGraph(D);|
  <immutable digraph with 4 vertices, 12 edges>
  !gapprompt@gap>| !gapinput@IsIsomorphicDigraph(D, dualD);|
  true
  !gapprompt@gap>| !gapinput@cube := Digraph([[2, 4, 5], [1, 3, 6], [2, 4, 7], [1, 3, 8], |
  !gapprompt@>| !gapinput@                    [1, 6, 8], [2, 5, 7], [3, 6, 8], [4, 5, 7]]);|
  <immutable digraph with 8 vertices, 24 edges>
  !gapprompt@gap>| !gapinput@oct := DualPlanarGraph(cube);;|
  !gapprompt@gap>| !gapinput@IsIsomorphicDigraph(oct, CompleteMultipartiteDigraph([2, 2, 2]));|
  true
\end{Verbatim}
 }

 }

 
\section{\textcolor{Chapter }{Hashing}}\logpage{[ 5, 8, 0 ]}
\hyperdef{L}{X7FF5E2D07D48647E}{}
{
 

\subsection{\textcolor{Chapter }{DigraphHash}}
\logpage{[ 5, 8, 1 ]}\nobreak
\hyperdef{L}{X7D97D990867B0149}{}
{\noindent\textcolor{FuncColor}{$\triangleright$\enspace\texttt{DigraphHash({\mdseries\slshape digraph})\index{DigraphHash@\texttt{DigraphHash}}
\label{DigraphHash}
}\hfill{\scriptsize (attribute)}}\\
\textbf{\indent Returns:\ }
The hash of a digraph.



 This function returns a hash of \mbox{\texttt{\mdseries\slshape digraph}}. 

 Note that the underlying hashing function is system dependent, and so the
value of this attribute is not guaranteed to be the same on different systems. }

 }

 }

 
\chapter{\textcolor{Chapter }{Properties of digraphs}}\label{Properties of digraphs}
\logpage{[ 6, 0, 0 ]}
\hyperdef{L}{X7ADDEFD478D470D5}{}
{
 
\section{\textcolor{Chapter }{Vertex properties}}\logpage{[ 6, 1, 0 ]}
\hyperdef{L}{X7D6A52AB7C69A9CA}{}
{
 

\subsection{\textcolor{Chapter }{DigraphHasAVertex}}
\logpage{[ 6, 1, 1 ]}\nobreak
\hyperdef{L}{X82DBFFF17E708803}{}
{\noindent\textcolor{FuncColor}{$\triangleright$\enspace\texttt{DigraphHasAVertex({\mdseries\slshape digraph})\index{DigraphHasAVertex@\texttt{DigraphHasAVertex}}
\label{DigraphHasAVertex}
}\hfill{\scriptsize (property)}}\\
\textbf{\indent Returns:\ }
\texttt{true} or \texttt{false}.



 Returns \texttt{true} if the digraph \mbox{\texttt{\mdseries\slshape digraph}} has at least one vertex, and \texttt{false} if it does not.

 See also \texttt{DigraphHasNoVertices} (\ref{DigraphHasNoVertices}). 

 If the argument \mbox{\texttt{\mdseries\slshape digraph}} is mutable, then the return value of this property is recomputed every time it
is called.

 
\begin{Verbatim}[commandchars=!@|,fontsize=\small,frame=single,label=Example]
  !gapprompt@gap>| !gapinput@D := Digraph([]);|
  <immutable empty digraph with 0 vertices>
  !gapprompt@gap>| !gapinput@DigraphHasAVertex(D);|
  false
  !gapprompt@gap>| !gapinput@D := Digraph([[]]);|
  <immutable empty digraph with 1 vertex>
  !gapprompt@gap>| !gapinput@DigraphHasAVertex(D);|
  true
  !gapprompt@gap>| !gapinput@D := Digraph([[], [1]]);|
  <immutable digraph with 2 vertices, 1 edge>
  !gapprompt@gap>| !gapinput@DigraphHasAVertex(D);|
  true
\end{Verbatim}
 }

 

\subsection{\textcolor{Chapter }{DigraphHasNoVertices}}
\logpage{[ 6, 1, 2 ]}\nobreak
\hyperdef{L}{X78E16B8B7DF55B59}{}
{\noindent\textcolor{FuncColor}{$\triangleright$\enspace\texttt{DigraphHasNoVertices({\mdseries\slshape digraph})\index{DigraphHasNoVertices@\texttt{DigraphHasNoVertices}}
\label{DigraphHasNoVertices}
}\hfill{\scriptsize (property)}}\\
\textbf{\indent Returns:\ }
\texttt{true} or \texttt{false}.



 Returns \texttt{true} if the digraph \mbox{\texttt{\mdseries\slshape digraph}} is the unique digraph with zero vertices, and \texttt{false} otherwise.

 See also \texttt{DigraphHasAVertex} (\ref{DigraphHasAVertex}). 

 If the argument \mbox{\texttt{\mdseries\slshape digraph}} is mutable, then the return value of this property is recomputed every time it
is called.

 
\begin{Verbatim}[commandchars=!@|,fontsize=\small,frame=single,label=Example]
  !gapprompt@gap>| !gapinput@D := Digraph([]);|
  <immutable empty digraph with 0 vertices>
  !gapprompt@gap>| !gapinput@DigraphHasNoVertices(D);|
  true
  !gapprompt@gap>| !gapinput@D := Digraph([[]]);|
  <immutable empty digraph with 1 vertex>
  !gapprompt@gap>| !gapinput@DigraphHasNoVertices(D);|
  false
  !gapprompt@gap>| !gapinput@D := Digraph([[], [1]]);|
  <immutable digraph with 2 vertices, 1 edge>
  !gapprompt@gap>| !gapinput@DigraphHasNoVertices(D);|
  false
\end{Verbatim}
 }

 }

 
\section{\textcolor{Chapter }{Edge properties}}\logpage{[ 6, 2, 0 ]}
\hyperdef{L}{X7F49388F8245F0E9}{}
{
 

\subsection{\textcolor{Chapter }{DigraphHasLoops}}
\logpage{[ 6, 2, 1 ]}\nobreak
\hyperdef{L}{X7D92935C7D535187}{}
{\noindent\textcolor{FuncColor}{$\triangleright$\enspace\texttt{DigraphHasLoops({\mdseries\slshape digraph})\index{DigraphHasLoops@\texttt{DigraphHasLoops}}
\label{DigraphHasLoops}
}\hfill{\scriptsize (property)}}\\
\textbf{\indent Returns:\ }
\texttt{true} or \texttt{false}.



 Returns \texttt{true} if the digraph \mbox{\texttt{\mdseries\slshape digraph}} has loops, and \texttt{false} if it does not. A loop is an edge with equal source and range. 

 If the argument \mbox{\texttt{\mdseries\slshape digraph}} is mutable, then the return value of this property is recomputed every time it
is called.

 
\begin{Verbatim}[commandchars=!@|,fontsize=\small,frame=single,label=Example]
  !gapprompt@gap>| !gapinput@D := Digraph([[1, 2], [2]]);|
  <immutable digraph with 2 vertices, 3 edges>
  !gapprompt@gap>| !gapinput@DigraphEdges(D);|
  [ [ 1, 1 ], [ 1, 2 ], [ 2, 2 ] ]
  !gapprompt@gap>| !gapinput@DigraphHasLoops(D);|
  true
  !gapprompt@gap>| !gapinput@D := Digraph([[2, 3], [1], [2]]);|
  <immutable digraph with 3 vertices, 4 edges>
  !gapprompt@gap>| !gapinput@DigraphEdges(D);|
  [ [ 1, 2 ], [ 1, 3 ], [ 2, 1 ], [ 3, 2 ] ]
  !gapprompt@gap>| !gapinput@DigraphHasLoops(D);|
  false
  !gapprompt@gap>| !gapinput@D := CompleteDigraph(IsMutableDigraph, 4);|
  <mutable digraph with 4 vertices, 12 edges>
  !gapprompt@gap>| !gapinput@DigraphHasLoops(D);|
  false
\end{Verbatim}
 }

 

\subsection{\textcolor{Chapter }{IsAntiSymmetricDigraph}}
\logpage{[ 6, 2, 2 ]}\nobreak
\hyperdef{L}{X7DB1BC2286FC08E2}{}
{\noindent\textcolor{FuncColor}{$\triangleright$\enspace\texttt{IsAntiSymmetricDigraph({\mdseries\slshape digraph})\index{IsAntiSymmetricDigraph@\texttt{IsAntiSymmetricDigraph}}
\label{IsAntiSymmetricDigraph}
}\hfill{\scriptsize (property)}}\\
\noindent\textcolor{FuncColor}{$\triangleright$\enspace\texttt{IsAntisymmetricDigraph({\mdseries\slshape digraph})\index{IsAntisymmetricDigraph@\texttt{IsAntisymmetricDigraph}}
\label{IsAntisymmetricDigraph}
}\hfill{\scriptsize (property)}}\\
\textbf{\indent Returns:\ }
\texttt{true} or \texttt{false}.



 This property is \texttt{true} if the digraph \mbox{\texttt{\mdseries\slshape digraph}} is antisymmetric, and \texttt{false} if it is not. 

 A digraph is \emph{antisymmetric} if whenever there is an edge with source \texttt{u} and range \texttt{v}, and an edge with source \texttt{v} and range \texttt{u}, then the vertices \texttt{u} and \texttt{v} are equal. 

 If the argument \mbox{\texttt{\mdseries\slshape digraph}} is mutable, then the return value of this property is recomputed every time it
is called.

 
\begin{Verbatim}[commandchars=!@|,fontsize=\small,frame=single,label=Example]
  !gapprompt@gap>| !gapinput@gr1 := Digraph([[2], [1, 3], [2, 3]]);|
  <immutable digraph with 3 vertices, 5 edges>
  !gapprompt@gap>| !gapinput@IsAntisymmetricDigraph(gr1);|
  false
  !gapprompt@gap>| !gapinput@DigraphEdges(gr1){[1, 2]};|
  [ [ 1, 2 ], [ 2, 1 ] ]
  !gapprompt@gap>| !gapinput@gr2 := Digraph([[1, 2], [3, 3], [1]]);|
  <immutable multidigraph with 3 vertices, 5 edges>
  !gapprompt@gap>| !gapinput@IsAntisymmetricDigraph(gr2);|
  true
  !gapprompt@gap>| !gapinput@DigraphEdges(gr2);|
  [ [ 1, 1 ], [ 1, 2 ], [ 2, 3 ], [ 2, 3 ], [ 3, 1 ] ]
\end{Verbatim}
 }

 

\subsection{\textcolor{Chapter }{IsBipartiteDigraph}}
\logpage{[ 6, 2, 3 ]}\nobreak
\hyperdef{L}{X860CFB0C8665F356}{}
{\noindent\textcolor{FuncColor}{$\triangleright$\enspace\texttt{IsBipartiteDigraph({\mdseries\slshape digraph})\index{IsBipartiteDigraph@\texttt{IsBipartiteDigraph}}
\label{IsBipartiteDigraph}
}\hfill{\scriptsize (property)}}\\
\textbf{\indent Returns:\ }
\texttt{true} or \texttt{false}.



 This property is \texttt{true} if the digraph \mbox{\texttt{\mdseries\slshape digraph}} is bipartite, and \texttt{false} if it is not. A digraph is bipartite if and only if the vertices of \mbox{\texttt{\mdseries\slshape digraph}} can be partitioned into two non\texttt{\symbol{45}}empty sets such that the
source and range of any edge of \mbox{\texttt{\mdseries\slshape digraph}} lie in distinct sets. Equivalently, a digraph is bipartite if and only if it
is 2\texttt{\symbol{45}}colorable; see \texttt{DigraphGreedyColouring} (\ref{DigraphGreedyColouring:for a digraph}). 

 See also \texttt{DigraphBicomponents} (\ref{DigraphBicomponents}). 

 If the argument \mbox{\texttt{\mdseries\slshape digraph}} is mutable, then the return value of this property is recomputed every time it
is called.

 
\begin{Verbatim}[commandchars=!@|,fontsize=\small,frame=single,label=Example]
  !gapprompt@gap>| !gapinput@D := ChainDigraph(4);|
  <immutable chain digraph with 4 vertices>
  !gapprompt@gap>| !gapinput@IsBipartiteDigraph(D);|
  true
  !gapprompt@gap>| !gapinput@D := CycleDigraph(3);|
  <immutable cycle digraph with 3 vertices>
  !gapprompt@gap>| !gapinput@IsBipartiteDigraph(D);|
  false
  !gapprompt@gap>| !gapinput@D := CompleteBipartiteDigraph(IsMutableDigraph, 5, 4);|
  <mutable digraph with 9 vertices, 40 edges>
  !gapprompt@gap>| !gapinput@IsBipartiteDigraph(D);|
  true
\end{Verbatim}
 }

 

\subsection{\textcolor{Chapter }{IsCompleteBipartiteDigraph}}
\logpage{[ 6, 2, 4 ]}\nobreak
\hyperdef{L}{X790A4DBF8533516E}{}
{\noindent\textcolor{FuncColor}{$\triangleright$\enspace\texttt{IsCompleteBipartiteDigraph({\mdseries\slshape digraph})\index{IsCompleteBipartiteDigraph@\texttt{IsCompleteBipartiteDigraph}}
\label{IsCompleteBipartiteDigraph}
}\hfill{\scriptsize (property)}}\\
\textbf{\indent Returns:\ }
\texttt{true} or \texttt{false}.



 Returns \texttt{true} if the digraph \mbox{\texttt{\mdseries\slshape digraph}} is a complete bipartite digraph, and \texttt{false} if it is not. 

 A digraph is a \emph{complete bipartite digraph} if it is bipartite, see \texttt{IsBipartiteDigraph} (\ref{IsBipartiteDigraph}), and there exists a unique edge with source \texttt{i} and range \texttt{j} if and only if \texttt{i} and \texttt{j} lie in different bicomponents of \mbox{\texttt{\mdseries\slshape digraph}}, see \texttt{DigraphBicomponents} (\ref{DigraphBicomponents}). 

 Equivalently, a bipartite digraph with bicomponents of size $m$ and $n$ is complete precisely when it has $2mn$ edges, none of which are multiple edges. 

 See also \texttt{CompleteBipartiteDigraph} (\ref{CompleteBipartiteDigraph}). 

 If the argument \mbox{\texttt{\mdseries\slshape digraph}} is mutable, then the return value of this property is recomputed every time it
is called.

 
\begin{Verbatim}[commandchars=!@|,fontsize=\small,frame=single,label=Example]
  !gapprompt@gap>| !gapinput@D := CycleDigraph(2);|
  <immutable cycle digraph with 2 vertices>
  !gapprompt@gap>| !gapinput@IsCompleteBipartiteDigraph(D);|
  true
  !gapprompt@gap>| !gapinput@D := CycleDigraph(4);|
  <immutable cycle digraph with 4 vertices>
  !gapprompt@gap>| !gapinput@IsBipartiteDigraph(D);|
  true
  !gapprompt@gap>| !gapinput@IsCompleteBipartiteDigraph(D);|
  false
  !gapprompt@gap>| !gapinput@D := CompleteBipartiteDigraph(IsMutableDigraph, 5, 4);|
  <mutable digraph with 9 vertices, 40 edges>
  !gapprompt@gap>| !gapinput@IsCompleteBipartiteDigraph(D);|
  true
\end{Verbatim}
 }

 

\subsection{\textcolor{Chapter }{IsCompleteDigraph}}
\logpage{[ 6, 2, 5 ]}\nobreak
\hyperdef{L}{X81F28D4D879FE3B2}{}
{\noindent\textcolor{FuncColor}{$\triangleright$\enspace\texttt{IsCompleteDigraph({\mdseries\slshape digraph})\index{IsCompleteDigraph@\texttt{IsCompleteDigraph}}
\label{IsCompleteDigraph}
}\hfill{\scriptsize (property)}}\\
\textbf{\indent Returns:\ }
\texttt{true} or \texttt{false}.



 Returns \texttt{true} if the digraph \mbox{\texttt{\mdseries\slshape digraph}} is complete, and \texttt{false} if it is not. 

 A digraph is \emph{complete} if it has no loops, and for all \emph{distinct} vertices \texttt{i} and \texttt{j}, there is exactly one edge with source \texttt{i} and range \texttt{j}. Equivalently, a digraph with $n$ vertices is complete precisely when it has $n(n - 1)$ edges, no loops, and no multiple edges. 

 If the argument \mbox{\texttt{\mdseries\slshape digraph}} is mutable, then the return value of this property is recomputed every time it
is called.

 
\begin{Verbatim}[commandchars=!@|,fontsize=\small,frame=single,label=Example]
  !gapprompt@gap>| !gapinput@D := Digraph([[2, 3], [1, 3], [1, 2]]);|
  <immutable digraph with 3 vertices, 6 edges>
  !gapprompt@gap>| !gapinput@IsCompleteDigraph(D);|
  true
  !gapprompt@gap>| !gapinput@D := Digraph([[2, 2], [1]]);|
  <immutable multidigraph with 2 vertices, 3 edges>
  !gapprompt@gap>| !gapinput@IsCompleteDigraph(D);|
  false
  !gapprompt@gap>| !gapinput@D := CompleteBipartiteDigraph(IsMutableDigraph, 5, 4);|
  <mutable digraph with 9 vertices, 40 edges>
  !gapprompt@gap>| !gapinput@IsCompleteDigraph(D);|
  false
\end{Verbatim}
 }

 

\subsection{\textcolor{Chapter }{IsCompleteMultipartiteDigraph}}
\logpage{[ 6, 2, 6 ]}\nobreak
\hyperdef{L}{X7FD1A15779FEC341}{}
{\noindent\textcolor{FuncColor}{$\triangleright$\enspace\texttt{IsCompleteMultipartiteDigraph({\mdseries\slshape digraph})\index{IsCompleteMultipartiteDigraph@\texttt{IsCompleteMultipartiteDigraph}}
\label{IsCompleteMultipartiteDigraph}
}\hfill{\scriptsize (property)}}\\
\textbf{\indent Returns:\ }
\texttt{true} or \texttt{false}.



 This property returns \texttt{true} if \mbox{\texttt{\mdseries\slshape digraph}} is a complete multipartite digraph, and \texttt{false} if not. 

 A digraph is a \emph{complete multipartite digraph} if and only if its vertices can be partitioned into at least two maximal
independent sets, where every possible edge between these independent sets
occurs in the digraph exactly once. 

 If the argument \mbox{\texttt{\mdseries\slshape digraph}} is mutable, then the return value of this property is recomputed every time it
is called.

 
\begin{Verbatim}[commandchars=!@|,fontsize=\small,frame=single,label=Example]
  !gapprompt@gap>| !gapinput@D := CompleteMultipartiteDigraph([2, 4, 6]);|
  <immutable complete multipartite digraph with 12 vertices, 88 edges>
  !gapprompt@gap>| !gapinput@IsCompleteMultipartiteDigraph(D);|
  true
  !gapprompt@gap>| !gapinput@D := CompleteBipartiteDigraph(IsMutableDigraph, 5, 4);|
  <mutable digraph with 9 vertices, 40 edges>
  !gapprompt@gap>| !gapinput@IsCompleteMultipartiteDigraph(D);|
  true
\end{Verbatim}
 }

 

\subsection{\textcolor{Chapter }{IsEmptyDigraph}}
\logpage{[ 6, 2, 7 ]}\nobreak
\hyperdef{L}{X7E894CCF7B1C27AE}{}
{\noindent\textcolor{FuncColor}{$\triangleright$\enspace\texttt{IsEmptyDigraph({\mdseries\slshape digraph})\index{IsEmptyDigraph@\texttt{IsEmptyDigraph}}
\label{IsEmptyDigraph}
}\hfill{\scriptsize (property)}}\\
\noindent\textcolor{FuncColor}{$\triangleright$\enspace\texttt{IsNullDigraph({\mdseries\slshape digraph})\index{IsNullDigraph@\texttt{IsNullDigraph}}
\label{IsNullDigraph}
}\hfill{\scriptsize (property)}}\\
\textbf{\indent Returns:\ }
\texttt{true} or \texttt{false}.



 Returns \texttt{true} if the digraph \mbox{\texttt{\mdseries\slshape digraph}} is empty, and \texttt{false} if it is not. A digraph is \emph{empty} if it has no edges.

 \texttt{IsNullDigraph} is a synonym for \texttt{IsEmptyDigraph}. See also \texttt{IsNonemptyDigraph} (\ref{IsNonemptyDigraph}). 

 If the argument \mbox{\texttt{\mdseries\slshape digraph}} is mutable, then the return value of this property is recomputed every time it
is called.

 
\begin{Verbatim}[commandchars=!@|,fontsize=\small,frame=single,label=Example]
  !gapprompt@gap>| !gapinput@D := Digraph([[], []]);|
  <immutable empty digraph with 2 vertices>
  !gapprompt@gap>| !gapinput@IsEmptyDigraph(D);|
  true
  !gapprompt@gap>| !gapinput@IsNullDigraph(D);|
  true
  !gapprompt@gap>| !gapinput@D := Digraph([[], [1]]);|
  <immutable digraph with 2 vertices, 1 edge>
  !gapprompt@gap>| !gapinput@IsEmptyDigraph(D);|
  false
  !gapprompt@gap>| !gapinput@IsNullDigraph(D);|
  false
\end{Verbatim}
 }

 

\subsection{\textcolor{Chapter }{IsEquivalenceDigraph}}
\logpage{[ 6, 2, 8 ]}\nobreak
\hyperdef{L}{X83762C257DED2751}{}
{\noindent\textcolor{FuncColor}{$\triangleright$\enspace\texttt{IsEquivalenceDigraph({\mdseries\slshape digraph})\index{IsEquivalenceDigraph@\texttt{IsEquivalenceDigraph}}
\label{IsEquivalenceDigraph}
}\hfill{\scriptsize (property)}}\\
\textbf{\indent Returns:\ }
\texttt{true} or \texttt{false}.



 A digraph is an equivalence digraph if and only if the digraph satisfies all
of \texttt{IsReflexiveDigraph} (\ref{IsReflexiveDigraph}), \texttt{IsSymmetricDigraph} (\ref{IsSymmetricDigraph}) and \texttt{IsTransitiveDigraph} (\ref{IsTransitiveDigraph}). A partial order \mbox{\texttt{\mdseries\slshape digraph}} corresponds to an equivalence relation. 

 If the argument \mbox{\texttt{\mdseries\slshape digraph}} is mutable, then the return value of this property is recomputed every time it
is called.

 
\begin{Verbatim}[commandchars=!@|,fontsize=\small,frame=single,label=Example]
  !gapprompt@gap>| !gapinput@D := Digraph([[1, 3], [2], [1, 3]]);|
  <immutable digraph with 3 vertices, 5 edges>
  !gapprompt@gap>| !gapinput@IsEquivalenceDigraph(D);|
  true
\end{Verbatim}
 }

 

\subsection{\textcolor{Chapter }{IsFunctionalDigraph}}
\logpage{[ 6, 2, 9 ]}\nobreak
\hyperdef{L}{X7E8F37E585DAED52}{}
{\noindent\textcolor{FuncColor}{$\triangleright$\enspace\texttt{IsFunctionalDigraph({\mdseries\slshape digraph})\index{IsFunctionalDigraph@\texttt{IsFunctionalDigraph}}
\label{IsFunctionalDigraph}
}\hfill{\scriptsize (property)}}\\
\textbf{\indent Returns:\ }
\texttt{true} or \texttt{false}.



 This property is \texttt{true} if the digraph \mbox{\texttt{\mdseries\slshape digraph}} is functional. 

 A digraph is \emph{functional} if every vertex is the source of a unique edge. 

 If the argument \mbox{\texttt{\mdseries\slshape digraph}} is mutable, then the return value of this property is recomputed every time it
is called.

 
\begin{Verbatim}[commandchars=!@|,fontsize=\small,frame=single,label=Example]
  !gapprompt@gap>| !gapinput@gr1 := Digraph([[3], [2], [2], [1], [6], [5]]);|
  <immutable digraph with 6 vertices, 6 edges>
  !gapprompt@gap>| !gapinput@IsFunctionalDigraph(gr1);|
  true
  !gapprompt@gap>| !gapinput@gr2 := Digraph([[1, 2], [1]]);|
  <immutable digraph with 2 vertices, 3 edges>
  !gapprompt@gap>| !gapinput@IsFunctionalDigraph(gr2);|
  false
  !gapprompt@gap>| !gapinput@gr3 := Digraph(3, [1, 2, 3], [2, 3, 1]);|
  <immutable digraph with 3 vertices, 3 edges>
  !gapprompt@gap>| !gapinput@IsFunctionalDigraph(gr3);|
  true
\end{Verbatim}
 }

 

\subsection{\textcolor{Chapter }{IsPermutationDigraph}}
\logpage{[ 6, 2, 10 ]}\nobreak
\hyperdef{L}{X793FB02C7C59D85B}{}
{\noindent\textcolor{FuncColor}{$\triangleright$\enspace\texttt{IsPermutationDigraph({\mdseries\slshape digraph})\index{IsPermutationDigraph@\texttt{IsPermutationDigraph}}
\label{IsPermutationDigraph}
}\hfill{\scriptsize (property)}}\\
\textbf{\indent Returns:\ }
\texttt{true} or \texttt{false}.



 This property is \texttt{true} if the digraph \mbox{\texttt{\mdseries\slshape digraph}} is functional and each node has only one in\texttt{\symbol{45}}neighbour. 

 If the argument \mbox{\texttt{\mdseries\slshape digraph}} is mutable, then the return value of this property is recomputed every time it
is called.

 
\begin{Verbatim}[commandchars=!@|,fontsize=\small,frame=single,label=Example]
  !gapprompt@gap>| !gapinput@gr1 := Digraph([[3], [2], [2], [1], [6], [5]]);|
  <immutable digraph with 6 vertices, 6 edges>
  !gapprompt@gap>| !gapinput@IsPermutationDigraph(gr1);|
  false
  !gapprompt@gap>| !gapinput@gr2 := Digraph([[1, 2], [1]]);|
  <immutable digraph with 2 vertices, 3 edges>
  !gapprompt@gap>| !gapinput@IsPermutationDigraph(gr2);|
  false
  !gapprompt@gap>| !gapinput@gr3 := Digraph(3, [1, 2, 3], [2, 3, 1]);|
  <immutable digraph with 3 vertices, 3 edges>
  !gapprompt@gap>| !gapinput@IsPermutationDigraph(gr3);|
  true
\end{Verbatim}
 }

 

\subsection{\textcolor{Chapter }{IsMultiDigraph}}
\logpage{[ 6, 2, 11 ]}\nobreak
\hyperdef{L}{X7BB84CFC7E8B2B26}{}
{\noindent\textcolor{FuncColor}{$\triangleright$\enspace\texttt{IsMultiDigraph({\mdseries\slshape digraph})\index{IsMultiDigraph@\texttt{IsMultiDigraph}}
\label{IsMultiDigraph}
}\hfill{\scriptsize (property)}}\\
\textbf{\indent Returns:\ }
\texttt{true} or \texttt{false}.



 A \emph{multidigraph} is one that has at least two edges with equal source and range.

 If the argument \mbox{\texttt{\mdseries\slshape digraph}} is mutable, then the return value of this property is recomputed every time it
is called.

 
\begin{Verbatim}[commandchars=!@|,fontsize=\small,frame=single,label=Example]
  !gapprompt@gap>| !gapinput@D := Digraph(["a", "b", "c"], ["a", "b", "b"], ["b", "c", "a"]);|
  <immutable digraph with 3 vertices, 3 edges>
  !gapprompt@gap>| !gapinput@IsMultiDigraph(D);|
  false
  !gapprompt@gap>| !gapinput@D := DigraphFromDigraph6String("&Bug");|
  <immutable digraph with 3 vertices, 6 edges>
  !gapprompt@gap>| !gapinput@IsDuplicateFree(DigraphEdges(D));|
  true
  !gapprompt@gap>| !gapinput@IsMultiDigraph(D);|
  false
  !gapprompt@gap>| !gapinput@D := Digraph([[1, 2, 3, 2], [2, 1], [3]]);|
  <immutable multidigraph with 3 vertices, 7 edges>
  !gapprompt@gap>| !gapinput@IsDuplicateFree(DigraphEdges(D));|
  false
  !gapprompt@gap>| !gapinput@IsMultiDigraph(D);|
  true
  !gapprompt@gap>| !gapinput@D := DigraphMutableCopy(D);|
  <mutable multidigraph with 3 vertices, 7 edges>
  !gapprompt@gap>| !gapinput@IsMultiDigraph(D);|
  true
\end{Verbatim}
 }

 

\subsection{\textcolor{Chapter }{IsNonemptyDigraph}}
\logpage{[ 6, 2, 12 ]}\nobreak
\hyperdef{L}{X7E14AE1F7CA23141}{}
{\noindent\textcolor{FuncColor}{$\triangleright$\enspace\texttt{IsNonemptyDigraph({\mdseries\slshape digraph})\index{IsNonemptyDigraph@\texttt{IsNonemptyDigraph}}
\label{IsNonemptyDigraph}
}\hfill{\scriptsize (property)}}\\
\textbf{\indent Returns:\ }
\texttt{true} or \texttt{false}.



 Returns \texttt{true} if the digraph \mbox{\texttt{\mdseries\slshape digraph}} is nonempty, and \texttt{false} if it is not. A digraph is \emph{nonempty} if it has at least one edge.

 See also \texttt{IsEmptyDigraph} (\ref{IsEmptyDigraph}). 

 If the argument \mbox{\texttt{\mdseries\slshape digraph}} is mutable, then the return value of this property is recomputed every time it
is called.

 
\begin{Verbatim}[commandchars=!@|,fontsize=\small,frame=single,label=Example]
  !gapprompt@gap>| !gapinput@D := Digraph([[], []]);|
  <immutable empty digraph with 2 vertices>
  !gapprompt@gap>| !gapinput@IsNonemptyDigraph(D);|
  false
  !gapprompt@gap>| !gapinput@D := Digraph([[], [1]]);|
  <immutable digraph with 2 vertices, 1 edge>
  !gapprompt@gap>| !gapinput@IsNonemptyDigraph(D);|
  true
\end{Verbatim}
 }

 

\subsection{\textcolor{Chapter }{IsReflexiveDigraph}}
\logpage{[ 6, 2, 13 ]}\nobreak
\hyperdef{L}{X8462D8E2792B23F6}{}
{\noindent\textcolor{FuncColor}{$\triangleright$\enspace\texttt{IsReflexiveDigraph({\mdseries\slshape digraph})\index{IsReflexiveDigraph@\texttt{IsReflexiveDigraph}}
\label{IsReflexiveDigraph}
}\hfill{\scriptsize (property)}}\\
\textbf{\indent Returns:\ }
\texttt{true} or \texttt{false}.



 This property is \texttt{true} if the digraph \mbox{\texttt{\mdseries\slshape digraph}} is reflexive, and \texttt{false} if it is not. A digraph is \emph{reflexive} if it has a loop at every vertex. 

 

 If the argument \mbox{\texttt{\mdseries\slshape digraph}} is mutable, then the return value of this property is recomputed every time it
is called.

 
\begin{Verbatim}[commandchars=!@|,fontsize=\small,frame=single,label=Example]
  !gapprompt@gap>| !gapinput@D := Digraph([[1, 2], [2]]);|
  <immutable digraph with 2 vertices, 3 edges>
  !gapprompt@gap>| !gapinput@IsReflexiveDigraph(D);|
  true
  !gapprompt@gap>| !gapinput@D := Digraph([[3, 1], [4, 2], [3], [2, 1]]);|
  <immutable digraph with 4 vertices, 7 edges>
  !gapprompt@gap>| !gapinput@IsReflexiveDigraph(D);|
  false
\end{Verbatim}
 }

 

\subsection{\textcolor{Chapter }{IsSymmetricDigraph}}
\logpage{[ 6, 2, 14 ]}\nobreak
\hyperdef{L}{X81B3EA7887219860}{}
{\noindent\textcolor{FuncColor}{$\triangleright$\enspace\texttt{IsSymmetricDigraph({\mdseries\slshape digraph})\index{IsSymmetricDigraph@\texttt{IsSymmetricDigraph}}
\label{IsSymmetricDigraph}
}\hfill{\scriptsize (property)}}\\
\textbf{\indent Returns:\ }
\texttt{true} or \texttt{false}.



 This property is \texttt{true} if the digraph \mbox{\texttt{\mdseries\slshape digraph}} is symmetric, and \texttt{false} if it is not.

 A \emph{symmetric digraph} is one where for each non\texttt{\symbol{45}}loop edge, having source \texttt{u} and range \texttt{v}, there is a corresponding edge with source \texttt{v} and range \texttt{u}. If there are \texttt{n} edges with source \texttt{u} and range \texttt{v}, then there must be precisely \texttt{n} edges with source \texttt{v} and range \texttt{u}. In other words, a symmetric digraph has a symmetric adjacency matrix \texttt{AdjacencyMatrix} (\ref{AdjacencyMatrix}). 

 If the argument \mbox{\texttt{\mdseries\slshape digraph}} is mutable, then the return value of this property is recomputed every time it
is called.

 
\begin{Verbatim}[commandchars=!@|,fontsize=\small,frame=single,label=Example]
  !gapprompt@gap>| !gapinput@gr1 := Digraph([[2], [1, 3], [2, 3]]);|
  <immutable digraph with 3 vertices, 5 edges>
  !gapprompt@gap>| !gapinput@IsSymmetricDigraph(gr1);|
  true
  !gapprompt@gap>| !gapinput@adj1 := AdjacencyMatrix(gr1);;|
  !gapprompt@gap>| !gapinput@Display(adj1);|
  [ [  0,  1,  0 ],
    [  1,  0,  1 ],
    [  0,  1,  1 ] ]
  !gapprompt@gap>| !gapinput@adj1 = TransposedMat(adj1);|
  true
  !gapprompt@gap>| !gapinput@gr1 = DigraphReverse(gr1);|
  true
  !gapprompt@gap>| !gapinput@gr2 := Digraph([[2, 3], [1, 3], [2, 3]]);|
  <immutable digraph with 3 vertices, 6 edges>
  !gapprompt@gap>| !gapinput@IsSymmetricDigraph(gr2);|
  false
  !gapprompt@gap>| !gapinput@adj2 := AdjacencyMatrix(gr2);;|
  !gapprompt@gap>| !gapinput@Display(adj2);|
  [ [  0,  1,  1 ],
    [  1,  0,  1 ],
    [  0,  1,  1 ] ]
  !gapprompt@gap>| !gapinput@adj2 = TransposedMat(adj2);|
  false
\end{Verbatim}
 }

 

\subsection{\textcolor{Chapter }{IsTournament}}
\logpage{[ 6, 2, 15 ]}\nobreak
\hyperdef{L}{X7DD8D1A185EBE865}{}
{\noindent\textcolor{FuncColor}{$\triangleright$\enspace\texttt{IsTournament({\mdseries\slshape digraph})\index{IsTournament@\texttt{IsTournament}}
\label{IsTournament}
}\hfill{\scriptsize (property)}}\\
\textbf{\indent Returns:\ }
\texttt{true} or \texttt{false}.



 This property is \texttt{true} if the digraph \mbox{\texttt{\mdseries\slshape digraph}} is a tournament, and \texttt{false} if it is not. 

 A tournament is an orientation of a complete (undirected) graph. Specifically,
a tournament is a digraph which has a unique directed edge (of some
orientation) between any pair of distinct vertices, and no loops. 

 If the argument \mbox{\texttt{\mdseries\slshape digraph}} is mutable, then the return value of this property is recomputed every time it
is called.

 
\begin{Verbatim}[commandchars=!@|,fontsize=\small,frame=single,label=Example]
  !gapprompt@gap>| !gapinput@D := Digraph([[2, 3, 4], [3, 4], [4], []]);|
  <immutable digraph with 4 vertices, 6 edges>
  !gapprompt@gap>| !gapinput@IsTournament(D);|
  true
  !gapprompt@gap>| !gapinput@D := Digraph([[2], [1], [3]]);|
  <immutable digraph with 3 vertices, 3 edges>
  !gapprompt@gap>| !gapinput@IsTournament(D);|
  false
  !gapprompt@gap>| !gapinput@D := CycleDigraph(IsMutableDigraph, 3);|
  <mutable digraph with 3 vertices, 3 edges>
  !gapprompt@gap>| !gapinput@IsTournament(D);|
  true
  !gapprompt@gap>| !gapinput@DigraphRemoveEdge(D, 1, 2);|
  <mutable digraph with 3 vertices, 2 edges>
  !gapprompt@gap>| !gapinput@IsTournament(D);|
  false
\end{Verbatim}
 }

 

\subsection{\textcolor{Chapter }{IsTransitiveDigraph}}
\logpage{[ 6, 2, 16 ]}\nobreak
\hyperdef{L}{X7F0835667F29F0C0}{}
{\noindent\textcolor{FuncColor}{$\triangleright$\enspace\texttt{IsTransitiveDigraph({\mdseries\slshape digraph})\index{IsTransitiveDigraph@\texttt{IsTransitiveDigraph}}
\label{IsTransitiveDigraph}
}\hfill{\scriptsize (property)}}\\
\textbf{\indent Returns:\ }
\texttt{true} or \texttt{false}.



 This property is \texttt{true} if the digraph \mbox{\texttt{\mdseries\slshape digraph}} is transitive, and \texttt{false} if it is not. A digraph is \emph{transitive} if whenever \texttt{[ i, j ]} and \texttt{[ j, k ]} are edges of the digraph, then \texttt{[ i, k ]} is also an edge of the digraph. 

 Let $n$ be the number of vertices of an arbitrary digraph, and let $m$ be the number of edges. For general digraphs, the methods used for this
property use a version of the Floyd\texttt{\symbol{45}}Warshall algorithm, and
have complexity $O(n^3)$. However for digraphs which are topologically sortable [\texttt{DigraphTopologicalSort} (\ref{DigraphTopologicalSort})], then methods with complexity $O(m + n + m \cdot n)$ will be used when appropriate. 

 If the argument \mbox{\texttt{\mdseries\slshape digraph}} is mutable, then the return value of this property is recomputed every time it
is called.

 
\begin{Verbatim}[commandchars=!@|,fontsize=\small,frame=single,label=Example]
  !gapprompt@gap>| !gapinput@D := Digraph([[1, 2], [3], [3]]);|
  <immutable digraph with 3 vertices, 4 edges>
  !gapprompt@gap>| !gapinput@IsTransitiveDigraph(D);|
  false
  !gapprompt@gap>| !gapinput@gr2 := Digraph([[1, 2, 3], [3], [3]]);|
  <immutable digraph with 3 vertices, 5 edges>
  !gapprompt@gap>| !gapinput@IsTransitiveDigraph(gr2);|
  true
  !gapprompt@gap>| !gapinput@gr2 = DigraphTransitiveClosure(D);|
  true
  !gapprompt@gap>| !gapinput@gr3 := Digraph([[1, 2, 2, 3], [3, 3], [3]]);|
  <immutable multidigraph with 3 vertices, 7 edges>
  !gapprompt@gap>| !gapinput@IsTransitiveDigraph(gr3);|
  true
\end{Verbatim}
 }

 }

 
\section{\textcolor{Chapter }{Edge Weights}}\logpage{[ 6, 3, 0 ]}
\hyperdef{L}{X85A9843E80A9B590}{}
{
 

\subsection{\textcolor{Chapter }{EdgeWeights}}
\logpage{[ 6, 3, 1 ]}\nobreak
\hyperdef{L}{X8639092D8689333D}{}
{\noindent\textcolor{FuncColor}{$\triangleright$\enspace\texttt{EdgeWeights({\mdseries\slshape digraph})\index{EdgeWeights@\texttt{EdgeWeights}}
\label{EdgeWeights}
}\hfill{\scriptsize (attribute)}}\\
\noindent\textcolor{FuncColor}{$\triangleright$\enspace\texttt{EdgeWeightsMutableCopy({\mdseries\slshape digraph})\index{EdgeWeightsMutableCopy@\texttt{EdgeWeightsMutableCopy}}
\label{EdgeWeightsMutableCopy}
}\hfill{\scriptsize (operation)}}\\
\textbf{\indent Returns:\ }
A list of lists of integers, floats or rationals.



 \texttt{EdgeWeights} returns the list of lists of edge weights of the edges of the digraph \mbox{\texttt{\mdseries\slshape digraph}}.

 More specifically, \texttt{weights[i][j]} is the weight given to the \texttt{j}th edge from vertex \texttt{i}, according to the ordering of edges given by \texttt{OutNeighbours(digraph)[i]}.

 The function \texttt{EdgeWeights} returns an immutable list of immutable lists, whereas the function \texttt{EdgeWeightsMutableCopy} returns a copy of \texttt{EdgeWeights} which is a mutable list of mutable lists.

 The edge weights of a digraph cannot be computed and must be set either using \texttt{SetEdgeWeights} or \texttt{EdgeWeightedDigraph} (\ref{EdgeWeightedDigraph}).

 
\begin{Verbatim}[commandchars=!@|,fontsize=\small,frame=single,label=Example]
  !gapprompt@gap>| !gapinput@gr := EdgeWeightedDigraph([[2], [3], [1]], [[5], [10], [15]]);|
  <immutable edge-weighted digraph with 3 vertices, 3 edges>
  !gapprompt@gap>| !gapinput@EdgeWeights(gr);|
  [ [ 5 ], [ 10 ], [ 15 ] ]
  !gapprompt@gap>| !gapinput@a := EdgeWeightsMutableCopy(gr);|
  [ [ 5 ], [ 10 ], [ 15 ] ]
  !gapprompt@gap>| !gapinput@a[1][1] := 100;|
  100
  !gapprompt@gap>| !gapinput@a;|
  [ [ 100 ], [ 10 ], [ 15 ] ]
  !gapprompt@gap>| !gapinput@b := EdgeWeights(gr);|
  [ [ 5 ], [ 10 ], [ 15 ] ]
  !gapprompt@gap>| !gapinput@b[1][1] := 534;|
  Error, List Assignment: <list> must be a mutable list
\end{Verbatim}
 }

 

\subsection{\textcolor{Chapter }{EdgeWeightedDigraph}}
\logpage{[ 6, 3, 2 ]}\nobreak
\hyperdef{L}{X7E14055A84A52A07}{}
{\noindent\textcolor{FuncColor}{$\triangleright$\enspace\texttt{EdgeWeightedDigraph({\mdseries\slshape digraph, weights})\index{EdgeWeightedDigraph@\texttt{EdgeWeightedDigraph}}
\label{EdgeWeightedDigraph}
}\hfill{\scriptsize (function)}}\\
\textbf{\indent Returns:\ }
A digraph or \texttt{fail}



 The argument \mbox{\texttt{\mdseries\slshape digraph}} may be a digraph or a list of lists of integers, floats or rationals.

 \mbox{\texttt{\mdseries\slshape weights}} must be a list of lists of integers, floats or rationals of an equal size and
shape to \texttt{OutNeighbours(digraph)}, otherwise it will fail.

 This will create a digraph and set the EdgeWeights to \mbox{\texttt{\mdseries\slshape weights}}.

 See \texttt{EdgeWeights} (\ref{EdgeWeights}). 
\begin{Verbatim}[commandchars=!@|,fontsize=\small,frame=single,label=Example]
  !gapprompt@gap>| !gapinput@g := EdgeWeightedDigraph(Digraph([[2], [1]]), [[5], [15]]);|
  <immutable edge-weighted digraph with 2 vertices, 2 edges>
  !gapprompt@gap>| !gapinput@g := EdgeWeightedDigraph([[2], [1]], [[5], [15]]);|
  <immutable edge-weighted digraph with 2 vertices, 2 edges>
  !gapprompt@gap>| !gapinput@EdgeWeights(g);|
  [ [ 5 ], [ 15 ] ]
\end{Verbatim}
 }

 

\subsection{\textcolor{Chapter }{EdgeWeightedDigraphTotalWeight}}
\logpage{[ 6, 3, 3 ]}\nobreak
\hyperdef{L}{X810B2DD7794AFBE8}{}
{\noindent\textcolor{FuncColor}{$\triangleright$\enspace\texttt{EdgeWeightedDigraphTotalWeight({\mdseries\slshape digraph})\index{EdgeWeightedDigraphTotalWeight@\texttt{EdgeWeightedDigraphTotalWeight}}
\label{EdgeWeightedDigraphTotalWeight}
}\hfill{\scriptsize (attribute)}}\\
\textbf{\indent Returns:\ }
An integer, float or rational.



 If \mbox{\texttt{\mdseries\slshape digraph}} is a digraph with edge weights, then this attribute returns the sum of the
weights of its edges.

 If the argument \mbox{\texttt{\mdseries\slshape digraph}} is mutable, then the return value of this attribute is recomputed every time
it is called.

 See \texttt{EdgeWeights} (\ref{EdgeWeights}). 
\begin{Verbatim}[commandchars=!@|,fontsize=\small,frame=single,label=Example]
  !gapprompt@gap>| !gapinput@D := EdgeWeightedDigraph([[2], [1], [1, 2]],|
  !gapprompt@>| !gapinput@                            [[12], [5], [6, 9]]);|
  <immutable edge-weighted digraph with 3 vertices, 4 edges>
  !gapprompt@gap>| !gapinput@EdgeWeightedDigraphTotalWeight(D);|
  32
\end{Verbatim}
 }

 

\subsection{\textcolor{Chapter }{EdgeWeightedDigraphMinimumSpanningTree}}
\logpage{[ 6, 3, 4 ]}\nobreak
\hyperdef{L}{X7E5060AB7FB4F31C}{}
{\noindent\textcolor{FuncColor}{$\triangleright$\enspace\texttt{EdgeWeightedDigraphMinimumSpanningTree({\mdseries\slshape digraph})\index{EdgeWeightedDigraphMinimumSpanningTree@\texttt{Edge}\-\texttt{Weighted}\-\texttt{Digraph}\-\texttt{Minimum}\-\texttt{Spanning}\-\texttt{Tree}}
\label{EdgeWeightedDigraphMinimumSpanningTree}
}\hfill{\scriptsize (attribute)}}\\
\textbf{\indent Returns:\ }
A digraph.



 If \mbox{\texttt{\mdseries\slshape digraph}} is a connected digraph with edge weights, then this attribute returns a
digraph which is a minimum spanning tree of \mbox{\texttt{\mdseries\slshape digraph}}.

 A \emph{spanning tree} of a digraph is a subdigraph with the same vertices but a subset of its edges
that form an undirected tree. It is \emph{minimum} if it has the smallest possible total weight for a spanning tree of that
digraph.

 If the argument \mbox{\texttt{\mdseries\slshape digraph}} is mutable, then the return value of this attribute is recomputed every time
it is called.

 See \texttt{EdgeWeights} (\ref{EdgeWeights}), \texttt{EdgeWeightedDigraphTotalWeight} (\ref{EdgeWeightedDigraphTotalWeight}) and \texttt{IsConnectedDigraph} (\ref{IsConnectedDigraph}). 
\begin{Verbatim}[commandchars=!@|,fontsize=\small,frame=single,label=Example]
  !gapprompt@gap>| !gapinput@D := EdgeWeightedDigraph([[2], [1], [1, 2]],|
  !gapprompt@>| !gapinput@                            [[12], [5], [6, 9]]);|
  <immutable edge-weighted digraph with 3 vertices, 4 edges>
  !gapprompt@gap>| !gapinput@T := EdgeWeightedDigraphMinimumSpanningTree(D);|
  <immutable edge-weighted digraph with 3 vertices, 2 edges>
  !gapprompt@gap>| !gapinput@EdgeWeights(T);|
  [ [  ], [ 5 ], [ 6 ] ]
\end{Verbatim}
 }

 

\subsection{\textcolor{Chapter }{EdgeWeightedDigraphShortestPaths (for a digraph)}}
\logpage{[ 6, 3, 5 ]}\nobreak
\hyperdef{L}{X7EDB9C9A840F6A84}{}
{\noindent\textcolor{FuncColor}{$\triangleright$\enspace\texttt{EdgeWeightedDigraphShortestPaths({\mdseries\slshape digraph})\index{EdgeWeightedDigraphShortestPaths@\texttt{EdgeWeightedDigraphShortestPaths}!for a digraph}
\label{EdgeWeightedDigraphShortestPaths:for a digraph}
}\hfill{\scriptsize (attribute)}}\\
\noindent\textcolor{FuncColor}{$\triangleright$\enspace\texttt{EdgeWeightedDigraphShortestPaths({\mdseries\slshape digraph, source})\index{EdgeWeightedDigraphShortestPaths@\texttt{EdgeWeightedDigraphShortestPaths}!for a digraph and a pos int}
\label{EdgeWeightedDigraphShortestPaths:for a digraph and a pos int}
}\hfill{\scriptsize (operation)}}\\
\textbf{\indent Returns:\ }
A record.



 If \mbox{\texttt{\mdseries\slshape digraph}} is an edge\texttt{\symbol{45}}weighted digraph, this attribute returns a
record describing the paths of lowest total weight (the \emph{shortest paths}) connecting each pair of vertices. If the optional argument \mbox{\texttt{\mdseries\slshape source}} is specified and is a vertex of \mbox{\texttt{\mdseries\slshape digraph}}, then the output will only contain information on paths originating from that
vertex. 

 In the two\texttt{\symbol{45}}argument form, the value returned is a record
containing three components: \texttt{distances}, \texttt{parents} and \texttt{edges}. Each of these is a list of integers with one entry for each vertex \texttt{v} as follows: 

 
\begin{itemize}
\item  \texttt{distances[v]} is the total weight of the shortest path from \mbox{\texttt{\mdseries\slshape source}} to \texttt{v}. 
\item  \texttt{parents[v]} is the final vertex before \texttt{v} on the shortest path from \mbox{\texttt{\mdseries\slshape source}} to \texttt{v}. 
\item  \texttt{edges[v]} is the index of the edge of lowest weight going from \texttt{parents[v]} to \texttt{v}. 
\end{itemize}
 Using these three components together, you can find the shortest edge weighted
path to all other vertices from a starting vertex. 

 If no path exists from \mbox{\texttt{\mdseries\slshape source}} to \texttt{v}, then \texttt{parents[v]} and \texttt{edges[v]} will both be \texttt{fail}. The distance from \mbox{\texttt{\mdseries\slshape source}} to itself is considered to be 0, and so both \texttt{parents[\mbox{\texttt{\mdseries\slshape source}}]} and \texttt{edges[\mbox{\texttt{\mdseries\slshape source}}]} are \texttt{fail}. Edge weights can have negative values, but there is currently no implemented
method for this operation if a negative\texttt{\symbol{45}}weighted cycle
exists. 

 In the one\texttt{\symbol{45}}argument form, the value returned is also a
record containing components \texttt{distances}, \texttt{parents} and \texttt{edges}, but each of these will instead be a list of lists in which the \texttt{i}th entry is the list that corresponds to paths starting at \texttt{i}. 

 For a simple way of finding the shortest path between two specific vertices,
see \texttt{EdgeWeightedDigraphShortestPath} (\ref{EdgeWeightedDigraphShortestPath}). See also the non\texttt{\symbol{45}}weighted operation \texttt{DigraphShortestPath} (\ref{DigraphShortestPath}). 

 
\begin{Verbatim}[commandchars=!@|,fontsize=\small,frame=single,label=Example]
  !gapprompt@gap>| !gapinput@D := EdgeWeightedDigraph([[2, 3], [4], [4], []],|
  !gapprompt@>| !gapinput@                            [[5, 1], [6], [11], []]);|
  <immutable edge-weighted digraph with 4 vertices, 4 edges>
  !gapprompt@gap>| !gapinput@EdgeWeightedDigraphShortestPaths(D, 1);|
  rec( distances := [ 0, 5, 1, 11 ], edges := [ fail, 1, 2, 1 ], 
    parents := [ fail, 1, 1, 2 ] )
  !gapprompt@gap>| !gapinput@D := EdgeWeightedDigraph([[2], [3], [1]], [[1], [2], [3]]);|
  <immutable edge-weighted digraph with 3 vertices, 3 edges>
  !gapprompt@gap>| !gapinput@EdgeWeightedDigraphShortestPaths(D);|
  rec( distances := [ [ 0, 1, 3 ], [ 5, 0, 2 ], [ 3, 4, 0 ] ], 
    edges := [ [ fail, 1, 1 ], [ 1, fail, 1 ], [ 1, 1, fail ] ], 
    parents := [ [ fail, 1, 1 ], [ 2, fail, 2 ], [ 3, 3, fail ] ] )
\end{Verbatim}
 }

 

\subsection{\textcolor{Chapter }{EdgeWeightedDigraphShortestPath}}
\logpage{[ 6, 3, 6 ]}\nobreak
\hyperdef{L}{X7DAF81B779BD9165}{}
{\noindent\textcolor{FuncColor}{$\triangleright$\enspace\texttt{EdgeWeightedDigraphShortestPath({\mdseries\slshape digraph, source, dest})\index{EdgeWeightedDigraphShortestPath@\texttt{EdgeWeightedDigraphShortestPath}}
\label{EdgeWeightedDigraphShortestPath}
}\hfill{\scriptsize (operation)}}\\
\textbf{\indent Returns:\ }
A pair of lists, or \texttt{fail}.



 If \mbox{\texttt{\mdseries\slshape digraph}} is an edge\texttt{\symbol{45}}weighted digraph with vertices \mbox{\texttt{\mdseries\slshape source}} and \mbox{\texttt{\mdseries\slshape dest}}, this operation returns a directed path from \mbox{\texttt{\mdseries\slshape source}} to \mbox{\texttt{\mdseries\slshape dest}} with the smallest possible total weight. The output is a pair of lists \texttt{[v, a]} of the form described in \texttt{DigraphPath} (\ref{DigraphPath}).

 If $\mbox{\texttt{\mdseries\slshape source}} = \mbox{\texttt{\mdseries\slshape dest}}$ or no path exists, then \texttt{fail} is returned.

 If \mbox{\texttt{\mdseries\slshape digraph}} contains a negative\texttt{\symbol{45}}weighted cycle, then there is currently
no applicable method for this attribute. 

 See \texttt{EdgeWeightedDigraphShortestPaths} (\ref{EdgeWeightedDigraphShortestPaths:for a digraph}). See also the non\texttt{\symbol{45}}weighted operation \texttt{DigraphShortestPath} (\ref{DigraphShortestPath}). 

 
\begin{Verbatim}[commandchars=!@|,fontsize=\small,frame=single,label=Example]
  !gapprompt@gap>| !gapinput@D := EdgeWeightedDigraph([[2, 3], [4], [4], []],|
  !gapprompt@>| !gapinput@                            [[5, 1], [6], [11], []]);|
  <immutable edge-weighted digraph with 4 vertices, 4 edges>
  !gapprompt@gap>| !gapinput@EdgeWeightedDigraphShortestPath(D, 1, 4);|
  [ [ 1, 2, 4 ], [ 1, 1 ] ]
  !gapprompt@gap>| !gapinput@EdgeWeightedDigraphShortestPath(D, 3, 2);|
  fail
\end{Verbatim}
 }

 

\subsection{\textcolor{Chapter }{DigraphMaximumFlow}}
\logpage{[ 6, 3, 7 ]}\nobreak
\hyperdef{L}{X869EFE9F818C6E22}{}
{\noindent\textcolor{FuncColor}{$\triangleright$\enspace\texttt{DigraphMaximumFlow({\mdseries\slshape digraph, start, destination})\index{DigraphMaximumFlow@\texttt{DigraphMaximumFlow}}
\label{DigraphMaximumFlow}
}\hfill{\scriptsize (attribute)}}\\
\textbf{\indent Returns:\ }
A list of lists of integers.



 If \mbox{\texttt{\mdseries\slshape digraph}} is an edge\texttt{\symbol{45}}weighted digraph with vertices \mbox{\texttt{\mdseries\slshape start}} and \mbox{\texttt{\mdseries\slshape destination}}, this returns a record representing the maximum flow from \mbox{\texttt{\mdseries\slshape start}} to \mbox{\texttt{\mdseries\slshape destination}} in the digraph. 

 A \emph{flow} is a function from the weighted edges of \mbox{\texttt{\mdseries\slshape digraph}} to the positive real numbers, such that: 
\begin{itemize}
\item  Each edge's flow is no more than its weight; 
\item  For each vertex other than \mbox{\texttt{\mdseries\slshape start}} and \mbox{\texttt{\mdseries\slshape destination}}, the sum of flows for all incoming edges is equal to the sum of flows for all
outgoing edges; 
\item  The sum of flows of edges leaving \mbox{\texttt{\mdseries\slshape start}} is equal to the sum of flows of edges entering \mbox{\texttt{\mdseries\slshape destination}} (this sum is denoted $M$). 
\end{itemize}
 A \emph{maximum flow} is a flow that maximises the value of $M$. 

 The flow is represented as a list of lists where each entry is a number
representing the flow on the edge in the corresponding position in \texttt{OutNeighbours(\mbox{\texttt{\mdseries\slshape digraph}})}. Note that the value $M$ of the flow can be found with \texttt{Sum(DigraphMaximumFlow(\mbox{\texttt{\mdseries\slshape digraph}}, \mbox{\texttt{\mdseries\slshape start}}, \mbox{\texttt{\mdseries\slshape destination}})[\mbox{\texttt{\mdseries\slshape start}}])}. 

 This attribute is computed by an implementation of the
push{\textendash}relabel maximum flow algorithm, which has time complexity $O(v^2 e)$ where $v$ is the number of vertices of the digraph, and $e$ is the number of edges. 

 See \texttt{EdgeWeights} (\ref{EdgeWeights}). 
\begin{Verbatim}[commandchars=!@|,fontsize=\small,frame=single,label=Example]
  !gapprompt@gap>| !gapinput@g := EdgeWeightedDigraph([[2, 2], [3], []], [[3, 2], [1], []]);|
  <immutable edge-weighted multidigraph with 3 vertices, 3 edges>
  !gapprompt@gap>| !gapinput@flow := DigraphMaximumFlow(g, 1, 3);|
  [ [ 1, 0 ], [ 1 ], [  ] ]
  !gapprompt@gap>| !gapinput@Sum(flow[1]);|
  1
\end{Verbatim}
 }

 

\subsection{\textcolor{Chapter }{DigraphMinimumCut}}
\logpage{[ 6, 3, 8 ]}\nobreak
\hyperdef{L}{X7FF8935482954667}{}
{\noindent\textcolor{FuncColor}{$\triangleright$\enspace\texttt{DigraphMinimumCut({\mdseries\slshape digraph, s, t})\index{DigraphMinimumCut@\texttt{DigraphMinimumCut}}
\label{DigraphMinimumCut}
}\hfill{\scriptsize (attribute)}}\\
\textbf{\indent Returns:\ }
A list of lists of integers.



 If \mbox{\texttt{\mdseries\slshape digraph}} is an edge\texttt{\symbol{45}}weighted digraph with distinct vertices \mbox{\texttt{\mdseries\slshape s}} and \mbox{\texttt{\mdseries\slshape t}}, this function returns a list of two lists representing the components of the
minimum $s$\texttt{\symbol{45}}$t$ cut of \mbox{\texttt{\mdseries\slshape digraph}}. 

 An \emph{$s$\texttt{\symbol{45}}$t$ cut} is a partition of the vertices $\{ S, T \}$ such that \mbox{\texttt{\mdseries\slshape s}} is in $S$ and \mbox{\texttt{\mdseries\slshape t}} is in $T$. The \emph{capacity} of an $s$\texttt{\symbol{45}}$t$ cut is the sum of the weights of every edge whose source is in $S$ and whose range is in $T$. The minimum $s$\texttt{\symbol{45}}$t$ cut is the $s$\texttt{\symbol{45}}$t$ cut whose capacity is the smallest possible.

 This attribute is computed by using \texttt{DigraphMaximumFlow} (\ref{DigraphMaximumFlow}) and the max\texttt{\symbol{45}}cut min\texttt{\symbol{45}}flow theorem.

 See \texttt{EdgeWeights} (\ref{EdgeWeights}). 
\begin{Verbatim}[commandchars=!@|,fontsize=\small,frame=single,label=Example]
  !gapprompt@gap>| !gapinput@g := EdgeWeightedDigraph([[2, 2], [3], []], [[3, 2], [1], []]);|
  <immutable edge-weighted multidigraph with 3 vertices, 3 edges>
  !gapprompt@gap>| !gapinput@DigraphMinimumCut(g, 1, 3);|
  [ [ 1, 2 ], [3] ]
\end{Verbatim}
 }

 

\subsection{\textcolor{Chapter }{DigraphMinimumCutSet}}
\logpage{[ 6, 3, 9 ]}\nobreak
\hyperdef{L}{X854C6CFF82C57FF8}{}
{\noindent\textcolor{FuncColor}{$\triangleright$\enspace\texttt{DigraphMinimumCutSet({\mdseries\slshape digraph, s, t})\index{DigraphMinimumCutSet@\texttt{DigraphMinimumCutSet}}
\label{DigraphMinimumCutSet}
}\hfill{\scriptsize (attribute)}}\\
\textbf{\indent Returns:\ }
A list of lists integers.



 If \mbox{\texttt{\mdseries\slshape digraph}} is an edge\texttt{\symbol{45}}weighted digraph with distinct vertices \mbox{\texttt{\mdseries\slshape s}} and \mbox{\texttt{\mdseries\slshape t}}, this function returns a list of lists of integers representing the minimum $s$\texttt{\symbol{45}}$t$ cut of \mbox{\texttt{\mdseries\slshape digraph}}. 

 An \emph{$s$\texttt{\symbol{45}}$t$ cut} is a partition of the vertices $\{ S, T \}$ such that \mbox{\texttt{\mdseries\slshape s}} is in $S$ and \mbox{\texttt{\mdseries\slshape t}} is in $T$. The \emph{cut set} corresponding to this cut is the set of edges whose source is in $S$ and whose range is in $T$. The minimum $s$\texttt{\symbol{45}}$t$ cut set is the $s$\texttt{\symbol{45}}$t$ cut set such that the sum of the weights of the edges in the cut set is the
smallest possible.

 This attribute is computed by using \texttt{DigraphMinimumCut} (\ref{DigraphMinimumCut}).

 See \texttt{EdgeWeights} (\ref{EdgeWeights}). 
\begin{Verbatim}[commandchars=!@|,fontsize=\small,frame=single,label=Example]
  !gapprompt@gap>| !gapinput@g := EdgeWeightedDigraph([[2, 2], [3], []], [[3, 2], [1], []]);|
  <immutable edge-weighted multidigraph with 3 vertices, 3 edges>
  !gapprompt@gap>| !gapinput@DigraphMinimumCutSet(g, 1, 3);|
  [ [ 2, 3 ] ]
\end{Verbatim}
 }

 

\subsection{\textcolor{Chapter }{RandomUniqueEdgeWeightedDigraph}}
\logpage{[ 6, 3, 10 ]}\nobreak
\hyperdef{L}{X8194E7078167A9E4}{}
{\noindent\textcolor{FuncColor}{$\triangleright$\enspace\texttt{RandomUniqueEdgeWeightedDigraph({\mdseries\slshape [filt, ]n[, p]})\index{RandomUniqueEdgeWeightedDigraph@\texttt{RandomUniqueEdgeWeightedDigraph}}
\label{RandomUniqueEdgeWeightedDigraph}
}\hfill{\scriptsize (operation)}}\\
\textbf{\indent Returns:\ }
An edge\texttt{\symbol{45}}weighted digraph.



 This operation returns a random edge\texttt{\symbol{45}}weighted digraph.

 Its behaviour is the same as that of \texttt{RandomDigraph} (\ref{RandomDigraph}) but the returned digraph will additionally have the \texttt{EdgeWeights} (\ref{EdgeWeights}) attribute populated with random unique weights from the set \texttt{[1 .. m]} where \texttt{m} is the number of edges in the digraph.

 If the optional first argument \mbox{\texttt{\mdseries\slshape filt}} is present, then this should specify the category or representation the
digraph being created will belong to. For example, if \mbox{\texttt{\mdseries\slshape filt}} is \texttt{IsMutableDigraph} (\ref{IsMutableDigraph}), then the digraph being created will be mutable, if \mbox{\texttt{\mdseries\slshape filt}} is \texttt{IsImmutableDigraph} (\ref{IsImmutableDigraph}), then the digraph will be immutable. If the optional first argument \mbox{\texttt{\mdseries\slshape filt}} is not present, then \texttt{IsImmutableDigraph} (\ref{IsImmutableDigraph}) is used by default.

 If \mbox{\texttt{\mdseries\slshape n}} is a non\texttt{\symbol{45}}negative integer, then the returned digraph will
have \mbox{\texttt{\mdseries\slshape n}} vertices. If the optional second argument \mbox{\texttt{\mdseries\slshape p}} is a float with value $0 \leq $ \mbox{\texttt{\mdseries\slshape  p }} $ \leq 1$, then an edge will exist between each pair of vertices with probability
approximately \mbox{\texttt{\mdseries\slshape p}}. If \mbox{\texttt{\mdseries\slshape p}} is not specified, then a random probability will be assumed (chosen with
uniform probability).

 For more information on the arguments and behaviour of this operation, see \texttt{RandomDigraph} (\ref{RandomDigraph}). 
\begin{Verbatim}[commandchars=!@|,fontsize=\small,frame=single,label=Example]
  !gapprompt@gap>| !gapinput@RandomUniqueEdgeWeightedDigraph(5);|
  <immutable edge-weighted digraph with 5 vertices, 21 edges>
  !gapprompt@gap>| !gapinput@RandomUniqueEdgeWeightedDigraph(5, 1 / 2);|
  <immutable edge-weighted digraph with 5 vertices, 14 edges>
  !gapprompt@gap>| !gapinput@RandomUniqueEdgeWeightedDigraph(IsEulerianDigraph, 5, 1 / 3);|
  <immutable edge-weighted digraph with 5 vertices, 6 edges>
\end{Verbatim}
 }

 }

 
\section{\textcolor{Chapter }{Orders}}\logpage{[ 6, 4, 0 ]}
\hyperdef{L}{X86424F167BD4F629}{}
{
 

\subsection{\textcolor{Chapter }{IsPreorderDigraph}}
\logpage{[ 6, 4, 1 ]}\nobreak
\hyperdef{L}{X8617726C7829F796}{}
{\noindent\textcolor{FuncColor}{$\triangleright$\enspace\texttt{IsPreorderDigraph({\mdseries\slshape digraph})\index{IsPreorderDigraph@\texttt{IsPreorderDigraph}}
\label{IsPreorderDigraph}
}\hfill{\scriptsize (property)}}\\
\noindent\textcolor{FuncColor}{$\triangleright$\enspace\texttt{IsQuasiorderDigraph({\mdseries\slshape digraph})\index{IsQuasiorderDigraph@\texttt{IsQuasiorderDigraph}}
\label{IsQuasiorderDigraph}
}\hfill{\scriptsize (property)}}\\
\textbf{\indent Returns:\ }
\texttt{true} or \texttt{false}.



 A digraph is a preorder digraph if and only if the digraph satisfies both \texttt{IsReflexiveDigraph} (\ref{IsReflexiveDigraph}) and \texttt{IsTransitiveDigraph} (\ref{IsTransitiveDigraph}). A preorder digraph (or quasiorder digraph) \mbox{\texttt{\mdseries\slshape digraph}} corresponds to the preorder relation $\leq$ defined by $x \leq y$ if and only if \texttt{[x, y]} is an edge of \mbox{\texttt{\mdseries\slshape digraph}}. 

 If the argument \mbox{\texttt{\mdseries\slshape digraph}} is mutable, then the return value of this property is recomputed every time it
is called.

 
\begin{Verbatim}[commandchars=!@|,fontsize=\small,frame=single,label=Example]
  !gapprompt@gap>| !gapinput@D := Digraph([[1], [2, 3], [2, 3]]);|
  <immutable digraph with 3 vertices, 5 edges>
  !gapprompt@gap>| !gapinput@IsPreorderDigraph(D);|
  true
  !gapprompt@gap>| !gapinput@D := Digraph([[1 .. 4], [1 .. 4], [1 .. 4], [1 .. 4]]);|
  <immutable digraph with 4 vertices, 16 edges>
  !gapprompt@gap>| !gapinput@IsPreorderDigraph(D);|
  true
  !gapprompt@gap>| !gapinput@D := Digraph([[2], [3], [4], [5], [1]]);|
  <immutable digraph with 5 vertices, 5 edges>
  !gapprompt@gap>| !gapinput@IsPreorderDigraph(D);|
  false
  !gapprompt@gap>| !gapinput@D := Digraph([[1], [1, 2], [2, 3]]);|
  <immutable digraph with 3 vertices, 5 edges>
  !gapprompt@gap>| !gapinput@IsQuasiorderDigraph(D);|
  false
\end{Verbatim}
 }

 

\subsection{\textcolor{Chapter }{IsPartialOrderDigraph}}
\logpage{[ 6, 4, 2 ]}\nobreak
\hyperdef{L}{X82BAE6D37D49A145}{}
{\noindent\textcolor{FuncColor}{$\triangleright$\enspace\texttt{IsPartialOrderDigraph({\mdseries\slshape digraph})\index{IsPartialOrderDigraph@\texttt{IsPartialOrderDigraph}}
\label{IsPartialOrderDigraph}
}\hfill{\scriptsize (property)}}\\
\textbf{\indent Returns:\ }
\texttt{true} or \texttt{false}.



 A digraph is a partial order digraph if and only if the digraph satisfies all
of \texttt{IsReflexiveDigraph} (\ref{IsReflexiveDigraph}), \texttt{IsAntisymmetricDigraph} (\ref{IsAntisymmetricDigraph}) and \texttt{IsTransitiveDigraph} (\ref{IsTransitiveDigraph}). A partial order \mbox{\texttt{\mdseries\slshape digraph}} corresponds to the partial order relation $\leq$ defined by $x \leq y$ if and only if \texttt{[x, y]} is an edge of \mbox{\texttt{\mdseries\slshape digraph}}. 

 If the argument \mbox{\texttt{\mdseries\slshape digraph}} is mutable, then the return value of this property is recomputed every time it
is called.

 
\begin{Verbatim}[commandchars=!@|,fontsize=\small,frame=single,label=Example]
  !gapprompt@gap>| !gapinput@D := Digraph([[1, 3], [2, 3], [3]]);|
  <immutable digraph with 3 vertices, 5 edges>
  !gapprompt@gap>| !gapinput@IsPartialOrderDigraph(D);|
  true
  !gapprompt@gap>| !gapinput@D := CycleDigraph(5);|
  <immutable cycle digraph with 5 vertices>
  !gapprompt@gap>| !gapinput@IsPartialOrderDigraph(D);|
  false
  !gapprompt@gap>| !gapinput@D := Digraph([[1, 1], [1, 1, 2], [3], [3, 3, 4, 4]]);|
  <immutable multidigraph with 4 vertices, 10 edges>
  !gapprompt@gap>| !gapinput@IsPartialOrderDigraph(D);|
  true
\end{Verbatim}
 }

 

\subsection{\textcolor{Chapter }{IsMeetSemilatticeDigraph}}
\logpage{[ 6, 4, 3 ]}\nobreak
\hyperdef{L}{X78D3E17B7F737516}{}
{\noindent\textcolor{FuncColor}{$\triangleright$\enspace\texttt{IsMeetSemilatticeDigraph({\mdseries\slshape digraph})\index{IsMeetSemilatticeDigraph@\texttt{IsMeetSemilatticeDigraph}}
\label{IsMeetSemilatticeDigraph}
}\hfill{\scriptsize (property)}}\\
\noindent\textcolor{FuncColor}{$\triangleright$\enspace\texttt{IsJoinSemilatticeDigraph({\mdseries\slshape digraph})\index{IsJoinSemilatticeDigraph@\texttt{IsJoinSemilatticeDigraph}}
\label{IsJoinSemilatticeDigraph}
}\hfill{\scriptsize (property)}}\\
\noindent\textcolor{FuncColor}{$\triangleright$\enspace\texttt{IsLatticeDigraph({\mdseries\slshape digraph})\index{IsLatticeDigraph@\texttt{IsLatticeDigraph}}
\label{IsLatticeDigraph}
}\hfill{\scriptsize (property)}}\\
\textbf{\indent Returns:\ }
\texttt{true} or \texttt{false}.



 \texttt{IsMeetSemilatticeDigraph} returns \texttt{true} if the digraph \mbox{\texttt{\mdseries\slshape digraph}} is a meet semilattice; \texttt{IsJoinSemilatticeDigraph} returns \texttt{true} if the digraph \mbox{\texttt{\mdseries\slshape digraph}} is a join semilattice; and \texttt{IsLatticeDigraph} returns \texttt{true} if the digraph \mbox{\texttt{\mdseries\slshape digraph}} is both a meet and a join semilattice. 

 For a partial order digraph \texttt{IsPartialOrderDigraph} (\ref{IsPartialOrderDigraph}) the corresponding partial order is the relation $\leq$, defined by $x \leq y$ if and only if \texttt{[x, y]} is an edge. A digraph is a \emph{meet semilattice} if it is a partial order and every pair of vertices has a greatest lower bound
(meet) with respect to the aforementioned relation. A \emph{join semilattice} is a partial order where every pair of vertices has a least upper bound (join)
with respect to the relation. 

 If the argument \mbox{\texttt{\mdseries\slshape digraph}} is mutable, then the return value of this property is recomputed every time it
is called.

 
\begin{Verbatim}[commandchars=!@|,fontsize=\small,frame=single,label=Example]
  !gapprompt@gap>| !gapinput@D := Digraph([[1, 3], [2, 3], [3]]);|
  <immutable digraph with 3 vertices, 5 edges>
  !gapprompt@gap>| !gapinput@IsMeetSemilatticeDigraph(D);|
  false
  !gapprompt@gap>| !gapinput@IsJoinSemilatticeDigraph(D);|
  true
  !gapprompt@gap>| !gapinput@IsLatticeDigraph(D);|
  false
  !gapprompt@gap>| !gapinput@D := Digraph([[1], [2], [1 .. 3]]);|
  <immutable digraph with 3 vertices, 5 edges>
  !gapprompt@gap>| !gapinput@IsJoinSemilatticeDigraph(D);|
  false
  !gapprompt@gap>| !gapinput@IsMeetSemilatticeDigraph(D);|
  true
  !gapprompt@gap>| !gapinput@IsLatticeDigraph(D);|
  false
  !gapprompt@gap>| !gapinput@D := Digraph([[1 .. 4], [2, 4], [3, 4], [4]]);|
  <immutable digraph with 4 vertices, 9 edges>
  !gapprompt@gap>| !gapinput@IsMeetSemilatticeDigraph(D);|
  true
  !gapprompt@gap>| !gapinput@IsJoinSemilatticeDigraph(D);|
  true
  !gapprompt@gap>| !gapinput@IsLatticeDigraph(D);|
  true
\end{Verbatim}
 }

 

\subsection{\textcolor{Chapter }{DigraphMeetTable}}
\logpage{[ 6, 4, 4 ]}\nobreak
\hyperdef{L}{X84F4711F7FA36848}{}
{\noindent\textcolor{FuncColor}{$\triangleright$\enspace\texttt{DigraphMeetTable({\mdseries\slshape digraph})\index{DigraphMeetTable@\texttt{DigraphMeetTable}}
\label{DigraphMeetTable}
}\hfill{\scriptsize (attribute)}}\\
\noindent\textcolor{FuncColor}{$\triangleright$\enspace\texttt{DigraphJoinTable({\mdseries\slshape digraph})\index{DigraphJoinTable@\texttt{DigraphJoinTable}}
\label{DigraphJoinTable}
}\hfill{\scriptsize (attribute)}}\\
\textbf{\indent Returns:\ }
A matrix or \texttt{fail}.



 \texttt{DigraphMeetTable} returns the \emph{meet table} of \mbox{\texttt{\mdseries\slshape digraph}} if \mbox{\texttt{\mdseries\slshape digraph}} is a meet semilattice digraph \texttt{IsMeetSemilatticeDigraph} (\ref{IsMeetSemilatticeDigraph}). Similarly, \texttt{DigraphJoinTable} returns the \emph{join table} of \mbox{\texttt{\mdseries\slshape digraph}} if \mbox{\texttt{\mdseries\slshape digraph}} is a join semilattice digraph. \texttt{IsJoinSemilatticeDigraph} (\ref{IsJoinSemilatticeDigraph}).

 When \mbox{\texttt{\mdseries\slshape digraph}} is a meet semilattice digraph with $n$ vertices, the \emph{meet table} of \mbox{\texttt{\mdseries\slshape digraph}} is the matrix $A$ such that the $(i, j)$ entry of $A$ is the meet of vertices $i$ and $j$.

 Similarly, when \mbox{\texttt{\mdseries\slshape digraph}} is a join semilattice digraph with $n$ vertices, the \emph{join table} of \mbox{\texttt{\mdseries\slshape digraph}} is the matrix $A$ such that the $(i, j)$ entry of $A$ is the join of vertices $i$ and $j$. Otherwise, each function returns \texttt{fail}. 
\begin{Verbatim}[commandchars=!@|,fontsize=\small,frame=single,label=Example]
  !gapprompt@gap>| !gapinput@D := Digraph([[1, 2, 3, 4], [2, 4], [3, 4], [4]]);|
  <immutable digraph with 4 vertices, 9 edges>
  !gapprompt@gap>| !gapinput@IsJoinSemilatticeDigraph(D);|
  true
  !gapprompt@gap>| !gapinput@Display(DigraphJoinTable(D));|
  [ [  1,  2,  3,  4 ],
    [  2,  2,  4,  4 ],
    [  3,  4,  3,  4 ],
    [  4,  4,  4,  4 ] ]
  !gapprompt@gap>| !gapinput@D := CycleDigraph(5);|
  <immutable cycle digraph with 5 vertices>
  !gapprompt@gap>| !gapinput@DigraphJoinTable(D);|
  fail
\end{Verbatim}
 }

 

\subsection{\textcolor{Chapter }{IsOrderIdeal (for a digraph and list)}}
\logpage{[ 6, 4, 5 ]}\nobreak
\hyperdef{L}{X786213437E99065B}{}
{\noindent\textcolor{FuncColor}{$\triangleright$\enspace\texttt{IsOrderIdeal({\mdseries\slshape D, subset})\index{IsOrderIdeal@\texttt{IsOrderIdeal}!for a digraph and list}
\label{IsOrderIdeal:for a digraph and list}
}\hfill{\scriptsize (operation)}}\\
\textbf{\indent Returns:\ }
\texttt{true} or \texttt{false}.



 This function returns \texttt{true} if the specified subset is "downwards" closed, i.e. contains every vertex less
than the given vertices in the order defined by \mbox{\texttt{\mdseries\slshape D}}. The function can only be used on digraphs satisfying \texttt{IsPartialOrderDigraph} (\ref{IsPartialOrderDigraph}) and will throw an error if passed a digraph that is not a partial order
digraph. 
\begin{Verbatim}[commandchars=!@|,fontsize=\small,frame=single,label=Example]
  !gapprompt@gap>| !gapinput@D1 := Digraph([[1, 3], [2, 3], [3]]);|
  <immutable digraph with 3 vertices, 5 edges>
  !gapprompt@gap>| !gapinput@IsOrderIdeal(D1, [1, 2, 3]);|
  true
  !gapprompt@gap>| !gapinput@D2 := DigraphDisjointUnion(D1, D1);|
  <immutable digraph with 6 vertices, 10 edges>
  !gapprompt@gap>| !gapinput@IsOrderIdeal(D2, [1, 2, 3]);|
  true
  !gapprompt@gap>| !gapinput@IsOrderIdeal(D2, [4, 5, 6]);|
  true
  !gapprompt@gap>| !gapinput@IsOrderIdeal(D2, [1, 5, 6]);|
  false
\end{Verbatim}
 }

 

\subsection{\textcolor{Chapter }{IsOrderFilter (for a digraph and a list)}}
\logpage{[ 6, 4, 6 ]}\nobreak
\hyperdef{L}{X7A1D8F4582F8EA53}{}
{\noindent\textcolor{FuncColor}{$\triangleright$\enspace\texttt{IsOrderFilter({\mdseries\slshape D, subset})\index{IsOrderFilter@\texttt{IsOrderFilter}!for a digraph and a list}
\label{IsOrderFilter:for a digraph and a list}
}\hfill{\scriptsize (operation)}}\\
\textbf{\indent Returns:\ }
 \texttt{true} or \texttt{false}.



 This function returns \texttt{true} if the specified subset is upwards closed. A subset is upwards closed if it
contains every vertex greater than the given vertices in the order defined by \mbox{\texttt{\mdseries\slshape D}}. The function can only be used on digraphs satisfying \texttt{IsPartialOrderDigraph} (\ref{IsPartialOrderDigraph}) and will throw an error if passed a digraph that is not a partial order
digraph. 
\begin{Verbatim}[commandchars=!@|,fontsize=\small,frame=single,label=Example]
  !gapprompt@gap>| !gapinput@D1 := DigraphByEdges(|
  !gapprompt@>| !gapinput@[[1, 2], [2, 3], [1, 3], [1, 1], [2, 2], [3, 3]]);|
  <immutable digraph with 3 vertices, 6 edges>
  !gapprompt@gap>| !gapinput@IsOrderFilter(D1, [2, 3]);|
  false
  !gapprompt@gap>| !gapinput@IsOrderFilter(D1, [1, 2]);|
  true    
      
\end{Verbatim}
 }

 

\subsection{\textcolor{Chapter }{IsUpperSemimodularDigraph}}
\logpage{[ 6, 4, 7 ]}\nobreak
\hyperdef{L}{X7DEDCD177E5AE824}{}
{\noindent\textcolor{FuncColor}{$\triangleright$\enspace\texttt{IsUpperSemimodularDigraph({\mdseries\slshape D})\index{IsUpperSemimodularDigraph@\texttt{IsUpperSemimodularDigraph}}
\label{IsUpperSemimodularDigraph}
}\hfill{\scriptsize (property)}}\\
\noindent\textcolor{FuncColor}{$\triangleright$\enspace\texttt{IsLowerSemimodularDigraph({\mdseries\slshape D})\index{IsLowerSemimodularDigraph@\texttt{IsLowerSemimodularDigraph}}
\label{IsLowerSemimodularDigraph}
}\hfill{\scriptsize (property)}}\\
\textbf{\indent Returns:\ }
\texttt{true} or \texttt{false}.



 \texttt{IsUpperSemimodularDigraph} returns \texttt{true} if the digraph \mbox{\texttt{\mdseries\slshape D}} represents an upper semimodular lattice and \texttt{false} if it does not. Similarly, \texttt{IsLowerSemimodularDigraph} returns \texttt{true} if \mbox{\texttt{\mdseries\slshape D}} represents a lower semimodular lattice and \texttt{false} if it does not. 

 In a lattice we say that a vertex \texttt{a} \emph{covers} a vertex \texttt{b} if \texttt{a} is greater than \texttt{b}, and there are no further vertices between \texttt{a} and \texttt{b}. A lattice is \emph{upper semimodular} if whenever the meet of \texttt{a} and \texttt{b} is covered by \texttt{a}, the join of \texttt{a} and \texttt{b} covers \texttt{b}. \emph{Lower semimodularity} is defined analogously. 

 See also \texttt{IsLatticeDigraph} (\ref{IsLatticeDigraph}), \texttt{NonUpperSemimodularPair} (\ref{NonUpperSemimodularPair}), and \texttt{NonLowerSemimodularPair} (\ref{NonLowerSemimodularPair}). If the argument \mbox{\texttt{\mdseries\slshape digraph}} is mutable, then the return value of this property is recomputed every time it
is called.

 
\begin{Verbatim}[commandchars=!|E,fontsize=\small,frame=single,label=Example]
  !gapprompt|gap>E !gapinput|D := DigraphFromDigraph6String(E
  !gapprompt|>E !gapinput|"&M~~sc`lYUZO__KIBboC_@h?U_?_GL?A_?c");E
  <immutable digraph with 14 vertices, 66 edges>
  !gapprompt|gap>E !gapinput|IsUpperSemimodularDigraph(D);E
  true
  !gapprompt|gap>E !gapinput|IsLowerSemimodularDigraph(D);E
  false
\end{Verbatim}
 }

 

\subsection{\textcolor{Chapter }{IsDistributiveLatticeDigraph}}
\logpage{[ 6, 4, 8 ]}\nobreak
\hyperdef{L}{X7B566B4A86E18F45}{}
{\noindent\textcolor{FuncColor}{$\triangleright$\enspace\texttt{IsDistributiveLatticeDigraph({\mdseries\slshape digraph})\index{IsDistributiveLatticeDigraph@\texttt{IsDistributiveLatticeDigraph}}
\label{IsDistributiveLatticeDigraph}
}\hfill{\scriptsize (property)}}\\
\textbf{\indent Returns:\ }
\texttt{true} or \texttt{false}.



 \texttt{IsDistributiveLatticeDigraph} returns \texttt{true} if the digraph \mbox{\texttt{\mdseries\slshape digraph}} is a distributive lattice digraph.

 A \emph{distributive lattice digraph} is a lattice digraph (\texttt{IsLatticeDigraph} (\ref{IsLatticeDigraph})) which is distributive. That is to say, the functions \texttt{PartialOrderDigraphMeetOfVertices} (\ref{PartialOrderDigraphMeetOfVertices}) and \texttt{PartialOrderDigraphJoinOfVertices} (\ref{PartialOrderDigraphJoinOfVertices}) distribute over each other.

 Equivalently, a distributive lattice digraph is a lattice digraph in which the \emph{lattice digraphs} representing $M3$ and $N5$ are not embeddable as lattices (see \href{https://en.wikipedia.org/wiki/Distributive_lattice} {\texttt{https://en.wikipedia.org/wiki/Distributive{\textunderscore}lattice}} and \texttt{IsLatticeEmbedding} (\ref{IsLatticeEmbedding:for digraphs and a permutation or transformation})).

  If the argument \mbox{\texttt{\mdseries\slshape digraph}} is mutable, then the return value of this property is recomputed every time it
is called.

 
\begin{Verbatim}[commandchars=!@|,fontsize=\small,frame=single,label=Example]
  !gapprompt@gap>| !gapinput@D := Digraph([[2, 3], [4], [4], []]);|
  <immutable digraph with 4 vertices, 4 edges>
  !gapprompt@gap>| !gapinput@D := DigraphReflexiveTransitiveClosure(D);|
  <immutable preorder digraph with 4 vertices, 9 edges>
  !gapprompt@gap>| !gapinput@IsDistributiveLatticeDigraph(D);|
  true
  !gapprompt@gap>| !gapinput@N5 := Digraph([[2, 4], [3], [5], [5], []]);|
  <immutable digraph with 5 vertices, 5 edges>
  !gapprompt@gap>| !gapinput@N5 := DigraphReflexiveTransitiveClosure(N5);|
  <immutable preorder digraph with 5 vertices, 13 edges>
  !gapprompt@gap>| !gapinput@IsDistributiveLatticeDigraph(N5);|
  false
\end{Verbatim}
 }

 

\subsection{\textcolor{Chapter }{IsModularLatticeDigraph}}
\logpage{[ 6, 4, 9 ]}\nobreak
\hyperdef{L}{X80659D5784790B52}{}
{\noindent\textcolor{FuncColor}{$\triangleright$\enspace\texttt{IsModularLatticeDigraph({\mdseries\slshape digraph})\index{IsModularLatticeDigraph@\texttt{IsModularLatticeDigraph}}
\label{IsModularLatticeDigraph}
}\hfill{\scriptsize (property)}}\\
\textbf{\indent Returns:\ }
\texttt{true} or \texttt{false}.



 \texttt{IsModularLatticeDigraph} returns \texttt{true} if the digraph \mbox{\texttt{\mdseries\slshape digraph}} is a modular lattice digraph.

 A \emph{modular lattice digraph} is a lattice digraph (\texttt{IsLatticeDigraph} (\ref{IsLatticeDigraph})) which is modular. That is to say, the lattice digraph representing $N5$ is not embeddable as a lattice (see \href{https://en.wikipedia.org/wiki/Modular_lattice} {\texttt{https://en.wikipedia.org/wiki/Modular{\textunderscore}lattice}} and \texttt{IsLatticeEmbedding} (\ref{IsLatticeEmbedding:for digraphs and a permutation or transformation})).

  If the argument \mbox{\texttt{\mdseries\slshape digraph}} is mutable, then the return value of this property is recomputed every time it
is called.

 
\begin{Verbatim}[commandchars=!@|,fontsize=\small,frame=single,label=Example]
  !gapprompt@gap>| !gapinput@D := ChainDigraph(5);|
  <immutable chain digraph with 5 vertices>
  !gapprompt@gap>| !gapinput@D := DigraphReflexiveTransitiveClosure(D);|
  <immutable preorder digraph with 5 vertices, 15 edges>
  !gapprompt@gap>| !gapinput@IsModularLatticeDigraph(D);|
  true
  !gapprompt@gap>| !gapinput@N5 := Digraph([[2, 3], [4], [4], []]);|
  <immutable digraph with 4 vertices, 4 edges>
  !gapprompt@gap>| !gapinput@DigraphReflexiveTransitiveClosure(N5);|
  <immutable preorder digraph with 4 vertices, 9 edges>
  !gapprompt@gap>| !gapinput@IsModularLatticeDigraph(N5);|
  false
\end{Verbatim}
 }

 }

 
\section{\textcolor{Chapter }{Regularity}}\logpage{[ 6, 5, 0 ]}
\hyperdef{L}{X7AAF896982EC22FA}{}
{
 

\subsection{\textcolor{Chapter }{IsInRegularDigraph}}
\logpage{[ 6, 5, 1 ]}\nobreak
\hyperdef{L}{X7A3F78B37E31D9E1}{}
{\noindent\textcolor{FuncColor}{$\triangleright$\enspace\texttt{IsInRegularDigraph({\mdseries\slshape digraph})\index{IsInRegularDigraph@\texttt{IsInRegularDigraph}}
\label{IsInRegularDigraph}
}\hfill{\scriptsize (property)}}\\
\textbf{\indent Returns:\ }
\texttt{true} or \texttt{false}.



 This property is \texttt{true} if there is an integer \texttt{n} such that for every vertex \texttt{v} of digraph \mbox{\texttt{\mdseries\slshape digraph}} there are exactly \texttt{n} edges terminating in \texttt{v}. See also \texttt{IsOutRegularDigraph} (\ref{IsOutRegularDigraph}) and \texttt{IsRegularDigraph} (\ref{IsRegularDigraph}). 

 If the argument \mbox{\texttt{\mdseries\slshape digraph}} is mutable, then the return value of this property is recomputed every time it
is called.

 
\begin{Verbatim}[commandchars=!@|,fontsize=\small,frame=single,label=Example]
  !gapprompt@gap>| !gapinput@IsInRegularDigraph(CompleteDigraph(4));|
  true
  !gapprompt@gap>| !gapinput@IsInRegularDigraph(ChainDigraph(4));|
  false
\end{Verbatim}
 }

 

\subsection{\textcolor{Chapter }{IsOutRegularDigraph}}
\logpage{[ 6, 5, 2 ]}\nobreak
\hyperdef{L}{X7F0A806E7BCB413C}{}
{\noindent\textcolor{FuncColor}{$\triangleright$\enspace\texttt{IsOutRegularDigraph({\mdseries\slshape digraph})\index{IsOutRegularDigraph@\texttt{IsOutRegularDigraph}}
\label{IsOutRegularDigraph}
}\hfill{\scriptsize (property)}}\\
\textbf{\indent Returns:\ }
\texttt{true} or \texttt{false}.



 This property is \texttt{true} if there is an integer \texttt{n} such that for every vertex \texttt{v} of digraph \mbox{\texttt{\mdseries\slshape digraph}} there are exactly \texttt{n} edges starting at \texttt{v}. 

 See also \texttt{IsInRegularDigraph} (\ref{IsInRegularDigraph}) and \texttt{IsRegularDigraph} (\ref{IsRegularDigraph}). 

 If the argument \mbox{\texttt{\mdseries\slshape digraph}} is mutable, then the return value of this property is recomputed every time it
is called.

 
\begin{Verbatim}[commandchars=!@|,fontsize=\small,frame=single,label=Example]
  !gapprompt@gap>| !gapinput@IsOutRegularDigraph(CompleteDigraph(4));|
  true
  !gapprompt@gap>| !gapinput@IsOutRegularDigraph(ChainDigraph(4));|
  false
\end{Verbatim}
 }

 

\subsection{\textcolor{Chapter }{IsRegularDigraph}}
\logpage{[ 6, 5, 3 ]}\nobreak
\hyperdef{L}{X7AF9DB1E7DB12306}{}
{\noindent\textcolor{FuncColor}{$\triangleright$\enspace\texttt{IsRegularDigraph({\mdseries\slshape digraph})\index{IsRegularDigraph@\texttt{IsRegularDigraph}}
\label{IsRegularDigraph}
}\hfill{\scriptsize (property)}}\\
\textbf{\indent Returns:\ }
\texttt{true} or \texttt{false}.



 This property is \texttt{true} if there is an integer \texttt{n} such that for every vertex \texttt{v} of digraph \mbox{\texttt{\mdseries\slshape digraph}} there are exactly \texttt{n} edges starting and terminating at \texttt{v}. In other words, the property is \texttt{true} if \mbox{\texttt{\mdseries\slshape digraph}} is both in\texttt{\symbol{45}}regular and and out\texttt{\symbol{45}}regular.
See also \texttt{IsInRegularDigraph} (\ref{IsInRegularDigraph}) and \texttt{IsOutRegularDigraph} (\ref{IsOutRegularDigraph}). 

 If the argument \mbox{\texttt{\mdseries\slshape digraph}} is mutable, then the return value of this property is recomputed every time it
is called.

 
\begin{Verbatim}[commandchars=!@|,fontsize=\small,frame=single,label=Example]
  !gapprompt@gap>| !gapinput@IsRegularDigraph(CompleteDigraph(4));|
  true
  !gapprompt@gap>| !gapinput@IsRegularDigraph(ChainDigraph(4));|
  false
\end{Verbatim}
 }

 

\subsection{\textcolor{Chapter }{IsDistanceRegularDigraph}}
\logpage{[ 6, 5, 4 ]}\nobreak
\hyperdef{L}{X7E2082AC7CAE59CD}{}
{\noindent\textcolor{FuncColor}{$\triangleright$\enspace\texttt{IsDistanceRegularDigraph({\mdseries\slshape digraph})\index{IsDistanceRegularDigraph@\texttt{IsDistanceRegularDigraph}}
\label{IsDistanceRegularDigraph}
}\hfill{\scriptsize (property)}}\\
\textbf{\indent Returns:\ }
\texttt{true} or \texttt{false}.



 If \mbox{\texttt{\mdseries\slshape digraph}} is a connected symmetric graph, this property returns \texttt{true} if for any two vertices \texttt{u} and \texttt{v} of \mbox{\texttt{\mdseries\slshape digraph}} and any two integers \texttt{i} and \texttt{j} between \texttt{0} and the diameter of \mbox{\texttt{\mdseries\slshape digraph}}, the number of vertices at distance \texttt{i} from \texttt{u} and distance \texttt{j} from \texttt{v} depends only on \texttt{i}, \texttt{j}, and the distance between vertices \texttt{u} and \texttt{v}.

 Alternatively, a distance regular graph is a graph for which there exist
integers \texttt{b{\textunderscore}i}, \texttt{c{\textunderscore}i}, and \texttt{i} such that for any two vertices \texttt{u}, \texttt{v} in \mbox{\texttt{\mdseries\slshape digraph}} which are distance \texttt{i} apart, there are exactly \texttt{b{\textunderscore}i} neighbors of \texttt{v} which are at distance \texttt{i \texttt{\symbol{45}} 1} away from \texttt{u}, and \texttt{c{\textunderscore}i} neighbors of \texttt{v} which are at distance \texttt{i + 1} away from \texttt{u}. This definition is used to check whether \mbox{\texttt{\mdseries\slshape digraph}} is distance regular.

 In the case where \mbox{\texttt{\mdseries\slshape digraph}} is not symmetric or not connected, the property is \texttt{false}. 

 If the argument \mbox{\texttt{\mdseries\slshape digraph}} is mutable, then the return value of this property is recomputed every time it
is called.

 
\begin{Verbatim}[commandchars=!@|,fontsize=\small,frame=single,label=Example]
  !gapprompt@gap>| !gapinput@D := DigraphSymmetricClosure(ChainDigraph(5));;|
  !gapprompt@gap>| !gapinput@IsDistanceRegularDigraph(D);|
  false
  !gapprompt@gap>| !gapinput@D := Digraph([[2, 3, 4], [1, 3, 4], [1, 2, 4], [1, 2, 3]]);|
  <immutable digraph with 4 vertices, 12 edges>
  !gapprompt@gap>| !gapinput@IsDistanceRegularDigraph(D);|
  true
\end{Verbatim}
 }

 }

 
\section{\textcolor{Chapter }{Connectivity and cycles}}\logpage{[ 6, 6, 0 ]}
\hyperdef{L}{X780DAEB37C5E07FD}{}
{
 

\subsection{\textcolor{Chapter }{IsAcyclicDigraph}}
\logpage{[ 6, 6, 1 ]}\nobreak
\hyperdef{L}{X7BD08D4478218F77}{}
{\noindent\textcolor{FuncColor}{$\triangleright$\enspace\texttt{IsAcyclicDigraph({\mdseries\slshape digraph})\index{IsAcyclicDigraph@\texttt{IsAcyclicDigraph}}
\label{IsAcyclicDigraph}
}\hfill{\scriptsize (property)}}\\
\textbf{\indent Returns:\ }
\texttt{true} or \texttt{false}.



 This property is \texttt{true} if the digraph \mbox{\texttt{\mdseries\slshape digraph}} is acyclic, and \texttt{false} if it is not. A digraph is \emph{acyclic} if every directed cycle on the digraph is trivial. See Section \ref{Definitions} for the definition of a directed cycle, and of a trivial directed cycle.

 The method used in this operation has complexity $O(m+n)$ where $m$ is the number of edges (counting multiple edges as one) and $n$ is the number of vertices in the digraph. 

 If the argument \mbox{\texttt{\mdseries\slshape digraph}} is mutable, then the return value of this property is recomputed every time it
is called.

 
\begin{Verbatim}[commandchars=!@J,fontsize=\small,frame=single,label=Example]
  !gapprompt@gap>J !gapinput@Petersen := Graph(SymmetricGroup(5), [[1, 2]], OnSets,J
  !gapprompt@>J !gapinput@function(x, y)J
  !gapprompt@>J !gapinput@  return IsEmpty(Intersection(x, y));J
  !gapprompt@>J !gapinput@end);;J
  !gapprompt@gap>J !gapinput@D := Digraph(Petersen);J
  <immutable digraph with 10 vertices, 30 edges>
  !gapprompt@gap>J !gapinput@IsAcyclicDigraph(D);J
  false
  !gapprompt@gap>J !gapinput@D := DigraphFromDiSparse6String(J
  !gapprompt@>J !gapinput@".b_OGCIDBaPGkULEbQHCeRIdrHcuZMfRyDAbPhTi|zF");J
  <immutable digraph with 35 vertices, 34 edges>
  !gapprompt@gap>J !gapinput@IsAcyclicDigraph(D);J
  true
  !gapprompt@gap>J !gapinput@IsAcyclicDigraph(ChainDigraph(10));J
  true
  !gapprompt@gap>J !gapinput@D := CompleteDigraph(IsMutableDigraph, 4);J
  <mutable digraph with 4 vertices, 12 edges>
  !gapprompt@gap>J !gapinput@IsAcyclicDigraph(D);J
  false
  !gapprompt@gap>J !gapinput@IsAcyclicDigraph(CycleDigraph(10));J
  false
\end{Verbatim}
 }

 

\subsection{\textcolor{Chapter }{IsChainDigraph}}
\logpage{[ 6, 6, 2 ]}\nobreak
\hyperdef{L}{X79563B297C220FB5}{}
{\noindent\textcolor{FuncColor}{$\triangleright$\enspace\texttt{IsChainDigraph({\mdseries\slshape digraph})\index{IsChainDigraph@\texttt{IsChainDigraph}}
\label{IsChainDigraph}
}\hfill{\scriptsize (property)}}\\
\textbf{\indent Returns:\ }
\texttt{true} or \texttt{false}.



 \texttt{IsChainDigraph} returns \texttt{true} if the digraph \mbox{\texttt{\mdseries\slshape digraph}} is isomorphic to the chain digraph with the same number of vertices as \mbox{\texttt{\mdseries\slshape digraph}}, and \texttt{false} if it is not; see \texttt{ChainDigraph} (\ref{ChainDigraph}).

 A digraph is a \emph{chain} if and only if it is a directed tree, in which every vertex has out degree at
most one; see \texttt{IsDirectedTree} (\ref{IsDirectedTree}) and \texttt{OutDegrees} (\ref{OutDegrees}). 

 If the argument \mbox{\texttt{\mdseries\slshape digraph}} is mutable, then the return value of this property is recomputed every time it
is called.

 
\begin{Verbatim}[commandchars=!@|,fontsize=\small,frame=single,label=Example]
  !gapprompt@gap>| !gapinput@D := Digraph([[1, 3], [2, 3], [3]]);|
  <immutable digraph with 3 vertices, 5 edges>
  !gapprompt@gap>| !gapinput@IsChainDigraph(D);|
  false
  !gapprompt@gap>| !gapinput@D := ChainDigraph(5);|
  <immutable chain digraph with 5 vertices>
  !gapprompt@gap>| !gapinput@IsChainDigraph(D);|
  true
  !gapprompt@gap>| !gapinput@D := DigraphReverse(D);|
  <immutable digraph with 5 vertices, 4 edges>
  !gapprompt@gap>| !gapinput@IsChainDigraph(D);|
  true
  !gapprompt@gap>| !gapinput@D := ChainDigraph(IsMutableDigraph, 5);|
  <mutable digraph with 5 vertices, 4 edges>
  !gapprompt@gap>| !gapinput@IsChainDigraph(D);|
  true
  !gapprompt@gap>| !gapinput@DigraphReverse(D);|
  <mutable digraph with 5 vertices, 4 edges>
  !gapprompt@gap>| !gapinput@IsChainDigraph(D);|
  true
\end{Verbatim}
 }

 

\subsection{\textcolor{Chapter }{IsConnectedDigraph}}
\logpage{[ 6, 6, 3 ]}\nobreak
\hyperdef{L}{X83C08C0B7EC1A91F}{}
{\noindent\textcolor{FuncColor}{$\triangleright$\enspace\texttt{IsConnectedDigraph({\mdseries\slshape digraph})\index{IsConnectedDigraph@\texttt{IsConnectedDigraph}}
\label{IsConnectedDigraph}
}\hfill{\scriptsize (property)}}\\
\textbf{\indent Returns:\ }
\texttt{true} or \texttt{false}.



 This property is \texttt{true} if the digraph \mbox{\texttt{\mdseries\slshape digraph}} is weakly connected and \texttt{false} if it is not. A digraph \mbox{\texttt{\mdseries\slshape digraph}} is \emph{weakly connected} if it is possible to travel from any vertex to any other vertex by traversing
edges in either direction (possibly against the orientation of some of them). 

 The method used in this function has complexity $O(m)$ if the digraph's \texttt{DigraphSource} (\ref{DigraphSource}) attribute is set, otherwise it has complexity $O(m+n)$ (where $m$ is the number of edges and $n$ is the number of vertices of the digraph). 

 If the argument \mbox{\texttt{\mdseries\slshape digraph}} is mutable, then the return value of this property is recomputed every time it
is called.

 
\begin{Verbatim}[commandchars=!@|,fontsize=\small,frame=single,label=Example]
  !gapprompt@gap>| !gapinput@D := Digraph([[2], [3], []]);;|
  !gapprompt@gap>| !gapinput@IsConnectedDigraph(D);|
  true
  !gapprompt@gap>| !gapinput@D := Digraph([[1, 3], [4], [3], []]);;|
  !gapprompt@gap>| !gapinput@IsConnectedDigraph(D);|
  false
  !gapprompt@gap>| !gapinput@D := Digraph(IsMutableDigraph, [[2], [3], []]);;|
  !gapprompt@gap>| !gapinput@IsConnectedDigraph(D);|
  true
  !gapprompt@gap>| !gapinput@D := Digraph(IsMutableDigraph, [[1, 3], [4], [3], []]);;|
  !gapprompt@gap>| !gapinput@IsConnectedDigraph(D);|
  false
\end{Verbatim}
 }

 

\subsection{\textcolor{Chapter }{IsBiconnectedDigraph}}
\logpage{[ 6, 6, 4 ]}\nobreak
\hyperdef{L}{X838FAF2D825977BE}{}
{\noindent\textcolor{FuncColor}{$\triangleright$\enspace\texttt{IsBiconnectedDigraph({\mdseries\slshape digraph})\index{IsBiconnectedDigraph@\texttt{IsBiconnectedDigraph}}
\label{IsBiconnectedDigraph}
}\hfill{\scriptsize (property)}}\\
\textbf{\indent Returns:\ }
\texttt{true} or \texttt{false}.



 A connected digraph is \emph{biconnected} if it is still connected (in the sense of \texttt{IsConnectedDigraph} (\ref{IsConnectedDigraph})) when any vertex is removed. If \mbox{\texttt{\mdseries\slshape D}} has at least 3 vertices, then \texttt{IsBiconnectedDigraph} implies \texttt{IsBridgelessDigraph} (\ref{IsBridgelessDigraph}); see \texttt{ArticulationPoints} (\ref{ArticulationPoints}) or \texttt{Bridges} (\ref{Bridges}) for a more detailed explanation. 

 \texttt{IsBiconnectedDigraph} returns \texttt{true} if the digraph \mbox{\texttt{\mdseries\slshape digraph}} is biconnected, and \texttt{false} if it is not. In particular, \texttt{IsBiconnectedDigraph} returns \texttt{false} if \mbox{\texttt{\mdseries\slshape digraph}} is not connected. 

 Multiple edges are ignored by this method. 

 The method used in this operation has complexity $O(m+n)$ where $m$ is the number of edges and $n$ is the number of vertices in the digraph. 

 See also \texttt{Bridges} (\ref{Bridges}), \texttt{ArticulationPoints} (\ref{ArticulationPoints}), and \texttt{IsBridgelessDigraph} (\ref{IsBridgelessDigraph}). 

 If the argument \mbox{\texttt{\mdseries\slshape digraph}} is mutable, then the return value of this property is recomputed every time it
is called.

 
\begin{Verbatim}[commandchars=!@|,fontsize=\small,frame=single,label=Example]
  !gapprompt@gap>| !gapinput@IsBiconnectedDigraph(Digraph([[1, 3], [2, 3], [3]]));|
  false
  !gapprompt@gap>| !gapinput@IsBiconnectedDigraph(CycleDigraph(5));|
  true
  !gapprompt@gap>| !gapinput@D := Digraph([[1, 1], [1, 1, 2], [3], [3, 3, 4, 4]]);;|
  !gapprompt@gap>| !gapinput@IsBiconnectedDigraph(D);|
  false
  !gapprompt@gap>| !gapinput@D := CompleteBipartiteDigraph(IsMutableDigraph, 5, 4);|
  <mutable digraph with 9 vertices, 40 edges>
  !gapprompt@gap>| !gapinput@IsBiconnectedDigraph(D);|
  true
\end{Verbatim}
 }

 

\subsection{\textcolor{Chapter }{IsBridgelessDigraph}}
\logpage{[ 6, 6, 5 ]}\nobreak
\hyperdef{L}{X7D2A6572820B7F24}{}
{\noindent\textcolor{FuncColor}{$\triangleright$\enspace\texttt{IsBridgelessDigraph({\mdseries\slshape digraph})\index{IsBridgelessDigraph@\texttt{IsBridgelessDigraph}}
\label{IsBridgelessDigraph}
}\hfill{\scriptsize (property)}}\\
\textbf{\indent Returns:\ }
\texttt{true} or \texttt{false}.



 A connected digraph is \emph{bridgeless} if it is still connected (in the sense of \texttt{IsConnectedDigraph} (\ref{IsConnectedDigraph})) when any edge is removed. If \mbox{\texttt{\mdseries\slshape digraph}} has at least 3 vertices, then \texttt{IsBiconnectedDigraph} (\ref{IsBiconnectedDigraph}) implies \texttt{IsBridgelessDigraph}; see \texttt{ArticulationPoints} (\ref{ArticulationPoints}) or \texttt{Bridges} (\ref{Bridges}) for a more detailed explanation. 

 \texttt{IsBridgelessDigraph} returns \texttt{true} if the digraph \mbox{\texttt{\mdseries\slshape digraph}} is bridgeless, and \texttt{false} if it is not. In particular, \texttt{IsBridgelessDigraph} returns \texttt{false} if \mbox{\texttt{\mdseries\slshape digraph}} is not connected. 

 Multiple edges are ignored by this method. 

 The method used in this operation has complexity $O(m+n)$ where $m$ is the number of edges and $n$ is the number of vertices in the digraph. 

 See also \texttt{Bridges} (\ref{Bridges}), \texttt{ArticulationPoints} (\ref{ArticulationPoints}), and \texttt{IsBiconnectedDigraph} (\ref{IsBiconnectedDigraph}). 

 If the argument \mbox{\texttt{\mdseries\slshape digraph}} is mutable, then the return value of this property is recomputed every time it
is called.

 
\begin{Verbatim}[commandchars=!@|,fontsize=\small,frame=single,label=Example]
  !gapprompt@gap>| !gapinput@IsBridgelessDigraph(Digraph([[1, 3], [2, 3], [3]]));|
  false
  !gapprompt@gap>| !gapinput@IsBridgelessDigraph(CycleDigraph(5));|
  true
  !gapprompt@gap>| !gapinput@D := Digraph([[1, 1], [1, 1, 2], [3], [3, 3, 4, 4]]);;|
  !gapprompt@gap>| !gapinput@IsBridgelessDigraph(D);|
  false
  !gapprompt@gap>| !gapinput@D := CompleteBipartiteDigraph(IsMutableDigraph, 5, 4);|
  <mutable digraph with 9 vertices, 40 edges>
  !gapprompt@gap>| !gapinput@IsBridgelessDigraph(D);|
  true
  !gapprompt@gap>| !gapinput@D := Digraph([[2, 5], [1, 3, 4, 5], [2, 4], [2, 3], [1, 2]]);|
  <immutable digraph with 5 vertices, 12 edges>
  !gapprompt@gap>| !gapinput@IsBridgelessDigraph(D);|
  true
  !gapprompt@gap>| !gapinput@IsBiconnectedDigraph(D);|
  false
  !gapprompt@gap>| !gapinput@D := Digraph([[2], [3], [4], [2]]);|
  <immutable digraph with 4 vertices, 4 edges>
  !gapprompt@gap>| !gapinput@IsBridgelessDigraph(D);|
  false
  !gapprompt@gap>| !gapinput@IsBiconnectedDigraph(D);|
  false
  !gapprompt@gap>| !gapinput@IsBridgelessDigraph(ChainDigraph(2));|
  false
  !gapprompt@gap>| !gapinput@IsBiconnectedDigraph(ChainDigraph(2));|
  true
\end{Verbatim}
 }

 

\subsection{\textcolor{Chapter }{IsStronglyConnectedDigraph}}
\logpage{[ 6, 6, 6 ]}\nobreak
\hyperdef{L}{X7B37B9467C68C208}{}
{\noindent\textcolor{FuncColor}{$\triangleright$\enspace\texttt{IsStronglyConnectedDigraph({\mdseries\slshape digraph})\index{IsStronglyConnectedDigraph@\texttt{IsStronglyConnectedDigraph}}
\label{IsStronglyConnectedDigraph}
}\hfill{\scriptsize (property)}}\\
\textbf{\indent Returns:\ }
\texttt{true} or \texttt{false}.



 This property is \texttt{true} if the digraph \mbox{\texttt{\mdseries\slshape digraph}} is strongly connected and \texttt{false} if it is not. 

 A digraph \mbox{\texttt{\mdseries\slshape digraph}} is \emph{strongly connected} if there is a directed path from every vertex to every other vertex. See
Section \ref{Definitions} for the definition of a directed path. 

 The method used in this operation is based on Gabow's Algorithm \cite{Gab00} and has complexity $O(m+n)$, where $m$ is the number of edges (counting multiple edges as one) and $n$ is the number of vertices in the digraph. 

 If the argument \mbox{\texttt{\mdseries\slshape digraph}} is mutable, then the return value of this property is recomputed every time it
is called.

 
\begin{Verbatim}[commandchars=!@|,fontsize=\small,frame=single,label=Example]
  !gapprompt@gap>| !gapinput@D := CycleDigraph(250000);|
  <immutable cycle digraph with 250000 vertices>
  !gapprompt@gap>| !gapinput@IsStronglyConnectedDigraph(D);|
  true
  !gapprompt@gap>| !gapinput@D := DigraphRemoveEdges(D, [[250000, 1]]);|
  <immutable digraph with 250000 vertices, 249999 edges>
  !gapprompt@gap>| !gapinput@IsStronglyConnectedDigraph(D);|
  false
  !gapprompt@gap>| !gapinput@D := CycleDigraph(IsMutableDigraph, 250000);|
  <mutable digraph with 250000 vertices, 250000 edges>
  !gapprompt@gap>| !gapinput@IsStronglyConnectedDigraph(D);|
  true
  !gapprompt@gap>| !gapinput@DigraphRemoveEdge(D, [250000, 1]);|
  <mutable digraph with 250000 vertices, 249999 edges>
  !gapprompt@gap>| !gapinput@IsStronglyConnectedDigraph(D);|
  false
\end{Verbatim}
 }

 

\subsection{\textcolor{Chapter }{IsAperiodicDigraph}}
\logpage{[ 6, 6, 7 ]}\nobreak
\hyperdef{L}{X80E883967EBE839E}{}
{\noindent\textcolor{FuncColor}{$\triangleright$\enspace\texttt{IsAperiodicDigraph({\mdseries\slshape digraph})\index{IsAperiodicDigraph@\texttt{IsAperiodicDigraph}}
\label{IsAperiodicDigraph}
}\hfill{\scriptsize (property)}}\\
\textbf{\indent Returns:\ }
\texttt{true} or \texttt{false}.



 This property is \texttt{true} if the digraph \mbox{\texttt{\mdseries\slshape digraph}} is aperiodic, i.e. if its \texttt{DigraphPeriod} (\ref{DigraphPeriod}) is equal to 1. Otherwise, the property is \texttt{false}. 

 If the argument \mbox{\texttt{\mdseries\slshape digraph}} is mutable, then the return value of this property is recomputed every time it
is called.

 
\begin{Verbatim}[commandchars=!@|,fontsize=\small,frame=single,label=Example]
  !gapprompt@gap>| !gapinput@D := Digraph([[6], [1], [2], [3], [4, 4], [5]]);|
  <immutable multidigraph with 6 vertices, 7 edges>
  !gapprompt@gap>| !gapinput@IsAperiodicDigraph(D);|
  false
  !gapprompt@gap>| !gapinput@D := Digraph([[2], [3, 5], [4], [5], [1, 2]]);|
  <immutable digraph with 5 vertices, 7 edges>
  !gapprompt@gap>| !gapinput@IsAperiodicDigraph(D);|
  true
  !gapprompt@gap>| !gapinput@D := Digraph(IsMutableDigraph, [[2], [3, 5], [4], [5], [1, 2]]);|
  <mutable digraph with 5 vertices, 7 edges>
  !gapprompt@gap>| !gapinput@IsAperiodicDigraph(D);|
  true
\end{Verbatim}
 }

 

\subsection{\textcolor{Chapter }{IsDirectedTree}}
\logpage{[ 6, 6, 8 ]}\nobreak
\hyperdef{L}{X7B46EA6C7B2DF2FB}{}
{\noindent\textcolor{FuncColor}{$\triangleright$\enspace\texttt{IsDirectedTree({\mdseries\slshape digraph})\index{IsDirectedTree@\texttt{IsDirectedTree}}
\label{IsDirectedTree}
}\hfill{\scriptsize (property)}}\\
\noindent\textcolor{FuncColor}{$\triangleright$\enspace\texttt{IsDirectedForest({\mdseries\slshape digraph})\index{IsDirectedForest@\texttt{IsDirectedForest}}
\label{IsDirectedForest}
}\hfill{\scriptsize (property)}}\\
\textbf{\indent Returns:\ }
\texttt{true} or \texttt{false}.



 The property \texttt{IsDirectedTree} returns \texttt{true} if the digraph \mbox{\texttt{\mdseries\slshape digraph}} is a directed tree, and the property \texttt{IsDirectedForest} returns \texttt{true} if \mbox{\texttt{\mdseries\slshape digraph}} is a directed forest; otherwise these properties return \texttt{false}. 

 A \emph{directed tree} is an acyclic digraph with precisely one source, without multiple edges, and
such that no two vertices share an out\texttt{\symbol{45}}neighbour. See and \texttt{IsAcyclicDigraph} (\ref{IsAcyclicDigraph}) and \texttt{DigraphSources} (\ref{DigraphSources}) for more information about these terms. 

 A \emph{directed forest} is a digraph with at least one vertex, each of whose connected components is a
directed tree. In other words, a directed forest is isomorphic to a disjoint
union of directed trees. In particular, every directed tree is a directed
forest. See \texttt{DigraphConnectedComponents} (\ref{DigraphConnectedComponents}) and \texttt{DigraphDisjointUnion} (\ref{DigraphDisjointUnion:for a list of digraphs}). 

 Please note that the empty digraph with zero vertices is not considered to be
a directed tree or directed forest, because it has no source. 

 If the argument \mbox{\texttt{\mdseries\slshape digraph}} is mutable, then the return value of this property is recomputed every time it
is called.

 
\begin{Verbatim}[commandchars=!@|,fontsize=\small,frame=single,label=Example]
  !gapprompt@gap>| !gapinput@D := Digraph([[3], [3], []]);|
  <immutable digraph with 3 vertices, 2 edges>
  !gapprompt@gap>| !gapinput@IsDirectedTree(D);|
  false
  !gapprompt@gap>| !gapinput@D := Digraph([[2], [3], []]);|
  <immutable digraph with 3 vertices, 2 edges>
  !gapprompt@gap>| !gapinput@IsDirectedTree(D);|
  true
  !gapprompt@gap>| !gapinput@IsDirectedForest(DigraphDisjointUnion(D, D));|
  true
  !gapprompt@gap>| !gapinput@D := Digraph([[2, 3], [6], [4, 5], [], [], []]);|
  <immutable digraph with 6 vertices, 5 edges>
  !gapprompt@gap>| !gapinput@IsDirectedTree(D);|
  true
\end{Verbatim}
 }

 

\subsection{\textcolor{Chapter }{IsUndirectedTree}}
\logpage{[ 6, 6, 9 ]}\nobreak
\hyperdef{L}{X80FC20FA7AC4BC2A}{}
{\noindent\textcolor{FuncColor}{$\triangleright$\enspace\texttt{IsUndirectedTree({\mdseries\slshape digraph})\index{IsUndirectedTree@\texttt{IsUndirectedTree}}
\label{IsUndirectedTree}
}\hfill{\scriptsize (property)}}\\
\noindent\textcolor{FuncColor}{$\triangleright$\enspace\texttt{IsUndirectedForest({\mdseries\slshape digraph})\index{IsUndirectedForest@\texttt{IsUndirectedForest}}
\label{IsUndirectedForest}
}\hfill{\scriptsize (property)}}\\
\textbf{\indent Returns:\ }
\texttt{true} or \texttt{false}.



 The property \texttt{IsUndirectedTree} returns \texttt{true} if the digraph \mbox{\texttt{\mdseries\slshape digraph}} is an undirected tree, and the property \texttt{IsUndirectedForest} returns \texttt{true} if \mbox{\texttt{\mdseries\slshape digraph}} is an undirected forest; otherwise, these properties return \texttt{false}. 

 An \emph{undirected tree} is a symmetric digraph without loops, in which for any pair of distinct
vertices \texttt{u} and \texttt{v}, there is exactly one directed path from \texttt{u} to \texttt{v}. See \texttt{IsSymmetricDigraph} (\ref{IsSymmetricDigraph}) and \texttt{DigraphHasLoops} (\ref{DigraphHasLoops}), and see Section \ref{Definitions} for the definition of directed path. This definition implies that an
undirected tree has no multiple edges. 

 An \emph{undirected forest} is a digraph, each of whose connected components is an undirected tree. In
other words, an undirected forest is isomorphic to a disjoint union of
undirected trees. See \texttt{DigraphConnectedComponents} (\ref{DigraphConnectedComponents}) and \texttt{DigraphDisjointUnion} (\ref{DigraphDisjointUnion:for a list of digraphs}). In particular, every undirected tree is an undirected forest. 

 Please note that the digraph with zero vertices is considered to be neither an
undirected tree nor an undirected forest. 

 If the argument \mbox{\texttt{\mdseries\slshape digraph}} is mutable, then the return value of this property is recomputed every time it
is called.

 
\begin{Verbatim}[commandchars=!@|,fontsize=\small,frame=single,label=Example]
  !gapprompt@gap>| !gapinput@D := Digraph([[3], [3], [1, 2]]);|
  <immutable digraph with 3 vertices, 4 edges>
  !gapprompt@gap>| !gapinput@IsUndirectedTree(D);|
  true
  !gapprompt@gap>| !gapinput@IsSymmetricDigraph(D) and not DigraphHasLoops(D);|
  true
  !gapprompt@gap>| !gapinput@D := Digraph([[3], [5], [1, 4], [3], [2]]);|
  <immutable digraph with 5 vertices, 6 edges>
  !gapprompt@gap>| !gapinput@IsConnectedDigraph(D);|
  false
  !gapprompt@gap>| !gapinput@IsUndirectedTree(D);|
  false
  !gapprompt@gap>| !gapinput@IsUndirectedForest(D);|
  true
  !gapprompt@gap>| !gapinput@D := Digraph([[1, 2], [1], [2]]);|
  <immutable digraph with 3 vertices, 4 edges>
  !gapprompt@gap>| !gapinput@IsUndirectedTree(D) or IsUndirectedForest(D);|
  false
  !gapprompt@gap>| !gapinput@IsSymmetricDigraph(D) or not DigraphHasLoops(D);|
  false
\end{Verbatim}
 }

 

\subsection{\textcolor{Chapter }{IsEulerianDigraph}}
\logpage{[ 6, 6, 10 ]}\nobreak
\hyperdef{L}{X79DFBB198525544E}{}
{\noindent\textcolor{FuncColor}{$\triangleright$\enspace\texttt{IsEulerianDigraph({\mdseries\slshape digraph})\index{IsEulerianDigraph@\texttt{IsEulerianDigraph}}
\label{IsEulerianDigraph}
}\hfill{\scriptsize (property)}}\\
\textbf{\indent Returns:\ }
\texttt{true} or \texttt{false}.



 This property returns true if the digraph \mbox{\texttt{\mdseries\slshape digraph}} is Eulerian. 

 A connected digraph is called \emph{Eulerian} if there exists a directed circuit on the digraph which includes every edge
exactly once. See Section \ref{Definitions} for the definition of a directed circuit. Note that the empty digraph with at
most one vertex is considered to be Eulerian. 

 If the argument \mbox{\texttt{\mdseries\slshape digraph}} is mutable, then the return value of this property is recomputed every time it
is called.

 
\begin{Verbatim}[commandchars=!@|,fontsize=\small,frame=single,label=Example]
  !gapprompt@gap>| !gapinput@D := Digraph([[]]);|
  <immutable empty digraph with 1 vertex>
  !gapprompt@gap>| !gapinput@IsEulerianDigraph(D);|
  true
  !gapprompt@gap>| !gapinput@D := Digraph([[2], []]);|
  <immutable digraph with 2 vertices, 1 edge>
  !gapprompt@gap>| !gapinput@IsEulerianDigraph(D);|
  false
  !gapprompt@gap>| !gapinput@D := Digraph([[3], [], [2]]);|
  <immutable digraph with 3 vertices, 2 edges>
  !gapprompt@gap>| !gapinput@IsEulerianDigraph(D);|
  false
  !gapprompt@gap>| !gapinput@D := Digraph([[2], [3], [1]]);|
  <immutable digraph with 3 vertices, 3 edges>
  !gapprompt@gap>| !gapinput@IsEulerianDigraph(D);|
  true
\end{Verbatim}
 }

 

\subsection{\textcolor{Chapter }{IsHamiltonianDigraph}}
\logpage{[ 6, 6, 11 ]}\nobreak
\hyperdef{L}{X79EFEBC37C2D262D}{}
{\noindent\textcolor{FuncColor}{$\triangleright$\enspace\texttt{IsHamiltonianDigraph({\mdseries\slshape digraph})\index{IsHamiltonianDigraph@\texttt{IsHamiltonianDigraph}}
\label{IsHamiltonianDigraph}
}\hfill{\scriptsize (property)}}\\
\textbf{\indent Returns:\ }
\texttt{true} or \texttt{false}.



 If \mbox{\texttt{\mdseries\slshape digraph}} is Hamiltonian, then this property returns \texttt{true}, and \texttt{false} if it is not. 

 A digraph with \texttt{n} vertices is \emph{Hamiltonian} if it has a directed cycle of length \texttt{n}. See Section \ref{Definitions} for the definition of a directed cycle. Note the empty digraphs on 0 and 1
vertices are considered to be Hamiltonian.

 The method used in this operation has the worst case complexity as \texttt{DigraphMonomorphism} (\ref{DigraphMonomorphism}). 

 If the argument \mbox{\texttt{\mdseries\slshape digraph}} is mutable, then the return value of this property is recomputed every time it
is called.

 
\begin{Verbatim}[commandchars=!@|,fontsize=\small,frame=single,label=Example]
  !gapprompt@gap>| !gapinput@g := Digraph([[]]);|
  <immutable empty digraph with 1 vertex>
  !gapprompt@gap>| !gapinput@IsHamiltonianDigraph(g);|
  true
  !gapprompt@gap>| !gapinput@g := Digraph([[2], [1]]);|
  <immutable digraph with 2 vertices, 2 edges>
  !gapprompt@gap>| !gapinput@IsHamiltonianDigraph(g);|
  true
  !gapprompt@gap>| !gapinput@g := Digraph([[3], [], [2]]);|
  <immutable digraph with 3 vertices, 2 edges>
  !gapprompt@gap>| !gapinput@IsHamiltonianDigraph(g);|
  false
  !gapprompt@gap>| !gapinput@g := Digraph([[2], [3], [1]]);|
  <immutable digraph with 3 vertices, 3 edges>
  !gapprompt@gap>| !gapinput@IsHamiltonianDigraph(g);|
  true
\end{Verbatim}
 }

 

\subsection{\textcolor{Chapter }{IsCycleDigraph}}
\logpage{[ 6, 6, 12 ]}\nobreak
\hyperdef{L}{X7E91320B8028F5D8}{}
{\noindent\textcolor{FuncColor}{$\triangleright$\enspace\texttt{IsCycleDigraph({\mdseries\slshape digraph})\index{IsCycleDigraph@\texttt{IsCycleDigraph}}
\label{IsCycleDigraph}
}\hfill{\scriptsize (property)}}\\
\textbf{\indent Returns:\ }
\texttt{true} or \texttt{false}.



 \texttt{IsCycleDigraph} returns \texttt{true} if the digraph \mbox{\texttt{\mdseries\slshape digraph}} is isomorphic to the cycle digraph with the same number of vertices as \mbox{\texttt{\mdseries\slshape digraph}}, and \texttt{false} if it is not; see \texttt{CycleDigraph} (\ref{CycleDigraph}).

 A digraph is a \emph{cycle} if and only if it is strongly connected and has the same number of edges as
vertices. 

 If the argument \mbox{\texttt{\mdseries\slshape digraph}} is mutable, then the return value of this property is recomputed every time it
is called.

 
\begin{Verbatim}[commandchars=!@|,fontsize=\small,frame=single,label=Example]
  !gapprompt@gap>| !gapinput@D := Digraph([[1, 3], [2, 3], [3]]);|
  <immutable digraph with 3 vertices, 5 edges>
  !gapprompt@gap>| !gapinput@IsCycleDigraph(D);|
  false
  !gapprompt@gap>| !gapinput@D := CycleDigraph(5);|
  <immutable cycle digraph with 5 vertices>
  !gapprompt@gap>| !gapinput@IsCycleDigraph(D);|
  true
  !gapprompt@gap>| !gapinput@D := OnDigraphs(D, (1, 2, 3));|
  <immutable digraph with 5 vertices, 5 edges>
  !gapprompt@gap>| !gapinput@D = CycleDigraph(5);|
  false
  !gapprompt@gap>| !gapinput@IsCycleDigraph(D);|
  true
\end{Verbatim}
 }

 }

 
\section{\textcolor{Chapter }{Planarity}}\label{Planarity}
\logpage{[ 6, 7, 0 ]}
\hyperdef{L}{X7E2305528492DDC0}{}
{
 

\subsection{\textcolor{Chapter }{IsPlanarDigraph}}
\logpage{[ 6, 7, 1 ]}\nobreak
\hyperdef{L}{X8606D415858C40AA}{}
{\noindent\textcolor{FuncColor}{$\triangleright$\enspace\texttt{IsPlanarDigraph({\mdseries\slshape digraph})\index{IsPlanarDigraph@\texttt{IsPlanarDigraph}}
\label{IsPlanarDigraph}
}\hfill{\scriptsize (property)}}\\
\textbf{\indent Returns:\ }
\texttt{true} or \texttt{false}.



 A \emph{planar} digraph is a digraph that can be embedded in the plane in such a way that its
edges do not intersect. A digraph is planar if and only if it does not have a
subdigraph that is homeomorphic to either the complete graph on \texttt{5} vertices or the complete bipartite graph with vertex sets of sizes \texttt{3} and \texttt{3}. 

 \texttt{IsPlanarDigraph} returns \texttt{true} if the digraph \mbox{\texttt{\mdseries\slshape digraph}} is planar and \texttt{false} if it is not. The directions and multiplicities of any edges in \mbox{\texttt{\mdseries\slshape digraph}} are ignored by \texttt{IsPlanarDigraph}. 

 See also \texttt{IsOuterPlanarDigraph} (\ref{IsOuterPlanarDigraph}). 

 This method uses the reference implementation in \href{https://github.com/graph-algorithms/edge-addition-planarity-suite} {edge-addition-planarity-suite}  by John Boyer of the algorithms described in \cite{BM06}. 
\begin{Verbatim}[commandchars=!@|,fontsize=\small,frame=single,label=Example]
  !gapprompt@gap>| !gapinput@IsPlanarDigraph(CompleteDigraph(4));|
  true
  !gapprompt@gap>| !gapinput@IsPlanarDigraph(CompleteDigraph(5));|
  false
  !gapprompt@gap>| !gapinput@IsPlanarDigraph(CompleteBipartiteDigraph(2, 3));|
  true
  !gapprompt@gap>| !gapinput@IsPlanarDigraph(CompleteBipartiteDigraph(3, 3));|
  false
  !gapprompt@gap>| !gapinput@IsPlanarDigraph(CompleteDigraph(IsMutableDigraph, 4));|
  true
  !gapprompt@gap>| !gapinput@IsPlanarDigraph(CompleteDigraph(IsMutableDigraph, 5));|
  false
  !gapprompt@gap>| !gapinput@IsPlanarDigraph(CompleteBipartiteDigraph(IsMutableDigraph, 2, 3));|
  true
  !gapprompt@gap>| !gapinput@IsPlanarDigraph(CompleteBipartiteDigraph(IsMutableDigraph, 3, 3));|
  false
\end{Verbatim}
 }

 

\subsection{\textcolor{Chapter }{IsOuterPlanarDigraph}}
\logpage{[ 6, 7, 2 ]}\nobreak
\hyperdef{L}{X8251E8B187E7F059}{}
{\noindent\textcolor{FuncColor}{$\triangleright$\enspace\texttt{IsOuterPlanarDigraph({\mdseries\slshape digraph})\index{IsOuterPlanarDigraph@\texttt{IsOuterPlanarDigraph}}
\label{IsOuterPlanarDigraph}
}\hfill{\scriptsize (property)}}\\
\textbf{\indent Returns:\ }
\texttt{true} or \texttt{false}.



 An \emph{outer planar} digraph is a digraph that can be embedded in the plane in such a way that its
edges do not intersect, and all vertices belong to the unbounded face of the
embedding. A digraph is outer planar if and only if it does not have a
subdigraph that is homeomorphic to either the complete graph on \texttt{4} vertices or the complete bipartite graph with vertex sets of sizes \texttt{2} and \texttt{3}. 

 \texttt{IsOuterPlanarDigraph} returns \texttt{true} if the digraph \mbox{\texttt{\mdseries\slshape digraph}} is outer planar and \texttt{false} if it is not. The directions and multiplicities of any edges in \mbox{\texttt{\mdseries\slshape digraph}} are ignored by \texttt{IsPlanarDigraph}. 

 See also \texttt{IsPlanarDigraph} (\ref{IsPlanarDigraph}). This method uses the reference implementation in \href{https://github.com/graph-algorithms/edge-addition-planarity-suite} {edge-addition-planarity-suite}  by John Boyer of the algorithms described in \cite{BM06}. 
\begin{Verbatim}[commandchars=!@|,fontsize=\small,frame=single,label=Example]
  !gapprompt@gap>| !gapinput@IsOuterPlanarDigraph(CompleteDigraph(4));|
  false
  !gapprompt@gap>| !gapinput@IsOuterPlanarDigraph(CompleteDigraph(5));|
  false
  !gapprompt@gap>| !gapinput@IsOuterPlanarDigraph(CompleteBipartiteDigraph(2, 3));|
  false
  !gapprompt@gap>| !gapinput@IsOuterPlanarDigraph(CompleteBipartiteDigraph(3, 3));|
  false
  !gapprompt@gap>| !gapinput@IsOuterPlanarDigraph(CycleDigraph(10));|
  true
  !gapprompt@gap>| !gapinput@IsOuterPlanarDigraph(CompleteDigraph(IsMutableDigraph, 4));|
  false
  !gapprompt@gap>| !gapinput@IsOuterPlanarDigraph(CompleteDigraph(IsMutableDigraph, 5));|
  false
  !gapprompt@gap>| !gapinput@IsOuterPlanarDigraph(CompleteBipartiteDigraph(IsMutableDigraph,|
  !gapprompt@>| !gapinput@                                                 2, 3));|
  false
  !gapprompt@gap>| !gapinput@IsOuterPlanarDigraph(CompleteBipartiteDigraph(IsMutableDigraph,|
  !gapprompt@>| !gapinput@                                                 3, 3));|
  false
  !gapprompt@gap>| !gapinput@IsOuterPlanarDigraph(CycleDigraph(IsMutableDigraph, 10));|
  true
  
\end{Verbatim}
 }

 }

 
\section{\textcolor{Chapter }{Homomorphisms and transformations}}\logpage{[ 6, 8, 0 ]}
\hyperdef{L}{X83DC8AD283C41326}{}
{
 

\subsection{\textcolor{Chapter }{IsDigraphCore}}
\logpage{[ 6, 8, 1 ]}\nobreak
\hyperdef{L}{X78C0B2637985648B}{}
{\noindent\textcolor{FuncColor}{$\triangleright$\enspace\texttt{IsDigraphCore({\mdseries\slshape digraph})\index{IsDigraphCore@\texttt{IsDigraphCore}}
\label{IsDigraphCore}
}\hfill{\scriptsize (property)}}\\
\textbf{\indent Returns:\ }
\texttt{true} or \texttt{false}.



 This property returns \texttt{true} if \mbox{\texttt{\mdseries\slshape digraph}} is a core, and \texttt{false} if it is not.

 A digraph \texttt{D} is a \emph{core} if and only if it has no proper subdigraphs \texttt{A} such that there exists a homomorphism from \texttt{D} to \texttt{A}. In other words, a digraph \texttt{D} is a core if and only if every endomorphism on \texttt{D} is an automorphism on \texttt{D}. 

 If the argument \mbox{\texttt{\mdseries\slshape digraph}} is mutable, then the return value of this property is recomputed every time it
is called.

 
\begin{Verbatim}[commandchars=!@|,fontsize=\small,frame=single,label=Example]
  !gapprompt@gap>| !gapinput@D := CompleteDigraph(6);|
  <immutable complete digraph with 6 vertices>
  !gapprompt@gap>| !gapinput@IsDigraphCore(D);|
  true
  !gapprompt@gap>| !gapinput@D := DigraphSymmetricClosure(CycleDigraph(6));|
  <immutable symmetric digraph with 6 vertices, 12 edges>
  !gapprompt@gap>| !gapinput@DigraphHomomorphism(D, CompleteDigraph(2));|
  Transformation( [ 1, 2, 1, 2, 1, 2 ] )
  !gapprompt@gap>| !gapinput@IsDigraphCore(D);|
  false
\end{Verbatim}
 }

 

\subsection{\textcolor{Chapter }{IsEdgeTransitive}}
\logpage{[ 6, 8, 2 ]}\nobreak
\hyperdef{L}{X8361CA6E8401FF26}{}
{\noindent\textcolor{FuncColor}{$\triangleright$\enspace\texttt{IsEdgeTransitive({\mdseries\slshape digraph})\index{IsEdgeTransitive@\texttt{IsEdgeTransitive}}
\label{IsEdgeTransitive}
}\hfill{\scriptsize (property)}}\\
\textbf{\indent Returns:\ }
\texttt{true} or \texttt{false}.



 If \mbox{\texttt{\mdseries\slshape digraph}} is a digraph without multiple edges, then \texttt{IsEdgeTransitive} returns \texttt{true} if \mbox{\texttt{\mdseries\slshape digraph}} is edge transitive, and \texttt{false} otherwise. A digraph is \emph{edge transitive} if its automorphism group acts transitively on its edges (via the action \texttt{OnPairs} (\textbf{Reference: OnPairs})). 

 If the argument \mbox{\texttt{\mdseries\slshape digraph}} is mutable, then the return value of this property is recomputed every time it
is called.

 
\begin{Verbatim}[commandchars=!@|,fontsize=\small,frame=single,label=Example]
  !gapprompt@gap>| !gapinput@IsEdgeTransitive(CompleteDigraph(2));|
  true
  !gapprompt@gap>| !gapinput@IsEdgeTransitive(ChainDigraph(3));|
  false
  !gapprompt@gap>| !gapinput@IsEdgeTransitive(Digraph([[2], [3, 3, 3], []]));|
  Error, the argument <D> must be a digraph with no multiple edges,
\end{Verbatim}
 }

 

\subsection{\textcolor{Chapter }{IsVertexTransitive}}
\logpage{[ 6, 8, 3 ]}\nobreak
\hyperdef{L}{X7B26E4447A1611E9}{}
{\noindent\textcolor{FuncColor}{$\triangleright$\enspace\texttt{IsVertexTransitive({\mdseries\slshape digraph})\index{IsVertexTransitive@\texttt{IsVertexTransitive}}
\label{IsVertexTransitive}
}\hfill{\scriptsize (property)}}\\
\textbf{\indent Returns:\ }
\texttt{true} or \texttt{false}.



 If \mbox{\texttt{\mdseries\slshape digraph}} is a digraph, then \texttt{IsVertexTransitive} returns \texttt{true} if \mbox{\texttt{\mdseries\slshape digraph}} is vertex transitive, and \texttt{false} otherwise. A digraph is \emph{vertex transitive} if its automorphism group acts transitively on its vertices. 

 If the argument \mbox{\texttt{\mdseries\slshape digraph}} is mutable, then the return value of this property is recomputed every time it
is called.

 
\begin{Verbatim}[commandchars=!@|,fontsize=\small,frame=single,label=Example]
  !gapprompt@gap>| !gapinput@IsVertexTransitive(CompleteDigraph(2));|
  true
  !gapprompt@gap>| !gapinput@IsVertexTransitive(ChainDigraph(3));|
  false
\end{Verbatim}
 }

 

\subsection{\textcolor{Chapter }{Is2EdgeTransitive}}
\logpage{[ 6, 8, 4 ]}\nobreak
\hyperdef{L}{X7FB8789A7BF847E6}{}
{\noindent\textcolor{FuncColor}{$\triangleright$\enspace\texttt{Is2EdgeTransitive({\mdseries\slshape digraph})\index{Is2EdgeTransitive@\texttt{Is2EdgeTransitive}}
\label{Is2EdgeTransitive}
}\hfill{\scriptsize (property)}}\\
\textbf{\indent Returns:\ }
\texttt{true} or \texttt{false}.



 If \mbox{\texttt{\mdseries\slshape digraph}} is a digraph without multiple edges, then \texttt{Is2EdgeTransitive} returns \texttt{true} if \mbox{\texttt{\mdseries\slshape digraph}} is 2\texttt{\symbol{45}}edge transitive, and \texttt{false} otherwise. If \mbox{\texttt{\mdseries\slshape digraph}} has multiple edges, then \texttt{Is2EdgeTransitive} returns an error. A digraph is \emph{2\texttt{\symbol{45}}edge transitive} if its automorphism group acts transitively on 2\texttt{\symbol{45}}edges via
the action \texttt{OnTuples} (\textbf{Reference: OnTuples}). A \emph{2\texttt{\symbol{45}}edge} in a digraph is a triple (u, v, w) of distinct vertices such that (u, v) and
(v, w) are edges. 

 If the argument \mbox{\texttt{\mdseries\slshape digraph}} is mutable, then the return value of this property is recomputed every time it
is called.

 
\begin{Verbatim}[commandchars=!@|,fontsize=\small,frame=single,label=Example]
  !gapprompt@gap>| !gapinput@Is2EdgeTransitive(CompleteDigraph(4));|
  true
  !gapprompt@gap>| !gapinput@Is2EdgeTransitive(DigraphByEdges([[1, 2], [2, 3], [3, 4]]));|
  false
  !gapprompt@gap>| !gapinput@Is2EdgeTransitive(CycleDigraph(5));|
  true
  !gapprompt@gap>| !gapinput@Is2EdgeTransitive(Digraph([[2], [3, 3, 3], []]));|
  Error, the argument <D> must be a digraph with no multiple edges,
\end{Verbatim}
 }

 }

 }

 
\chapter{\textcolor{Chapter }{Homomorphisms}}\label{Homomorphisms}
\logpage{[ 7, 0, 0 ]}
\hyperdef{L}{X84975388859F203D}{}
{
  
\section{\textcolor{Chapter }{Acting on digraphs}}\logpage{[ 7, 1, 0 ]}
\hyperdef{L}{X8200192B80AD2071}{}
{
 

\subsection{\textcolor{Chapter }{OnDigraphs (for a digraph and a perm)}}
\logpage{[ 7, 1, 1 ]}\nobreak
\hyperdef{L}{X808972017C486F1F}{}
{\noindent\textcolor{FuncColor}{$\triangleright$\enspace\texttt{OnDigraphs({\mdseries\slshape digraph, perm})\index{OnDigraphs@\texttt{OnDigraphs}!for a digraph and a perm}
\label{OnDigraphs:for a digraph and a perm}
}\hfill{\scriptsize (operation)}}\\
\noindent\textcolor{FuncColor}{$\triangleright$\enspace\texttt{OnDigraphs({\mdseries\slshape digraph, trans})\index{OnDigraphs@\texttt{OnDigraphs}!for a digraph and a transformation}
\label{OnDigraphs:for a digraph and a transformation}
}\hfill{\scriptsize (operation)}}\\
\textbf{\indent Returns:\ }
A digraph.



 If \mbox{\texttt{\mdseries\slshape digraph}} is a digraph, and the second argument \mbox{\texttt{\mdseries\slshape perm}} is a \emph{permutation} of the vertices of \mbox{\texttt{\mdseries\slshape digraph}}, then this operation returns a digraph constructed by relabelling the
vertices of \mbox{\texttt{\mdseries\slshape digraph}} according to \mbox{\texttt{\mdseries\slshape perm}}. Note that for an automorphism \texttt{f} of a digraph, we have \texttt{OnDigraphs(\mbox{\texttt{\mdseries\slshape digraph}}, f) = }\mbox{\texttt{\mdseries\slshape digraph}}. 

 If the second argument is a \emph{transformation} \mbox{\texttt{\mdseries\slshape trans}} of the vertices of \mbox{\texttt{\mdseries\slshape digraph}}, then this operation returns a digraph constructed by transforming the source
and range of each edge according to \mbox{\texttt{\mdseries\slshape trans}}. Thus a vertex which does not appear in the image of \mbox{\texttt{\mdseries\slshape trans}} will be isolated in the returned digraph, and the returned digraph may contain
multiple edges, even if \mbox{\texttt{\mdseries\slshape digraph}} does not. If \mbox{\texttt{\mdseries\slshape trans}} is mathematically a permutation, then the result coincides with \texttt{OnDigraphs(\mbox{\texttt{\mdseries\slshape digraph}}, AsPermutation(\mbox{\texttt{\mdseries\slshape trans}}))}. 

 Note: \texttt{OnDigraphs} for a digraph and a permutation or transformation can also be used via the \texttt{\texttt{\symbol{92}}\texttt{\symbol{94}}} operator. 

 The \texttt{DigraphVertexLabels} (\ref{DigraphVertexLabels}) of \mbox{\texttt{\mdseries\slshape digraph}} will not be retained in the returned digraph. 

 If \mbox{\texttt{\mdseries\slshape digraph}} belongs to \texttt{IsMutableDigraph} (\ref{IsMutableDigraph}), then relabelling of the vertices is performed directly on \mbox{\texttt{\mdseries\slshape digraph}}. If \mbox{\texttt{\mdseries\slshape digraph}} belongs to \texttt{IsImmutableDigraph} (\ref{IsImmutableDigraph}), an immutable copy of \mbox{\texttt{\mdseries\slshape digraph}} with the vertices relabelled is returned.

 
\begin{Verbatim}[commandchars=!@|,fontsize=\small,frame=single,label=Example]
  !gapprompt@gap>| !gapinput@D := Digraph([[3], [1, 3, 5], [1], [1, 2, 4], [2, 3, 5]]);|
  <immutable digraph with 5 vertices, 11 edges>
  !gapprompt@gap>| !gapinput@new := OnDigraphs(D, (1, 2));|
  <immutable digraph with 5 vertices, 11 edges>
  !gapprompt@gap>| !gapinput@OutNeighbours(new);|
  [ [ 2, 3, 5 ], [ 3 ], [ 2 ], [ 2, 1, 4 ], [ 1, 3, 5 ] ]
  !gapprompt@gap>| !gapinput@D := Digraph([[2], [], [2]]);|
  <immutable digraph with 3 vertices, 2 edges>
  !gapprompt@gap>| !gapinput@t := Transformation([1, 2, 1]);;|
  !gapprompt@gap>| !gapinput@new := OnDigraphs(D, t);|
  <immutable multidigraph with 3 vertices, 2 edges>
  !gapprompt@gap>| !gapinput@OutNeighbours(new);|
  [ [ 2, 2 ], [  ], [  ] ]
  !gapprompt@gap>| !gapinput@ForAll(DigraphEdges(D),|
  !gapprompt@>| !gapinput@ e -> IsDigraphEdge(new, [e[1] ^ t, e[2] ^ t]));|
  true
\end{Verbatim}
 }

 

\subsection{\textcolor{Chapter }{\texttt{\symbol{92}}\texttt{\symbol{94}} (for a digraph and a permutation or transformation)}}
\logpage{[ 7, 1, 2 ]}\nobreak
\hyperdef{L}{X787FD1B479258B12}{}
{\noindent\textcolor{FuncColor}{$\triangleright$\enspace\texttt{\texttt{\symbol{92}}\texttt{\symbol{94}}({\mdseries\slshape digraph, permOrTrans})\index{\texttt{\symbol{92}}\texttt{\symbol{94}}@\texttt{\texttt{\symbol{92}}\texttt{\symbol{94}}}!for a digraph and a permutation or transformation}
\label{bSlash^:for a digraph and a permutation or transformation}
}\hfill{\scriptsize (operation)}}\\
\textbf{\indent Returns:\ }
A digraph.



 The \texttt{\texttt{\symbol{94}}} operator acts on digraphs with either a permutation or a transformation. For a
digraph \mbox{\texttt{\mdseries\slshape digraph}} and a permutation \mbox{\texttt{\mdseries\slshape perm}}, \texttt{digraph \texttt{\symbol{94}} perm} gives the same result as \texttt{OnDigraphs(digraph, perm)} {\textemdash} a digraph whose vertices have been relabelled according to the
permutation. 

 Similarly, if the second argument is a transformation \mbox{\texttt{\mdseries\slshape trans}}, then \texttt{digraph \texttt{\symbol{94}} trans} calls \texttt{OnDigraphs(digraph, trans)}. This allows a simple way to apply vertex relabelling or transformations
directly using the \texttt{\texttt{\symbol{94}}} symbol. 

 
\begin{Verbatim}[commandchars=!@|,fontsize=\small,frame=single,label=Example]
  !gapprompt@gap>| !gapinput@D := CycleDigraph(5);|
  <immutable cycle digraph with 5 vertices>
  !gapprompt@gap>| !gapinput@p := (1, 5)(2, 4);;|
  !gapprompt@gap>| !gapinput@D ^ p = DigraphReverse(D);|
  true
  !gapprompt@gap>| !gapinput@OnDigraphs(D, p) = D ^ p;|
  true
  !gapprompt@gap>| !gapinput@idp := ();;|
  !gapprompt@gap>| !gapinput@D ^ idp = D;|
  true
  !gapprompt@gap>| !gapinput@q := (1, 2, 3, 4, 5);;|
  !gapprompt@gap>| !gapinput@(D ^ q) ^ (q ^ -1) = D;|
  true
  !gapprompt@gap>| !gapinput@t := Transformation([2, 3, 4, 5, 1]);;|
  !gapprompt@gap>| !gapinput@D ^ t = OnDigraphs(D, t);|
  true
  !gapprompt@gap>| !gapinput@idt := Transformation([1, 2, 3, 4, 5]);;|
  !gapprompt@gap>| !gapinput@D ^ idt = D;|
  true
  !gapprompt@gap>| !gapinput@M := DigraphMutableCopy(D);;|
  !gapprompt@gap>| !gapinput@M ^ p = OnDigraphs(M, p);|
  true
\end{Verbatim}
 }

 

\subsection{\textcolor{Chapter }{OnMultiDigraphs}}
\logpage{[ 7, 1, 3 ]}\nobreak
\hyperdef{L}{X7AFBA7608498F9CE}{}
{\noindent\textcolor{FuncColor}{$\triangleright$\enspace\texttt{OnMultiDigraphs({\mdseries\slshape digraph, pair})\index{OnMultiDigraphs@\texttt{OnMultiDigraphs}}
\label{OnMultiDigraphs}
}\hfill{\scriptsize (operation)}}\\
\noindent\textcolor{FuncColor}{$\triangleright$\enspace\texttt{OnMultiDigraphs({\mdseries\slshape digraph, perm1, perm2})\index{OnMultiDigraphs@\texttt{OnMultiDigraphs}!for a digraph, perm, and perm}
\label{OnMultiDigraphs:for a digraph, perm, and perm}
}\hfill{\scriptsize (operation)}}\\
\textbf{\indent Returns:\ }
A digraph.



 If \mbox{\texttt{\mdseries\slshape digraph}} is a digraph, and \mbox{\texttt{\mdseries\slshape pair}} is a pair consisting of a permutation of the vertices and a permutation of the
edges of \mbox{\texttt{\mdseries\slshape digraph}}, then this operation returns a digraph constructed by relabelling the
vertices and edges of \mbox{\texttt{\mdseries\slshape digraph}} according to \mbox{\texttt{\mdseries\slshape perm[1]}} and \mbox{\texttt{\mdseries\slshape perm[2]}}, respectively. 

 In its second form, \texttt{OnMultiDigraphs} returns a digraph with vertices and edges permuted by \mbox{\texttt{\mdseries\slshape perm1}} and \mbox{\texttt{\mdseries\slshape perm2}}, respectively. 

 Note that \texttt{OnDigraphs(\mbox{\texttt{\mdseries\slshape digraph}}, perm)=OnMultiDigraphs(\mbox{\texttt{\mdseries\slshape digraph}}, [perm, ()])} where \texttt{perm} is a permutation of the vertices of \mbox{\texttt{\mdseries\slshape digraph}}. If you are only interested in the action of a permutation on the vertices of
a digraph, then you can use \texttt{OnDigraphs} instead of \texttt{OnMultiDigraphs}. If \mbox{\texttt{\mdseries\slshape digraph}} belongs to \texttt{IsMutableDigraph} (\ref{IsMutableDigraph}), then relabelling of the vertices is performed directly on \mbox{\texttt{\mdseries\slshape digraph}}. If \mbox{\texttt{\mdseries\slshape digraph}} belongs to \texttt{IsImmutableDigraph} (\ref{IsImmutableDigraph}), an immutable copy of \mbox{\texttt{\mdseries\slshape digraph}} with the vertices relabelled is returned.

 
\begin{Verbatim}[commandchars=!@|,fontsize=\small,frame=single,label=Example]
  !gapprompt@gap>| !gapinput@D1 := Digraph([|
  !gapprompt@>| !gapinput@[3, 6, 3], [], [3], [9, 10], [9], [],  [], [10, 4, 10], [], []]);|
  <immutable multidigraph with 10 vertices, 10 edges>
  !gapprompt@gap>| !gapinput@p := BlissCanonicalLabelling(D1);|
  [ (1,7)(3,6)(4,10)(5,9), () ]
  !gapprompt@gap>| !gapinput@D2 := OnMultiDigraphs(D1, p);|
  <immutable multidigraph with 10 vertices, 10 edges>
  !gapprompt@gap>| !gapinput@OutNeighbours(D2);|
  [ [  ], [  ], [  ], [  ], [  ], [ 6 ], [ 6, 3, 6 ], [ 4, 10, 4 ], 
    [ 5 ], [ 5, 4 ] ]
\end{Verbatim}
 }

 

\subsection{\textcolor{Chapter }{OnTuplesDigraphs (for a list of digraphs and a perm)}}
\logpage{[ 7, 1, 4 ]}\nobreak
\hyperdef{L}{X7F6A97E6870CDDA3}{}
{\noindent\textcolor{FuncColor}{$\triangleright$\enspace\texttt{OnTuplesDigraphs({\mdseries\slshape list, perm})\index{OnTuplesDigraphs@\texttt{OnTuplesDigraphs}!for a list of digraphs and a perm}
\label{OnTuplesDigraphs:for a list of digraphs and a perm}
}\hfill{\scriptsize (operation)}}\\
\noindent\textcolor{FuncColor}{$\triangleright$\enspace\texttt{OnSetsDigraphs({\mdseries\slshape list, perm})\index{OnSetsDigraphs@\texttt{OnSetsDigraphs}!for a set of digraphs and a perm}
\label{OnSetsDigraphs:for a set of digraphs and a perm}
}\hfill{\scriptsize (operation)}}\\
\textbf{\indent Returns:\ }
A list or set of digraphs.



 If \mbox{\texttt{\mdseries\slshape list}} is a list of digraphs, and \mbox{\texttt{\mdseries\slshape perm}} is a \emph{permutation} of the vertices of the digraphs in \mbox{\texttt{\mdseries\slshape list}}, then \texttt{OnTuplesDigraphs} returns a new list constructed by applying \mbox{\texttt{\mdseries\slshape perm}} via \texttt{OnDigraphs} (\ref{OnDigraphs:for a digraph and a perm}) to a copy (with the same mutability) of each entry of \mbox{\texttt{\mdseries\slshape list}} in turn. 

 More precisely, \texttt{OnTuplesDigraphs(\mbox{\texttt{\mdseries\slshape list}},\mbox{\texttt{\mdseries\slshape perm}})} is a list of length \texttt{Length(\mbox{\texttt{\mdseries\slshape list}})}, whose \texttt{i}\texttt{\symbol{45}}th entry is \texttt{OnDigraphs(DigraphCopy(\mbox{\texttt{\mdseries\slshape list}}[i]), \mbox{\texttt{\mdseries\slshape perm}})}. 

 If \mbox{\texttt{\mdseries\slshape list}} is moreover a \textsf{GAP} set (i.e. a duplicate\texttt{\symbol{45}}free sorted list), then \texttt{OnSetsDigraphs} returns the sorted output of \texttt{OnTuplesDigraphs}, which is therefore again a set. 
\begin{Verbatim}[commandchars=!@|,fontsize=\small,frame=single,label=Example]
  !gapprompt@gap>| !gapinput@list := [CycleDigraph(IsMutableDigraph, 6),|
  !gapprompt@>| !gapinput@            DigraphReverse(CycleDigraph(6))];|
  [ <mutable digraph with 6 vertices, 6 edges>, 
    <immutable digraph with 6 vertices, 6 edges> ]
  !gapprompt@gap>| !gapinput@p := (1, 6)(2, 5)(3, 4);;|
  !gapprompt@gap>| !gapinput@result_tuples := OnTuplesDigraphs(list, p);|
  [ <mutable digraph with 6 vertices, 6 edges>, 
    <immutable digraph with 6 vertices, 6 edges> ]
  !gapprompt@gap>| !gapinput@result_tuples[2] = OnDigraphs(list[2], p);|
  true
  !gapprompt@gap>| !gapinput@result_tuples = list;|
  false
  !gapprompt@gap>| !gapinput@result_tuples = Reversed(list);|
  true
  !gapprompt@gap>| !gapinput@result_sets := OnSetsDigraphs(list, p);|
  [ <immutable digraph with 6 vertices, 6 edges>, 
    <mutable digraph with 6 vertices, 6 edges> ]
  !gapprompt@gap>| !gapinput@result_sets = list;|
  true
\end{Verbatim}
 }

 }

 
\section{\textcolor{Chapter }{Isomorphisms and canonical labellings}}\label{Isomorphisms and canonical labellings}
\logpage{[ 7, 2, 0 ]}
\hyperdef{L}{X7E48B9F87A0F22D4}{}
{
  From version 0.11.0 of \textsf{Digraphs} it is possible to use either \href{http://www.tcs.tkk.fi/Software/bliss/} {bliss}  or \href{https://pallini.di.uniroma1.it} {nauty}  (via \textsf{NautyTracesInterface}) to calculate canonical labellings and automorphism groups of digraphs; see \cite{JK07} and \cite{MP14} for more details about \href{http://www.tcs.tkk.fi/Software/bliss/} {bliss}  and \href{https://pallini.di.uniroma1.it} {nauty} , respectively. 

\subsection{\textcolor{Chapter }{DigraphsUseNauty}}
\logpage{[ 7, 2, 1 ]}\nobreak
\hyperdef{L}{X83E593F3855B122E}{}
{\noindent\textcolor{FuncColor}{$\triangleright$\enspace\texttt{DigraphsUseNauty({\mdseries\slshape })\index{DigraphsUseNauty@\texttt{DigraphsUseNauty}}
\label{DigraphsUseNauty}
}\hfill{\scriptsize (function)}}\\
\noindent\textcolor{FuncColor}{$\triangleright$\enspace\texttt{DigraphsUseBliss({\mdseries\slshape })\index{DigraphsUseBliss@\texttt{DigraphsUseBliss}}
\label{DigraphsUseBliss}
}\hfill{\scriptsize (function)}}\\
\textbf{\indent Returns:\ }
Nothing.



 These functions can be used to specify whether \href{https://pallini.di.uniroma1.it} {nauty}  or \href{http://www.tcs.tkk.fi/Software/bliss/} {bliss}  should be used by default by \textsf{Digraphs}. If \textsf{NautyTracesInterface} is not available, then these functions do nothing. Otherwise, by calling \texttt{DigraphsUseNauty} subsequent computations will default to using \href{https://pallini.di.uniroma1.it} {nauty}  rather than \href{http://www.tcs.tkk.fi/Software/bliss/} {bliss} , where possible. 

 You can call these functions at any point in a \textsf{GAP} session, as many times as you like, it is guaranteed that existing digraphs
remain valid, and that comparison of existing digraphs and newly created
digraphs via \texttt{IsIsomorphicDigraph} (\ref{IsIsomorphicDigraph:for digraphs}), \texttt{IsIsomorphicDigraph} (\ref{IsIsomorphicDigraph:for digraphs and homogeneous lists}), \texttt{IsomorphismDigraphs} (\ref{IsomorphismDigraphs:for digraphs}), and \texttt{IsomorphismDigraphs} (\ref{IsomorphismDigraphs:for digraphs and homogeneous lists}) are also valid.

 It is also possible to compute the automorphism group of a specific digraph
using both \href{https://pallini.di.uniroma1.it} {nauty}  and \href{http://www.tcs.tkk.fi/Software/bliss/} {bliss}  using \texttt{NautyAutomorphismGroup} (\ref{NautyAutomorphismGroup}) and \texttt{BlissAutomorphismGroup} (\ref{BlissAutomorphismGroup:for a digraph}), respectively. }

 

\subsection{\textcolor{Chapter }{AutomorphismGroup (for a digraph)}}
\logpage{[ 7, 2, 2 ]}\nobreak
\hyperdef{L}{X858C32127A190175}{}
{\noindent\textcolor{FuncColor}{$\triangleright$\enspace\texttt{AutomorphismGroup({\mdseries\slshape digraph})\index{AutomorphismGroup@\texttt{AutomorphismGroup}!for a digraph}
\label{AutomorphismGroup:for a digraph}
}\hfill{\scriptsize (attribute)}}\\
\textbf{\indent Returns:\ }
A permutation group.



 If \mbox{\texttt{\mdseries\slshape digraph}} is a digraph, then this attribute contains the group of automorphisms of \mbox{\texttt{\mdseries\slshape digraph}}. An \emph{automorphism} of \mbox{\texttt{\mdseries\slshape digraph}} is an isomorphism from \mbox{\texttt{\mdseries\slshape digraph}} to itself. See \texttt{IsomorphismDigraphs} (\ref{IsomorphismDigraphs:for digraphs}) for more information about isomorphisms of digraphs. 

 If \mbox{\texttt{\mdseries\slshape digraph}} is not a multidigraph then the automorphism group is returned as a group of
permutations on the set of vertices of \mbox{\texttt{\mdseries\slshape digraph}}. 

 If \mbox{\texttt{\mdseries\slshape digraph}} is a multidigraph then the automorphism group is returned as the direct
product of a group of permutations on the set of vertices of \mbox{\texttt{\mdseries\slshape digraph}} with a group of permutations on the set of edges of \mbox{\texttt{\mdseries\slshape digraph}}. These groups can be accessed using \texttt{Projection} (\textbf{Reference: Projection for a domain and a positive integer}) on the returned group.

 By default, the automorphism group is found using \href{http://www.tcs.tkk.fi/Software/bliss/} {bliss}  by Tommi Junttila and Petteri Kaski. If \textsf{NautyTracesInterface} is available, then \href{https://pallini.di.uniroma1.it} {nauty}  by Brendan Mckay and Adolfo Piperno can be used instead; see \texttt{BlissAutomorphismGroup} (\ref{BlissAutomorphismGroup:for a digraph}), \texttt{NautyAutomorphismGroup} (\ref{NautyAutomorphismGroup}), \texttt{DigraphsUseBliss} (\ref{DigraphsUseBliss}), and \texttt{DigraphsUseNauty} (\ref{DigraphsUseNauty}). 

 If the argument \mbox{\texttt{\mdseries\slshape digraph}} is mutable, then the return value of this attribute is recomputed every time
it is called.

 
\begin{Verbatim}[commandchars=!@|,fontsize=\small,frame=single,label=Example]
  !gapprompt@gap>| !gapinput@johnson := DigraphFromGraph6String("E}lw");|
  <immutable symmetric digraph with 6 vertices, 24 edges>
  !gapprompt@gap>| !gapinput@G := AutomorphismGroup(johnson);|
  Group([ (3,4), (2,3)(4,5), (1,2)(5,6) ])
  !gapprompt@gap>| !gapinput@cycle := CycleDigraph(9);|
  <immutable cycle digraph with 9 vertices>
  !gapprompt@gap>| !gapinput@G := AutomorphismGroup(cycle);|
  Group([ (1,2,3,4,5,6,7,8,9) ])
  !gapprompt@gap>| !gapinput@IsCyclic(G) and Size(G) = 9;|
  true
\end{Verbatim}
 }

 

\subsection{\textcolor{Chapter }{BlissAutomorphismGroup (for a digraph)}}
\logpage{[ 7, 2, 3 ]}\nobreak
\hyperdef{L}{X7E7B0D88865A89F6}{}
{\noindent\textcolor{FuncColor}{$\triangleright$\enspace\texttt{BlissAutomorphismGroup({\mdseries\slshape digraph})\index{BlissAutomorphismGroup@\texttt{BlissAutomorphismGroup}!for a digraph}
\label{BlissAutomorphismGroup:for a digraph}
}\hfill{\scriptsize (attribute)}}\\
\noindent\textcolor{FuncColor}{$\triangleright$\enspace\texttt{BlissAutomorphismGroup({\mdseries\slshape digraph, vertex{\textunderscore}colours})\index{BlissAutomorphismGroup@\texttt{BlissAutomorphismGroup}!for a digraph and homogeneous list}
\label{BlissAutomorphismGroup:for a digraph and homogeneous list}
}\hfill{\scriptsize (operation)}}\\
\noindent\textcolor{FuncColor}{$\triangleright$\enspace\texttt{BlissAutomorphismGroup({\mdseries\slshape digraph, vertex{\textunderscore}colours, edge{\textunderscore}colours})\index{BlissAutomorphismGroup@\texttt{BlissAutomorphismGroup}!for a digraph, homogeneous list, and list}
\label{BlissAutomorphismGroup:for a digraph, homogeneous list, and list}
}\hfill{\scriptsize (operation)}}\\
\textbf{\indent Returns:\ }
A permutation group.



 If \mbox{\texttt{\mdseries\slshape digraph}} is a digraph, then this attribute contains the group of automorphisms of \mbox{\texttt{\mdseries\slshape digraph}} as calculated using \href{http://www.tcs.tkk.fi/Software/bliss/} {bliss}  by Tommi Junttila and Petteri Kaski. 

 The attribute \texttt{AutomorphismGroup} (\ref{AutomorphismGroup:for a digraph}) and operation \texttt{AutomorphismGroup} (\ref{AutomorphismGroup:for a digraph and a homogeneous list}) returns the value of either \texttt{BlissAutomorphismGroup} or \texttt{NautyAutomorphismGroup} (\ref{NautyAutomorphismGroup}). These groups are, of course, equal but their generating sets may differ. 

 The attribute \texttt{AutomorphismGroup} (\ref{AutomorphismGroup:for a digraph, homogeneous list, and list}) returns the value of \texttt{BlissAutomorphismGroup} as it is not implemented for \href{https://pallini.di.uniroma1.it} {nauty}  The requirements for the optional arguments \mbox{\texttt{\mdseries\slshape vertex{\textunderscore}colours}} and \mbox{\texttt{\mdseries\slshape edge{\textunderscore}colours}} are documented in \texttt{AutomorphismGroup} (\ref{AutomorphismGroup:for a digraph, homogeneous list, and list}). See also \texttt{DigraphsUseBliss} (\ref{DigraphsUseBliss}), and \texttt{DigraphsUseNauty} (\ref{DigraphsUseNauty}). 

 If the argument \mbox{\texttt{\mdseries\slshape digraph}} is mutable, then the return value of this attribute is recomputed every time
it is called.

 
\begin{Verbatim}[commandchars=!@|,fontsize=\small,frame=single,label=Example]
  !gapprompt@gap>| !gapinput@G := BlissAutomorphismGroup(JohnsonDigraph(5, 2));;|
  !gapprompt@gap>| !gapinput@IsSymmetricGroup(G);|
  true
  !gapprompt@gap>| !gapinput@Size(G);|
  120
\end{Verbatim}
 }

 

\subsection{\textcolor{Chapter }{NautyAutomorphismGroup}}
\logpage{[ 7, 2, 4 ]}\nobreak
\hyperdef{L}{X857758B18144C0CD}{}
{\noindent\textcolor{FuncColor}{$\triangleright$\enspace\texttt{NautyAutomorphismGroup({\mdseries\slshape digraph[, vert{\textunderscore}colours]})\index{NautyAutomorphismGroup@\texttt{NautyAutomorphismGroup}}
\label{NautyAutomorphismGroup}
}\hfill{\scriptsize (attribute)}}\\
\textbf{\indent Returns:\ }
A permutation group.



 If \mbox{\texttt{\mdseries\slshape digraph}} is a digraph, then this attribute contains the group of automorphisms of \mbox{\texttt{\mdseries\slshape digraph}} as calculated using \href{https://pallini.di.uniroma1.it} {nauty}  by Brendan Mckay and Adolfo Piperno via \textsf{NautyTracesInterface}.  The attribute \texttt{AutomorphismGroup} (\ref{AutomorphismGroup:for a digraph}) and operation \texttt{AutomorphismGroup} (\ref{AutomorphismGroup:for a digraph and a homogeneous list}) returns the value of either \texttt{NautyAutomorphismGroup} or \texttt{BlissAutomorphismGroup} (\ref{BlissAutomorphismGroup:for a digraph}). These groups are, of course, equal but their generating sets may differ.

 See also \texttt{DigraphsUseBliss} (\ref{DigraphsUseBliss}), and \texttt{DigraphsUseNauty} (\ref{DigraphsUseNauty}). 

 If the argument \mbox{\texttt{\mdseries\slshape digraph}} is mutable, then the return value of this attribute is recomputed every time
it is called.

 
\begin{Verbatim}[commandchars=!@|,fontsize=\small,frame=single,label=Example]
  !gapprompt@gap>| !gapinput@NautyAutomorphismGroup(JohnsonDigraph(5, 2));|
  Group([ (3,4)(6,7)(8,9), (2,3)(5,6)(9,10), (2,5)(3,6)(4,7),
   (1,2)(6,8)(7,9) ])
\end{Verbatim}
 }

 

\subsection{\textcolor{Chapter }{AutomorphismGroup (for a digraph and a homogeneous list)}}
\logpage{[ 7, 2, 5 ]}\nobreak
\hyperdef{L}{X877732B1783C391B}{}
{\noindent\textcolor{FuncColor}{$\triangleright$\enspace\texttt{AutomorphismGroup({\mdseries\slshape digraph, vert{\textunderscore}colours})\index{AutomorphismGroup@\texttt{AutomorphismGroup}!for a digraph and a homogeneous list}
\label{AutomorphismGroup:for a digraph and a homogeneous list}
}\hfill{\scriptsize (operation)}}\\
\textbf{\indent Returns:\ }
A permutation group.



 This operation computes the automorphism group of a
vertex\texttt{\symbol{45}}coloured digraph. A
vertex\texttt{\symbol{45}}coloured digraph can be specified by its underlying
digraph \mbox{\texttt{\mdseries\slshape digraph}} and its colouring \mbox{\texttt{\mdseries\slshape vert{\textunderscore}colours}}. Let \texttt{n} be the number of vertices of \mbox{\texttt{\mdseries\slshape digraph}}. The colouring \mbox{\texttt{\mdseries\slshape vert{\textunderscore}colours}} may have one of the following two forms: 
\begin{itemize}
\item  a list of \texttt{n} integers, where \mbox{\texttt{\mdseries\slshape vert{\textunderscore}colours}}\texttt{[i]} is the colour of vertex \texttt{i}, using the colours \texttt{[1 .. m]} for some \texttt{m {\textless}= n}; or 
\item  a list of non\texttt{\symbol{45}}empty disjoint lists whose union is \texttt{DigraphVertices(\mbox{\texttt{\mdseries\slshape digraph}})}, such that \mbox{\texttt{\mdseries\slshape vert{\textunderscore}colours}}\texttt{[i]} is the list of all vertices with colour \texttt{i}. 
\end{itemize}
 The \emph{automorphism group} of a coloured digraph \mbox{\texttt{\mdseries\slshape digraph}} with colouring \mbox{\texttt{\mdseries\slshape vert{\textunderscore}colours}} is the group consisting of its automorphisms; an \emph{automorphism} of \mbox{\texttt{\mdseries\slshape digraph}} is an isomorphism of coloured digraphs from \mbox{\texttt{\mdseries\slshape digraph}} to itself. This group is equal to the subgroup of \texttt{AutomorphismGroup(\mbox{\texttt{\mdseries\slshape digraph}})} consisting of those automorphisms that preserve the colouring specified by \mbox{\texttt{\mdseries\slshape vert{\textunderscore}colours}}. See \texttt{AutomorphismGroup} (\ref{AutomorphismGroup:for a digraph}), and see \texttt{IsomorphismDigraphs} (\ref{IsomorphismDigraphs:for digraphs and homogeneous lists}) for more information about isomorphisms of coloured digraphs. 

 If \mbox{\texttt{\mdseries\slshape digraph}} is not a multidigraph then the automorphism group is returned as a group of
permutations on the set of vertices of \mbox{\texttt{\mdseries\slshape digraph}}. 

 If \mbox{\texttt{\mdseries\slshape digraph}} is a multidigraph then the automorphism group is returned as the direct
product of a group of permutations on the set of vertices of \mbox{\texttt{\mdseries\slshape digraph}} with a group of permutations on the set of edges of \mbox{\texttt{\mdseries\slshape digraph}}. These groups can be accessed using \texttt{Projection} (\textbf{Reference: Projection for a domain and a positive integer}) on the returned group.

 By default, the automorphism group is found using \href{http://www.tcs.tkk.fi/Software/bliss/} {bliss}  by Tommi Junttila and Petteri Kaski. If \textsf{NautyTracesInterface} is available, then \href{https://pallini.di.uniroma1.it} {nauty}  by Brendan Mckay and Adolfo Piperno can be used instead; see \texttt{BlissAutomorphismGroup} (\ref{BlissAutomorphismGroup:for a digraph and homogeneous list}), \texttt{NautyAutomorphismGroup} (\ref{NautyAutomorphismGroup}), \texttt{DigraphsUseBliss} (\ref{DigraphsUseBliss}), and \texttt{DigraphsUseNauty} (\ref{DigraphsUseNauty}). 
\begin{Verbatim}[commandchars=!@|,fontsize=\small,frame=single,label=Example]
  !gapprompt@gap>| !gapinput@cycle := CycleDigraph(9);|
  <immutable cycle digraph with 9 vertices>
  !gapprompt@gap>| !gapinput@G := AutomorphismGroup(cycle);;|
  !gapprompt@gap>| !gapinput@IsCyclic(G) and Size(G) = 9;|
  true
  !gapprompt@gap>| !gapinput@colours := [[1, 4, 7], [2, 5, 8], [3, 6, 9]];;|
  !gapprompt@gap>| !gapinput@H := AutomorphismGroup(cycle, colours);;|
  !gapprompt@gap>| !gapinput@Size(H);|
  3
  !gapprompt@gap>| !gapinput@H = AutomorphismGroup(cycle, [1, 2, 3, 1, 2, 3, 1, 2, 3]);|
  true
  !gapprompt@gap>| !gapinput@H = SubgroupByProperty(G, p -> OnTuplesSets(colours, p) = colours);|
  true
  !gapprompt@gap>| !gapinput@IsTrivial(AutomorphismGroup(cycle, [1, 1, 2, 2, 2, 2, 2, 2, 2]));|
  true
\end{Verbatim}
 }

 

\subsection{\textcolor{Chapter }{AutomorphismGroup (for a digraph, homogeneous list, and list)}}
\logpage{[ 7, 2, 6 ]}\nobreak
\hyperdef{L}{X7DA2C4FE837FFE01}{}
{\noindent\textcolor{FuncColor}{$\triangleright$\enspace\texttt{AutomorphismGroup({\mdseries\slshape digraph, vert{\textunderscore}colours, edge{\textunderscore}colours})\index{AutomorphismGroup@\texttt{AutomorphismGroup}!for a digraph, homogeneous list, and list}
\label{AutomorphismGroup:for a digraph, homogeneous list, and list}
}\hfill{\scriptsize (operation)}}\\
\textbf{\indent Returns:\ }
A permutation group.



 This operation computes the automorphism group of a vertex\texttt{\symbol{45}}
and/or edge\texttt{\symbol{45}}coloured digraph. A coloured digraph can be
specified by its underlying digraph \mbox{\texttt{\mdseries\slshape digraph}} and colourings \mbox{\texttt{\mdseries\slshape vert{\textunderscore}colours}}, \mbox{\texttt{\mdseries\slshape edge{\textunderscore}colours}}. Let \texttt{n} be the number of vertices of \mbox{\texttt{\mdseries\slshape digraph}}. The colourings must have the following forms: 
\begin{itemize}
\item  \mbox{\texttt{\mdseries\slshape vert{\textunderscore}colours}} must be \texttt{fail} or a list of \texttt{n} integers, where \mbox{\texttt{\mdseries\slshape vert{\textunderscore}colours}}\texttt{[i]} is the colour of vertex \texttt{i}, using the colours \texttt{[1 .. m]} for some \texttt{m {\textless}= n}; 
\item  \mbox{\texttt{\mdseries\slshape edge{\textunderscore}colours}} must be \texttt{fail} or a list of \texttt{n} lists of integers of the same shape as \texttt{OutNeighbours(digraph)}, where \mbox{\texttt{\mdseries\slshape edge{\textunderscore}colours}}\texttt{[i][j]} is the colour of the edge \texttt{OutNeighbours(digraph)[i][j]}, using the colours \texttt{[1 .. k]} for some \texttt{k {\textless}= n}; 
\end{itemize}
 Giving \mbox{\texttt{\mdseries\slshape vert{\textunderscore}colours}} [\mbox{\texttt{\mdseries\slshape edge{\textunderscore}colours}}] as \texttt{fail} is equivalent to setting all vertices [edges] to be the same colour. 

 Unlike \texttt{AutomorphismGroup} (\ref{AutomorphismGroup:for a digraph}), it is possible to obtain the automorphism group of an
edge\texttt{\symbol{45}}coloured multidigraph (see \texttt{IsMultiDigraph} (\ref{IsMultiDigraph})) when no two edges share the same source, range, and colour. The \emph{automorphism group} of a vertex/edge\texttt{\symbol{45}}coloured digraph \mbox{\texttt{\mdseries\slshape digraph}} with colouring \mbox{\texttt{\mdseries\slshape c}} is the group consisting of its vertex/edge\texttt{\symbol{45}}colour
preserving automorphisms; an \emph{automorphism} of \mbox{\texttt{\mdseries\slshape digraph}} is an isomorphism of vertex/edge\texttt{\symbol{45}}coloured digraphs from \mbox{\texttt{\mdseries\slshape digraph}} to itself. This group is equal to the subgroup of \texttt{AutomorphismGroup(\mbox{\texttt{\mdseries\slshape digraph}})} consisting of those automorphisms that preserve the colouring specified by \mbox{\texttt{\mdseries\slshape colours}}. See \texttt{AutomorphismGroup} (\ref{AutomorphismGroup:for a digraph}), and see \texttt{IsomorphismDigraphs} (\ref{IsomorphismDigraphs:for digraphs and homogeneous lists}) for more information about isomorphisms of coloured digraphs. 

 If \mbox{\texttt{\mdseries\slshape digraph}} is not a multidigraph then the automorphism group is returned as a group of
permutations on the set of vertices of \mbox{\texttt{\mdseries\slshape digraph}}. 

 If \mbox{\texttt{\mdseries\slshape digraph}} is a multidigraph then the automorphism group is returned as the direct
product of a group of permutations on the set of vertices of \mbox{\texttt{\mdseries\slshape digraph}} with a group of permutations on the set of edges of \mbox{\texttt{\mdseries\slshape digraph}}. These groups can be accessed using \texttt{Projection} (\textbf{Reference: Projection for a domain and a positive integer}) on the returned group.

 By default, the automorphism group is found using \href{http://www.tcs.tkk.fi/Software/bliss/} {bliss}  by Tommi Junttila and Petteri Kaski. If \textsf{NautyTracesInterface} is available, then \href{https://pallini.di.uniroma1.it} {nauty}  by Brendan Mckay and Adolfo Piperno can be used instead; see \texttt{BlissAutomorphismGroup} (\ref{BlissAutomorphismGroup:for a digraph, homogeneous list, and list}), \texttt{NautyAutomorphismGroup} (\ref{NautyAutomorphismGroup}), \texttt{DigraphsUseBliss} (\ref{DigraphsUseBliss}), and \texttt{DigraphsUseNauty} (\ref{DigraphsUseNauty}). 
\begin{Verbatim}[commandchars=!@|,fontsize=\small,frame=single,label=Example]
  !gapprompt@gap>| !gapinput@cycle := CycleDigraph(12);|
  <immutable cycle digraph with 12 vertices>
  !gapprompt@gap>| !gapinput@vert_colours := List([1 .. 12], x -> x mod 3 + 1);;|
  !gapprompt@gap>| !gapinput@edge_colours := List([1 .. 12], x -> [x mod 2 + 1]);;|
  !gapprompt@gap>| !gapinput@Size(AutomorphismGroup(cycle));|
  12
  !gapprompt@gap>| !gapinput@Size(AutomorphismGroup(cycle, vert_colours));|
  4
  !gapprompt@gap>| !gapinput@Size(AutomorphismGroup(cycle, fail, edge_colours));|
  6
  !gapprompt@gap>| !gapinput@Size(AutomorphismGroup(cycle, vert_colours, edge_colours));|
  2
  !gapprompt@gap>| !gapinput@IsTrivial(AutomorphismGroup(cycle,|
  !gapprompt@>| !gapinput@vert_colours, List([1 .. 12], x -> [x mod 4 + 1])));|
  true
\end{Verbatim}
 }

 

\subsection{\textcolor{Chapter }{BlissCanonicalLabelling (for a digraph)}}
\logpage{[ 7, 2, 7 ]}\nobreak
\hyperdef{L}{X7D676FE67A6684FF}{}
{\noindent\textcolor{FuncColor}{$\triangleright$\enspace\texttt{BlissCanonicalLabelling({\mdseries\slshape digraph})\index{BlissCanonicalLabelling@\texttt{BlissCanonicalLabelling}!for a digraph}
\label{BlissCanonicalLabelling:for a digraph}
}\hfill{\scriptsize (attribute)}}\\
\noindent\textcolor{FuncColor}{$\triangleright$\enspace\texttt{NautyCanonicalLabelling({\mdseries\slshape digraph})\index{NautyCanonicalLabelling@\texttt{NautyCanonicalLabelling}!for a digraph}
\label{NautyCanonicalLabelling:for a digraph}
}\hfill{\scriptsize (attribute)}}\\
\textbf{\indent Returns:\ }
A permutation, or a list of two permutations.



 A function $\rho$ that maps a digraph to a digraph is a \emph{canonical representative map} if the following two conditions hold for all digraphs $G$ and $H$: 

 
\begin{itemize}
\item  $\rho(G)$ and $G$ are isomorphic as digraphs; and 
\item  $\rho(G)=\rho(H)$ if and only if $G$ and $H$ are isomorphic as digraphs. 
\end{itemize}
 A \emph{canonical labelling} of a digraph $G$ (under $\rho$) is an isomorphism of $G$ onto its \emph{canonical representative}, $\rho(G)$. See \texttt{IsomorphismDigraphs} (\ref{IsomorphismDigraphs:for digraphs}) for more information about isomorphisms of digraphs. 

 \texttt{BlissCanonicalLabelling} returns a canonical labelling of the digraph \mbox{\texttt{\mdseries\slshape digraph}} found using \href{http://www.tcs.tkk.fi/Software/bliss/} {bliss}  by Tommi Junttila and Petteri Kaski. \texttt{NautyCanonicalLabelling} returns a canonical labelling of the digraph \mbox{\texttt{\mdseries\slshape digraph}} found using \href{https://pallini.di.uniroma1.it} {nauty}  by Brendan McKay and Adolfo Piperno. Note that the canonical labellings
returned by \href{http://www.tcs.tkk.fi/Software/bliss/} {bliss}  and \href{https://pallini.di.uniroma1.it} {nauty}  are not usually the same (and may depend of the version used).

 \texttt{BlissCanonicalLabelling} can only be computed if \mbox{\texttt{\mdseries\slshape digraph}} has no multiple edges; see \texttt{IsMultiDigraph} (\ref{IsMultiDigraph}). 

 
\begin{Verbatim}[commandchars=!@|,fontsize=\small,frame=single,label=Example]
  !gapprompt@gap>| !gapinput@digraph1 := DigraphFromDiSparse6String(".ImNS_AiB?qRN");|
  <immutable digraph with 10 vertices, 8 edges>
  !gapprompt@gap>| !gapinput@BlissCanonicalLabelling(digraph1);|
  (1,9,5,7)(3,6,4,10)
  !gapprompt@gap>| !gapinput@p := (1, 2, 7, 5)(3, 9)(6, 10, 8);;|
  !gapprompt@gap>| !gapinput@digraph2 := OnDigraphs(digraph1, p);|
  <immutable digraph with 10 vertices, 8 edges>
  !gapprompt@gap>| !gapinput@digraph1 = digraph2;|
  false
  !gapprompt@gap>| !gapinput@OnDigraphs(digraph1, BlissCanonicalLabelling(digraph1)) =|
  !gapprompt@>| !gapinput@   OnDigraphs(digraph2, BlissCanonicalLabelling(digraph2));|
  true
\end{Verbatim}
 }

 

\subsection{\textcolor{Chapter }{BlissCanonicalLabelling (for a digraph and a list)}}
\logpage{[ 7, 2, 8 ]}\nobreak
\hyperdef{L}{X87DA265D803DB337}{}
{\noindent\textcolor{FuncColor}{$\triangleright$\enspace\texttt{BlissCanonicalLabelling({\mdseries\slshape digraph, colours})\index{BlissCanonicalLabelling@\texttt{BlissCanonicalLabelling}!for a digraph and a list}
\label{BlissCanonicalLabelling:for a digraph and a list}
}\hfill{\scriptsize (operation)}}\\
\noindent\textcolor{FuncColor}{$\triangleright$\enspace\texttt{NautyCanonicalLabelling({\mdseries\slshape digraph, colours})\index{NautyCanonicalLabelling@\texttt{NautyCanonicalLabelling}!for a digraph and a list}
\label{NautyCanonicalLabelling:for a digraph and a list}
}\hfill{\scriptsize (operation)}}\\
\textbf{\indent Returns:\ }
A permutation.



 A function $\rho$ that maps a coloured digraph to a coloured digraph is a \emph{canonical representative map} if the following two conditions hold for all coloured digraphs $G$ and $H$: 
\begin{itemize}
\item  $\rho(G)$ and $G$ are isomorphic as coloured digraphs; and 
\item  $\rho(G)=\rho(H)$ if and only if $G$ and $H$ are isomorphic as coloured digraphs. 
\end{itemize}
 A \emph{canonical labelling} of a coloured digraph $G$ (under $\rho$) is an isomorphism of $G$ onto its \emph{canonical representative}, $\rho(G)$. See \texttt{IsomorphismDigraphs} (\ref{IsomorphismDigraphs:for digraphs and homogeneous lists}) for more information about isomorphisms of coloured digraphs. 

 A coloured digraph can be specified by its underlying digraph \mbox{\texttt{\mdseries\slshape digraph}} and its colouring \mbox{\texttt{\mdseries\slshape colours}}. Let \texttt{n} be the number of vertices of \mbox{\texttt{\mdseries\slshape digraph}}. The colouring \mbox{\texttt{\mdseries\slshape colours}} may have one of the following two forms: 

 
\begin{itemize}
\item  a list of \texttt{n} integers, where \mbox{\texttt{\mdseries\slshape colours}}\texttt{[i]} is the colour of vertex \texttt{i}, using the colours \texttt{[1 .. m]} for some \texttt{m {\textless}= n}; or 
\item  a list of non\texttt{\symbol{45}}empty disjoint lists whose union is \texttt{DigraphVertices(\mbox{\texttt{\mdseries\slshape digraph}})}, such that \mbox{\texttt{\mdseries\slshape colours}}\texttt{[i]} is the list of all vertices with colour \texttt{i}. 
\end{itemize}
 If \mbox{\texttt{\mdseries\slshape digraph}} and \mbox{\texttt{\mdseries\slshape colours}} together form a coloured digraph, \texttt{BlissCanonicalLabelling} returns a canonical labelling of the digraph \mbox{\texttt{\mdseries\slshape digraph}} found using \href{http://www.tcs.tkk.fi/Software/bliss/} {bliss}  by Tommi Junttila and Petteri Kaski. Similarly, \texttt{NautyCanonicalLabelling} returns a canonical labelling of the digraph \mbox{\texttt{\mdseries\slshape digraph}} found using \href{https://pallini.di.uniroma1.it} {nauty}  by Brendan McKay and Adolfo Piperno. Note that the canonical labellings
returned by \href{http://www.tcs.tkk.fi/Software/bliss/} {bliss}  and \href{https://pallini.di.uniroma1.it} {nauty}  are not usually the same (and may depend of the version used).

 \texttt{BlissCanonicalLabelling} can only be computed if \mbox{\texttt{\mdseries\slshape digraph}} has no multiple edges; see \texttt{IsMultiDigraph} (\ref{IsMultiDigraph}). The canonical labelling of \mbox{\texttt{\mdseries\slshape digraph}} is given as a permutation of its vertices. The canonical representative of \mbox{\texttt{\mdseries\slshape digraph}} can be created from \mbox{\texttt{\mdseries\slshape digraph}} and its canonical labelling \texttt{p} by using the operation \texttt{OnDigraphs} (\ref{OnDigraphs:for a digraph and a perm}): 
\begin{Verbatim}[commandchars=!@|,fontsize=\small,frame=single,label=Example]
  !gapprompt@gap>| !gapinput@OnDigraphs(digraph, p);|
\end{Verbatim}
 The colouring of the canonical representative can easily be constructed. A
vertex \texttt{v} (in \mbox{\texttt{\mdseries\slshape digraph}}) has colour \texttt{i} if and only if the vertex \texttt{v \texttt{\symbol{94}} p} (in the canonical representative) has colour \texttt{i}, where \texttt{p} is the permutation of the canonical labelling that acts on the vertices of \mbox{\texttt{\mdseries\slshape digraph}}. In particular, if \mbox{\texttt{\mdseries\slshape colours}} has the first form that is described above, then the colouring of the
canonical representative is given by: 
\begin{Verbatim}[commandchars=!@|,fontsize=\small,frame=single,label=Example]
  !gapprompt@gap>| !gapinput@List(DigraphVertices(digraph), i -> colours[i / p]);|
\end{Verbatim}
 On the other hand, if \mbox{\texttt{\mdseries\slshape colours}} has the second form above, then the canonical representative has colouring: 
\begin{Verbatim}[commandchars=!@|,fontsize=\small,frame=single,label=Example]
  !gapprompt@gap>| !gapinput@OnTuplesSets(colours, p);|
\end{Verbatim}
 

 
\begin{Verbatim}[commandchars=!@|,fontsize=\small,frame=single,label=Example]
  !gapprompt@gap>| !gapinput@digraph := DigraphFromDiSparse6String(".ImNS_AiB?qRN");|
  <immutable digraph with 10 vertices, 8 edges>
  !gapprompt@gap>| !gapinput@colours := [[1, 2, 8, 9, 10], [3, 4, 5, 6, 7]];;|
  !gapprompt@gap>| !gapinput@p := BlissCanonicalLabelling(digraph, colours);|
  (1,5,8,4,10,3,9)(6,7)
  !gapprompt@gap>| !gapinput@OnDigraphs(digraph, p);|
  <immutable digraph with 10 vertices, 8 edges>
  !gapprompt@gap>| !gapinput@OnTuplesSets(colours, p);|
  [ [ 1, 2, 3, 4, 5 ], [ 6, 7, 8, 9, 10 ] ]
  !gapprompt@gap>| !gapinput@colours := [1, 1, 1, 1, 2, 3, 1, 3, 2, 1];;|
  !gapprompt@gap>| !gapinput@p := BlissCanonicalLabelling(digraph, colours);|
  (1,6,9,7)(3,4,5,8,10)
  !gapprompt@gap>| !gapinput@OnDigraphs(digraph, p);|
  <immutable digraph with 10 vertices, 8 edges>
  !gapprompt@gap>| !gapinput@List(DigraphVertices(digraph), i -> colours[i / p]);|
  [ 1, 1, 1, 1, 1, 1, 2, 2, 3, 3 ]
\end{Verbatim}
 }

 

\subsection{\textcolor{Chapter }{BlissCanonicalDigraph}}
\logpage{[ 7, 2, 9 ]}\nobreak
\hyperdef{L}{X877B1D377EC197D7}{}
{\noindent\textcolor{FuncColor}{$\triangleright$\enspace\texttt{BlissCanonicalDigraph({\mdseries\slshape digraph})\index{BlissCanonicalDigraph@\texttt{BlissCanonicalDigraph}}
\label{BlissCanonicalDigraph}
}\hfill{\scriptsize (attribute)}}\\
\noindent\textcolor{FuncColor}{$\triangleright$\enspace\texttt{NautyCanonicalDigraph({\mdseries\slshape digraph})\index{NautyCanonicalDigraph@\texttt{NautyCanonicalDigraph}}
\label{NautyCanonicalDigraph}
}\hfill{\scriptsize (attribute)}}\\
\textbf{\indent Returns:\ }
A digraph.



 The attribute \texttt{BlissCanonicalLabelling} returns the canonical representative found by applying \texttt{BlissCanonicalLabelling} (\ref{BlissCanonicalLabelling:for a digraph}). The digraph returned is canonical in the sense that 
\begin{itemize}
\item  \texttt{BlissCanonicalDigraph(\mbox{\texttt{\mdseries\slshape digraph}})} and \mbox{\texttt{\mdseries\slshape digraph}} are isomorphic as digraphs; and 
\item  If \texttt{gr} is any digraph then \texttt{BlissCanonicalDigraph(gr)} and \texttt{BlissCanonicalDigraph(\mbox{\texttt{\mdseries\slshape digraph}})} are equal if and only if \texttt{gr} and \mbox{\texttt{\mdseries\slshape digraph}} are isomorphic as digraphs. 
\end{itemize}
 Analogously, the attribute \texttt{NautyCanonicalLabelling} returns the canonical representative found by applying \texttt{NautyCanonicalLabelling} (\ref{NautyCanonicalLabelling:for a digraph}). 

 If the argument \mbox{\texttt{\mdseries\slshape digraph}} is mutable, then the return value of this attribute is recomputed every time
it is called.

 
\begin{Verbatim}[commandchars=!@|,fontsize=\small,frame=single,label=Example]
  !gapprompt@gap>| !gapinput@digraph := Digraph([[1], [2, 3], [3], [1, 2, 3]]);|
  <immutable digraph with 4 vertices, 7 edges>
  !gapprompt@gap>| !gapinput@canon := BlissCanonicalDigraph(digraph);|
  <immutable digraph with 4 vertices, 7 edges>
  !gapprompt@gap>| !gapinput@OutNeighbours(canon);|
  [ [ 1 ], [ 2 ], [ 3, 2 ], [ 1, 3, 2 ] ]
\end{Verbatim}
 }

 

\subsection{\textcolor{Chapter }{DigraphGroup}}
\logpage{[ 7, 2, 10 ]}\nobreak
\hyperdef{L}{X803ACEDA7BBAC5B3}{}
{\noindent\textcolor{FuncColor}{$\triangleright$\enspace\texttt{DigraphGroup({\mdseries\slshape digraph})\index{DigraphGroup@\texttt{DigraphGroup}}
\label{DigraphGroup}
}\hfill{\scriptsize (attribute)}}\\
\textbf{\indent Returns:\ }
A permutation group.



 If \mbox{\texttt{\mdseries\slshape digraph}} is immutable and was created knowing a subgroup of its automorphism group,
then this group is stored in the attribute \texttt{DigraphGroup}. If \mbox{\texttt{\mdseries\slshape digraph}} is mutable, or was not created knowing a subgroup of its automorphism group,
then \texttt{DigraphGroup} returns the entire automorphism group of \mbox{\texttt{\mdseries\slshape digraph}}. Note that if \mbox{\texttt{\mdseries\slshape digraph}} is mutable, then the automorphism group is recomputed every time this function
is called. 

 Note that certain other constructor operations such as \texttt{CayleyDigraph} (\ref{CayleyDigraph}), \texttt{BipartiteDoubleDigraph} (\ref{BipartiteDoubleDigraph}), and \texttt{DoubleDigraph} (\ref{DoubleDigraph}), may not require a group as one of the arguments, but use the standard
constructor method using a group, and hence set the \texttt{DigraphGroup} attribute for the resulting digraph. 
\begin{Verbatim}[commandchars=!@|,fontsize=\small,frame=single,label=Example]
  !gapprompt@gap>| !gapinput@n := 4;;|
  !gapprompt@gap>| !gapinput@adj := function(x, y)|
  !gapprompt@>| !gapinput@     return (((x - y) mod n) = 1) or (((x - y) mod n) = n - 1);|
  !gapprompt@>| !gapinput@   end;;|
  !gapprompt@gap>| !gapinput@group := CyclicGroup(IsPermGroup, n);|
  Group([ (1,2,3,4) ])
  !gapprompt@gap>| !gapinput@D := Digraph(IsMutableDigraph, group, [1 .. n], \^, adj);|
  <mutable digraph with 4 vertices, 8 edges>
  !gapprompt@gap>| !gapinput@HasDigraphGroup(D);|
  false
  !gapprompt@gap>| !gapinput@DigraphGroup(D);|
  Group([ (2,4), (1,2)(3,4) ])
  !gapprompt@gap>| !gapinput@AutomorphismGroup(D);|
  Group([ (2,4), (1,2)(3,4) ])
  !gapprompt@gap>| !gapinput@D := Digraph(group, [1 .. n], \^, adj);|
  <immutable digraph with 4 vertices, 8 edges>
  !gapprompt@gap>| !gapinput@HasDigraphGroup(D);|
  true
  !gapprompt@gap>| !gapinput@DigraphGroup(D);|
  Group([ (1,2,3,4) ])
  !gapprompt@gap>| !gapinput@D := DoubleDigraph(D);|
  <immutable digraph with 8 vertices, 32 edges>
  !gapprompt@gap>| !gapinput@HasDigraphGroup(D);|
  true
  !gapprompt@gap>| !gapinput@DigraphGroup(D);|
  Group([ (1,2,3,4)(5,6,7,8), (1,5)(2,6)(3,7)(4,8) ])
  !gapprompt@gap>| !gapinput@AutomorphismGroup(D) =|
  !gapprompt@>| !gapinput@Group([(6, 8), (5, 7), (4, 6), (3, 5), (2, 4),|
  !gapprompt@>| !gapinput@       (1, 2)(3, 4)(5, 6)(7, 8)]);|
  true
  !gapprompt@gap>| !gapinput@D := Digraph([[2, 3], [], []]);|
  <immutable digraph with 3 vertices, 2 edges>
  !gapprompt@gap>| !gapinput@HasDigraphGroup(D);|
  false
  !gapprompt@gap>| !gapinput@HasAutomorphismGroup(D);|
  false
  !gapprompt@gap>| !gapinput@DigraphGroup(D);|
  Group([ (2,3) ])
  !gapprompt@gap>| !gapinput@HasAutomorphismGroup(D);|
  true
  !gapprompt@gap>| !gapinput@group := DihedralGroup(8);|
  <pc group of size 8 with 3 generators>
  !gapprompt@gap>| !gapinput@D := CayleyDigraph(group);|
  <immutable digraph with 8 vertices, 24 edges>
  !gapprompt@gap>| !gapinput@HasDigraphGroup(D);|
  true
  !gapprompt@gap>| !gapinput@GeneratorsOfGroup(DigraphGroup(D));|
  [ (1,2)(3,5)(4,6)(7,8), (1,7,4,3)(2,5,6,8), (1,4)(2,6)(3,7)(5,8) ]
\end{Verbatim}
 }

 

\subsection{\textcolor{Chapter }{DigraphOrbits}}
\logpage{[ 7, 2, 11 ]}\nobreak
\hyperdef{L}{X8096C5287E459279}{}
{\noindent\textcolor{FuncColor}{$\triangleright$\enspace\texttt{DigraphOrbits({\mdseries\slshape digraph})\index{DigraphOrbits@\texttt{DigraphOrbits}}
\label{DigraphOrbits}
}\hfill{\scriptsize (attribute)}}\\
\textbf{\indent Returns:\ }
 An immutable list of lists of integers. 



 \texttt{DigraphOrbits} returns the orbits of the action of the \texttt{DigraphGroup} (\ref{DigraphGroup}) on the set of vertices of \mbox{\texttt{\mdseries\slshape digraph}}. 
\begin{Verbatim}[commandchars=!@|,fontsize=\small,frame=single,label=Example]
  !gapprompt@gap>| !gapinput@G := Group([(2, 3)(7, 8, 9), (1, 2, 3)(4, 5, 6)(8, 9)]);;|
  !gapprompt@gap>| !gapinput@D := EdgeOrbitsDigraph(G, [1, 2]);|
  <immutable digraph with 9 vertices, 6 edges>
  !gapprompt@gap>| !gapinput@DigraphOrbits(D);|
  [ [ 1, 2, 3 ], [ 4, 5, 6 ], [ 7, 8, 9 ] ]
  !gapprompt@gap>| !gapinput@D := DigraphMutableCopy(D);|
  <mutable digraph with 9 vertices, 6 edges>
  !gapprompt@gap>| !gapinput@DigraphOrbits(D);|
  [ [ 1, 2, 3 ], [ 4, 5, 6, 7, 8, 9 ] ]
\end{Verbatim}
 }

 

\subsection{\textcolor{Chapter }{DigraphOrbitReps}}
\logpage{[ 7, 2, 12 ]}\nobreak
\hyperdef{L}{X8386028782F2D3FF}{}
{\noindent\textcolor{FuncColor}{$\triangleright$\enspace\texttt{DigraphOrbitReps({\mdseries\slshape digraph})\index{DigraphOrbitReps@\texttt{DigraphOrbitReps}}
\label{DigraphOrbitReps}
}\hfill{\scriptsize (attribute)}}\\
\textbf{\indent Returns:\ }
 An immutable list of integers. 



 \texttt{DigraphOrbitReps} returns a list of orbit representatives of the action of the \texttt{DigraphGroup} (\ref{DigraphGroup}) on the set of vertices of \mbox{\texttt{\mdseries\slshape digraph}}. 
\begin{Verbatim}[commandchars=!|B,fontsize=\small,frame=single,label=Example]
  !gapprompt|gap>B !gapinput|D := CayleyDigraph(AlternatingGroup(4));B
  <immutable digraph with 12 vertices, 24 edges>
  !gapprompt|gap>B !gapinput|DigraphOrbitReps(D);B
  [ 1 ]
  !gapprompt|gap>B !gapinput|D := DigraphMutableCopy(D);B
  <mutable digraph with 12 vertices, 24 edges>
  !gapprompt|gap>B !gapinput|DigraphOrbitReps(D);B
  [ 1 ]
  !gapprompt|gap>B !gapinput|D := DigraphFromDigraph6String("&IGO??S?`?_@?a?CK?O");B
  <immutable digraph with 10 vertices, 14 edges>
  !gapprompt|gap>B !gapinput|DigraphOrbitReps(D);B
  [ 1, 2, 3, 4, 5, 6, 7, 8, 9, 10 ]
  !gapprompt|gap>B !gapinput|DigraphOrbitReps(DigraphMutableCopy(D));B
  [ 1, 2, 3, 4, 5, 6, 7, 8, 9, 10 ]
\end{Verbatim}
 }

 

\subsection{\textcolor{Chapter }{DigraphSchreierVector}}
\logpage{[ 7, 2, 13 ]}\nobreak
\hyperdef{L}{X8657604E87A25E5F}{}
{\noindent\textcolor{FuncColor}{$\triangleright$\enspace\texttt{DigraphSchreierVector({\mdseries\slshape digraph})\index{DigraphSchreierVector@\texttt{DigraphSchreierVector}}
\label{DigraphSchreierVector}
}\hfill{\scriptsize (attribute)}}\\
\textbf{\indent Returns:\ }
 An immutable list of integers. 



 \texttt{DigraphSchreierVector} returns the so\texttt{\symbol{45}}called \emph{Schreier vector} of the action of the \texttt{DigraphGroup} (\ref{DigraphGroup}) on the set of vertices of \mbox{\texttt{\mdseries\slshape digraph}}. The Schreier vector is a list \texttt{sch} of integers with length \texttt{DigraphNrVertices(\mbox{\texttt{\mdseries\slshape digraph}})} where: 
\begin{description}
\item[{\texttt{sch[i] {\textless} 0:}}]  implies that \texttt{i} is an orbit representative and \texttt{DigraphOrbitReps(\mbox{\texttt{\mdseries\slshape digraph}})[\texttt{\symbol{45}}sch[i]] = i}. 
\item[{\texttt{sch[i] {\textgreater} 0:}}]  implies that \texttt{i / gens[sch[i]]} is one step closer to the root (or representative) of the tree, where \texttt{gens} is the generators of \texttt{DigraphGroup(\mbox{\texttt{\mdseries\slshape digraph}})}. 
\end{description}
 
\begin{Verbatim}[commandchars=!@|,fontsize=\small,frame=single,label=Example]
  !gapprompt@gap>| !gapinput@n := 4;;|
  !gapprompt@gap>| !gapinput@adj := function(x, y)|
  !gapprompt@>| !gapinput@     return (((x - y) mod n) = 1) or (((x - y) mod n) = n - 1);|
  !gapprompt@>| !gapinput@   end;;|
  !gapprompt@gap>| !gapinput@group := CyclicGroup(IsPermGroup, n);|
  Group([ (1,2,3,4) ])
  !gapprompt@gap>| !gapinput@D := Digraph(IsMutableDigraph, group, [1 .. n], \^, adj);|
  <mutable digraph with 4 vertices, 8 edges>
  !gapprompt@gap>| !gapinput@sch := DigraphSchreierVector(D);|
  [ -1, 2, 2, 1 ]
  !gapprompt@gap>| !gapinput@D := CayleyDigraph(AlternatingGroup(4));|
  <immutable digraph with 12 vertices, 24 edges>
  !gapprompt@gap>| !gapinput@sch := DigraphSchreierVector(D);|
  [ -1, 2, 2, 1, 1, 2, 1, 2, 2, 1, 1, 1 ]
  !gapprompt@gap>| !gapinput@DigraphOrbitReps(D);|
  [ 1 ]
  !gapprompt@gap>| !gapinput@gens := GeneratorsOfGroup(DigraphGroup(D));|
  [ (1,7,5)(2,10,9)(3,4,11)(6,8,12), (1,3,2)(4,5,6)(7,9,8)(10,11,12) ]
  !gapprompt@gap>| !gapinput@10 / gens[sch[10]];|
  2
  !gapprompt@gap>| !gapinput@7 / gens[sch[7]];|
  1
  !gapprompt@gap>| !gapinput@5 / gens[sch[5]];|
  7
\end{Verbatim}
 }

 

\subsection{\textcolor{Chapter }{DigraphStabilizer}}
\logpage{[ 7, 2, 14 ]}\nobreak
\hyperdef{L}{X78913684795FB256}{}
{\noindent\textcolor{FuncColor}{$\triangleright$\enspace\texttt{DigraphStabilizer({\mdseries\slshape digraph, v})\index{DigraphStabilizer@\texttt{DigraphStabilizer}}
\label{DigraphStabilizer}
}\hfill{\scriptsize (operation)}}\\
\textbf{\indent Returns:\ }
 A permutation group. 



 \texttt{DigraphStabilizer} returns the stabilizer of the vertex \mbox{\texttt{\mdseries\slshape v}} under of the action of the \texttt{DigraphGroup} (\ref{DigraphGroup}) on the set of vertices of \mbox{\texttt{\mdseries\slshape digraph}}. 
\begin{Verbatim}[commandchars=!@|,fontsize=\small,frame=single,label=Example]
  !gapprompt@gap>| !gapinput@D := DigraphFromDigraph6String("&GYHPQgWTIIPW");|
  <immutable digraph with 8 vertices, 24 edges>
  !gapprompt@gap>| !gapinput@DigraphStabilizer(D, 8);|
  Group(())
  !gapprompt@gap>| !gapinput@DigraphStabilizer(D, 2);|
  Group(())
  !gapprompt@gap>| !gapinput@D := DigraphMutableCopy(D);|
  <mutable digraph with 8 vertices, 24 edges>
  !gapprompt@gap>| !gapinput@DigraphStabilizer(D, 8);|
  Group(())
  !gapprompt@gap>| !gapinput@DigraphStabilizer(D, 2);|
  Group(())
\end{Verbatim}
 }

 

\subsection{\textcolor{Chapter }{IsIsomorphicDigraph (for digraphs)}}
\logpage{[ 7, 2, 15 ]}\nobreak
\hyperdef{L}{X7B4F24B283C9EE28}{}
{\noindent\textcolor{FuncColor}{$\triangleright$\enspace\texttt{IsIsomorphicDigraph({\mdseries\slshape digraph1, digraph2})\index{IsIsomorphicDigraph@\texttt{IsIsomorphicDigraph}!for digraphs}
\label{IsIsomorphicDigraph:for digraphs}
}\hfill{\scriptsize (operation)}}\\
\textbf{\indent Returns:\ }
\texttt{true} or \texttt{false}.



 This operation returns \texttt{true} if there exists an isomorphism from the digraph \mbox{\texttt{\mdseries\slshape digraph1}} to the digraph \mbox{\texttt{\mdseries\slshape digraph2}}. See \texttt{IsomorphismDigraphs} (\ref{IsomorphismDigraphs:for digraphs}) for more information about isomorphisms of digraphs. 

 By default, an isomorphism is found using the canonical labellings of the
digraphs obtained from \href{http://www.tcs.tkk.fi/Software/bliss/} {bliss}  by Tommi Junttila and Petteri Kaski. If \textsf{NautyTracesInterface} is available, then \href{https://pallini.di.uniroma1.it} {nauty}  by Brendan Mckay and Adolfo Piperno can be used instead; see \texttt{DigraphsUseBliss} (\ref{DigraphsUseBliss}), and \texttt{DigraphsUseNauty} (\ref{DigraphsUseNauty}). 
\begin{Verbatim}[commandchars=!@|,fontsize=\small,frame=single,label=Example]
  !gapprompt@gap>| !gapinput@digraph1 := CycleDigraph(4);|
  <immutable cycle digraph with 4 vertices>
  !gapprompt@gap>| !gapinput@digraph2 := CycleDigraph(5);|
  <immutable cycle digraph with 5 vertices>
  !gapprompt@gap>| !gapinput@IsIsomorphicDigraph(digraph1, digraph2);|
  false
  !gapprompt@gap>| !gapinput@digraph2 := DigraphReverse(digraph1);|
  <immutable digraph with 4 vertices, 4 edges>
  !gapprompt@gap>| !gapinput@IsIsomorphicDigraph(digraph1, digraph2);|
  true
  !gapprompt@gap>| !gapinput@digraph1 := Digraph([[3], [], []]);|
  <immutable digraph with 3 vertices, 1 edge>
  !gapprompt@gap>| !gapinput@digraph2 := Digraph([[], [], [2]]);|
  <immutable digraph with 3 vertices, 1 edge>
  !gapprompt@gap>| !gapinput@IsIsomorphicDigraph(digraph1, digraph2);|
  true
\end{Verbatim}
 }

 

\subsection{\textcolor{Chapter }{IsIsomorphicDigraph (for digraphs and homogeneous lists)}}
\logpage{[ 7, 2, 16 ]}\nobreak
\hyperdef{L}{X7D4A1FA8868DA930}{}
{\noindent\textcolor{FuncColor}{$\triangleright$\enspace\texttt{IsIsomorphicDigraph({\mdseries\slshape digraph1, digraph2, colours1, colours2})\index{IsIsomorphicDigraph@\texttt{IsIsomorphicDigraph}!for digraphs and homogeneous lists}
\label{IsIsomorphicDigraph:for digraphs and homogeneous lists}
}\hfill{\scriptsize (operation)}}\\
\textbf{\indent Returns:\ }
\texttt{true} or \texttt{false}.



 This operation tests for isomorphism of coloured digraphs. A coloured digraph
can be specified by its underlying digraph \mbox{\texttt{\mdseries\slshape digraph1}} and its colouring \mbox{\texttt{\mdseries\slshape colours1}}. Let \texttt{n} be the number of vertices of \mbox{\texttt{\mdseries\slshape digraph1}}. The colouring \mbox{\texttt{\mdseries\slshape colours1}} may have one of the following two forms: 
\begin{itemize}
\item  a list of \texttt{n} integers, where \mbox{\texttt{\mdseries\slshape colours}}\texttt{[i]} is the colour of vertex \texttt{i}, using the colours \texttt{[1 .. m]} for some \texttt{m {\textless}= n}; or 
\item  a list of non\texttt{\symbol{45}}empty disjoint lists whose union is \texttt{DigraphVertices(\mbox{\texttt{\mdseries\slshape digraph}})}, such that \mbox{\texttt{\mdseries\slshape colours}}\texttt{[i]} is the list of all vertices with colour \texttt{i}. 
\end{itemize}
 If \mbox{\texttt{\mdseries\slshape digraph1}} and \mbox{\texttt{\mdseries\slshape digraph2}} are digraphs without multiple edges, and \mbox{\texttt{\mdseries\slshape colours1}} and \mbox{\texttt{\mdseries\slshape colours2}} are colourings of \mbox{\texttt{\mdseries\slshape digraph1}} and \mbox{\texttt{\mdseries\slshape digraph2}}, respectively, then this operation returns \texttt{true} if there exists an isomorphism between these two coloured digraphs. See \texttt{IsomorphismDigraphs} (\ref{IsomorphismDigraphs:for digraphs and homogeneous lists}) for more information about isomorphisms of coloured digraphs. 

 By default, an isomorphism is found using the canonical labellings of the
digraphs obtained from \href{http://www.tcs.tkk.fi/Software/bliss/} {bliss}  by Tommi Junttila and Petteri Kaski. If \textsf{NautyTracesInterface} is available, then \href{https://pallini.di.uniroma1.it} {nauty}  by Brendan Mckay and Adolfo Piperno can be used instead; see \texttt{DigraphsUseBliss} (\ref{DigraphsUseBliss}), and \texttt{DigraphsUseNauty} (\ref{DigraphsUseNauty}). 
\begin{Verbatim}[commandchars=!@|,fontsize=\small,frame=single,label=Example]
  !gapprompt@gap>| !gapinput@digraph1 := ChainDigraph(4);|
  <immutable chain digraph with 4 vertices>
  !gapprompt@gap>| !gapinput@digraph2 := ChainDigraph(3);|
  <immutable chain digraph with 3 vertices>
  !gapprompt@gap>| !gapinput@IsIsomorphicDigraph(digraph1, digraph2,|
  !gapprompt@>| !gapinput@ [[1, 4], [2, 3]], [[1, 2], [3]]);|
  false
  !gapprompt@gap>| !gapinput@digraph2 := DigraphReverse(digraph1);|
  <immutable digraph with 4 vertices, 3 edges>
  !gapprompt@gap>| !gapinput@IsIsomorphicDigraph(digraph1, digraph2,|
  !gapprompt@>| !gapinput@ [1, 1, 1, 1], [1, 1, 1, 1]);|
  true
  !gapprompt@gap>| !gapinput@IsIsomorphicDigraph(digraph1, digraph2,|
  !gapprompt@>| !gapinput@ [1, 2, 2, 1], [1, 2, 2, 1]);|
  true
  !gapprompt@gap>| !gapinput@IsIsomorphicDigraph(digraph1, digraph2,|
  !gapprompt@>| !gapinput@ [1, 1, 2, 2], [1, 1, 2, 2]);|
  false
\end{Verbatim}
 }

 

\subsection{\textcolor{Chapter }{IsomorphismDigraphs (for digraphs)}}
\logpage{[ 7, 2, 17 ]}\nobreak
\hyperdef{L}{X8219FE6C839D9457}{}
{\noindent\textcolor{FuncColor}{$\triangleright$\enspace\texttt{IsomorphismDigraphs({\mdseries\slshape digraph1, digraph2})\index{IsomorphismDigraphs@\texttt{IsomorphismDigraphs}!for digraphs}
\label{IsomorphismDigraphs:for digraphs}
}\hfill{\scriptsize (operation)}}\\
\textbf{\indent Returns:\ }
 A permutation, or a pair of permutations, or \texttt{fail}.



 This operation returns an isomorphism between the digraphs \mbox{\texttt{\mdseries\slshape digraph1}} and \mbox{\texttt{\mdseries\slshape digraph2}} if one exists, else this operation returns \texttt{fail}. 

 An \emph{isomorphism} from a digraph \mbox{\texttt{\mdseries\slshape digraph1}} to a digraph \mbox{\texttt{\mdseries\slshape digraph2}} is a bijection \texttt{p} from the vertices of \mbox{\texttt{\mdseries\slshape digraph1}} to the vertices of \mbox{\texttt{\mdseries\slshape digraph2}} with the following property: for all vertices \texttt{i} and \texttt{j} of \mbox{\texttt{\mdseries\slshape digraph1}}, \texttt{[i, j]} is an edge of \mbox{\texttt{\mdseries\slshape digraph1}} if and only if \texttt{[i \texttt{\symbol{94}} p, j \texttt{\symbol{94}} p]} is an edge of \mbox{\texttt{\mdseries\slshape digraph2}}. 

 If there exists such an isomorphism, then this operation returns one. The form
of this isomorphism is a permutation \texttt{p} of the vertices of \mbox{\texttt{\mdseries\slshape digraph1}} such that 

 \texttt{OnDigraphs(\mbox{\texttt{\mdseries\slshape digraph1}}, p) = digraph2}. By default, an isomorphism is found using the canonical labellings of the
digraphs obtained from \href{http://www.tcs.tkk.fi/Software/bliss/} {bliss}  by Tommi Junttila and Petteri Kaski. If \textsf{NautyTracesInterface} is available, then \href{https://pallini.di.uniroma1.it} {nauty}  by Brendan Mckay and Adolfo Piperno can be used instead; see \texttt{DigraphsUseBliss} (\ref{DigraphsUseBliss}), and \texttt{DigraphsUseNauty} (\ref{DigraphsUseNauty}). 
\begin{Verbatim}[commandchars=!@|,fontsize=\small,frame=single,label=Example]
  !gapprompt@gap>| !gapinput@digraph1 := CycleDigraph(4);|
  <immutable cycle digraph with 4 vertices>
  !gapprompt@gap>| !gapinput@digraph2 := CycleDigraph(5);|
  <immutable cycle digraph with 5 vertices>
  !gapprompt@gap>| !gapinput@IsomorphismDigraphs(digraph1, digraph2);|
  fail
  !gapprompt@gap>| !gapinput@digraph1 := CompleteBipartiteDigraph(10, 5);|
  <immutable complete bipartite digraph with bicomponent sizes 10 and 5>
  !gapprompt@gap>| !gapinput@digraph2 := CompleteBipartiteDigraph(5, 10);|
  <immutable complete bipartite digraph with bicomponent sizes 5 and 10>
  !gapprompt@gap>| !gapinput@p := IsomorphismDigraphs(digraph1, digraph2);|
  (1,6,11)(2,7,12)(3,8,13)(4,9,14)(5,10,15)
  !gapprompt@gap>| !gapinput@OnDigraphs(digraph1, p) = digraph2;|
  true
\end{Verbatim}
 }

 

\subsection{\textcolor{Chapter }{IsomorphismDigraphs (for digraphs and homogeneous lists)}}
\logpage{[ 7, 2, 18 ]}\nobreak
\hyperdef{L}{X7ED93C0F86D9D34F}{}
{\noindent\textcolor{FuncColor}{$\triangleright$\enspace\texttt{IsomorphismDigraphs({\mdseries\slshape digraph1, digraph2, colours1, colours2})\index{IsomorphismDigraphs@\texttt{IsomorphismDigraphs}!for digraphs and homogeneous lists}
\label{IsomorphismDigraphs:for digraphs and homogeneous lists}
}\hfill{\scriptsize (operation)}}\\
\textbf{\indent Returns:\ }
 A permutation, or \texttt{fail}.



 This operation searches for an isomorphism between coloured digraphs. A
coloured digraph can be specified by its underlying digraph \mbox{\texttt{\mdseries\slshape digraph1}} and its colouring \mbox{\texttt{\mdseries\slshape colours1}}. Let \texttt{n} be the number of vertices of \mbox{\texttt{\mdseries\slshape digraph1}}. The colouring \mbox{\texttt{\mdseries\slshape colours1}} may have one of the following two forms: 
\begin{itemize}
\item  a list of \texttt{n} integers, where \mbox{\texttt{\mdseries\slshape colours}}\texttt{[i]} is the colour of vertex \texttt{i}, using the colours \texttt{[1 .. m]} for some \texttt{m {\textless}= n}; or 
\item  a list of non\texttt{\symbol{45}}empty disjoint lists whose union is \texttt{DigraphVertices(\mbox{\texttt{\mdseries\slshape digraph}})}, such that \mbox{\texttt{\mdseries\slshape colours}}\texttt{[i]} is the list of all vertices with colour \texttt{i}. 
\end{itemize}
 An \emph{isomorphism} between coloured digraphs is an isomorphism between the underlying digraphs
that preserves the colourings. See \texttt{IsomorphismDigraphs} (\ref{IsomorphismDigraphs:for digraphs}) for more information about isomorphisms of digraphs. More precisely, let \texttt{f} be an isomorphism of digraphs from the digraph \mbox{\texttt{\mdseries\slshape digraph1}} (with colouring \mbox{\texttt{\mdseries\slshape colours1}}) to the digraph \mbox{\texttt{\mdseries\slshape digraph2}} (with colouring \mbox{\texttt{\mdseries\slshape colours2}}), and let \texttt{p} be the permutation of the vertices of \mbox{\texttt{\mdseries\slshape digraph1}} that corresponds to \texttt{f}. Then \texttt{f} preserves the colourings of \mbox{\texttt{\mdseries\slshape digraph1}} and \mbox{\texttt{\mdseries\slshape digraph2}} {\textendash} and hence is an isomorphism of coloured digraphs {\textendash}
if \texttt{\mbox{\texttt{\mdseries\slshape colours1}}[i] = \mbox{\texttt{\mdseries\slshape colours2}}[i \texttt{\symbol{94}} p]} for all vertices \texttt{i} in \mbox{\texttt{\mdseries\slshape digraph1}}. 

 This operation returns such an isomorphism if one exists, else it returns \texttt{fail}. 

 By default, an isomorphism is found using the canonical labellings of the
digraphs obtained from \href{http://www.tcs.tkk.fi/Software/bliss/} {bliss}  by Tommi Junttila and Petteri Kaski. If \textsf{NautyTracesInterface} is available, then \href{https://pallini.di.uniroma1.it} {nauty}  by Brendan Mckay and Adolfo Piperno can be used instead; see \texttt{DigraphsUseBliss} (\ref{DigraphsUseBliss}), and \texttt{DigraphsUseNauty} (\ref{DigraphsUseNauty}). 
\begin{Verbatim}[commandchars=!@|,fontsize=\small,frame=single,label=Example]
  !gapprompt@gap>| !gapinput@digraph1 := ChainDigraph(4);|
  <immutable chain digraph with 4 vertices>
  !gapprompt@gap>| !gapinput@digraph2 := ChainDigraph(3);|
  <immutable chain digraph with 3 vertices>
  !gapprompt@gap>| !gapinput@IsomorphismDigraphs(digraph1, digraph2,|
  !gapprompt@>| !gapinput@ [[1, 4], [2, 3]], [[1, 2], [3]]);|
  fail
  !gapprompt@gap>| !gapinput@digraph2 := DigraphReverse(digraph1);|
  <immutable digraph with 4 vertices, 3 edges>
  !gapprompt@gap>| !gapinput@colours1 := [1, 1, 1, 1];;|
  !gapprompt@gap>| !gapinput@colours2 := [1, 1, 1, 1];;|
  !gapprompt@gap>| !gapinput@p := IsomorphismDigraphs(digraph1, digraph2, colours1, colours2);|
  (1,4)(2,3)
  !gapprompt@gap>| !gapinput@OnDigraphs(digraph1, p) = digraph2;|
  true
  !gapprompt@gap>| !gapinput@List(DigraphVertices(digraph1), i -> colours1[i ^ p]) = colours2;|
  true
  !gapprompt@gap>| !gapinput@colours1 := [1, 1, 2, 2];;|
  !gapprompt@gap>| !gapinput@colours2 := [2, 2, 1, 1];;|
  !gapprompt@gap>| !gapinput@p := IsomorphismDigraphs(digraph1, digraph2, colours1, colours2);|
  (1,4)(2,3)
  !gapprompt@gap>| !gapinput@OnDigraphs(digraph1, p) = digraph2;|
  true
  !gapprompt@gap>| !gapinput@List(DigraphVertices(digraph1), i -> colours1[i ^ p]) = colours2;|
  true
  !gapprompt@gap>| !gapinput@IsomorphismDigraphs(digraph1, digraph2,|
  !gapprompt@>| !gapinput@ [1, 1, 2, 2], [1, 1, 2, 2]);|
  fail
\end{Verbatim}
 }

 

\subsection{\textcolor{Chapter }{RepresentativeOutNeighbours}}
\logpage{[ 7, 2, 19 ]}\nobreak
\hyperdef{L}{X7E1C806D81DFE15E}{}
{\noindent\textcolor{FuncColor}{$\triangleright$\enspace\texttt{RepresentativeOutNeighbours({\mdseries\slshape digraph})\index{RepresentativeOutNeighbours@\texttt{RepresentativeOutNeighbours}}
\label{RepresentativeOutNeighbours}
}\hfill{\scriptsize (attribute)}}\\
\textbf{\indent Returns:\ }
An immutable list of lists.



 This function returns the list \texttt{out} of \emph{out\texttt{\symbol{45}}neighbours} of each representative of the orbits of the action of \texttt{DigraphGroup} (\ref{DigraphGroup}) on the vertex set of the digraph \mbox{\texttt{\mdseries\slshape digraph}}. 

 More specifically, if \texttt{reps} is the list of orbit representatives, then a vertex \texttt{j} appears in \texttt{out[i]} each time there exists an edge with source \texttt{reps[i]} and range \texttt{j} in \mbox{\texttt{\mdseries\slshape digraph}}. 

 If \texttt{DigraphGroup} (\ref{DigraphGroup}) is trivial, then \texttt{OutNeighbours} (\ref{OutNeighbours}) is returned. 
\begin{Verbatim}[commandchars=!@|,fontsize=\small,frame=single,label=Example]
  !gapprompt@gap>| !gapinput@D := Digraph([|
  !gapprompt@>| !gapinput@ [2, 1, 3, 4, 5], [3, 5], [2], [1, 2, 3, 5], [1, 2, 3, 4]]);|
  <immutable digraph with 5 vertices, 16 edges>
  !gapprompt@gap>| !gapinput@DigraphGroup(D);|
  Group(())
  !gapprompt@gap>| !gapinput@RepresentativeOutNeighbours(D);|
  [ [ 2, 1, 3, 4, 5 ], [ 3, 5 ], [ 2 ], [ 1, 2, 3, 5 ], [ 1, 2, 3, 4 ] ]
  !gapprompt@gap>| !gapinput@D := Digraph(IsMutableDigraph, [|
  !gapprompt@>| !gapinput@ [2, 1, 3, 4, 5], [3, 5], [2], [1, 2, 3, 5], [1, 2, 3, 4]]);|
  <mutable digraph with 5 vertices, 16 edges>
  !gapprompt@gap>| !gapinput@DigraphGroup(D);|
  Group(())
  !gapprompt@gap>| !gapinput@RepresentativeOutNeighbours(D);|
  [ [ 2, 1, 3, 4, 5 ], [ 3, 5 ], [ 2 ], [ 1, 2, 3, 5 ], [ 1, 2, 3, 4 ] ]
  !gapprompt@gap>| !gapinput@D := DigraphFromDigraph6String("&GYHPQgWTIIPW");|
  <immutable digraph with 8 vertices, 24 edges>
  !gapprompt@gap>| !gapinput@G := DigraphGroup(D);;|
  !gapprompt@gap>| !gapinput@GeneratorsOfGroup(G);|
  [ (1,2)(3,4)(5,6)(7,8), (1,3,2,4)(5,7,6,8), (1,5)(2,6)(3,8)(4,7) ]
  !gapprompt@gap>| !gapinput@Set(RepresentativeOutNeighbours(D), Set);|
  [ [ 2, 3, 5 ] ]
\end{Verbatim}
 }

 

\subsection{\textcolor{Chapter }{IsDigraphIsomorphism (for digraphs and transformation or permutation)}}
\logpage{[ 7, 2, 20 ]}\nobreak
\hyperdef{L}{X80183D4A7C51365A}{}
{\noindent\textcolor{FuncColor}{$\triangleright$\enspace\texttt{IsDigraphIsomorphism({\mdseries\slshape src, ran, x})\index{IsDigraphIsomorphism@\texttt{IsDigraphIsomorphism}!for digraphs and transformation or permutation}
\label{IsDigraphIsomorphism:for digraphs and transformation or permutation}
}\hfill{\scriptsize (operation)}}\\
\noindent\textcolor{FuncColor}{$\triangleright$\enspace\texttt{IsDigraphIsomorphism({\mdseries\slshape src, ran, x, col1, col2})\index{IsDigraphIsomorphism@\texttt{IsDigraphIsomorphism}}
\label{IsDigraphIsomorphism}
}\hfill{\scriptsize (operation)}}\\
\noindent\textcolor{FuncColor}{$\triangleright$\enspace\texttt{IsDigraphAutomorphism({\mdseries\slshape digraph, x})\index{IsDigraphAutomorphism@\texttt{IsDigraphAutomorphism}!for a digraph and a transformation or permutation}
\label{IsDigraphAutomorphism:for a digraph and a transformation or permutation}
}\hfill{\scriptsize (operation)}}\\
\noindent\textcolor{FuncColor}{$\triangleright$\enspace\texttt{IsDigraphAutomorphism({\mdseries\slshape digraph, x, col})\index{IsDigraphAutomorphism@\texttt{IsDigraphAutomorphism}}
\label{IsDigraphAutomorphism}
}\hfill{\scriptsize (operation)}}\\
\textbf{\indent Returns:\ }
\texttt{true} or \texttt{false}.



 \texttt{IsDigraphIsomorphism} returns \texttt{true} if the permutation or transformation \mbox{\texttt{\mdseries\slshape x}} is an isomorphism from the digraph \mbox{\texttt{\mdseries\slshape src}} to the digraph \mbox{\texttt{\mdseries\slshape ran}}. 

 \texttt{IsDigraphAutomorphism} returns \texttt{true} if the permutation or transformation \mbox{\texttt{\mdseries\slshape x}} is an automorphism of the digraph \mbox{\texttt{\mdseries\slshape digraph}}. 

 A permutation or transformation \mbox{\texttt{\mdseries\slshape x}} is an \emph{isomorphism} from a digraph \mbox{\texttt{\mdseries\slshape src}} to a digraph \mbox{\texttt{\mdseries\slshape ran}} if the following hold: 
\begin{itemize}
\item  \mbox{\texttt{\mdseries\slshape x}} is a bijection from the vertices of \mbox{\texttt{\mdseries\slshape src}} to those of \mbox{\texttt{\mdseries\slshape ran}}; 
\item  \texttt{[u \texttt{\symbol{94}} \mbox{\texttt{\mdseries\slshape x}}, v \texttt{\symbol{94}} \mbox{\texttt{\mdseries\slshape x}}]} is an edge of \mbox{\texttt{\mdseries\slshape ran}} if and only if \texttt{[u, v]} is an edge of \mbox{\texttt{\mdseries\slshape src}}; and 
\item  \mbox{\texttt{\mdseries\slshape x}} fixes every \texttt{i} which is not a vertex of \mbox{\texttt{\mdseries\slshape src}}. 
\end{itemize}
 See also \texttt{AutomorphismGroup} (\ref{AutomorphismGroup:for a digraph}).

 If \mbox{\texttt{\mdseries\slshape col1}} and \mbox{\texttt{\mdseries\slshape col2}}, or \mbox{\texttt{\mdseries\slshape col}}, are given, then they must represent vertex colourings; see \texttt{AutomorphismGroup} (\ref{AutomorphismGroup:for a digraph and a homogeneous list}) for details of the permissible values for these arguments. The homomorphism
must then also have the property: 
\begin{itemize}
\item  \texttt{col1[i] = col2[i \texttt{\symbol{94}} x]} for all vertices \texttt{i} of \mbox{\texttt{\mdseries\slshape src}}, for \texttt{IsDigraphIsomorphism}. 
\item  \texttt{col[i] = col[i \texttt{\symbol{94}} x]} for all vertices \texttt{i} of \mbox{\texttt{\mdseries\slshape digraph}}, for \texttt{IsDigraphAutomorphism}. 
\end{itemize}
 For some digraphs, it can be faster to use \texttt{IsDigraphAutomorphism} than to test membership in the automorphism group of \mbox{\texttt{\mdseries\slshape digraph}}. 
\begin{Verbatim}[commandchars=!@|,fontsize=\small,frame=single,label=Example]
  !gapprompt@gap>| !gapinput@src := Digraph([[1], [1, 2], [1, 3]]);|
  <immutable digraph with 3 vertices, 5 edges>
  !gapprompt@gap>| !gapinput@IsDigraphAutomorphism(src, (1, 2, 3));|
  false
  !gapprompt@gap>| !gapinput@IsDigraphAutomorphism(src, (2, 3));|
  true
  !gapprompt@gap>| !gapinput@IsDigraphAutomorphism(src, (2, 3), [2, 1, 1]);|
  true
  !gapprompt@gap>| !gapinput@IsDigraphAutomorphism(src, (2, 3), [2, 2, 1]);|
  false
  !gapprompt@gap>| !gapinput@IsDigraphAutomorphism(src, (2, 3)(4, 5));|
  true
  !gapprompt@gap>| !gapinput@IsDigraphAutomorphism(src, (1, 4));|
  false
  !gapprompt@gap>| !gapinput@IsDigraphAutomorphism(src, ());|
  true
  !gapprompt@gap>| !gapinput@ran := Digraph([[2, 1], [2], [2, 3]]);|
  <immutable digraph with 3 vertices, 5 edges>
  !gapprompt@gap>| !gapinput@IsDigraphIsomorphism(src, ran, (1, 2));|
  true
  !gapprompt@gap>| !gapinput@IsDigraphIsomorphism(ran, src, (1, 2));|
  true
  !gapprompt@gap>| !gapinput@IsDigraphIsomorphism(ran, src, (1, 2));|
  true
  !gapprompt@gap>| !gapinput@IsDigraphIsomorphism(src, Digraph([[3], [1, 3], [2]]), (1, 2, 3));|
  false
  !gapprompt@gap>| !gapinput@IsDigraphIsomorphism(src, ran, (1, 2), [1, 2, 3], [2, 1, 3]);|
  true
  !gapprompt@gap>| !gapinput@IsDigraphIsomorphism(src, ran, (1, 2), [1, 2, 2], [2, 1, 3]);|
  false
\end{Verbatim}
 }

 

\subsection{\textcolor{Chapter }{IsDigraphColouring}}
\logpage{[ 7, 2, 21 ]}\nobreak
\hyperdef{L}{X7F16B8B3825A627A}{}
{\noindent\textcolor{FuncColor}{$\triangleright$\enspace\texttt{IsDigraphColouring({\mdseries\slshape digraph, list})\index{IsDigraphColouring@\texttt{IsDigraphColouring}}
\label{IsDigraphColouring}
}\hfill{\scriptsize (operation)}}\\
\noindent\textcolor{FuncColor}{$\triangleright$\enspace\texttt{IsDigraphColouring({\mdseries\slshape digraph, t})\index{IsDigraphColouring@\texttt{IsDigraphColouring}!for a transformation}
\label{IsDigraphColouring:for a transformation}
}\hfill{\scriptsize (operation)}}\\
\textbf{\indent Returns:\ }
 \texttt{true} or \texttt{false}. 



 The operation \texttt{IsDigraphColouring} verifies whether or not the list \mbox{\texttt{\mdseries\slshape list}} describes a proper colouring of the digraph \mbox{\texttt{\mdseries\slshape digraph}}. 

 A list \mbox{\texttt{\mdseries\slshape list}} describes a \emph{proper colouring} of a digraph \mbox{\texttt{\mdseries\slshape digraph}} if \mbox{\texttt{\mdseries\slshape list}} consists of positive integers, the length of \mbox{\texttt{\mdseries\slshape list}} equals the number of vertices in \mbox{\texttt{\mdseries\slshape digraph}}, and for any vertices \texttt{u, v} of \mbox{\texttt{\mdseries\slshape digraph}} if \texttt{u} and \texttt{v} are adjacent, then \texttt{\mbox{\texttt{\mdseries\slshape list}}[u] {\textgreater}{\textless} \mbox{\texttt{\mdseries\slshape list}}[v]}. 

 A transformation \mbox{\texttt{\mdseries\slshape t}} describes a proper colouring of a digraph \mbox{\texttt{\mdseries\slshape digraph}}, if \texttt{ImageListOfTransformation(\mbox{\texttt{\mdseries\slshape t}}, DigraphNrVertices(\mbox{\texttt{\mdseries\slshape digraph}}))} is a proper colouring of \mbox{\texttt{\mdseries\slshape digraph}}. 

 See also \texttt{IsDigraphHomomorphism} (\ref{IsDigraphHomomorphism}). 
\begin{Verbatim}[commandchars=!@|,fontsize=\small,frame=single,label=Example]
  !gapprompt@gap>| !gapinput@D := JohnsonDigraph(5, 3);|
  <immutable symmetric digraph with 10 vertices, 60 edges>
  !gapprompt@gap>| !gapinput@IsDigraphColouring(D, [1, 2, 3, 3, 2, 1, 4, 5, 6, 7]);|
  true
  !gapprompt@gap>| !gapinput@IsDigraphColouring(D, [1, 2, 3, 3, 2, 1, 2, 5, 6, 7]);|
  false
  !gapprompt@gap>| !gapinput@IsDigraphColouring(D, [1, 2, 3, 3, 2, 1, 2, 5, 6, -1]);|
  false
  !gapprompt@gap>| !gapinput@IsDigraphColouring(D, [1, 2, 3]);|
  false
  !gapprompt@gap>| !gapinput@IsDigraphColouring(D, IdentityTransformation);|
  true
\end{Verbatim}
 }

 

\subsection{\textcolor{Chapter }{MaximalCommonSubdigraph}}
\logpage{[ 7, 2, 22 ]}\nobreak
\hyperdef{L}{X7D4378F27B49C9AC}{}
{\noindent\textcolor{FuncColor}{$\triangleright$\enspace\texttt{MaximalCommonSubdigraph({\mdseries\slshape D1, D2})\index{MaximalCommonSubdigraph@\texttt{MaximalCommonSubdigraph}}
\label{MaximalCommonSubdigraph}
}\hfill{\scriptsize (operation)}}\\
\textbf{\indent Returns:\ }
A list containing a digraph and two transformations.



 If \mbox{\texttt{\mdseries\slshape D1}} and \mbox{\texttt{\mdseries\slshape D2}} are digraphs without multiple edges, then \texttt{MaximalCommonSubdigraph} returns a maximal common subgraph \texttt{M} of \mbox{\texttt{\mdseries\slshape D1}} and \mbox{\texttt{\mdseries\slshape D2}} with the maximum number of vertices. So \texttt{M} is a digraph which embeds into both \mbox{\texttt{\mdseries\slshape D1}} and \mbox{\texttt{\mdseries\slshape D2}} and has the largest number of vertices among such digraphs. It returns a list \texttt{[M, t1, t2]} where \texttt{M} is the maximal common subdigraph and \texttt{t1, t2} are transformations embedding \texttt{M} into \mbox{\texttt{\mdseries\slshape D1}} and \mbox{\texttt{\mdseries\slshape D2}} respectively. 
\begin{Verbatim}[commandchars=!@|,fontsize=\small,frame=single,label=Example]
  !gapprompt@gap>| !gapinput@MaximalCommonSubdigraph(PetersenGraph(), CompleteDigraph(10));|
  [ <immutable digraph with 2 vertices, 2 edges>,
    IdentityTransformation, IdentityTransformation ]
  !gapprompt@gap>| !gapinput@MaximalCommonSubdigraph(PetersenGraph(),|
  !gapprompt@>| !gapinput@DigraphSymmetricClosure(CycleDigraph(5)));|
  [ <immutable digraph with 5 vertices, 10 edges>,
    IdentityTransformation, IdentityTransformation ]
  !gapprompt@gap>| !gapinput@MaximalCommonSubdigraph(NullDigraph(0), CompleteDigraph(10));|
  [ <immutable empty digraph with 0 vertices>, IdentityTransformation,
    IdentityTransformation ]
\end{Verbatim}
 }

 

\subsection{\textcolor{Chapter }{MinimalCommonSuperdigraph}}
\logpage{[ 7, 2, 23 ]}\nobreak
\hyperdef{L}{X7B99C75B8021FCDA}{}
{\noindent\textcolor{FuncColor}{$\triangleright$\enspace\texttt{MinimalCommonSuperdigraph({\mdseries\slshape D1, D2})\index{MinimalCommonSuperdigraph@\texttt{MinimalCommonSuperdigraph}}
\label{MinimalCommonSuperdigraph}
}\hfill{\scriptsize (operation)}}\\
\textbf{\indent Returns:\ }
A list containing a digraph and two transformations.



 If \mbox{\texttt{\mdseries\slshape D1}} and \mbox{\texttt{\mdseries\slshape D2}} are digraphs without multiple edges, then \texttt{MinimalCommonSuperdigraph} returns a minimal common superdigraph \texttt{M} of \mbox{\texttt{\mdseries\slshape D1}} and \mbox{\texttt{\mdseries\slshape D2}} with the minimum number of vertices. So \texttt{M} is a digraph into which both \mbox{\texttt{\mdseries\slshape D1}} and \mbox{\texttt{\mdseries\slshape D2}} embed and has the smallest number of vertices among such digraphs. It returns
a list \texttt{[M, t1, t2]} where \texttt{M} is the minimal common superdigraph and \texttt{t1, t2} are transformations embedding \mbox{\texttt{\mdseries\slshape D1}} and \mbox{\texttt{\mdseries\slshape D2}} respectively into \texttt{M}. 
\begin{Verbatim}[commandchars=!@|,fontsize=\small,frame=single,label=Example]
  !gapprompt@gap>| !gapinput@MinimalCommonSuperdigraph(PetersenGraph(), CompleteDigraph(10));|
  [ <immutable digraph with 18 vertices, 118 edges>,
    IdentityTransformation,
    Transformation( [ 1, 2, 11, 12, 13, 14, 15, 16, 17, 18, 11, 12, 13,
        14, 15, 16, 17, 18 ] ) ]
  !gapprompt@gap>| !gapinput@MinimalCommonSuperdigraph(PetersenGraph(),|
  !gapprompt@>| !gapinput@DigraphSymmetricClosure(CycleDigraph(5)));|
  [ <immutable digraph with 10 vertices, 30 edges>,
    IdentityTransformation, IdentityTransformation ]
  !gapprompt@gap>| !gapinput@MinimalCommonSuperdigraph(NullDigraph(0), CompleteDigraph(10));|
  [ <immutable digraph with 10 vertices, 90 edges>,
    IdentityTransformation, IdentityTransformation ]
\end{Verbatim}
 }

 }

 
\section{\textcolor{Chapter }{Homomorphisms of digraphs}}\logpage{[ 7, 3, 0 ]}
\hyperdef{L}{X7E1A02FE8384C03C}{}
{
 The following methods exist to find homomorphisms between digraphs. If an
argument to one of these methods is a digraph with multiple edges, then the
multiplicity of edges will be ignored in order to perform the calculation; the
digraph will be treated as if it has no multiple edges. 

\subsection{\textcolor{Chapter }{HomomorphismDigraphsFinder}}
\logpage{[ 7, 3, 1 ]}\nobreak
\hyperdef{L}{X79ABF0E783FD67C7}{}
{\noindent\textcolor{FuncColor}{$\triangleright$\enspace\texttt{HomomorphismDigraphsFinder({\mdseries\slshape D1, D2, hook, user{\textunderscore}param, max{\textunderscore}results, hint, injective, image, partial{\textunderscore}map, colors1, colors2[, order, aut{\textunderscore}grp]})\index{HomomorphismDigraphsFinder@\texttt{HomomorphismDigraphsFinder}}
\label{HomomorphismDigraphsFinder}
}\hfill{\scriptsize (function)}}\\
\textbf{\indent Returns:\ }
The argument \mbox{\texttt{\mdseries\slshape user{\textunderscore}param}}.



 This function finds homomorphisms from the digraph \mbox{\texttt{\mdseries\slshape D1}} to the digraph \mbox{\texttt{\mdseries\slshape D2}} subject to the conditions imposed by the other arguments as described below.

 If \texttt{f} and \texttt{g} are homomorphisms found by \texttt{HomomorphismDigraphsFinder}, then \texttt{f} cannot be obtained from \texttt{g} by right multiplying by an automorphism of \mbox{\texttt{\mdseries\slshape D2}} in \mbox{\texttt{\mdseries\slshape aut{\textunderscore}grp}}. 
\begin{description}
\item[{\mbox{\texttt{\mdseries\slshape hook}}}]  This argument should be a function or \texttt{fail}.

 If \mbox{\texttt{\mdseries\slshape hook}} is a function, then it must have two arguments \mbox{\texttt{\mdseries\slshape user{\textunderscore}param}} (see below) and a transformation \texttt{t}. The function \texttt{\mbox{\texttt{\mdseries\slshape hook}}(\mbox{\texttt{\mdseries\slshape user{\textunderscore}param}}, t)} is called every time a new homomorphism \texttt{t} is found by \texttt{HomomorphismDigraphsFinder}. If the function returns \texttt{true}, then \texttt{HomomorphismDigraphsFinder} stops and does not find any further homomorphisms. This feature might be
useful if you are searching for a homomorphism that satisfies some condition
that you cannot specify via the other arguments to \texttt{HomomorphismDigraphsFinder}. 

 If \mbox{\texttt{\mdseries\slshape hook}} is \texttt{fail}, then a default function is used which simply adds every new homomorphism
found by \texttt{HomomorphismDigraphsFinder} to \mbox{\texttt{\mdseries\slshape user{\textunderscore}param}}, which must be a mutable list in this case. 
\item[{\mbox{\texttt{\mdseries\slshape user{\textunderscore}param}}}]  If \mbox{\texttt{\mdseries\slshape hook}} is a function, then \mbox{\texttt{\mdseries\slshape user{\textunderscore}param}} can be any \textsf{GAP} object. The object \mbox{\texttt{\mdseries\slshape user{\textunderscore}param}} is used as the first argument of the function \mbox{\texttt{\mdseries\slshape hook}}. For example, \mbox{\texttt{\mdseries\slshape user{\textunderscore}param}} might be a transformation semigroup, and \texttt{\mbox{\texttt{\mdseries\slshape hook}}(\mbox{\texttt{\mdseries\slshape user{\textunderscore}param}}, t)} might set \mbox{\texttt{\mdseries\slshape user{\textunderscore}param}} to be the closure of \mbox{\texttt{\mdseries\slshape user{\textunderscore}param}} and \texttt{t}. 

 If the value of \mbox{\texttt{\mdseries\slshape hook}} is \texttt{fail}, then the value of \mbox{\texttt{\mdseries\slshape user{\textunderscore}param}} must be a mutable list. 
\item[{\mbox{\texttt{\mdseries\slshape max{\textunderscore}results}}}]  This argument should be a positive integer or \texttt{infinity}. \texttt{HomomorphismDigraphsFinder} will return after it has found \mbox{\texttt{\mdseries\slshape max{\textunderscore}results}} homomorphisms or the search is complete, whichever happens first. 
\item[{\mbox{\texttt{\mdseries\slshape hint}}}]  This argument should be a positive integer or \texttt{fail}. 

 If \mbox{\texttt{\mdseries\slshape hint}} is a positive integer, then only homorphisms of rank \mbox{\texttt{\mdseries\slshape hint}} are found.

 If \mbox{\texttt{\mdseries\slshape hint}} is \texttt{fail}, then no restriction is put on the rank of homomorphisms found. 
\item[{\mbox{\texttt{\mdseries\slshape injective}}}]  This argument should be \texttt{0}, \texttt{1}, or \texttt{2}. If it is \texttt{2}, then only embeddings are found, if it is \texttt{1}, then only injective homomorphisms are found, and if it is \texttt{0} there are no restrictions imposed by this argument. 

 For backwards compatibility, \mbox{\texttt{\mdseries\slshape injective}} can also be \texttt{false} or \texttt{true} which correspond to the values \texttt{0} and \texttt{1} described in the previous paragraph, respectively. 
\item[{\mbox{\texttt{\mdseries\slshape image}}}]  This argument should be a subset of the vertices of the graph \mbox{\texttt{\mdseries\slshape D2}}. \texttt{HomomorphismDigraphsFinder} only finds homomorphisms from \mbox{\texttt{\mdseries\slshape D1}} to the subgraph of \mbox{\texttt{\mdseries\slshape D2}} induced by the vertices \mbox{\texttt{\mdseries\slshape image}}.

 The returned homomorphisms (if any) are still "up to the action" of the group
specified by \mbox{\texttt{\mdseries\slshape aut{\textunderscore}grp}} (which is the entire automorphism group by default). This might generate
unexpected results. For example, if \mbox{\texttt{\mdseries\slshape D1}} has automorphism group where one orbit consists of, say, \texttt{1} and \texttt{2}, then \texttt{HomomorphismDigraphsFinder} will only attempt to find homomorphisms mapping \texttt{1} to \texttt{1}, and if there are no such homomorphisms with image set equal to \mbox{\texttt{\mdseries\slshape image}}, then no homomorphisms will be returned (even if there is a homomorphism from \mbox{\texttt{\mdseries\slshape D1}} to \mbox{\texttt{\mdseries\slshape D2}} mapping \texttt{1} to \texttt{2}). To ensure that that \textsc{all} homomorphisms with image set equal to \mbox{\texttt{\mdseries\slshape image}} are considered it is necessary for the last argument \mbox{\texttt{\mdseries\slshape aut{\textunderscore}grp}} to be the trivial permutation group. 
\item[{\mbox{\texttt{\mdseries\slshape partial{\textunderscore}map}}}]  This argument should be a partial map from \mbox{\texttt{\mdseries\slshape D1}} to \mbox{\texttt{\mdseries\slshape D2}}, that is, a (not necessarily dense) list of vertices of the digraph \mbox{\texttt{\mdseries\slshape D2}} of length no greater than the number vertices in the digraph \mbox{\texttt{\mdseries\slshape D1}}. \texttt{HomomorphismDigraphsFinder} only finds homomorphisms extending \mbox{\texttt{\mdseries\slshape partial{\textunderscore}map}} (if any). 
\item[{\mbox{\texttt{\mdseries\slshape colors1}}}]  This should be a list representing possible colours of vertices in the digraph \mbox{\texttt{\mdseries\slshape D1}}; see \texttt{AutomorphismGroup} (\ref{AutomorphismGroup:for a digraph and a homogeneous list}) for details of the permissible values for this argument. 
\item[{\mbox{\texttt{\mdseries\slshape colors2}}}]  This should be a list representing possible colours of vertices in the digraph \mbox{\texttt{\mdseries\slshape D2}}; see \texttt{AutomorphismGroup} (\ref{AutomorphismGroup:for a digraph and a homogeneous list}) for details of the permissible values for this argument. 
\item[{\mbox{\texttt{\mdseries\slshape order}}}]  The optional argument \mbox{\texttt{\mdseries\slshape order}} specifies the order the vertices in \mbox{\texttt{\mdseries\slshape D1}} appear in the search for homomorphisms. The value of this parameter can have a
large impact on the runtime of the function. It seems in many cases to be a
good idea for this to be the \texttt{DigraphWelshPowellOrder} (\ref{DigraphWelshPowellOrder}), i.e. vertices ordered from highest to lowest degree. 
\item[{\mbox{\texttt{\mdseries\slshape aut{\textunderscore}grp}}}]  The optional argument \mbox{\texttt{\mdseries\slshape aut{\textunderscore}grp}} should be a subgroup of the automorphism group of \mbox{\texttt{\mdseries\slshape D2}}. This function returns unique representatives of the homomorphisms found up
to right multiplication by \mbox{\texttt{\mdseries\slshape aut{\textunderscore}grp}}. If this argument is not specific, it defaults to the full automorphism group
of \mbox{\texttt{\mdseries\slshape D2}}, which may be costly to calculate. 
\end{description}
 
\begin{Verbatim}[commandchars=@|B,fontsize=\small,frame=single,label=Example]
  @gapprompt|gap>B @gapinput|D := ChainDigraph(10);B
  <immutable chain digraph with 10 vertices>
  @gapprompt|gap>B @gapinput|D := DigraphSymmetricClosure(D);B
  <immutable symmetric digraph with 10 vertices, 18 edges>
  @gapprompt|gap>B @gapinput|HomomorphismDigraphsFinder(D, D, fail, [], infinity, 2, 0,B
  @gapprompt|>B @gapinput|[3, 4], [], fail, fail);B
  #I  WARNING you are trying to find homomorphisms by specifying a subset
  of the vertices of the target digraph. This might lead to unexpected
  results! If this happens, try passing Group(()) as the last argument.
  Please see the documentation of HomomorphismDigraphsFinder for details.
  [ Transformation( [ 3, 4, 3, 4, 3, 4, 3, 4, 3, 4 ] ),
    Transformation( [ 4, 3, 4, 3, 4, 3, 4, 3, 4, 3 ] ) ]
  @gapprompt|gap>B @gapinput|D2 := CompleteDigraph(6);;B
  @gapprompt|gap>B @gapinput|HomomorphismDigraphsFinder(D, D2, fail, [], 1, fail, 0,B
  @gapprompt|>B @gapinput|[1 .. 6], [1, 2, 1], fail, fail);B
  [ Transformation( [ 1, 2, 1, 3, 4, 5, 6, 1, 2, 1 ] ) ]
  @gapprompt|gap>B @gapinput|func := function(user_param, t)B
  @gapprompt|>B @gapinput|Add(user_param, t * user_param[1]);B
  @gapprompt|>B @gapinput|end;;B
  @gapprompt|gap>B @gapinput|HomomorphismDigraphsFinder(D, D2, func, [Transformation([2, 2])],B
  @gapprompt|>B @gapinput|3, fail, 0, [1 .. 6], [1, 2, 1], fail, fail);B
  [ Transformation( [ 2, 2 ] ),
    Transformation( [ 2, 2, 2, 3, 4, 5, 6, 2, 2, 2 ] ),
    Transformation( [ 2, 2, 2, 3, 4, 5, 6, 2, 2, 3 ] ),
    Transformation( [ 2, 2, 2, 3, 4, 5, 6, 2, 2, 4 ] ) ]
  @gapprompt|gap>B @gapinput|HomomorphismDigraphsFinder(NullDigraph(2), NullDigraph(3), fail,B
  @gapprompt|>B @gapinput|[], infinity, fail, 1, [1, 2, 3], fail, fail, fail, fail,B
  @gapprompt|>B @gapinput|Group(()));B
  [ IdentityTransformation, Transformation( [ 1, 3, 3 ] ),
    Transformation( [ 2, 1 ] ), Transformation( [ 2, 3, 3 ] ),
    Transformation( [ 3, 1, 3 ] ), Transformation( [ 3, 2, 3 ] ) ]
  @gapprompt|gap>B @gapinput|HomomorphismDigraphsFinder(NullDigraph(2), NullDigraph(3), fail,B
  @gapprompt|>B @gapinput|[], infinity, fail, 1, [1, 2, 3], fail, fail, fail, fail,B
  @gapprompt|>B @gapinput|Group((1, 2)));B
  [ IdentityTransformation, Transformation( [ 1, 3, 3 ] ),
    Transformation( [ 3, 1, 3 ] ) ]
\end{Verbatim}
 }

 

\subsection{\textcolor{Chapter }{DigraphHomomorphism}}
\logpage{[ 7, 3, 2 ]}\nobreak
\hyperdef{L}{X85E9B019877AD7FE}{}
{\noindent\textcolor{FuncColor}{$\triangleright$\enspace\texttt{DigraphHomomorphism({\mdseries\slshape digraph1, digraph2})\index{DigraphHomomorphism@\texttt{DigraphHomomorphism}}
\label{DigraphHomomorphism}
}\hfill{\scriptsize (operation)}}\\
\textbf{\indent Returns:\ }
 A transformation, or \texttt{fail}.



 A homomorphism from \mbox{\texttt{\mdseries\slshape digraph1}} to \mbox{\texttt{\mdseries\slshape digraph2}} is a mapping from the vertex set of \mbox{\texttt{\mdseries\slshape digraph1}} to a subset of the vertices of \mbox{\texttt{\mdseries\slshape digraph2}}, such that every pair of vertices \texttt{[i,j]} which has an edge \texttt{i\texttt{\symbol{45}}{\textgreater}j} is mapped to a pair of vertices \texttt{[a,b]} which has an edge \texttt{a\texttt{\symbol{45}}{\textgreater}b}. Note that non\texttt{\symbol{45}}adjacent vertices can still be mapped to
adjacent vertices. 

 \texttt{DigraphHomomorphism} returns a single homomorphism between \mbox{\texttt{\mdseries\slshape digraph1}} and \mbox{\texttt{\mdseries\slshape digraph2}} if it exists, otherwise it returns \texttt{fail}. 
\begin{Verbatim}[commandchars=!@|,fontsize=\small,frame=single,label=Example]
  !gapprompt@gap>| !gapinput@gr1 := ChainDigraph(3);;|
  !gapprompt@gap>| !gapinput@gr2 := Digraph([[3, 5], [2], [3, 1], [], [4]]);|
  <immutable digraph with 5 vertices, 6 edges>
  !gapprompt@gap>| !gapinput@DigraphHomomorphism(gr1, gr1);|
  IdentityTransformation
  !gapprompt@gap>| !gapinput@map := DigraphHomomorphism(gr1, gr2);|
  Transformation( [ 3, 1, 5, 4, 5 ] )
  !gapprompt@gap>| !gapinput@IsDigraphHomomorphism(gr1, gr2, map);|
  true
\end{Verbatim}
 }

 

\subsection{\textcolor{Chapter }{HomomorphismsDigraphs}}
\logpage{[ 7, 3, 3 ]}\nobreak
\hyperdef{L}{X8653EDA183E06D05}{}
{\noindent\textcolor{FuncColor}{$\triangleright$\enspace\texttt{HomomorphismsDigraphs({\mdseries\slshape digraph1, digraph2})\index{HomomorphismsDigraphs@\texttt{HomomorphismsDigraphs}}
\label{HomomorphismsDigraphs}
}\hfill{\scriptsize (operation)}}\\
\noindent\textcolor{FuncColor}{$\triangleright$\enspace\texttt{HomomorphismsDigraphsRepresentatives({\mdseries\slshape digraph1, digraph2})\index{HomomorphismsDigraphsRepresentatives@\texttt{Homomorphisms}\-\texttt{Digraphs}\-\texttt{Representatives}}
\label{HomomorphismsDigraphsRepresentatives}
}\hfill{\scriptsize (operation)}}\\
\textbf{\indent Returns:\ }
 A list of transformations.



 \texttt{HomomorphismsDigraphsRepresentatives} finds every \texttt{DigraphHomomorphism} (\ref{DigraphHomomorphism}) between \mbox{\texttt{\mdseries\slshape digraph1}} and \mbox{\texttt{\mdseries\slshape digraph2}}, up to right multiplication by an element of the \texttt{AutomorphismGroup} (\ref{AutomorphismGroup:for a digraph}) of \mbox{\texttt{\mdseries\slshape digraph2}}. In other words, every homomorphism \texttt{f} between \mbox{\texttt{\mdseries\slshape digraph1}} and \mbox{\texttt{\mdseries\slshape digraph2}} can be written as the composition \texttt{f = g * x}, where \texttt{g} is one of the \texttt{HomomorphismsDigraphsRepresentatives} and \texttt{x} is an automorphism of \mbox{\texttt{\mdseries\slshape digraph2}}. 

 \texttt{HomomorphismsDigraphs} returns all homomorphisms between \mbox{\texttt{\mdseries\slshape digraph1}} and \mbox{\texttt{\mdseries\slshape digraph2}}. 
\begin{Verbatim}[commandchars=!@|,fontsize=\small,frame=single,label=Example]
  !gapprompt@gap>| !gapinput@gr1 := ChainDigraph(3);;|
  !gapprompt@gap>| !gapinput@gr2 := Digraph([[3, 5], [2], [3, 1], [], [4]]);|
  <immutable digraph with 5 vertices, 6 edges>
  !gapprompt@gap>| !gapinput@HomomorphismsDigraphs(gr1, gr2);|
  [ Transformation( [ 1, 3, 1 ] ), Transformation( [ 1, 3, 3 ] ),
    Transformation( [ 1, 5, 4, 4, 5 ] ), Transformation( [ 2, 2, 2 ] ),
    Transformation( [ 3, 1, 3 ] ), Transformation( [ 3, 1, 5, 4, 5 ] ),
    Transformation( [ 3, 3, 1 ] ), Transformation( [ 3, 3, 3 ] ) ]
  !gapprompt@gap>| !gapinput@HomomorphismsDigraphsRepresentatives(gr1, CompleteDigraph(3));|
  [ Transformation( [ 2, 1 ] ), Transformation( [ 2, 1, 2 ] ) ]
\end{Verbatim}
 }

 

\subsection{\textcolor{Chapter }{DigraphMonomorphism}}
\logpage{[ 7, 3, 4 ]}\nobreak
\hyperdef{L}{X82D0FCD87D47ACA8}{}
{\noindent\textcolor{FuncColor}{$\triangleright$\enspace\texttt{DigraphMonomorphism({\mdseries\slshape digraph1, digraph2})\index{DigraphMonomorphism@\texttt{DigraphMonomorphism}}
\label{DigraphMonomorphism}
}\hfill{\scriptsize (operation)}}\\
\textbf{\indent Returns:\ }
A transformation, or \texttt{fail}.



 \texttt{DigraphMonomorphism} returns a single \emph{injective} \texttt{DigraphHomomorphism} (\ref{DigraphHomomorphism}) between \mbox{\texttt{\mdseries\slshape digraph1}} and \mbox{\texttt{\mdseries\slshape digraph2}} if one exists, otherwise it returns \texttt{fail}. 
\begin{Verbatim}[commandchars=!@|,fontsize=\small,frame=single,label=Example]
  !gapprompt@gap>| !gapinput@gr1 := ChainDigraph(3);;|
  !gapprompt@gap>| !gapinput@gr2 := Digraph([[3, 5], [2], [3, 1], [], [4]]);|
  <immutable digraph with 5 vertices, 6 edges>
  !gapprompt@gap>| !gapinput@DigraphMonomorphism(gr1, gr1);|
  IdentityTransformation
  !gapprompt@gap>| !gapinput@DigraphMonomorphism(gr1, gr2);|
  Transformation( [ 3, 1, 5, 4, 5 ] )
\end{Verbatim}
 }

 

\subsection{\textcolor{Chapter }{MonomorphismsDigraphs}}
\logpage{[ 7, 3, 5 ]}\nobreak
\hyperdef{L}{X84B4BDB984C2A9A8}{}
{\noindent\textcolor{FuncColor}{$\triangleright$\enspace\texttt{MonomorphismsDigraphs({\mdseries\slshape digraph1, digraph2})\index{MonomorphismsDigraphs@\texttt{MonomorphismsDigraphs}}
\label{MonomorphismsDigraphs}
}\hfill{\scriptsize (operation)}}\\
\noindent\textcolor{FuncColor}{$\triangleright$\enspace\texttt{MonomorphismsDigraphsRepresentatives({\mdseries\slshape digraph1, digraph2})\index{MonomorphismsDigraphsRepresentatives@\texttt{Monomorphisms}\-\texttt{Digraphs}\-\texttt{Representatives}}
\label{MonomorphismsDigraphsRepresentatives}
}\hfill{\scriptsize (operation)}}\\
\textbf{\indent Returns:\ }
A list of transformations.



 These operations behave the same as \texttt{HomomorphismsDigraphs} (\ref{HomomorphismsDigraphs}) and \texttt{HomomorphismsDigraphsRepresentatives} (\ref{HomomorphismsDigraphsRepresentatives}), except they only return \emph{injective} homomorphisms. 
\begin{Verbatim}[commandchars=!@|,fontsize=\small,frame=single,label=Example]
  !gapprompt@gap>| !gapinput@gr1 := ChainDigraph(3);;|
  !gapprompt@gap>| !gapinput@gr2 := Digraph([[3, 5], [2], [3, 1], [], [4]]);|
  <immutable digraph with 5 vertices, 6 edges>
  !gapprompt@gap>| !gapinput@MonomorphismsDigraphs(gr1, gr2);|
  [ Transformation( [ 1, 5, 4, 4, 5 ] ),
    Transformation( [ 3, 1, 5, 4, 5 ] ) ]
  !gapprompt@gap>| !gapinput@MonomorphismsDigraphsRepresentatives(gr1, CompleteDigraph(3));|
  [ Transformation( [ 2, 1 ] ) ]
\end{Verbatim}
 }

 

\subsection{\textcolor{Chapter }{DigraphEpimorphism}}
\logpage{[ 7, 3, 6 ]}\nobreak
\hyperdef{L}{X7CB5AD9F861684FD}{}
{\noindent\textcolor{FuncColor}{$\triangleright$\enspace\texttt{DigraphEpimorphism({\mdseries\slshape digraph1, digraph2})\index{DigraphEpimorphism@\texttt{DigraphEpimorphism}}
\label{DigraphEpimorphism}
}\hfill{\scriptsize (operation)}}\\
\textbf{\indent Returns:\ }
A transformation, or \texttt{fail}.



 \texttt{DigraphEpimorphism} returns a single \emph{surjective} \texttt{DigraphHomomorphism} (\ref{DigraphHomomorphism}) between \mbox{\texttt{\mdseries\slshape digraph1}} and \mbox{\texttt{\mdseries\slshape digraph2}} if one exists, otherwise it returns \texttt{fail}. 
\begin{Verbatim}[commandchars=!@|,fontsize=\small,frame=single,label=Example]
  !gapprompt@gap>| !gapinput@gr1 := DigraphReverse(ChainDigraph(4));|
  <immutable digraph with 4 vertices, 3 edges>
  !gapprompt@gap>| !gapinput@gr2 := DigraphRemoveEdge(CompleteDigraph(3), [1, 2]);|
  <immutable digraph with 3 vertices, 5 edges>
  !gapprompt@gap>| !gapinput@DigraphEpimorphism(gr2, gr1);|
  fail
  !gapprompt@gap>| !gapinput@DigraphEpimorphism(gr1, gr2);|
  Transformation( [ 3, 1, 2, 3 ] )
\end{Verbatim}
 }

 

\subsection{\textcolor{Chapter }{EpimorphismsDigraphs}}
\logpage{[ 7, 3, 7 ]}\nobreak
\hyperdef{L}{X7DFB1F5D873937B3}{}
{\noindent\textcolor{FuncColor}{$\triangleright$\enspace\texttt{EpimorphismsDigraphs({\mdseries\slshape digraph1, digraph2})\index{EpimorphismsDigraphs@\texttt{EpimorphismsDigraphs}}
\label{EpimorphismsDigraphs}
}\hfill{\scriptsize (operation)}}\\
\noindent\textcolor{FuncColor}{$\triangleright$\enspace\texttt{EpimorphismsDigraphsRepresentatives({\mdseries\slshape digraph1, digraph2})\index{EpimorphismsDigraphsRepresentatives@\texttt{EpimorphismsDigraphsRepresentatives}}
\label{EpimorphismsDigraphsRepresentatives}
}\hfill{\scriptsize (operation)}}\\
\textbf{\indent Returns:\ }
A list of transformations.



 These operations behave the same as \texttt{HomomorphismsDigraphs} (\ref{HomomorphismsDigraphs}) and \texttt{HomomorphismsDigraphsRepresentatives} (\ref{HomomorphismsDigraphsRepresentatives}), except they only return \emph{surjective} homomorphisms. 
\begin{Verbatim}[commandchars=!@|,fontsize=\small,frame=single,label=Example]
  !gapprompt@gap>| !gapinput@gr1 := DigraphReverse(ChainDigraph(4));|
  <immutable digraph with 4 vertices, 3 edges>
  !gapprompt@gap>| !gapinput@gr2 := DigraphSymmetricClosure(CycleDigraph(3));|
  <immutable symmetric digraph with 3 vertices, 6 edges>
  !gapprompt@gap>| !gapinput@EpimorphismsDigraphsRepresentatives(gr1, gr2);|
  [ Transformation( [ 3, 1, 2, 1 ] ), Transformation( [ 3, 1, 2, 3 ] ),
    Transformation( [ 2, 1, 2, 3 ] ) ]
  !gapprompt@gap>| !gapinput@EpimorphismsDigraphs(gr1, gr2);|
  [ Transformation( [ 1, 2, 1, 3 ] ), Transformation( [ 1, 2, 3, 1 ] ),
    Transformation( [ 1, 2, 3, 2 ] ), Transformation( [ 1, 3, 1, 2 ] ),
    Transformation( [ 1, 3, 2, 1 ] ), Transformation( [ 1, 3, 2, 3 ] ),
    Transformation( [ 2, 1, 2, 3 ] ), Transformation( [ 2, 1, 3, 1 ] ),
    Transformation( [ 2, 1, 3, 2 ] ), Transformation( [ 2, 3, 1, 2 ] ),
    Transformation( [ 2, 3, 1, 3 ] ), Transformation( [ 2, 3, 2, 1 ] ),
    Transformation( [ 3, 1, 2, 1 ] ), Transformation( [ 3, 1, 2, 3 ] ),
    Transformation( [ 3, 1, 3, 2 ] ), Transformation( [ 3, 2, 1, 2 ] ),
    Transformation( [ 3, 2, 1, 3 ] ), Transformation( [ 3, 2, 3, 1 ] ) ]
\end{Verbatim}
 }

 

\subsection{\textcolor{Chapter }{DigraphEmbedding}}
\logpage{[ 7, 3, 8 ]}\nobreak
\hyperdef{L}{X7AD04CC186E86CCA}{}
{\noindent\textcolor{FuncColor}{$\triangleright$\enspace\texttt{DigraphEmbedding({\mdseries\slshape digraph1, digraph2})\index{DigraphEmbedding@\texttt{DigraphEmbedding}}
\label{DigraphEmbedding}
}\hfill{\scriptsize (operation)}}\\
\textbf{\indent Returns:\ }
 A transformation, or \texttt{fail}.



 An embedding of a digraph \mbox{\texttt{\mdseries\slshape digraph1}} into another digraph \mbox{\texttt{\mdseries\slshape digraph2}} is a \texttt{DigraphMonomorphism} (\ref{DigraphMonomorphism}) from \mbox{\texttt{\mdseries\slshape digraph1}} to \mbox{\texttt{\mdseries\slshape digraph2}} which has the additional property that a pair of vertices \texttt{[i, j]} which have no edge \texttt{i \texttt{\symbol{45}}{\textgreater} j} in \mbox{\texttt{\mdseries\slshape digraph1}} are mapped to a pair of vertices \texttt{[a, b]} which have no edge \texttt{a\texttt{\symbol{45}}{\textgreater}b} in \mbox{\texttt{\mdseries\slshape digraph2}}.

 In other words, an embedding \texttt{t} is an isomorphism from \mbox{\texttt{\mdseries\slshape digraph1}} to the \texttt{InducedSubdigraph} (\ref{InducedSubdigraph}) of \mbox{\texttt{\mdseries\slshape digraph2}} on the image of \texttt{t}. 

 \texttt{DigraphEmbedding} returns a single embedding if one exists, otherwise it returns \texttt{fail}. 
\begin{Verbatim}[commandchars=!@|,fontsize=\small,frame=single,label=Example]
  !gapprompt@gap>| !gapinput@gr := ChainDigraph(3);|
  <immutable chain digraph with 3 vertices>
  !gapprompt@gap>| !gapinput@DigraphEmbedding(gr, CompleteDigraph(4));|
  fail
  !gapprompt@gap>| !gapinput@DigraphEmbedding(gr, Digraph([[3], [1, 4], [1], [3]]));|
  Transformation( [ 2, 4, 3, 4 ] )
\end{Verbatim}
 }

 \index{subgraph isomorphism@subgraph isomorphism} 

\subsection{\textcolor{Chapter }{EmbeddingsDigraphs}}
\logpage{[ 7, 3, 9 ]}\nobreak
\hyperdef{L}{X82EBDF137EEDC5FD}{}
{\noindent\textcolor{FuncColor}{$\triangleright$\enspace\texttt{EmbeddingsDigraphs({\mdseries\slshape D1, D2})\index{EmbeddingsDigraphs@\texttt{EmbeddingsDigraphs}}
\label{EmbeddingsDigraphs}
}\hfill{\scriptsize (operation)}}\\
\noindent\textcolor{FuncColor}{$\triangleright$\enspace\texttt{EmbeddingsDigraphsRepresentatives({\mdseries\slshape D1, D2})\index{EmbeddingsDigraphsRepresentatives@\texttt{EmbeddingsDigraphsRepresentatives}}
\label{EmbeddingsDigraphsRepresentatives}
}\hfill{\scriptsize (operation)}}\\
\textbf{\indent Returns:\ }
A list of transformations.



 These operations behave the same as \texttt{HomomorphismsDigraphs} (\ref{HomomorphismsDigraphs}) and \texttt{HomomorphismsDigraphsRepresentatives} (\ref{HomomorphismsDigraphsRepresentatives}), except they only return embeddings of \mbox{\texttt{\mdseries\slshape D1}} into \mbox{\texttt{\mdseries\slshape D2}}.

 See also \texttt{IsDigraphEmbedding} (\ref{IsDigraphEmbedding}). 
\begin{Verbatim}[commandchars=!@|,fontsize=\small,frame=single,label=Example]
  !gapprompt@gap>| !gapinput@D1 := NullDigraph(2);|
  <immutable empty digraph with 2 vertices>
  !gapprompt@gap>| !gapinput@D2 := CycleDigraph(5);|
  <immutable cycle digraph with 5 vertices>
  !gapprompt@gap>| !gapinput@EmbeddingsDigraphsRepresentatives(D1, D2);|
  [ Transformation( [ 1, 3, 3 ] ), Transformation( [ 1, 4, 3, 4 ] ) ]
  !gapprompt@gap>| !gapinput@EmbeddingsDigraphs(D1, D2);|
  [ Transformation( [ 1, 3, 3 ] ), Transformation( [ 1, 4, 3, 4 ] ),
    Transformation( [ 2, 4, 4, 5, 1 ] ),
    Transformation( [ 2, 5, 4, 5, 1 ] ),
    Transformation( [ 3, 1, 5, 1, 2 ] ),
    Transformation( [ 3, 5, 5, 1, 2 ] ),
    Transformation( [ 4, 1, 1, 2, 3 ] ),
    Transformation( [ 4, 2, 1, 2, 3 ] ),
    Transformation( [ 5, 2, 2, 3, 4 ] ),
    Transformation( [ 5, 3, 2, 3, 4 ] ) ]
\end{Verbatim}
 }

 

\subsection{\textcolor{Chapter }{IsDigraphHomomorphism (for digraphs and a permutation or transformation)}}
\logpage{[ 7, 3, 10 ]}\nobreak
\hyperdef{L}{X80FE9FCB79718A61}{}
{\noindent\textcolor{FuncColor}{$\triangleright$\enspace\texttt{IsDigraphHomomorphism({\mdseries\slshape src, ran, x})\index{IsDigraphHomomorphism@\texttt{IsDigraphHomomorphism}!for digraphs and a permutation or transformation}
\label{IsDigraphHomomorphism:for digraphs and a permutation or transformation}
}\hfill{\scriptsize (operation)}}\\
\noindent\textcolor{FuncColor}{$\triangleright$\enspace\texttt{IsDigraphHomomorphism({\mdseries\slshape src, ran, x, col1, col2})\index{IsDigraphHomomorphism@\texttt{IsDigraphHomomorphism}}
\label{IsDigraphHomomorphism}
}\hfill{\scriptsize (operation)}}\\
\noindent\textcolor{FuncColor}{$\triangleright$\enspace\texttt{IsDigraphEpimorphism({\mdseries\slshape src, ran, x})\index{IsDigraphEpimorphism@\texttt{IsDigraphEpimorphism}!for digraphs and a permutation or transformation}
\label{IsDigraphEpimorphism:for digraphs and a permutation or transformation}
}\hfill{\scriptsize (operation)}}\\
\noindent\textcolor{FuncColor}{$\triangleright$\enspace\texttt{IsDigraphEpimorphism({\mdseries\slshape src, ran, x, col1, col2})\index{IsDigraphEpimorphism@\texttt{IsDigraphEpimorphism}}
\label{IsDigraphEpimorphism}
}\hfill{\scriptsize (operation)}}\\
\noindent\textcolor{FuncColor}{$\triangleright$\enspace\texttt{IsDigraphMonomorphism({\mdseries\slshape src, ran, x})\index{IsDigraphMonomorphism@\texttt{IsDigraphMonomorphism}!for digraphs and a permutation or transformation}
\label{IsDigraphMonomorphism:for digraphs and a permutation or transformation}
}\hfill{\scriptsize (operation)}}\\
\noindent\textcolor{FuncColor}{$\triangleright$\enspace\texttt{IsDigraphMonomorphism({\mdseries\slshape src, ran, x, col1, col2})\index{IsDigraphMonomorphism@\texttt{IsDigraphMonomorphism}}
\label{IsDigraphMonomorphism}
}\hfill{\scriptsize (operation)}}\\
\noindent\textcolor{FuncColor}{$\triangleright$\enspace\texttt{IsDigraphEndomorphism({\mdseries\slshape digraph, x})\index{IsDigraphEndomorphism@\texttt{IsDigraphEndomorphism}!for digraphs and a permutation or transformation}
\label{IsDigraphEndomorphism:for digraphs and a permutation or transformation}
}\hfill{\scriptsize (operation)}}\\
\noindent\textcolor{FuncColor}{$\triangleright$\enspace\texttt{IsDigraphEndomorphism({\mdseries\slshape digraph, x, col})\index{IsDigraphEndomorphism@\texttt{IsDigraphEndomorphism}}
\label{IsDigraphEndomorphism}
}\hfill{\scriptsize (operation)}}\\
\textbf{\indent Returns:\ }
\texttt{true} or \texttt{false}.



 \texttt{IsDigraphHomomorphism} returns \texttt{true} if the permutation or transformation \mbox{\texttt{\mdseries\slshape x}} is a homomorphism from the digraph \mbox{\texttt{\mdseries\slshape src}} to the digraph \mbox{\texttt{\mdseries\slshape ran}}. 

 \texttt{IsDigraphEpimorphism} returns \texttt{true} if the permutation or transformation \mbox{\texttt{\mdseries\slshape x}} is a surjective homomorphism from the digraph \mbox{\texttt{\mdseries\slshape src}} to the digraph \mbox{\texttt{\mdseries\slshape ran}}. 

 \texttt{IsDigraphMonomorphism} returns \texttt{true} if the permutation or transformation \mbox{\texttt{\mdseries\slshape x}} is an injective homomorphism from the digraph \mbox{\texttt{\mdseries\slshape src}} to the digraph \mbox{\texttt{\mdseries\slshape ran}}. 

 \texttt{IsDigraphEndomorphism} returns \texttt{true} if the permutation or transformation \mbox{\texttt{\mdseries\slshape x}} is an endomorphism of the digraph \mbox{\texttt{\mdseries\slshape digraph}}. 

 A permutation or transformation \mbox{\texttt{\mdseries\slshape x}} is a \emph{homomorphism} from a digraph \mbox{\texttt{\mdseries\slshape src}} to a digraph \mbox{\texttt{\mdseries\slshape ran}} if the following hold: 
\begin{itemize}
\item  \texttt{[u \texttt{\symbol{94}} \mbox{\texttt{\mdseries\slshape x}}, v \texttt{\symbol{94}} \mbox{\texttt{\mdseries\slshape x}}]} is an edge of \mbox{\texttt{\mdseries\slshape ran}} whenever \texttt{[u, v]} is an edge of \mbox{\texttt{\mdseries\slshape src}}; and 
\item  \mbox{\texttt{\mdseries\slshape x}} maps the vertices of \mbox{\texttt{\mdseries\slshape src}} to a subset of the vertices of \mbox{\texttt{\mdseries\slshape ran}}, i.e. \texttt{IsSubset(DigraphVertices(\mbox{\texttt{\mdseries\slshape ran}}), OnSets(DigraphVertices(\mbox{\texttt{\mdseries\slshape src}}), \mbox{\texttt{\mdseries\slshape x}}))} is \texttt{true}. 
\end{itemize}
 Note that if \texttt{i} is any integer greater than \texttt{DigraphNrVertice(\mbox{\texttt{\mdseries\slshape src}})}, then the action of \mbox{\texttt{\mdseries\slshape x}} on \texttt{i} is ignored by this function. One consequence of this is that distinct
transformations or permutations might represent the same homomorphism. For
example, if \mbox{\texttt{\mdseries\slshape src}} and \mbox{\texttt{\mdseries\slshape ran}} are \texttt{CycleDigraph(2)}, then both the permutations \texttt{(1, 2)} and \texttt{(1, 2)(3, 4)} represent the same automorphism of \mbox{\texttt{\mdseries\slshape src}}. 

 See also \texttt{GeneratorsOfEndomorphismMonoid} (\ref{GeneratorsOfEndomorphismMonoid}).

 If \mbox{\texttt{\mdseries\slshape col1}} and \mbox{\texttt{\mdseries\slshape col2}}, or \mbox{\texttt{\mdseries\slshape col}}, are given, then they must represent vertex colourings; see \texttt{AutomorphismGroup} (\ref{AutomorphismGroup:for a digraph and a homogeneous list}) for details of the permissible values for these argument. The homomorphism
must then also have the property: 
\begin{itemize}
\item  \texttt{col[i] = col[i \texttt{\symbol{94}} x]} for all vertices \texttt{i} of \mbox{\texttt{\mdseries\slshape digraph}}, in the case of \texttt{IsDigraphEndomorphism}.
\item  \texttt{col1[i] = col2[i \texttt{\symbol{94}} x]} for all vertices \texttt{i} of \mbox{\texttt{\mdseries\slshape src}}, in the cases of the other operations.
\end{itemize}
 See also \texttt{DigraphsRespectsColouring} (\ref{DigraphsRespectsColouring}). 
\begin{Verbatim}[commandchars=!@|,fontsize=\small,frame=single,label=Example]
  !gapprompt@gap>| !gapinput@src := Digraph([[1], [1, 2], [1, 3]]);|
  <immutable digraph with 3 vertices, 5 edges>
  !gapprompt@gap>| !gapinput@ran := Digraph([[1], [1, 2]]);|
  <immutable digraph with 2 vertices, 3 edges>
  !gapprompt@gap>| !gapinput@IsDigraphHomomorphism(src, ran, Transformation([1, 2, 2]));|
  true
  !gapprompt@gap>| !gapinput@IsDigraphHomomorphism(src, ran, Transformation([2, 1, 2]));|
  false
  !gapprompt@gap>| !gapinput@IsDigraphHomomorphism(src, ran, Transformation([3, 3, 3]));|
  false
  !gapprompt@gap>| !gapinput@IsDigraphHomomorphism(src, src, Transformation([3, 3, 3]));|
  true
  !gapprompt@gap>| !gapinput@IsDigraphHomomorphism(src, ran, Transformation([1, 2, 2]),|
  !gapprompt@>| !gapinput@                         [1, 2, 2], [1, 2]);|
  true
  !gapprompt@gap>| !gapinput@IsDigraphHomomorphism(src, ran, Transformation([1, 2, 2]),|
  !gapprompt@>| !gapinput@                         [2, 1, 1], [1, 2]);|
  false
  !gapprompt@gap>| !gapinput@IsDigraphEndomorphism(src, Transformation([3, 3, 3]));|
  true
  !gapprompt@gap>| !gapinput@IsDigraphEndomorphism(src, Transformation([3, 3, 3]), [1, 1, 1]);|
  true
  !gapprompt@gap>| !gapinput@IsDigraphEndomorphism(src, Transformation([3, 3, 3]), [1, 1, 2]);|
  false
  !gapprompt@gap>| !gapinput@IsDigraphEpimorphism(src, ran, Transformation([3, 3, 3]));|
  false
  !gapprompt@gap>| !gapinput@IsDigraphMonomorphism(src, ran, Transformation([1, 2, 2]));|
  false
  !gapprompt@gap>| !gapinput@IsDigraphEpimorphism(src, ran, Transformation([1, 2, 2]));|
  true
  !gapprompt@gap>| !gapinput@IsDigraphMonomorphism(ran, src, ());|
  true
\end{Verbatim}
 }

 

\subsection{\textcolor{Chapter }{IsDigraphEmbedding (for digraphs and a permutation or transformation)}}
\logpage{[ 7, 3, 11 ]}\nobreak
\hyperdef{L}{X7D6B5A0887903810}{}
{\noindent\textcolor{FuncColor}{$\triangleright$\enspace\texttt{IsDigraphEmbedding({\mdseries\slshape src, ran, x})\index{IsDigraphEmbedding@\texttt{IsDigraphEmbedding}!for digraphs and a permutation or transformation}
\label{IsDigraphEmbedding:for digraphs and a permutation or transformation}
}\hfill{\scriptsize (operation)}}\\
\noindent\textcolor{FuncColor}{$\triangleright$\enspace\texttt{IsDigraphEmbedding({\mdseries\slshape src, ran, x, col1, col2})\index{IsDigraphEmbedding@\texttt{IsDigraphEmbedding}}
\label{IsDigraphEmbedding}
}\hfill{\scriptsize (operation)}}\\
\textbf{\indent Returns:\ }
\texttt{true} or \texttt{false}.



 \texttt{IsDigraphEmbedding} returns \texttt{true} if the permutation or transformation \mbox{\texttt{\mdseries\slshape x}} is a embedding of the digraph \mbox{\texttt{\mdseries\slshape src}} into the digraph \mbox{\texttt{\mdseries\slshape ran}}, while respecting the colourings \mbox{\texttt{\mdseries\slshape col1}} and \mbox{\texttt{\mdseries\slshape col2}} if given. 

 A permutation or transformation \mbox{\texttt{\mdseries\slshape x}} is a \emph{embedding} of a digraph \mbox{\texttt{\mdseries\slshape src}} into a digraph \mbox{\texttt{\mdseries\slshape ran}} if \mbox{\texttt{\mdseries\slshape x}} is a monomorphism from \mbox{\texttt{\mdseries\slshape src}} to \mbox{\texttt{\mdseries\slshape ran}} and the inverse of \mbox{\texttt{\mdseries\slshape x}} is a monomorphism from the subdigraph of \mbox{\texttt{\mdseries\slshape ran}} induced by the image of \mbox{\texttt{\mdseries\slshape x}} to \mbox{\texttt{\mdseries\slshape src}}. See also \texttt{IsDigraphHomomorphism} (\ref{IsDigraphHomomorphism}).

 
\begin{Verbatim}[commandchars=!@|,fontsize=\small,frame=single,label=Example]
  !gapprompt@gap>| !gapinput@src := Digraph([[1], [1, 2]]);|
  <immutable digraph with 2 vertices, 3 edges>
  !gapprompt@gap>| !gapinput@ran := Digraph([[1], [1, 2], [1, 3]]);|
  <immutable digraph with 3 vertices, 5 edges>
  !gapprompt@gap>| !gapinput@IsDigraphMonomorphism(src, ran, ());|
  true
  !gapprompt@gap>| !gapinput@IsDigraphEmbedding(src, ran, ());|
  true
  !gapprompt@gap>| !gapinput@IsDigraphEmbedding(src, ran, (), [2, 1], [2, 1, 1]);|
  true
  !gapprompt@gap>| !gapinput@IsDigraphEmbedding(src, ran, (), [2, 1], [1, 2, 1]);|
  false
  !gapprompt@gap>| !gapinput@ran := Digraph([[1, 2], [1, 2], [1, 3]]);|
  <immutable digraph with 3 vertices, 6 edges>
  !gapprompt@gap>| !gapinput@IsDigraphMonomorphism(src, ran, IdentityTransformation);|
  true
  !gapprompt@gap>| !gapinput@IsDigraphEmbedding(src, ran, IdentityTransformation);|
  false
\end{Verbatim}
 }

 \index{subgraph isomorphism@subgraph isomorphism} 

\subsection{\textcolor{Chapter }{SubdigraphsMonomorphisms}}
\logpage{[ 7, 3, 12 ]}\nobreak
\hyperdef{L}{X87E60FB17BC444EF}{}
{\noindent\textcolor{FuncColor}{$\triangleright$\enspace\texttt{SubdigraphsMonomorphisms({\mdseries\slshape digraph1, digraph2})\index{SubdigraphsMonomorphisms@\texttt{SubdigraphsMonomorphisms}}
\label{SubdigraphsMonomorphisms}
}\hfill{\scriptsize (operation)}}\\
\noindent\textcolor{FuncColor}{$\triangleright$\enspace\texttt{SubdigraphsMonomorphismsRepresentatives({\mdseries\slshape digraph1, digraph2})\index{SubdigraphsMonomorphismsRepresentatives@\texttt{Subdigraphs}\-\texttt{Monomorphisms}\-\texttt{Representatives}}
\label{SubdigraphsMonomorphismsRepresentatives}
}\hfill{\scriptsize (operation)}}\\
\textbf{\indent Returns:\ }
A list of transformations.



 These operations behave the same as \texttt{HomomorphismsDigraphs} (\ref{HomomorphismsDigraphs}) and \texttt{HomomorphismsDigraphsRepresentatives} (\ref{HomomorphismsDigraphsRepresentatives}), except they only return \emph{injective} homomorphisms with the following property: the (not necessarily induced)
subdigraphs defined by the images of these monomorphisms are all of the
subdigraphs of \mbox{\texttt{\mdseries\slshape digraph2}} that are isomorphic to \mbox{\texttt{\mdseries\slshape digraph1}}. Note that the subdigraphs of the previous sentence are those obtained by
applying the corresponding monomorphism to the vertices and the edges of \mbox{\texttt{\mdseries\slshape digraph1}}, and are therefore possibly strictly contained in the induced subdigraph on
the same vertex set. 
\begin{Verbatim}[commandchars=!@|,fontsize=\small,frame=single,label=Example]
  !gapprompt@gap>| !gapinput@Set(SubdigraphsMonomorphisms(CompleteBipartiteDigraph(2, 2),|
  !gapprompt@>| !gapinput@CompleteDigraph(4)));|
  [ Transformation( [ 1, 3, 2 ] ), Transformation( [ 2, 3, 1 ] ),
    Transformation( [ 3, 4, 2, 1 ] ) ]
  !gapprompt@gap>| !gapinput@SubdigraphsMonomorphismsRepresentatives(|
  !gapprompt@>| !gapinput@CompleteBipartiteDigraph(2, 2), CompleteDigraph(4));|
  [ Transformation( [ 1, 3, 2 ] ) ]
\end{Verbatim}
 }

 

\subsection{\textcolor{Chapter }{DigraphsRespectsColouring}}
\logpage{[ 7, 3, 13 ]}\nobreak
\hyperdef{L}{X7F80709A80CBFA98}{}
{\noindent\textcolor{FuncColor}{$\triangleright$\enspace\texttt{DigraphsRespectsColouring({\mdseries\slshape src, ran, x, col1, col2})\index{DigraphsRespectsColouring@\texttt{DigraphsRespectsColouring}}
\label{DigraphsRespectsColouring}
}\hfill{\scriptsize (operation)}}\\
\textbf{\indent Returns:\ }
 \texttt{true} or \texttt{false}. 



 The operation \texttt{DigraphsRespectsColouring} verifies whether or not the permutation or transformation \mbox{\texttt{\mdseries\slshape x}} respects the vertex colourings \mbox{\texttt{\mdseries\slshape col1}} and \mbox{\texttt{\mdseries\slshape col2}} of the digraphs \mbox{\texttt{\mdseries\slshape src}} and \mbox{\texttt{\mdseries\slshape range}}. That is, \texttt{DigraphsRespectsColouring} returns \texttt{true} if and only if for all vertices \texttt{i} of \mbox{\texttt{\mdseries\slshape src}}, \texttt{col1[i] = col2[i \texttt{\symbol{94}} x]}. 

 
\begin{Verbatim}[commandchars=!@|,fontsize=\small,frame=single,label=Example]
  !gapprompt@gap>| !gapinput@src := Digraph([[1], [1, 2]]);|
  <immutable digraph with 2 vertices, 3 edges>
  !gapprompt@gap>| !gapinput@ran := Digraph([[1], [1, 2], [1, 3]]);|
  <immutable digraph with 3 vertices, 5 edges>
  !gapprompt@gap>| !gapinput@DigraphsRespectsColouring(src, ran, (1, 2), [2, 1], [1, 2, 1]);|
  true
  !gapprompt@gap>| !gapinput@DigraphsRespectsColouring(src, ran, (1, 2), [2, 1], [2, 1, 1]);|
  false
\end{Verbatim}
 }

 

\subsection{\textcolor{Chapter }{GeneratorsOfEndomorphismMonoid}}
\logpage{[ 7, 3, 14 ]}\nobreak
\hyperdef{L}{X7E93B268823F6478}{}
{\noindent\textcolor{FuncColor}{$\triangleright$\enspace\texttt{GeneratorsOfEndomorphismMonoid({\mdseries\slshape digraph[, colors][, limit]})\index{GeneratorsOfEndomorphismMonoid@\texttt{GeneratorsOfEndomorphismMonoid}}
\label{GeneratorsOfEndomorphismMonoid}
}\hfill{\scriptsize (function)}}\\
\noindent\textcolor{FuncColor}{$\triangleright$\enspace\texttt{GeneratorsOfEndomorphismMonoidAttr({\mdseries\slshape digraph})\index{GeneratorsOfEndomorphismMonoidAttr@\texttt{GeneratorsOfEndomorphismMonoidAttr}}
\label{GeneratorsOfEndomorphismMonoidAttr}
}\hfill{\scriptsize (attribute)}}\\
\textbf{\indent Returns:\ }
 A list of transformations.



 An endomorphism of \mbox{\texttt{\mdseries\slshape digraph}} is a homomorphism \texttt{DigraphHomomorphism} (\ref{DigraphHomomorphism}) from \mbox{\texttt{\mdseries\slshape digraph}} back to itself. \texttt{GeneratorsOfEndomorphismMonoid}, called with a single argument, returns a generating set for the monoid of
all endomorphisms of \mbox{\texttt{\mdseries\slshape digraph}}. If \mbox{\texttt{\mdseries\slshape digraph}} belongs to \texttt{IsImmutableDigraph} (\ref{IsImmutableDigraph}), then the value of \texttt{GeneratorsOfEndomorphismMonoid} will not be recomputed on future calls.

 If the \mbox{\texttt{\mdseries\slshape colors}} argument is specified, then \texttt{GeneratorsOfEndomorphismMonoid} will return a generating set for the monoid of endomorphisms which respect the
given colouring. The colouring \mbox{\texttt{\mdseries\slshape colors}} can be in one of two forms: 

 
\begin{itemize}
\item  A list of positive integers of size the number of vertices of \mbox{\texttt{\mdseries\slshape digraph}}, where \mbox{\texttt{\mdseries\slshape colors}}\texttt{[i]} is the colour of vertex \texttt{i}. 
\item  A list of lists, such that \mbox{\texttt{\mdseries\slshape colors}}\texttt{[i]} is a list of all vertices with colour \texttt{i}. 
\end{itemize}
 If the \mbox{\texttt{\mdseries\slshape limit}} argument is specified, then it will return only the first \mbox{\texttt{\mdseries\slshape limit}} homomorphisms, where \mbox{\texttt{\mdseries\slshape limit}} must be a positive integer or \texttt{infinity}. 

 
\begin{Verbatim}[commandchars=!@|,fontsize=\small,frame=single,label=Example]
  !gapprompt@gap>| !gapinput@gr := Digraph(List([1 .. 3], x -> [1 .. 3]));;|
  !gapprompt@gap>| !gapinput@GeneratorsOfEndomorphismMonoid(gr);|
  [ Transformation( [ 1, 3, 2 ] ), Transformation( [ 2, 1 ] ),
    IdentityTransformation, Transformation( [ 1, 2, 1 ] ),
    Transformation( [ 1, 2, 2 ] ), Transformation( [ 1, 1, 2 ] ),
    Transformation( [ 1, 1, 1 ] ) ]
  !gapprompt@gap>| !gapinput@GeneratorsOfEndomorphismMonoid(gr, 3);|
  [ Transformation( [ 1, 3, 2 ] ), Transformation( [ 2, 1 ] ),
    IdentityTransformation ]
  !gapprompt@gap>| !gapinput@gr := CompleteDigraph(3);;|
  !gapprompt@gap>| !gapinput@GeneratorsOfEndomorphismMonoid(gr);|
  [ Transformation( [ 2, 3, 1 ] ), Transformation( [ 2, 1 ] ),
    IdentityTransformation ]
  !gapprompt@gap>| !gapinput@GeneratorsOfEndomorphismMonoid(gr, [1, 2, 2]);|
  [ Transformation( [ 1, 3, 2 ] ), IdentityTransformation ]
  !gapprompt@gap>| !gapinput@GeneratorsOfEndomorphismMonoid(gr, [[1], [2, 3]]);|
  [ Transformation( [ 1, 3, 2 ] ), IdentityTransformation ]
\end{Verbatim}
 }

 

\subsection{\textcolor{Chapter }{DigraphColouring (for a digraph and a number of colours)}}
\logpage{[ 7, 3, 15 ]}\nobreak
\hyperdef{L}{X86D0A74B83D6B9A7}{}
{\noindent\textcolor{FuncColor}{$\triangleright$\enspace\texttt{DigraphColouring({\mdseries\slshape digraph, n})\index{DigraphColouring@\texttt{DigraphColouring}!for a digraph and a number of colours}
\label{DigraphColouring:for a digraph and a number of colours}
}\hfill{\scriptsize (operation)}}\\
\textbf{\indent Returns:\ }
 A transformation, or \texttt{fail}.



 A \emph{proper colouring} of a digraph is a labelling of its vertices in such a way that adjacent
vertices have different labels. A \emph{proper \texttt{n}\texttt{\symbol{45}}colouring} is a proper colouring that uses exactly \texttt{n} colours. Equivalently, a proper (\texttt{n}\texttt{\symbol{45}})colouring of a digraph can be defined to be a \texttt{DigraphEpimorphism} (\ref{DigraphEpimorphism}) from a digraph onto the complete digraph (with \texttt{n} vertices); see \texttt{CompleteDigraph} (\ref{CompleteDigraph}). Note that a digraph with loops (\texttt{DigraphHasLoops} (\ref{DigraphHasLoops})) does not have a proper \texttt{n}\texttt{\symbol{45}}colouring for any value \texttt{n}. 

 If \mbox{\texttt{\mdseries\slshape digraph}} is a digraph and \mbox{\texttt{\mdseries\slshape n}} is a non\texttt{\symbol{45}}negative integer, then \texttt{DigraphColouring(\mbox{\texttt{\mdseries\slshape digraph}}, \mbox{\texttt{\mdseries\slshape n}})} returns an epimorphism from \mbox{\texttt{\mdseries\slshape digraph}} onto the complete digraph with \mbox{\texttt{\mdseries\slshape n}} vertices if one exists, else it returns \texttt{fail}. 

 See also \texttt{DigraphGreedyColouring} (\ref{DigraphGreedyColouring:for a digraph}) and 

 Note that a digraph with at least two vertices has a
2\texttt{\symbol{45}}colouring if and only if it is bipartite, see \texttt{IsBipartiteDigraph} (\ref{IsBipartiteDigraph}). 
\begin{Verbatim}[commandchars=!@|,fontsize=\small,frame=single,label=Example]
  !gapprompt@gap>| !gapinput@DigraphColouring(CompleteDigraph(5), 4);|
  fail
  !gapprompt@gap>| !gapinput@DigraphColouring(ChainDigraph(10), 1);|
  fail
  !gapprompt@gap>| !gapinput@D := ChainDigraph(10);;|
  !gapprompt@gap>| !gapinput@t := DigraphColouring(D, 2);|
  Transformation( [ 1, 2, 1, 2, 1, 2, 1, 2, 1, 2 ] )
  !gapprompt@gap>| !gapinput@IsDigraphColouring(D, t);|
  true
  !gapprompt@gap>| !gapinput@DigraphGreedyColouring(D);|
  Transformation( [ 2, 1, 2, 1, 2, 1, 2, 1, 2, 1 ] )
\end{Verbatim}
 }

 

\subsection{\textcolor{Chapter }{DigraphGreedyColouring (for a digraph and vertex order)}}
\logpage{[ 7, 3, 16 ]}\nobreak
\hyperdef{L}{X7AB7200D831013C1}{}
{\noindent\textcolor{FuncColor}{$\triangleright$\enspace\texttt{DigraphGreedyColouring({\mdseries\slshape digraph, order})\index{DigraphGreedyColouring@\texttt{DigraphGreedyColouring}!for a digraph and vertex order}
\label{DigraphGreedyColouring:for a digraph and vertex order}
}\hfill{\scriptsize (operation)}}\\
\noindent\textcolor{FuncColor}{$\triangleright$\enspace\texttt{DigraphGreedyColouring({\mdseries\slshape digraph, func})\index{DigraphGreedyColouring@\texttt{DigraphGreedyColouring}!for a digraph and vertex order function}
\label{DigraphGreedyColouring:for a digraph and vertex order function}
}\hfill{\scriptsize (operation)}}\\
\noindent\textcolor{FuncColor}{$\triangleright$\enspace\texttt{DigraphGreedyColouring({\mdseries\slshape digraph})\index{DigraphGreedyColouring@\texttt{DigraphGreedyColouring}!for a digraph}
\label{DigraphGreedyColouring:for a digraph}
}\hfill{\scriptsize (attribute)}}\\
\textbf{\indent Returns:\ }
 A transformation, or \texttt{fail}.



 A \emph{proper colouring} of a digraph is a labelling of its vertices in such a way that adjacent
vertices have different labels. Note that a digraph with loops (\texttt{DigraphHasLoops} (\ref{DigraphHasLoops})) does not have any proper colouring. 

 If \mbox{\texttt{\mdseries\slshape digraph}} is a digraph and \mbox{\texttt{\mdseries\slshape order}} is a dense list consisting of all of the vertices of \mbox{\texttt{\mdseries\slshape digraph}} (in any order), then \texttt{DigraphGreedyColouring} uses a greedy algorithm with the specified order to obtain some proper
colouring of \mbox{\texttt{\mdseries\slshape digraph}}, which may not use the minimal number of colours. 

 If \mbox{\texttt{\mdseries\slshape digraph}} is a digraph and \mbox{\texttt{\mdseries\slshape func}} is a function whose argument is a digraph, and that returns a dense list \mbox{\texttt{\mdseries\slshape order}}, then \texttt{DigraphGreedyColouring(\mbox{\texttt{\mdseries\slshape digraph}}, \mbox{\texttt{\mdseries\slshape func}})} returns \texttt{DigraphGreedyColouring(\mbox{\texttt{\mdseries\slshape digraph}}, \mbox{\texttt{\mdseries\slshape func}}(\mbox{\texttt{\mdseries\slshape digraph}}))}. 

 If the optional second argument (either a list or a function), is not
specified, then \texttt{DigraphWelshPowellOrder} (\ref{DigraphWelshPowellOrder}) is used by default. 

 See also \texttt{DigraphColouring} (\ref{DigraphColouring:for a digraph and a number of colours}). 

 
\begin{Verbatim}[commandchars=!@|,fontsize=\small,frame=single,label=Example]
  !gapprompt@gap>| !gapinput@DigraphGreedyColouring(ChainDigraph(10));|
  Transformation( [ 2, 1, 2, 1, 2, 1, 2, 1, 2, 1 ] )
  !gapprompt@gap>| !gapinput@DigraphGreedyColouring(ChainDigraph(10), [1 .. 10]);|
  Transformation( [ 1, 2, 1, 2, 1, 2, 1, 2, 1, 2 ] )
\end{Verbatim}
 }

 

\subsection{\textcolor{Chapter }{DigraphWelshPowellOrder}}
\logpage{[ 7, 3, 17 ]}\nobreak
\hyperdef{L}{X7F9CB3B27B9590DB}{}
{\noindent\textcolor{FuncColor}{$\triangleright$\enspace\texttt{DigraphWelshPowellOrder({\mdseries\slshape digraph})\index{DigraphWelshPowellOrder@\texttt{DigraphWelshPowellOrder}}
\label{DigraphWelshPowellOrder}
}\hfill{\scriptsize (attribute)}}\\
\textbf{\indent Returns:\ }
 A list of the vertices.



 \texttt{DigraphWelshPowellOrder} returns a list of all of the vertices of the digraph \mbox{\texttt{\mdseries\slshape digraph}} ordered according to the sum of the number of out\texttt{\symbol{45}} and
in\texttt{\symbol{45}}neighbours, from highest to lowest. 

 
\begin{Verbatim}[commandchars=!@|,fontsize=\small,frame=single,label=Example]
  !gapprompt@gap>| !gapinput@DigraphWelshPowellOrder(Digraph([[4], [9], [9], [],|
  !gapprompt@>| !gapinput@                                    [4, 6, 9], [1], [], [],|
  !gapprompt@>| !gapinput@                                    [4, 5], [4, 5]]));|
  [ 5, 9, 4, 1, 6, 10, 2, 3, 7, 8 ]
\end{Verbatim}
 }

 

\subsection{\textcolor{Chapter }{ChromaticNumber}}
\logpage{[ 7, 3, 18 ]}\nobreak
\hyperdef{L}{X807D49358529A37C}{}
{\noindent\textcolor{FuncColor}{$\triangleright$\enspace\texttt{ChromaticNumber({\mdseries\slshape digraph})\index{ChromaticNumber@\texttt{ChromaticNumber}}
\label{ChromaticNumber}
}\hfill{\scriptsize (attribute)}}\\
\textbf{\indent Returns:\ }
 A non\texttt{\symbol{45}}negative integer.



 A \emph{proper colouring} of a digraph is a labelling of its vertices in such a way that adjacent
vertices have different labels. Equivalently, a proper digraph colouring can
be defined to be a \texttt{DigraphEpimorphism} (\ref{DigraphEpimorphism}) from a digraph onto a complete digraph. 

 If \mbox{\texttt{\mdseries\slshape digraph}} is a digraph without loops (see \texttt{DigraphHasLoops} (\ref{DigraphHasLoops}), then \texttt{ChromaticNumber} returns the least non\texttt{\symbol{45}}negative integer \texttt{n} such that there is a proper colouring of \mbox{\texttt{\mdseries\slshape digraph}} with \texttt{n} colours. In other words, for a digraph with at least one vertex, \texttt{ChromaticNumber} returns the least number \texttt{n} such that \texttt{DigraphColouring(\mbox{\texttt{\mdseries\slshape digraph}}, n)} does not return \texttt{fail}. See \texttt{DigraphColouring} (\ref{DigraphColouring:for a digraph and a number of colours}). 

 It is possible to select the algorithm to compute the chromatic number via the
use of value options. The permitted algorithms and values to pass as options
are: 
\begin{itemize}
\item \texttt{lawler} \texttt{\symbol{45}} Lawler's Algorithm \cite{Law1976}
\item \texttt{byskov} \texttt{\symbol{45}} Byskov's Algorithm \cite{Bys2002}
\item \texttt{zykov} \texttt{\symbol{45}} Zykov's Algorithm \cite{Corneil1973}
\item \texttt{christofides} \texttt{\symbol{45}} Christofides's Algorithm \cite{Wang1974}
\end{itemize}
 
\begin{Verbatim}[commandchars=!@|,fontsize=\small,frame=single,label=Example]
  !gapprompt@gap>| !gapinput@ChromaticNumber(NullDigraph(10));|
  1
  !gapprompt@gap>| !gapinput@ChromaticNumber(CompleteDigraph(10));|
  10
  !gapprompt@gap>| !gapinput@ChromaticNumber(CompleteBipartiteDigraph(5, 5));|
  2
  !gapprompt@gap>| !gapinput@ChromaticNumber(Digraph([[], [3], [5], [2, 3], [4]]));|
  3
  !gapprompt@gap>| !gapinput@ChromaticNumber(NullDigraph(0));|
  0
  !gapprompt@gap>| !gapinput@D := PetersenGraph(IsMutableDigraph);|
  <mutable digraph with 10 vertices, 30 edges>
  !gapprompt@gap>| !gapinput@ChromaticNumber(D);|
  3
  !gapprompt@gap>| !gapinput@ChromaticNumber(CompleteDigraph(10) : lawler);|
  10
  !gapprompt@gap>| !gapinput@ChromaticNumber(CompleteDigraph(10) : byskov);|
  10
\end{Verbatim}
 }

 

\subsection{\textcolor{Chapter }{DigraphCore}}
\logpage{[ 7, 3, 19 ]}\nobreak
\hyperdef{L}{X863487EF7D64EF56}{}
{\noindent\textcolor{FuncColor}{$\triangleright$\enspace\texttt{DigraphCore({\mdseries\slshape D})\index{DigraphCore@\texttt{DigraphCore}}
\label{DigraphCore}
}\hfill{\scriptsize (attribute)}}\\
\textbf{\indent Returns:\ }
A list of positive integers.



 If \mbox{\texttt{\mdseries\slshape D}} is a digraph, then \texttt{DigraphCore} returns a list of vertices corresponding to the \texttt{core} of \mbox{\texttt{\mdseries\slshape D}}. In particular, the subdigraph of \mbox{\texttt{\mdseries\slshape D}} induced by this list is isomorphic to the core of \mbox{\texttt{\mdseries\slshape D}}.

 The \emph{core} of a digraph \texttt{D} is the minimal subdigraph \mbox{\texttt{\mdseries\slshape C}} of \texttt{D} which is a homomorphic image of \texttt{D}. The core of a digraph is unique up to isomorphism. 
\begin{Verbatim}[commandchars=!@|,fontsize=\small,frame=single,label=Example]
  !gapprompt@gap>| !gapinput@D := DigraphSymmetricClosure(CycleDigraph(8));|
  <immutable symmetric digraph with 8 vertices, 16 edges>
  !gapprompt@gap>| !gapinput@DigraphCore(D);|
  [ 1, 2 ]
  !gapprompt@gap>| !gapinput@D := PetersenGraph();|
  <immutable digraph with 10 vertices, 30 edges>
  !gapprompt@gap>| !gapinput@DigraphCore(D);|
  [ 1 .. 10 ]
  !gapprompt@gap>| !gapinput@D := Digraph(IsMutableDigraph, [[3], [3], [4], [5], [2]]);|
  <mutable digraph with 5 vertices, 5 edges>
  !gapprompt@gap>| !gapinput@DigraphCore(D);|
  [ 2, 3, 4, 5 ]
\end{Verbatim}
 }

 

\subsection{\textcolor{Chapter }{LatticeDigraphEmbedding}}
\logpage{[ 7, 3, 20 ]}\nobreak
\hyperdef{L}{X7EEA340C81FCA52D}{}
{\noindent\textcolor{FuncColor}{$\triangleright$\enspace\texttt{LatticeDigraphEmbedding({\mdseries\slshape L1, L2})\index{LatticeDigraphEmbedding@\texttt{LatticeDigraphEmbedding}}
\label{LatticeDigraphEmbedding}
}\hfill{\scriptsize (operation)}}\\
\textbf{\indent Returns:\ }
A transformation, or \texttt{fail}.



 If \mbox{\texttt{\mdseries\slshape L1}} and \mbox{\texttt{\mdseries\slshape L2}} are lattice digraphs (\texttt{IsLatticeDigraph} (\ref{IsLatticeDigraph}) returns \texttt{true}, then \texttt{LatticeDigraphEmbedding} returns a single \emph{injective} \texttt{DigraphHomomorphism} (\ref{DigraphHomomorphism}) between \mbox{\texttt{\mdseries\slshape L1}} and \mbox{\texttt{\mdseries\slshape L2}}, with the property that it is a \emph{lattice homomorphism}. If no such homomorphism exists, \texttt{fail} is returned.

 A \emph{lattice homomorphism} is a digraph homomorphism which respects meets and joins of every pair of
vertices. Note that every injective lattice homomorphism \texttt{map} is an embedding, in the sense that the inverse of \texttt{map} is a lattice homomorphism also. 
\begin{Verbatim}[commandchars=!@|,fontsize=\small,frame=single,label=Example]
  !gapprompt@gap>| !gapinput@D := DigraphReflexiveTransitiveClosure(ChainDigraph(5));|
  <immutable preorder digraph with 5 vertices, 15 edges>
  !gapprompt@gap>| !gapinput@L1 := DigraphReflexiveTransitiveClosure(ChainDigraph(5));|
  <immutable preorder digraph with 5 vertices, 15 edges>
  !gapprompt@gap>| !gapinput@L2 := DigraphReflexiveTransitiveClosure(ChainDigraph(6));|
  <immutable preorder digraph with 6 vertices, 21 edges>
  !gapprompt@gap>| !gapinput@LatticeDigraphEmbedding(L1, L2);|
  IdentityTransformation
  !gapprompt@gap>| !gapinput@LatticeDigraphEmbedding(L2, L1);|
  fail
\end{Verbatim}
 }

 

\subsection{\textcolor{Chapter }{IsLatticeHomomorphism (for digraphs and a permutation or transformation)}}
\logpage{[ 7, 3, 21 ]}\nobreak
\hyperdef{L}{X7E183D4979D324EF}{}
{\noindent\textcolor{FuncColor}{$\triangleright$\enspace\texttt{IsLatticeHomomorphism({\mdseries\slshape L1, L2, map})\index{IsLatticeHomomorphism@\texttt{IsLatticeHomomorphism}!for digraphs and a permutation or transformation}
\label{IsLatticeHomomorphism:for digraphs and a permutation or transformation}
}\hfill{\scriptsize (operation)}}\\
\noindent\textcolor{FuncColor}{$\triangleright$\enspace\texttt{IsLatticeEpimorphism({\mdseries\slshape L1, L2, map})\index{IsLatticeEpimorphism@\texttt{IsLatticeEpimorphism}!for digraphs and a permutation or transformation}
\label{IsLatticeEpimorphism:for digraphs and a permutation or transformation}
}\hfill{\scriptsize (operation)}}\\
\noindent\textcolor{FuncColor}{$\triangleright$\enspace\texttt{IsLatticeEmbedding({\mdseries\slshape L1, L2, map})\index{IsLatticeEmbedding@\texttt{IsLatticeEmbedding}!for digraphs and a permutation or transformation}
\label{IsLatticeEmbedding:for digraphs and a permutation or transformation}
}\hfill{\scriptsize (operation)}}\\
\noindent\textcolor{FuncColor}{$\triangleright$\enspace\texttt{IsLatticeMonomorphism({\mdseries\slshape L1, L2, map})\index{IsLatticeMonomorphism@\texttt{IsLatticeMonomorphism}!for digraphs and a permutation or transformation}
\label{IsLatticeMonomorphism:for digraphs and a permutation or transformation}
}\hfill{\scriptsize (operation)}}\\
\noindent\textcolor{FuncColor}{$\triangleright$\enspace\texttt{IsLatticeEndomorphism({\mdseries\slshape L, map})\index{IsLatticeEndomorphism@\texttt{IsLatticeEndomorphism}!for digraphs and a permutation or transformation}
\label{IsLatticeEndomorphism:for digraphs and a permutation or transformation}
}\hfill{\scriptsize (operation)}}\\
\textbf{\indent Returns:\ }
\texttt{true} or \texttt{false}.



 Each of the function described in this section (except \texttt{IsLatticeEndomorphism}) takes a pair of digraphs \mbox{\texttt{\mdseries\slshape L1}} and \mbox{\texttt{\mdseries\slshape L2}}, and a transformation \mbox{\texttt{\mdseries\slshape map}}, returning \texttt{true} if \mbox{\texttt{\mdseries\slshape map}} is a \emph{lattice homomorphism} from \mbox{\texttt{\mdseries\slshape L1}} to \mbox{\texttt{\mdseries\slshape L2}}, and \texttt{false} otherwise. If \mbox{\texttt{\mdseries\slshape L1}} or \mbox{\texttt{\mdseries\slshape L2}} is not a lattice, then \texttt{false} is returned.

 A transformation or permutation \mbox{\texttt{\mdseries\slshape map}} is a \emph{lattice homomorphism} if \mbox{\texttt{\mdseries\slshape map}} respects meets and joins of every pair of vertices, and \mbox{\texttt{\mdseries\slshape map}} fixes every \texttt{i} which is not a vertex of \mbox{\texttt{\mdseries\slshape L1}}. 

 \texttt{IsLatticeHomomorphism} returns \texttt{true} if the permutation or transformation \mbox{\texttt{\mdseries\slshape map}} is a lattice homomorphism from the lattice digraph \mbox{\texttt{\mdseries\slshape L1}} to the lattice digraph \mbox{\texttt{\mdseries\slshape L2}}. 

 \texttt{IsLatticeEpimorphism} returns \texttt{true} if the permutation or transformation \mbox{\texttt{\mdseries\slshape map}} is a surjective lattice homomorphism from the lattice digraph \mbox{\texttt{\mdseries\slshape L1}} to the lattice digraph \mbox{\texttt{\mdseries\slshape L2}}. 

 \texttt{IsLatticeEmbedding} returns \texttt{true} if the permutation or transformation \mbox{\texttt{\mdseries\slshape map}} is an injective lattice homomorphism from the lattice digraph \mbox{\texttt{\mdseries\slshape L1}} to the lattice digraph \mbox{\texttt{\mdseries\slshape L2}}. The function \texttt{IsLatticeMonomorphism} is a synonym of \texttt{IsLatticeEmbedding}. 

 \texttt{IsLatticeEndomorphism} returns \texttt{true} if the permutation or transformation \mbox{\texttt{\mdseries\slshape map}} is an lattice endomorphism of the lattice digraph \mbox{\texttt{\mdseries\slshape L}}. 
\begin{Verbatim}[commandchars=!@|,fontsize=\small,frame=single,label=Example]
  !gapprompt@gap>| !gapinput@G := Digraph([[2, 4], [3, 7], [6], [5, 7], [6], [], [6]]);|
  <immutable digraph with 7 vertices, 9 edges>
  !gapprompt@gap>| !gapinput@D := DigraphRemoveVertex(G, 7);|
  <immutable digraph with 6 vertices, 6 edges>
  !gapprompt@gap>| !gapinput@G := DigraphReflexiveTransitiveClosure(G);|
  <immutable preorder digraph with 7 vertices, 22 edges>
  !gapprompt@gap>| !gapinput@D := DigraphReflexiveTransitiveClosure(D);|
  <immutable preorder digraph with 6 vertices, 17 edges>
  !gapprompt@gap>| !gapinput@IsDigraphEmbedding(D, G, IdentityTransformation);|
  true
  !gapprompt@gap>| !gapinput@IsLatticeHomomorphism(D, G, IdentityTransformation);|
  false
  !gapprompt@gap>| !gapinput@D := Digraph([[2, 3], [4], [4], []]);|
  <immutable digraph with 4 vertices, 4 edges>
  !gapprompt@gap>| !gapinput@G := Digraph([[2, 3], [4], [4], [5], []]);|
  <immutable digraph with 5 vertices, 5 edges>
  !gapprompt@gap>| !gapinput@D := DigraphReflexiveTransitiveClosure(D);|
  <immutable preorder digraph with 4 vertices, 9 edges>
  !gapprompt@gap>| !gapinput@G := DigraphReflexiveTransitiveClosure(G);|
  <immutable preorder digraph with 5 vertices, 14 edges>
  !gapprompt@gap>| !gapinput@IsLatticeEmbedding(D, G, IdentityTransformation);|
  true
  !gapprompt@gap>| !gapinput@IsLatticeMonomorphism(D, G, IdentityTransformation);|
  true
  !gapprompt@gap>| !gapinput@f := Transformation([1, 2, 3, 4, 4]);|
  Transformation( [ 1, 2, 3, 4, 4 ] )
  !gapprompt@gap>| !gapinput@IsLatticeEpimorphism(G, D, f);|
  true
  !gapprompt@gap>| !gapinput@IsLatticeEndomorphism(D, (2, 3));|
  true
\end{Verbatim}
 }

 }

 }

 
\chapter{\textcolor{Chapter }{Finding cliques and independent sets}}\label{Finding cliques and independent sets}
\logpage{[ 8, 0, 0 ]}
\hyperdef{L}{X7A7E0B957932DF44}{}
{
  In \textsf{Digraphs}, a \emph{clique} of a digraph is a set of mutually adjacent vertices of the digraph, and an \emph{independent set} is a set of mutually non\texttt{\symbol{45}}adjacent vertices of the digraph.
A \emph{maximal clique} is a clique which is not properly contained in another clique, and a \emph{maximal independent set} is an independent set which is not properly contained in another independent
set. Using this definition in \textsf{Digraphs}, cliques and independent sets are both permitted, but not required, to
contain vertices at which there is a loop. Another name for a clique is a \emph{complete subgraph}. 

 \textsf{Digraphs} provides extensive functionality for computing cliques and independent sets of
a digraph, whether maximal or not. The fundamental algorithm used in most of
the methods in \textsf{Digraphs} to calculate cliques and independent sets is a version of the
Bron\texttt{\symbol{45}}Kerbosch algorithm. Calculating the cliques and
independent sets of a digraph is a well\texttt{\symbol{45}}known and hard
problem, and searching for cliques or independent sets in a digraph can be
very lengthy, even for a digraph with a small number of vertices. \textsf{Digraphs} uses several strategies to increase the performance of these calculations. 

 From the definition of cliques and independent sets, it follows that the
presence of loops and multiple edges in a digraph is irrelevant to the
existence of cliques and independent sets in the digraph. See \texttt{DigraphHasLoops} (\ref{DigraphHasLoops}) and \texttt{IsMultiDigraph} (\ref{IsMultiDigraph}) for more information about these properties. Therefore given a digraph \mbox{\texttt{\mdseries\slshape digraph}}, the cliques and independent sets of \mbox{\texttt{\mdseries\slshape digraph}} are equal to the cliques and independent sets of the digraph: 
\begin{itemize}
\item  \texttt{DigraphRemoveLoops(DigraphRemoveAllMultipleEdges(}\mbox{\texttt{\mdseries\slshape digraph}}\texttt{))}. 
\end{itemize}
 See \texttt{DigraphRemoveLoops} (\ref{DigraphRemoveLoops}) and \texttt{DigraphRemoveAllMultipleEdges} (\ref{DigraphRemoveAllMultipleEdges}) for more information about these attributes. Furthermore, the cliques of this
digraph are equal to the cliques of the digraph formed by removing any edge \texttt{[u,v]} for which the corresponding reverse edge \texttt{[v,u]} is not present. Therefore, overall, the cliques of \mbox{\texttt{\mdseries\slshape digraph}} are equal to the cliques of the symmetric digraph: 
\begin{itemize}
\item  \texttt{MaximalSymmetricSubdigraphWithoutLoops(}\mbox{\texttt{\mdseries\slshape digraph}}\texttt{)}. 
\end{itemize}
 See \texttt{MaximalSymmetricSubdigraphWithoutLoops} (\ref{MaximalSymmetricSubdigraphWithoutLoops}) for more information about this. The \texttt{AutomorphismGroup} (\ref{AutomorphismGroup:for a digraph}) of this symmetric digraph contains the automorphism group of \mbox{\texttt{\mdseries\slshape digraph}} as a subgroup. By performing the search for maximal cliques with the help of
this larger automorphism group to reduce the search space, the computation
time may be reduced. The functions and attributes which return representatives
of cliques of \mbox{\texttt{\mdseries\slshape digraph}} will return orbit representatives of cliques under the action of the
automorphism group of the \emph{maximal symmetric subdigraph without loops} on sets of vertices.

 The independent sets of a digraph are equal to the independent sets of the \texttt{DigraphSymmetricClosure} (\ref{DigraphSymmetricClosure}). Therefore, overall, the independent sets of \mbox{\texttt{\mdseries\slshape digraph}} are equal to the independent sets of the symmetric digraph: 
\begin{itemize}
\item  \texttt{DigraphSymmetricClosure(DigraphRemoveLoops(DigraphRemoveAllMultipleEdges( }\mbox{\texttt{\mdseries\slshape digraph}}\texttt{)))}. 
\end{itemize}
 Again, the automorphism group of this symmetric digraph contains the
automorphism group of \mbox{\texttt{\mdseries\slshape digraph}}. Since a search for independent sets is equal to a search for cliques in the \texttt{DigraphDual} (\ref{DigraphDual}), the methods used in \textsf{Digraphs} usually transform a search for independent sets into a search for cliques, as
described above. The functions and attributes which return representatives of
independent sets of \mbox{\texttt{\mdseries\slshape digraph}} will return orbit representatives of independent sets under the action of the
automorphism group of the \emph{symmetric closure} of the digraph formed by removing loops and multiple edges.

 Please note that in \textsf{Digraphs}, cliques and independent sets are not required to be maximal. Some authors
use the word clique to mean \emph{maximal} clique, and some authors use the phrase independent set to mean \emph{maximal} independent set, but please note that \textsf{Digraphs} does not use this definition. 

 
\section{\textcolor{Chapter }{Finding cliques}}\logpage{[ 8, 1, 0 ]}
\hyperdef{L}{X788474E38446039B}{}
{
 

\subsection{\textcolor{Chapter }{IsClique}}
\logpage{[ 8, 1, 1 ]}\nobreak
\hyperdef{L}{X80670E3E7C36183F}{}
{\noindent\textcolor{FuncColor}{$\triangleright$\enspace\texttt{IsClique({\mdseries\slshape digraph, l})\index{IsClique@\texttt{IsClique}}
\label{IsClique}
}\hfill{\scriptsize (operation)}}\\
\noindent\textcolor{FuncColor}{$\triangleright$\enspace\texttt{IsMaximalClique({\mdseries\slshape digraph, l})\index{IsMaximalClique@\texttt{IsMaximalClique}}
\label{IsMaximalClique}
}\hfill{\scriptsize (operation)}}\\
\textbf{\indent Returns:\ }
\texttt{true} or \texttt{false}.



 If \mbox{\texttt{\mdseries\slshape digraph}} is a digraph and \mbox{\texttt{\mdseries\slshape l}} is a duplicate\texttt{\symbol{45}}free list of vertices of \mbox{\texttt{\mdseries\slshape digraph}}, then \texttt{IsClique(}\mbox{\texttt{\mdseries\slshape digraph}}\texttt{,}\mbox{\texttt{\mdseries\slshape l}}\texttt{)} returns \texttt{true} if \mbox{\texttt{\mdseries\slshape l}} is a \emph{clique} of \mbox{\texttt{\mdseries\slshape digraph}} and \texttt{false} if it is not. Similarly, \texttt{IsMaximalClique(}\mbox{\texttt{\mdseries\slshape digraph}}\texttt{,}\mbox{\texttt{\mdseries\slshape l}}\texttt{)} returns \texttt{true} if \mbox{\texttt{\mdseries\slshape l}} is a \emph{maximal clique} of \mbox{\texttt{\mdseries\slshape digraph}} and \texttt{false} if it is not. 

 A \emph{clique} of a digraph is a set of mutually adjacent vertices of the digraph. A \emph{maximal clique} is a clique that is not properly contained in another clique. A clique is
permitted, but not required, to contain vertices at which there is a loop. 
\begin{Verbatim}[commandchars=!@|,fontsize=\small,frame=single,label=Example]
  !gapprompt@gap>| !gapinput@D := CompleteDigraph(4);;|
  !gapprompt@gap>| !gapinput@IsClique(D, [1, 3, 2]);|
  true
  !gapprompt@gap>| !gapinput@IsMaximalClique(D, [1, 3, 2]);|
  false
  !gapprompt@gap>| !gapinput@IsMaximalClique(D, DigraphVertices(D));|
  true
  !gapprompt@gap>| !gapinput@D := Digraph([[1, 2, 4, 4], [1, 3, 4], [2, 1], [1, 2]]);|
  <immutable multidigraph with 4 vertices, 11 edges>
  !gapprompt@gap>| !gapinput@IsClique(D, [2, 3, 4]);|
  false
  !gapprompt@gap>| !gapinput@IsMaximalClique(D, [1, 2, 4]);|
  true
  !gapprompt@gap>| !gapinput@D := CompleteDigraph(IsMutableDigraph, 4);;|
  !gapprompt@gap>| !gapinput@IsClique(D, [1, 3, 2]);|
  true
\end{Verbatim}
 }

 

\subsection{\textcolor{Chapter }{CliquesFinder}}
\logpage{[ 8, 1, 2 ]}\nobreak
\hyperdef{L}{X7A3BC9FF8169ABCD}{}
{\noindent\textcolor{FuncColor}{$\triangleright$\enspace\texttt{CliquesFinder({\mdseries\slshape digraph, hook, user{\textunderscore}param, limit, include, exclude, max, size, reps})\index{CliquesFinder@\texttt{CliquesFinder}}
\label{CliquesFinder}
}\hfill{\scriptsize (function)}}\\
\textbf{\indent Returns:\ }
The argument \mbox{\texttt{\mdseries\slshape user{\textunderscore}param}}.



 This function finds cliques of the digraph \mbox{\texttt{\mdseries\slshape digraph}} subject to the conditions imposed by the other arguments as described below.
Note that a clique is represented by the immutable list of the vertices that
it contains. 

 Let \texttt{G} denote the automorphism group of the maximal symmetric subdigraph of \mbox{\texttt{\mdseries\slshape digraph}} without loops (see \texttt{AutomorphismGroup} (\ref{AutomorphismGroup:for a digraph}) and \texttt{MaximalSymmetricSubdigraphWithoutLoops} (\ref{MaximalSymmetricSubdigraphWithoutLoops})). 
\begin{description}
\item[{\mbox{\texttt{\mdseries\slshape hook}}}]  This argument should be a function or \texttt{fail}.

 If \mbox{\texttt{\mdseries\slshape hook}} is a function, then it should have two arguments \mbox{\texttt{\mdseries\slshape user{\textunderscore}param}} (see below) and a clique \texttt{c}. The function \texttt{\mbox{\texttt{\mdseries\slshape hook}}(\mbox{\texttt{\mdseries\slshape user{\textunderscore}param}}, c)} is called every time a new clique \texttt{c} is found by \texttt{CliquesFinder}.

 If \mbox{\texttt{\mdseries\slshape hook}} is \texttt{fail}, then a default function is used that simply adds every new clique found by \texttt{CliquesFinder} to \mbox{\texttt{\mdseries\slshape user{\textunderscore}param}}, which must be a list in this case. 
\item[{\mbox{\texttt{\mdseries\slshape user{\textunderscore}param}}}]  If \mbox{\texttt{\mdseries\slshape hook}} is a function, then \mbox{\texttt{\mdseries\slshape user{\textunderscore}param}} can be any \textsf{GAP} object. The object \mbox{\texttt{\mdseries\slshape user{\textunderscore}param}} is used as the first argument for the function \mbox{\texttt{\mdseries\slshape hook}}. For example, \mbox{\texttt{\mdseries\slshape user{\textunderscore}param}} might be a list, and \texttt{\mbox{\texttt{\mdseries\slshape hook}}(\mbox{\texttt{\mdseries\slshape user{\textunderscore}param}}, c)} might add the size of the clique \texttt{c} to the list \mbox{\texttt{\mdseries\slshape user{\textunderscore}param}}. 

 If the value of \mbox{\texttt{\mdseries\slshape hook}} is \texttt{fail}, then the value of \mbox{\texttt{\mdseries\slshape user{\textunderscore}param}} must be a list. 
\item[{\mbox{\texttt{\mdseries\slshape limit}}}]  This argument should be a positive integer or \texttt{infinity}. \texttt{CliquesFinder} will return after it has found \mbox{\texttt{\mdseries\slshape limit}} cliques or the search is complete. 
\item[{\mbox{\texttt{\mdseries\slshape include}} and \mbox{\texttt{\mdseries\slshape exclude}}}]  These arguments should each be a (possibly empty)
duplicate\texttt{\symbol{45}}free list of vertices of \mbox{\texttt{\mdseries\slshape digraph}} (i.e. positive integers less than the number of vertices of \mbox{\texttt{\mdseries\slshape digraph}}). 

 \texttt{CliquesFinder} will only look for cliques containing all of the vertices in \mbox{\texttt{\mdseries\slshape include}} and containing none of the vertices in \mbox{\texttt{\mdseries\slshape exclude}}. 

 Note that the search may be much more efficient if each of these lists is
invariant under the action of \texttt{G} on sets of vertices. 
\item[{\mbox{\texttt{\mdseries\slshape max}}}]  This argument should be \texttt{true} or \texttt{false}. If \mbox{\texttt{\mdseries\slshape max}} is true then \texttt{CliquesFinder} will only search for \emph{maximal} cliques. If \texttt{max} is \texttt{false} then non\texttt{\symbol{45}}maximal cliques may be found. 
\item[{\mbox{\texttt{\mdseries\slshape size}}}]  This argument should be \texttt{fail} or a positive integer. If \mbox{\texttt{\mdseries\slshape size}} is a positive integer then \texttt{CliquesFinder} will only search for cliques that contain precisely \mbox{\texttt{\mdseries\slshape size}} vertices. If \mbox{\texttt{\mdseries\slshape size}} is \texttt{fail} then cliques of any size may be found. 
\item[{\mbox{\texttt{\mdseries\slshape reps}}}]  This argument should be \texttt{true} or \texttt{false}.

 If \mbox{\texttt{\mdseries\slshape reps}} is \texttt{true} then the arguments \mbox{\texttt{\mdseries\slshape include}} and \mbox{\texttt{\mdseries\slshape exclude}} are each required to be invariant under the action of \texttt{G} on sets of vertices. In this case, \texttt{CliquesFinder} will find representatives of the orbits of the desired cliques under the
action of \texttt{G}, \emph{although representatives may be returned that are in the same orbit}. If \mbox{\texttt{\mdseries\slshape reps}} is false then \texttt{CliquesFinder} will not take this into consideration.

 For a digraph such that \texttt{G} is non\texttt{\symbol{45}}trivial, the search for clique representatives can
be much more efficient than the search for all cliques. 
\end{description}
 This function uses a version of the Bron\texttt{\symbol{45}}Kerbosch
algorithm. 
\begin{Verbatim}[commandchars=!@|,fontsize=\small,frame=single,label=Example]
  !gapprompt@gap>| !gapinput@D := CompleteDigraph(5);|
  <immutable complete digraph with 5 vertices>
  !gapprompt@gap>| !gapinput@user_param := [];;|
  !gapprompt@gap>| !gapinput@f := function(a, b)  # Calculate size of clique|
  !gapprompt@>| !gapinput@  AddSet(user_param, Size(b));|
  !gapprompt@>| !gapinput@end;;|
  !gapprompt@gap>| !gapinput@CliquesFinder(D, f, user_param, infinity, [], [], false, fail,|
  !gapprompt@>| !gapinput@                 true);|
  [ 1, 2, 3, 4, 5 ]
  !gapprompt@gap>| !gapinput@CliquesFinder(D, fail, [], 5, [2, 4], [3], false, fail, false);|
  [ [ 2, 4 ], [ 1, 2, 4 ], [ 2, 4, 5 ], [ 1, 2, 4, 5 ] ]
  !gapprompt@gap>| !gapinput@CliquesFinder(D, fail, [], 2, [2, 4], [3], false, fail, false);|
  [ [ 2, 4 ], [ 1, 2, 4 ] ]
  !gapprompt@gap>| !gapinput@CliquesFinder(D, fail, [], infinity, [], [], true, 5, false);|
  [ [ 1, 2, 3, 4, 5 ] ]
  !gapprompt@gap>| !gapinput@CliquesFinder(D, fail, [], infinity, [1, 3], [], false, 3, false);|
  [ [ 1, 2, 3 ], [ 1, 3, 4 ], [ 1, 3, 5 ] ]
  !gapprompt@gap>| !gapinput@CliquesFinder(D, fail, [], infinity, [1, 3], [], true, 3, false);|
  [  ]
  !gapprompt@gap>| !gapinput@D := CompleteDigraph(IsMutableDigraph, 5);|
  <mutable digraph with 5 vertices, 20 edges>
  !gapprompt@gap>| !gapinput@user_param := [];;|
  !gapprompt@gap>| !gapinput@f := function(a, b)  # Calculate size of clique|
  !gapprompt@>| !gapinput@  AddSet(user_param, Size(b));|
  !gapprompt@>| !gapinput@end;;|
  !gapprompt@gap>| !gapinput@CliquesFinder(D, f, user_param, infinity, [], [], false, fail,|
  !gapprompt@>| !gapinput@                 true);|
  [ 1, 2, 3, 4, 5 ]
  !gapprompt@gap>| !gapinput@CliquesFinder(D, fail, [], 5, [2, 4], [3], false, fail, false);|
  [ [ 2, 4 ], [ 1, 2, 4 ], [ 2, 4, 5 ], [ 1, 2, 4, 5 ] ]
  !gapprompt@gap>| !gapinput@CliquesFinder(D, fail, [], 2, [2, 4], [3], false, fail, false);|
  [ [ 2, 4 ], [ 1, 2, 4 ] ]
  !gapprompt@gap>| !gapinput@CliquesFinder(D, fail, [], infinity, [], [], true, 5, false);|
  [ [ 1, 2, 3, 4, 5 ] ]
  !gapprompt@gap>| !gapinput@CliquesFinder(D, fail, [], infinity, [1, 3], [], false, 3, false);|
  [ [ 1, 2, 3 ], [ 1, 3, 4 ], [ 1, 3, 5 ] ]
  !gapprompt@gap>| !gapinput@CliquesFinder(D, fail, [], infinity, [1, 3], [], true, 3, false);|
  [  ]
\end{Verbatim}
 }

 

\subsection{\textcolor{Chapter }{DigraphClique}}
\logpage{[ 8, 1, 3 ]}\nobreak
\hyperdef{L}{X789854AC7B466402}{}
{\noindent\textcolor{FuncColor}{$\triangleright$\enspace\texttt{DigraphClique({\mdseries\slshape digraph[, include[, exclude[, size]]]})\index{DigraphClique@\texttt{DigraphClique}}
\label{DigraphClique}
}\hfill{\scriptsize (function)}}\\
\noindent\textcolor{FuncColor}{$\triangleright$\enspace\texttt{DigraphMaximalClique({\mdseries\slshape digraph[, include[, exclude[, size]]]})\index{DigraphMaximalClique@\texttt{DigraphMaximalClique}}
\label{DigraphMaximalClique}
}\hfill{\scriptsize (function)}}\\
\textbf{\indent Returns:\ }
An immutable list of positive integers, or \texttt{fail}.



 If \mbox{\texttt{\mdseries\slshape digraph}} is a digraph, then these functions returns a clique of \mbox{\texttt{\mdseries\slshape digraph}} if one exists that satisfies the arguments, else it returns \texttt{fail}. A clique is defined by the set of vertices that it contains; see \texttt{IsClique} (\ref{IsClique}) and \texttt{IsMaximalClique} (\ref{IsMaximalClique}).

 The optional arguments \mbox{\texttt{\mdseries\slshape include}} and \mbox{\texttt{\mdseries\slshape exclude}} must each be a (possibly empty) duplicate\texttt{\symbol{45}}free list of
vertices of \mbox{\texttt{\mdseries\slshape digraph}}, and the optional argument \mbox{\texttt{\mdseries\slshape size}} must be a positive integer. By default, \mbox{\texttt{\mdseries\slshape include}} and \mbox{\texttt{\mdseries\slshape exclude}} are empty. These functions will search for a clique of \mbox{\texttt{\mdseries\slshape digraph}} that includes the vertices of \mbox{\texttt{\mdseries\slshape include}} but does not include any vertices in \mbox{\texttt{\mdseries\slshape exclude}}; if the argument \mbox{\texttt{\mdseries\slshape size}} is supplied, then additionally the clique will be required to contain
precisely \mbox{\texttt{\mdseries\slshape size}} vertices.

 If \mbox{\texttt{\mdseries\slshape include}} is not a clique, then these functions return \texttt{fail}. Otherwise, the functions behave in the following way, depending on the
number of arguments: 
\begin{description}
\item[{One or two arguments}]  If one or two arguments are supplied, then \texttt{DigraphClique} and \texttt{DigraphMaximalClique} greedily enlarge the clique \mbox{\texttt{\mdseries\slshape include}} until it cannot continue. The result is guaranteed to be a maximal clique.
This may or may not return an answer more quickly than using \texttt{DigraphMaximalCliques} (\ref{DigraphMaximalCliques}) with a limit of 1. 
\item[{Three arguments}]  If three arguments are supplied, then \texttt{DigraphClique} greedily enlarges the clique \mbox{\texttt{\mdseries\slshape include}} until it cannot continue, although this clique may not be maximal.

 Given three arguments, \texttt{DigraphMaximalClique} returns the maximal clique returned by \texttt{DigraphMaximalCliques(}\mbox{\texttt{\mdseries\slshape digraph}}\texttt{, }\mbox{\texttt{\mdseries\slshape include}}\texttt{, }\mbox{\texttt{\mdseries\slshape exclude}}\texttt{, 1)} if one exists, else \texttt{fail}. 
\item[{Four arguments}]  If four arguments are supplied, then \texttt{DigraphClique} returns the clique returned by \texttt{DigraphCliques(}\mbox{\texttt{\mdseries\slshape digraph}}\texttt{, }\mbox{\texttt{\mdseries\slshape include}}\texttt{, }\mbox{\texttt{\mdseries\slshape exclude}}\texttt{, 1, }\mbox{\texttt{\mdseries\slshape size}}\texttt{)} if one exists, else \texttt{fail}. This clique may not be maximal.

 Given four arguments, \texttt{DigraphMaximalClique} returns the maximal clique returned by \texttt{DigraphMaximalCliques(}\mbox{\texttt{\mdseries\slshape digraph}}\texttt{, }\mbox{\texttt{\mdseries\slshape include}}\texttt{, }\mbox{\texttt{\mdseries\slshape exclude}}\texttt{, 1, }\mbox{\texttt{\mdseries\slshape size}}\texttt{)} if one exists, else \texttt{fail}. 
\end{description}
 
\begin{Verbatim}[commandchars=!@|,fontsize=\small,frame=single,label=Example]
  !gapprompt@gap>| !gapinput@D := Digraph([[2, 3, 4], [1, 3], [1, 2], [1, 5], []]);|
  <immutable digraph with 5 vertices, 9 edges>
  !gapprompt@gap>| !gapinput@IsSymmetricDigraph(D);|
  false
  !gapprompt@gap>| !gapinput@DigraphClique(D);|
  [ 5 ]
  !gapprompt@gap>| !gapinput@DigraphMaximalClique(D);|
  [ 5 ]
  !gapprompt@gap>| !gapinput@DigraphClique(D, [1, 2]);|
  [ 1, 2, 3 ]
  !gapprompt@gap>| !gapinput@DigraphMaximalClique(D, [1, 3]);|
  [ 1, 3, 2 ]
  !gapprompt@gap>| !gapinput@DigraphClique(D, [1], [2]);|
  [ 1, 4 ]
  !gapprompt@gap>| !gapinput@DigraphMaximalClique(D, [1], [3, 4]);|
  fail
  !gapprompt@gap>| !gapinput@DigraphClique(D, [1, 5]);|
  fail
  !gapprompt@gap>| !gapinput@DigraphClique(D, [], [], 2);|
  [ 1, 2 ]
  !gapprompt@gap>| !gapinput@D := Digraph(IsMutableDigraph,|
  !gapprompt@>| !gapinput@                [[2, 3, 4], [1, 3], [1, 2], [1, 5], []]);|
  <mutable digraph with 5 vertices, 9 edges>
  !gapprompt@gap>| !gapinput@IsSymmetricDigraph(D);|
  false
  !gapprompt@gap>| !gapinput@DigraphClique(D);|
  [ 5 ]
\end{Verbatim}
 }

 

\subsection{\textcolor{Chapter }{DigraphMaximalCliques}}
\logpage{[ 8, 1, 4 ]}\nobreak
\hyperdef{L}{X84EA2F9482B8D4AF}{}
{\noindent\textcolor{FuncColor}{$\triangleright$\enspace\texttt{DigraphMaximalCliques({\mdseries\slshape digraph[, include[, exclude[, limit[, size]]]]})\index{DigraphMaximalCliques@\texttt{DigraphMaximalCliques}}
\label{DigraphMaximalCliques}
}\hfill{\scriptsize (function)}}\\
\noindent\textcolor{FuncColor}{$\triangleright$\enspace\texttt{DigraphMaximalCliquesReps({\mdseries\slshape digraph[, include[, exclude[, limit[, size]]]]})\index{DigraphMaximalCliquesReps@\texttt{DigraphMaximalCliquesReps}}
\label{DigraphMaximalCliquesReps}
}\hfill{\scriptsize (function)}}\\
\noindent\textcolor{FuncColor}{$\triangleright$\enspace\texttt{DigraphCliques({\mdseries\slshape digraph[, include[, exclude[, limit[, size]]]]})\index{DigraphCliques@\texttt{DigraphCliques}}
\label{DigraphCliques}
}\hfill{\scriptsize (function)}}\\
\noindent\textcolor{FuncColor}{$\triangleright$\enspace\texttt{DigraphCliquesReps({\mdseries\slshape digraph[, include[, exclude[, limit[, size]]]]})\index{DigraphCliquesReps@\texttt{DigraphCliquesReps}}
\label{DigraphCliquesReps}
}\hfill{\scriptsize (function)}}\\
\noindent\textcolor{FuncColor}{$\triangleright$\enspace\texttt{DigraphMaximalCliquesAttr({\mdseries\slshape digraph})\index{DigraphMaximalCliquesAttr@\texttt{DigraphMaximalCliquesAttr}}
\label{DigraphMaximalCliquesAttr}
}\hfill{\scriptsize (attribute)}}\\
\noindent\textcolor{FuncColor}{$\triangleright$\enspace\texttt{DigraphMaximalCliquesRepsAttr({\mdseries\slshape digraph})\index{DigraphMaximalCliquesRepsAttr@\texttt{DigraphMaximalCliquesRepsAttr}}
\label{DigraphMaximalCliquesRepsAttr}
}\hfill{\scriptsize (attribute)}}\\
\textbf{\indent Returns:\ }
An immutable list of lists of positive integers.



 If \mbox{\texttt{\mdseries\slshape digraph}} is digraph, then these functions and attributes use \texttt{CliquesFinder} (\ref{CliquesFinder}) to return cliques of \mbox{\texttt{\mdseries\slshape digraph}}. A clique is defined by the set of vertices that it contains; see \texttt{IsClique} (\ref{IsClique}) and \texttt{IsMaximalClique} (\ref{IsMaximalClique}).

 The optional arguments \mbox{\texttt{\mdseries\slshape include}} and \mbox{\texttt{\mdseries\slshape exclude}} must each be a (possibly empty) list of vertices of \mbox{\texttt{\mdseries\slshape digraph}}, the optional argument \mbox{\texttt{\mdseries\slshape limit}} must be either a positive integer or \texttt{infinity}, and the optional argument \mbox{\texttt{\mdseries\slshape size}} must be a positive integer. If not specified, then \mbox{\texttt{\mdseries\slshape include}} and \mbox{\texttt{\mdseries\slshape exclude}} are chosen to be empty lists, and \mbox{\texttt{\mdseries\slshape limit}} is set to \texttt{infinity}. 

 The functions will return as many suitable cliques as possible, up to the
number \mbox{\texttt{\mdseries\slshape limit}}. These functions will find cliques that contain all of the vertices of \mbox{\texttt{\mdseries\slshape include}} but do not contain any of the vertices of \mbox{\texttt{\mdseries\slshape exclude}}. The argument \mbox{\texttt{\mdseries\slshape size}} restricts the search to those cliques that contain precisely \mbox{\texttt{\mdseries\slshape size}} vertices. If the function or attribute has \texttt{Maximal} in its name, then only maximal cliques will be returned; otherwise
non\texttt{\symbol{45}}maximal cliques may be returned. 

 Let \texttt{G} denote the automorphism group of maximal symmetric subdigraph of \mbox{\texttt{\mdseries\slshape digraph}} without loops (see \texttt{AutomorphismGroup} (\ref{AutomorphismGroup:for a digraph}) and \texttt{MaximalSymmetricSubdigraphWithoutLoops} (\ref{MaximalSymmetricSubdigraphWithoutLoops})). 
\begin{description}
\item[{Distinct cliques}]  \texttt{DigraphMaximalCliques} and \texttt{DigraphCliques} each return a duplicate\texttt{\symbol{45}}free list of at most \mbox{\texttt{\mdseries\slshape limit}} cliques of \mbox{\texttt{\mdseries\slshape digraph}} that satisfy the arguments.

 The computation may be significantly faster if \mbox{\texttt{\mdseries\slshape include}} and \mbox{\texttt{\mdseries\slshape exclude}} are invariant under the action of \texttt{G} on sets of vertices. 
\item[{Orbit representatives of cliques}]  To use \texttt{DigraphMaximalCliquesReps} or \texttt{DigraphCliquesReps}, the arguments \mbox{\texttt{\mdseries\slshape include}} and \mbox{\texttt{\mdseries\slshape exclude}} must each be invariant under the action of \texttt{G} on sets of vertices.

 If this is the case, then \texttt{DigraphMaximalCliquesReps} and \texttt{DigraphCliquesReps} each return a duplicate\texttt{\symbol{45}}free list of at most \mbox{\texttt{\mdseries\slshape limit}} orbits representatives (under the action of \texttt{G} on sets vertices) of cliques of \mbox{\texttt{\mdseries\slshape digraph}} that satisfy the arguments.

 The representatives are not guaranteed to be in distinct orbits. However, if
fewer than \mbox{\texttt{\mdseries\slshape lim}} results are returned, then there will be at least one representative from each
orbit of maximal cliques. 
\end{description}
 
\begin{Verbatim}[commandchars=!@|,fontsize=\small,frame=single,label=Example]
  !gapprompt@gap>| !gapinput@D := Digraph([|
  !gapprompt@>| !gapinput@[2, 3], [1, 3], [1, 2, 4], [3, 5, 6], [4, 6], [4, 5]]);|
  <immutable digraph with 6 vertices, 14 edges>
  !gapprompt@gap>| !gapinput@IsSymmetricDigraph(D);|
  true
  !gapprompt@gap>| !gapinput@G := AutomorphismGroup(D);|
  Group([ (5,6), (1,2), (1,5)(2,6)(3,4) ])
  !gapprompt@gap>| !gapinput@AsSet(DigraphMaximalCliques(D));|
  [ [ 1, 2, 3 ], [ 3, 4 ], [ 4, 5, 6 ] ]
  !gapprompt@gap>| !gapinput@DigraphMaximalCliquesReps(D);|
  [ [ 1, 2, 3 ], [ 3, 4 ] ]
  !gapprompt@gap>| !gapinput@Orbit(G, [1, 2, 3], OnSets);|
  [ [ 1, 2, 3 ], [ 4, 5, 6 ] ]
  !gapprompt@gap>| !gapinput@Orbit(G, [3, 4], OnSets);|
  [ [ 3, 4 ] ]
  !gapprompt@gap>| !gapinput@DigraphMaximalCliquesReps(D, [3, 4], [], 1);|
  [ [ 3, 4 ] ]
  !gapprompt@gap>| !gapinput@DigraphMaximalCliques(D, [1, 2], [5, 6], 1, 2);|
  [  ]
  !gapprompt@gap>| !gapinput@DigraphCliques(D, [1], [5, 6], infinity, 2);|
  [ [ 1, 2 ], [ 1, 3 ] ]
  !gapprompt@gap>| !gapinput@D := Digraph(IsMutableDigraph, [|
  !gapprompt@>| !gapinput@[2, 3], [1, 3], [1, 2, 4], [3, 5, 6], [4, 6], [4, 5]]);|
  <mutable digraph with 6 vertices, 14 edges>
  !gapprompt@gap>| !gapinput@IsSymmetricDigraph(D);|
  true
  !gapprompt@gap>| !gapinput@G := AutomorphismGroup(D);|
  Group([ (5,6), (1,2), (1,5)(2,6)(3,4) ])
  !gapprompt@gap>| !gapinput@AsSet(DigraphMaximalCliques(D));|
  [ [ 1, 2, 3 ], [ 3, 4 ], [ 4, 5, 6 ] ]
\end{Verbatim}
 }

 

\subsection{\textcolor{Chapter }{CliqueNumber}}
\logpage{[ 8, 1, 5 ]}\nobreak
\hyperdef{L}{X78427A8B81FEB457}{}
{\noindent\textcolor{FuncColor}{$\triangleright$\enspace\texttt{CliqueNumber({\mdseries\slshape digraph})\index{CliqueNumber@\texttt{CliqueNumber}}
\label{CliqueNumber}
}\hfill{\scriptsize (attribute)}}\\
\textbf{\indent Returns:\ }
A non\texttt{\symbol{45}}negative integer.



 If \mbox{\texttt{\mdseries\slshape digraph}} is a digraph, then \texttt{CliqueNumber(\mbox{\texttt{\mdseries\slshape digraph}})} returns the largest integer \texttt{n} such that \mbox{\texttt{\mdseries\slshape digraph}} contains a clique with \texttt{n} vertices as an induced subdigraph. 

 A \emph{clique} of a digraph is a set of mutually adjacent vertices of the digraph. Loops and
multiple edges are ignored for the purpose of determining the clique number of
a digraph. 
\begin{Verbatim}[commandchars=!@|,fontsize=\small,frame=single,label=Example]
  !gapprompt@gap>| !gapinput@D := CompleteDigraph(4);;|
  !gapprompt@gap>| !gapinput@CliqueNumber(D);|
  4
  !gapprompt@gap>| !gapinput@D := Digraph([[1, 2, 4, 4], [1, 3, 4], [2, 1], [1, 2]]);|
  <immutable multidigraph with 4 vertices, 11 edges>
  !gapprompt@gap>| !gapinput@CliqueNumber(D);|
  3
  !gapprompt@gap>| !gapinput@D := CompleteDigraph(IsMutableDigraph, 4);;|
  !gapprompt@gap>| !gapinput@CliqueNumber(D);|
  4
\end{Verbatim}
 }

 }

 
\section{\textcolor{Chapter }{Finding independent sets}}\logpage{[ 8, 2, 0 ]}
\hyperdef{L}{X7E31253287708348}{}
{
 

\subsection{\textcolor{Chapter }{IsIndependentSet}}
\logpage{[ 8, 2, 1 ]}\nobreak
\hyperdef{L}{X790733B67B0BC894}{}
{\noindent\textcolor{FuncColor}{$\triangleright$\enspace\texttt{IsIndependentSet({\mdseries\slshape digraph, l})\index{IsIndependentSet@\texttt{IsIndependentSet}}
\label{IsIndependentSet}
}\hfill{\scriptsize (operation)}}\\
\noindent\textcolor{FuncColor}{$\triangleright$\enspace\texttt{IsMaximalIndependentSet({\mdseries\slshape digraph, l})\index{IsMaximalIndependentSet@\texttt{IsMaximalIndependentSet}}
\label{IsMaximalIndependentSet}
}\hfill{\scriptsize (operation)}}\\
\textbf{\indent Returns:\ }
\texttt{true} or \texttt{false}.



 If \mbox{\texttt{\mdseries\slshape digraph}} is a digraph and \mbox{\texttt{\mdseries\slshape l}} is a duplicate\texttt{\symbol{45}}free list of vertices of \mbox{\texttt{\mdseries\slshape digraph}}, then \texttt{IsIndependentSet(}\mbox{\texttt{\mdseries\slshape digraph}}\texttt{,}\mbox{\texttt{\mdseries\slshape l}}\texttt{)} returns \texttt{true} if \mbox{\texttt{\mdseries\slshape l}} is an \emph{independent set} of \mbox{\texttt{\mdseries\slshape digraph}} and \texttt{false} if it is not. Similarly, \texttt{IsMaximalIndependentSet(}\mbox{\texttt{\mdseries\slshape digraph}}\texttt{,}\mbox{\texttt{\mdseries\slshape l}}\texttt{)} returns \texttt{true} if \mbox{\texttt{\mdseries\slshape l}} is a \emph{maximal independent set} of \mbox{\texttt{\mdseries\slshape digraph}} and \texttt{false} if it is not. 

 An \emph{independent set} of a digraph is a set of mutually non\texttt{\symbol{45}}adjacent vertices of
the digraph. A \emph{maximal independent set} is an independent set that is not properly contained in another independent
set. An independent set is permitted, but not required, to contain vertices at
which there is a loop. 
\begin{Verbatim}[commandchars=!@|,fontsize=\small,frame=single,label=Example]
  !gapprompt@gap>| !gapinput@D := CycleDigraph(4);;|
  !gapprompt@gap>| !gapinput@IsIndependentSet(D, [1]);|
  true
  !gapprompt@gap>| !gapinput@IsMaximalIndependentSet(D, [1]);|
  false
  !gapprompt@gap>| !gapinput@IsIndependentSet(D, [1, 4, 3]);|
  false
  !gapprompt@gap>| !gapinput@IsIndependentSet(D, [2, 4]);|
  true
  !gapprompt@gap>| !gapinput@IsMaximalIndependentSet(D, [2, 4]);|
  true
  !gapprompt@gap>| !gapinput@D := CycleDigraph(IsMutableDigraph, 4);;|
  !gapprompt@gap>| !gapinput@IsIndependentSet(D, [1]);|
  true
\end{Verbatim}
 }

 

\subsection{\textcolor{Chapter }{DigraphIndependentSet}}
\logpage{[ 8, 2, 2 ]}\nobreak
\hyperdef{L}{X84350678782BFAB7}{}
{\noindent\textcolor{FuncColor}{$\triangleright$\enspace\texttt{DigraphIndependentSet({\mdseries\slshape digraph[, include[, exclude[, size]]]})\index{DigraphIndependentSet@\texttt{DigraphIndependentSet}}
\label{DigraphIndependentSet}
}\hfill{\scriptsize (function)}}\\
\noindent\textcolor{FuncColor}{$\triangleright$\enspace\texttt{DigraphMaximalIndependentSet({\mdseries\slshape digraph[, include[, exclude[, size]]]})\index{DigraphMaximalIndependentSet@\texttt{DigraphMaximalIndependentSet}}
\label{DigraphMaximalIndependentSet}
}\hfill{\scriptsize (function)}}\\
\textbf{\indent Returns:\ }
An immutable list of positive integers, or \texttt{fail}.



 If \mbox{\texttt{\mdseries\slshape digraph}} is a digraph, then these functions returns an independent set of \mbox{\texttt{\mdseries\slshape digraph}} if one exists that satisfies the arguments, else it returns \texttt{fail}. An independent set is defined by the set of vertices that it contains; see \texttt{IsIndependentSet} (\ref{IsIndependentSet}) and \texttt{IsMaximalIndependentSet} (\ref{IsMaximalIndependentSet}).

 The optional arguments \mbox{\texttt{\mdseries\slshape include}} and \mbox{\texttt{\mdseries\slshape exclude}} must each be a (possibly empty) duplicate\texttt{\symbol{45}}free list of
vertices of \mbox{\texttt{\mdseries\slshape digraph}}, and the optional argument \mbox{\texttt{\mdseries\slshape size}} must be a positive integer. By default, \mbox{\texttt{\mdseries\slshape include}} and \mbox{\texttt{\mdseries\slshape exclude}} are empty. These functions will search for an independent set of \mbox{\texttt{\mdseries\slshape digraph}} that includes the vertices of \mbox{\texttt{\mdseries\slshape include}} but does not include any vertices in \mbox{\texttt{\mdseries\slshape exclude}}; if the argument \mbox{\texttt{\mdseries\slshape size}} is supplied, then additionally the independent set will be required to contain
precisely \mbox{\texttt{\mdseries\slshape size}} vertices.

 If \mbox{\texttt{\mdseries\slshape include}} is not an independent set, then these functions return \texttt{fail}. Otherwise, the functions behave in the following way, depending on the
number of arguments: 
\begin{description}
\item[{One or two arguments}]  If one or two arguments are supplied, then \texttt{DigraphIndependentSet} and \texttt{DigraphMaximalIndependentSet} greedily enlarge the independent set \mbox{\texttt{\mdseries\slshape include}} until it cannot continue. The result is guaranteed to be a maximal independent
set. This may or may not return an answer more quickly than using \texttt{DigraphMaximalIndependentSets} (\ref{DigraphMaximalIndependentSets}) with a limit of 1. 
\item[{Three arguments}]  If three arguments are supplied, then \texttt{DigraphIndependentSet} greedily enlarges the independent set \mbox{\texttt{\mdseries\slshape include}} until it cannot continue, although this independent set may not be maximal.

 Given three arguments, \texttt{DigraphMaximalIndependentSet} returns the maximal independent set returned by \texttt{DigraphMaximalIndependentSets(}\mbox{\texttt{\mdseries\slshape digraph}}\texttt{, }\mbox{\texttt{\mdseries\slshape include}}\texttt{, }\mbox{\texttt{\mdseries\slshape exclude}}\texttt{, 1)} if one exists, else \texttt{fail}. 
\item[{Four arguments}]  If four arguments are supplied, then \texttt{DigraphIndependentSet} returns the independent set returned by \texttt{DigraphIndependentSets(}\mbox{\texttt{\mdseries\slshape digraph}}\texttt{, }\mbox{\texttt{\mdseries\slshape include}}\texttt{, }\mbox{\texttt{\mdseries\slshape exclude}}\texttt{, 1, }\mbox{\texttt{\mdseries\slshape size}}\texttt{)} if one exists, else \texttt{fail}. This independent set may not be maximal.

 Given four arguments, \texttt{DigraphMaximalIndependentSet} returns the maximal independent set returned by \texttt{DigraphMaximalIndependentSets(}\mbox{\texttt{\mdseries\slshape digraph}}\texttt{, }\mbox{\texttt{\mdseries\slshape include}}\texttt{, }\mbox{\texttt{\mdseries\slshape exclude}}\texttt{, 1, }\mbox{\texttt{\mdseries\slshape size}}\texttt{)} if one exists, else \texttt{fail}. 
\end{description}
 
\begin{Verbatim}[commandchars=!@|,fontsize=\small,frame=single,label=Example]
  !gapprompt@gap>| !gapinput@D := ChainDigraph(6);|
  <immutable chain digraph with 6 vertices>
  !gapprompt@gap>| !gapinput@DigraphIndependentSet(D);|
  [ 6, 4, 2 ]
  !gapprompt@gap>| !gapinput@DigraphMaximalIndependentSet(D);|
  [ 6, 4, 2 ]
  !gapprompt@gap>| !gapinput@DigraphIndependentSet(D, [2, 4]);|
  [ 2, 4, 6 ]
  !gapprompt@gap>| !gapinput@DigraphMaximalIndependentSet(D, [1, 3]);|
  [ 1, 3, 6 ]
  !gapprompt@gap>| !gapinput@DigraphIndependentSet(D, [2, 4], [6]);|
  [ 2, 4 ]
  !gapprompt@gap>| !gapinput@DigraphMaximalIndependentSet(D, [2, 4], [6]);|
  fail
  !gapprompt@gap>| !gapinput@DigraphIndependentSet(D, [1], [], 2);|
  [ 1, 3 ]
  !gapprompt@gap>| !gapinput@DigraphMaximalIndependentSet(D, [1], [], 2);|
  fail
  !gapprompt@gap>| !gapinput@DigraphMaximalIndependentSet(D, [1], [], 3);|
  [ 1, 3, 5 ]
  !gapprompt@gap>| !gapinput@D := ChainDigraph(IsMutableDigraph, 6);|
  <mutable digraph with 6 vertices, 5 edges>
  !gapprompt@gap>| !gapinput@DigraphIndependentSet(D);|
  [ 6, 4, 2 ]
  !gapprompt@gap>| !gapinput@DigraphMaximalIndependentSet(D);|
  [ 6, 4, 2 ]
  !gapprompt@gap>| !gapinput@DigraphIndependentSet(D, [2, 4]);|
  [ 2, 4, 6 ]
  !gapprompt@gap>| !gapinput@DigraphMaximalIndependentSet(D, [1, 3]);|
  [ 1, 3, 6 ]
  !gapprompt@gap>| !gapinput@DigraphIndependentSet(D, [2, 4], [6]);|
  [ 2, 4 ]
  !gapprompt@gap>| !gapinput@DigraphMaximalIndependentSet(D, [2, 4], [6]);|
  fail
  !gapprompt@gap>| !gapinput@DigraphIndependentSet(D, [1], [], 2);|
  [ 1, 3 ]
  !gapprompt@gap>| !gapinput@DigraphMaximalIndependentSet(D, [1], [], 2);|
  fail
  !gapprompt@gap>| !gapinput@DigraphMaximalIndependentSet(D, [1], [], 3);|
  [ 1, 3, 5 ]
\end{Verbatim}
 }

 

\subsection{\textcolor{Chapter }{DigraphMaximalIndependentSets}}
\logpage{[ 8, 2, 3 ]}\nobreak
\hyperdef{L}{X7D2C78EB81A798A9}{}
{\noindent\textcolor{FuncColor}{$\triangleright$\enspace\texttt{DigraphMaximalIndependentSets({\mdseries\slshape digraph[, include[, exclude[, limit[, size]]]]})\index{DigraphMaximalIndependentSets@\texttt{DigraphMaximalIndependentSets}}
\label{DigraphMaximalIndependentSets}
}\hfill{\scriptsize (function)}}\\
\noindent\textcolor{FuncColor}{$\triangleright$\enspace\texttt{DigraphMaximalIndependentSetsReps({\mdseries\slshape digraph[, include[, exclude[, limit[, size]]]]})\index{DigraphMaximalIndependentSetsReps@\texttt{DigraphMaximalIndependentSetsReps}}
\label{DigraphMaximalIndependentSetsReps}
}\hfill{\scriptsize (function)}}\\
\noindent\textcolor{FuncColor}{$\triangleright$\enspace\texttt{DigraphIndependentSets({\mdseries\slshape digraph[, include[, exclude[, limit[, size]]]]})\index{DigraphIndependentSets@\texttt{DigraphIndependentSets}}
\label{DigraphIndependentSets}
}\hfill{\scriptsize (function)}}\\
\noindent\textcolor{FuncColor}{$\triangleright$\enspace\texttt{DigraphIndependentSetsReps({\mdseries\slshape digraph[, include[, exclude[, limit[, size]]]]})\index{DigraphIndependentSetsReps@\texttt{DigraphIndependentSetsReps}}
\label{DigraphIndependentSetsReps}
}\hfill{\scriptsize (function)}}\\
\noindent\textcolor{FuncColor}{$\triangleright$\enspace\texttt{DigraphMaximalIndependentSetsAttr({\mdseries\slshape digraph})\index{DigraphMaximalIndependentSetsAttr@\texttt{DigraphMaximalIndependentSetsAttr}}
\label{DigraphMaximalIndependentSetsAttr}
}\hfill{\scriptsize (attribute)}}\\
\noindent\textcolor{FuncColor}{$\triangleright$\enspace\texttt{DigraphMaximalIndependentSetsRepsAttr({\mdseries\slshape digraph})\index{DigraphMaximalIndependentSetsRepsAttr@\texttt{Digraph}\-\texttt{Maximal}\-\texttt{Independent}\-\texttt{Sets}\-\texttt{Reps}\-\texttt{Attr}}
\label{DigraphMaximalIndependentSetsRepsAttr}
}\hfill{\scriptsize (attribute)}}\\
\textbf{\indent Returns:\ }
An immutable list of lists of positive integers.



 If \mbox{\texttt{\mdseries\slshape digraph}} is digraph, then these functions and attributes use \texttt{CliquesFinder} (\ref{CliquesFinder}) to return independent sets of \mbox{\texttt{\mdseries\slshape digraph}}. An independent set is defined by the set of vertices that it contains; see \texttt{IsMaximalIndependentSet} (\ref{IsMaximalIndependentSet}) and \texttt{IsIndependentSet} (\ref{IsIndependentSet}).

 The optional arguments \mbox{\texttt{\mdseries\slshape include}} and \mbox{\texttt{\mdseries\slshape exclude}} must each be a (possibly empty) list of vertices of \mbox{\texttt{\mdseries\slshape digraph}}, the optional argument \mbox{\texttt{\mdseries\slshape limit}} must be either a positive integer or \texttt{infinity}, and the optional argument \mbox{\texttt{\mdseries\slshape size}} must be a positive integer. If not specified, then \mbox{\texttt{\mdseries\slshape include}} and \mbox{\texttt{\mdseries\slshape exclude}} are chosen to be empty lists, and \mbox{\texttt{\mdseries\slshape limit}} is set to \texttt{infinity}. 

 The functions will return as many suitable independent sets as possible, up to
the number \mbox{\texttt{\mdseries\slshape limit}}. These functions will find independent sets that contain all of the vertices
of \mbox{\texttt{\mdseries\slshape include}} but do not contain any of the vertices of \mbox{\texttt{\mdseries\slshape exclude}} The argument \mbox{\texttt{\mdseries\slshape size}} restricts the search to those cliques that contain precisely \mbox{\texttt{\mdseries\slshape size}} vertices. If the function or attribute has \texttt{Maximal} in its name, then only maximal independent sets will be returned; otherwise
non\texttt{\symbol{45}}maximal independent sets may be returned. 

 Let \texttt{G} denote the \texttt{AutomorphismGroup} (\ref{AutomorphismGroup:for a digraph}) of the \texttt{DigraphSymmetricClosure} (\ref{DigraphSymmetricClosure}) of the digraph formed from \mbox{\texttt{\mdseries\slshape digraph}} by removing loops and ignoring the multiplicity of edges. 
\begin{description}
\item[{Distinct independent sets}]  \texttt{DigraphMaximalIndependentSets} and \texttt{DigraphIndependentSets} each return a duplicate\texttt{\symbol{45}}free list of at most \mbox{\texttt{\mdseries\slshape limit}} independent sets of \mbox{\texttt{\mdseries\slshape digraph}} that satisfy the arguments.

 The computation may be significantly faster if \mbox{\texttt{\mdseries\slshape include}} and \mbox{\texttt{\mdseries\slshape exclude}} are invariant under the action of \texttt{G} on sets of vertices. 
\item[{Representatives of distinct orbits of independent sets}]  To use \texttt{DigraphMaximalIndependentSetsReps} or \texttt{DigraphIndependentSetsReps}, the arguments \mbox{\texttt{\mdseries\slshape include}} and \mbox{\texttt{\mdseries\slshape exclude}} must each be invariant under the action of \texttt{G} on sets of vertices.

 If this is the case, then \texttt{DigraphMaximalIndependentSetsReps} and \texttt{DigraphIndependentSetsReps} each return a list of at most \mbox{\texttt{\mdseries\slshape limit}} orbits representatives (under the action of \texttt{G} on sets of vertices) of independent sets of \mbox{\texttt{\mdseries\slshape digraph}} that satisfy the arguments. 

 The representatives are not guaranteed to be in distinct orbits. However, if \mbox{\texttt{\mdseries\slshape lim}} is not specified, or fewer than \mbox{\texttt{\mdseries\slshape lim}} results are returned, then there will be at least one representative from each
orbit of maximal independent sets. 
\end{description}
 
\begin{Verbatim}[commandchars=!@|,fontsize=\small,frame=single,label=Example]
  !gapprompt@gap>| !gapinput@D := CycleDigraph(5);|
  <immutable cycle digraph with 5 vertices>
  !gapprompt@gap>| !gapinput@DigraphMaximalIndependentSetsReps(D);|
  [ [ 1, 3 ] ]
  !gapprompt@gap>| !gapinput@DigraphIndependentSetsReps(D);|
  [ [ 1 ], [ 1, 3 ] ]
  !gapprompt@gap>| !gapinput@Set(DigraphMaximalIndependentSets(D));|
  [ [ 1, 3 ], [ 1, 4 ], [ 2, 4 ], [ 2, 5 ], [ 3, 5 ] ]
  !gapprompt@gap>| !gapinput@DigraphMaximalIndependentSets(D, [1]);|
  [ [ 1, 3 ], [ 1, 4 ] ]
  !gapprompt@gap>| !gapinput@DigraphIndependentSets(D, [], [4, 5]);|
  [ [ 1 ], [ 2 ], [ 3 ], [ 1, 3 ] ]
  !gapprompt@gap>| !gapinput@DigraphIndependentSets(D, [], [4, 5], 1, 2);|
  [ [ 1, 3 ] ]
  !gapprompt@gap>| !gapinput@D := CycleDigraph(IsMutableDigraph, 5);|
  <mutable digraph with 5 vertices, 5 edges>
  !gapprompt@gap>| !gapinput@DigraphMaximalIndependentSetsReps(D);|
  [ [ 1, 3 ] ]
  !gapprompt@gap>| !gapinput@DigraphIndependentSetsReps(D);|
  [ [ 1 ], [ 1, 3 ] ]
  !gapprompt@gap>| !gapinput@Set(DigraphMaximalIndependentSets(D));|
  [ [ 1, 3 ], [ 1, 4 ], [ 2, 4 ], [ 2, 5 ], [ 3, 5 ] ]
  !gapprompt@gap>| !gapinput@DigraphMaximalIndependentSets(D, [1]);|
  [ [ 1, 3 ], [ 1, 4 ] ]
  !gapprompt@gap>| !gapinput@DigraphIndependentSets(D, [], [4, 5]);|
  [ [ 1 ], [ 2 ], [ 3 ], [ 1, 3 ] ]
  !gapprompt@gap>| !gapinput@DigraphIndependentSets(D, [], [4, 5], 1, 2);|
  [ [ 1, 3 ] ]
\end{Verbatim}
 }

 }

 }

 
\chapter{\textcolor{Chapter }{Visualising and IO}}\label{Visualising and IO}
\logpage{[ 9, 0, 0 ]}
\hyperdef{L}{X81714DBC80CC4BCB}{}
{
 
\section{\textcolor{Chapter }{Visualising a digraph}}\logpage{[ 9, 1, 0 ]}
\hyperdef{L}{X7CF67F8278D9469C}{}
{
 

\subsection{\textcolor{Chapter }{Splash}}
\logpage{[ 9, 1, 1 ]}\nobreak
\hyperdef{L}{X83B3318784E78415}{}
{\noindent\textcolor{FuncColor}{$\triangleright$\enspace\texttt{Splash({\mdseries\slshape str[, opts]})\index{Splash@\texttt{Splash}}
\label{Splash}
}\hfill{\scriptsize (function)}}\\
\textbf{\indent Returns:\ }
Nothing.



 This function attempts to convert the string \mbox{\texttt{\mdseries\slshape str}} into a pdf document and open this document, i.e. to splash it all over your
monitor.

 The string \mbox{\texttt{\mdseries\slshape str}} must correspond to a valid \texttt{dot} or \texttt{LaTeX} text file and you must have have \texttt{GraphViz} and \texttt{pdflatex} installed on your computer. For details about these file formats, see \href{https://www.latex-project.org} {\texttt{https://www.latex\texttt{\symbol{45}}project.org}} and \href{https://www.graphviz.org} {\texttt{https://www.graphviz.org}}.

 This function is provided to allow convenient, immediate viewing of the
pictures produced by the function \texttt{DotDigraph} (\ref{DotDigraph}).

 The optional second argument \mbox{\texttt{\mdseries\slshape opts}} should be a record with components corresponding to various options, given
below. 
\begin{description}
\item[{path}]  this should be a string representing the path to the directory where you want \texttt{Splash} to do its work. The default value of this option is \texttt{"\texttt{\symbol{126}}/"}. 
\item[{directory}]  this should be a string representing the name of the directory in \texttt{path} where you want \texttt{Splash} to do its work. This function will create this directory if does not already
exist. 

 The default value of this option is \texttt{"tmp.viz"} if the option \texttt{path} is present, and the result of \texttt{DirectoryTemporary} (\textbf{Reference: DirectoryTemporary}) is used otherwise. 
\item[{filename}]  this should be a string representing the name of the file where \mbox{\texttt{\mdseries\slshape str}} will be written. The default value of this option is \texttt{"vizpicture"}. 
\item[{viewer}]  this should be a string representing the name of the program which should open
the files produced by \texttt{GraphViz} or \texttt{pdflatex}. 
\item[{type}]  this option can be used to specify that the string \mbox{\texttt{\mdseries\slshape str}} contains a {\LaTeX} or \texttt{dot} document. You can specify this option in \mbox{\texttt{\mdseries\slshape str}} directly by making the first line \texttt{"\%latex"} or \texttt{"//dot"}. There is no default value for this option, this option must be specified in \mbox{\texttt{\mdseries\slshape str}} or in \mbox{\texttt{\mdseries\slshape opt.type}}. 
\item[{engine}]  this option can be used to specify the GraphViz engine to use to render a \texttt{dot} document. The valid choices are \texttt{"dot"}, \texttt{"neato"}, \texttt{"circo"}, \texttt{"twopi"}, \texttt{"fdp"}, \texttt{"sfdp"}, and \texttt{"patchwork"}. Please refer to the GraphViz documentation for details on these engines. The
default value for this option is \texttt{"dot"}, and it must be specified in \mbox{\texttt{\mdseries\slshape opt.engine}}. 
\item[{filetype}]  this should be a string representing the type of file which \texttt{Splash} should produce. For {\LaTeX} files, this option is ignored and the default value \texttt{"pdf"} is used. 
\end{description}
 This function was written by Attila Egri\texttt{\symbol{45}}Nagy and Manuel
Delgado with some minor changes by J. D. Mitchell.

 
\begin{Verbatim}[commandchars=!@|,fontsize=\small,frame=single,label=Example]
  !gapprompt@gap>| !gapinput@Splash(DotDigraph(RandomDigraph(4)));|
\end{Verbatim}
 }

 

\subsection{\textcolor{Chapter }{DotDigraph}}
\logpage{[ 9, 1, 2 ]}\nobreak
\hyperdef{L}{X7F9B99C478EE093A}{}
{\noindent\textcolor{FuncColor}{$\triangleright$\enspace\texttt{DotDigraph({\mdseries\slshape digraph})\index{DotDigraph@\texttt{DotDigraph}}
\label{DotDigraph}
}\hfill{\scriptsize (attribute)}}\\
\noindent\textcolor{FuncColor}{$\triangleright$\enspace\texttt{DotColoredDigraph({\mdseries\slshape digraph, vert, edge})\index{DotColoredDigraph@\texttt{DotColoredDigraph}}
\label{DotColoredDigraph}
}\hfill{\scriptsize (operation)}}\\
\noindent\textcolor{FuncColor}{$\triangleright$\enspace\texttt{DotVertexLabelledDigraph({\mdseries\slshape digraph})\index{DotVertexLabelledDigraph@\texttt{DotVertexLabelledDigraph}}
\label{DotVertexLabelledDigraph}
}\hfill{\scriptsize (operation)}}\\
\noindent\textcolor{FuncColor}{$\triangleright$\enspace\texttt{DotVertexColoredDigraph({\mdseries\slshape digraph, vert})\index{DotVertexColoredDigraph@\texttt{DotVertexColoredDigraph}}
\label{DotVertexColoredDigraph}
}\hfill{\scriptsize (operation)}}\\
\noindent\textcolor{FuncColor}{$\triangleright$\enspace\texttt{DotEdgeColoredDigraph({\mdseries\slshape digraph, edge})\index{DotEdgeColoredDigraph@\texttt{DotEdgeColoredDigraph}}
\label{DotEdgeColoredDigraph}
}\hfill{\scriptsize (operation)}}\\
\textbf{\indent Returns:\ }
A string.



 \texttt{DotDigraph} produces a graphical representation of the digraph \mbox{\texttt{\mdseries\slshape digraph}}. Vertices are displayed as circles, numbered consistently with \mbox{\texttt{\mdseries\slshape digraph}}. Edges are displayed as arrowed lines between vertices, with the arrowhead of
each line pointing towards the range of the edge.

 \texttt{DotColoredDigraph} differs from \texttt{DotDigraph} only in that the values in given in the two lists are used to color the
vertices and edges of the graph when displayed. The list for vertex colours
should be a list of length equal to the number of vertices, containing strings
that are accepted by the graphviz software, which is the one used for graph
representation. The list for edge colours should be a list of lists with the
same shape of the outneighbours of the digraph that contains strings that
correspond to colours accepted by the graphviz software. If the lists are not
the appropriate size, or have holes then the function will return an error.

 \texttt{DotVertexColoredDigraph} differs from \texttt{DotDigraph} only in that the values in given in the list are used to color the vertices of
the graph when displayed. The list for vertex colours should be a list of
length equal to the number of vertices, containing strings that are accepted
by the graphviz software, which is the one used for graph representation. If
the list is not the appropriate size, or has holes then the function will
return an error.

 \texttt{DotEdgeColoredDigraph} differs from \texttt{DotDigraph} only in that the values in given in the list are used to color the vertices
and edges of the graph when displayed. The list for edge colours should be a
list of lists with the same shape of the outneighbours of the digraph that
contains strings that correspond to colours accepted by the graphviz software.
If the list is not the appropriate size, or has holes then the function will
return an error.

 \texttt{DotVertexLabelledDigraph} differs from \texttt{DotDigraph} only in that the values in \texttt{DigraphVertexLabels} (\ref{DigraphVertexLabels}) are used to label the vertices in the produced picture rather than the numbers \texttt{1} to the number of vertices of the digraph. 

 The output is in \texttt{dot} format (also known as \texttt{GraphViz}) format. For details about this file format, and information about how to
display or edit this format see \href{https://www.graphviz.org} {\texttt{https://www.graphviz.org}}. 

 The string returned by \texttt{DotDigraph} or \texttt{DotVertexLabelledDigraph} can be written to a file using the command \texttt{FileString} (\textbf{GAPDoc: FileString}).

 
\begin{Verbatim}[commandchars=!@|,fontsize=\small,frame=single,label=Example]
  !gapprompt@gap>| !gapinput@D := CompleteDigraph(4);|
  <immutable complete digraph with 4 vertices>
  !gapprompt@gap>| !gapinput@vertcolors := ["blue", "green", "red", "yellow"];;|
  !gapprompt@gap>| !gapinput@Print(DotVertexLabelledDigraph(D));|
  //dot
  digraph hgn{
  node [shape=circle]
  1 [label="1"]
  2 [label="2"]
  3 [label="3"]
  4 [label="4"]
  1 -> 2
  1 -> 3
  1 -> 4
  2 -> 1
  2 -> 3
  2 -> 4
  3 -> 1
  3 -> 2
  3 -> 4
  4 -> 1
  4 -> 2
  4 -> 3
  }
  !gapprompt@gap>| !gapinput@Print(DotVertexColoredDigraph(D, vertcolors));|
  //dot
  digraph hgn{
  node [shape=circle]
  1[color=blue, style=filled]
  2[color=green, style=filled]
  3[color=red, style=filled]
  4[color=yellow, style=filled]
  1 -> 2
  1 -> 3
  1 -> 4
  2 -> 1
  2 -> 3
  2 -> 4
  3 -> 1
  3 -> 2
  3 -> 4
  4 -> 1
  4 -> 2
  4 -> 3
  }
  !gapprompt@gap>| !gapinput@colors := ["lightblue", "pink", "purple"];;|
  !gapprompt@gap>| !gapinput@edgecolors := ListWithIdenticalEntries(4, colors);;|
  !gapprompt@gap>| !gapinput@Print(DotEdgeColoredDigraph(D, edgecolors));|
  //dot
  digraph hgn{
  node [shape=circle]
  1
  2
  3
  4
  1 -> 2[color=lightblue]
  1 -> 3[color=pink]
  1 -> 4[color=purple]
  2 -> 1[color=lightblue]
  2 -> 3[color=pink]
  2 -> 4[color=purple]
  3 -> 1[color=lightblue]
  3 -> 2[color=pink]
  3 -> 4[color=purple]
  4 -> 1[color=lightblue]
  4 -> 2[color=pink]
  4 -> 3[color=purple]
  }
  !gapprompt@gap>| !gapinput@Print(DotColoredDigraph(D, vertcolors, edgecolors));|
  //dot
  digraph hgn{
  node [shape=circle]
  1[color=blue, style=filled]
  2[color=green, style=filled]
  3[color=red, style=filled]
  4[color=yellow, style=filled]
  1 -> 2[color=lightblue]
  1 -> 3[color=pink]
  1 -> 4[color=purple]
  2 -> 1[color=lightblue]
  2 -> 3[color=pink]
  2 -> 4[color=purple]
  3 -> 1[color=lightblue]
  3 -> 2[color=pink]
  3 -> 4[color=purple]
  4 -> 1[color=lightblue]
  4 -> 2[color=pink]
  4 -> 3[color=purple]
  }
  !gapprompt@gap>| !gapinput@FileString("k4.dot", DotDigraph(D));|
  133
  !gapprompt@gap>| !gapinput@D := Digraph([[2, 3], [1, 3], [1]]);|
  <immutable digraph with 3 vertices, 5 edges>
  !gapprompt@gap>| !gapinput@vertcolors := ["blue", "red", "green"];;|
  !gapprompt@gap>| !gapinput@edgecolors := [["orange", "yellow"], ["orange", |
  !gapprompt@>| !gapinput@"pink"], ["yellow"]];;|
  !gapprompt@gap>| !gapinput@Print(DotColoredDigraph(D, vertcolors, edgecolors));|
  //dot
  digraph hgn{
  node [shape=circle]
  1[color=blue, style=filled]
  2[color=red, style=filled]
  3[color=green, style=filled]
  1 -> 2[color=orange]
  1 -> 3[color=yellow]
  2 -> 1[color=orange]
  2 -> 3[color=pink]
  3 -> 1[color=yellow]
  }
  !gapprompt@gap>| !gapinput@FileString("digraph.dot", DotDigraph(D));|
  82
\end{Verbatim}
 }

 

\subsection{\textcolor{Chapter }{DotSymmetricDigraph}}
\logpage{[ 9, 1, 3 ]}\nobreak
\hyperdef{L}{X78DA31D1869D2F94}{}
{\noindent\textcolor{FuncColor}{$\triangleright$\enspace\texttt{DotSymmetricDigraph({\mdseries\slshape digraph})\index{DotSymmetricDigraph@\texttt{DotSymmetricDigraph}}
\label{DotSymmetricDigraph}
}\hfill{\scriptsize (attribute)}}\\
\noindent\textcolor{FuncColor}{$\triangleright$\enspace\texttt{DotSymmetricColoredDigraph({\mdseries\slshape digraph, vert, edge})\index{DotSymmetricColoredDigraph@\texttt{DotSymmetricColoredDigraph}}
\label{DotSymmetricColoredDigraph}
}\hfill{\scriptsize (operation)}}\\
\noindent\textcolor{FuncColor}{$\triangleright$\enspace\texttt{DotSymmetricVertexColoredDigraph({\mdseries\slshape digraph, vert})\index{DotSymmetricVertexColoredDigraph@\texttt{DotSymmetricVertexColoredDigraph}}
\label{DotSymmetricVertexColoredDigraph}
}\hfill{\scriptsize (operation)}}\\
\noindent\textcolor{FuncColor}{$\triangleright$\enspace\texttt{DotSymmetricEdgeColoredDigraph({\mdseries\slshape digraph, edge})\index{DotSymmetricEdgeColoredDigraph@\texttt{DotSymmetricEdgeColoredDigraph}}
\label{DotSymmetricEdgeColoredDigraph}
}\hfill{\scriptsize (operation)}}\\
\textbf{\indent Returns:\ }
A string.



 This function produces a graphical representation of the symmetric digraph \mbox{\texttt{\mdseries\slshape digraph}}. \texttt{DotSymmetricDigraph} will return an error if \mbox{\texttt{\mdseries\slshape digraph}} is not a symmetric digraph. See \texttt{IsSymmetricDigraph} (\ref{IsSymmetricDigraph}).

 The function \texttt{DotSymmetricColoredDigraph} differs from \texttt{DotDigraph} only in that the values given in the two lists are used to color the vertices
and edges of the graph when displayed. The list for vertex colours should be a
list of length equal to the number of vertices, containing strings that are
accepted by the graphviz software, which is the one used for graph
representation. The list for edge colours should be a list of lists with the
same shape of the outneighbours of the digraph that contains strings that
correspond to colours accepted by the graphviz software. If the list is not
the appropriate size, or has holes then the function will return an error. 

 The function \texttt{DotSymmetricVertexColoredDigraph} differs from \texttt{DotDigraph} only in that the values in given in the list is used to color the vertices of
the graph when displayed. The list for vertex colours should be a list of
length equal to the number of vertices, containing strings that are accepted
by the graphviz software, which is the one used for graph representation. If
the list is not the appropriate size, or has holes then the function will
return an error. 

 The function \texttt{DotSymmetricEdgeColoredDigraph} differs from \texttt{DotDigraph} only in that the values given in the list are used to color the edges of the
graph when displayed. The list for edge colours should be a list of lists with
the same shape of the outneighbours, containing strings that are accepted by
the graphviz software, which is the one used for graph representation. If the
list is not the appropriate size, or has holes then the function will return
an error. 

 Vertices are displayed as circles, numbered consistently with \mbox{\texttt{\mdseries\slshape digraph}}. Since \mbox{\texttt{\mdseries\slshape digraph}} is symmetric, for every non\texttt{\symbol{45}}loop edge there is a
complementary edge with opposite source and range. \texttt{DotSymmetricDigraph} displays each pair of complementary edges as a single line between the
relevant vertices, with no arrowhead.

 The output is in \texttt{dot} format (also known as \texttt{GraphViz}) format. For details about this file format, and information about how to
display or edit this format see \href{https://www.graphviz.org} {\texttt{https://www.graphviz.org}}. 

 The string returned by \texttt{DotSymmetricDigraph} can be written to a file using the command \texttt{FileString} (\textbf{GAPDoc: FileString}).

 
\begin{Verbatim}[commandchars=!@|,fontsize=\small,frame=single,label=Example]
  
  !gapprompt@gap>| !gapinput@D := Digraph([[2], [1, 3], [2]]);|
  <immutable digraph with 3 vertices, 4 edges>
  !gapprompt@gap>| !gapinput@IsSymmetricDigraph(D);|
  true
  !gapprompt@gap>| !gapinput@vertcolors := ["blue", "pink", "purple"];;|
  !gapprompt@gap>| !gapinput@edgecolors := [["green"], ["green", "red"], ["red"]];;|
  !gapprompt@gap>| !gapinput@Print(DotSymmetricColoredDigraph(D, vertcolors, edgecolors));|
  //dot
  graph hgn{
  node [shape=circle]
  
  1[color=blue, style=filled]
  2[color=pink, style=filled]
  3[color=purple, style=filled]
  1 -- 2[color=green]
  2 -- 3[color=red]
  }
  !gapprompt@gap>| !gapinput@Print(DotSymmetricVertexColoredDigraph(D, vertcolors));|
  //dot
  graph hgn{
  node [shape=circle]
  
  1[color=blue, style=filled]
  2[color=pink, style=filled]
  3[color=purple, style=filled]
  1 -- 2
  2 -- 3
  }
  
  !gapprompt@gap>| !gapinput@Print(DotSymmetricEdgeColoredDigraph(D, edgecolors));|
  //dot
  graph hgn{
  node [shape=circle]
  
  1
  2
  3
  1 -- 2[color=green]
  2 -- 3[color=red]
  }
\end{Verbatim}
 }

 

\subsection{\textcolor{Chapter }{DotPartialOrderDigraph}}
\logpage{[ 9, 1, 4 ]}\nobreak
\hyperdef{L}{X84120D447E2CAC23}{}
{\noindent\textcolor{FuncColor}{$\triangleright$\enspace\texttt{DotPartialOrderDigraph({\mdseries\slshape digraph})\index{DotPartialOrderDigraph@\texttt{DotPartialOrderDigraph}}
\label{DotPartialOrderDigraph}
}\hfill{\scriptsize (attribute)}}\\
\textbf{\indent Returns:\ }
A string.



 This function produces a graphical representation of a partial order digraph \mbox{\texttt{\mdseries\slshape digraph}}. \texttt{DotPartialOrderDigraph} will return an error if \mbox{\texttt{\mdseries\slshape digraph}} is not a partial order digraph. See \texttt{IsPartialOrderDigraph} (\ref{IsPartialOrderDigraph}).

 Since \mbox{\texttt{\mdseries\slshape digraph}} is a partial order, it is both reflexive and transitive. The output of \texttt{DotPartialOrderDigraph} is the \texttt{DotDigraph} (\ref{DotDigraph}) of the \texttt{DigraphReflexiveTransitiveReduction} (\ref{DigraphReflexiveTransitiveReduction}) of \mbox{\texttt{\mdseries\slshape digraph}}.

 The output is in \texttt{dot} format (also known as \texttt{GraphViz}) format. For details about this file format, and information about how to
display or edit this format see \href{https://www.graphviz.org} {\texttt{https://www.graphviz.org}}. 

 The string returned by \texttt{DotPartialOrderDigraph} can be written to a file using the command \texttt{FileString} (\textbf{GAPDoc: FileString}).

 
\begin{Verbatim}[commandchars=!@|,fontsize=\small,frame=single,label=Example]
  !gapprompt@gap>| !gapinput@poset := Digraph([[1, 4], [2], [2, 3, 4], [4]]);|
  <immutable digraph with 4 vertices, 7 edges>
  !gapprompt@gap>| !gapinput@IsPartialOrderDigraph(poset);|
  true
  !gapprompt@gap>| !gapinput@FileString("poset.dot", DotPartialOrderDigraph(poset));|
  70
\end{Verbatim}
 }

 

\subsection{\textcolor{Chapter }{DotPreorderDigraph}}
\logpage{[ 9, 1, 5 ]}\nobreak
\hyperdef{L}{X7F5653C8854FAEA6}{}
{\noindent\textcolor{FuncColor}{$\triangleright$\enspace\texttt{DotPreorderDigraph({\mdseries\slshape digraph})\index{DotPreorderDigraph@\texttt{DotPreorderDigraph}}
\label{DotPreorderDigraph}
}\hfill{\scriptsize (attribute)}}\\
\noindent\textcolor{FuncColor}{$\triangleright$\enspace\texttt{DotQuasiorderDigraph({\mdseries\slshape digraph})\index{DotQuasiorderDigraph@\texttt{DotQuasiorderDigraph}}
\label{DotQuasiorderDigraph}
}\hfill{\scriptsize (attribute)}}\\
\textbf{\indent Returns:\ }
A string.



 This function produces a graphical representation of a preorder digraph \mbox{\texttt{\mdseries\slshape digraph}}. \texttt{DotPreorderDigraph} will return an error if \mbox{\texttt{\mdseries\slshape digraph}} is not a preorder digraph. See \texttt{IsPreorderDigraph} (\ref{IsPreorderDigraph}).

 A preorder digraph is reflexive and transitive but in general it is not
anti\texttt{\symbol{45}}symmetric and may have strongly connected components
containing more than one vertex. The \texttt{QuotientDigraph} (\ref{QuotientDigraph}) \mbox{\texttt{\mdseries\slshape Q}} obtained by forming the quotient of \mbox{\texttt{\mdseries\slshape digraph}} by the partition of its vertices into the strongly connected components
satisfies \texttt{IsPartialOrderDigraph} (\ref{IsPartialOrderDigraph}). Thus every vertex of \mbox{\texttt{\mdseries\slshape Q}} corresponds to a strongly connected component of \mbox{\texttt{\mdseries\slshape digraph}}. The output of \texttt{DotPreorderDigraph} displays the \texttt{DigraphReflexiveTransitiveReduction} (\ref{DigraphReflexiveTransitiveReduction}) of \mbox{\texttt{\mdseries\slshape Q}} with vertices displayed as rounded rectangles labelled by all of the vertices
of \mbox{\texttt{\mdseries\slshape digraph}} in the corresponding strongly connected component. 

 The output is in \texttt{dot} format (also known as \texttt{GraphViz}) format. For details about this file format, and information about how to
display or edit this format see \href{https://www.graphviz.org} {\texttt{https://www.graphviz.org}}. 

 The string returned by \texttt{DotPreorderDigraph} can be written to a file using the command \texttt{FileString} (\textbf{GAPDoc: FileString}).

 
\begin{Verbatim}[commandchars=!@|,fontsize=\small,frame=single,label=Example]
  !gapprompt@gap>| !gapinput@preset := Digraph([[1, 2, 4, 5], [1, 2, 4, 5], [3, 4], [4], |
  !gapprompt@>| !gapinput@[1, 2, 4, 5]]);|
  <immutable digraph with 5 vertices, 15 edges>
  !gapprompt@gap>| !gapinput@IsPreorderDigraph(preset);|
  true
  !gapprompt@gap>| !gapinput@FileString("preset.dot", DotPreorderDigraph(preset));|
  179
\end{Verbatim}
 }

 

\subsection{\textcolor{Chapter }{DotHighlightedDigraph}}
\logpage{[ 9, 1, 6 ]}\nobreak
\hyperdef{L}{X808D12CF81E2C19F}{}
{\noindent\textcolor{FuncColor}{$\triangleright$\enspace\texttt{DotHighlightedDigraph({\mdseries\slshape digraph, verts[, colour1, colour2]})\index{DotHighlightedDigraph@\texttt{DotHighlightedDigraph}}
\label{DotHighlightedDigraph}
}\hfill{\scriptsize (operation)}}\\
\textbf{\indent Returns:\ }
A string.



 \texttt{DotHighlightedDigraph} produces a graphical representation of the digraph \mbox{\texttt{\mdseries\slshape digraph}}, where the vertices in the list \mbox{\texttt{\mdseries\slshape verts}}, and edges between them, are drawn with colour \mbox{\texttt{\mdseries\slshape colour1}} and all other vertices and edges in \mbox{\texttt{\mdseries\slshape digraph}} are drawn with colour \mbox{\texttt{\mdseries\slshape colour2}}. If \mbox{\texttt{\mdseries\slshape colour1}} and \mbox{\texttt{\mdseries\slshape colour2}} are not given then \texttt{DotHighlightedDigraph} uses black and grey respectively. 

 Note that \texttt{DotHighlightedDigraph} does not validate the colours \mbox{\texttt{\mdseries\slshape colour1}} and \mbox{\texttt{\mdseries\slshape colour2}} \texttt{\symbol{45}} consult the GraphViz documentation to see what is
available. See \texttt{DotDigraph} (\ref{DotDigraph}) for more details on the output.

 
\begin{Verbatim}[commandchars=!@|,fontsize=\small,frame=single,label=Example]
  !gapprompt@gap>| !gapinput@digraph := Digraph([[2, 3], [2], [1, 3]]);|
  <immutable digraph with 3 vertices, 5 edges>
  !gapprompt@gap>| !gapinput@FileString("my_digraph.dot", DotHighlightedDigraph(digraph, |
  !gapprompt@>| !gapinput@[1, 2], "red", "black"));|
  264
\end{Verbatim}
 }

 }

 
\section{\textcolor{Chapter }{Reading and writing digraphs to a file}}\logpage{[ 9, 2, 0 ]}
\hyperdef{L}{X81C9B1DC791C0BBD}{}
{
 This section describes different ways to store and read graphs from a file in
the \textsf{Digraphs} package. 
\begin{description}
\item[{ Graph6 }]  \emph{ Graph6 } is a graph data format for storing undirected graphs with no multiple edges
nor loops of size up to $ 2^{36} - 1 $ in printable characters. The format consists of two parts. The first part uses
one to eight bytes to store the number of vertices. And the second part is the
upper half of the adjacency matrix converted into ASCII characters. For a more
detail description see \href{ http://cs.anu.edu.au/~bdm/data/formats.txt} {Graph6}. 
\item[{ Sparse6 }]  \emph{ Sparse6 } is a graph data format for storing undirected graphs with possibly multiple
edges or loops. The maximal number of vertices allowed is $ 2^{36} - 1 $. The format consists of two parts. The first part uses one to eight bytes to
store the number of vertices. And the second part only stores information
about the edges. Therefore, the \emph{ Sparse6 } format return a more compact encoding then \emph{ Graph6 } for sparse graph, i.e. graphs where the number of edges is much less than the
number of vertices squared. For a more detail description see \href{ http://cs.anu.edu.au/~bdm/data/formats.txt} {Sparse6}. 
\item[{ Digraph6 }]  \emph{Digraph6} is a new format based on \emph{ Graph6 }, but designed for digraphs. The entire adjacency matrix is stored, and
therefore there is support for directed edges and
single\texttt{\symbol{45}}vertex loops. However, multiple edges are not
supported. 
\item[{ DiSparse6 }]  \emph{DiSparse6} is a new format based on \emph{ Sparse6 }, but designed for digraphs. In this format the list of edges is partitioned
into increasing and decreasing edges, depending whether the edge has its
source bigger than the range. Then both sets of edges are written separately
in \emph{ Sparse6 } format with a separation symbol in between. 
\item[{ dreadnaut }]  \emph{dreadnaut} is a graph data format designed for directed and undirected graphs. The format
supports loops but multiple edges are ignored. The format consists of an
initial section that defines the graph's structural properties, such as the
number of vertices, the starting value for vertices, and whether the graph is
directed. This is followed by a list of edges. For more information and
examples of the format see \href{http://pallini.di.uniroma1.it/Guide.html} {nauty and Traces User's Guide}. 
\item[{ DIMACS }]  \emph{DIMACS} is a graph data format that can be used for symmetric digraphs. For a more
detailed description, see \texttt{WriteDIMACSDigraph} (\ref{WriteDIMACSDigraph}) 
\end{description}
 

\subsection{\textcolor{Chapter }{String}}
\logpage{[ 9, 2, 1 ]}\nobreak
\hyperdef{L}{X81FB5BE27903EC32}{}
{\noindent\textcolor{FuncColor}{$\triangleright$\enspace\texttt{String({\mdseries\slshape digraph})\index{String@\texttt{String}}
\label{String}
}\hfill{\scriptsize (attribute)}}\\
\noindent\textcolor{FuncColor}{$\triangleright$\enspace\texttt{PrintString({\mdseries\slshape digraph})\index{PrintString@\texttt{PrintString}}
\label{PrintString}
}\hfill{\scriptsize (operation)}}\\
\textbf{\indent Returns:\ }
A string.



 Returns a string \texttt{string} such that \texttt{EvalString(string)} is equal to \mbox{\texttt{\mdseries\slshape digraph}}, and has the same mutability. See \texttt{EvalString} (\textbf{Reference: EvalString}). 

 The methods installed for \texttt{String} make some attempts to ensure that \texttt{string} has as short a length as possible, but there may exist shorter strings that
also evaluate to \mbox{\texttt{\mdseries\slshape digraph}}. 

 It is possible that \texttt{string} may contain escaped special characters. To obtain a representation of \mbox{\texttt{\mdseries\slshape digraph}} that can be entered as GAP input, please use \texttt{Print} (\textbf{Reference: Print}). Note that \texttt{Print} for a digraph delegates to \texttt{PrintString}, which delegates to \texttt{String}. 
\begin{Verbatim}[commandchars=!|B,fontsize=\small,frame=single,label=Example]
  !gapprompt|gap>B !gapinput|D := CycleDigraph(3);B
  <immutable cycle digraph with 3 vertices>
  !gapprompt|gap>B !gapinput|Print(D);B
  CycleDigraph(3);
  !gapprompt|gap>B !gapinput|G := PetersenGraph(IsMutableDigraph);B
  <mutable digraph with 10 vertices, 30 edges>
  !gapprompt|gap>B !gapinput|String(G);B
  "DigraphFromGraph6String(IsMutableDigraph, \"IheA@GUAo\");"
  !gapprompt|gap>B !gapinput|Print(last);B
  DigraphFromGraph6String(IsMutableDigraph, "IheA@GUAo");
  !gapprompt|gap>B !gapinput|DigraphFromGraph6String(IsMutableDigraph, "IheA@GUAo");B
  <mutable digraph with 10 vertices, 30 edges>
\end{Verbatim}
 }

 

\subsection{\textcolor{Chapter }{DigraphFromGraph6String}}
\logpage{[ 9, 2, 2 ]}\nobreak
\hyperdef{L}{X873D6EB8836A3C61}{}
{\noindent\textcolor{FuncColor}{$\triangleright$\enspace\texttt{DigraphFromGraph6String({\mdseries\slshape [filt, ]str})\index{DigraphFromGraph6String@\texttt{DigraphFromGraph6String}}
\label{DigraphFromGraph6String}
}\hfill{\scriptsize (operation)}}\\
\noindent\textcolor{FuncColor}{$\triangleright$\enspace\texttt{DigraphFromDigraph6String({\mdseries\slshape [filt, ]str})\index{DigraphFromDigraph6String@\texttt{DigraphFromDigraph6String}}
\label{DigraphFromDigraph6String}
}\hfill{\scriptsize (operation)}}\\
\noindent\textcolor{FuncColor}{$\triangleright$\enspace\texttt{DigraphFromSparse6String({\mdseries\slshape [filt, ]str})\index{DigraphFromSparse6String@\texttt{DigraphFromSparse6String}}
\label{DigraphFromSparse6String}
}\hfill{\scriptsize (operation)}}\\
\noindent\textcolor{FuncColor}{$\triangleright$\enspace\texttt{DigraphFromDiSparse6String({\mdseries\slshape [filt, ]str})\index{DigraphFromDiSparse6String@\texttt{DigraphFromDiSparse6String}}
\label{DigraphFromDiSparse6String}
}\hfill{\scriptsize (operation)}}\\
\textbf{\indent Returns:\ }
A digraph.



 If \mbox{\texttt{\mdseries\slshape str}} is a string encoding a graph in Graph6, Digraph6, Sparse6 or DiSparse6 format,
then the corresponding function returns a digraph. In the case of either
Graph6 or Sparse6, formats which do not support directed edges, this will be a
digraph such that for every edge, the edge going in the opposite direction is
also present.

 Each of these functions takes an optional first argument \mbox{\texttt{\mdseries\slshape filt}}, which should be either \texttt{IsMutableDigraph} (\ref{IsMutableDigraph}) or \texttt{IsImmutableDigraph} (\ref{IsImmutableDigraph}), and which specifies whether the output digraph shall be mutable or
immutable. If no first argument is provided, then an immutable digraph is
returned by default. 
\begin{Verbatim}[commandchars=!@|,fontsize=\small,frame=single,label=Example]
  !gapprompt@gap>| !gapinput@DigraphFromGraph6String("?");|
  <immutable empty digraph with 0 vertices>
  !gapprompt@gap>| !gapinput@DigraphFromGraph6String("C]");|
  <immutable symmetric digraph with 4 vertices, 8 edges>
  !gapprompt@gap>| !gapinput@DigraphFromGraph6String("H?AAEM{");|
  <immutable symmetric digraph with 9 vertices, 22 edges>
  !gapprompt@gap>| !gapinput@DigraphFromDigraph6String("&?");|
  <immutable empty digraph with 0 vertices>
  !gapprompt@gap>| !gapinput@DigraphFromDigraph6String(IsMutableDigraph, "&DOOOW?");|
  <mutable digraph with 5 vertices, 5 edges>
  !gapprompt@gap>| !gapinput@DigraphFromDigraph6String("&CQFG");|
  <immutable digraph with 4 vertices, 6 edges>
  !gapprompt@gap>| !gapinput@DigraphFromDigraph6String("&IM[SrKLc~lhesbU[F_");|
  <immutable digraph with 10 vertices, 51 edges>
  !gapprompt@gap>| !gapinput@DigraphFromDiSparse6String(".CaWBGA?b");|
  <immutable multidigraph with 4 vertices, 9 edges>
\end{Verbatim}
 }

 

\subsection{\textcolor{Chapter }{Graph6String}}
\logpage{[ 9, 2, 3 ]}\nobreak
\hyperdef{L}{X782A052979EA2376}{}
{\noindent\textcolor{FuncColor}{$\triangleright$\enspace\texttt{Graph6String({\mdseries\slshape digraph})\index{Graph6String@\texttt{Graph6String}}
\label{Graph6String}
}\hfill{\scriptsize (operation)}}\\
\noindent\textcolor{FuncColor}{$\triangleright$\enspace\texttt{Digraph6String({\mdseries\slshape digraph})\index{Digraph6String@\texttt{Digraph6String}}
\label{Digraph6String}
}\hfill{\scriptsize (operation)}}\\
\noindent\textcolor{FuncColor}{$\triangleright$\enspace\texttt{Sparse6String({\mdseries\slshape digraph})\index{Sparse6String@\texttt{Sparse6String}}
\label{Sparse6String}
}\hfill{\scriptsize (operation)}}\\
\noindent\textcolor{FuncColor}{$\triangleright$\enspace\texttt{DiSparse6String({\mdseries\slshape digraph})\index{DiSparse6String@\texttt{DiSparse6String}}
\label{DiSparse6String}
}\hfill{\scriptsize (operation)}}\\
\textbf{\indent Returns:\ }
A string.



 These four functions return a highly compressed string fully describing the
digraph \mbox{\texttt{\mdseries\slshape digraph}}. 

 Graph6 and Digraph6 are formats best used on small, dense graphs, if
applicable. For larger, sparse graphs use \emph{Sparse6} and \emph{Disparse6} (this latter also preserves multiple edges). 

 See \texttt{WriteDigraphs} (\ref{WriteDigraphs}). 
\begin{Verbatim}[commandchars=!@|,fontsize=\small,frame=single,label=Example]
  !gapprompt@gap>| !gapinput@gr := Digraph([[2, 3], [1], [1]]);|
  <immutable digraph with 3 vertices, 4 edges>
  !gapprompt@gap>| !gapinput@Sparse6String(gr);|
  ":Bc"
  !gapprompt@gap>| !gapinput@DiSparse6String(gr);|
  ".Bc{f"
\end{Verbatim}
 }

 

\subsection{\textcolor{Chapter }{DigraphFromDreadnautString}}
\logpage{[ 9, 2, 4 ]}\nobreak
\hyperdef{L}{X87DD8B1E867760A5}{}
{\noindent\textcolor{FuncColor}{$\triangleright$\enspace\texttt{DigraphFromDreadnautString({\mdseries\slshape s})\index{DigraphFromDreadnautString@\texttt{DigraphFromDreadnautString}}
\label{DigraphFromDreadnautString}
}\hfill{\scriptsize (operation)}}\\
\noindent\textcolor{FuncColor}{$\triangleright$\enspace\texttt{DreadnautString({\mdseries\slshape digraph[, partition]})\index{DreadnautString@\texttt{DreadnautString}}
\label{DreadnautString}
}\hfill{\scriptsize (operation)}}\\


 These operations read or write a single digraph to or from a string in the
dreadnaut format, as appropriate. 

 \texttt{DigraphFromDreadnautString} expects the argument \mbox{\texttt{\mdseries\slshape s}} to be a string, containing a single graph in dreadnaut format, and returns a
digraph.

 If the vertices are not already 1\texttt{\symbol{45}}indexed, their labels
will be reindexed by subtracting a constant offset so that the smallest label
becomes 1, with the relative ordering of vertex labels being preserved. If a
partition is explicitly specified, each vertex will be assigned a label
corresponding to its part in the partition. Should the graph be undirected,
the symmetric closure of the graph will be returned. See \texttt{DigraphVertexLabels} (\ref{DigraphVertexLabels}) and \texttt{DigraphSymmetricClosure} (\ref{DigraphSymmetricClosure}).

 \texttt{DreadnautString} expects a digraph \mbox{\texttt{\mdseries\slshape digraph}} and returns a string in dreadnaut format. An optional second argument \mbox{\texttt{\mdseries\slshape partition}} may be provided, which should be a list of length equal to the number of
vertices in \mbox{\texttt{\mdseries\slshape digraph}}. Each entry in the list assigns the corresponding vertex to a group, allowing
you to specify a vertex partition. For example, if the digraph has three
vertices and \mbox{\texttt{\mdseries\slshape partition}} is \texttt{["a", "b", "a"]}, this would mean that the first and third vertices belong to the same part,
and the second vertex belongs to a separate part. This partition will appear
in the resulting string as a line of the form \texttt{f = [...]}.

 Note that \texttt{WriteDigraphs} (\ref{WriteDigraphs}) does not support specifying a partition when writing a digraph to a file in
dreadnaut format. To include a partition when writing to a file, use \texttt{DreadnautString} to generate the string, then write it to a file using standard I/O functions.
If \mbox{\texttt{\mdseries\slshape digraph}} has multiple edges, then \texttt{DreadnautString} returns a digraph constructed from \mbox{\texttt{\mdseries\slshape digraph}} by removing all multiple edges.

 A detailed description of commands and options in dreadnaut format can be
found at \href{http://pallini.di.uniroma1.it/Guide.html} {nauty and Traces User's Guide}. Of those commands, the following are supported: \texttt{n}, \texttt{g} (and all available subcommands), \texttt{{\textunderscore}} (underscore), \texttt{{\textunderscore}{\textunderscore}} (double underscore), \texttt{f}, \texttt{\$=\#}, \texttt{\$\$}, \texttt{d, \texttt{\symbol{45}}d}, \texttt{"...", !} 
\begin{Verbatim}[commandchars=!@A,fontsize=\small,frame=single,label=Example]
  !gapprompt@gap>A !gapinput@gr := Digraph([[2], [1, 3, 4], [2, 4], [2, 3]]);A
  <immutable digraph with 4 vertices, 8 edges>
  !gapprompt@gap>A !gapinput@s := DreadnautString(gr);;A
  !gapprompt@gap>A !gapinput@DigraphFromDreadnautString(s) = gr;A
  true
  !gapprompt@gap>A !gapinput@DreadnautString(gr, ["a", "b", "b", "a"]);A
  "n=4 $=1 d g\n1 : 2;\n2 : 1 3 4;\n3 : 2 4;\n4 : 2 3.\nf = [1 4 | 2 3]"
\end{Verbatim}
 }

 

\subsection{\textcolor{Chapter }{DIMACSString}}
\logpage{[ 9, 2, 5 ]}\nobreak
\hyperdef{L}{X7923733D86767CE1}{}
{\noindent\textcolor{FuncColor}{$\triangleright$\enspace\texttt{DIMACSString({\mdseries\slshape digraph})\index{DIMACSString@\texttt{DIMACSString}}
\label{DIMACSString}
}\hfill{\scriptsize (operation)}}\\
\noindent\textcolor{FuncColor}{$\triangleright$\enspace\texttt{DigraphFromDIMACSString({\mdseries\slshape s})\index{DigraphFromDIMACSString@\texttt{DigraphFromDIMACSString}}
\label{DigraphFromDIMACSString}
}\hfill{\scriptsize (operation)}}\\
\textbf{\indent Returns:\ }
A string.



 \texttt{DIMACSString} takes a single symmetric digraph \mbox{\texttt{\mdseries\slshape digraph}}, and returns a string representation of \mbox{\texttt{\mdseries\slshape digraph}} in the DIMACS format. 

 \texttt{DigraphFromDIMACSString} takes such a string and returns the single symmetric digraph which it
describes. 

 For more information on the DIMACS format, see \texttt{WriteDIMACSDigraph} (\ref{WriteDIMACSDigraph}).

 These functions support file\texttt{\symbol{45}}based DIMACS encoding and
decoding through \texttt{WriteDigraphs} (\ref{WriteDigraphs}) and \texttt{ReadDigraphs} (\ref{ReadDigraphs}), for consistency with other encoders and decoders. Alternatively, \texttt{WriteDIMACSDigraph} (\ref{WriteDIMACSDigraph}) and \texttt{ReadDIMACSDigraph} (\ref{ReadDIMACSDigraph}) may be used for direct DIMACS I/O. 
\begin{Verbatim}[commandchars=!@|,fontsize=\small,frame=single,label=Example]
  !gapprompt@gap>| !gapinput@gr := Digraph([[2], [1, 3, 4], [2, 4], [2, 3]]);|
  <immutable digraph with 4 vertices, 8 edges>
  !gapprompt@gap>| !gapinput@DIMACSString(gr);|
  "p edge 4 4\ne 1 2\ne 2 3\ne 2 4\ne 3 4\nn 1 1\nn 2 2\nn 3 3\nn 4 4"
  !gapprompt@gap>| !gapinput@DigraphFromDIMACSString(last);|
  <immutable digraph with 4 vertices, 8 edges>
\end{Verbatim}
 }

 

\subsection{\textcolor{Chapter }{DigraphFile}}
\logpage{[ 9, 2, 6 ]}\nobreak
\hyperdef{L}{X7845ACDA7D4D333D}{}
{\noindent\textcolor{FuncColor}{$\triangleright$\enspace\texttt{DigraphFile({\mdseries\slshape filename[, coder][, mode]})\index{DigraphFile@\texttt{DigraphFile}}
\label{DigraphFile}
}\hfill{\scriptsize (function)}}\\
\textbf{\indent Returns:\ }
An IO file object.



 If \mbox{\texttt{\mdseries\slshape filename}} is a string representing the name of a file, then \texttt{DigraphFile} returns an \href{https://gap-packages.github.io/io} {IO}  package file object for that file. 

 If the optional argument \mbox{\texttt{\mdseries\slshape coder}} is specified and is a function which either encodes a digraph as a string, or
decodes a string into a digraph, then this function will be used when reading
or writing to the returned file object. If the optional argument \mbox{\texttt{\mdseries\slshape coder}} is not specified, then the encoding of the digraphs in the returned file
object must be specified in the the file extension. The file extension must be
one of: \texttt{.g6}, \texttt{.s6}, \texttt{.d6}, \texttt{.ds6}, \texttt{.txt}, \texttt{.p}, or \texttt{.pickle}; more details of these file formats is given below. 

 If the optional argument \mbox{\texttt{\mdseries\slshape mode}} is specified, then it must be one of: \texttt{"w"} (for write), \texttt{"a"} (for append), or \texttt{"r"} (for read). If \mbox{\texttt{\mdseries\slshape mode}} is not specified, then \texttt{"r"} is used by default. 

 If \mbox{\texttt{\mdseries\slshape filename}} ends in one of: \texttt{.gz}, \texttt{.bz2}, or \texttt{.xz}, then the digraphs which are read from, or written to, the returned file
object are decompressed, or compressed, appropriately. 

 The file object returned by \texttt{DigraphFile} can be given as the first argument for either of the functions \texttt{ReadDigraphs} (\ref{ReadDigraphs}) or \texttt{WriteDigraphs} (\ref{WriteDigraphs}). The purpose of this is to reduce the overhead of recreating the file object
inside the functions \texttt{ReadDigraphs} (\ref{ReadDigraphs}) or \texttt{WriteDigraphs} (\ref{WriteDigraphs}) when, for example, reading or writing many digraphs in a loop. 

 The currently supported file formats, and associated filename extensions, are: 
\begin{description}
\item[{graph6 (.g6)}]  A standard and widely\texttt{\symbol{45}}used format for undirected graphs,
with no support for loops or multiple edges. Only symmetric graphs are allowed
\texttt{\symbol{45}}\texttt{\symbol{45}} each edge is combined with its
converse edge to produce a single undirected edge. This format is best used
for "dense" graphs \texttt{\symbol{45}}\texttt{\symbol{45}} those with many
edges per vertex. 
\item[{sparse6 (.s6)}]  Unlike graph6, sparse6 has support for loops and multiple edges. However, its
use is still limited to symmetric graphs. This format is
better\texttt{\symbol{45}}suited to "sparse" graphs
\texttt{\symbol{45}}\texttt{\symbol{45}} those with few edges per vertex. 
\item[{digraph6 (.d6)}]  This format is based on graph6, but stores direction information
\texttt{\symbol{45}} therefore is not limited to symmetric graphs. Loops are
allowed, but multiple edges are not. Best compression with "dense" graphs. 
\item[{disparse6 (.ds6)}]  Any type of digraph can be encoded in disparse6: directions, loops, and
multiple edges are all allowed. Similar to sparse6, this has the best
compression rate with "sparse" graphs. 
\item[{plain text (.txt)}]  This is a human\texttt{\symbol{45}}readable format which stores graphs in the
form \texttt{0 7 0 8 1 7 2 8 3 8 4 8 5 8 6 8} i.e. pairs of vertices describing edges in a graph. More specifically, the
vertices making up one edge must be separated by a single space, and pairs of
vertices must be separated by two spaces. 

 See \texttt{ReadPlainTextDigraph} (\ref{ReadPlainTextDigraph}) for a more flexible way to store digraphs in a plain text file. 

 
\item[{pickled (\texttt{.p} or \texttt{.pickle})}]  Digraphs are pickled using the \href{https://gap-packages.github.io/io} {IO}  package. This is particularly good when the \texttt{DigraphGroup} (\ref{DigraphGroup}) is non\texttt{\symbol{45}}trivial. 
\item[{dreadnaut (.dre)}]  A graph format designed for directed and undirected graphs. The format
supports loops but multiple edges are ignored. The format consists of an
initial section that defines the graph's structural properties, such as the
number of vertices, the starting value for vertices, and whether the graph is
directed. This is followed by a list of edges. For more information and
examples of the format see \href{http://pallini.di.uniroma1.it/Guide.html} {nauty and Traces User's Guide}.

 
\item[{ DIMACS (.dimacs) }]  A graph format that can be used for symmetric digraphs. For a more detailed
description, see \texttt{WriteDIMACSDigraph} (\ref{WriteDIMACSDigraph}) 
\end{description}
 
\begin{Verbatim}[commandchars=!@|,fontsize=\small,frame=single,label=Example]
  !gapprompt@gap>| !gapinput@filename := Concatenation(DIGRAPHS_Dir(), "/tst/out/man.d6.gz");;|
  !gapprompt@gap>| !gapinput@file := DigraphFile(filename, "w");;|
  !gapprompt@gap>| !gapinput@for i in [1 .. 10] do|
  !gapprompt@>| !gapinput@WriteDigraphs(file, Digraph([[1, 3], [2], [1, 2]]));|
  !gapprompt@>| !gapinput@od;|
  !gapprompt@gap>| !gapinput@IO_Close(file);;|
  !gapprompt@gap>| !gapinput@file := DigraphFile(filename, "r");;|
  !gapprompt@gap>| !gapinput@ReadDigraphs(file, 9);|
  <immutable digraph with 3 vertices, 5 edges>
  !gapprompt@gap>| !gapinput@IO_Close(file);;|
\end{Verbatim}
 }

 

\subsection{\textcolor{Chapter }{ReadDigraphs}}
\logpage{[ 9, 2, 7 ]}\nobreak
\hyperdef{L}{X7E8142007ED0CB9F}{}
{\noindent\textcolor{FuncColor}{$\triangleright$\enspace\texttt{ReadDigraphs({\mdseries\slshape filename[, decoder][, n]})\index{ReadDigraphs@\texttt{ReadDigraphs}}
\label{ReadDigraphs}
}\hfill{\scriptsize (function)}}\\
\textbf{\indent Returns:\ }
A digraph, or a list of digraphs.



 If \mbox{\texttt{\mdseries\slshape filename}} is a string containing the name of a file containing encoded digraphs or an \href{https://gap-packages.github.io/io} {IO}  file object created using \texttt{DigraphFile} (\ref{DigraphFile}), then \texttt{ReadDigraphs} returns the digraphs encoded in the file as a list. Note that if \mbox{\texttt{\mdseries\slshape filename}} is a compressed file, which has been compressed appropriately to give a
filename extension of \texttt{.gz}, \texttt{.bz2}, or \texttt{.xz}, then \texttt{ReadDigraphs} can read \mbox{\texttt{\mdseries\slshape filename}} without it first needing to be decompressed. 

 If the optional argument \mbox{\texttt{\mdseries\slshape decoder}} is specified and is a function which decodes a string into a digraph, then \texttt{ReadDigraphs} will use \mbox{\texttt{\mdseries\slshape decoder}} to decode the digraphs contained in \mbox{\texttt{\mdseries\slshape filename}}.

 If the optional argument \mbox{\texttt{\mdseries\slshape n}} is specified, then \texttt{ReadDigraphs} returns the \mbox{\texttt{\mdseries\slshape n}}th digraph encoded in the file \mbox{\texttt{\mdseries\slshape filename}}. 

 If the optional argument \mbox{\texttt{\mdseries\slshape decoder}} is not specified, then \texttt{ReadDigraphs} will deduce which decoder to use based on the filename extension of \mbox{\texttt{\mdseries\slshape filename}} (after removing the compression\texttt{\symbol{45}}related filename extensions \texttt{.gz}, \texttt{.bz2}, and \texttt{.xz}). For example, if the filename extension is \texttt{.g6}, then \texttt{ReadDigraphs} will use the graph6 decoder \texttt{DigraphFromGraph6String} (\ref{DigraphFromGraph6String}). 

 The currently supported file formats, and associated filename extensions, are: 
\begin{description}
\item[{graph6 (.g6)}]  A standard and widely\texttt{\symbol{45}}used format for undirected graphs,
with no support for loops or multiple edges. Only symmetric graphs are allowed
\texttt{\symbol{45}}\texttt{\symbol{45}} each edge is combined with its
converse edge to produce a single undirected edge. This format is best used
for "dense" graphs \texttt{\symbol{45}}\texttt{\symbol{45}} those with many
edges per vertex. 
\item[{sparse6 (.s6)}]  Unlike graph6, sparse6 has support for loops and multiple edges. However, its
use is still limited to symmetric graphs. This format is
better\texttt{\symbol{45}}suited to "sparse" graphs
\texttt{\symbol{45}}\texttt{\symbol{45}} those with few edges per vertex. 
\item[{digraph6 (.d6)}]  This format is based on graph6, but stores direction information
\texttt{\symbol{45}} therefore is not limited to symmetric graphs. Loops are
allowed, but multiple edges are not. Best compression with "dense" graphs. 
\item[{disparse6 (.ds6)}]  Any type of digraph can be encoded in disparse6: directions, loops, and
multiple edges are all allowed. Similar to sparse6, this has the best
compression rate with "sparse" graphs. 
\item[{plain text (.txt)}]  This is a human\texttt{\symbol{45}}readable format which stores graphs in the
form \texttt{0 7 0 8 1 7 2 8 3 8 4 8 5 8 6 8} i.e. pairs of vertices describing edges in a graph. More specifically, the
vertices making up one edge must be separated by a single space, and pairs of
vertices must be separated by two spaces. 

 See \texttt{ReadPlainTextDigraph} (\ref{ReadPlainTextDigraph}) for a more flexible way to store digraphs in a plain text file. 

  
\item[{pickled (\texttt{.p} or \texttt{.pickle})}]  Digraphs are pickled using the \href{https://gap-packages.github.io/io} {IO}  package. This is particularly good when the \texttt{DigraphGroup} (\ref{DigraphGroup}) is non\texttt{\symbol{45}}trivial. 

 
\item[{dreadnaut (.dre)}]  A graph format designed for directed and undirected graphs. The format
supports loops but multiple edges are ignored. The format consists of an
initial section that defines the graph's structural properties, such as the
number of vertices, the starting value for vertices, and whether the graph is
directed. This is followed by a list of edges. For more information and
examples of the format see \href{http://pallini.di.uniroma1.it/Guide.html} {nauty and Traces User's Guide}. 

 
\item[{ DIMACS (.dimacs) }]  A graph format that can be used for symmetric digraphs. For a more detailed
description, see \texttt{WriteDIMACSDigraph} (\ref{WriteDIMACSDigraph}) 
\end{description}
 
\begin{Verbatim}[commandchars=!@|,fontsize=\small,frame=single,label=Example]
  !gapprompt@gap>| !gapinput@ReadDigraphs(|
  !gapprompt@>| !gapinput@Concatenation(DIGRAPHS_Dir(), "/data/graph5.g6.gz"), 10);|
  <immutable symmetric digraph with 5 vertices, 8 edges>
  !gapprompt@gap>| !gapinput@ReadDigraphs(|
  !gapprompt@>| !gapinput@Concatenation(DIGRAPHS_Dir(), "/data/graph5.g6.gz"), 17);|
  <immutable symmetric digraph with 5 vertices, 12 edges>
  !gapprompt@gap>| !gapinput@ReadDigraphs(|
  !gapprompt@>| !gapinput@Concatenation(DIGRAPHS_Dir(), "/data/tree9.4.txt"));|
  [ <immutable digraph with 9 vertices, 8 edges>, 
    <immutable digraph with 9 vertices, 8 edges>, 
    <immutable digraph with 9 vertices, 8 edges>, 
    <immutable digraph with 9 vertices, 8 edges>, 
    <immutable digraph with 9 vertices, 8 edges>, 
    <immutable digraph with 9 vertices, 8 edges>, 
    <immutable digraph with 9 vertices, 8 edges>, 
    <immutable digraph with 9 vertices, 8 edges>, 
    <immutable digraph with 9 vertices, 8 edges>, 
    <immutable digraph with 9 vertices, 8 edges>, 
    <immutable digraph with 9 vertices, 8 edges>, 
    <immutable digraph with 9 vertices, 8 edges>, 
    <immutable digraph with 9 vertices, 8 edges>, 
    <immutable digraph with 9 vertices, 8 edges> ]
\end{Verbatim}
 }

 

\subsection{\textcolor{Chapter }{WriteDigraphs}}
\logpage{[ 9, 2, 8 ]}\nobreak
\hyperdef{L}{X82A4F0767C9EFC87}{}
{\noindent\textcolor{FuncColor}{$\triangleright$\enspace\texttt{WriteDigraphs({\mdseries\slshape filename, digraphs[, encoder][, mode]})\index{WriteDigraphs@\texttt{WriteDigraphs}}
\label{WriteDigraphs}
}\hfill{\scriptsize (function)}}\\


 If \mbox{\texttt{\mdseries\slshape digraphs}} is a list of digraphs or a digraph and \mbox{\texttt{\mdseries\slshape filename}} is a string or an \href{https://gap-packages.github.io/io} {IO}  file object created using \texttt{DigraphFile} (\ref{DigraphFile}), then \texttt{WriteDigraphs} writes the digraphs to the file represented by \mbox{\texttt{\mdseries\slshape filename}}. If the supplied filename ends in one of the extensions \texttt{.gz}, \texttt{.bz2}, or \texttt{.xz}, then the file will be compressed appropriately. Excluding these extensions,
if the file ends with an extension in the list below, the corresponding graph
format will be used to encode it. If such an extension is not included, an
appropriate format will be chosen intelligently, and an extension appended, to
minimise file size. 

  For more verbose information on the progress of the function, set the info
level of \mbox{\texttt{\mdseries\slshape InfoDigraphs}} to 1 or higher, using \texttt{SetInfoLevel}.

 The currently supported file formats are: 
\begin{description}
\item[{graph6 (.g6)}]  A standard and widely\texttt{\symbol{45}}used format for undirected graphs,
with no support for loops or multiple edges. Only symmetric graphs are allowed
\texttt{\symbol{45}}\texttt{\symbol{45}} each edge is combined with its
converse edge to produce a single undirected edge. This format is best used
for "dense" graphs \texttt{\symbol{45}}\texttt{\symbol{45}} those with many
edges per vertex. 
\item[{sparse6 (.s6)}]  Unlike graph6, sparse6 has support for loops and multiple edges. However, its
use is still limited to symmetric graphs. This format is
better\texttt{\symbol{45}}suited to "sparse" graphs
\texttt{\symbol{45}}\texttt{\symbol{45}} those with few edges per vertex. 
\item[{digraph6 (.d6)}]  This format is based on graph6, but stores direction information
\texttt{\symbol{45}} therefore is not limited to symmetric graphs. Loops are
allowed, but multiple edges are not. Best compression with "dense" graphs. 
\item[{disparse6 (.ds6)}]  Any type of digraph can be encoded in disparse6: directions, loops, and
multiple edges are all allowed. Similar to sparse6, this has the best
compression rate with "sparse" graphs. 
\item[{plain text (.txt)}]  This is a human\texttt{\symbol{45}}readable format which stores graphs in the
form \texttt{0 7 0 8 1 7 2 8 3 8 4 8 5 8 6 8} i.e. pairs of vertices describing edges in a graph. More specifically, the
vertices making up one edge must be separated by a single space, and pairs of
vertices must be separated by two spaces. 

 See \texttt{ReadPlainTextDigraph} (\ref{ReadPlainTextDigraph}) for a more flexible way to store digraphs in a plain text file. 

  
\item[{pickled (\texttt{.p} or \texttt{.pickle})}]  Digraphs are pickled using the \href{https://gap-packages.github.io/io} {IO}  package. This is particularly good when the \texttt{DigraphGroup} (\ref{DigraphGroup}) is non\texttt{\symbol{45}}trivial. 
\item[{dreadnaut (.dre)}]  A graph format designed for directed and undirected graphs. The format
supports loops but multiple edges are ignored. The format consists of an
initial section that defines the graph's structural properties, such as the
number of vertices, the starting value for vertices, and whether the graph is
directed. This is followed by a list of edges. For more information and
examples of the format see \href{http://pallini.di.uniroma1.it/Guide.html} {nauty and Traces User's Guide}. 

 
\item[{ DIMACS (.dimacs) }]  A graph format that can be used for symmetric digraphs. For a more detailed
description, see \texttt{WriteDIMACSDigraph} (\ref{WriteDIMACSDigraph}) 
\end{description}
 
\begin{Verbatim}[commandchars=!@|,fontsize=\small,frame=single,label=Example]
  !gapprompt@gap>| !gapinput@grs := [];;|
  !gapprompt@gap>| !gapinput@grs[1] := Digraph([]);|
  <immutable empty digraph with 0 vertices>
  !gapprompt@gap>| !gapinput@grs[2] := Digraph([[1, 3], [2], [1, 2]]);|
  <immutable digraph with 3 vertices, 5 edges>
  !gapprompt@gap>| !gapinput@grs[3] := Digraph([|
  !gapprompt@>| !gapinput@[6, 7], [6, 9], [1, 3, 4, 5, 8, 9],|
  !gapprompt@>| !gapinput@[1, 2, 3, 4, 5, 6, 7, 10], [1, 5, 6, 7, 10], [2, 4, 5, 9, 10],|
  !gapprompt@>| !gapinput@[3, 4, 5, 6, 7, 8, 9, 10], [1, 3, 5, 7, 8, 9], [1, 2, 5],|
  !gapprompt@>| !gapinput@[1, 2, 4, 6, 7, 8]]);|
  <immutable digraph with 10 vertices, 51 edges>
  !gapprompt@gap>| !gapinput@filename := Concatenation(DIGRAPHS_Dir(), "/tst/out/man.d6.gz");;|
  !gapprompt@gap>| !gapinput@WriteDigraphs(filename, grs, "w");|
  IO_OK
  !gapprompt@gap>| !gapinput@ReadDigraphs(filename);|
  [ <immutable empty digraph with 0 vertices>, 
    <immutable digraph with 3 vertices, 5 edges>, 
    <immutable digraph with 10 vertices, 51 edges> ]
\end{Verbatim}
 }

 

\subsection{\textcolor{Chapter }{IteratorFromDigraphFile}}
\logpage{[ 9, 2, 9 ]}\nobreak
\hyperdef{L}{X86D531B287733561}{}
{\noindent\textcolor{FuncColor}{$\triangleright$\enspace\texttt{IteratorFromDigraphFile({\mdseries\slshape filename[, decoder]})\index{IteratorFromDigraphFile@\texttt{IteratorFromDigraphFile}}
\label{IteratorFromDigraphFile}
}\hfill{\scriptsize (function)}}\\
\textbf{\indent Returns:\ }
An iterator.



 If \mbox{\texttt{\mdseries\slshape filename}} is a string representing the name of a file containing encoded digraphs, then \texttt{IteratorFromDigraphFile} returns an iterator for which the value of \texttt{NextIterator} (\textbf{Reference: NextIterator}) is the next digraph encoded in the file. 

 If the optional argument \mbox{\texttt{\mdseries\slshape decoder}} is specified and is a function which decodes a string into a digraph, then \texttt{IteratorFromDigraphFile} will use \mbox{\texttt{\mdseries\slshape decoder}} to decode the digraphs contained in \mbox{\texttt{\mdseries\slshape filename}}. 

 The purpose of this function is to easily allow looping over digraphs encoded
in a file when loading all of the encoded digraphs would require too much
memory. 

 To see what file types are available, see \texttt{WriteDigraphs} (\ref{WriteDigraphs}). 
\begin{Verbatim}[commandchars=!@|,fontsize=\small,frame=single,label=Example]
  !gapprompt@gap>| !gapinput@filename := Concatenation(DIGRAPHS_Dir(), "/tst/out/man.d6.gz");;|
  !gapprompt@gap>| !gapinput@file := DigraphFile(filename, "w");;|
  !gapprompt@gap>| !gapinput@for i in [1 .. 10] do|
  !gapprompt@>| !gapinput@  WriteDigraphs(file, Digraph([[1, 3], [2], [1, 2]]));|
  !gapprompt@>| !gapinput@od;|
  !gapprompt@gap>| !gapinput@IO_Close(file);;|
  !gapprompt@gap>| !gapinput@iter := IteratorFromDigraphFile(filename);|
  <iterator>
  !gapprompt@gap>| !gapinput@for x in iter do od;|
\end{Verbatim}
 }

 

\subsection{\textcolor{Chapter }{DigraphPlainTextLineEncoder}}
\logpage{[ 9, 2, 10 ]}\nobreak
\hyperdef{L}{X7AAA35867B5D8BF4}{}
{\noindent\textcolor{FuncColor}{$\triangleright$\enspace\texttt{DigraphPlainTextLineEncoder({\mdseries\slshape delimiter1[, delimiter2], offset})\index{DigraphPlainTextLineEncoder@\texttt{DigraphPlainTextLineEncoder}}
\label{DigraphPlainTextLineEncoder}
}\hfill{\scriptsize (function)}}\\
\noindent\textcolor{FuncColor}{$\triangleright$\enspace\texttt{DigraphPlainTextLineDecoder({\mdseries\slshape delimiter1[, delimiter2], offset})\index{DigraphPlainTextLineDecoder@\texttt{DigraphPlainTextLineDecoder}}
\label{DigraphPlainTextLineDecoder}
}\hfill{\scriptsize (operation)}}\\
\textbf{\indent Returns:\ }
A string.



 These two functions return a function which encodes or decodes a digraph in a
plain text format.

 \mbox{\texttt{\mdseries\slshape DigraphPlainTextLineEncoder}} returns a function which takes a single digraph as an argument. The function
returns a string describing the edges of that digraph; each edge is written as
a pair of integers separated by the string \mbox{\texttt{\mdseries\slshape delimiter2}}, and the edges themselves are separated by the string \mbox{\texttt{\mdseries\slshape delimiter1}}. \mbox{\texttt{\mdseries\slshape DigraphPlainTextLineDecoder}} returns the corresponding decoder function, which takes a string argument in
this format and returns a digraph.

 If only one delimiter is passed as an argument to \mbox{\texttt{\mdseries\slshape DigraphPlainTextLineDecoder}}, it will return a function which decodes a single edge, returning its
contents as a list of integers.

 The argument \mbox{\texttt{\mdseries\slshape offset}} should be an integer, which will describe a number to be added to each vertex
before it is encoded, or after it is decoded. This may be used, for example,
to label vertices starting at 0 instead of 1.

 Note that the number of vertices of a digraph is not stored, and so vertices
which are not connected to any edge may be lost. 
\begin{Verbatim}[commandchars=!@|,fontsize=\small,frame=single,label=Example]
  !gapprompt@gap>| !gapinput@gr := Digraph([[2, 3], [1], [1]]);|
  <immutable digraph with 3 vertices, 4 edges>
  !gapprompt@gap>| !gapinput@enc := DigraphPlainTextLineEncoder("  ", " ", -1);;|
  !gapprompt@gap>| !gapinput@dec := DigraphPlainTextLineDecoder("  ", " ", 1);;|
  !gapprompt@gap>| !gapinput@enc(gr);|
  "0 1  0 2  1 0  2 0"
  !gapprompt@gap>| !gapinput@dec(last);|
  <immutable digraph with 3 vertices, 4 edges>
\end{Verbatim}
 }

 

\subsection{\textcolor{Chapter }{TournamentLineDecoder}}
\logpage{[ 9, 2, 11 ]}\nobreak
\hyperdef{L}{X7808D99079295C22}{}
{\noindent\textcolor{FuncColor}{$\triangleright$\enspace\texttt{TournamentLineDecoder({\mdseries\slshape str})\index{TournamentLineDecoder@\texttt{TournamentLineDecoder}}
\label{TournamentLineDecoder}
}\hfill{\scriptsize (operation)}}\\
\textbf{\indent Returns:\ }
A digraph.



 This function takes a string \mbox{\texttt{\mdseries\slshape str}}, decodes it, and then returns the tournament [see \texttt{IsTournament} (\ref{IsTournament})] which it defines, according to the following rules. 

 The characters of the string \mbox{\texttt{\mdseries\slshape str}} represent the entries in the upper triangle of a tournament's adjacency
matrix. The number of vertices \texttt{n} will be detected from the length of the string and will be as large as
possible. 

 The first character represents the possible edge \texttt{1 \texttt{\symbol{45}}{\textgreater} 2}, the second represents \texttt{1 \texttt{\symbol{45}}{\textgreater} 3} and so on until \texttt{1 \texttt{\symbol{45}}{\textgreater} n}; then the following character represents \texttt{2 \texttt{\symbol{45}}{\textgreater} 3}, and so on up to the character which represents the edge \texttt{n\texttt{\symbol{45}}1 \texttt{\symbol{45}}{\textgreater} n}. 

 If a character of the string with corresponding edge \texttt{i \texttt{\symbol{45}}{\textgreater} j} is equal to \texttt{1}, then the edge \texttt{i \texttt{\symbol{45}}{\textgreater} j} is present in the tournament. Otherwise, the edge \texttt{i \texttt{\symbol{45}}{\textgreater} j} is present instead. In this way, all the possible edges are encoded
one\texttt{\symbol{45}}by\texttt{\symbol{45}}one. 
\begin{Verbatim}[commandchars=!@|,fontsize=\small,frame=single,label=Example]
  !gapprompt@gap>| !gapinput@gr := TournamentLineDecoder("100001");|
  <immutable digraph with 4 vertices, 6 edges>
  !gapprompt@gap>| !gapinput@OutNeighbours(gr);|
  [ [ 2 ], [  ], [ 1, 2, 4 ], [ 1, 2 ] ]
\end{Verbatim}
 }

 

\subsection{\textcolor{Chapter }{AdjacencyMatrixUpperTriangleLineDecoder}}
\logpage{[ 9, 2, 12 ]}\nobreak
\hyperdef{L}{X7DF22D1687D096D6}{}
{\noindent\textcolor{FuncColor}{$\triangleright$\enspace\texttt{AdjacencyMatrixUpperTriangleLineDecoder({\mdseries\slshape str})\index{AdjacencyMatrixUpperTriangleLineDecoder@\texttt{Adjacency}\-\texttt{Matrix}\-\texttt{Upper}\-\texttt{Triangle}\-\texttt{Line}\-\texttt{Decoder}}
\label{AdjacencyMatrixUpperTriangleLineDecoder}
}\hfill{\scriptsize (operation)}}\\
\textbf{\indent Returns:\ }
A digraph.



 This function takes a string \mbox{\texttt{\mdseries\slshape str}}, decodes it, and then returns the topologically sorted digraph [see \texttt{DigraphTopologicalSort} (\ref{DigraphTopologicalSort})] which it defines, according to the following rules. 

 The characters of the string \mbox{\texttt{\mdseries\slshape str}} represent the entries in the upper triangle of a digraph's adjacency matrix.
The number of vertices \texttt{n} will be detected from the length of the string and will be as large as
possible. 

 The first character represents the possible edge \texttt{1 \texttt{\symbol{45}}{\textgreater} 2}, the second represents \texttt{1 \texttt{\symbol{45}}{\textgreater} 3} and so on until \texttt{1 \texttt{\symbol{45}}{\textgreater} n}; then the following character represents \texttt{2 \texttt{\symbol{45}}{\textgreater} 3}, and so on up to the character which represents the edge \texttt{n\texttt{\symbol{45}}1 \texttt{\symbol{45}}{\textgreater} n}. If a character of the string with corresponding edge \texttt{i \texttt{\symbol{45}}{\textgreater} j} is equal to \texttt{1}, then this edge is present in the digraph. Otherwise, it is not present. In
this way, all the possible edges are encoded
one\texttt{\symbol{45}}by\texttt{\symbol{45}}one. 

 In particular, note that there exists no edge \texttt{[i, j]} if $j \leq i$. In order words, the digraph will be topologically sorted. 
\begin{Verbatim}[commandchars=!@|,fontsize=\small,frame=single,label=Example]
  !gapprompt@gap>| !gapinput@gr := AdjacencyMatrixUpperTriangleLineDecoder("100001");|
  <immutable digraph with 4 vertices, 2 edges>
  !gapprompt@gap>| !gapinput@OutNeighbours(gr);|
  [ [ 2 ], [  ], [ 4 ], [  ] ]
  !gapprompt@gap>| !gapinput@gr := AdjacencyMatrixUpperTriangleLineDecoder("111111x111");|
  <immutable digraph with 5 vertices, 9 edges>
  !gapprompt@gap>| !gapinput@OutNeighbours(gr);|
  [ [ 2, 3, 4, 5 ], [ 3, 4 ], [ 4, 5 ], [ 5 ], [  ] ]
\end{Verbatim}
 }

 

\subsection{\textcolor{Chapter }{TCodeDecoder}}
\logpage{[ 9, 2, 13 ]}\nobreak
\hyperdef{L}{X875393338486D77D}{}
{\noindent\textcolor{FuncColor}{$\triangleright$\enspace\texttt{TCodeDecoder({\mdseries\slshape str})\index{TCodeDecoder@\texttt{TCodeDecoder}}
\label{TCodeDecoder}
}\hfill{\scriptsize (operation)}}\\
\textbf{\indent Returns:\ }
A digraph.



 If \mbox{\texttt{\mdseries\slshape str}} is a string consisting of at least two non\texttt{\symbol{45}}negative
integers separated by spaces, then this function will attempt to return the
digraph which it defines as a TCode string. 

 The first integer of the string defines the number of vertices \texttt{v} in the digraph, and the second defines the number of edges \texttt{e}. The following \texttt{2e} integers should be vertex numbers in the range \texttt{[0 .. v\texttt{\symbol{45}}1]}. These integers are read in pairs and define the digraph's edges. This
function will return an error if \mbox{\texttt{\mdseries\slshape str}} has fewer than \texttt{2e+2} entries. 

 Note that the vertex numbers will be incremented by 1 in the digraph returned.
Hence the string fragment \texttt{0 6} will describe the edge \texttt{[1,7]}. 
\begin{Verbatim}[commandchars=!@|,fontsize=\small,frame=single,label=Example]
  !gapprompt@gap>| !gapinput@gr := TCodeDecoder("3 2 0 2 2 1");|
  <immutable digraph with 3 vertices, 2 edges>
  !gapprompt@gap>| !gapinput@OutNeighbours(gr);|
  [ [ 3 ], [  ], [ 2 ] ]
  !gapprompt@gap>| !gapinput@gr := TCodeDecoder("12 3 0 10 5 2 8 8");|
  <immutable digraph with 12 vertices, 3 edges>
  !gapprompt@gap>| !gapinput@OutNeighbours(gr);|
  [ [ 11 ], [  ], [  ], [  ], [  ], [ 3 ], [  ], [  ], [ 9 ], [  ], 
    [  ], [  ] ]
\end{Verbatim}
 }

 

\subsection{\textcolor{Chapter }{PlainTextString}}
\logpage{[ 9, 2, 14 ]}\nobreak
\hyperdef{L}{X8540FC21796A793F}{}
{\noindent\textcolor{FuncColor}{$\triangleright$\enspace\texttt{PlainTextString({\mdseries\slshape digraph})\index{PlainTextString@\texttt{PlainTextString}}
\label{PlainTextString}
}\hfill{\scriptsize (operation)}}\\
\noindent\textcolor{FuncColor}{$\triangleright$\enspace\texttt{DigraphFromPlainTextString({\mdseries\slshape s})\index{DigraphFromPlainTextString@\texttt{DigraphFromPlainTextString}}
\label{DigraphFromPlainTextString}
}\hfill{\scriptsize (operation)}}\\
\textbf{\indent Returns:\ }
A string.



 \mbox{\texttt{\mdseries\slshape PlainTextString}} takes a single digraph, and returns a string describing the edges of that
digraph. \mbox{\texttt{\mdseries\slshape DigraphFromPlainTextString}} takes such a string and returns the digraph which it describes. Each edge is
written as a pair of integers separated by a single space. The edges
themselves are separated by a double space. Vertex numbers are reduced by 1
when they are encoded, so that vertices in the string are labelled starting at
0.

 Note that the number of vertices of a digraph is not stored, and so vertices
which are not connected to any edge may be lost.

 The operation \texttt{DigraphFromPlainTextString} takes an optional first argument \texttt{IsMutableDigraph} (\ref{IsMutableDigraph}) or \texttt{IsImmutableDigraph} (\ref{IsImmutableDigraph}), which specifies whether the output digraph shall be mutable or immutable. If
no first argument is provided, then an immutable digraph is returned by
default. 
\begin{Verbatim}[commandchars=!@|,fontsize=\small,frame=single,label=Example]
  !gapprompt@gap>| !gapinput@gr := Digraph([[2, 3], [1], [1]]);|
  <immutable digraph with 3 vertices, 4 edges>
  !gapprompt@gap>| !gapinput@PlainTextString(gr);|
  "0 1  0 2  1 0  2 0"
  !gapprompt@gap>| !gapinput@DigraphFromPlainTextString(last);|
  <immutable digraph with 3 vertices, 4 edges>
\end{Verbatim}
 }

 

\subsection{\textcolor{Chapter }{WritePlainTextDigraph}}
\logpage{[ 9, 2, 15 ]}\nobreak
\hyperdef{L}{X81BD23697857B244}{}
{\noindent\textcolor{FuncColor}{$\triangleright$\enspace\texttt{WritePlainTextDigraph({\mdseries\slshape filename, digraph, delimiter, offset})\index{WritePlainTextDigraph@\texttt{WritePlainTextDigraph}}
\label{WritePlainTextDigraph}
}\hfill{\scriptsize (function)}}\\
\noindent\textcolor{FuncColor}{$\triangleright$\enspace\texttt{ReadPlainTextDigraph({\mdseries\slshape filename, delimiter, offset, ignore})\index{ReadPlainTextDigraph@\texttt{ReadPlainTextDigraph}}
\label{ReadPlainTextDigraph}
}\hfill{\scriptsize (operation)}}\\


 These functions write and read a single digraph in a
human\texttt{\symbol{45}}readable plain text format as follows: each line
contains a single edge, and each edge is written as a pair of integers
separated by the string \mbox{\texttt{\mdseries\slshape delimiter}}.

 \mbox{\texttt{\mdseries\slshape filename}} should be the name of a file which will be written to or read from, and \mbox{\texttt{\mdseries\slshape offset}} should be an integer which is added to each vertex number as it is written or
read. For example, if \texttt{WritePlainTextDigraph} is called with \mbox{\texttt{\mdseries\slshape offset}} \texttt{\texttt{\symbol{45}}1}, then the vertices will be numbered in the file starting from 0 instead of 1
\texttt{\symbol{45}} \texttt{ReadPlainTextDigraph} would then need to be called with \mbox{\texttt{\mdseries\slshape offset}} \texttt{1} to convert back to the original graph.

 \mbox{\texttt{\mdseries\slshape ignore}} should be a list of characters which will be ignored when reading the graph. 
\begin{Verbatim}[commandchars=!@|,fontsize=\small,frame=single,label=Example]
  !gapprompt@gap>| !gapinput@gr := Digraph([[1, 2, 3], [1, 1], [2]]);|
  <immutable multidigraph with 3 vertices, 6 edges>
  !gapprompt@gap>| !gapinput@filename := Concatenation(DIGRAPHS_Dir(), "/tst/out/plain.txt");;|
  !gapprompt@gap>| !gapinput@WritePlainTextDigraph(filename, gr, ",", -1);|
  !gapprompt@gap>| !gapinput@ReadPlainTextDigraph(filename, ",", 1, ['/', '%']);|
  <immutable multidigraph with 3 vertices, 6 edges>
\end{Verbatim}
 }

 

\subsection{\textcolor{Chapter }{WriteDIMACSDigraph}}
\logpage{[ 9, 2, 16 ]}\nobreak
\hyperdef{L}{X804FDF6A7ED64049}{}
{\noindent\textcolor{FuncColor}{$\triangleright$\enspace\texttt{WriteDIMACSDigraph({\mdseries\slshape filename, digraph})\index{WriteDIMACSDigraph@\texttt{WriteDIMACSDigraph}}
\label{WriteDIMACSDigraph}
}\hfill{\scriptsize (operation)}}\\
\noindent\textcolor{FuncColor}{$\triangleright$\enspace\texttt{ReadDIMACSDigraph({\mdseries\slshape filename})\index{ReadDIMACSDigraph@\texttt{ReadDIMACSDigraph}}
\label{ReadDIMACSDigraph}
}\hfill{\scriptsize (operation)}}\\


 These operations write or read the single symmetric digraph \mbox{\texttt{\mdseries\slshape digraph}} to or from a file in DIMACS format, as appropriate. The operation \texttt{WriteDIMACSDigraph} records the vertices and edges of \mbox{\texttt{\mdseries\slshape digraph}}. The vertex labels of \mbox{\texttt{\mdseries\slshape digraph}} will be recorded only if they are integers. See \texttt{IsSymmetricDigraph} (\ref{IsSymmetricDigraph}) and \texttt{DigraphVertexLabels} (\ref{DigraphVertexLabels}).

 The first argument \mbox{\texttt{\mdseries\slshape filename}} should be the name of the file which will be written to or read from. A file
can contain one symmetric digraph in DIMACS format. If \mbox{\texttt{\mdseries\slshape filename}} ends in one of \texttt{.gz}, \texttt{.bz2}, or \texttt{.xz}, then the file is compressed, or decompressed, appropriately. 

 The DIMACS format is described as follows. Each line in the DIMACS file has
one of four types: 
\begin{itemize}
\item  A line beginning with \texttt{c} and followed by any number of characters is a comment line, and is ignored. 
\item  A line beginning with \texttt{p} defines the numbers of vertices and edges the digraph. This line has the
format \texttt{p edge {\textless}nr{\textunderscore}vertices{\textgreater}
{\textless}nr{\textunderscore}edges{\textgreater}}, where \texttt{{\textless}nr{\textunderscore}vertices{\textgreater}} and \texttt{{\textless}nr{\textunderscore}edges{\textgreater}} are replaced by the relevant integers. There must be exactly one such line in
the file, and it must occur before any of the following kinds of line.

 Although it is required to be present, the value of \texttt{{\textless}nr{\textunderscore}edges{\textgreater}} will be ignored. The correct number of edges will be deduced from the rest of
the information in the file. 
\item  A line of the form \texttt{e {\textless}v{\textgreater} {\textless}w{\textgreater}}, where \texttt{{\textless}v{\textgreater}} and \texttt{{\textless}w{\textgreater}} are integers in the range \texttt{[1 .. {\textless}nr{\textunderscore}vertices{\textgreater}]}, specifies that there is a (symmetric) edge in the digraph between the
vertices \texttt{{\textless}v{\textgreater}} and \texttt{{\textless}w{\textgreater}}. A symmetric edge only needs to be defined once; an additional line \texttt{e {\textless}v{\textgreater} {\textless}w{\textgreater}}, or \texttt{e {\textless}w{\textgreater} {\textless}v{\textgreater}}, will be interpreted as an additional, multiple, edge. Loops are permitted. 
\item  A line of the form \texttt{n {\textless}v{\textgreater} {\textless}label{\textgreater}}, where \texttt{{\textless}v{\textgreater}} is an integer in the range \texttt{[1 .. {\textless}nr{\textunderscore}vertices{\textgreater}]} and \texttt{{\textless}label{\textgreater}} is an integer, signifies that the vertex \texttt{{\textless}v{\textgreater}} has the label \texttt{{\textless}label{\textgreater}} in the digraph. If a label is not specified for a vertex, then \texttt{ReadDIMACSDigraph} will assign the label \texttt{1}, according to the DIMACS specification. 
\end{itemize}
 A detailed definition of the DIMACS format can be found in \ref{DIMACS} \ref{DIMACS}. 

 
\begin{Verbatim}[commandchars=!@|,fontsize=\small,frame=single,label=Example]
  !gapprompt@gap>| !gapinput@gr := Digraph([[2], [1, 3, 4], [2, 4], [2, 3]]);|
  <immutable digraph with 4 vertices, 8 edges>
  !gapprompt@gap>| !gapinput@filename := Concatenation(DIGRAPHS_Dir(),|
  !gapprompt@>| !gapinput@                             "/tst/out/dimacs.dimacs");;|
  !gapprompt@gap>| !gapinput@WriteDIMACSDigraph(filename, gr);;|
  !gapprompt@gap>| !gapinput@ReadDIMACSDigraph(filename);|
  <immutable digraph with 4 vertices, 8 edges>
\end{Verbatim}
 }

 

\subsection{\textcolor{Chapter }{WholeFileEncoders}}
\logpage{[ 9, 2, 17 ]}\nobreak
\hyperdef{L}{X7BC595AF841B78A9}{}
{\noindent\textcolor{FuncColor}{$\triangleright$\enspace\texttt{WholeFileEncoders\index{WholeFileEncoders@\texttt{WholeFileEncoders}}
\label{WholeFileEncoders}
}\hfill{\scriptsize (global variable)}}\\
\noindent\textcolor{FuncColor}{$\triangleright$\enspace\texttt{WholeFileDecoders\index{WholeFileDecoders@\texttt{WholeFileDecoders}}
\label{WholeFileDecoders}
}\hfill{\scriptsize (global variable)}}\\
\noindent\textcolor{FuncColor}{$\triangleright$\enspace\texttt{IsWholeFileEncoder({\mdseries\slshape encoder})\index{IsWholeFileEncoder@\texttt{IsWholeFileEncoder}}
\label{IsWholeFileEncoder}
}\hfill{\scriptsize (function)}}\\
\noindent\textcolor{FuncColor}{$\triangleright$\enspace\texttt{IsWholeFileDecoder({\mdseries\slshape decoder})\index{IsWholeFileDecoder@\texttt{IsWholeFileDecoder}}
\label{IsWholeFileDecoder}
}\hfill{\scriptsize (function)}}\\


 \texttt{WholeFileEncoders} and \texttt{WholeFileDecoders} are hashsets containing functions that encode and decode an entire file. These
are functions designed for graph formats where graphs are declared across
multiple lines, such as dreadnaut.

 \texttt{IsWholeFileEncoder} and \texttt{IsWholeFileDecoder} are functions that check whether a given argument belongs to the hashsets \texttt{WholeFileEncoders} and \texttt{WholeFileDecoders}, respectively. 

 }

 }

 }

   {\nobreakspace} 

\appendix


\chapter{\textcolor{Chapter }{ Grape to Digraphs Command Map }}\logpage{[ "A", 0, 0 ]}
\hyperdef{L}{X800DAF137BA8E25B}{}
{
  Below is a table of \href{https://gap-packages.github.io/grape} {GRAPE}  commands with the \textsf{Digraphs} counterparts. The sections in this chapter correspond to the chapters in the \href{https://gap-packages.github.io/grape} {GRAPE}  manual. 
\section{\textcolor{Chapter }{ Functions to construct and modify graphs }}\logpage{[ "A", 1, 0 ]}
\hyperdef{L}{X7CC8A9DB86F9F0F4}{}
{
   \emph{ The table in this section contains more information when viewed in html
format. }  \begin{center}
\begin{tabular}{|l|l|l|}\hline
 \href{https://gap-packages.github.io/grape} {GRAPE}  command &
 \textsf{Digraphs} command &
  \\
\hline
 \texttt{Graph} &
 \texttt{Digraph} (\ref{Digraph}) &
  \\
\hline
 \texttt{EdgeOrbitsGraph} &
 \texttt{EdgeOrbitsDigraph} (\ref{EdgeOrbitsDigraph}) &
  \\
\hline
 \texttt{NullGraph} &
 \texttt{NullDigraph} (\ref{NullDigraph})  &
  \\
\hline
 \texttt{CompleteGraph} &
 \texttt{CompleteDigraph} (\ref{CompleteDigraph}) &
  \\
\hline
 \texttt{JohnsonGraph} &
 \texttt{JohnsonDigraph} (\ref{JohnsonDigraph}) &
  \\
\hline
 \texttt{CayleyGraph} &
 \texttt{CayleyDigraph} (\ref{CayleyDigraph}) &
  \\
\hline
 \texttt{AddEdgeOrbit} &
 \texttt{DigraphAddEdgeOrbit} (\ref{DigraphAddEdgeOrbit}) &
  \\
\hline
 \texttt{RemoveEdgeOrbit} &
 \texttt{DigraphRemoveEdgeOrbit} (\ref{DigraphRemoveEdgeOrbit}) &
  \\
\hline
 \texttt{AssignVertexNames} &
 \texttt{SetDigraphVertexLabels} (\ref{SetDigraphVertexLabels})  &
  \\
\hline
\end{tabular}\\[2mm]
\end{center}

 }

 
\section{\textcolor{Chapter }{ Functions to inspect graphs, vertices and edges }}\logpage{[ "A", 2, 0 ]}
\hyperdef{L}{X7C38A6397D2770FB}{}
{
   \emph{ The table in this section contains more information when viewed in html
format. }  \begin{center}
\begin{tabular}{|l|l|l|}\hline
 \href{https://gap-packages.github.io/grape} {GRAPE}  command &
 \textsf{Digraphs} command &
  \\
\hline
 \texttt{IsGraph} &
 \texttt{IsDigraph} (\ref{IsDigraph}) &
  \\
\hline
 \texttt{OrderGraph} &
 \texttt{DigraphNrVertices} (\ref{DigraphNrVertices}) &
  \\
\hline
 \texttt{IsVertex(\mbox{\texttt{\mdseries\slshape graph}}, \mbox{\texttt{\mdseries\slshape v}})} &
 \texttt{\mbox{\texttt{\mdseries\slshape v}} in DigraphVertices(\mbox{\texttt{\mdseries\slshape digraph}})} &
  \\
\hline
 \texttt{VertexName} &
 \texttt{DigraphVertexLabel} (\ref{DigraphVertexLabel}) &
  \\
\hline
 \texttt{VertexNames} &
 \texttt{DigraphVertexLabels} (\ref{DigraphVertexLabels}) &
  \\
\hline
 \texttt{Vertices} &
 \texttt{DigraphVertices} (\ref{DigraphVertices}) &
  \\
\hline
 \texttt{VertexDegree} &
 \texttt{OutDegreeOfVertex} (\ref{OutDegreeOfVertex})  &
  \\
\hline
 \texttt{VertexDegrees} &
 \texttt{OutDegreeSet} (\ref{OutDegreeSet})  &
  \\
\hline
 \texttt{IsLoopy} &
 \texttt{DigraphHasLoops} (\ref{DigraphHasLoops}) &
  \\
\hline
 \texttt{IsSimpleGraph} &
   \texttt{IsSymmetricDigraph} (\ref{IsSymmetricDigraph})  &
  \\
\hline
 \texttt{Adjacency} &
 \texttt{OutNeighboursOfVertex} (\ref{OutNeighboursOfVertex})  &
  \\
\hline
 \texttt{IsEdge} &
 \texttt{IsDigraphEdge} (\ref{IsDigraphEdge:for digraph and list}) &
  \\
\hline
 \texttt{DirectedEdges} &
 \texttt{DigraphEdges} (\ref{DigraphEdges}) &
  \\
\hline
 \texttt{UndirectedEdges} &
 None &
  \\
\hline
 \texttt{Distance} &
 \texttt{DigraphShortestDistance} (\ref{DigraphShortestDistance:for a digraph and two vertices}) &
  \\
\hline
 \texttt{Diameter} &
 \texttt{DigraphDiameter} (\ref{DigraphDiameter}) &
  \\
\hline
 \texttt{Girth} &
 \texttt{DigraphUndirectedGirth} (\ref{DigraphUndirectedGirth}) &
  \\
\hline
 \texttt{IsConnectedGraph} &
 \texttt{IsStronglyConnectedDigraph} (\ref{IsStronglyConnectedDigraph}) &
  \\
\hline
 \texttt{IsBipartite} &
 \texttt{IsBipartiteDigraph} (\ref{IsBipartiteDigraph})  &
  \\
\hline
 \texttt{IsNullGraph} &
 \texttt{IsNullDigraph} (\ref{IsNullDigraph})  &
  \\
\hline
 \texttt{IsCompleteGraph} &
 \texttt{IsCompleteDigraph} (\ref{IsCompleteDigraph}) &
  \\
\hline
\end{tabular}\\[2mm]
\end{center}

 }

 
\section{\textcolor{Chapter }{ Functions to determine regularity properties of graphs }}\logpage{[ "A", 3, 0 ]}
\hyperdef{L}{X78A5A34783C70CD7}{}
{
   \emph{ The table in this section contains more information when viewed in html
format. }  \begin{center}
\begin{tabular}{|l|l|l|}\hline
 \href{https://gap-packages.github.io/grape} {GRAPE}  command &
 \textsf{Digraphs} command &
  \\
\hline
 \texttt{IsRegularGraph} &
 \texttt{IsOutRegularDigraph} (\ref{IsOutRegularDigraph})  &
  \\
\hline
 \texttt{LocalParameters} &
 None &
  \\
\hline
 \texttt{GlobalParameters} &
 None &
  \\
\hline
 \texttt{IsDistanceRegular} &
 \texttt{IsDistanceRegularDigraph} (\ref{IsDistanceRegularDigraph}) &
  \\
\hline
 \texttt{CollapsedAdjacencyMat} &
 None &
  \\
\hline
 \texttt{OrbitalGraphColadjMats} &
 None &
  \\
\hline
 \texttt{VertexTransitiveDRGs} &
 None &
  \\
\hline
\end{tabular}\\[2mm]
\end{center}

 }

 
\section{\textcolor{Chapter }{ Some special vertex subsets of a graph }}\logpage{[ "A", 4, 0 ]}
\hyperdef{L}{X7C74978D85F7AF5C}{}
{
   \emph{ The table in this section contains more information when viewed in html
format. }  \begin{center}
\begin{tabular}{|l|l|l|}\hline
 \href{https://gap-packages.github.io/grape} {GRAPE}  command &
 \textsf{Digraphs} command &
  \\
\hline
 \texttt{ConnectedComponent} &
 \texttt{DigraphConnectedComponent} (\ref{DigraphConnectedComponent}) &
  \\
\hline
 \texttt{ConnectedComponents} &
 \texttt{DigraphConnectedComponents} (\ref{DigraphConnectedComponents}) &
  \\
\hline
 \texttt{Bicomponents} &
 \texttt{DigraphBicomponents} (\ref{DigraphBicomponents}) &
  \\
\hline
 \texttt{DistanceSet} &
 \texttt{DigraphDistanceSet} (\ref{DigraphDistanceSet:for a digraph, a pos int, and an int}) &
  \\
\hline
 \texttt{Layers} &
 \texttt{DigraphLayers} (\ref{DigraphLayers}) &
  \\
\hline
 \texttt{IndependentSet} &
 \texttt{DigraphIndependentSet} (\ref{DigraphIndependentSet}) &
  \\
\hline
\end{tabular}\\[2mm]
\end{center}

 }

 
\section{\textcolor{Chapter }{ Functions to construct new graphs from old }}\logpage{[ "A", 5, 0 ]}
\hyperdef{L}{X7E51B92F7E89E266}{}
{
   \emph{ The table in this section contains more information when viewed in html
format. }  \begin{center}
\begin{tabular}{|l|l|l|}\hline
 \href{https://gap-packages.github.io/grape} {GRAPE}  command &
 \textsf{Digraphs} command &
  \\
\hline
 \texttt{InducedSubgraph} &
 \texttt{InducedSubdigraph} (\ref{InducedSubdigraph}) &
  \\
\hline
 \texttt{DistanceSetInduced} &
 None &
  \\
\hline
 \texttt{DistanceGraph} &
 \texttt{DistanceDigraph} (\ref{DistanceDigraph:for digraph and list}) &
  \\
\hline
 \texttt{ComplementGraph} &
 \texttt{DigraphDual} (\ref{DigraphDual}) &
  \\
\hline
 \texttt{PointGraph} &
 None &
  \\
\hline
 \texttt{EdgeGraph} &
 \texttt{EdgeUndirectedDigraph} (\ref{EdgeUndirectedDigraph}) &
  \\
\hline
 \texttt{SwitchedGraph} &
 None &
  \\
\hline
 \texttt{UnderlyingGraph} &
 \texttt{DigraphSymmetricClosure} (\ref{DigraphSymmetricClosure}) &
  \\
\hline
 \texttt{QuotientGraph} &
 \texttt{QuotientDigraph} (\ref{QuotientDigraph}) &
  \\
\hline
 \texttt{BipartiteDouble} &
 \texttt{BipartiteDoubleDigraph} (\ref{BipartiteDoubleDigraph}) &
  \\
\hline
 \texttt{GeodesicsGraph} &
 None &
  \\
\hline
 \texttt{CollapsedIndependentOrbitsGraph} &
 None &
  \\
\hline
 \texttt{CollapsedCompleteOrbitsGraph} &
 None &
  \\
\hline
 \texttt{NewGroupGraph} &
 None &
  \\
\hline
\end{tabular}\\[2mm]
\end{center}

 }

 
\section{\textcolor{Chapter }{ Vertex\texttt{\symbol{45}}Colouring and Complete Subgraphs }}\logpage{[ "A", 6, 0 ]}
\hyperdef{L}{X7A61E3FC7A2BB50E}{}
{
   \emph{ The table in this section contains more information when viewed in html
format. }  \begin{center}
\begin{tabular}{|l|l|l|}\hline
 \href{https://gap-packages.github.io/grape} {GRAPE}  command &
 \textsf{Digraphs} command &
  \\
\hline
 \texttt{VertexColouring} &
 \texttt{DigraphGreedyColouring} (\ref{DigraphGreedyColouring:for a digraph}) &
  \\
\hline
 \texttt{CompleteSubgraphs} &
 \texttt{DigraphCliques} (\ref{DigraphCliques}) &
  \\
\hline
 \texttt{CompleteSubgraphsOfGivenSize} &
 \texttt{DigraphCliques} (\ref{DigraphCliques}) &
  \\
\hline
\end{tabular}\\[2mm]
\end{center}

 }

 
\section{\textcolor{Chapter }{ Automorphism groups and isomorphism testing for graphs }}\logpage{[ "A", 7, 0 ]}
\hyperdef{L}{X7DFC61EC7CBDBC7E}{}
{
   \emph{ The table in this section contains more information when viewed in html
format. }  \begin{center}
\begin{tabular}{|l|l|l|}\hline
 \href{https://gap-packages.github.io/grape} {GRAPE}  command &
 \textsf{Digraphs} command &
  \\
\hline
 \texttt{AutGroupGraph} &
 \texttt{AutomorphismGroup} (\ref{AutomorphismGroup:for a digraph}) &
  \\
\hline
 \texttt{GraphIsomorphism} &
 \texttt{IsomorphismDigraphs} (\ref{IsomorphismDigraphs:for digraphs}) &
  \\
\hline
 \texttt{IsIsomorphicGraph} &
 \texttt{IsIsomorphicDigraph} (\ref{IsIsomorphicDigraph:for digraphs}) &
  \\
\hline
 \texttt{GraphIsomorphismClassRepresentatives} &
 None &
  \\
\hline
\end{tabular}\\[2mm]
\end{center}

 }

 }


\chapter{\textcolor{Chapter }{ DIMACS: Graph Format for Clique and Coloring Problems }}\label{DIMACS}
\logpage{[ "B", 0, 0 ]}
\hyperdef{L}{X819C7A688246B572}{}
{
  
\section{\textcolor{Chapter }{ Note from the Digraphs authors }}\logpage{[ "B", 1, 0 ]}
\hyperdef{L}{X7D4AB14683444D77}{}
{
  The contents of this appendix were originally available in a PostScript file \href{http://mat.gsia.cmu.edu/COLOR/general/ccformat.ps} {on the Carnegie Mellon University website}. This file is still \href{https://web.archive.org/web/20220120085844/http://mat.gsia.cmu.edu/COLOR/general/ccformat.ps} {accessible through the Wayback Machine}. We reproduce its contents here for convenience, without adjustments beyond
minor re\texttt{\symbol{45}}formatting. }

 
\section{\textcolor{Chapter }{ Preamble }}\logpage{[ "B", 2, 0 ]}
\hyperdef{L}{X7E738C697CA72B1F}{}
{
  \textsc{Last revision of this document: May 08, 1993.} 

 \emph{ This paper outlines a suggested graph format. If you have comments on this or
other formats or you have information you think should be included, please
send a note to } \href{mailto://challenge@dimacs.rutgers.edu} {\texttt{challenge@dimacs.rutgers.edu}}. }

 
\section{\textcolor{Chapter }{ Introduction }}\logpage{[ "B", 3, 0 ]}
\hyperdef{L}{X7DFB63A97E67C0A1}{}
{
  One purpose of the DIMACS Challenge is to ease the effort required to test and
compare algorithms and heuristics by providing a common testbed of instances
and analysis tools. To facilitate this effort, a standard format must be
chosen for the problems addressed. This document outlines a format for graphs
that is suitable for those looking at graph coloring and finding cliques in
graphs. This format is a flexible format suitable for many types of graph and
network problems. This format was also the format chosen for the First
Computational Challenge on network flows and matchings. 

 This document describes three problems: unweighted clique, weighted clique,
and graph coloring. A separate format is used for satisfiability. }

 
\section{\textcolor{Chapter }{ File Formats for Graph Problems }}\logpage{[ "B", 4, 0 ]}
\hyperdef{L}{X80AD820D85A4C8DA}{}
{
  This section describes a standard file format for graph inputs and outputs.
There is no requirement that participants follow these specifications;
however, compatible implementations will be able to make full use of DIMACS
support tools. (Some tools assume that output is appended to input in a single
file.) Participants are welcome to develop translation programs to convert
instances to and from more convenient, or more compact, representations; the
Unix \textsc{awk} facility is recommended as especially suitable for this task. All files
contain ASCII characters. Input and output files contain several types of \emph{lines}, described below. A line is terminated with an
end\texttt{\symbol{45}}of\texttt{\symbol{45}}line character. Fields in each
line are separated by at least one blank space. Each line begins with a
one\texttt{\symbol{45}}character designator to identify the line type. 
\subsection{\textcolor{Chapter }{ Input Files }}\logpage{[ "B", 4, 1 ]}
\hyperdef{L}{X792906427BCCD4D1}{}
{
  An input file contains all the information about a graph needed to define
either a clique problem or a coloring problem. Some information may be
included that is not relevant to one problem (for instance, node weights are
not needed for coloring problem) so that information may be ignored. In this
format, nodes are numbered from 1 up to $n$. There are $m$ edges in the graph. Files are assumed to be well\texttt{\symbol{45}}formed and
internally consistent: node identifier values are valid, nodes are defined
uniquely, exactly $m$ edges are defined, and so forth. A input checker will be made available to
ensure compatibility with this standard. 
\begin{description}
\item[{Comments}]  Comment lines give human\texttt{\symbol{45}}readable information about the
file and are ignored by programs. Comment lines can appear anywhere in the
file. Each comment line begins with a lowercase character \textsc{c}. 
\begin{Verbatim}[commandchars=!@|,fontsize=\small,frame=single,label=Example]
  c This is an example of a comment line.
\end{Verbatim}
 
\item[{ Problem line }]  There is one problem line per input file. The problem line must appear before
any node or arc descriptor lines. For network instances, the problem line has
the following format. 
\begin{Verbatim}[commandchars=!@|,fontsize=\small,frame=single,label=Example]
  p FORMAT NODES EDGES
\end{Verbatim}
 The lowercase character \texttt{p} signifies that this is the problem line. The \texttt{FORMAT} field is for consistency with the previous Challenge, and should contain the
word $edge$. The \texttt{NODES} field contains an integer value specifying $n$, the number of nodes in the graph. The \texttt{EDGES} field contains an integer value specifying $m$, the number of edges in the graph. 
\item[{Node Descriptors}]  For this Challenge, a node descriptor is required only for the weighted clique
problem. These lines will give the weight assigned to a node in the clique.
There is one node descriptor line for each node, with the following format.
Nodes without a descriptor will take on a default value of 1. 
\begin{Verbatim}[commandchars=!@|,fontsize=\small,frame=single,label=Example]
  n ID VALUE
\end{Verbatim}
 The lowercase character \texttt{n} signifies that this is a node descriptor line. The \texttt{ID} field gives a node identification number, an integer between 1 and $n$. The \texttt{VALUE} gives the objective value for having this node in the clique. This value is
assumed to be integer and can be either positive or negative (or zero). 
\item[{Edge Descriptors}]  There is one edge descriptor line for each edge the graph, each with the
following format. Each edge $(v,w)$ appears exactly once in the input file and is not repeated as $(w,v)$. 
\begin{Verbatim}[commandchars=!@|,fontsize=\small,frame=single,label=Example]
  e W  V
\end{Verbatim}
 The lowercase character \texttt{e} signifies that this is an edge descriptor line. For an edge $(w,v)$ the fields \texttt{W} and \texttt{V} specify its endpoints. 
\item[{Optional Descriptors}]  In addition to the required information, there can be additional pieces of
information about a graph. This will typically define the parameters used to
generate the graph or otherwise define generator\texttt{\symbol{45}}specific
information. The following list may be added to as interesting problem
generators are decided on: 
\begin{description}
\item[{Geometric Descriptors}]  One common method to generate or display graphs is to have the nodes be
embedded in some space and to have the edges be included according to some
function of the distance between nodes according to some metric. The node
information can be defined by a dimension descriptor and a vertex embedding
descriptor. 
\begin{Verbatim}[commandchars=!@|,fontsize=\small,frame=single,label=Example]
  d DIM METRIC
\end{Verbatim}
 is the dimension descriptor. \texttt{DIM} is an integer giving the number of dimensions of the space, while \texttt{METRIC} is a string representing the metric for the space. \texttt{METRIC} is a string that can take a number of forms. \textsc{Lp} (i.e. \textsc{L1}, \textsc{L2}, \textsc{L122}, and so on) denotes the $\ell_p$ norm where the distance between two nodes embedded at $(x_1,x_2,\ldots,x_d)$ and $(y_1,y_2,\ldots y_d)$ is $\left(\sum_{i=1}^d |x_i-y_i|^p \right)^{1/p}$. The string \textsc{LINF} is used to denote the $\ell_\infty$ norm. L2S denotes the squared euclidean norm (which can be less susceptible to
computer differences in round\texttt{\symbol{45}}off and accuracy issues). 
\begin{Verbatim}[commandchars=!@|,fontsize=\small,frame=single,label=Example]
  v  X1  X2  X3  ... XD
\end{Verbatim}
 The lowercase character \texttt{v} signifies that this is a vertex embedding descriptor line. The fields \texttt{X1, X2, ..., XD} give the \texttt{d} coordinate values for the vertex. Note that these lines must appear after the \texttt{d} descriptor. 
\item[{Parameter Descriptors}]  The parameter descriptors are used to give other information about how the
graph was generated. The lines are generator\texttt{\symbol{45}}specific, and
as such it is not expected that most codes will use most (or any) of them.
They are included only to aid those codes specifically designed to attack
specially structured problems. The general form of the parameter descriptor
is: 
\begin{Verbatim}[commandchars=!@|,fontsize=\small,frame=single,label=Example]
  x PARAM VALUE
\end{Verbatim}
 The lowercase character \texttt{x} signifies that this is a parameter descriptor line. The \texttt{PARAM} field is a string that gives the name of the parameter, while the \texttt{VALUE} field is a numeric value that gives the corresponding value. The following \texttt{PARAM} values have been defined: \begin{center}
\begin{tabular}{|l|l|}\hline
\texttt{PARAM}&
Description (Geometric Graphs)\\
\hline
\hline
\texttt{MINLENGTH}&
 Edge included only if length greater than or equal to \texttt{VALUE} \\
\hline
\texttt{MAXLENGTH}&
 Edge included only if length less than or equal to \texttt{VALUE} \\
\hline
\end{tabular}\\[2mm]
\end{center}

 Note that this information is in addition to the required edge descriptors. 
\end{description}
 
\end{description}
 }

 
\subsection{\textcolor{Chapter }{ Output Files }}\logpage{[ "B", 4, 2 ]}
\hyperdef{L}{X7B0FD9A97D715330}{}
{
  Every algorithm or heuristic should create an output file. This output file
should consist of one or more of the following lines, depending on the type of
algorithm and problem being solved. 
\begin{description}
\item[{Solution Line}]  Format: 
\begin{Verbatim}[commandchars=!@|,fontsize=\small,frame=single,label=Example]
  s TYPE SOLUTION
\end{Verbatim}
 The lowercase character \texttt{s} signifies that this is a solution line. The \texttt{TYPE} field denotes the type of solution contained in the file. This should be one
of the following strings: $col$ denotes a graph coloring, $clq$ denotes a maximum weighted clique, and $cqu$ denotes a maximum unweighted clique (one that has ignored the \texttt{n} descriptor lines). The \texttt{SOLUTION} field contains an integer corresponding to the solution value. This is the
clique size for unweighted clique, clique value for weighted clique, or number
of colors used for graph coloring. 
\item[{Bound Line}]  Format: 
\begin{Verbatim}[commandchars=!@|,fontsize=\small,frame=single,label=Example]
  b BOUND
\end{Verbatim}
 The lowercase character \texttt{b} signifies that this is a bound on the the solution. The \texttt{BOUND} field contains an integer value that gives a bound on the solution value. This
bound is an upper bound on the maximum clique value for cliques and weighted
clique and a lower bound on the number of colors needed for coloring the
graph. 
\item[{Clique Line}]  Format: 
\begin{Verbatim}[commandchars=!@|,fontsize=\small,frame=single,label=Example]
  v V
\end{Verbatim}
 The lowercase character \texttt{v} signifies that this is a clique vertex line. The \texttt{V} field gives the node number for the node in the clique. There will be one
clique line for each node in the clique. 
\item[{Label Line}]  Format: 
\begin{Verbatim}[commandchars=!@|,fontsize=\small,frame=single,label=Example]
  l V N
\end{Verbatim}
 The lowercase character \texttt{l} signifies that this is a label line, generally used for graph coloring. The \texttt{V} field gives the node number for the node in the clique while the \texttt{N} field gives the corresponding label. There will be one label line for each
node in the graph. 
\end{description}
 }

 }

 }

\def\bibname{References\logpage{[ "Bib", 0, 0 ]}
\hyperdef{L}{X7A6F98FD85F02BFE}{}
}

\bibliographystyle{alpha}
\bibliography{digraphs.bib}

\addcontentsline{toc}{chapter}{References}

\def\indexname{Index\logpage{[ "Ind", 0, 0 ]}
\hyperdef{L}{X83A0356F839C696F}{}
}

\cleardoublepage
\phantomsection
\addcontentsline{toc}{chapter}{Index}


\printindex

\immediate\write\pagenrlog{["Ind", 0, 0], \arabic{page},}
\newpage
\immediate\write\pagenrlog{["End"], \arabic{page}];}
\immediate\closeout\pagenrlog
\end{document}
